% Options for packages loaded elsewhere
% Options for packages loaded elsewhere
\PassOptionsToPackage{unicode}{hyperref}
\PassOptionsToPackage{hyphens}{url}
\PassOptionsToPackage{dvipsnames,svgnames,x11names}{xcolor}
%
\documentclass[
  11pt,
  letterpaper,
  DIV=11,
  numbers=noendperiod]{scrreprt}
\usepackage{xcolor}
\usepackage[margin=1in]{geometry}
\usepackage{amsmath,amssymb}
\setcounter{secnumdepth}{5}
\usepackage{iftex}
\ifPDFTeX
  \usepackage[T1]{fontenc}
  \usepackage[utf8]{inputenc}
  \usepackage{textcomp} % provide euro and other symbols
\else % if luatex or xetex
  \usepackage{unicode-math} % this also loads fontspec
  \defaultfontfeatures{Scale=MatchLowercase}
  \defaultfontfeatures[\rmfamily]{Ligatures=TeX,Scale=1}
\fi
\usepackage{lmodern}
\ifPDFTeX\else
  % xetex/luatex font selection
\fi
% Use upquote if available, for straight quotes in verbatim environments
\IfFileExists{upquote.sty}{\usepackage{upquote}}{}
\IfFileExists{microtype.sty}{% use microtype if available
  \usepackage[]{microtype}
  \UseMicrotypeSet[protrusion]{basicmath} % disable protrusion for tt fonts
}{}
\makeatletter
\@ifundefined{KOMAClassName}{% if non-KOMA class
  \IfFileExists{parskip.sty}{%
    \usepackage{parskip}
  }{% else
    \setlength{\parindent}{0pt}
    \setlength{\parskip}{6pt plus 2pt minus 1pt}}
}{% if KOMA class
  \KOMAoptions{parskip=half}}
\makeatother
% Make \paragraph and \subparagraph free-standing
\makeatletter
\ifx\paragraph\undefined\else
  \let\oldparagraph\paragraph
  \renewcommand{\paragraph}{
    \@ifstar
      \xxxParagraphStar
      \xxxParagraphNoStar
  }
  \newcommand{\xxxParagraphStar}[1]{\oldparagraph*{#1}\mbox{}}
  \newcommand{\xxxParagraphNoStar}[1]{\oldparagraph{#1}\mbox{}}
\fi
\ifx\subparagraph\undefined\else
  \let\oldsubparagraph\subparagraph
  \renewcommand{\subparagraph}{
    \@ifstar
      \xxxSubParagraphStar
      \xxxSubParagraphNoStar
  }
  \newcommand{\xxxSubParagraphStar}[1]{\oldsubparagraph*{#1}\mbox{}}
  \newcommand{\xxxSubParagraphNoStar}[1]{\oldsubparagraph{#1}\mbox{}}
\fi
\makeatother

\usepackage{color}
\usepackage{fancyvrb}
\newcommand{\VerbBar}{|}
\newcommand{\VERB}{\Verb[commandchars=\\\{\}]}
\DefineVerbatimEnvironment{Highlighting}{Verbatim}{commandchars=\\\{\}}
% Add ',fontsize=\small' for more characters per line
\usepackage{framed}
\definecolor{shadecolor}{RGB}{241,243,245}
\newenvironment{Shaded}{\begin{snugshade}}{\end{snugshade}}
\newcommand{\AlertTok}[1]{\textcolor[rgb]{0.68,0.00,0.00}{#1}}
\newcommand{\AnnotationTok}[1]{\textcolor[rgb]{0.37,0.37,0.37}{#1}}
\newcommand{\AttributeTok}[1]{\textcolor[rgb]{0.40,0.45,0.13}{#1}}
\newcommand{\BaseNTok}[1]{\textcolor[rgb]{0.68,0.00,0.00}{#1}}
\newcommand{\BuiltInTok}[1]{\textcolor[rgb]{0.00,0.23,0.31}{#1}}
\newcommand{\CharTok}[1]{\textcolor[rgb]{0.13,0.47,0.30}{#1}}
\newcommand{\CommentTok}[1]{\textcolor[rgb]{0.37,0.37,0.37}{#1}}
\newcommand{\CommentVarTok}[1]{\textcolor[rgb]{0.37,0.37,0.37}{\textit{#1}}}
\newcommand{\ConstantTok}[1]{\textcolor[rgb]{0.56,0.35,0.01}{#1}}
\newcommand{\ControlFlowTok}[1]{\textcolor[rgb]{0.00,0.23,0.31}{\textbf{#1}}}
\newcommand{\DataTypeTok}[1]{\textcolor[rgb]{0.68,0.00,0.00}{#1}}
\newcommand{\DecValTok}[1]{\textcolor[rgb]{0.68,0.00,0.00}{#1}}
\newcommand{\DocumentationTok}[1]{\textcolor[rgb]{0.37,0.37,0.37}{\textit{#1}}}
\newcommand{\ErrorTok}[1]{\textcolor[rgb]{0.68,0.00,0.00}{#1}}
\newcommand{\ExtensionTok}[1]{\textcolor[rgb]{0.00,0.23,0.31}{#1}}
\newcommand{\FloatTok}[1]{\textcolor[rgb]{0.68,0.00,0.00}{#1}}
\newcommand{\FunctionTok}[1]{\textcolor[rgb]{0.28,0.35,0.67}{#1}}
\newcommand{\ImportTok}[1]{\textcolor[rgb]{0.00,0.46,0.62}{#1}}
\newcommand{\InformationTok}[1]{\textcolor[rgb]{0.37,0.37,0.37}{#1}}
\newcommand{\KeywordTok}[1]{\textcolor[rgb]{0.00,0.23,0.31}{\textbf{#1}}}
\newcommand{\NormalTok}[1]{\textcolor[rgb]{0.00,0.23,0.31}{#1}}
\newcommand{\OperatorTok}[1]{\textcolor[rgb]{0.37,0.37,0.37}{#1}}
\newcommand{\OtherTok}[1]{\textcolor[rgb]{0.00,0.23,0.31}{#1}}
\newcommand{\PreprocessorTok}[1]{\textcolor[rgb]{0.68,0.00,0.00}{#1}}
\newcommand{\RegionMarkerTok}[1]{\textcolor[rgb]{0.00,0.23,0.31}{#1}}
\newcommand{\SpecialCharTok}[1]{\textcolor[rgb]{0.37,0.37,0.37}{#1}}
\newcommand{\SpecialStringTok}[1]{\textcolor[rgb]{0.13,0.47,0.30}{#1}}
\newcommand{\StringTok}[1]{\textcolor[rgb]{0.13,0.47,0.30}{#1}}
\newcommand{\VariableTok}[1]{\textcolor[rgb]{0.07,0.07,0.07}{#1}}
\newcommand{\VerbatimStringTok}[1]{\textcolor[rgb]{0.13,0.47,0.30}{#1}}
\newcommand{\WarningTok}[1]{\textcolor[rgb]{0.37,0.37,0.37}{\textit{#1}}}

\usepackage{longtable,booktabs,array}
\newcounter{none} % for unnumbered tables
\usepackage{calc} % for calculating minipage widths
% Correct order of tables after \paragraph or \subparagraph
\usepackage{etoolbox}
\makeatletter
\patchcmd\longtable{\par}{\if@noskipsec\mbox{}\fi\par}{}{}
\makeatother
% Allow footnotes in longtable head/foot
\IfFileExists{footnotehyper.sty}{\usepackage{footnotehyper}}{\usepackage{footnote}}
\makesavenoteenv{longtable}
\usepackage{graphicx}
\makeatletter
\newsavebox\pandoc@box
\newcommand*\pandocbounded[1]{% scales image to fit in text height/width
  \sbox\pandoc@box{#1}%
  \Gscale@div\@tempa{\textheight}{\dimexpr\ht\pandoc@box+\dp\pandoc@box\relax}%
  \Gscale@div\@tempb{\linewidth}{\wd\pandoc@box}%
  \ifdim\@tempb\p@<\@tempa\p@\let\@tempa\@tempb\fi% select the smaller of both
  \ifdim\@tempa\p@<\p@\scalebox{\@tempa}{\usebox\pandoc@box}%
  \else\usebox{\pandoc@box}%
  \fi%
}
% Set default figure placement to htbp
\def\fps@figure{htbp}
\makeatother


% definitions for citeproc citations
\NewDocumentCommand\citeproctext{}{}
\NewDocumentCommand\citeproc{mm}{%
  \begingroup\def\citeproctext{#2}\cite{#1}\endgroup}
\makeatletter
 % allow citations to break across lines
 \let\@cite@ofmt\@firstofone
 % avoid brackets around text for \cite:
 \def\@biblabel#1{}
 \def\@cite#1#2{{#1\if@tempswa , #2\fi}}
\makeatother
\newlength{\cslhangindent}
\setlength{\cslhangindent}{1.5em}
\newlength{\csllabelwidth}
\setlength{\csllabelwidth}{3em}
\newenvironment{CSLReferences}[2] % #1 hanging-indent, #2 entry-spacing
 {\begin{list}{}{%
  \setlength{\itemindent}{0pt}
  \setlength{\leftmargin}{0pt}
  \setlength{\parsep}{0pt}
  % turn on hanging indent if param 1 is 1
  \ifodd #1
   \setlength{\leftmargin}{\cslhangindent}
   \setlength{\itemindent}{-1\cslhangindent}
  \fi
  % set entry spacing
  \setlength{\itemsep}{#2\baselineskip}}}
 {\end{list}}
\usepackage{calc}
\newcommand{\CSLBlock}[1]{\hfill\break\parbox[t]{\linewidth}{\strut\ignorespaces#1\strut}}
\newcommand{\CSLLeftMargin}[1]{\parbox[t]{\csllabelwidth}{\strut#1\strut}}
\newcommand{\CSLRightInline}[1]{\parbox[t]{\linewidth - \csllabelwidth}{\strut#1\strut}}
\newcommand{\CSLIndent}[1]{\hspace{\cslhangindent}#1}



\setlength{\emergencystretch}{3em} % prevent overfull lines

\providecommand{\tightlist}{%
  \setlength{\itemsep}{0pt}\setlength{\parskip}{0pt}}



 


\KOMAoption{captions}{tableheading}
\makeatletter
\@ifpackageloaded{tcolorbox}{}{\usepackage[skins,breakable]{tcolorbox}}
\@ifpackageloaded{fontawesome5}{}{\usepackage{fontawesome5}}
\definecolor{quarto-callout-color}{HTML}{909090}
\definecolor{quarto-callout-note-color}{HTML}{0758E5}
\definecolor{quarto-callout-important-color}{HTML}{CC1914}
\definecolor{quarto-callout-warning-color}{HTML}{EB9113}
\definecolor{quarto-callout-tip-color}{HTML}{00A047}
\definecolor{quarto-callout-caution-color}{HTML}{FC5300}
\definecolor{quarto-callout-color-frame}{HTML}{acacac}
\definecolor{quarto-callout-note-color-frame}{HTML}{4582ec}
\definecolor{quarto-callout-important-color-frame}{HTML}{d9534f}
\definecolor{quarto-callout-warning-color-frame}{HTML}{f0ad4e}
\definecolor{quarto-callout-tip-color-frame}{HTML}{02b875}
\definecolor{quarto-callout-caution-color-frame}{HTML}{fd7e14}
\makeatother
\makeatletter
\@ifpackageloaded{bookmark}{}{\usepackage{bookmark}}
\makeatother
\makeatletter
\@ifpackageloaded{caption}{}{\usepackage{caption}}
\AtBeginDocument{%
\ifdefined\contentsname
  \renewcommand*\contentsname{Table of contents}
\else
  \newcommand\contentsname{Table of contents}
\fi
\ifdefined\listfigurename
  \renewcommand*\listfigurename{List of Figures}
\else
  \newcommand\listfigurename{List of Figures}
\fi
\ifdefined\listtablename
  \renewcommand*\listtablename{List of Tables}
\else
  \newcommand\listtablename{List of Tables}
\fi
\ifdefined\figurename
  \renewcommand*\figurename{Figure}
\else
  \newcommand\figurename{Figure}
\fi
\ifdefined\tablename
  \renewcommand*\tablename{Table}
\else
  \newcommand\tablename{Table}
\fi
}
\@ifpackageloaded{float}{}{\usepackage{float}}
\floatstyle{ruled}
\@ifundefined{c@chapter}{\newfloat{codelisting}{h}{lop}}{\newfloat{codelisting}{h}{lop}[chapter]}
\floatname{codelisting}{Listing}
\newcommand*\listoflistings{\listof{codelisting}{List of Listings}}
\makeatother
\makeatletter
\makeatother
\makeatletter
\@ifpackageloaded{caption}{}{\usepackage{caption}}
\@ifpackageloaded{subcaption}{}{\usepackage{subcaption}}
\makeatother
\usepackage{bookmark}
\IfFileExists{xurl.sty}{\usepackage{xurl}}{} % add URL line breaks if available
\urlstyle{same}
\hypersetup{
  pdftitle={Mathematical Foundations of Natural Language Processing},
  pdfauthor={He Wang},
  colorlinks=true,
  linkcolor={blue},
  filecolor={Maroon},
  citecolor={Blue},
  urlcolor={Blue},
  pdfcreator={LaTeX via pandoc}}


\title{Mathematical Foundations of Natural Language Processing}
\usepackage{etoolbox}
\makeatletter
\providecommand{\subtitle}[1]{% add subtitle to \maketitle
  \apptocmd{\@title}{\par {\large #1 \par}}{}{}
}
\makeatother
\subtitle{A Clean Mathematical Introduction for Applied Mathematics and
Statistics Students}
\author{He Wang}
\date{2026-06-21}
\begin{document}
\maketitle

\renewcommand*\contentsname{Table of contents}
{
\setcounter{tocdepth}{2}
\tableofcontents
}

\bookmarksetup{startatroot}

\chapter{Mathematical Foundations of Natural Language
Processing}\label{mathematical-foundations-of-natural-language-processing}

From Tokens and Probability to Transformers, Retrieval, and Responsible
AI

\hfill\break

Natural language processing is the study of how language can be
represented, modeled, learned, retrieved, generated, and evaluated by
computational systems. This book develops NLP from a mathematical point
of view, with enough Python computation to make the ideas concrete and
testable.

This book is designed for students in applied mathematics, statistics,
data science, computer science, and related quantitative programs who
want to understand modern NLP beyond package usage. The central theme is
that language data can be studied through familiar mathematical
structures: finite sets, vectors, matrices, probability distributions,
optimization problems, graphs, geometry, and topology.

\section{What makes this book
mathematical?}\label{what-makes-this-book-mathematical}

Rather than beginning with software libraries, the book begins with
objects and operations.

{\def\LTcaptype{none} % do not increment counter
\begin{longtable}[]{@{}
  >{\raggedright\arraybackslash}p{(\linewidth - 4\tabcolsep) * \real{0.3333}}
  >{\raggedright\arraybackslash}p{(\linewidth - 4\tabcolsep) * \real{0.3333}}
  >{\raggedright\arraybackslash}p{(\linewidth - 4\tabcolsep) * \real{0.3333}}@{}}
\toprule\noalign{}
\begin{minipage}[b]{\linewidth}\raggedright
NLP idea
\end{minipage} & \begin{minipage}[b]{\linewidth}\raggedright
Mathematical viewpoint
\end{minipage} & \begin{minipage}[b]{\linewidth}\raggedright
Computational object
\end{minipage} \\
\midrule\noalign{}
\endhead
\bottomrule\noalign{}
\endlastfoot
Tokens & Elements of a finite alphabet & strings, integer ids \\
Documents & Empirical measures or vectors & token lists, count
vectors \\
Vocabulary & Finite set with indexing & dictionaries, arrays \\
Language model & Conditional probability distribution & tables or neural
networks \\
Embedding & Map into a vector space & dense matrix rows \\
Similarity & Inner product, norm, or metric & cosine score, distance
matrix \\
Attention & Similarity-weighted averaging operator & query-key-value
matrices \\
Transformer & Nonlinear operator on sequences & stacked attention and
feedforward blocks \\
Retrieval & Nearest-neighbor search in representation space & ranked
document lists \\
Evaluation & Statistical estimation under uncertainty & metrics,
confidence intervals \\
Responsible NLP & Error analysis across populations and contexts &
audits, calibration, bias diagnostics \\
\end{longtable}
}

The goal is not only to run an NLP model. The goal is to understand what
the model is doing, what assumptions it makes, when it works, and how to
diagnose its failures.

\section{Intended audience}\label{intended-audience}

The book assumes basic familiarity with:

\begin{itemize}
\tightlist
\item
  linear algebra: vectors, matrices, inner products, projections,
  eigenvalues, and singular value decomposition;
\item
  probability: conditional probability, likelihood, entropy,
  expectation, and random variables;
\item
  statistics: classification, estimation, empirical risk, validation,
  and uncertainty;
\item
  Python: NumPy, pandas, matplotlib, and basic machine learning
  workflows.
\end{itemize}

No prior NLP background is required. The computer labs are written for
independent study and include the necessary background before
programming begins.

\section{How to use this book}\label{how-to-use-this-book}

Each chapter is organized around four layers.

\begin{enumerate}
\def\labelenumi{\arabic{enumi}.}
\tightlist
\item
  \textbf{Conceptual layer:} What is the NLP problem and why does it
  matter?
\item
  \textbf{Mathematical layer:} What objects, maps, distributions, or
  optimization problems define the method?
\item
  \textbf{Computational layer:} How can the idea be implemented in
  Python on a small transparent example?
\item
  \textbf{Critical layer:} What can go wrong, and how should we evaluate
  or audit the method?
\end{enumerate}

A recommended weekly workflow is:

\begin{enumerate}
\def\labelenumi{\arabic{enumi}.}
\tightlist
\item
  read the chapter overview and mathematical explanation;
\item
  work through the examples by hand for a small toy corpus;
\item
  open the corresponding lab in Google Colab;
\item
  complete the guided exercises;
\item
  write a short reflection on what the model assumes and where it may
  fail.
\end{enumerate}

\section{Book structure}\label{book-structure}

\subsection{Part I. Language as Data}\label{part-i.-language-as-data}

Part I introduces the basic data structures of NLP: texts, tokens,
vocabularies, documents, and document vectors.

{\def\LTcaptype{none} % do not increment counter
\begin{longtable}[]{@{}
  >{\raggedleft\arraybackslash}p{(\linewidth - 4\tabcolsep) * \real{0.4000}}
  >{\raggedright\arraybackslash}p{(\linewidth - 4\tabcolsep) * \real{0.3000}}
  >{\raggedright\arraybackslash}p{(\linewidth - 4\tabcolsep) * \real{0.3000}}@{}}
\toprule\noalign{}
\begin{minipage}[b]{\linewidth}\raggedleft
Chapter
\end{minipage} & \begin{minipage}[b]{\linewidth}\raggedright
Topic
\end{minipage} & \begin{minipage}[b]{\linewidth}\raggedright
Lab
\end{minipage} \\
\midrule\noalign{}
\endhead
\bottomrule\noalign{}
\endlastfoot
1 & \href{chapters/chapter-01-what-is-nlp.qmd}{What Is Natural Language
Processing?} &
\href{https://colab.research.google.com/github/wanghemath/Book-MathNLP/blob/main/labs/chapter-01-what-is-nlp.ipynb}{Open
Lab 1 in Colab} \\
2 & \href{chapters/chapter-02-text-tokens-vocabularies.qmd}{Text,
Tokens, and Vocabularies} &
\href{https://colab.research.google.com/github/wanghemath/Book-MathNLP/blob/main/labs/chapter-02-text-tokens-vocabularies.ipynb}{Open
Lab 2 in Colab} \\
3 & \href{chapters/chapter-03-documents-as-vectors.qmd}{Documents as
Vectors} &
\href{https://colab.research.google.com/github/wanghemath/Book-MathNLP/blob/main/labs/chapter-03-documents-as-vectors.ipynb}{Open
Lab 3 in Colab} \\
\end{longtable}
}

\subsection{Part II. Probability Models for
Language}\label{part-ii.-probability-models-for-language}

Part II develops probabilistic language modeling, smoothing, likelihood,
classification, and uncertainty.

{\def\LTcaptype{none} % do not increment counter
\begin{longtable}[]{@{}
  >{\raggedleft\arraybackslash}p{(\linewidth - 4\tabcolsep) * \real{0.4000}}
  >{\raggedright\arraybackslash}p{(\linewidth - 4\tabcolsep) * \real{0.3000}}
  >{\raggedright\arraybackslash}p{(\linewidth - 4\tabcolsep) * \real{0.3000}}@{}}
\toprule\noalign{}
\begin{minipage}[b]{\linewidth}\raggedleft
Chapter
\end{minipage} & \begin{minipage}[b]{\linewidth}\raggedright
Topic
\end{minipage} & \begin{minipage}[b]{\linewidth}\raggedright
Lab
\end{minipage} \\
\midrule\noalign{}
\endhead
\bottomrule\noalign{}
\endlastfoot
4 & \href{chapters/chapter-04-probability-and-language.qmd}{Probability
and Language} &
\href{https://colab.research.google.com/github/wanghemath/Book-MathNLP/blob/main/labs/chapter-04-probability-and-language.ipynb}{Open
Lab 4 in Colab} \\
5 & \href{chapters/chapter-05-ngram-language-models.qmd}{N-gram Language
Models} &
\href{https://colab.research.google.com/github/wanghemath/Book-MathNLP/blob/main/labs/chapter-05-ngram-language-models.ipynb}{Open
Lab 5 in Colab} \\
6 & \href{chapters/chapter-06-naive-bayes-text-classification.qmd}{Naive
Bayes Text Classification} &
\href{https://colab.research.google.com/github/wanghemath/Book-MathNLP/blob/main/labs/chapter-06-naive-bayes-text-classification.ipynb}{Open
Lab 6 in Colab} \\
\end{longtable}
}

\subsection{Part III. Linear Algebra of
Language}\label{part-iii.-linear-algebra-of-language}

Part III treats language as geometry: vector spaces, similarity, matrix
factorization, embeddings, and representation diagnostics.

{\def\LTcaptype{none} % do not increment counter
\begin{longtable}[]{@{}
  >{\raggedleft\arraybackslash}p{(\linewidth - 4\tabcolsep) * \real{0.4000}}
  >{\raggedright\arraybackslash}p{(\linewidth - 4\tabcolsep) * \real{0.3000}}
  >{\raggedright\arraybackslash}p{(\linewidth - 4\tabcolsep) * \real{0.3000}}@{}}
\toprule\noalign{}
\begin{minipage}[b]{\linewidth}\raggedleft
Chapter
\end{minipage} & \begin{minipage}[b]{\linewidth}\raggedright
Topic
\end{minipage} & \begin{minipage}[b]{\linewidth}\raggedright
Lab
\end{minipage} \\
\midrule\noalign{}
\endhead
\bottomrule\noalign{}
\endlastfoot
7 & \href{chapters/chapter-07-vector-space-models.qmd}{Vector Space
Models} &
\href{https://colab.research.google.com/github/wanghemath/Book-MathNLP/blob/main/labs/chapter-07-vector-space-models.ipynb}{Open
Lab 7 in Colab} \\
8 & \href{chapters/chapter-08-word-embeddings.qmd}{Word Embeddings} &
\href{https://colab.research.google.com/github/wanghemath/Book-MathNLP/blob/main/labs/chapter-08-word-embeddings.ipynb}{Open
Lab 8 in Colab} \\
9 & \href{chapters/chapter-09-geometry-of-embeddings.qmd}{Geometry of
Embeddings} &
\href{https://colab.research.google.com/github/wanghemath/Book-MathNLP/blob/main/labs/chapter-09-geometry-of-embeddings.ipynb}{Open
Lab 9 in Colab} \\
\end{longtable}
}

\subsection{Part IV. Neural Networks for
Sequences}\label{part-iv.-neural-networks-for-sequences}

Part IV introduces neural computation for text, including embeddings,
recurrent networks, and gated memory.

{\def\LTcaptype{none} % do not increment counter
\begin{longtable}[]{@{}
  >{\raggedleft\arraybackslash}p{(\linewidth - 4\tabcolsep) * \real{0.4000}}
  >{\raggedright\arraybackslash}p{(\linewidth - 4\tabcolsep) * \real{0.3000}}
  >{\raggedright\arraybackslash}p{(\linewidth - 4\tabcolsep) * \real{0.3000}}@{}}
\toprule\noalign{}
\begin{minipage}[b]{\linewidth}\raggedleft
Chapter
\end{minipage} & \begin{minipage}[b]{\linewidth}\raggedright
Topic
\end{minipage} & \begin{minipage}[b]{\linewidth}\raggedright
Lab
\end{minipage} \\
\midrule\noalign{}
\endhead
\bottomrule\noalign{}
\endlastfoot
10 & \href{chapters/chapter-10-neural-networks-for-text.qmd}{Neural
Networks for Text} &
\href{https://colab.research.google.com/github/wanghemath/Book-MathNLP/blob/main/labs/chapter-10-neural-networks-for-text.ipynb}{Open
Lab 10 in Colab} \\
11 & \href{chapters/chapter-11-recurrent-neural-networks.qmd}{Recurrent
Neural Networks} &
\href{https://colab.research.google.com/github/wanghemath/Book-MathNLP/blob/main/labs/chapter-11-recurrent-neural-networks.ipynb}{Open
Lab 11 in Colab} \\
12 & \href{chapters/chapter-12-lstm-gru-models.qmd}{LSTM and GRU Models}
&
\href{https://colab.research.google.com/github/wanghemath/Book-MathNLP/blob/main/labs/chapter-12-lstm-gru-models.ipynb}{Open
Lab 12 in Colab} \\
\end{longtable}
}

\subsection{Part V. Attention and
Transformers}\label{part-v.-attention-and-transformers}

Part V develops the mathematical core of modern NLP: attention,
self-attention, positional encoding, and transformers.

{\def\LTcaptype{none} % do not increment counter
\begin{longtable}[]{@{}
  >{\raggedleft\arraybackslash}p{(\linewidth - 4\tabcolsep) * \real{0.4000}}
  >{\raggedright\arraybackslash}p{(\linewidth - 4\tabcolsep) * \real{0.3000}}
  >{\raggedright\arraybackslash}p{(\linewidth - 4\tabcolsep) * \real{0.3000}}@{}}
\toprule\noalign{}
\begin{minipage}[b]{\linewidth}\raggedleft
Chapter
\end{minipage} & \begin{minipage}[b]{\linewidth}\raggedright
Topic
\end{minipage} & \begin{minipage}[b]{\linewidth}\raggedright
Lab
\end{minipage} \\
\midrule\noalign{}
\endhead
\bottomrule\noalign{}
\endlastfoot
13 & \href{chapters/chapter-13-attention-mechanisms.qmd}{Attention
Mechanisms} &
\href{https://colab.research.google.com/github/wanghemath/Book-MathNLP/blob/main/labs/chapter-13-attention-mechanisms.ipynb}{Open
Lab 13 in Colab} \\
14 & \href{chapters/chapter-14-self-attention.qmd}{Self-Attention} &
\href{https://colab.research.google.com/github/wanghemath/Book-MathNLP/blob/main/labs/chapter-14-self-attention.ipynb}{Open
Lab 14 in Colab} \\
15 & \href{chapters/chapter-15-transformers.qmd}{Transformers} &
\href{https://colab.research.google.com/github/wanghemath/Book-MathNLP/blob/main/labs/chapter-15-transformers.ipynb}{Open
Lab 15 in Colab} \\
16 & \href{chapters/chapter-16-positional-encoding.qmd}{Positional
Encoding and Sequence Geometry} &
\href{https://colab.research.google.com/github/wanghemath/Book-MathNLP/blob/main/labs/chapter-16-positional-encoding.ipynb}{Open
Lab 16 in Colab} \\
\end{longtable}
}

\subsection{Part VI. Large Language Models and
Retrieval}\label{part-vi.-large-language-models-and-retrieval}

Part VI studies autoregressive models, training objectives,
retrieval-augmented generation, and evaluation.

{\def\LTcaptype{none} % do not increment counter
\begin{longtable}[]{@{}
  >{\raggedleft\arraybackslash}p{(\linewidth - 4\tabcolsep) * \real{0.4000}}
  >{\raggedright\arraybackslash}p{(\linewidth - 4\tabcolsep) * \real{0.3000}}
  >{\raggedright\arraybackslash}p{(\linewidth - 4\tabcolsep) * \real{0.3000}}@{}}
\toprule\noalign{}
\begin{minipage}[b]{\linewidth}\raggedleft
Chapter
\end{minipage} & \begin{minipage}[b]{\linewidth}\raggedright
Topic
\end{minipage} & \begin{minipage}[b]{\linewidth}\raggedright
Lab
\end{minipage} \\
\midrule\noalign{}
\endhead
\bottomrule\noalign{}
\endlastfoot
17 &
\href{chapters/chapter-17-autoregressive-language-models.qmd}{Autoregressive
Language Models} &
\href{https://colab.research.google.com/github/wanghemath/Book-MathNLP/blob/main/labs/chapter-17-autoregressive-language-models.ipynb}{Open
Lab 17 in Colab} \\
18 &
\href{chapters/chapter-18-training-large-language-models.qmd}{Training
Large Language Models} &
\href{https://colab.research.google.com/github/wanghemath/Book-MathNLP/blob/main/labs/chapter-18-training-large-language-models.ipynb}{Open
Lab 18 in Colab} \\
19 & \href{chapters/chapter-19-embeddings-retrieval-rag.qmd}{Embeddings,
Retrieval, and RAG} &
\href{https://colab.research.google.com/github/wanghemath/Book-MathNLP/blob/main/labs/chapter-19-embeddings-retrieval-rag.ipynb}{Open
Lab 19 in Colab} \\
20 & \href{chapters/chapter-20-evaluation.qmd}{Evaluation} &
\href{https://colab.research.google.com/github/wanghemath/Book-MathNLP/blob/main/labs/chapter-20-evaluation-of-language-models.ipynb}{Open
Lab 20 in Colab} \\
\end{longtable}
}

\subsection{Part VII. Advanced Mathematical
Perspectives}\label{part-vii.-advanced-mathematical-perspectives}

Part VII connects NLP to information theory, optimization, graphs,
topology, geometry, and responsible AI.

{\def\LTcaptype{none} % do not increment counter
\begin{longtable}[]{@{}
  >{\raggedleft\arraybackslash}p{(\linewidth - 4\tabcolsep) * \real{0.4000}}
  >{\raggedright\arraybackslash}p{(\linewidth - 4\tabcolsep) * \real{0.3000}}
  >{\raggedright\arraybackslash}p{(\linewidth - 4\tabcolsep) * \real{0.3000}}@{}}
\toprule\noalign{}
\begin{minipage}[b]{\linewidth}\raggedleft
Chapter
\end{minipage} & \begin{minipage}[b]{\linewidth}\raggedright
Topic
\end{minipage} & \begin{minipage}[b]{\linewidth}\raggedright
Lab
\end{minipage} \\
\midrule\noalign{}
\endhead
\bottomrule\noalign{}
\endlastfoot
21 & \href{chapters/chapter-21-information-theory.qmd}{Information
Theory} &
\href{https://colab.research.google.com/github/wanghemath/Book-MathNLP/blob/main/labs/chapter-21-information-theory.ipynb}{Open
Lab 21 in Colab} \\
22 & \href{chapters/chapter-22-optimization-for-nlp.qmd}{Optimization
for NLP} &
\href{https://colab.research.google.com/github/wanghemath/Book-MathNLP/blob/main/labs/chapter-22-optimization-for-nlp.ipynb}{Open
Lab 22 in Colab} \\
23 & \href{chapters/chapter-23-graphs-and-language.qmd}{Graphs and
Language} &
\href{https://colab.research.google.com/github/wanghemath/Book-MathNLP/blob/main/labs/chapter-23-graphs-and-language.ipynb}{Open
Lab 23 in Colab} \\
24 & \href{chapters/chapter-24-topology-geometry-language.qmd}{Topology,
Geometry, and Language} &
\href{https://colab.research.google.com/github/wanghemath/Book-MathNLP/blob/main/labs/chapter-24-topology-geometry-language.ipynb}{Open
Lab 24 in Colab} \\
25 & \href{chapters/chapter-25-ethics-bias-responsibility.qmd}{Ethics,
Bias, and Responsibility} &
\href{https://colab.research.google.com/github/wanghemath/Book-MathNLP/blob/main/labs/chapter-25-ethics-bias-responsibility.ipynb}{Open
Lab 25 in Colab} \\
\end{longtable}
}

\section{Mathematical spine of the
course}\label{mathematical-spine-of-the-course}

A useful way to read the book is to follow the mathematical objects as
they become more expressive.

\begin{Shaded}
\begin{Highlighting}[]
\NormalTok{text}
\NormalTok{  {-}\textgreater{} tokens}
\NormalTok{  {-}\textgreater{} count vectors}
\NormalTok{  {-}\textgreater{} probability distributions}
\NormalTok{  {-}\textgreater{} sparse and dense vector spaces}
\NormalTok{  {-}\textgreater{} neural sequence maps}
\NormalTok{  {-}\textgreater{} attention operators}
\NormalTok{  {-}\textgreater{} transformer blocks}
\NormalTok{  {-}\textgreater{} autoregressive probability models}
\NormalTok{  {-}\textgreater{} retrieval systems}
\NormalTok{  {-}\textgreater{} evaluated and audited NLP systems}
\end{Highlighting}
\end{Shaded}

The same mathematical ideas appear repeatedly:

\begin{itemize}
\tightlist
\item
  \textbf{Linearity:} document-term matrices, embeddings, projections,
  attention scores, and low-rank adaptation.
\item
  \textbf{Probability:} likelihood, language modeling, entropy,
  calibration, decoding, and preference learning.
\item
  \textbf{Optimization:} cross-entropy minimization, stochastic
  gradients, AdamW, clipping, and regularization.
\item
  \textbf{Geometry:} similarity, clustering, anisotropy, hubness,
  positional encoding, and representation diagnostics.
\item
  \textbf{Graphs:} co-occurrence networks, dependency structures,
  knowledge graphs, message passing, and retrieval links.
\item
  \textbf{Evaluation:} metrics, uncertainty, subgroup analysis,
  factuality, robustness, and responsible deployment.
\end{itemize}

\section{Computer labs}\label{computer-labs}

The labs are stored in the \texttt{labs/} folder and are designed for
Google Colab. Each lab includes:

\begin{itemize}
\tightlist
\item
  background knowledge before programming;
\item
  small transparent datasets;
\item
  runnable Python examples;
\item
  guided exercises;
\item
  AI-assisted reflection prompts;
\item
  a mini-project for independent study.
\end{itemize}

Start from the \href{labs-summary.qmd}{Labs Summary} page, or open any
lab directly from the chapter table above.

\section{Suggested projects}\label{suggested-projects}

Students can use the later chapters as starting points for final
projects, for example:

\begin{itemize}
\tightlist
\item
  compare TF-IDF, dense embeddings, and hybrid retrieval on a small
  course-document corpus;
\item
  audit a text classifier for subgroup performance and calibration;
\item
  build a transparent retrieval-augmented question-answering prototype;
\item
  study geometry and topology of sentence embeddings from different
  domains;
\item
  implement and evaluate decoding methods for a small autoregressive
  model;
\item
  compare graph-based and vector-based representations of language data.
\end{itemize}

See the \href{final-projects.qmd}{Final Projects} page for more ideas.

\section{Repository}\label{repository}

The GitHub repository for this book is:

\url{https://github.com/wanghemath/Book-MathNLP}

When the repository is published with GitHub Pages, the book can be
rendered with Quarto and the labs can be opened directly in Google
Colab.

\bookmarksetup{startatroot}

\chapter{Introduction}\label{introduction}

Why Natural Language Processing Is a Mathematical Subject

\hfill\break

Natural language processing is often introduced through applications:
search engines, translation systems, chatbots, summarizers, speech
interfaces, and large language models. This book takes a different but
complementary route. It asks: \textbf{what mathematical objects are
being manipulated when a machine processes language?}

The answer is richer than a single method or model. Language can be
represented as strings, finite sets, sequences, sparse vectors,
probability distributions, matrices, graphs, geometric point clouds, and
learned nonlinear maps. Modern NLP is powerful because these
representations can be connected: tokens become vectors, vectors become
probability models, probability models become neural networks, neural
networks become attention-based sequence operators, and these operators
become components of retrieval and generation systems.

This book develops NLP through that mathematical chain.

\section{Goals of the book}\label{goals-of-the-book}

The main goal is to help students understand NLP as a mathematical
discipline, not only as a collection of software tools. By the end of
the book, students should be able to:

\begin{enumerate}
\def\labelenumi{\arabic{enumi}.}
\tightlist
\item
  represent text using tokens, vocabularies, vectors, matrices, and
  probability distributions;
\item
  derive and implement basic probabilistic language models;
\item
  understand text classification through Bayes rules, likelihoods, and
  decision boundaries;
\item
  explain document-term matrices, TF-IDF, cosine similarity, SVD, and
  embedding geometry;
\item
  implement the core computations behind attention and transformer
  blocks;
\item
  interpret autoregressive language modeling through conditional
  probability and cross-entropy;
\item
  build small retrieval and retrieval-augmented generation systems;
\item
  evaluate NLP systems using metrics, uncertainty, calibration, and
  subgroup analysis;
\item
  connect NLP to information theory, optimization, graphs, geometry,
  topology, and ethics.
\end{enumerate}

The book is intentionally mathematical, but it is also computational.
Each major idea is paired with a transparent Python implementation,
usually on a small corpus that students can inspect directly.

\section{Why a mathematical
approach?}\label{why-a-mathematical-approach}

Many NLP systems are large, but their core ideas can be studied on small
examples. A sentence such as

\begin{quote}
language models learn patterns from text
\end{quote}

can already illustrate much of the mathematical structure of NLP.

It can be treated as:

{\def\LTcaptype{none} % do not increment counter
\begin{longtable}[]{@{}
  >{\raggedright\arraybackslash}p{(\linewidth - 4\tabcolsep) * \real{0.3333}}
  >{\raggedright\arraybackslash}p{(\linewidth - 4\tabcolsep) * \real{0.3333}}
  >{\raggedright\arraybackslash}p{(\linewidth - 4\tabcolsep) * \real{0.3333}}@{}}
\toprule\noalign{}
\begin{minipage}[b]{\linewidth}\raggedright
Viewpoint
\end{minipage} & \begin{minipage}[b]{\linewidth}\raggedright
Mathematical object
\end{minipage} & \begin{minipage}[b]{\linewidth}\raggedright
Example question
\end{minipage} \\
\midrule\noalign{}
\endhead
\bottomrule\noalign{}
\endlastfoot
string processing & sequence of characters & How should text be
normalized? \\
tokenization & finite sequence from a vocabulary & What counts as a word
or subword? \\
vectorization & vector in \(\mathbb R^d\) & Which coordinates represent
the document? \\
probability & empirical or learned distribution & How surprising is the
next token? \\
geometry & point in a metric or inner product space & Which texts are
similar? \\
neural modeling & parameterized nonlinear map & How are features learned
from data? \\
attention & weighted averaging operator & Which tokens should influence
each other? \\
evaluation & statistical estimator & How reliable is the measured
performance? \\
responsibility & decision under social constraints & Who is harmed by
errors or bias? \\
\end{longtable}
}

A mathematical perspective helps students see common structure across
many models. For example, cosine similarity in TF-IDF retrieval,
nearest-neighbor search in embedding space, attention scores in
transformers, and reranking in retrieval-augmented generation all
involve comparing representations. Cross-entropy appears in language
modeling, classification, calibration, and preference learning. Matrix
factorization appears in latent semantic analysis, embedding learning,
low-rank adaptation, and representation diagnostics.

\section{The central pipeline}\label{the-central-pipeline}

A simplified NLP pipeline is:

\begin{Shaded}
\begin{Highlighting}[]
\NormalTok{raw text}
\NormalTok{  {-}\textgreater{} tokens}
\NormalTok{  {-}\textgreater{} vocabulary indices}
\NormalTok{  {-}\textgreater{} vectors or embeddings}
\NormalTok{  {-}\textgreater{} model scores}
\NormalTok{  {-}\textgreater{} probability distributions}
\NormalTok{  {-}\textgreater{} predictions, retrieved documents, or generated text}
\NormalTok{  {-}\textgreater{} evaluation and auditing}
\end{Highlighting}
\end{Shaded}

Each arrow hides mathematical choices. Tokenization determines the
sample space. Vectorization determines what information is preserved.
The model class determines which dependencies can be learned. The loss
function determines what the system is optimized to do. The evaluation
protocol determines what failures are visible.

One guiding theme of the book is therefore:

\begin{tcolorbox}[enhanced jigsaw, arc=.35mm, bottomrule=.15mm, bottomtitle=1mm, breakable, colback=white, colbacktitle=quarto-callout-note-color!10!white, colframe=quarto-callout-note-color-frame, coltitle=black, left=2mm, leftrule=.75mm, opacityback=0, opacitybacktitle=0.6, rightrule=.15mm, title=\textcolor{quarto-callout-note-color}{\faInfo}\hspace{0.5em}{Note}, titlerule=0mm, toprule=.15mm, toptitle=1mm]

Every NLP system is a sequence of mathematical modeling decisions. Good
NLP requires understanding those decisions, not merely running the final
model.

\end{tcolorbox}

\section{Prerequisites}\label{prerequisites}

The book is written for advanced undergraduate students, master's
students, or beginning PhD students in applied mathematics, statistics,
data science, computer science, computational linguistics, or related
fields.

Useful background includes:

\begin{itemize}
\tightlist
\item
  \textbf{linear algebra:} vectors, matrices, inner products, norms,
  projections, eigenvalues, and singular value decomposition;
\item
  \textbf{probability:} conditional probability, likelihood,
  expectation, entropy, Bayes rule, and random variables;
\item
  \textbf{statistics and machine learning:} classification, validation,
  loss functions, empirical risk, and basic evaluation metrics;
\item
  \textbf{calculus and optimization:} gradients, gradient descent,
  stochastic gradients, and regularization;
\item
  \textbf{Python:} basic programming with lists, dictionaries, NumPy
  arrays, pandas data frames, and matplotlib plots.
\end{itemize}

Students do not need prior NLP experience. The computer labs are
designed for independent study and introduce the needed background
before programming begins.

\section{Notation}\label{notation}

The following notation is used throughout the book.

{\def\LTcaptype{none} % do not increment counter
\begin{longtable}[]{@{}
  >{\raggedright\arraybackslash}p{(\linewidth - 2\tabcolsep) * \real{0.5000}}
  >{\raggedright\arraybackslash}p{(\linewidth - 2\tabcolsep) * \real{0.5000}}@{}}
\toprule\noalign{}
\begin{minipage}[b]{\linewidth}\raggedright
Symbol
\end{minipage} & \begin{minipage}[b]{\linewidth}\raggedright
Meaning
\end{minipage} \\
\midrule\noalign{}
\endhead
\bottomrule\noalign{}
\endlastfoot
\(V\) & vocabulary set \\
\(|V|\) & vocabulary size \\
\(x_1, x_2, \ldots, x_T\) & token sequence of length \(T\) \\
\(d\) & embedding dimension or feature dimension \\
\(X \in \mathbb R^{n \times d}\) & data matrix with \(n\) examples and
\(d\) features \\
\(p(x)\) & probability of token, document, or event \(x\) \\
\(p(x_t \mid x_{<t})\) & autoregressive next-token probability \\
\(\theta\) & model parameters \\
\(\ell(\theta)\) & loss function or objective \\
\(q, k, v\) & query, key, and value vectors in attention \\
\(Q, K, V\) & query, key, and value matrices \\
\(A\) & attention weight matrix or adjacency matrix, depending on
context \\
\end{longtable}
}

For a probability distribution \(p\) over a finite set \(V\), entropy is
written as

\[
H(p) = -\sum_{x \in V} p(x) \log p(x).
\]

For two distributions \(p\) and \(q\), cross-entropy is

\[
H(p,q) = -\sum_{x \in V} p(x) \log q(x),
\]

and KL divergence is

\[
D_{\mathrm{KL}}(p \Vert q) = \sum_{x \in V} p(x) \log \frac{p(x)}{q(x)}.
\]

For vectors \(u,v \in \mathbb R^d\), cosine similarity is

\[
\operatorname{cos}(u,v) = \frac{u \cdot v}{\|u\|\,\|v\|}.
\]

The attention operation is usually written as

\[
\operatorname{Attention}(Q,K,V) = \operatorname{softmax}\left(\frac{QK^T}{\sqrt{d_k}}\right)V.
\]

These formulas will appear repeatedly in different contexts.

\section{Computational philosophy}\label{computational-philosophy}

The labs avoid treating NLP models as black boxes. Instead, students are
asked to build small versions of the main algorithms. The goal is not to
compete with industrial-scale models. The goal is to understand the
computations well enough to inspect, modify, and evaluate them.

For example, students will implement:

\begin{itemize}
\tightlist
\item
  tokenization and vocabulary construction;
\item
  document-term matrices and TF-IDF;
\item
  n-gram language models and smoothing;
\item
  Naive Bayes classification;
\item
  cosine retrieval and SVD-based latent spaces;
\item
  word embeddings from co-occurrence statistics;
\item
  attention and self-attention from matrix operations;
\item
  transformer-style blocks;
\item
  autoregressive decoding methods;
\item
  retrieval metrics and transparent RAG pipelines;
\item
  fairness, calibration, and bias diagnostics.
\end{itemize}

Whenever possible, the book uses small datasets created inside the
notebook. This makes the mathematical structure visible. Students can
print the vocabulary, inspect the matrix entries, compute probabilities
by hand, and compare manual calculations with Python output.

\section{How to use the computer
labs}\label{how-to-use-the-computer-labs}

The labs are located in the \texttt{labs/} folder and are designed to
open directly in Google Colab. Each lab follows the same
independent-study structure:

\begin{enumerate}
\def\labelenumi{\arabic{enumi}.}
\tightlist
\item
  \textbf{Background before programming.} The lab explains the core
  mathematical idea before code appears.
\item
  \textbf{Small transparent data.} Students work with short corpora or
  toy examples before moving to larger settings.
\item
  \textbf{Guided implementation.} Important formulas are translated into
  Python step by step.
\item
  \textbf{Diagnostics.} Students inspect intermediate objects such as
  vocabularies, matrices, probabilities, hidden states, attention
  weights, or retrieval rankings.
\item
  \textbf{Exercises.} Each lab contains tasks that require modification,
  experimentation, and interpretation.
\item
  \textbf{Mini-project.} Each lab ends with an open-ended extension
  suitable for independent study.
\item
  \textbf{AI-assisted reflection.} Students are encouraged to use AI
  tools to explain, critique, and test their reasoning, but not to
  replace their own mathematical work.
\end{enumerate}

A productive workflow is:

\begin{Shaded}
\begin{Highlighting}[]
\NormalTok{read the chapter}
\NormalTok{  {-}\textgreater{} open the corresponding Colab lab}
\NormalTok{  {-}\textgreater{} run every cell}
\NormalTok{  {-}\textgreater{} change the toy data}
\NormalTok{  {-}\textgreater{} answer the exercises}
\NormalTok{  {-}\textgreater{} write a short reflection}
\NormalTok{  {-}\textgreater{} connect the result back to the formulas}
\end{Highlighting}
\end{Shaded}

\section{Responsible use of AI tools}\label{responsible-use-of-ai-tools}

Because NLP is closely connected to modern generative AI, this book
treats AI tools as objects of study and as possible learning assistants.
Students may use AI systems to ask for explanations, generate test
examples, or compare interpretations. However, they should verify all
claims through computation, derivation, or source checking.

A useful rule is:

\begin{tcolorbox}[enhanced jigsaw, arc=.35mm, bottomrule=.15mm, bottomtitle=1mm, breakable, colback=white, colbacktitle=quarto-callout-important-color!10!white, colframe=quarto-callout-important-color-frame, coltitle=black, left=2mm, leftrule=.75mm, opacityback=0, opacitybacktitle=0.6, rightrule=.15mm, title=\textcolor{quarto-callout-important-color}{\faExclamation}\hspace{0.5em}{Important}, titlerule=0mm, toprule=.15mm, toptitle=1mm]

Use AI to ask better questions, not to avoid mathematical reasoning.

\end{tcolorbox}

For each AI-assisted prompt in the labs, students should record:

\begin{itemize}
\tightlist
\item
  the question they asked;
\item
  the answer they received;
\item
  what they verified independently;
\item
  what was unclear, incomplete, or incorrect;
\item
  how the result connects to the mathematical model.
\end{itemize}

This habit is especially important in later chapters on evaluation,
retrieval-augmented generation, fairness, and responsibility.

\section{Chapter roadmap}\label{chapter-roadmap}

The book has seven parts.

\subsection{Part I. Language as Data}\label{part-i.-language-as-data-1}

Chapters 1--3 introduce text as data. Students learn what NLP studies,
how text becomes tokens and vocabularies, and how documents become
vectors. These chapters establish the finite-set and vector-space
foundations of the book.

\subsection{Part II. Probability Models for
Language}\label{part-ii.-probability-models-for-language-1}

Chapters 4--6 develop probability models. Students study empirical
distributions, likelihood, entropy, n-gram language models, smoothing,
and Naive Bayes classification. These chapters show how language
modeling begins with conditional probability.

\subsection{Part III. Linear Algebra of
Language}\label{part-iii.-linear-algebra-of-language-1}

Chapters 7--9 develop vector space methods, embeddings, and geometry.
Students study document-term matrices, TF-IDF, SVD, word embeddings,
cosine similarity, anisotropy, hubness, clustering, and semantic
directions.

\subsection{Part IV. Neural Networks for
Sequences}\label{part-iv.-neural-networks-for-sequences-1}

Chapters 10--12 introduce neural networks for text. Students study
embedding lookup, pooling, recurrent neural networks, backpropagation
through time, LSTM cells, GRU cells, and gated memory.

\subsection{Part V. Attention and
Transformers}\label{part-v.-attention-and-transformers-1}

Chapters 13--16 develop attention, self-attention, transformer blocks,
and positional encoding. These chapters form the mathematical core of
modern NLP architectures.

\subsection{Part VI. Large Language Models and
Retrieval}\label{part-vi.-large-language-models-and-retrieval-1}

Chapters 17--20 study autoregressive language models, training
objectives, retrieval, RAG, and evaluation. Students learn how
next-token prediction, decoding, fine-tuning, retrieval, calibration,
and uncertainty interact in modern systems.

\subsection{Part VII. Advanced Mathematical
Perspectives}\label{part-vii.-advanced-mathematical-perspectives-1}

Chapters 21--25 connect NLP to information theory, optimization, graphs,
topology, geometry, ethics, bias, and responsibility. These chapters
invite students to view NLP as part of a larger mathematical and social
ecosystem.

\section{Learning outcomes by mathematical
theme}\label{learning-outcomes-by-mathematical-theme}

The following table summarizes the mathematical themes of the book.

{\def\LTcaptype{none} % do not increment counter
\begin{longtable}[]{@{}
  >{\raggedright\arraybackslash}p{(\linewidth - 4\tabcolsep) * \real{0.3000}}
  >{\raggedleft\arraybackslash}p{(\linewidth - 4\tabcolsep) * \real{0.4000}}
  >{\raggedright\arraybackslash}p{(\linewidth - 4\tabcolsep) * \real{0.3000}}@{}}
\toprule\noalign{}
\begin{minipage}[b]{\linewidth}\raggedright
Theme
\end{minipage} & \begin{minipage}[b]{\linewidth}\raggedleft
Chapters
\end{minipage} & \begin{minipage}[b]{\linewidth}\raggedright
What students learn
\end{minipage} \\
\midrule\noalign{}
\endhead
\bottomrule\noalign{}
\endlastfoot
finite structures & 1--3 & tokens, vocabularies, documents, indexing \\
probability & 4--6, 17--21 & likelihood, entropy, language modeling,
calibration \\
linear algebra & 3, 7--9, 13--16 & matrices, similarity, SVD,
embeddings, attention \\
optimization & 10--12, 18, 22 & losses, gradients, SGD, AdamW, clipping,
regularization \\
geometry & 8--9, 16, 24 & distances, neighborhoods, directions,
intrinsic structure \\
graphs & 19, 23 & co-occurrence graphs, knowledge graphs, random walks,
message passing \\
evaluation & 20, 25 & metrics, uncertainty, subgroup analysis, audits \\
responsibility & 20, 25 & fairness, bias, factuality, privacy, human
impact \\
\end{longtable}
}

\section{Suggested study questions}\label{suggested-study-questions}

Students should return to these questions throughout the course:

\begin{enumerate}
\def\labelenumi{\arabic{enumi}.}
\tightlist
\item
  What is the sample space?
\item
  What information is preserved by the representation?
\item
  What information is lost?
\item
  What probability distribution is being estimated or optimized?
\item
  What geometry does the representation create?
\item
  What loss function defines success during training?
\item
  What metric defines success during evaluation?
\item
  What kinds of examples are likely to fail?
\item
  Who is affected by the errors?
\item
  What evidence supports the final conclusion?
\end{enumerate}

These questions help connect early chapters on tokens and vectors with
later chapters on transformers, retrieval, and responsible AI.

\section{A first example: language as counting and
prediction}\label{a-first-example-language-as-counting-and-prediction}

Consider a tiny corpus:

\begin{Shaded}
\begin{Highlighting}[]
\NormalTok{language models learn from text}
\NormalTok{language models predict text}
\NormalTok{retrieval systems use text}
\end{Highlighting}
\end{Shaded}

From this corpus, we can construct a vocabulary

\[
V = \{\text{from}, \text{language}, \text{learn}, \text{models}, \text{predict}, \text{retrieval}, \text{systems}, \text{text}, \text{use}\}.
\]

Each document can be represented as a vector in \(\mathbb R^{|V|}\). A
language model can estimate probabilities such as

\[
p(\text{text} \mid \text{predict}).
\]

A retrieval system can compare a query vector with document vectors. A
neural model can learn dense embeddings for the same tokens. An
attention model can compute which earlier words influence a later
prediction.

Thus a simple corpus already contains the central mathematical objects
of the book: sets, vectors, matrices, probabilities, optimization
problems, and evaluation criteria.

\section{What this book is not}\label{what-this-book-is-not}

This book is not a catalog of every NLP library or a survey of every
current model. Modern NLP changes quickly, and software APIs change even
faster. Instead, this book focuses on durable mathematical principles.
Students who understand these principles will be better prepared to
learn new architectures, evaluate new systems, and use future tools
responsibly.

\section{Final perspective}\label{final-perspective}

Natural language is human, contextual, and social. Mathematical models
of language are useful because they simplify language into structures
that can be computed with. But every simplification has consequences. A
responsible mathematical approach therefore asks both technical and
critical questions:

\begin{itemize}
\tightlist
\item
  What is the model?
\item
  What does the model optimize?
\item
  What does the model ignore?
\item
  How do we know whether it works?
\item
  What happens when it fails?
\end{itemize}

The purpose of this book is to give students the mathematical and
computational tools to answer these questions carefully.

\bookmarksetup{startatroot}

\chapter{Learning Roadmap}\label{learning-roadmap}

A Guided Path Through Mathematical Natural Language Processing

\hfill\break

This roadmap explains how to move through \emph{Mathematical Foundations
of Natural Language Processing}. It connects the 25 chapters, 25 Google
Colab labs, mathematical prerequisites, learning milestones, and
possible project directions.

Natural language processing can look like a collection of separate
tools: tokenizers, classifiers, embeddings, transformers, retrieval
systems, and evaluation metrics. The purpose of this book is to show
that these tools are connected by a small number of mathematical ideas.

The main journey is:

\begin{Shaded}
\begin{Highlighting}[]
\NormalTok{raw text}
\NormalTok{  {-}\textgreater{} tokens and vocabularies}
\NormalTok{  {-}\textgreater{} vectors and matrices}
\NormalTok{  {-}\textgreater{} probability models}
\NormalTok{  {-}\textgreater{} embeddings and geometry}
\NormalTok{  {-}\textgreater{} neural sequence maps}
\NormalTok{  {-}\textgreater{} attention and transformers}
\NormalTok{  {-}\textgreater{} language models and retrieval}
\NormalTok{  {-}\textgreater{} evaluation, auditing, and responsible use}
\end{Highlighting}
\end{Shaded}

Students should not try to memorize every method. Instead, they should
ask the same questions in every chapter:

\begin{enumerate}
\def\labelenumi{\arabic{enumi}.}
\tightlist
\item
  What is the mathematical object?
\item
  What map or probability distribution is being learned?
\item
  What loss function or evaluation criterion is used?
\item
  What assumptions are made?
\item
  What can fail, and how can we diagnose it?
\end{enumerate}

\section{Suggested preparation}\label{suggested-preparation}

Before starting, students should be comfortable with the following
ideas. The book and labs review these ideas when needed, but independent
study is easier if students can recognize them.

{\def\LTcaptype{none} % do not increment counter
\begin{longtable}[]{@{}
  >{\raggedright\arraybackslash}p{(\linewidth - 4\tabcolsep) * \real{0.3333}}
  >{\raggedright\arraybackslash}p{(\linewidth - 4\tabcolsep) * \real{0.3333}}
  >{\raggedright\arraybackslash}p{(\linewidth - 4\tabcolsep) * \real{0.3333}}@{}}
\toprule\noalign{}
\begin{minipage}[b]{\linewidth}\raggedright
Area
\end{minipage} & \begin{minipage}[b]{\linewidth}\raggedright
Useful background
\end{minipage} & \begin{minipage}[b]{\linewidth}\raggedright
Where it appears
\end{minipage} \\
\midrule\noalign{}
\endhead
\bottomrule\noalign{}
\endlastfoot
Linear algebra & vectors, matrices, inner products, norms, projections,
SVD & document vectors, embeddings, attention, transformers \\
Probability & conditional probability, likelihood, expectation, entropy
& language models, Naive Bayes, decoding, evaluation \\
Optimization & gradients, loss functions, stochastic gradient descent &
neural networks, transformer training, fine-tuning \\
Statistics & train/test split, classification metrics, calibration,
uncertainty & evaluation, model comparison, auditing \\
Graphs & adjacency matrices, random walks, Laplacians & dependency
graphs, PageRank, graph neural networks \\
Geometry and topology & distances, neighborhoods, connected components,
persistence intuition & embedding geometry, representation
diagnostics \\
Python & lists, dictionaries, NumPy arrays, pandas, plotting & all
labs \\
\end{longtable}
}

Helpful companion pages:

\begin{itemize}
\tightlist
\item
  \href{math-background.qmd}{Mathematical Background}
\item
  \href{notion.qmd}{Notation and Core Notions}
\item
  \href{python-reference.qmd}{Python Reference}
\item
  \href{labs-summary.qmd}{Computer Labs}
\item
  \href{references.qmd}{References}
\end{itemize}

\section{The seven-part structure}\label{the-seven-part-structure}

The book is organized into seven parts. Each part adds a new
mathematical language for NLP.

{\def\LTcaptype{none} % do not increment counter
\begin{longtable}[]{@{}
  >{\raggedleft\arraybackslash}p{(\linewidth - 6\tabcolsep) * \real{0.2857}}
  >{\raggedleft\arraybackslash}p{(\linewidth - 6\tabcolsep) * \real{0.2857}}
  >{\raggedright\arraybackslash}p{(\linewidth - 6\tabcolsep) * \real{0.2143}}
  >{\raggedright\arraybackslash}p{(\linewidth - 6\tabcolsep) * \real{0.2143}}@{}}
\toprule\noalign{}
\begin{minipage}[b]{\linewidth}\raggedleft
Part
\end{minipage} & \begin{minipage}[b]{\linewidth}\raggedleft
Chapters
\end{minipage} & \begin{minipage}[b]{\linewidth}\raggedright
Main viewpoint
\end{minipage} & \begin{minipage}[b]{\linewidth}\raggedright
Main computational object
\end{minipage} \\
\midrule\noalign{}
\endhead
\bottomrule\noalign{}
\endlastfoot
I & 1--3 & language as data & tokens, vocabularies, document-term
matrices \\
II & 4--6 & language as probability & empirical distributions,
likelihoods, classifiers \\
III & 7--9 & language as geometry & vector spaces, embeddings,
similarity maps \\
IV & 10--12 & language as sequence computation & neural maps, hidden
states, gated memory \\
V & 13--16 & language as attention & query-key-value operators,
positional geometry \\
VI & 17--20 & language models and retrieval & autoregressive models, RAG
systems, evaluation \\
VII & 21--25 & advanced mathematical perspectives & information,
optimization, graphs, topology, ethics \\
\end{longtable}
}

\section{Part I. Language as Data}\label{part-i.-language-as-data-2}

\textbf{Goal:} Convert raw text into mathematical objects.

In this part, students learn that NLP begins with representation. A
sentence is not directly a vector or probability distribution. It must
first be segmented, normalized, indexed, and counted.

{\def\LTcaptype{none} % do not increment counter
\begin{longtable}[]{@{}
  >{\raggedleft\arraybackslash}p{(\linewidth - 4\tabcolsep) * \real{0.4000}}
  >{\raggedright\arraybackslash}p{(\linewidth - 4\tabcolsep) * \real{0.3000}}
  >{\raggedright\arraybackslash}p{(\linewidth - 4\tabcolsep) * \real{0.3000}}@{}}
\toprule\noalign{}
\begin{minipage}[b]{\linewidth}\raggedleft
Chapter
\end{minipage} & \begin{minipage}[b]{\linewidth}\raggedright
Lab
\end{minipage} & \begin{minipage}[b]{\linewidth}\raggedright
Core questions
\end{minipage} \\
\midrule\noalign{}
\endhead
\bottomrule\noalign{}
\endlastfoot
1 & \href{chapters/chapter-01-what-is-nlp.qmd}{What Is Natural Language
Processing?} /
\href{https://colab.research.google.com/github/wanghemath/Book-MathNLP/blob/main/labs/chapter-01-what-is-nlp.ipynb}{Lab
1} & What is an NLP pipeline? How does text become data? \\
2 & \href{chapters/chapter-02-text-tokens-vocabularies.qmd}{Text,
Tokens, and Vocabularies} /
\href{https://colab.research.google.com/github/wanghemath/Book-MathNLP/blob/main/labs/chapter-02-text-tokens-vocabularies.ipynb}{Lab
2} & What is a token? How do vocabulary choices change the data? \\
3 & \href{chapters/chapter-03-documents-as-vectors.qmd}{Documents as
Vectors} /
\href{https://colab.research.google.com/github/wanghemath/Book-MathNLP/blob/main/labs/chapter-03-documents-as-vectors.ipynb}{Lab
3} & How do count vectors, TF-IDF, and cosine similarity represent
documents? \\
\end{longtable}
}

\textbf{Checkpoint after Part I:} Students should be able to build a
small vocabulary, construct a document-term matrix, compute cosine
similarity, and explain why preprocessing choices affect downstream
models.

\section{Part II. Probability Models for
Language}\label{part-ii.-probability-models-for-language-2}

\textbf{Goal:} Model language using finite probability distributions.

This part introduces the probabilistic view of NLP. A language model
assigns probabilities to tokens or token sequences. A classifier assigns
probabilities or scores to labels. Smoothing, likelihood, and entropy
become central tools.

{\def\LTcaptype{none} % do not increment counter
\begin{longtable}[]{@{}
  >{\raggedleft\arraybackslash}p{(\linewidth - 4\tabcolsep) * \real{0.4000}}
  >{\raggedright\arraybackslash}p{(\linewidth - 4\tabcolsep) * \real{0.3000}}
  >{\raggedright\arraybackslash}p{(\linewidth - 4\tabcolsep) * \real{0.3000}}@{}}
\toprule\noalign{}
\begin{minipage}[b]{\linewidth}\raggedleft
Chapter
\end{minipage} & \begin{minipage}[b]{\linewidth}\raggedright
Lab
\end{minipage} & \begin{minipage}[b]{\linewidth}\raggedright
Core questions
\end{minipage} \\
\midrule\noalign{}
\endhead
\bottomrule\noalign{}
\endlastfoot
4 & \href{chapters/chapter-04-probability-and-language.qmd}{Probability
and Language} /
\href{https://colab.research.google.com/github/wanghemath/Book-MathNLP/blob/main/labs/chapter-04-probability-and-language.ipynb}{Lab
4} & What is an empirical token distribution? What are entropy and KL
divergence? \\
5 & \href{chapters/chapter-05-ngram-language-models.qmd}{N-gram Language
Models} /
\href{https://colab.research.google.com/github/wanghemath/Book-MathNLP/blob/main/labs/chapter-05-ngram-language-models.ipynb}{Lab
5} & How does the Markov assumption create an \(n\)-gram model? \\
6 & \href{chapters/chapter-06-naive-bayes-text-classification.qmd}{Naive
Bayes Text Classification} /
\href{https://colab.research.google.com/github/wanghemath/Book-MathNLP/blob/main/labs/chapter-06-naive-bayes-text-classification.ipynb}{Lab
6} & How does Bayes rule become a text classifier? Why do
log-probabilities matter? \\
\end{longtable}
}

\textbf{Checkpoint after Part II:} Students should be able to estimate
token probabilities, compute log-likelihood and perplexity, use
smoothing, and implement multinomial Naive Bayes from scratch.

\section{Part III. Linear Algebra of
Language}\label{part-iii.-linear-algebra-of-language-2}

\textbf{Goal:} Understand language as geometry in high-dimensional
vector spaces.

This part moves from sparse document vectors to dense word embeddings.
Students learn that similarity, projection, low-rank approximation, and
clustering are not merely computational tricks; they are geometric
operations.

{\def\LTcaptype{none} % do not increment counter
\begin{longtable}[]{@{}
  >{\raggedleft\arraybackslash}p{(\linewidth - 4\tabcolsep) * \real{0.4000}}
  >{\raggedright\arraybackslash}p{(\linewidth - 4\tabcolsep) * \real{0.3000}}
  >{\raggedright\arraybackslash}p{(\linewidth - 4\tabcolsep) * \real{0.3000}}@{}}
\toprule\noalign{}
\begin{minipage}[b]{\linewidth}\raggedleft
Chapter
\end{minipage} & \begin{minipage}[b]{\linewidth}\raggedright
Lab
\end{minipage} & \begin{minipage}[b]{\linewidth}\raggedright
Core questions
\end{minipage} \\
\midrule\noalign{}
\endhead
\bottomrule\noalign{}
\endlastfoot
7 & \href{chapters/chapter-07-vector-space-models.qmd}{Vector Space
Models} /
\href{https://colab.research.google.com/github/wanghemath/Book-MathNLP/blob/main/labs/chapter-07-vector-space-models.ipynb}{Lab
7} & How do document vectors support search and retrieval? \\
8 & \href{chapters/chapter-08-word-embeddings.qmd}{Word Embeddings} /
\href{https://colab.research.google.com/github/wanghemath/Book-MathNLP/blob/main/labs/chapter-08-word-embeddings.ipynb}{Lab
8} & How do co-occurrence, PMI, SVD, and negative sampling produce
embeddings? \\
9 & \href{chapters/chapter-09-geometry-of-embeddings.qmd}{Geometry of
Embeddings} /
\href{https://colab.research.google.com/github/wanghemath/Book-MathNLP/blob/main/labs/chapter-09-geometry-of-embeddings.ipynb}{Lab
9} & What do distance, angle, anisotropy, hubness, and clustering
reveal? \\
\end{longtable}
}

\textbf{Checkpoint after Part III:} Students should be able to compute
cosine nearest neighbors, interpret SVD as dimensionality reduction,
visualize embeddings, and diagnose geometric artifacts such as hubness
or anisotropy.

\section{Part IV. Neural Networks for
Sequences}\label{part-iv.-neural-networks-for-sequences-2}

\textbf{Goal:} Move from fixed vector representations to learned
sequence maps.

This part introduces neural networks for text. Students first study
embedding lookup and pooling, then learn why order-sensitive models such
as RNNs and gated RNNs were historically important.

{\def\LTcaptype{none} % do not increment counter
\begin{longtable}[]{@{}
  >{\raggedleft\arraybackslash}p{(\linewidth - 4\tabcolsep) * \real{0.4000}}
  >{\raggedright\arraybackslash}p{(\linewidth - 4\tabcolsep) * \real{0.3000}}
  >{\raggedright\arraybackslash}p{(\linewidth - 4\tabcolsep) * \real{0.3000}}@{}}
\toprule\noalign{}
\begin{minipage}[b]{\linewidth}\raggedleft
Chapter
\end{minipage} & \begin{minipage}[b]{\linewidth}\raggedright
Lab
\end{minipage} & \begin{minipage}[b]{\linewidth}\raggedright
Core questions
\end{minipage} \\
\midrule\noalign{}
\endhead
\bottomrule\noalign{}
\endlastfoot
10 & \href{chapters/chapter-10-neural-networks-for-text.qmd}{Neural
Networks for Text} /
\href{https://colab.research.google.com/github/wanghemath/Book-MathNLP/blob/main/labs/chapter-10-neural-networks-for-text.ipynb}{Lab
10} & How do embeddings, pooling, softmax, and cross-entropy form a
neural text classifier? \\
11 & \href{chapters/chapter-11-recurrent-neural-networks.qmd}{Recurrent
Neural Networks} /
\href{https://colab.research.google.com/github/wanghemath/Book-MathNLP/blob/main/labs/chapter-11-recurrent-neural-networks.ipynb}{Lab
11} & How does recurrence model order? What is backpropagation through
time? \\
12 & \href{chapters/chapter-12-lstm-gru-models.qmd}{LSTM and GRU Models}
/
\href{https://colab.research.google.com/github/wanghemath/Book-MathNLP/blob/main/labs/chapter-12-lstm-gru-models.ipynb}{Lab
12} & How do gates help preserve or forget information? \\
\end{longtable}
}

\textbf{Checkpoint after Part IV:} Students should be able to explain
hidden states, recurrent updates, vanishing gradients, gating, and the
role of nonlinear sequence maps in NLP.

\section{Part V. Attention and
Transformers}\label{part-v.-attention-and-transformers-2}

\textbf{Goal:} Build the mathematical core of modern transformer models.

This part studies attention as a matrix operator. Students learn how
queries, keys, values, masks, positional encodings, residual
connections, and normalization combine into a transformer block.

{\def\LTcaptype{none} % do not increment counter
\begin{longtable}[]{@{}
  >{\raggedleft\arraybackslash}p{(\linewidth - 4\tabcolsep) * \real{0.4000}}
  >{\raggedright\arraybackslash}p{(\linewidth - 4\tabcolsep) * \real{0.3000}}
  >{\raggedright\arraybackslash}p{(\linewidth - 4\tabcolsep) * \real{0.3000}}@{}}
\toprule\noalign{}
\begin{minipage}[b]{\linewidth}\raggedleft
Chapter
\end{minipage} & \begin{minipage}[b]{\linewidth}\raggedright
Lab
\end{minipage} & \begin{minipage}[b]{\linewidth}\raggedright
Core questions
\end{minipage} \\
\midrule\noalign{}
\endhead
\bottomrule\noalign{}
\endlastfoot
13 & \href{chapters/chapter-13-attention-mechanisms.qmd}{Attention
Mechanisms} /
\href{https://colab.research.google.com/github/wanghemath/Book-MathNLP/blob/main/labs/chapter-13-attention-mechanisms.ipynb}{Lab
13} & How do attention weights arise from query-key similarities? \\
14 & \href{chapters/chapter-14-self-attention.qmd}{Self-Attention} /
\href{https://colab.research.google.com/github/wanghemath/Book-MathNLP/blob/main/labs/chapter-14-self-attention.ipynb}{Lab
14} & How does each token become context-dependent? Why do masks
matter? \\
15 & \href{chapters/chapter-15-transformers.qmd}{Transformers} /
\href{https://colab.research.google.com/github/wanghemath/Book-MathNLP/blob/main/labs/chapter-15-transformers.ipynb}{Lab
15} & How do attention, feedforward maps, residuals, and layer
normalization form a block? \\
16 & \href{chapters/chapter-16-positional-encoding.qmd}{Positional
Encoding and Sequence Geometry} /
\href{https://colab.research.google.com/github/wanghemath/Book-MathNLP/blob/main/labs/chapter-16-positional-encoding.ipynb}{Lab
16} & How is order represented in a model built from set-like attention
operations? \\
\end{longtable}
}

\textbf{Checkpoint after Part V:} Students should be able to implement
scaled dot-product attention, explain masking, derive the shapes of
transformer matrices, and describe why positional information is
necessary.

\section{Part VI. Large Language Models and
Retrieval}\label{part-vi.-large-language-models-and-retrieval-2}

\textbf{Goal:} Connect transformer mathematics to language modeling,
training, retrieval, and evaluation.

This part focuses on systems that generate and retrieve language.
Students study autoregressive factorization, training objectives,
supervised fine-tuning, preference learning, retrieval-augmented
generation, and model evaluation.

{\def\LTcaptype{none} % do not increment counter
\begin{longtable}[]{@{}
  >{\raggedleft\arraybackslash}p{(\linewidth - 4\tabcolsep) * \real{0.4000}}
  >{\raggedright\arraybackslash}p{(\linewidth - 4\tabcolsep) * \real{0.3000}}
  >{\raggedright\arraybackslash}p{(\linewidth - 4\tabcolsep) * \real{0.3000}}@{}}
\toprule\noalign{}
\begin{minipage}[b]{\linewidth}\raggedleft
Chapter
\end{minipage} & \begin{minipage}[b]{\linewidth}\raggedright
Lab
\end{minipage} & \begin{minipage}[b]{\linewidth}\raggedright
Core questions
\end{minipage} \\
\midrule\noalign{}
\endhead
\bottomrule\noalign{}
\endlastfoot
17 &
\href{chapters/chapter-17-autoregressive-language-models.qmd}{Autoregressive
Language Models} /
\href{https://colab.research.google.com/github/wanghemath/Book-MathNLP/blob/main/labs/chapter-17-autoregressive-language-models.ipynb}{Lab
17} & How does the chain rule define next-token prediction? \\
18 &
\href{chapters/chapter-18-training-large-language-models.qmd}{Training
Large Language Models} /
\href{https://colab.research.google.com/github/wanghemath/Book-MathNLP/blob/main/labs/chapter-18-training-large-language-models.ipynb}{Lab
18} & How do pretraining, SFT, preference losses, and low-rank
adaptation fit together? \\
19 & \href{chapters/chapter-19-embeddings-retrieval-rag.qmd}{Embeddings,
Retrieval, and RAG} /
\href{https://colab.research.google.com/github/wanghemath/Book-MathNLP/blob/main/labs/chapter-19-embeddings-retrieval-rag.ipynb}{Lab
19} & How do sparse, dense, hybrid, and reranked retrieval support
generation? \\
20 & \href{chapters/chapter-20-evaluation.qmd}{Evaluation} /
\href{https://colab.research.google.com/github/wanghemath/Book-MathNLP/blob/main/labs/chapter-20-evaluation-of-language-models.ipynb}{Lab
20} & How should NLP systems be evaluated statistically and
critically? \\
\end{longtable}
}

\textbf{Checkpoint after Part VI:} Students should be able to compute
perplexity, implement decoding rules, build a small retrieval pipeline,
evaluate retrieval and classification metrics, and explain why
evaluation requires uncertainty and subgroup analysis.

\section{Part VII. Advanced Mathematical
Perspectives}\label{part-vii.-advanced-mathematical-perspectives-2}

\textbf{Goal:} Use deeper mathematical viewpoints to understand and
audit NLP systems.

The final part returns to the mathematical foundations and uses them as
diagnostic tools. Information theory measures uncertainty and
divergence. Optimization explains training dynamics. Graphs model
relations. Topology studies shape. Ethics and responsibility ask whether
models should be trusted in context.

{\def\LTcaptype{none} % do not increment counter
\begin{longtable}[]{@{}
  >{\raggedleft\arraybackslash}p{(\linewidth - 4\tabcolsep) * \real{0.4000}}
  >{\raggedright\arraybackslash}p{(\linewidth - 4\tabcolsep) * \real{0.3000}}
  >{\raggedright\arraybackslash}p{(\linewidth - 4\tabcolsep) * \real{0.3000}}@{}}
\toprule\noalign{}
\begin{minipage}[b]{\linewidth}\raggedleft
Chapter
\end{minipage} & \begin{minipage}[b]{\linewidth}\raggedright
Lab
\end{minipage} & \begin{minipage}[b]{\linewidth}\raggedright
Core questions
\end{minipage} \\
\midrule\noalign{}
\endhead
\bottomrule\noalign{}
\endlastfoot
21 & \href{chapters/chapter-21-information-theory.qmd}{Information
Theory} /
\href{https://colab.research.google.com/github/wanghemath/Book-MathNLP/blob/main/labs/chapter-21-information-theory.ipynb}{Lab
21} & What do entropy, cross-entropy, KL divergence, mutual information,
and PMI measure? \\
22 & \href{chapters/chapter-22-optimization-for-nlp.qmd}{Optimization
for NLP} /
\href{https://colab.research.google.com/github/wanghemath/Book-MathNLP/blob/main/labs/chapter-22-optimization-for-nlp.ipynb}{Lab
22} & How do gradients, SGD, Adam, AdamW, clipping, and masked losses
train NLP models? \\
23 & \href{chapters/chapter-23-graphs-and-language.qmd}{Graphs and
Language} /
\href{https://colab.research.google.com/github/wanghemath/Book-MathNLP/blob/main/labs/chapter-23-graphs-and-language.ipynb}{Lab
23} & How do language relations become graphs and message-passing
systems? \\
24 & \href{chapters/chapter-24-topology-geometry-language.qmd}{Topology,
Geometry, and Language} /
\href{https://colab.research.google.com/github/wanghemath/Book-MathNLP/blob/main/labs/chapter-24-topology-geometry-language.ipynb}{Lab
24} & What can neighborhoods, connected components, Betti numbers, and
persistence reveal? \\
25 & \href{chapters/chapter-25-ethics-bias-responsibility.qmd}{Ethics,
Bias, and Responsibility} /
\href{https://colab.research.google.com/github/wanghemath/Book-MathNLP/blob/main/labs/chapter-25-ethics-bias-responsibility.ipynb}{Lab
25} & How can fairness, calibration, privacy, factuality, and
accountability be audited? \\
\end{longtable}
}

\textbf{Checkpoint after Part VII:} Students should be able to connect
advanced mathematical tools to practical NLP diagnostics and explain why
responsible use is part of the mathematical modeling workflow.

\section{Three possible learning
pathways}\label{three-possible-learning-pathways}

Different students may use the book in different ways.

\subsection{Pathway A. Mathematical foundations
first}\label{pathway-a.-mathematical-foundations-first}

This pathway is best for students who want a rigorous conceptual
foundation before studying large models.

\begin{Shaded}
\begin{Highlighting}[]
\NormalTok{Chapters 1{-}{-}9}
\NormalTok{then 21{-}{-}22}
\NormalTok{then 10{-}{-}20}
\NormalTok{then 23{-}{-}25}
\end{Highlighting}
\end{Shaded}

This path emphasizes tokens, vectors, probability, geometry, information
theory, and optimization before modern architectures.

\subsection{Pathway B. Modern NLP systems
first}\label{pathway-b.-modern-nlp-systems-first}

This pathway is best for students who want to reach transformers, LLMs,
and RAG quickly.

\begin{Shaded}
\begin{Highlighting}[]
\NormalTok{Chapters 1{-}{-}3}
\NormalTok{then 10}
\NormalTok{then 13{-}{-}20}
\NormalTok{then 4{-}{-}9 and 21{-}{-}25 as deeper foundations}
\end{Highlighting}
\end{Shaded}

This path introduces the main system components early, then returns to
the mathematical foundations in more depth.

\subsection{Pathway C. Full semester
course}\label{pathway-c.-full-semester-course}

This pathway is best for a 14-week graduate course.

{\def\LTcaptype{none} % do not increment counter
\begin{longtable}[]{@{}
  >{\raggedleft\arraybackslash}p{(\linewidth - 6\tabcolsep) * \real{0.3077}}
  >{\raggedright\arraybackslash}p{(\linewidth - 6\tabcolsep) * \real{0.2308}}
  >{\raggedright\arraybackslash}p{(\linewidth - 6\tabcolsep) * \real{0.2308}}
  >{\raggedright\arraybackslash}p{(\linewidth - 6\tabcolsep) * \real{0.2308}}@{}}
\toprule\noalign{}
\begin{minipage}[b]{\linewidth}\raggedleft
Week
\end{minipage} & \begin{minipage}[b]{\linewidth}\raggedright
Main chapters
\end{minipage} & \begin{minipage}[b]{\linewidth}\raggedright
Labs
\end{minipage} & \begin{minipage}[b]{\linewidth}\raggedright
Theme
\end{minipage} \\
\midrule\noalign{}
\endhead
\bottomrule\noalign{}
\endlastfoot
1 & 1--2 & 1--2 & NLP pipeline, tokenization, vocabularies \\
2 & 3--4 & 3--4 & document vectors, probability distributions \\
3 & 5--6 & 5--6 & \(n\)-gram models, Naive Bayes \\
4 & 7--8 & 7--8 & vector space retrieval, embeddings \\
5 & 9--10 & 9--10 & embedding geometry, neural text models \\
6 & 11--12 & 11--12 & RNNs, LSTMs, GRUs \\
7 & 13--14 & 13--14 & attention and self-attention \\
8 & 15--16 & 15--16 & transformers and positional encoding \\
9 & 17--18 & 17--18 & autoregressive LMs and training objectives \\
10 & 19--20 & 19--20 & retrieval, RAG, evaluation \\
11 & 21--22 & 21--22 & information theory and optimization \\
12 & 23--24 & 23--24 & graphs, topology, geometry \\
13 & 25 & 25 & ethics, bias, responsibility \\
14 & project presentations & project notebooks & synthesis and
communication \\
\end{longtable}
}

\section{Independent-study workflow}\label{independent-study-workflow}

For each chapter, students should complete the following cycle.

\begin{enumerate}
\def\labelenumi{\arabic{enumi}.}
\tightlist
\item
  \textbf{Read the chapter.} Focus on definitions, equations, and
  examples.
\item
  \textbf{Reproduce one calculation by hand.} For example, compute a
  cosine similarity, one row of attention, or one Naive Bayes log-score.
\item
  \textbf{Open the Colab lab.} Run all cells once before changing
  anything.
\item
  \textbf{Modify one variable.} Change a smoothing parameter, embedding
  dimension, threshold, temperature, or graph radius.
\item
  \textbf{Explain the effect.} Write two or three sentences connecting
  the result to the mathematics.
\item
  \textbf{Complete the exercises.} Use the hints only after attempting
  the problem.
\item
  \textbf{Complete the mini-project.} Extend the lab to a new toy
  dataset or new diagnostic question.
\item
  \textbf{Write a reflection.} State one assumption, one limitation, and
  one possible improvement.
\end{enumerate}

\section{Milestones and evidence of
learning}\label{milestones-and-evidence-of-learning}

The following milestones can be used for self-assessment or course
grading.

{\def\LTcaptype{none} % do not increment counter
\begin{longtable}[]{@{}
  >{\raggedright\arraybackslash}p{(\linewidth - 2\tabcolsep) * \real{0.5000}}
  >{\raggedright\arraybackslash}p{(\linewidth - 2\tabcolsep) * \real{0.5000}}@{}}
\toprule\noalign{}
\begin{minipage}[b]{\linewidth}\raggedright
Milestone
\end{minipage} & \begin{minipage}[b]{\linewidth}\raggedright
Evidence
\end{minipage} \\
\midrule\noalign{}
\endhead
\bottomrule\noalign{}
\endlastfoot
Representation fluency & Build token lists, vocabularies, and
document-term matrices without templates \\
Probabilistic reasoning & Compute likelihood, log-likelihood, entropy,
KL divergence, and perplexity \\
Geometric reasoning & Use cosine similarity, projections, SVD,
clustering, and visualization appropriately \\
Neural modeling & Explain embeddings, softmax, cross-entropy, gradients,
recurrence, and gating \\
Transformer literacy & Implement attention and identify the role of
masks, heads, residuals, and positions \\
Retrieval literacy & Compare sparse, dense, hybrid, and reranked
retrieval systems \\
Evaluation literacy & Report metrics with uncertainty and examine
subgroup behavior \\
Responsible modeling & Identify bias, distribution shift, factuality
risk, and deployment limitations \\
Communication & Explain mathematical NLP ideas clearly in writing and
code comments \\
\end{longtable}
}

\section{Suggested final project
directions}\label{suggested-final-project-directions}

A good final project should combine mathematics, computation, and
critical evaluation. Possible directions include:

\begin{enumerate}
\def\labelenumi{\arabic{enumi}.}
\tightlist
\item
  \textbf{Transparent retrieval system:} compare TF-IDF, dense toy
  embeddings, hybrid retrieval, and reranking on a course-document
  corpus.
\item
  \textbf{Small language model study:} train and evaluate an \(n\)-gram
  or small neural next-token model with different smoothing or decoding
  choices.
\item
  \textbf{Embedding geometry audit:} compare nearest neighbors, PCA
  plots, anisotropy, hubness, and semantic directions across corpora.
\item
  \textbf{RAG reliability audit:} build a small RAG system and evaluate
  retrieval quality, citation support, and answer faithfulness.
\item
  \textbf{Fairness and calibration study:} evaluate a text classifier
  across subgroups and study threshold, calibration, and uncertainty
  tradeoffs.
\item
  \textbf{Graph representation project:} build a co-occurrence,
  dependency, or knowledge graph and analyze PageRank, communities, or
  message passing.
\item
  \textbf{Topology of embeddings:} build neighborhood graphs or
  Vietoris--Rips summaries for embeddings and compare geometric shape
  across labels or domains.
\end{enumerate}

A strong project report should include:

\begin{itemize}
\tightlist
\item
  a clear research question;
\item
  the mathematical representation used;
\item
  the dataset and preprocessing choices;
\item
  the model or algorithm;
\item
  quantitative results;
\item
  visualizations or diagnostic tables;
\item
  limitations and responsible-use discussion.
\end{itemize}

\section{Recommended repository
organization}\label{recommended-repository-organization}

The GitHub repository should have a structure similar to:

\begin{Shaded}
\begin{Highlighting}[]
\NormalTok{Book{-}MathNLP/}
\NormalTok{  index.qmd}
\NormalTok{  introduction.qmd}
\NormalTok{  math{-}background.qmd}
\NormalTok{  notion.qmd}
\NormalTok{  python{-}reference.qmd}
\NormalTok{  labs{-}summary.qmd}
\NormalTok{  roadmap.qmd}
\NormalTok{  references.qmd}
\NormalTok{  references.bib}
\NormalTok{  chapters/}
\NormalTok{    chapter{-}01{-}what{-}is{-}nlp.qmd}
\NormalTok{    ...}
\NormalTok{    chapter{-}25{-}ethics{-}bias{-}responsibility.qmd}
\NormalTok{  labs/}
\NormalTok{    chapter{-}01{-}what{-}is{-}nlp.ipynb}
\NormalTok{    ...}
\NormalTok{    chapter{-}25{-}ethics{-}bias{-}responsibility.ipynb}
\NormalTok{  assets/}
\NormalTok{    style.scss}
\NormalTok{    interactive.js}
\end{Highlighting}
\end{Shaded}

\section{Closing perspective}\label{closing-perspective}

The central message of the book is that modern NLP is not magic. It is a
layered mathematical system. Text is discretized into tokens. Tokens
become vectors. Vectors define probabilities, similarities, graphs, and
neural maps. Models are trained by optimization and judged by
statistical evidence. Responsible use requires understanding not only
how a model works, but also where it fails.

Students who complete this roadmap should be able to read modern NLP
papers and systems with a clearer question in mind:

\begin{quote}
What are the objects, maps, probabilities, losses, and assumptions
behind this language system?
\end{quote}

\bookmarksetup{startatroot}

\chapter{Notation and Core Notions}\label{notation-and-core-notions}

A compact reference for mathematical NLP

\hfill\break

This page collects the main notation, conventions, and core notions used
throughout the book. It is designed as a quick reference for students
moving between chapters, computer labs, and projects.

Natural language processing combines discrete objects, continuous
representations, probability models, optimization algorithms, and
evaluation protocols. The same symbol may appear in several contexts, so
this page explains the standard meaning used in this book and notes when
a symbol has context-dependent interpretations.

\section{1. General conventions}\label{general-conventions}

Unless otherwise stated:

\begin{itemize}
\tightlist
\item
  vectors are column vectors;
\item
  matrices are written with capital letters;
\item
  finite sets are written with capital letters such as \(V\), \(D\), and
  \(C\);
\item
  random variables are written with capital letters such as \(X\) and
  \(Y\);
\item
  observed values are written with lowercase letters such as \(x\) and
  \(y\);
\item
  logarithms are natural logarithms unless base \(2\) is explicitly
  stated;
\item
  probabilities are defined over finite sets in most chapters;
\item
  all examples use small corpora so that the mathematical structure is
  visible.
\end{itemize}

A recurring theme is that an NLP pipeline is a composition of maps:

\[
\text{text} \longrightarrow \text{tokens} \longrightarrow \text{indices} \longrightarrow \text{vectors} \longrightarrow \text{scores} \longrightarrow \text{probabilities} \longrightarrow \text{decisions}.
\]

Each arrow represents a modeling choice.

\section{2. Text, tokens, and
vocabularies}\label{text-tokens-and-vocabularies}

{\def\LTcaptype{none} % do not increment counter
\begin{longtable}[]{@{}
  >{\raggedright\arraybackslash}p{(\linewidth - 4\tabcolsep) * \real{0.3333}}
  >{\raggedright\arraybackslash}p{(\linewidth - 4\tabcolsep) * \real{0.3333}}
  >{\raggedright\arraybackslash}p{(\linewidth - 4\tabcolsep) * \real{0.3333}}@{}}
\toprule\noalign{}
\begin{minipage}[b]{\linewidth}\raggedright
Symbol
\end{minipage} & \begin{minipage}[b]{\linewidth}\raggedright
Meaning
\end{minipage} & \begin{minipage}[b]{\linewidth}\raggedright
Typical chapter use
\end{minipage} \\
\midrule\noalign{}
\endhead
\bottomrule\noalign{}
\endlastfoot
\(s\) & raw text string & preprocessing and tokenization \\
\(x_t\) & token at position \(t\) & sequences and language models \\
\(x_{1:T}\) & token sequence \(x_1,x_2,\ldots,x_T\) & sequence
modeling \\
\(T\) & sequence length & RNNs, attention, transformers \\
\(V\) & vocabulary set & tokenization and vectorization \\
\(|V|\) & vocabulary size & language models and embeddings \\
\(w\) & a word, subword, or token type & vocabulary-level quantities \\
\(\iota\) & indexing map from tokens to integers & token ids \\
\(\langle\mathrm{unk}\rangle\) & unknown token & out-of-vocabulary
handling \\
\(\langle\mathrm{bos}\rangle\) & beginning-of-sequence token & language
modeling \\
\(\langle\mathrm{eos}\rangle\) & end-of-sequence token & language
modeling \\
\(\langle\mathrm{pad}\rangle\) & padding token & batching
variable-length sequences \\
\end{longtable}
}

A tokenizer maps a string to a finite sequence:

\[
\tau(s) = (x_1,x_2,\ldots,x_T), \qquad x_t \in V.
\]

The vocabulary is not just a list of words. It determines the sample
space for probability models and the coordinate system for count
vectors.

\begin{tcolorbox}[enhanced jigsaw, arc=.35mm, bottomrule=.15mm, bottomtitle=1mm, breakable, colback=white, colbacktitle=quarto-callout-note-color!10!white, colframe=quarto-callout-note-color-frame, coltitle=black, left=2mm, leftrule=.75mm, opacityback=0, opacitybacktitle=0.6, rightrule=.15mm, title=\textcolor{quarto-callout-note-color}{\faInfo}\hspace{0.5em}{Note}, titlerule=0mm, toprule=.15mm, toptitle=1mm]

Changing the tokenizer changes the mathematical problem. A word
tokenizer, character tokenizer, and subword tokenizer produce different
vocabularies, different sequence lengths, and different probability
distributions.

\end{tcolorbox}

\section{3. Documents and corpora}\label{documents-and-corpora}

{\def\LTcaptype{none} % do not increment counter
\begin{longtable}[]{@{}ll@{}}
\toprule\noalign{}
Symbol & Meaning \\
\midrule\noalign{}
\endhead
\bottomrule\noalign{}
\endlastfoot
\(d_i\) & document \(i\) \\
\(D\) & document collection or corpus \\
\(n\) & number of documents or examples \\
\(m\) & vocabulary size, often \(m=|V|\) \\
\(c(w,d)\) & count of token \(w\) in document \(d\) \\
\(\mathrm{tf}(w,d)\) & term frequency of token \(w\) in document
\(d\) \\
\(\mathrm{df}(w)\) & document frequency of token \(w\) \\
\(\mathrm{idf}(w)\) & inverse document frequency of token \(w\) \\
\end{longtable}
}

A corpus can be represented as a list of token sequences:

\[
D = \{x^{(1)}_{1:T_1}, x^{(2)}_{1:T_2}, \ldots, x^{(n)}_{1:T_n}\}.
\]

A document-term matrix is usually written as

\[
X \in \mathbb R^{n \times m},
\]

where row \(i\) represents document \(d_i\) and column \(j\) represents
token \(w_j\).

\section{4. Vectors, matrices, and
geometry}\label{vectors-matrices-and-geometry}

{\def\LTcaptype{none} % do not increment counter
\begin{longtable}[]{@{}
  >{\raggedright\arraybackslash}p{(\linewidth - 2\tabcolsep) * \real{0.5000}}
  >{\raggedright\arraybackslash}p{(\linewidth - 2\tabcolsep) * \real{0.5000}}@{}}
\toprule\noalign{}
\begin{minipage}[b]{\linewidth}\raggedright
Symbol
\end{minipage} & \begin{minipage}[b]{\linewidth}\raggedright
Meaning
\end{minipage} \\
\midrule\noalign{}
\endhead
\bottomrule\noalign{}
\endlastfoot
\(x \in \mathbb R^d\) & feature vector or embedding \\
\(d\) & feature dimension or embedding dimension \\
\(X \in \mathbb R^{n \times d}\) & data matrix with \(n\) examples and
\(d\) features \\
\(E \in \mathbb R^{|V| \times d}\) & embedding matrix \\
\(e_w\) & embedding vector for token \(w\) \\
\(\lVert x \rVert\) & Euclidean norm of \(x\) \\
\(x^\top y\) & inner product of vectors \(x\) and \(y\) \\
\(\cos(x,y)\) & cosine similarity \\
\(G = XX^\top\) & Gram matrix of pairwise inner products \\
\end{longtable}
}

The cosine similarity between two nonzero vectors is

\[
\cos(x,y) = \frac{x^\top y}{\lVert x \rVert\lVert y \rVert}.
\]

This score appears in TF-IDF retrieval, embedding search, clustering,
nearest neighbors, and some attention interpretations.

\subsection{Matrix factorization}\label{matrix-factorization}

For a matrix \(X \in \mathbb R^{n \times d}\), the singular value
decomposition is

\[
X = U\Sigma V^\top.
\]

In NLP, low-rank approximations are used to reveal latent structure:

\[
X_k = U_k\Sigma_k V_k^\top.
\]

The same idea appears in latent semantic analysis, embedding geometry,
representation diagnostics, and low-rank adaptation.

\section{5. Probability and language
models}\label{probability-and-language-models}

{\def\LTcaptype{none} % do not increment counter
\begin{longtable}[]{@{}ll@{}}
\toprule\noalign{}
Symbol & Meaning \\
\midrule\noalign{}
\endhead
\bottomrule\noalign{}
\endlastfoot
\(p(x)\) & probability of event or token \(x\) \\
\(p(x,y)\) & joint probability \\
\(p(y \mid x)\) & conditional probability \\
\(\hat p\) & empirical distribution estimated from data \\
\(\theta\) & model parameters \\
\(p_\theta\) & probability distribution defined by parameters
\(\theta\) \\
\(\ell(\theta)\) & loss or negative log-likelihood \\
\(L(\theta)\) & objective function \\
\end{longtable}
}

A language model assigns probabilities to token sequences. The
autoregressive factorization is

\[
p_\theta(x_{1:T}) = \prod_{t=1}^T p_\theta(x_t \mid x_{<t}),
\]

where \(x_{<t} = x_1,\ldots,x_{t-1}\).

The negative log-likelihood for one sequence is

\[
-\log p_\theta(x_{1:T}) = -\sum_{t=1}^T \log p_\theta(x_t \mid x_{<t}).
\]

Per-token average negative log-likelihood is cross-entropy when the data
distribution is fixed.

\section{6. Information theory}\label{information-theory}

{\def\LTcaptype{none} % do not increment counter
\begin{longtable}[]{@{}
  >{\raggedright\arraybackslash}p{(\linewidth - 2\tabcolsep) * \real{0.5000}}
  >{\raggedright\arraybackslash}p{(\linewidth - 2\tabcolsep) * \real{0.5000}}@{}}
\toprule\noalign{}
\begin{minipage}[b]{\linewidth}\raggedright
Symbol
\end{minipage} & \begin{minipage}[b]{\linewidth}\raggedright
Meaning
\end{minipage} \\
\midrule\noalign{}
\endhead
\bottomrule\noalign{}
\endlastfoot
\(H(p)\) & entropy of distribution \(p\) \\
\(H(p,q)\) & cross-entropy between true distribution \(p\) and model
\(q\) \\
\(D_{\mathrm{KL}}(p\Vert q)\) & Kullback--Leibler divergence from \(p\)
to \(q\) \\
\(I(X;Y)\) & mutual information between \(X\) and \(Y\) \\
\(\mathrm{PMI}(x,y)\) & pointwise mutual information \\
\(\mathrm{PPL}\) & perplexity \\
\end{longtable}
}

Entropy measures average uncertainty:

\[
H(p) = -\sum_{x \in V} p(x)\log p(x).
\]

Cross-entropy measures average code length when data come from \(p\) but
are coded using \(q\):

\[
H(p,q) = -\sum_{x \in V} p(x)\log q(x).
\]

KL divergence is the excess cross-entropy:

\[
D_{\mathrm{KL}}(p\Vert q) = \sum_{x \in V} p(x)\log\frac{p(x)}{q(x)}.
\]

Perplexity is often used to report language-model quality:

\[
\mathrm{PPL} = \exp\{H\}.
\]

Lower perplexity usually means the model assigns higher probability to
the observed tokens, but it does not automatically imply better
factuality, usefulness, fairness, or safety.

\section{7. Classification notation}\label{classification-notation}

{\def\LTcaptype{none} % do not increment counter
\begin{longtable}[]{@{}ll@{}}
\toprule\noalign{}
Symbol & Meaning \\
\midrule\noalign{}
\endhead
\bottomrule\noalign{}
\endlastfoot
\(y\) & class label \\
\(C\) & set of classes \\
\(K\) & number of classes \\
\(p(y \mid x)\) & posterior class probability \\
\(\hat y\) & predicted class \\
\(W\) & classifier weight matrix \\
\(b\) & bias vector \\
\(z\) & logits or unnormalized scores \\
\end{longtable}
}

A common linear classifier computes logits

\[
z = Wx + b,
\]

then converts logits to probabilities using softmax:

\[
\mathrm{softmax}(z)_j = \frac{\exp(z_j)}{\sum_{k=1}^K \exp(z_k)}.
\]

The cross-entropy loss for a true class \(y\) is

\[
\ell(z,y) = -\log \mathrm{softmax}(z)_y.
\]

\section{8. Neural network notation}\label{neural-network-notation}

{\def\LTcaptype{none} % do not increment counter
\begin{longtable}[]{@{}ll@{}}
\toprule\noalign{}
Symbol & Meaning \\
\midrule\noalign{}
\endhead
\bottomrule\noalign{}
\endlastfoot
\(h_t\) & hidden state at position \(t\) \\
\(c_t\) & LSTM cell state at position \(t\) \\
\(\sigma\) & sigmoid function \\
\(\tanh\) & hyperbolic tangent nonlinearity \\
\(\phi\) & generic nonlinear activation \\
\(W\), \(U\), \(b\) & neural-network parameters \\
\(\nabla_\theta L\) & gradient of objective with respect to
parameters \\
\(\eta\) & learning rate \\
\end{longtable}
}

A simple recurrent neural network has the form

\[
h_t = \phi(W_x x_t + W_h h_{t-1} + b).
\]

A feedforward layer has the form

\[
h = \phi(Wx+b).
\]

Training usually means solving an empirical risk minimization problem:

\[
\min_\theta \frac{1}{n}\sum_{i=1}^n \ell(f_\theta(x_i),y_i).
\]

\section{9. Attention and
transformers}\label{attention-and-transformers}

{\def\LTcaptype{none} % do not increment counter
\begin{longtable}[]{@{}ll@{}}
\toprule\noalign{}
Symbol & Meaning \\
\midrule\noalign{}
\endhead
\bottomrule\noalign{}
\endlastfoot
\(q\) & query vector \\
\(k\) & key vector \\
\(v\) & value vector \\
\(Q\) & query matrix \\
\(K\) & key matrix \\
\(V\) & value matrix \\
\(A\) & attention weight matrix \\
\(d_k\) & key/query dimension \\
\(M\) & attention mask \\
\(P\) & positional encoding matrix \\
\end{longtable}
}

Scaled dot-product attention is

\[
\mathrm{Attention}(Q,K,V) = \mathrm{softmax}\left(\frac{QK^\top}{\sqrt{d_k}} + M\right)V.
\]

The matrix

\[
A = \mathrm{softmax}\left(\frac{QK^\top}{\sqrt{d_k}} + M\right)
\]

contains attention weights. Each row is usually a probability
distribution over source positions.

\begin{tcolorbox}[enhanced jigsaw, arc=.35mm, bottomrule=.15mm, bottomtitle=1mm, breakable, colback=white, colbacktitle=quarto-callout-warning-color!10!white, colframe=quarto-callout-warning-color-frame, coltitle=black, left=2mm, leftrule=.75mm, opacityback=0, opacitybacktitle=0.6, rightrule=.15mm, title=\textcolor{quarto-callout-warning-color}{\faExclamationTriangle}\hspace{0.5em}{Warning}, titlerule=0mm, toprule=.15mm, toptitle=1mm]

Attention weights are useful diagnostics, but they are not automatically
faithful explanations of model behavior. They show one internal
weighting pattern, not the full causal reason for a prediction.

\end{tcolorbox}

A transformer block combines attention, residual connections,
normalization, and feedforward layers. A simplified encoder block is

\[
Z_1 = \mathrm{LayerNorm}(X + \mathrm{SelfAttention}(X)),
\]

\[
Z_2 = \mathrm{LayerNorm}(Z_1 + \mathrm{FFN}(Z_1)).
\]

\section{10. Retrieval and RAG
notation}\label{retrieval-and-rag-notation}

{\def\LTcaptype{none} % do not increment counter
\begin{longtable}[]{@{}
  >{\raggedright\arraybackslash}p{(\linewidth - 2\tabcolsep) * \real{0.5000}}
  >{\raggedright\arraybackslash}p{(\linewidth - 2\tabcolsep) * \real{0.5000}}@{}}
\toprule\noalign{}
\begin{minipage}[b]{\linewidth}\raggedright
Symbol
\end{minipage} & \begin{minipage}[b]{\linewidth}\raggedright
Meaning
\end{minipage} \\
\midrule\noalign{}
\endhead
\bottomrule\noalign{}
\endlastfoot
\(q\) & user query, or query vector depending on context \\
\(d_i\) & document or chunk \(i\) \\
\(r(q,d_i)\) & retrieval score for query \(q\) and document \(d_i\) \\
\(\mathrm{top}_k(q)\) & top \(k\) retrieved documents for query \(q\) \\
\(\mathrm{rel}(q,d)\) & relevance label \\
\(\mathrm{Recall@}k\) & fraction of relevant items retrieved in top
\(k\) \\
\(\mathrm{Precision@}k\) & fraction of top \(k\) results that are
relevant \\
\(\mathrm{MRR}\) & mean reciprocal rank \\
\(\mathrm{nDCG}\) & normalized discounted cumulative gain \\
\end{longtable}
}

A retrieval system ranks documents by a score:

\[
r(q,d_i) = \mathrm{sim}(\psi(q),\psi(d_i)),
\]

where \(\psi\) is a representation map and \(\mathrm{sim}\) is a
similarity function.

A retrieval-augmented generation system can be summarized as

\[
\text{query} \longrightarrow \text{retrieve context} \longrightarrow \text{generate answer conditioned on context}.
\]

When evaluating RAG, separate the retrieval problem from the generation
problem. A wrong answer may be caused by poor retrieval, poor context
use, hallucination, or ambiguity in the question.

\section{11. Graph notation}\label{graph-notation}

{\def\LTcaptype{none} % do not increment counter
\begin{longtable}[]{@{}ll@{}}
\toprule\noalign{}
Symbol & Meaning \\
\midrule\noalign{}
\endhead
\bottomrule\noalign{}
\endlastfoot
\(G=(N,E)\) & graph with node set \(N\) and edge set \(E\) \\
\(A\) & adjacency matrix \\
\(D\) & degree matrix, or corpus depending on context \\
\(L=D-A\) & graph Laplacian \\
\(\tilde A\) & normalized or self-loop augmented adjacency matrix \\
\(N(i)\) & neighborhood of node \(i\) \\
\(h_i\) & node representation \\
\end{longtable}
}

A graph convolution layer often has the form

\[
H' = \phi(\tilde AHW),
\]

where \(H\) is a matrix of node features, \(W\) is a learnable weight
matrix, and \(\tilde A\) controls message passing.

Graph views of language include dependency graphs, word co-occurrence
graphs, document-term bipartite graphs, citation graphs, and knowledge
graphs.

\section{12. Topology and geometry
notation}\label{topology-and-geometry-notation}

{\def\LTcaptype{none} % do not increment counter
\begin{longtable}[]{@{}
  >{\raggedright\arraybackslash}p{(\linewidth - 2\tabcolsep) * \real{0.5000}}
  >{\raggedright\arraybackslash}p{(\linewidth - 2\tabcolsep) * \real{0.5000}}@{}}
\toprule\noalign{}
\begin{minipage}[b]{\linewidth}\raggedright
Symbol
\end{minipage} & \begin{minipage}[b]{\linewidth}\raggedright
Meaning
\end{minipage} \\
\midrule\noalign{}
\endhead
\bottomrule\noalign{}
\endlastfoot
\(X\) & finite point cloud or data matrix \\
\(d(x,y)\) & distance between points \(x\) and \(y\) \\
\(\epsilon\) & scale parameter \\
\(G_\epsilon\) & neighborhood graph at scale \(\epsilon\) \\
\(\beta_0\) & number of connected components \\
\(\beta_1\) & number of one-dimensional loops \\
\(\mathrm{VR}_\epsilon(X)\) & Vietoris--Rips complex at scale
\(\epsilon\) \\
\end{longtable}
}

A Vietoris--Rips complex at scale \(\epsilon\) contains a simplex
whenever all its vertices are pairwise within distance \(\epsilon\).

For NLP, topology is used as a diagnostic lens: it can help compare
embedding spaces, detect clustering, study loops or branches in
representation trajectories, and summarize how local neighborhoods
change with scale.

\section{13. Evaluation and responsibility
notation}\label{evaluation-and-responsibility-notation}

{\def\LTcaptype{none} % do not increment counter
\begin{longtable}[]{@{}ll@{}}
\toprule\noalign{}
Symbol & Meaning \\
\midrule\noalign{}
\endhead
\bottomrule\noalign{}
\endlastfoot
\(\mathrm{TP}\) & true positives \\
\(\mathrm{FP}\) & false positives \\
\(\mathrm{TN}\) & true negatives \\
\(\mathrm{FN}\) & false negatives \\
\(\mathrm{Acc}\) & accuracy \\
\(\mathrm{Prec}\) & precision \\
\(\mathrm{Rec}\) & recall \\
\(F_1\) & harmonic mean of precision and recall \\
\(\mathrm{ECE}\) & expected calibration error \\
\(g\) & group or subgroup variable \\
\(\Delta_g\) & performance gap across groups \\
\end{longtable}
}

Precision and recall are

\[
\mathrm{Precision} = \frac{\mathrm{TP}}{\mathrm{TP}+\mathrm{FP}},
\]

\[
\mathrm{Recall} = \frac{\mathrm{TP}}{\mathrm{TP}+\mathrm{FN}}.
\]

The \(F_1\) score is

\[
F_1 = \frac{2\cdot \mathrm{Precision}\cdot \mathrm{Recall}}{\mathrm{Precision}+\mathrm{Recall}}.
\]

Calibration asks whether predicted probabilities match empirical
frequencies. Fairness and responsibility require checking not only
global averages but also subgroup errors, uncertainty, data provenance,
privacy, and possible misuse.

\section{14. Common overloaded symbols}\label{common-overloaded-symbols}

Some symbols are used differently in different contexts. The local
meaning should be clear from the chapter.

{\def\LTcaptype{none} % do not increment counter
\begin{longtable}[]{@{}
  >{\raggedright\arraybackslash}p{(\linewidth - 2\tabcolsep) * \real{0.5000}}
  >{\raggedright\arraybackslash}p{(\linewidth - 2\tabcolsep) * \real{0.5000}}@{}}
\toprule\noalign{}
\begin{minipage}[b]{\linewidth}\raggedright
Symbol
\end{minipage} & \begin{minipage}[b]{\linewidth}\raggedright
Possible meanings
\end{minipage} \\
\midrule\noalign{}
\endhead
\bottomrule\noalign{}
\endlastfoot
\(X\) & data matrix, random variable, point cloud \\
\(D\) & corpus, degree matrix, dataset \\
\(A\) & attention matrix, adjacency matrix, event \\
\(q\) & probability distribution, query text, query vector \\
\(d\) & document, dimension, distance \\
\(V\) & vocabulary, value matrix in attention when written as \(V\) \\
\(K\) & number of classes, key matrix, number of clusters \\
\(L\) & loss function, graph Laplacian, sequence length in some
contexts \\
\end{longtable}
}

When ambiguity is possible, the text will specify the meaning
explicitly.

\section{15. Chapter-by-chapter notation
map}\label{chapter-by-chapter-notation-map}

{\def\LTcaptype{none} % do not increment counter
\begin{longtable}[]{@{}
  >{\raggedright\arraybackslash}p{(\linewidth - 2\tabcolsep) * \real{0.5000}}
  >{\raggedright\arraybackslash}p{(\linewidth - 2\tabcolsep) * \real{0.5000}}@{}}
\toprule\noalign{}
\begin{minipage}[b]{\linewidth}\raggedright
Chapters
\end{minipage} & \begin{minipage}[b]{\linewidth}\raggedright
Main notation to know
\end{minipage} \\
\midrule\noalign{}
\endhead
\bottomrule\noalign{}
\endlastfoot
1--3 & \(s\), \(x_t\), \(V\), \(D\), \(X\), term counts, TF-IDF \\
4--6 & \(p\), \(\hat p\), likelihood, Bayes rule, class labels \\
7--9 & vectors, norms, cosine similarity, Gram matrices, SVD,
embeddings \\
10--12 & neural parameters, hidden states, gates, gradients \\
13--16 & \(Q\), \(K\), \(V\), attention matrices, masks, positional
encodings \\
17--20 & autoregressive probabilities, cross-entropy, perplexity,
retrieval metrics \\
21--22 & entropy, KL divergence, optimization objectives, SGD, AdamW \\
23--24 & adjacency matrices, Laplacians, point clouds, Betti numbers \\
25 & calibration, subgroup metrics, fairness gaps, privacy
diagnostics \\
\end{longtable}
}

\section{16. Reading formulas in this
book}\label{reading-formulas-in-this-book}

Many formulas should be read both algebraically and computationally. For
example,

\[
z = Wx+b
\]

means all of the following at once:

\begin{enumerate}
\def\labelenumi{\arabic{enumi}.}
\tightlist
\item
  algebraically, a linear map followed by a translation;
\item
  geometrically, a projection onto score directions;
\item
  statistically, a model for class logits;
\item
  computationally, a matrix multiplication plus broadcasting;
\item
  pedagogically, a place where dimensions must match.
\end{enumerate}

When studying a new formula, ask:

\begin{itemize}
\tightlist
\item
  What are the dimensions of all objects?
\item
  Which quantities are fixed data and which are parameters?
\item
  Which operation creates the output?
\item
  What loss or metric will evaluate the output?
\item
  What failure modes are hidden by the notation?
\end{itemize}

\section{17. Python naming conventions used in
labs}\label{python-naming-conventions-used-in-labs}

The labs use Python variable names that mirror the mathematical
notation.

{\def\LTcaptype{none} % do not increment counter
\begin{longtable}[]{@{}ll@{}}
\toprule\noalign{}
Math notation & Python variable examples \\
\midrule\noalign{}
\endhead
\bottomrule\noalign{}
\endlastfoot
\(V\) & \texttt{vocab}, \texttt{vocab\_list}, \texttt{token\_to\_id} \\
\(X\) & \texttt{X}, \texttt{X\_counts}, \texttt{X\_tfidf},
\texttt{X\_emb} \\
\(y\) & \texttt{y}, \texttt{labels}, \texttt{y\_true},
\texttt{y\_pred} \\
\(p\) & \texttt{p}, \texttt{probs}, \texttt{p\_hat} \\
\(Q,K,V\) & \texttt{Q}, \texttt{K}, \texttt{V\_values} \\
\(A\) & \texttt{A}, \texttt{attention}, \texttt{adjacency} \\
\(\theta\) & \texttt{params}, \texttt{weights},
\texttt{model\_params} \\
\(L\) & \texttt{loss}, \texttt{objective} \\
\end{longtable}
}

In code, \texttt{V} can conflict with the value matrix in attention. For
this reason, labs often use \texttt{vocab} for the vocabulary and
\texttt{V\_values} for the value matrix.

\section{18. Minimal checklist for
students}\label{minimal-checklist-for-students}

Before starting a chapter or lab, check that you can identify:

\begin{enumerate}
\def\labelenumi{\arabic{enumi}.}
\tightlist
\item
  the input object;
\item
  the output object;
\item
  the representation space;
\item
  the probability distribution or scoring function;
\item
  the loss, metric, or diagnostic;
\item
  the assumptions made by the representation;
\item
  one possible failure mode.
\end{enumerate}

This checklist turns notation into understanding. The goal is not to
memorize symbols, but to recognize the mathematical structure behind NLP
systems.

\bookmarksetup{startatroot}

\chapter{Mathematical Background}\label{mathematical-background}

Core tools for studying natural language processing

\hfill\break

Natural language processing uses many kinds of mathematics at once. A
tokenizer creates finite sets. A document becomes a vector. A language
model becomes a conditional probability distribution. A transformer
becomes a composition of matrix maps, nonlinearities, normalizations,
and optimization steps. This page collects the main background ideas
used throughout the book.

This background chapter is not meant to replace a full course in linear
algebra, probability, statistics, optimization, graph theory, or
topology. Instead, it gives a compact map of the mathematical objects
that appear repeatedly in NLP. Students can read this page before
starting the book, and then return to the relevant sections while
working through the chapters and labs.

\section{How to use this page}\label{how-to-use-this-page}

Use this page as a reference while reading the book.

{\def\LTcaptype{none} % do not increment counter
\begin{longtable}[]{@{}
  >{\raggedright\arraybackslash}p{(\linewidth - 2\tabcolsep) * \real{0.5000}}
  >{\raggedright\arraybackslash}p{(\linewidth - 2\tabcolsep) * \real{0.5000}}@{}}
\toprule\noalign{}
\begin{minipage}[b]{\linewidth}\raggedright
When you see this in NLP
\end{minipage} & \begin{minipage}[b]{\linewidth}\raggedright
Review this mathematical topic
\end{minipage} \\
\midrule\noalign{}
\endhead
\bottomrule\noalign{}
\endlastfoot
token ids, vocabularies, text normalization & finite sets and maps \\
count vectors, TF-IDF, embeddings & vectors, norms, inner products,
matrices \\
cosine similarity, nearest neighbors & geometry and metrics \\
document-term matrices, latent semantic analysis & matrix factorization
and SVD \\
n-gram models, Naive Bayes, language models & probability and
likelihood \\
cross-entropy, KL divergence, perplexity & information theory \\
neural networks and transformers & optimization and matrix calculus \\
attention & softmax, weighted averages, bilinear scores \\
retrieval and RAG & similarity search, ranking, evaluation metrics \\
dependency graphs, knowledge graphs & graph theory and linear algebra on
graphs \\
representation diagnostics & geometry, topology, statistics, and
fairness audits \\
\end{longtable}
}

A useful study habit is to ask four questions whenever a new NLP method
appears.

\begin{enumerate}
\def\labelenumi{\arabic{enumi}.}
\tightlist
\item
  What is the input set?
\item
  What mathematical representation is assigned to each input?
\item
  What function, probability distribution, or optimization problem is
  being used?
\item
  What metric or loss tells us whether the method worked?
\end{enumerate}

\section{1. Finite sets, maps, and
indexing}\label{finite-sets-maps-and-indexing}

Most NLP systems begin by converting text into discrete symbols.

A \textbf{vocabulary} is a finite set

\[
V = \{w_1,w_2,\ldots,w_m\}.
\]

An \textbf{indexing map} assigns each token to an integer:

\[
\iota: V \to \{0,1,\ldots,m-1\}.
\]

After indexing, a sentence such as

\begin{Shaded}
\begin{Highlighting}[]
\NormalTok{language models learn patterns}
\end{Highlighting}
\end{Shaded}

may become a sequence of token ids such as

\begin{Shaded}
\begin{Highlighting}[]
\NormalTok{[12, 7, 35, 19]}
\end{Highlighting}
\end{Shaded}

This is mathematically important because neural networks do not operate
directly on English words. They operate on arrays of numbers.

\subsection{Key ideas}\label{key-ideas}

{\def\LTcaptype{none} % do not increment counter
\begin{longtable}[]{@{}
  >{\raggedright\arraybackslash}p{(\linewidth - 2\tabcolsep) * \real{0.5000}}
  >{\raggedright\arraybackslash}p{(\linewidth - 2\tabcolsep) * \real{0.5000}}@{}}
\toprule\noalign{}
\begin{minipage}[b]{\linewidth}\raggedright
Concept
\end{minipage} & \begin{minipage}[b]{\linewidth}\raggedright
Meaning in NLP
\end{minipage} \\
\midrule\noalign{}
\endhead
\bottomrule\noalign{}
\endlastfoot
finite set & vocabulary, label set, document collection \\
sequence & ordered list of tokens \\
map & tokenizer, embedding lookup, classifier \\
equivalence relation & normalization that treats variants as the same
token \\
composition & pipeline: text \(\to\) tokens \(\to\) vectors \(\to\)
prediction \\
\end{longtable}
}

\subsection{Common notation}\label{common-notation}

{\def\LTcaptype{none} % do not increment counter
\begin{longtable}[]{@{}ll@{}}
\toprule\noalign{}
Symbol & Meaning \\
\midrule\noalign{}
\endhead
\bottomrule\noalign{}
\endlastfoot
\(V\) & vocabulary \\
\(m=|V|\) & vocabulary size \\
\(x=(x_1,\ldots,x_T)\) & token sequence of length \(T\) \\
\(y\) & label or target output \\
\(D\) & document collection or dataset \\
\end{longtable}
}

\section{2. Vectors and matrices}\label{vectors-and-matrices}

Once words and documents are indexed, they are commonly represented as
vectors.

A vector in \(\mathbb{R}^d\) is an ordered list of real numbers:

\[
\mathbf{x} = (x_1,x_2,\ldots,x_d).
\]

A document-term vector might record how often each word appears in a
document. If the vocabulary is

\[
V=\{\text{data},\text{language},\text{model},\text{probability}\},
\]

then the document

\begin{Shaded}
\begin{Highlighting}[]
\NormalTok{language model language}
\end{Highlighting}
\end{Shaded}

has count vector

\[
\mathbf{x}=(0,2,1,0).
\]

A collection of \(n\) documents can be represented by a document-term
matrix

\[
X \in \mathbb{R}^{n \times m},
\]

where row \(i\) represents document \(i\) and column \(j\) represents
token \(w_j\).

\subsection{Inner products and norms}\label{inner-products-and-norms}

The dot product of two vectors is

\[
\mathbf{x}\cdot \mathbf{y}=\sum_{j=1}^d x_jy_j.
\]

The Euclidean norm is

\[
\|\mathbf{x}\|_2=\sqrt{\sum_{j=1}^d x_j^2}.
\]

The cosine similarity is

\[
\cos(\mathbf{x},\mathbf{y})=
\frac{\mathbf{x}\cdot \mathbf{y}}{\|\mathbf{x}\|_2\|\mathbf{y}\|_2}.
\]

Cosine similarity is widely used in NLP retrieval because it compares
vector direction rather than raw magnitude. This is useful when two
documents have similar topic proportions but different lengths.

\subsection{Important matrix
operations}\label{important-matrix-operations}

{\def\LTcaptype{none} % do not increment counter
\begin{longtable}[]{@{}
  >{\raggedright\arraybackslash}p{(\linewidth - 4\tabcolsep) * \real{0.3333}}
  >{\raggedright\arraybackslash}p{(\linewidth - 4\tabcolsep) * \real{0.3333}}
  >{\raggedright\arraybackslash}p{(\linewidth - 4\tabcolsep) * \real{0.3333}}@{}}
\toprule\noalign{}
\begin{minipage}[b]{\linewidth}\raggedright
Operation
\end{minipage} & \begin{minipage}[b]{\linewidth}\raggedright
Formula
\end{minipage} & \begin{minipage}[b]{\linewidth}\raggedright
NLP interpretation
\end{minipage} \\
\midrule\noalign{}
\endhead
\bottomrule\noalign{}
\endlastfoot
matrix-vector product & \(A\mathbf{x}\) & linear transformation of a
representation \\
matrix-matrix product & \(AB\) & composition of many linear maps \\
transpose & \(A^T\) & switch rows and columns \\
Gram matrix & \(XX^T\) & pairwise document similarities \\
covariance-like matrix & \(X^TX\) & feature-feature relationships \\
diagonal scaling & \(XD\) & TF-IDF reweighting \\
\end{longtable}
}

\section{3. Linear transformations and
embeddings}\label{linear-transformations-and-embeddings}

An \textbf{embedding} maps a discrete object into a vector space. For
example, a word embedding is a function

\[
e: V \to \mathbb{R}^d.
\]

If the vocabulary has \(m\) words, an embedding table is a matrix

\[
E \in \mathbb{R}^{m\times d}.
\]

The row \(E_{j,:}\) is the vector representation of token \(w_j\).

The embedding lookup for token id \(j\) can be written as

\[
\mathbf{e}_j = E_{j,:}.
\]

Equivalently, if \(\mathbf{u}_j\) is the one-hot vector for token \(j\),
then

\[
\mathbf{e}_j = \mathbf{u}_j^T E.
\]

This identity explains why embedding lookup is a computationally
efficient version of multiplying by a large matrix.

\subsection{What embeddings try to
preserve}\label{what-embeddings-try-to-preserve}

Embeddings are useful when geometric relationships encode linguistic
relationships.

{\def\LTcaptype{none} % do not increment counter
\begin{longtable}[]{@{}ll@{}}
\toprule\noalign{}
Geometric relation & Linguistic interpretation \\
\midrule\noalign{}
\endhead
\bottomrule\noalign{}
\endlastfoot
small distance & similar meaning or usage \\
large dot product & high compatibility or relevance \\
common direction & shared semantic or syntactic change \\
cluster & group of related words or documents \\
subspace & latent topic, style, or feature \\
\end{longtable}
}

Embeddings are not perfect. They may encode frequency artifacts, social
bias, domain shift, and spurious correlations. This is why the book
repeatedly connects representation learning with diagnostics and
evaluation.

\section{4. Projection, orthogonality, and least
squares}\label{projection-orthogonality-and-least-squares}

A projection finds the closest point in a subspace. If \(U\) is a
subspace of \(\mathbb{R}^d\), the projection of \(\mathbf{x}\) onto
\(U\) is the vector \(\mathbf{p}\in U\) that minimizes

\[
\|\mathbf{x}-\mathbf{p}\|_2.
\]

Projection appears in NLP when we reduce dimension, remove a direction
from embeddings, or fit linear models.

If \(A\in \mathbb{R}^{n\times d}\) and \(\mathbf{b}\in \mathbb{R}^n\),
the least-squares problem is

\[
\min_{\mathbf{w}\in \mathbb{R}^d}\|A\mathbf{w}-\mathbf{b}\|_2^2.
\]

The normal equation is

\[
A^TA\mathbf{w}=A^T\mathbf{b}.
\]

In modern machine learning, we usually solve such problems with
numerical optimization rather than explicitly forming the inverse.

\section{5. Singular value
decomposition}\label{singular-value-decomposition}

The singular value decomposition, or SVD, is one of the most important
matrix tools in this book. For a matrix \(X\in \mathbb{R}^{n\times m}\),
the SVD has the form

\[
X = U\Sigma V^T.
\]

Here:

\begin{itemize}
\tightlist
\item
  columns of \(U\) describe directions among documents;
\item
  columns of \(V\) describe directions among terms or features;
\item
  diagonal entries of \(\Sigma\) are singular values, ordered from
  largest to smallest.
\end{itemize}

A rank-\(k\) approximation is

\[
X_k = U_k\Sigma_k V_k^T.
\]

This is the best rank-\(k\) approximation to \(X\) in Frobenius norm:

\[
X_k = \arg\min_{\operatorname{rank}(Y)\leq k}\|X-Y\|_F.
\]

In NLP, SVD appears in latent semantic analysis, dimensionality
reduction, visualization, denoising, and embedding diagnostics.

\section{6. Probability distributions}\label{probability-distributions}

A probability distribution on a finite set \(V\) assigns a number
\(p(w)\) to each \(w\in V\) such that

\[
p(w)\geq 0,\qquad \sum_{w\in V}p(w)=1.
\]

In NLP, finite distributions appear everywhere:

\begin{itemize}
\tightlist
\item
  token frequency distributions;
\item
  next-token distributions;
\item
  class posterior probabilities;
\item
  attention weights;
\item
  retrieval scores after normalization.
\end{itemize}

\subsection{Conditional probability}\label{conditional-probability}

The conditional probability of \(y\) given \(x\) is

\[
P(y\mid x)=\frac{P(x,y)}{P(x)}.
\]

A classifier estimates \(P(y\mid x)\). A language model estimates

\[
P(x_t\mid x_1,\ldots,x_{t-1}).
\]

The chain rule gives

\[
P(x_1,\ldots,x_T)=\prod_{t=1}^T P(x_t\mid x_1,
\ldots,x_{t-1}).
\]

This identity is the mathematical basis of autoregressive language
modeling.

\section{7. Likelihood, log-likelihood, and maximum
likelihood}\label{likelihood-log-likelihood-and-maximum-likelihood}

Suppose a model with parameter \(\theta\) assigns probability
\(p_\theta(x_i)\) to each training example \(x_i\). The likelihood is

\[
L(\theta)=\prod_{i=1}^n p_\theta(x_i).
\]

The log-likelihood is

\[
\ell(\theta)=\sum_{i=1}^n \log p_\theta(x_i).
\]

Maximum likelihood estimation chooses

\[
\hat\theta=\arg\max_\theta \ell(\theta).
\]

Equivalently, many machine learning algorithms minimize the negative
log-likelihood:

\[
\hat\theta=\arg\min_\theta -\sum_{i=1}^n \log p_\theta(x_i).
\]

In language modeling, this becomes the token-level cross-entropy loss.

\section{8. Entropy, cross-entropy, KL divergence, and
perplexity}\label{entropy-cross-entropy-kl-divergence-and-perplexity}

Information theory provides the language for uncertainty and prediction
quality.

The entropy of a distribution \(p\) on \(V\) is

\[
H(p)=-\sum_{w\in V}p(w)\log p(w).
\]

Entropy measures uncertainty. A concentrated distribution has low
entropy. A nearly uniform distribution has high entropy.

The cross-entropy from \(p\) to \(q\) is

\[
H(p,q)=-\sum_{w\in V}p(w)\log q(w).
\]

The KL divergence from \(p\) to \(q\) is

\[
D_{\mathrm{KL}}(p\|q)=\sum_{w\in V}p(w)\log\frac{p(w)}{q(w)}.
\]

The identity

\[
H(p,q)=H(p)+D_{\mathrm{KL}}(p\|q)
\]

shows that cross-entropy is minimized when \(q=p\).

If average token cross-entropy is measured using natural logarithms,
perplexity is

\[
\operatorname{PPL}=\exp(H).
\]

Perplexity can be interpreted as an effective number of choices per
token.

\section{9. Bayes rule and
classification}\label{bayes-rule-and-classification}

Bayes rule states

\[
P(y\mid x)=\frac{P(x\mid y)P(y)}{P(x)}.
\]

For classification, the denominator \(P(x)\) is the same for all
classes, so the Bayes decision rule is

\[
\hat y = \arg\max_y P(x\mid y)P(y).
\]

Naive Bayes assumes that features are conditionally independent given
the class. For a count vector

\[
\mathbf{x}=(x_1,\ldots,x_m),
\]

the multinomial Naive Bayes score has the form

\[
\log P(y)+\sum_{j=1}^m x_j\log P(w_j\mid y).
\]

This formula is simple, interpretable, and surprisingly strong for many
text classification problems.

\section{10. Softmax and attention
weights}\label{softmax-and-attention-weights}

The softmax function converts real-valued scores into probabilities.
Given scores

\[
\mathbf{z}=(z_1,\ldots,z_m),
\]

the softmax distribution is

\[
\operatorname{softmax}(\mathbf{z})_j=
\frac{e^{z_j}}{\sum_{k=1}^m e^{z_k}}.
\]

For numerical stability, compute softmax as

\[
\operatorname{softmax}(\mathbf{z})_j=
\frac{e^{z_j-c}}{\sum_{k=1}^m e^{z_k-c}},
\qquad c=\max_k z_k.
\]

In attention, the score between query \(\mathbf{q}\) and key
\(\mathbf{k}\) is often

\[
\frac{\mathbf{q}\cdot \mathbf{k}}{\sqrt{d}}.
\]

The attention output is a weighted average of value vectors:

\[
\operatorname{Attention}(Q,K,V)=
\operatorname{softmax}\left(\frac{QK^T}{\sqrt{d}}\right)V.
\]

This formula is one of the central mathematical expressions in modern
NLP.

\section{11. Differentiation, gradients, and
backpropagation}\label{differentiation-gradients-and-backpropagation}

Training a neural model means minimizing a loss function. If

\[
L(\theta)
\]

is a loss depending on parameters \(\theta\), the gradient

\[
\nabla_\theta L
\]

points in the direction of steepest increase. Gradient descent updates
parameters by

\[
\theta_{t+1}=\theta_t-\eta\nabla_\theta L(\theta_t),
\]

where \(\eta>0\) is the learning rate.

Backpropagation is repeated use of the chain rule. If

\[
z=f(y),\qquad y=g(x),
\]

then

\[
\frac{dz}{dx}=\frac{dz}{dy}\frac{dy}{dx}.
\]

In a neural network, every layer contributes one part of this chain.

\subsection{Common losses}\label{common-losses}

{\def\LTcaptype{none} % do not increment counter
\begin{longtable}[]{@{}
  >{\raggedright\arraybackslash}p{(\linewidth - 4\tabcolsep) * \real{0.3333}}
  >{\raggedright\arraybackslash}p{(\linewidth - 4\tabcolsep) * \real{0.3333}}
  >{\raggedright\arraybackslash}p{(\linewidth - 4\tabcolsep) * \real{0.3333}}@{}}
\toprule\noalign{}
\begin{minipage}[b]{\linewidth}\raggedright
Task
\end{minipage} & \begin{minipage}[b]{\linewidth}\raggedright
Common loss
\end{minipage} & \begin{minipage}[b]{\linewidth}\raggedright
Mathematical form
\end{minipage} \\
\midrule\noalign{}
\endhead
\bottomrule\noalign{}
\endlastfoot
binary classification & logistic loss & \(-y\log p-(1-y)\log(1-p)\) \\
multiclass classification & cross-entropy & \(-\log p_y\) \\
language modeling & token cross-entropy &
\(-\sum_t \log p(x_t\mid x_{<t})\) \\
regression & squared error & \(\|y-\hat y\|_2^2\) \\
preference learning & pairwise logistic loss &
\(-\log\sigma(s^+-s^-)\) \\
\end{longtable}
}

\section{12. Optimization algorithms}\label{optimization-algorithms}

Full-batch gradient descent uses all training data for each update.
Stochastic gradient descent uses small random batches. For mini-batch
\(B_t\), the update is

\[
\theta_{t+1}=\theta_t-\eta\frac{1}{|B_t|}\sum_{i\in B_t}\nabla_\theta L_i(\theta_t).
\]

Momentum smooths updates:

\[
v_{t+1}=\beta v_t + \nabla L(\theta_t),
\]

\[
\theta_{t+1}=\theta_t-\eta v_{t+1}.
\]

Adam and AdamW adapt learning rates using running estimates of first and
second moments of gradients. These methods are standard in training
transformer models.

\subsection{Practical optimization
issues}\label{practical-optimization-issues}

{\def\LTcaptype{none} % do not increment counter
\begin{longtable}[]{@{}
  >{\raggedright\arraybackslash}p{(\linewidth - 4\tabcolsep) * \real{0.3333}}
  >{\raggedright\arraybackslash}p{(\linewidth - 4\tabcolsep) * \real{0.3333}}
  >{\raggedright\arraybackslash}p{(\linewidth - 4\tabcolsep) * \real{0.3333}}@{}}
\toprule\noalign{}
\begin{minipage}[b]{\linewidth}\raggedright
Issue
\end{minipage} & \begin{minipage}[b]{\linewidth}\raggedright
Meaning
\end{minipage} & \begin{minipage}[b]{\linewidth}\raggedright
Common response
\end{minipage} \\
\midrule\noalign{}
\endhead
\bottomrule\noalign{}
\endlastfoot
learning rate too large & loss diverges & reduce learning rate \\
learning rate too small & training is slow & increase learning rate or
use schedule \\
exploding gradients & updates become unstable & gradient clipping \\
overfitting & train loss improves but validation worsens &
regularization, early stopping, more data \\
padding tokens affect loss & model learns from fake positions & loss
masks \\
class imbalance & rare classes are ignored & reweighting or threshold
adjustment \\
\end{longtable}
}

\section{13. Sequences and recurrence}\label{sequences-and-recurrence}

A text is not only a bag of words. It is an ordered sequence

\[
x_1,x_2,\ldots,x_T.
\]

A recurrent neural network updates a hidden state by

\[
\mathbf{h}_t=\phi(W_x\mathbf{x}_t+W_h\mathbf{h}_{t-1}+\mathbf{b}),
\]

where \(\phi\) is a nonlinear activation function.

The hidden state is a summary of previous tokens. The challenge is that
long chains of matrix multiplications can produce vanishing or exploding
gradients. LSTMs and GRUs address this issue with gates that control
memory flow.

\section{14. Transformers and residual
computation}\label{transformers-and-residual-computation}

A transformer block combines several operations:

\begin{enumerate}
\def\labelenumi{\arabic{enumi}.}
\tightlist
\item
  self-attention;
\item
  residual connection;
\item
  layer normalization;
\item
  feedforward network;
\item
  another residual connection;
\item
  another layer normalization.
\end{enumerate}

A simplified encoder block can be written as

\[
Y = \operatorname{LayerNorm}(X + \operatorname{SelfAttention}(X)),
\]

\[
Z = \operatorname{LayerNorm}(Y + \operatorname{FFN}(Y)).
\]

The residual terms help preserve information and improve gradient flow.
Layer normalization stabilizes the scale of activations.

\section{15. Positional encoding}\label{positional-encoding}

Self-attention alone is permutation equivariant: if token order is
permuted, outputs are permuted in the same way. To model word order,
transformers add position information.

A token representation often has the form

\[
\mathbf{z}_t = \mathbf{e}_{x_t}+\mathbf{p}_t,
\]

where \(\mathbf{e}_{x_t}\) is the token embedding and \(\mathbf{p}_t\)
is a position vector.

Sinusoidal positional encodings use functions such as

\[
\sin\left(\frac{t}{10000^{2j/d}}\right),
\qquad
\cos\left(\frac{t}{10000^{2j/d}}\right).
\]

Relative position methods modify attention scores according to token
distance. Rotary positional embedding, or RoPE, rotates query and key
vectors by position-dependent angles.

\section{16. Graphs and language}\label{graphs-and-language}

A graph consists of vertices and edges:

\[
G=(V,E).
\]

Graphs appear in NLP in many ways:

\begin{itemize}
\tightlist
\item
  dependency parse graphs;
\item
  word co-occurrence graphs;
\item
  document similarity graphs;
\item
  knowledge graphs;
\item
  retrieval citation graphs;
\item
  conversation graphs.
\end{itemize}

The adjacency matrix \(A\) is defined by

\[
A_{ij}=1
\]

if there is an edge from node \(i\) to node \(j\), and \(0\) otherwise.
Weighted graphs allow nonnegative edge weights.

The degree matrix \(D\) has diagonal entries

\[
D_{ii}=\sum_j A_{ij}.
\]

The graph Laplacian is

\[
L=D-A.
\]

Graph neural networks update node features by mixing information from
neighbors. A simple graph convolution has the form

\[
H^{(\ell+1)}=\sigma(\widetilde A H^{(\ell)}W^{(\ell)}),
\]

where \(\widetilde A\) is a normalized adjacency matrix.

\section{17. Metrics, ranking, and retrieval
evaluation}\label{metrics-ranking-and-retrieval-evaluation}

Retrieval systems rank documents according to a score. If \(q\) is a
query and \(d\) is a document, a retrieval score might be

\[
s(q,d)=\cos(\mathbf{q},\mathbf{d}).
\]

The output is an ordered list of documents. Evaluation must respect
ranking.

{\def\LTcaptype{none} % do not increment counter
\begin{longtable}[]{@{}
  >{\raggedright\arraybackslash}p{(\linewidth - 2\tabcolsep) * \real{0.5000}}
  >{\raggedright\arraybackslash}p{(\linewidth - 2\tabcolsep) * \real{0.5000}}@{}}
\toprule\noalign{}
\begin{minipage}[b]{\linewidth}\raggedright
Metric
\end{minipage} & \begin{minipage}[b]{\linewidth}\raggedright
What it measures
\end{minipage} \\
\midrule\noalign{}
\endhead
\bottomrule\noalign{}
\endlastfoot
precision@k & fraction of top-\(k\) documents that are relevant \\
recall@k & fraction of all relevant documents found in top \(k\) \\
MRR & reciprocal rank of the first relevant document \\
MAP & average precision across relevant documents \\
nDCG & relevance-weighted ranking quality with discount for lower
ranks \\
\end{longtable}
}

For retrieval-augmented generation, evaluation must check both retrieval
and generation. A good answer may fail because the retriever missed the
right evidence, or because the generator ignored the evidence.

\section{18. Statistics and
uncertainty}\label{statistics-and-uncertainty}

Evaluation metrics are estimates computed on finite datasets. They have
uncertainty.

If a model has accuracy \(\hat p\) on \(n\) independent examples, an
approximate standard error is

\[
\sqrt{\frac{\hat p(1-\hat p)}{n}}.
\]

For complex metrics, bootstrap resampling is often more useful. A
bootstrap procedure repeatedly samples the evaluation set with
replacement, recomputes the metric, and uses the resulting empirical
distribution to estimate uncertainty.

\subsection{Confusion matrix}\label{confusion-matrix}

For binary classification, the confusion matrix records:

{\def\LTcaptype{none} % do not increment counter
\begin{longtable}[]{@{}lrr@{}}
\toprule\noalign{}
& predicted positive & predicted negative \\
\midrule\noalign{}
\endhead
\bottomrule\noalign{}
\endlastfoot
actual positive & TP & FN \\
actual negative & FP & TN \\
\end{longtable}
}

The most common metrics are

\[
\operatorname{precision}=\frac{TP}{TP+FP},
\]

\[
\operatorname{recall}=\frac{TP}{TP+FN},
\]

\[
F_1=\frac{2\cdot \operatorname{precision}\cdot \operatorname{recall}}
{\operatorname{precision}+\operatorname{recall}}.
\]

\section{19. Geometry and topology of
representations}\label{geometry-and-topology-of-representations}

A set of embeddings can be treated as a finite metric space. Given
points

\[
\mathbf{x}_1,\ldots,\mathbf{x}_n\in \mathbb{R}^d,
\]

we can compute pairwise distances

\[
d_{ij}=\|\mathbf{x}_i-\mathbf{x}_j\|_2.
\]

A neighborhood graph connects points whose distance is below a scale
\(\epsilon\). As \(\epsilon\) changes, the graph may gain or lose
connected components.

Topology studies properties such as connected components and loops.
Persistent homology studies how such properties change across scale.

In NLP, these ideas can help diagnose representation spaces. For
example:

\begin{itemize}
\tightlist
\item
  Are embeddings organized into clusters?
\item
  Are there hubs that appear near too many points?
\item
  Does the representation collapse into a narrow cone?
\item
  Are trajectories through generated text smooth or unstable?
\item
  Do different groups or domains occupy different regions?
\end{itemize}

\section{20. Responsible mathematical
modeling}\label{responsible-mathematical-modeling}

Responsible NLP is not separate from mathematics. It uses mathematical
tools to ask whether systems behave reliably and fairly.

Examples include:

{\def\LTcaptype{none} % do not increment counter
\begin{longtable}[]{@{}
  >{\raggedright\arraybackslash}p{(\linewidth - 2\tabcolsep) * \real{0.5000}}
  >{\raggedright\arraybackslash}p{(\linewidth - 2\tabcolsep) * \real{0.5000}}@{}}
\toprule\noalign{}
\begin{minipage}[b]{\linewidth}\raggedright
Concern
\end{minipage} & \begin{minipage}[b]{\linewidth}\raggedright
Mathematical diagnostic
\end{minipage} \\
\midrule\noalign{}
\endhead
\bottomrule\noalign{}
\endlastfoot
unequal error rates & subgroup precision, recall, false positive rate \\
overconfidence & calibration curves, ECE, Brier score \\
distribution shift & train-test divergence, importance weighting \\
embedding bias & association scores, projection diagnostics \\
hallucination in RAG & evidence coverage and answer support \\
privacy leakage & randomized response, membership inference risk \\
unreliable automation & abstention, uncertainty thresholds, human
review \\
\end{longtable}
}

A model is not responsible simply because it has high average accuracy.
The errors must be understood by context, subgroup, uncertainty, and
consequences.

\section{Chapter-by-chapter background
map}\label{chapter-by-chapter-background-map}

{\def\LTcaptype{none} % do not increment counter
\begin{longtable}[]{@{}
  >{\raggedright\arraybackslash}p{(\linewidth - 2\tabcolsep) * \real{0.5000}}
  >{\raggedright\arraybackslash}p{(\linewidth - 2\tabcolsep) * \real{0.5000}}@{}}
\toprule\noalign{}
\begin{minipage}[b]{\linewidth}\raggedright
Chapters
\end{minipage} & \begin{minipage}[b]{\linewidth}\raggedright
Main mathematical tools
\end{minipage} \\
\midrule\noalign{}
\endhead
\bottomrule\noalign{}
\endlastfoot
1--3 & finite sets, token sequences, vectors, document-term matrices,
cosine similarity \\
4--6 & probability distributions, Bayes rule, likelihood, entropy,
smoothing \\
7--9 & vector spaces, TF-IDF scaling, Gram matrices, SVD, embedding
geometry \\
10--12 & softmax, cross-entropy, gradients, recurrence, gated memory \\
13--16 & attention matrices, masking, positional encoding, residual
maps, normalization \\
17--20 & chain rule for probabilities, perplexity, retrieval metrics,
calibration, bootstrap uncertainty \\
21--25 & information theory, optimization, graph Laplacians, topology,
fairness and responsibility \\
\end{longtable}
}

\section{Minimal Python notation}\label{minimal-python-notation}

The labs mostly use NumPy notation. The table below connects
mathematical symbols to Python objects.

{\def\LTcaptype{none} % do not increment counter
\begin{longtable}[]{@{}
  >{\raggedright\arraybackslash}p{(\linewidth - 2\tabcolsep) * \real{0.5000}}
  >{\raggedright\arraybackslash}p{(\linewidth - 2\tabcolsep) * \real{0.5000}}@{}}
\toprule\noalign{}
\begin{minipage}[b]{\linewidth}\raggedright
Mathematical object
\end{minipage} & \begin{minipage}[b]{\linewidth}\raggedright
Python representation
\end{minipage} \\
\midrule\noalign{}
\endhead
\bottomrule\noalign{}
\endlastfoot
scalar \(a\) & \texttt{float} \\
vector \(\mathbf{x}\in\mathbb{R}^d\) & one-dimensional NumPy array with
shape \texttt{(d,)} \\
matrix \(X\in\mathbb{R}^{n\times d}\) & two-dimensional NumPy array with
shape \texttt{(n,\ d)} \\
token sequence \((x_1,\ldots,x_T)\) & Python list of strings or ids \\
vocabulary map & Python dictionary \\
probability vector & nonnegative array summing to one \\
dataset & list, pandas DataFrame, or pair of arrays \\
model parameters & arrays to be optimized \\
\end{longtable}
}

Common NumPy operations:

\begin{Shaded}
\begin{Highlighting}[]
\NormalTok{x }\OperatorTok{@}\NormalTok{ y          }\CommentTok{\# dot product for vectors}
\NormalTok{X }\OperatorTok{@}\NormalTok{ W          }\CommentTok{\# matrix multiplication}
\NormalTok{X.T            }\CommentTok{\# transpose}
\NormalTok{np.linalg.norm(x)}
\NormalTok{np.exp(z)}
\NormalTok{np.log(p)}
\NormalTok{np.argmax(scores)}
\NormalTok{np.linalg.svd(X)}
\end{Highlighting}
\end{Shaded}

\section{A compact formula sheet}\label{a-compact-formula-sheet}

{\def\LTcaptype{none} % do not increment counter
\begin{longtable}[]{@{}
  >{\raggedright\arraybackslash}p{(\linewidth - 2\tabcolsep) * \real{0.5000}}
  >{\raggedright\arraybackslash}p{(\linewidth - 2\tabcolsep) * \real{0.5000}}@{}}
\toprule\noalign{}
\begin{minipage}[b]{\linewidth}\raggedright
Topic
\end{minipage} & \begin{minipage}[b]{\linewidth}\raggedright
Formula
\end{minipage} \\
\midrule\noalign{}
\endhead
\bottomrule\noalign{}
\endlastfoot
cosine similarity &
\(\cos(\mathbf{x},\mathbf{y})=\frac{\mathbf{x}\cdot\mathbf{y}}{\|\mathbf{x}\|\|\mathbf{y}\|}\) \\
softmax & \(p_j=\frac{e^{z_j}}{\sum_k e^{z_k}}\) \\
cross-entropy & \(H(p,q)=-\sum_j p_j\log q_j\) \\
KL divergence &
\(D_{\mathrm{KL}}(p\|q)=\sum_jp_j\log\frac{p_j}{q_j}\) \\
perplexity & \(\operatorname{PPL}=\exp(H)\) \\
Bayes rule & \(P(y\mid x)=\frac{P(x\mid y)P(y)}{P(x)}\) \\
gradient descent & \(\theta_{t+1}=\theta_t-\eta\nabla L(\theta_t)\) \\
attention & \(\operatorname{softmax}(QK^T/\sqrt{d})V\) \\
graph Laplacian & \(L=D-A\) \\
rank-\(k\) SVD approximation & \(X_k=U_k\Sigma_kV_k^T\) \\
\end{longtable}
}

\section{Self-check questions}\label{self-check-questions}

\begin{enumerate}
\def\labelenumi{\arabic{enumi}.}
\tightlist
\item
  Why does an embedding table have one row per vocabulary item?
\item
  Why is cosine similarity often preferred over raw dot product for
  document retrieval?
\item
  What is the difference between entropy and cross-entropy?
\item
  Why do we usually optimize log-likelihood instead of likelihood?
\item
  What does softmax do to a vector of real-valued scores?
\item
  Why does a transformer need positional information?
\item
  What is the difference between model training loss and evaluation
  metric?
\item
  Why can a high average score hide subgroup failures?
\item
  How can a graph represent language data?
\item
  What does it mean to study embeddings as a finite metric space?
\end{enumerate}

\section{Suggested review path}\label{suggested-review-path}

Students who need a quick review before the book can follow this
sequence.

\begin{enumerate}
\def\labelenumi{\arabic{enumi}.}
\tightlist
\item
  Review finite sets, maps, vectors, matrices, and cosine similarity.
\item
  Work through probability distributions, conditional probability, and
  Bayes rule.
\item
  Study entropy, cross-entropy, KL divergence, and perplexity.
\item
  Review gradients, softmax, and cross-entropy loss.
\item
  Learn the attention formula and identify the shapes of \(Q\), \(K\),
  and \(V\).
\item
  Return to graph and topology sections later, when reaching Chapters 23
  and 24.
\end{enumerate}

The goal is not to memorize every formula at once. The goal is to
recognize the mathematical structure behind each NLP method and to
connect the theory with the Python labs.

\part{Part I. Language as Data}

\chapter{What Is Natural Language
Processing?}\label{what-is-natural-language-processing}

Language tasks as mathematical maps

\hfill\break

Natural language processing studies how to represent, analyze, predict,
and generate human language using mathematical and computational models.
This chapter introduces NLP from a clean mathematical point of view. The
goal is not to begin with software tools, but to identify the
mathematical objects behind those tools: sets, sequences, functions,
probability distributions, vector spaces, loss functions, and
estimators.

\section{Learning goals}\label{learning-goals}

After studying this chapter, you should be able to:

\begin{itemize}
\tightlist
\item
  describe text data as finite sequences over a vocabulary;
\item
  distinguish representation, model, loss, training, and evaluation;
\item
  express common NLP tasks as mathematical maps;
\item
  explain the difference between rule-based, statistical, and neural
  approaches;
\item
  identify the main mathematical objects used throughout the book.
\end{itemize}

\section{1.1 Why NLP needs mathematics}\label{why-nlp-needs-mathematics}

Human language is not naturally numerical. A sentence such as

\begin{quote}
The model predicts the next token.
\end{quote}

contains words, order, meaning, context, and grammar. A computer model
cannot directly operate on this sentence until we choose a mathematical
representation. This choice is the first modeling decision in NLP.

At the most basic level, NLP converts language into objects that can be
studied using mathematics. For example:

{\def\LTcaptype{none} % do not increment counter
\begin{longtable}[]{@{}ll@{}}
\toprule\noalign{}
Language object & Mathematical object \\
\midrule\noalign{}
\endhead
\bottomrule\noalign{}
\endlastfoot
word & element of a finite set \\
sentence & finite sequence \\
document & sequence, vector, or probability distribution \\
corpus & finite sample of documents \\
meaning & representation in a vector space or probability model \\
classification & function from texts to labels \\
generation & sampling from a conditional distribution \\
translation & map between sequence spaces \\
retrieval & nearest-neighbor problem in a metric space \\
\end{longtable}
}

The point of this book is that modern NLP is not a collection of
disconnected algorithms. It is a sequence of mathematical constructions.

\begin{tcolorbox}[enhanced jigsaw, arc=.35mm, bottomrule=.15mm, bottomtitle=1mm, breakable, colback=white, colbacktitle=quarto-callout-note-color!10!white, colframe=quarto-callout-note-color-frame, coltitle=black, left=2mm, leftrule=.75mm, opacityback=0, opacitybacktitle=0.6, rightrule=.15mm, title=\textcolor{quarto-callout-note-color}{\faInfo}\hspace{0.5em}{Central viewpoint}, titlerule=0mm, toprule=.15mm, toptitle=1mm]

Before naming an NLP algorithm, identify the mathematical object:

\begin{itemize}
\tightlist
\item
  What is the input space?
\item
  What is the output space?
\item
  Is the model a function, a probability distribution, an estimator, or
  an optimization problem?
\item
  What loss function measures success?
\item
  What data are used to estimate the unknown parameters?
\end{itemize}

\end{tcolorbox}

For MA Applied Mathematics and MS Statistics students, this viewpoint
should feel natural. In statistics, we ask what is observed, what is
random, what is estimated, and what criterion defines a good estimate.
In applied mathematics, we ask what map is being constructed, what space
it acts on, and what structure is preserved or lost. NLP uses both
habits.

\section{1.2 Text as discrete structured
data}\label{text-as-discrete-structured-data}

Let \(V\) be a vocabulary, meaning a finite set of allowed tokens. A
token may be a word, character, subword unit, punctuation mark, or
special symbol. A sentence of length \(T\) is represented as a finite
sequence

\[
x = (w_1,w_2,\ldots,w_T), \qquad w_t \in V.
\]

The length \(T\) may vary from one sentence to another. If \(V\) has
size \(d\), then each token belongs to a set with \(d\) possible values.
The space of all token sequences of length \(T\) is

\[
V^T = \{(w_1,\ldots,w_T): w_t \in V \text{ for } t=1,\ldots,T\}.
\]

The space of all finite token sequences is

\[
V^* = \bigcup_{T \geq 0} V^T.
\]

Here \(V^*\) is not a vector space. It is a discrete sequence space.
This is one reason NLP is mathematically interesting: language begins as
symbolic, variable-length, discrete data, while most learning algorithms
require fixed-dimensional numerical inputs.

\begin{tcolorbox}[enhanced jigsaw, arc=.35mm, bottomrule=.15mm, bottomtitle=1mm, breakable, colback=white, colbacktitle=quarto-callout-important-color!10!white, colframe=quarto-callout-important-color-frame, coltitle=black, left=2mm, leftrule=.75mm, opacityback=0, opacitybacktitle=0.6, rightrule=.15mm, title=\textcolor{quarto-callout-important-color}{\faExclamation}\hspace{0.5em}{Important distinction}, titlerule=0mm, toprule=.15mm, toptitle=1mm]

A sentence is not automatically a vector. It becomes a vector only after
we choose a representation map.

\end{tcolorbox}

A representation map has the form

\[
\phi: V^* \longrightarrow R^d,
\]

where \(R^d\) is a numerical feature space. The map \(\phi\) might count
words, compute TF-IDF weights, average word embeddings, or use a
transformer to produce a contextual embedding. Different choices of
\(\phi\) preserve different information.

\section{1.3 Corpus, sample, and data-generating
process}\label{corpus-sample-and-data-generating-process}

A corpus is a finite collection of texts:

\[
D = \{x_1,x_2,\ldots,x_n\}, \qquad x_i \in V^*.
\]

In supervised NLP, each text may also have a label:

\[
D = \{(x_i,y_i)\}_{i=1}^n,
\]

where \(y_i\) may represent sentiment, topic, spam/not spam, author,
language, or another response variable.

From a statistical point of view, we often imagine that the observed
data are samples from an unknown distribution. For supervised learning,
one writes

\[
(X,Y) \sim P,
\]

where \(X\) is a random text and \(Y\) is a random label. The training
set is then idealized as

\[
(X_1,Y_1),\ldots,(X_n,Y_n) \overset{\text{iid}}{\sim} P.
\]

This iid assumption is useful but often imperfect. Real text corpora may
contain duplicated documents, time trends, domain shifts, author
dependence, and sampling bias. Nevertheless, the notation helps clarify
what machine learning is trying to do: construct a model from finite
data that performs well on future examples.

\subsection{Empirical distribution of a
corpus}\label{empirical-distribution-of-a-corpus}

Given a corpus \(D=\{x_1,\ldots,x_n\}\), the empirical distribution
assigns probability mass

\[
\widehat P_D(x)=\frac{1}{n}\sum_{i=1}^n 1\{x_i=x\}.
\]

In large language modeling, the data distribution is usually not
directly a distribution over whole documents only. It is also used to
produce many local prediction problems. For example, a sentence

\[
(w_1,w_2,\ldots,w_T)
\]

can be turned into training examples of the form

\[
(w_1,\ldots,w_{t-1}) \mapsto w_t.
\]

This is the next-token prediction view that will appear repeatedly in
later chapters.

\section{1.4 NLP tasks as mathematical
maps}\label{nlp-tasks-as-mathematical-maps}

A clean way to understand an NLP task is to write down its input space
and output space.

\subsection{Text classification}\label{text-classification}

In text classification, the input is a text sequence and the output is a
label. If \(C\) is a finite set of classes, then a classifier is a map

\[
f: V^* \longrightarrow C.
\]

For binary sentiment classification, one may use

\[
C = \{0,1\},
\]

where \(0\) means negative and \(1\) means positive. A probabilistic
classifier instead estimates

\[
p_\theta(y \mid x), \qquad y \in C.
\]

The predicted class is often

\[
\widehat y = \arg\max_{y\in C} p_\theta(y\mid x).
\]

Thus classification is either a function problem or a conditional
probability estimation problem.

\subsection{Sequence labeling}\label{sequence-labeling}

In sequence labeling, each input token receives a label. If

\[
x=(w_1,\ldots,w_T),
\]

then the output is another sequence

\[
y=(y_1,\ldots,y_T), \qquad y_t \in L,
\]

where \(L\) is a label set. Examples include part-of-speech tagging and
named entity recognition. The task is a map

\[
f: V^T \longrightarrow L^T.
\]

The key mathematical feature is that the output has the same length as
the input.

\subsection{Language modeling}\label{language-modeling}

A language model assigns probabilities to token sequences. For a
sentence \(x=(w_1,\ldots,w_T)\), a language model aims to estimate

\[
P(w_1,\ldots,w_T).
\]

Using the chain rule of probability,

\[
P(w_1,\ldots,w_T)
=
\prod_{t=1}^T P(w_t \mid w_1,\ldots,w_{t-1}).
\]

This identity is exact. A model becomes necessary because the true
conditional probabilities are unknown and the number of possible
histories is enormous.

\subsection{Machine translation}\label{machine-translation}

Machine translation maps a sequence in one language to a sequence in
another. Let \(V_S\) be the source vocabulary and \(V_T\) the target
vocabulary. A translation system may be written as

\[
f: V_S^* \longrightarrow V_T^*.
\]

A probabilistic translation model estimates

\[
p_\theta(y\mid x), \qquad x\in V_S^*,\quad y\in V_T^*.
\]

The predicted translation is often idealized as

\[
\widehat y = \arg\max_{y\in V_T^*} p_\theta(y\mid x).
\]

In practice this maximization is approximate because the output space is
huge.

\subsection{Information retrieval and semantic
search}\label{information-retrieval-and-semantic-search}

In retrieval, the input is a query \(q\) and a collection of documents
\(D=\{d_1,\ldots,d_n\}\). The goal is to rank documents by relevance. A
simple mathematical form is a scoring function

\[
s: V^* \times V^* \longrightarrow R,
\]

where \(s(q,d_i)\) measures how relevant document \(d_i\) is to query
\(q\). A retrieval system returns documents ordered by

\[
s(q,d_i).
\]

In vector search, one first embeds queries and documents into \(R^k\):

\[
\phi(q), \phi(d_i) \in R^k,
\]

and then ranks documents by inner product, cosine similarity, or
distance.

\section{1.5 Representation, model, loss, and
evaluation}\label{representation-model-loss-and-evaluation}

A common source of confusion is to mix together representation, model,
training, and evaluation. They are different mathematical components.

\subsection{Representation}\label{representation}

A representation is a map

\[
\phi: V^* \longrightarrow R^d.
\]

For example, the bag-of-words representation maps a document to a count
vector. If the vocabulary is

\[
V = \{v_1,\ldots,v_d\},
\]

then the \(j\)th component of \(\phi(x)\) may be

\[
\phi_j(x) = \#\{t: w_t=v_j\}.
\]

This representation ignores word order. The sentence

\begin{quote}
dogs chase cats
\end{quote}

and the sentence

\begin{quote}
cats chase dogs
\end{quote}

have the same bag-of-words vector if they use the same words once each.
The representation is simple and useful, but it loses syntactic
structure.

\subsection{Model}\label{model}

A model is a parameterized family of functions or distributions. For
classification, one may use

\[
f_\theta: R^d \longrightarrow C.
\]

Combined with a representation \(\phi\), the complete classifier is

\[
x \mapsto f_\theta(\phi(x)).
\]

A probabilistic classifier may instead use

\[
p_\theta(y\mid \phi(x)).
\]

The parameter \(\theta\) is estimated from training data.

\subsection{Loss function}\label{loss-function}

A loss function measures error. For supervised learning, a typical loss
has the form

\[
\ell(f_\theta(\phi(x)),y).
\]

The empirical risk is

\[
\widehat R_n(\theta)
=
\frac{1}{n}\sum_{i=1}^n \ell(f_\theta(\phi(x_i)),y_i).
\]

Training often means solving the optimization problem

\[
\widehat\theta
\in
\arg\min_{\theta} \widehat R_n(\theta).
\]

For probabilistic models, a common choice is the negative log-likelihood
loss:

\[
\ell(\theta;x,y)=-\log p_\theta(y\mid x).
\]

\subsection{Evaluation}\label{evaluation}

Evaluation measures performance on data not used for training. If
\((X,Y)\) is a future example from the target distribution, the
population risk is

\[
R(\theta)=E[\ell(f_\theta(\phi(X)),Y)].
\]

Since \(P\) is unknown, we estimate performance using validation or test
data. This connects NLP directly to statistical estimation.

\begin{tcolorbox}[enhanced jigsaw, arc=.35mm, bottomrule=.15mm, bottomtitle=1mm, breakable, colback=white, colbacktitle=quarto-callout-tip-color!10!white, colframe=quarto-callout-tip-color-frame, coltitle=black, left=2mm, leftrule=.75mm, opacityback=0, opacitybacktitle=0.6, rightrule=.15mm, title=\textcolor{quarto-callout-tip-color}{\faLightbulb}\hspace{0.5em}{A four-part checklist}, titlerule=0mm, toprule=.15mm, toptitle=1mm]

For any NLP system, ask:

\begin{enumerate}
\def\labelenumi{\arabic{enumi}.}
\tightlist
\item
  What is the representation \(\phi\)?
\item
  What is the model \(f_\theta\) or \(p_\theta\)?
\item
  What loss function is minimized?
\item
  What evaluation metric is reported?
\end{enumerate}

\end{tcolorbox}

\section{1.6 Three historical
approaches}\label{three-historical-approaches}

NLP methods are often grouped into rule-based, statistical, and neural
approaches. These are not mutually exclusive. Many modern systems
combine all three.

\subsection{Rule-based NLP}\label{rule-based-nlp}

Rule-based NLP uses explicit linguistic or symbolic rules. A simple rule
for detecting a question might be:

\[
f(x)=
\begin{cases}
1, & \text{if the final token is ?},\\
0, & \text{otherwise.}
\end{cases}
\]

Rule-based systems can be interpretable, but they may be brittle. They
often fail when language is ambiguous, informal, noisy, or
domain-specific.

\subsection{Statistical NLP}\label{statistical-nlp}

Statistical NLP treats language as random data. Instead of writing all
rules by hand, one estimates probabilities or parameters from a corpus.
A bigram language model, for example, approximates

\[
P(w_t\mid w_1,\ldots,w_{t-1})
\approx
P(w_t\mid w_{t-1}).
\]

The conditional probability can be estimated by counts:

\[
\widehat P(b\mid a)
=
\frac{\#(a,b)}{\#(a)},
\]

where \(\#(a,b)\) is the number of times token \(a\) is followed by
token \(b\) in the corpus. This is a maximum likelihood estimate under a
Markov assumption.

Classical statistical NLP is developed in detail in texts such as
Manning and Schuetze (1999) and Jurafsky and Martin (2024).

\subsection{Neural NLP}\label{neural-nlp}

Neural NLP represents tokens and texts by vectors and learns nonlinear
maps from data. A simplified neural classifier can be written as

\[
p_\theta(y\mid x)
=
\operatorname{softmax}(W h_\theta(x)+b),
\]

where \(h_\theta(x)\) is a learned representation of the input text. In
transformer models, \(h_\theta(x)\) is produced by repeated attention
and feedforward layers.

The important mathematical shift is that feature construction is no
longer entirely hand-designed. The representation itself becomes part of
the learned model.

\section{1.7 A small running example}\label{a-small-running-example}

Consider the tiny corpus

\[
D = \{x_1,x_2,x_3\},
\]

where

\[
\begin{aligned}
x_1 &= \text{machine learning models learn from data},\\
x_2 &= \text{language models predict the next token},\\
x_3 &= \text{statistics studies data uncertainty and models}.
\end{aligned}
\]

After tokenization, one possible vocabulary is

\[
V = \{\text{and},\text{data},\text{from},\text{language},\text{learn},\text{learning},\text{machine},\text{models},\text{next},\text{predict},\text{statistics},\text{studies},\text{the},\text{token},\text{uncertainty}\}.
\]

The vocabulary size is \(d=15\). Each document can be represented by a
count vector in \(R^{15}\).

For example, if the vocabulary is ordered as above, then \(x_1\) has
count vector

\[
\phi(x_1)
= (0,1,1,0,1,1,1,1,0,0,0,0,0,0,0)^T.
\]

The word \texttt{models} appears in all three documents, while
\texttt{machine} appears only in the first. Thus the count
representation already gives a simple numerical view of the corpus.

The term-document matrix is

\[
X =
\begin{bmatrix}
| & | & | \\
\phi(x_1) & \phi(x_2) & \phi(x_3) \\
| & | & |
\end{bmatrix}
\in R^{15\times 3}.
\]

This matrix is the beginning of linear algebraic NLP. Search,
clustering, latent semantic analysis, and embeddings can all be
introduced as ways to analyze or transform this matrix.

\subsection{A short Python check}\label{a-short-python-check}

The following code is not needed for the mathematical definition, but it
confirms the construction of the vocabulary and count vectors.

\protect\phantomsection\label{ch01-count-vectors}
\begin{Shaded}
\begin{Highlighting}[]
\NormalTok{corpus }\OperatorTok{=}\NormalTok{ [}
    \StringTok{"machine learning models learn from data"}\NormalTok{,}
    \StringTok{"language models predict the next token"}\NormalTok{,}
    \StringTok{"statistics studies data uncertainty and models"}\NormalTok{,}
\NormalTok{]}

\NormalTok{tokenized }\OperatorTok{=}\NormalTok{ [doc.split() }\ControlFlowTok{for}\NormalTok{ doc }\KeywordTok{in}\NormalTok{ corpus]}
\NormalTok{vocab }\OperatorTok{=} \BuiltInTok{sorted}\NormalTok{(}\BuiltInTok{set}\NormalTok{(token }\ControlFlowTok{for}\NormalTok{ doc }\KeywordTok{in}\NormalTok{ tokenized }\ControlFlowTok{for}\NormalTok{ token }\KeywordTok{in}\NormalTok{ doc))}

\NormalTok{count\_vectors }\OperatorTok{=}\NormalTok{ []}
\ControlFlowTok{for}\NormalTok{ doc }\KeywordTok{in}\NormalTok{ tokenized:}
\NormalTok{    count\_vectors.append([doc.count(token) }\ControlFlowTok{for}\NormalTok{ token }\KeywordTok{in}\NormalTok{ vocab])}

\BuiltInTok{print}\NormalTok{(}\StringTok{"Vocabulary size:"}\NormalTok{, }\BuiltInTok{len}\NormalTok{(vocab))}
\BuiltInTok{print}\NormalTok{(}\StringTok{"Vocabulary:"}\NormalTok{, vocab)}
\BuiltInTok{print}\NormalTok{(}\StringTok{"Count vectors:"}\NormalTok{)}
\ControlFlowTok{for}\NormalTok{ vector }\KeywordTok{in}\NormalTok{ count\_vectors:}
    \BuiltInTok{print}\NormalTok{(vector)}
\end{Highlighting}
\end{Shaded}

\begin{verbatim}
Vocabulary size: 15
Vocabulary: ['and', 'data', 'from', 'language', 'learn', 'learning', 'machine', 'models', 'next', 'predict', 'statistics', 'studies', 'the', 'token', 'uncertainty']
Count vectors:
[0, 1, 1, 0, 1, 1, 1, 1, 0, 0, 0, 0, 0, 0, 0]
[0, 0, 0, 1, 0, 0, 0, 1, 1, 1, 0, 0, 1, 1, 0]
[1, 1, 0, 0, 0, 0, 0, 1, 0, 0, 1, 1, 0, 0, 1]
\end{verbatim}

The code implements exactly the map

\[
\phi_j(x)=\#\{t:w_t=v_j\}.
\]

Later chapters will replace this simple count map by TF-IDF, low-rank
embeddings, neural embeddings, and transformer representations.

\section{1.8 What makes language mathematically
difficult?}\label{what-makes-language-mathematically-difficult}

NLP is difficult not because text cannot be turned into numbers, but
because the numerical representation must preserve useful structure.
Several mathematical difficulties appear immediately.

\subsection{Variable length}\label{variable-length}

Texts have different lengths. One sentence may contain five tokens,
while another may contain five hundred. Many mathematical models expect
fixed-dimensional vectors. Therefore we need maps that handle
variable-length inputs:

\[
\phi: V^* \longrightarrow R^d.
\]

Bag-of-words does this by ignoring order. Recurrent networks do this by
updating hidden states. Transformers do this by operating on token
sequences and producing contextual vectors.

\subsection{Discreteness}\label{discreteness}

Tokens are discrete. There is no natural meaning to arithmetic such as

\[
\text{cat} + \text{the}.
\]

Embedding methods introduce a learned map

\[
e: V \longrightarrow R^k,
\]

so that tokens become vectors. Once tokens are embedded, we can use
inner products, distances, projections, and optimization.

\subsection{Context dependence}\label{context-dependence}

The meaning of a word depends on context. For example, the word
\texttt{model} may refer to a mathematical model, a machine learning
model, or a physical object. A fixed word embedding assigns one vector
to each token type:

\[
e(\text{model}) \in R^k.
\]

A contextual embedding assigns a vector to a token occurrence inside a
sequence:

\[
h_t = h_t(w_1,\ldots,w_T) \in R^k.
\]

Thus the same token can receive different vectors in different
sentences.

\subsection{Combinatorial explosion}\label{combinatorial-explosion}

If the vocabulary has size \(d\), then the number of possible sequences
of length \(T\) is

\[
|V^T| = d^T.
\]

Even for moderate \(d\) and \(T\), this number is enormous. This is why
language modeling cannot simply store probabilities for every possible
sentence. A model must share statistical strength across related
contexts.

\subsection{Ambiguity}\label{ambiguity}

Language is ambiguous at many levels. A sentence can have multiple
possible meanings, labels, or continuations. Probability is therefore
natural. Instead of predicting one continuation with certainty, a
language model estimates

\[
p_\theta(w_t\mid w_1,\ldots,w_{t-1}).
\]

Generation then samples from this distribution or chooses tokens
according to a decoding rule.

\section{1.9 The mathematical pipeline of
NLP}\label{the-mathematical-pipeline-of-nlp}

A typical NLP system can be written as the composition

\[
V^*
\xrightarrow{\phi}
R^d
\xrightarrow{f_\theta}
Y.
\]

Here:

\begin{itemize}
\tightlist
\item
  \(V^*\) is the symbolic input space;
\item
  \(\phi\) is a representation map;
\item
  \(R^d\) is a numerical feature or embedding space;
\item
  \(f_\theta\) is a learned model;
\item
  \(Y\) is the output space.
\end{itemize}

For a probabilistic system, the output may be a distribution rather than
a single prediction:

\[
V^*
\xrightarrow{\phi}
R^d
\xrightarrow{p_\theta}
\Delta(Y),
\]

where \(\Delta(Y)\) denotes the set of probability distributions on
\(Y\).

If \(Y\) is finite with \(K\) classes, then

\[
\Delta(Y)=\left\{p\in R^K: p_j\geq 0,\ \sum_{j=1}^K p_j=1\right\}.
\]

This simplex will appear often because classification and next-token
prediction both output probability vectors.

\section{1.10 A first view of next-token
prediction}\label{a-first-view-of-next-token-prediction}

Modern large language models are often trained by next-token prediction.
Given a context

\[
c_t=(w_1,\ldots,w_{t-1}),
\]

the model outputs a probability distribution over the next token:

\[
p_\theta(\cdot\mid c_t) \in \Delta(V).
\]

If the true next token is \(w_t\), the negative log-likelihood loss is

\[
\ell_t(\theta) = -\log p_\theta(w_t\mid w_1,\ldots,w_{t-1}).
\]

For a sequence \((w_1,\ldots,w_T)\), the total loss is

\[
L(\theta)
=
-\sum_{t=1}^T \log p_\theta(w_t\mid w_1,\ldots,w_{t-1}).
\]

For a corpus, we sum this quantity over all training sequences. This is
simply maximum likelihood estimation for a conditional probability
model.

\begin{tcolorbox}[enhanced jigsaw, arc=.35mm, bottomrule=.15mm, bottomtitle=1mm, breakable, colback=white, colbacktitle=quarto-callout-note-color!10!white, colframe=quarto-callout-note-color-frame, coltitle=black, left=2mm, leftrule=.75mm, opacityback=0, opacitybacktitle=0.6, rightrule=.15mm, title=\textcolor{quarto-callout-note-color}{\faInfo}\hspace{0.5em}{Statistical interpretation}, titlerule=0mm, toprule=.15mm, toptitle=1mm]

Next-token prediction is not magic. It is conditional probability
estimation on a very large discrete space, using a high-dimensional
nonlinear parameterization.

\end{tcolorbox}

This perspective connects language models to classical statistics. The
novelty is not that the objective is unfamiliar. The novelty is the
scale of the data, the dimension of the parameter space, and the
expressiveness of the model class.

\section{1.11 Example: sentiment classification as empirical risk
minimization}\label{example-sentiment-classification-as-empirical-risk-minimization}

Suppose we have labeled texts

\[
D = \{(x_i,y_i)\}_{i=1}^n,
\]

where \(y_i\in\{0,1\}\). Choose a representation

\[
\phi: V^* \to R^d.
\]

A logistic text classifier has the form

\[
p_\theta(Y=1\mid x)
=
\sigma(\beta_0 + \beta^T\phi(x)),
\]

where

\[
\sigma(z)=\frac{1}{1+e^{-z}}
\]

and \(\theta=(\beta_0,\beta)\). The negative log-likelihood is

\[
\widehat R_n(\theta)
=
-
\frac{1}{n}
\sum_{i=1}^n
\left[
 y_i\log p_\theta(Y=1\mid x_i)
+
(1-y_i)\log(1-p_\theta(Y=1\mid x_i))
\right].
\]

Training the classifier means solving

\[
\widehat\theta
\in
\arg\min_\theta \widehat R_n(\theta).
\]

This example shows the basic structure of many NLP models:

\[
\text{text}
\longrightarrow
\text{features}
\longrightarrow
\text{probability}
\longrightarrow
\text{loss}
\longrightarrow
\text{optimization}.
\]

The later neural models are more complicated, but they still follow this
structure.

\section{1.12 Interactive HTML: representation and a linear
classifier}\label{interactive-html-representation-and-a-linear-classifier}

The following HTML component illustrates the simplest version of the NLP
pipeline

\[
\text{text}\longrightarrow \phi(x)\longrightarrow \beta_0+\beta^T\phi(x)\longrightarrow \sigma(\beta_0+\beta^T\phi(x)).
\]

Here the representation has only two coordinates: the count of a
positive token and the count of a negative token. This is intentionally
small so that the geometry is visible.

The HTML version of this chapter contains an interactive Plotly.js
component that lets the reader adjust feature counts and linear
classifier coefficients. In PDF output, the same idea is studied through
the logistic-classifier formulas and Python lab.

\section{1.13 What this book
emphasizes}\label{what-this-book-emphasizes}

This book is designed for students who already have or are building
mathematical maturity. It therefore emphasizes five principles.

\subsection{Principle 1: Start from the
space}\label{principle-1-start-from-the-space}

Always identify the space of objects being modeled. Are we working with
\(V\), \(V^T\), \(V^*\), \(R^d\), a probability simplex, or a space of
matrices?

\subsection{Principle 2: Track
dimensions}\label{principle-2-track-dimensions}

Dimension checking prevents many conceptual and programming errors. For
example, in attention,

\[
QK^T
\]

is meaningful only when the matrix dimensions agree. Later chapters will
carefully track the shapes of query, key, and value matrices.

\subsection{Principle 3: Separate representation from
prediction}\label{principle-3-separate-representation-from-prediction}

The representation map \(\phi\) and the prediction function \(f_\theta\)
play different roles. A linear classifier on a powerful representation
may outperform a nonlinear classifier on a weak representation.

\subsection{Principle 4: Interpret probabilities
conditionally}\label{principle-4-interpret-probabilities-conditionally}

Language models are conditional probability models. The conditioning
information matters. The expressions

\[
P(w_t), \qquad P(w_t\mid w_{t-1}), \qquad P(w_t\mid w_1,\ldots,w_{t-1})
\]

represent different assumptions.

\subsection{Principle 5: Understand loss
functions}\label{principle-5-understand-loss-functions}

A model is trained to do what its loss function rewards. If the loss is
next-token cross-entropy, the model is optimized for next-token
prediction. Other capabilities may emerge, but the training objective
remains central.

\section{1.14 Chapter summary}\label{chapter-summary}

Natural language processing begins with symbolic, variable-length
sequences and turns them into mathematical objects. The most important
objects in this chapter are:

{\def\LTcaptype{none} % do not increment counter
\begin{longtable}[]{@{}
  >{\raggedright\arraybackslash}p{(\linewidth - 2\tabcolsep) * \real{0.5000}}
  >{\raggedright\arraybackslash}p{(\linewidth - 2\tabcolsep) * \real{0.5000}}@{}}
\toprule\noalign{}
\begin{minipage}[b]{\linewidth}\raggedright
Concept
\end{minipage} & \begin{minipage}[b]{\linewidth}\raggedright
Mathematical form
\end{minipage} \\
\midrule\noalign{}
\endhead
\bottomrule\noalign{}
\endlastfoot
vocabulary & finite set \(V\) \\
sentence & sequence \((w_1,\ldots,w_T)\in V^T\) \\
all finite texts & sequence space \(V^*\) \\
corpus & finite sample \(D=\{x_1,\ldots,x_n\}\) \\
representation & map \(\phi:V^*\to R^d\) \\
classifier & map \(f:V^*\to C\) \\
probabilistic classifier & conditional distribution
\(p_\theta(y\mid x)\) \\
language model & distribution \(P(w_1,\ldots,w_T)\) \\
next-token model & \(p_\theta(w_t\mid w_1,\ldots,w_{t-1})\) \\
training & empirical risk minimization \\
evaluation & estimation of out-of-sample performance \\
\end{longtable}
}

The central lesson is simple: an NLP system is a mathematical pipeline.
Once the pipeline is written clearly, the role of linear algebra,
probability, statistics, optimization, and geometry becomes visible.

\section{1.15 Exercises}\label{exercises}

\subsection{Conceptual exercises}\label{conceptual-exercises}

\begin{enumerate}
\def\labelenumi{\arabic{enumi}.}
\item
  Let \(V=\{a,b,c\}\). List all elements of \(V^2\). How many elements
  are in \(V^T\) for general \(T\)?
\item
  Explain why \(V^*\) is not naturally a vector space. What extra
  structure is needed before linear algebra can be used?
\item
  For each task below, identify the input space and output space.

  \begin{itemize}
  \tightlist
  \item
    spam classification;
  \item
    named entity recognition;
  \item
    machine translation;
  \item
    document retrieval;
  \item
    next-token prediction.
  \end{itemize}
\item
  Give an example of two sentences with the same bag-of-words
  representation but different meanings.
\item
  Explain the difference between a representation map \(\phi\) and a
  prediction model \(f_\theta\).
\end{enumerate}

\subsection{Mathematical exercises}\label{mathematical-exercises}

\begin{enumerate}
\def\labelenumi{\arabic{enumi}.}
\setcounter{enumi}{5}
\item
  Suppose \(|V|=10{,}000\). How many possible sequences of length \(5\)
  are there? Why does this create difficulty for direct probability
  estimation?
\item
  Let \(C=\{1,2,3\}\). Write the probability simplex \(\Delta(C)\)
  explicitly as a subset of \(R^3\).
\item
  Let \(D=\{x_1,x_2,x_3,x_4\}\) be a corpus. Write the empirical
  distribution \(\widehat P_D\). What probability does it assign to a
  document that does not appear in the corpus?
\item
  Consider a binary classifier with predicted probability
  \(p_\theta(Y=1\mid x)=0.8\). Compute the negative log-likelihood loss
  when \(y=1\) and when \(y=0\).
\item
  Let \(x=(w_1,w_2,w_3,w_4)\). Use the chain rule to factor
  \(P(w_1,w_2,w_3,w_4)\) into conditional probabilities.
\end{enumerate}

\subsection{Applied exercises}\label{applied-exercises}

\begin{enumerate}
\def\labelenumi{\arabic{enumi}.}
\setcounter{enumi}{10}
\item
  Choose a short paragraph from a textbook. Define a vocabulary and
  construct its bag-of-words vector by hand.
\item
  Consider the task of helping students find university resources from
  official documents. Write this as a retrieval problem. What is the
  query space? What is the document space? What is the scoring function?
\item
  For a sentiment classification system, list one possible
  representation, one possible model, one possible loss function, and
  one possible evaluation metric.
\item
  Find a case where word order matters. Explain why a bag-of-words model
  would fail to distinguish the meanings.
\item
  Describe one way in which a corpus may fail to be an iid sample from
  the target distribution.
\end{enumerate}

\section{1.16 Python lab connection}\label{python-lab-connection}

The associated lab for this chapter is planned as:

\begin{Shaded}
\begin{Highlighting}[]
\NormalTok{labs/chapter{-}01{-}what{-}is{-}nlp.ipynb}
\end{Highlighting}
\end{Shaded}

The lab should be designed for independent study. It should include:

\begin{itemize}
\tightlist
\item
  text as raw strings;
\item
  tokenization by simple splitting;
\item
  vocabulary construction;
\item
  count-vector representation;
\item
  a small term-document matrix;
\item
  examples showing what information is lost by bag-of-words;
\item
  a final reflection connecting representation choices to modeling
  assumptions.
\end{itemize}

\section{1.17 Notes and references}\label{notes-and-references}

This chapter gives a mathematical orientation to NLP rather than a
historical survey. Classical statistical NLP is treated extensively in
Manning and Schuetze (1999). A broad modern introduction is given in
Jurafsky and Martin (2024). Vector space models and retrieval are
closely connected to information retrieval; see Manning et al. (2008).
Neural language modeling begins with ideas such as Bengio et al. (2003)
and leads eventually to attention-based architectures such as Vaswani et
al. (2017).

\chapter{Text, Tokens, and
Vocabularies}\label{text-tokens-and-vocabularies-1}

From raw strings to finite mathematical objects

\hfill\break

Natural language begins as text, but mathematical NLP begins only after
text has been converted into well-defined objects. This chapter develops
the first conversion step: characters, words, subwords, tokens,
vocabularies, and integer sequences. The goal is to make tokenization
mathematically explicit, because every later model in NLP depends on
this first representation map.

\section{Learning goals}\label{learning-goals-1}

After studying this chapter, you should be able to:

\begin{enumerate}
\def\labelenumi{\arabic{enumi}.}
\tightlist
\item
  model raw text as strings over an alphabet;
\item
  define tokenization as a map from strings to finite sequences of
  symbols or indices;
\item
  construct vocabularies from a corpus and represent tokens by one-hot
  vectors;
\item
  distinguish character, word, and subword tokenization;
\item
  explain the out-of-vocabulary problem and why modern language models
  use subword units;
\item
  build document-token count vectors and document-term matrices;
\item
  describe vocabulary construction as a statistical estimation and
  compression problem;
\item
  use Python to implement simple tokenizers and vocabulary maps;
\item
  use AI tools critically to inspect tokenization decisions.
\end{enumerate}

\section{Why tokenization is
mathematical}\label{why-tokenization-is-mathematical}

A computer does not directly understand a sentence such as

\begin{quote}
The model predicts the next token.
\end{quote}

Before a statistical model, neural network, or transformer can operate
on it, the sentence must be changed into a finite mathematical object.
For example, it might become

\[
(\text{the},\text{model},\text{predicts},\text{the},\text{next},\text{token}),
\]

or, after assigning integer IDs to tokens,

\[
(4, 18, 31, 4, 25, 7).
\]

This conversion is not a minor engineering detail. It defines the sample
space on which later probability models live. It also defines the
coordinates of the vectors and matrices used in learning algorithms. A
different tokenizer gives a different dataset.

\textbf{Central idea.} Tokenization is a representation map. It changes
an object from one mathematical space to another:

\[
\text{raw text} \longrightarrow \text{sequence of tokens} \longrightarrow \text{sequence of integers}.
\]

Once this map is fixed, we can define vectors, matrices, probability
distributions, loss functions, and neural networks on top of the
resulting representation.

For MA Applied Mathematics and MS Statistics students, the right
question is not only \emph{which tokenizer should we use?} but also:

\begin{itemize}
\tightlist
\item
  What is the domain and codomain of the tokenizer?
\item
  What information is preserved or discarded?
\item
  What random variables are being defined?
\item
  What statistical estimates depend on the vocabulary?
\item
  How does the choice of tokens change the geometry of the data?
\end{itemize}

\section{Text as strings over an
alphabet}\label{text-as-strings-over-an-alphabet}

Let \(A\) be a finite alphabet. In elementary formal language theory,
the set of all finite strings over \(A\) is denoted by

\[
A^* = \{\epsilon\} \cup A \cup A^2 \cup A^3 \cup \cdots,
\]

where \(\epsilon\) is the empty string. A raw text document is an
element

\[
x \in A^*.
\]

The alphabet \(A\) may be a character alphabet, a Unicode alphabet, a
byte alphabet, or another finite set of primitive symbols.

\subsection{Example: character
alphabet}\label{example-character-alphabet}

If

\[
A = \{a,b,c,\ldots,z,\text{space}\},
\]

then the sentence

\begin{Shaded}
\begin{Highlighting}[]
\NormalTok{machine learning}
\end{Highlighting}
\end{Shaded}

is a string in \(A^*\). Its length is the number of characters,
including the space.

\subsection{Corpus}\label{corpus}

A corpus is a finite collection of documents:

\[
D = \{x_1,x_2,\ldots,x_n\},
\]

where each \(x_i \in A^*\). In many statistical settings, we view
documents as observations sampled from an unknown distribution \(P\) on
\(A^*\):

\[
x_1,x_2,\ldots,x_n \sim P.
\]

This is already a modeling assumption. Real documents are not always
independent and identically distributed. For example, news articles from
the same week, posts from the same author, and pages from the same
textbook have dependence. Still, the i.i.d. viewpoint is a useful
starting point for defining empirical distributions.

\section{Normalization as a preprocessing
map}\label{normalization-as-a-preprocessing-map}

Before tokenization, one often applies a normalization function

\[
N:A^* \to A^*.
\]

Examples include lowercasing, removing punctuation, replacing multiple
spaces by one space, and converting accented characters into a standard
form.

For instance,

\[
N(\text{``Models, models!''}) = \text{``models models''}.
\]

Normalization is a many-to-one map. It usually discards information. If
uppercase and lowercase forms are collapsed, then we can no longer
distinguish
\texttt{Apple\textquotesingle{}\textquotesingle{}\ as\ a\ company\ from}apple'\,'
as a fruit by capitalization alone.

\textbf{Information loss.} A preprocessing map \(N\) is
information-preserving only if it is injective. Most normalizers used in
NLP are not injective. Therefore, preprocessing is a modeling decision,
not merely a cleaning step.

\section{Tokenization as a map}\label{tokenization-as-a-map}

A tokenizer maps a string to a finite sequence of tokens. Let \(V\) be a
vocabulary, meaning a finite set of allowed tokens:

\[
V = \{v_1,v_2,\ldots,v_d\}.
\]

A tokenizer is a function

\[
\tau:A^* \to V^*.
\]

For a document \(x\), the output is

\[
\tau(x)=(t_1,t_2,\ldots,t_m),
\]

where each \(t_j \in V\). The length \(m\) depends on both \(x\) and the
tokenizer.

After assigning each token an integer ID, we also have an index map

\[
\iota:V \to \{1,2,\ldots,d\}.
\]

The composition

\[
\iota \circ \tau:A^* \to \{1,2,\ldots,d\}^*
\]

is the form most directly used by neural networks.

\subsection{Small example}\label{small-example}

Let

\[
V = \{\text{data},\text{language},\text{learning},\text{model},\text{models},\text{token}\}.
\]

If

\[
x=\text{``language models learn from data''},
\]

then a simple word tokenizer might produce

\[
\tau(x)=(\text{language},\text{models},\text{learn},\text{from},\text{data}).
\]

But if \(\text{learn}\) and \(\text{from}\) are not in \(V\), then the
tokenizer must decide what to do. It might replace them by a special
unknown token:

\[
(\text{language},\text{models},\text{[UNK]},\text{[UNK]},\text{data}).
\]

This is the out-of-vocabulary problem.

\section{Three levels of
tokenization}\label{three-levels-of-tokenization}

There are three common mathematical levels of tokenization.

{\def\LTcaptype{none} % do not increment counter
\begin{longtable}[]{@{}
  >{\raggedright\arraybackslash}p{(\linewidth - 6\tabcolsep) * \real{0.2308}}
  >{\raggedleft\arraybackslash}p{(\linewidth - 6\tabcolsep) * \real{0.3077}}
  >{\raggedright\arraybackslash}p{(\linewidth - 6\tabcolsep) * \real{0.2308}}
  >{\raggedright\arraybackslash}p{(\linewidth - 6\tabcolsep) * \real{0.2308}}@{}}
\toprule\noalign{}
\begin{minipage}[b]{\linewidth}\raggedright
Tokenization level
\end{minipage} & \begin{minipage}[b]{\linewidth}\raggedleft
Basic unit
\end{minipage} & \begin{minipage}[b]{\linewidth}\raggedright
Strength
\end{minipage} & \begin{minipage}[b]{\linewidth}\raggedright
Weakness
\end{minipage} \\
\midrule\noalign{}
\endhead
\bottomrule\noalign{}
\endlastfoot
Character & characters & no unknown words & long sequences \\
Word & words & interpretable & large vocabulary and unknown words \\
Subword & pieces of words & balances vocabulary size and coverage & less
directly interpretable \\
\end{longtable}
}

\subsection{Character tokenization}\label{character-tokenization}

A character tokenizer uses individual characters as tokens. If

\begin{Shaded}
\begin{Highlighting}[]
\NormalTok{data}
\end{Highlighting}
\end{Shaded}

is tokenized at the character level, then

\[
\tau(x)=(d,a,t,a).
\]

This approach has a small vocabulary but produces long sequences. It
also makes semantic units hard to identify.

\subsection{Word tokenization}\label{word-tokenization}

A word tokenizer splits text into word-like units:

\[
\tau(\text{``data science''})=(\text{data},\text{science}).
\]

This is interpretable, but it creates a very large vocabulary. Rare
words, misspellings, names, and new technical terms cause difficulty.

\subsection{Subword tokenization}\label{subword-tokenization}

A subword tokenizer uses pieces that may be smaller than words but
larger than characters. For example,

\[
\text{``tokenization''} \mapsto (\text{token},\text{ization})
\]

or

\[
\text{``unbelievable''} \mapsto (\text{un},\text{believ},\text{able}).
\]

Modern language models typically use subword or byte-level tokenization
because it controls vocabulary size while still representing rare and
new words.

\section{One-hot representation}\label{one-hot-representation}

Once a vocabulary \(V=\{v_1,\ldots,v_d\}\) is fixed, each token \(v_j\)
can be represented by the \(j\)th standard basis vector in
\(\mathbb{R}^d\):

\[
e_j = (0,\ldots,0,1,0,\ldots,0)^T.
\]

The one-hot map is

\[
\phi:V \to \mathbb{R}^d,
\qquad
\phi(v_j)=e_j.
\]

For example, if

\[
V=(\text{data},\text{language},\text{model},\text{token}),
\]

then

\[
\phi(\text{model})=(0,0,1,0)^T.
\]

One-hot vectors are simple but do not encode similarity. In a one-hot
representation,

\[
\langle e_i,e_j\rangle = 0 \qquad \text{for } i\ne j.
\]

Thus the representation initially treats all distinct tokens as
orthogonal. Word embeddings, studied later, replace this rigid geometry
by learned vectors whose inner products carry semantic and statistical
information.

\section{Count vectors and the bag-of-words
map}\label{count-vectors-and-the-bag-of-words-map}

A tokenizer preserves order:

\[
(t_1,t_2,\ldots,t_m).
\]

Many classical NLP methods then discard order and keep only counts.
Given a vocabulary \(V=\{v_1,\ldots,v_d\}\), define the count vector

\[
c(x)=(c_1(x),\ldots,c_d(x))^T \in \mathbb{N}^d,
\]

where

\[
c_j(x)=\#\{k: t_k=v_j\}.
\]

The map

\[
x \mapsto c(x)
\]

is called the bag-of-words representation. It is useful because it turns
documents into vectors, but it loses word order. For example,

\begin{Shaded}
\begin{Highlighting}[]
\NormalTok{models use data}
\end{Highlighting}
\end{Shaded}

and

\begin{Shaded}
\begin{Highlighting}[]
\NormalTok{data use models}
\end{Highlighting}
\end{Shaded}

have the same bag-of-words vector under a word tokenizer.

\textbf{Representation tradeoff.} Sequences preserve order but are
harder to model. Count vectors are easy to analyze using linear algebra
and statistics, but they discard order.

\section{Document-term matrix}\label{document-term-matrix}

For a corpus \(D=\{x_1,\ldots,x_n\}\) and vocabulary
\(V=\{v_1,\ldots,v_d\}\), the document-term matrix is

\[
X \in \mathbb{R}^{n\times d},
\]

where

\[
X_{ij}=c_j(x_i).
\]

Each row is a document vector, and each column corresponds to a token.
This matrix is usually sparse because most documents use only a small
fraction of the vocabulary.

The document-term matrix is one of the central objects in classical NLP,
information retrieval, text classification, and topic modeling. It is
also the bridge from text to linear algebra.

\section{A small corpus by hand}\label{a-small-corpus-by-hand}

Consider the corpus

\[
\begin{aligned}
x_1 &= \text{``data models learn''},\\
x_2 &= \text{``language models predict tokens''},\\
x_3 &= \text{``statistics studies data''}.
\end{aligned}
\]

Suppose the vocabulary is

\[
V=(\text{data},\text{language},\text{learn},\text{models},\text{predict},\text{statistics},\text{studies},\text{tokens}).
\]

Then the document-term matrix is

\[
X=
\begin{bmatrix}
1&0&1&1&0&0&0&0\\
0&1&0&1&1&0&0&1\\
1&0&0&0&0&1&1&0
\end{bmatrix}.
\]

This matrix is not just a data table. It is the object on which linear
algebraic operations such as inner products, projections, low-rank
approximations, and clustering will be performed.

\section{Python example: tokenization and vocabulary
construction}\label{python-example-tokenization-and-vocabulary-construction}

The following example builds a simple word tokenizer using only Python
standard libraries. This is intentionally transparent. Industrial
tokenizers are more complicated, but the mathematical structure is the
same.

\begin{Shaded}
\begin{Highlighting}[]
\ImportTok{import}\NormalTok{ re}
\ImportTok{from}\NormalTok{ collections }\ImportTok{import}\NormalTok{ Counter}

\NormalTok{corpus }\OperatorTok{=}\NormalTok{ [}
    \StringTok{"Data models learn from data."}\NormalTok{,}
    \StringTok{"Language models predict the next token."}\NormalTok{,}
    \StringTok{"Statistics studies data, uncertainty, and models."}
\NormalTok{]}

\KeywordTok{def}\NormalTok{ normalize(text):}
    \CommentTok{"""Lowercase and keep alphabetic tokens only."""}
\NormalTok{    text }\OperatorTok{=}\NormalTok{ text.lower()}
\NormalTok{    text }\OperatorTok{=}\NormalTok{ re.sub(}\VerbatimStringTok{r"}\PreprocessorTok{[\^{}a{-}z}\DecValTok{\textbackslash{}s}\PreprocessorTok{]}\VerbatimStringTok{"}\NormalTok{, }\StringTok{" "}\NormalTok{, text)}
\NormalTok{    text }\OperatorTok{=}\NormalTok{ re.sub(}\VerbatimStringTok{r"}\DecValTok{\textbackslash{}s}\OperatorTok{+}\VerbatimStringTok{"}\NormalTok{, }\StringTok{" "}\NormalTok{, text).strip()}
    \ControlFlowTok{return}\NormalTok{ text}

\KeywordTok{def}\NormalTok{ word\_tokenize(text):}
    \ControlFlowTok{return}\NormalTok{ normalize(text).split()}

\NormalTok{tokenized }\OperatorTok{=}\NormalTok{ [word\_tokenize(doc) }\ControlFlowTok{for}\NormalTok{ doc }\KeywordTok{in}\NormalTok{ corpus]}
\NormalTok{tokenized}
\end{Highlighting}
\end{Shaded}

\protect\phantomsection\label{py-ch02-simple-tokenizer}
\begin{verbatim}
[['data', 'models', 'learn', 'from', 'data'],
 ['language', 'models', 'predict', 'the', 'next', 'token'],
 ['statistics', 'studies', 'data', 'uncertainty', 'and', 'models']]
\end{verbatim}

\protect\phantomsection\label{py-ch02-vocabulary}
\begin{Shaded}
\begin{Highlighting}[]
\NormalTok{all\_tokens }\OperatorTok{=}\NormalTok{ [token }\ControlFlowTok{for}\NormalTok{ doc }\KeywordTok{in}\NormalTok{ tokenized }\ControlFlowTok{for}\NormalTok{ token }\KeywordTok{in}\NormalTok{ doc]}
\NormalTok{counts }\OperatorTok{=}\NormalTok{ Counter(all\_tokens)}

\NormalTok{vocab }\OperatorTok{=} \BuiltInTok{sorted}\NormalTok{(counts)}
\NormalTok{token\_to\_id }\OperatorTok{=}\NormalTok{ \{token: j }\ControlFlowTok{for}\NormalTok{ j, token }\KeywordTok{in} \BuiltInTok{enumerate}\NormalTok{(vocab)\}}
\NormalTok{id\_to\_token }\OperatorTok{=}\NormalTok{ \{j: token }\ControlFlowTok{for}\NormalTok{ token, j }\KeywordTok{in}\NormalTok{ token\_to\_id.items()\}}

\BuiltInTok{print}\NormalTok{(}\StringTok{"Vocabulary size:"}\NormalTok{, }\BuiltInTok{len}\NormalTok{(vocab))}
\BuiltInTok{print}\NormalTok{(token\_to\_id)}
\end{Highlighting}
\end{Shaded}

\begin{verbatim}
Vocabulary size: 13
{'and': 0, 'data': 1, 'from': 2, 'language': 3, 'learn': 4, 'models': 5, 'next': 6, 'predict': 7, 'statistics': 8, 'studies': 9, 'the': 10, 'token': 11, 'uncertainty': 12}
\end{verbatim}

\begin{Shaded}
\begin{Highlighting}[]
\NormalTok{integer\_sequences }\OperatorTok{=}\NormalTok{ [}
\NormalTok{    [token\_to\_id[token] }\ControlFlowTok{for}\NormalTok{ token }\KeywordTok{in}\NormalTok{ doc]}
    \ControlFlowTok{for}\NormalTok{ doc }\KeywordTok{in}\NormalTok{ tokenized}
\NormalTok{]}

\NormalTok{integer\_sequences}
\end{Highlighting}
\end{Shaded}

\protect\phantomsection\label{py-ch02-integer-sequences}
\begin{verbatim}
[[1, 5, 4, 2, 1], [3, 5, 7, 10, 6, 11], [8, 9, 1, 12, 0, 5]]
\end{verbatim}

The variable \texttt{integer\_sequences} is an explicit realization of
the mathematical map

\[
\iota\circ\tau:A^*\to \{0,1,\ldots,d-1\}^*.
\]

\section{Python example: document-term
matrix}\label{python-example-document-term-matrix}

\begin{Shaded}
\begin{Highlighting}[]
\ImportTok{import}\NormalTok{ pandas }\ImportTok{as}\NormalTok{ pd}

\NormalTok{rows }\OperatorTok{=}\NormalTok{ []}
\ControlFlowTok{for}\NormalTok{ doc }\KeywordTok{in}\NormalTok{ tokenized:}
\NormalTok{    doc\_counts }\OperatorTok{=}\NormalTok{ Counter(doc)}
\NormalTok{    rows.append([doc\_counts[token] }\ControlFlowTok{for}\NormalTok{ token }\KeywordTok{in}\NormalTok{ vocab])}

\NormalTok{X }\OperatorTok{=}\NormalTok{ pd.DataFrame(rows, columns}\OperatorTok{=}\NormalTok{vocab)}
\NormalTok{X}
\end{Highlighting}
\end{Shaded}

\protect\phantomsection\label{py-ch02-document-term-matrix}
{\def\LTcaptype{none} % do not increment counter
\begin{longtable}[]{@{}llllllllllllll@{}}
\toprule\noalign{}
& and & data & from & language & learn & models & next & predict &
statistics & studies & the & token & uncertainty \\
\midrule\noalign{}
\endhead
\bottomrule\noalign{}
\endlastfoot
0 & 0 & 2 & 1 & 0 & 1 & 1 & 0 & 0 & 0 & 0 & 0 & 0 & 0 \\
1 & 0 & 0 & 0 & 1 & 0 & 1 & 1 & 1 & 0 & 0 & 1 & 1 & 0 \\
2 & 1 & 1 & 0 & 0 & 0 & 1 & 0 & 0 & 1 & 1 & 0 & 0 & 1 \\
\end{longtable}
}

The row sums are the document lengths after tokenization.

\begin{Shaded}
\begin{Highlighting}[]
\NormalTok{X.}\BuiltInTok{sum}\NormalTok{(axis}\OperatorTok{=}\DecValTok{1}\NormalTok{)}
\end{Highlighting}
\end{Shaded}

\protect\phantomsection\label{py-ch02-row-sums}
\begin{verbatim}
0    5
1    6
2    6
dtype: int64
\end{verbatim}

The column sums are the token frequencies over the corpus.

\begin{Shaded}
\begin{Highlighting}[]
\NormalTok{X.}\BuiltInTok{sum}\NormalTok{(axis}\OperatorTok{=}\DecValTok{0}\NormalTok{).sort\_values(ascending}\OperatorTok{=}\VariableTok{False}\NormalTok{)}
\end{Highlighting}
\end{Shaded}

\protect\phantomsection\label{py-ch02-column-sums}
\begin{verbatim}
data           3
models         3
and            1
from           1
language       1
learn          1
next           1
predict        1
statistics     1
studies        1
the            1
token          1
uncertainty    1
dtype: int64
\end{verbatim}

\section{Frequency thresholding and vocabulary
selection}\label{frequency-thresholding-and-vocabulary-selection}

In real corpora, the vocabulary can be very large. A common method is to
keep only tokens whose frequency is at least a threshold \(r\):

\[
V_r = \{v: f(v)\ge r\},
\]

where \(f(v)\) is the corpus frequency of token \(v\).

This creates a smaller vocabulary but discards rare tokens. Rare tokens
may be noise, but they may also be important names, mathematical terms,
or technical concepts.

The threshold \(r\) therefore controls a bias-variance-like tradeoff:

\begin{itemize}
\tightlist
\item
  large vocabulary: more expressive, but sparse and harder to estimate;
\item
  small vocabulary: more stable, but less informative.
\end{itemize}

\textbf{Vocabulary selection as estimation.} If \(p(v)\) is the unknown
probability of token \(v\) and \(\hat p(v)=f(v)/M\) is its empirical
frequency from \(M\) observed tokens, then low-frequency tokens have
high relative uncertainty. Removing rare tokens is a form of
regularization.

\section{Interactive HTML: vocabulary threshold
explorer}\label{interactive-html-vocabulary-threshold-explorer}

The following interactive component uses Plotly.js to show how the
vocabulary changes as the frequency threshold changes. It is designed to
render in the HTML version of the Quarto book.

The mathematical point of the explorer is that the dimension \(d\) of
the vector space \(\mathbb{R}^d\) is not intrinsic to the corpus. It is
chosen by the vocabulary construction rule.

\section{Unknown tokens}\label{unknown-tokens}

Suppose a word tokenizer is trained on a corpus whose vocabulary is
\(V\). When a new document contains a token not in \(V\), one classical
solution is to use a special token \texttt{{[}UNK{]}}.

Mathematically, extend the vocabulary to

\[
V' = V\cup \{\text{[UNK]}\}.
\]

Then define

\[
\pi(v)=
\begin{cases}
v, & v\in V,\\
\text{[UNK]}, & v\notin V.
\end{cases}
\]

This projection collapses all unseen tokens into a single symbol. The
price is severe information loss. For example, all of the following may
become \texttt{{[}UNK{]}}:

\begin{Shaded}
\begin{Highlighting}[]
\NormalTok{ChatGPT}
\NormalTok{Noetherian}
\NormalTok{eigenvector}
\NormalTok{Boston}
\NormalTok{mispeledword}
\end{Highlighting}
\end{Shaded}

The unknown token solves the finite-vocabulary problem but destroys
distinctions among rare and unseen terms.

\section{Subword tokenization as a
remedy}\label{subword-tokenization-as-a-remedy}

Subword tokenization reduces the unknown-token problem by building words
from smaller units. If the vocabulary contains common pieces such as

\[
\text{eigen}, \quad \text{vector}, \quad \text{ization}, \quad \text{pre}, \quad \text{train},
\]

then rare words may be represented as sequences of meaningful or
semi-meaningful pieces.

For example,

\[
\text{eigenvectorization}
\mapsto
(\text{eigen},\text{vector},\text{ization}).
\]

The word may not be in the vocabulary, but its parts may be.

\section{A clean view of byte-pair
encoding}\label{a-clean-view-of-byte-pair-encoding}

Byte-pair encoding, often abbreviated BPE, is a common subword
vocabulary construction method. The exact implementations used in
production systems vary, but the core mathematical idea is simple.

Start with a corpus represented as sequences of basic symbols, often
characters or bytes. Repeatedly merge the most frequent adjacent pair
into a new token.

\subsection{Abstract BPE algorithm}\label{abstract-bpe-algorithm}

Let the current tokenized corpus be a finite multiset of sequences over
a vocabulary \(V_k\). At step \(k\):

\begin{enumerate}
\def\labelenumi{\arabic{enumi}.}
\item
  count all adjacent token pairs \((a,b)\) in the corpus;
\item
  find a most frequent pair

  \[
  (a_k,b_k)\in \arg\max_{(a,b)} f_k(a,b);
  \]
\item
  create a new token

  \[
  z_k = a_k b_k;
  \]
\item
  replace occurrences of \((a_k,b_k)\) by \(z_k\);
\item
  update the vocabulary

  \[
  V_{k+1}=V_k\cup\{z_k\}.
  \]
\end{enumerate}

After \(K\) merges, the vocabulary has grown from basic symbols to
useful subwords.

\textbf{BPE as greedy compression.} BPE greedily chooses frequent
adjacent pairs. It is not optimizing a simple global convex objective,
but it can be interpreted as building a compact code for frequent local
patterns.

\section{Python example: a tiny BPE
demonstration}\label{python-example-a-tiny-bpe-demonstration}

The following code implements a small educational version of BPE. It is
not intended to match production tokenizers. Its purpose is to reveal
the mathematical structure.

\protect\phantomsection\label{py-ch02-mini-bpe}
\begin{Shaded}
\begin{Highlighting}[]
\ImportTok{from}\NormalTok{ collections }\ImportTok{import}\NormalTok{ Counter}

\KeywordTok{def}\NormalTok{ word\_as\_chars(word):}
    \CommentTok{\# The end{-}of{-}word marker helps distinguish "low" from "lower".}
    \ControlFlowTok{return} \BuiltInTok{tuple}\NormalTok{(}\BuiltInTok{list}\NormalTok{(word) }\OperatorTok{+}\NormalTok{ [}\StringTok{"\textless{}/w\textgreater{}"}\NormalTok{])}

\NormalTok{words }\OperatorTok{=}\NormalTok{ [}\StringTok{"low"}\NormalTok{, }\StringTok{"lower"}\NormalTok{, }\StringTok{"newest"}\NormalTok{, }\StringTok{"widest"}\NormalTok{, }\StringTok{"low"}\NormalTok{, }\StringTok{"lower"}\NormalTok{, }\StringTok{"low"}\NormalTok{]}
\NormalTok{corpus\_tokens }\OperatorTok{=}\NormalTok{ [word\_as\_chars(w) }\ControlFlowTok{for}\NormalTok{ w }\KeywordTok{in}\NormalTok{ words]}


\KeywordTok{def}\NormalTok{ pair\_counts(sequences):}
\NormalTok{    counts }\OperatorTok{=}\NormalTok{ Counter()}
    \ControlFlowTok{for}\NormalTok{ seq }\KeywordTok{in}\NormalTok{ sequences:}
        \ControlFlowTok{for}\NormalTok{ a, b }\KeywordTok{in} \BuiltInTok{zip}\NormalTok{(seq[:}\OperatorTok{{-}}\DecValTok{1}\NormalTok{], seq[}\DecValTok{1}\NormalTok{:]):}
\NormalTok{            counts[(a, b)] }\OperatorTok{+=} \DecValTok{1}
    \ControlFlowTok{return}\NormalTok{ counts}


\KeywordTok{def}\NormalTok{ merge\_pair(sequence, pair):}
\NormalTok{    a, b }\OperatorTok{=}\NormalTok{ pair}
\NormalTok{    merged }\OperatorTok{=}\NormalTok{ []}
\NormalTok{    i }\OperatorTok{=} \DecValTok{0}
    \ControlFlowTok{while}\NormalTok{ i }\OperatorTok{\textless{}} \BuiltInTok{len}\NormalTok{(sequence):}
        \ControlFlowTok{if}\NormalTok{ i }\OperatorTok{\textless{}} \BuiltInTok{len}\NormalTok{(sequence) }\OperatorTok{{-}} \DecValTok{1} \KeywordTok{and}\NormalTok{ sequence[i] }\OperatorTok{==}\NormalTok{ a }\KeywordTok{and}\NormalTok{ sequence[i }\OperatorTok{+} \DecValTok{1}\NormalTok{] }\OperatorTok{==}\NormalTok{ b:}
\NormalTok{            merged.append(a }\OperatorTok{+}\NormalTok{ b)}
\NormalTok{            i }\OperatorTok{+=} \DecValTok{2}
        \ControlFlowTok{else}\NormalTok{:}
\NormalTok{            merged.append(sequence[i])}
\NormalTok{            i }\OperatorTok{+=} \DecValTok{1}
    \ControlFlowTok{return} \BuiltInTok{tuple}\NormalTok{(merged)}


\KeywordTok{def}\NormalTok{ bpe\_step(sequences):}
\NormalTok{    counts }\OperatorTok{=}\NormalTok{ pair\_counts(sequences)}
\NormalTok{    best\_pair, best\_count }\OperatorTok{=}\NormalTok{ counts.most\_common(}\DecValTok{1}\NormalTok{)[}\DecValTok{0}\NormalTok{]}
\NormalTok{    new\_sequences }\OperatorTok{=}\NormalTok{ [merge\_pair(seq, best\_pair) }\ControlFlowTok{for}\NormalTok{ seq }\KeywordTok{in}\NormalTok{ sequences]}
    \ControlFlowTok{return}\NormalTok{ best\_pair, best\_count, new\_sequences}

\NormalTok{sequences }\OperatorTok{=}\NormalTok{ corpus\_tokens}
\ControlFlowTok{for}\NormalTok{ step }\KeywordTok{in} \BuiltInTok{range}\NormalTok{(}\DecValTok{5}\NormalTok{):}
\NormalTok{    best\_pair, best\_count, sequences }\OperatorTok{=}\NormalTok{ bpe\_step(sequences)}
    \BuiltInTok{print}\NormalTok{(}\SpecialStringTok{f"Step }\SpecialCharTok{\{}\NormalTok{step }\OperatorTok{+} \DecValTok{1}\SpecialCharTok{\}}\SpecialStringTok{: merge }\SpecialCharTok{\{}\NormalTok{best\_pair}\SpecialCharTok{\}}\SpecialStringTok{ with frequency }\SpecialCharTok{\{}\NormalTok{best\_count}\SpecialCharTok{\}}\SpecialStringTok{"}\NormalTok{)}
    \BuiltInTok{print}\NormalTok{(sequences[:}\DecValTok{3}\NormalTok{])}
\end{Highlighting}
\end{Shaded}

\begin{verbatim}
Step 1: merge ('l', 'o') with frequency 5
[('lo', 'w', '</w>'), ('lo', 'w', 'e', 'r', '</w>'), ('n', 'e', 'w', 'e', 's', 't', '</w>')]
Step 2: merge ('lo', 'w') with frequency 5
[('low', '</w>'), ('low', 'e', 'r', '</w>'), ('n', 'e', 'w', 'e', 's', 't', '</w>')]
Step 3: merge ('low', '</w>') with frequency 3
[('low</w>',), ('low', 'e', 'r', '</w>'), ('n', 'e', 'w', 'e', 's', 't', '</w>')]
Step 4: merge ('low', 'e') with frequency 2
[('low</w>',), ('lowe', 'r', '</w>'), ('n', 'e', 'w', 'e', 's', 't', '</w>')]
Step 5: merge ('lowe', 'r') with frequency 2
[('low</w>',), ('lower', '</w>'), ('n', 'e', 'w', 'e', 's', 't', '</w>')]
\end{verbatim}

Notice that BPE is data-dependent. The same word may be segmented
differently if the training corpus changes.

\section{Tokenization and probability
spaces}\label{tokenization-and-probability-spaces}

Tokenization changes the probability space of language modeling. If
tokens are words, then a language model estimates probabilities such as

\[
P(W_t=\text{models}\mid W_1,\ldots,W_{t-1}).
\]

If tokens are subwords, then the random variables are subword tokens:

\[
P(T_t=\text{mod}\mid T_1,\ldots,T_{t-1}).
\]

These are not the same statistical problem. The vocabulary determines
the categorical sample space.

For a vocabulary of size \(d\), a next-token model outputs a probability
vector

\[
p_t=(p_{t1},\ldots,p_{td})\in \Delta^{d-1},
\]

where

\[
\Delta^{d-1}=\left\{p\in\mathbb{R}^d: p_j\ge 0,\ \sum_{j=1}^d p_j=1\right\}.
\]

Thus changing the tokenizer changes \(d\), the dimension of the output
simplex.

\section{Empirical token
distribution}\label{empirical-token-distribution}

Given a tokenized corpus with total number of tokens \(M\), define the
empirical token distribution

\[
\hat p_j = \frac{f_j}{M},
\]

where \(f_j\) is the frequency of token \(v_j\).

The vector

\[
\hat p=(\hat p_1,\ldots,\hat p_d)
\]

belongs to the probability simplex \(\Delta^{d-1}\). It is a simple but
important statistical summary of a corpus.

The empirical entropy is

\[
\widehat H = -\sum_{j=1}^d \hat p_j \log \hat p_j,
\]

where terms with \(\hat p_j=0\) are omitted. A corpus whose token
frequencies are evenly spread has larger entropy than one dominated by a
few common tokens.

\protect\phantomsection\label{py-ch02-empirical-distribution}
\begin{Shaded}
\begin{Highlighting}[]
\ImportTok{import}\NormalTok{ math}

\NormalTok{M }\OperatorTok{=} \BuiltInTok{sum}\NormalTok{(counts.values())}
\NormalTok{empirical\_p }\OperatorTok{=}\NormalTok{ \{token: freq }\OperatorTok{/}\NormalTok{ M }\ControlFlowTok{for}\NormalTok{ token, freq }\KeywordTok{in}\NormalTok{ counts.items()\}}
\NormalTok{entropy }\OperatorTok{=} \OperatorTok{{-}}\BuiltInTok{sum}\NormalTok{(p }\OperatorTok{*}\NormalTok{ math.log(p) }\ControlFlowTok{for}\NormalTok{ p }\KeywordTok{in}\NormalTok{ empirical\_p.values())}

\BuiltInTok{print}\NormalTok{(}\StringTok{"Total tokens:"}\NormalTok{, M)}
\BuiltInTok{print}\NormalTok{(}\StringTok{"Empirical entropy:"}\NormalTok{, }\BuiltInTok{round}\NormalTok{(entropy, }\DecValTok{4}\NormalTok{))}
\BuiltInTok{print}\NormalTok{(}\BuiltInTok{sorted}\NormalTok{(empirical\_p.items(), key}\OperatorTok{=}\KeywordTok{lambda}\NormalTok{ x: }\OperatorTok{{-}}\NormalTok{x[}\DecValTok{1}\NormalTok{])[:}\DecValTok{5}\NormalTok{])}
\end{Highlighting}
\end{Shaded}

\begin{verbatim}
Total tokens: 17
Empirical entropy: 2.4455
[('data', 0.17647058823529413), ('models', 0.17647058823529413), ('learn', 0.058823529411764705), ('from', 0.058823529411764705), ('language', 0.058823529411764705)]
\end{verbatim}

\section{Zipf-like behavior}\label{zipf-like-behavior}

In large natural language corpora, token frequencies tend to be highly
skewed. A small number of tokens occur very frequently, while many
tokens occur rarely. This qualitative pattern is often called Zipf-like
behavior.

If tokens are ranked by decreasing frequency, the frequency of the token
with rank \(r\) often roughly follows a power-law shape:

\[
f_r \approx C r^{-\alpha}
\]

for constants \(C>0\) and \(\alpha>0\) over a range of ranks.

This matters because rare tokens are not an exception. They are a
structural feature of language data.

\section{Vocabulary growth and Heaps-like
behavior}\label{vocabulary-growth-and-heaps-like-behavior}

Another common empirical pattern is that the number of distinct observed
tokens grows as more text is collected. If \(M\) is the number of
observed tokens and \(V(M)\) is the number of distinct token types, then
one often sees approximate behavior of the form

\[
V(M) \approx K M^\beta,
\]

where \(0<\beta<1\).

The implication is important: collecting more data keeps revealing new
tokens. Therefore a fixed word vocabulary will always face rare and
unseen words.

\section{Sequence length versus vocabulary
size}\label{sequence-length-versus-vocabulary-size}

Tokenization involves a fundamental tradeoff:

\[
\text{small vocabulary} \quad \Longleftrightarrow \quad \text{longer sequences},
\]

\[
\text{large vocabulary} \quad \Longleftrightarrow \quad \text{shorter sequences}.
\]

Character tokenization has a very small vocabulary but long sequences.
Word tokenization has shorter sequences but a huge vocabulary. Subword
tokenization lives between these extremes.

This matters computationally. In a transformer, self-attention has a
cost that grows roughly quadratically with sequence length \(m\):

\[
\text{attention cost} \propto m^2.
\]

A tokenizer that produces longer sequences can substantially increase
computational cost.

\section{Padding and truncation}\label{padding-and-truncation}

Neural network batches often require sequences of equal length. If a
batch contains token sequences

\[
(t_{i1},\ldots,t_{im_i}), \qquad i=1,\ldots,B,
\]

then a maximum length \(L\) is chosen. Shorter sequences are padded
using a special token \texttt{{[}PAD{]}}, and longer sequences are
truncated.

For example,

\[
(7,3,8) \mapsto (7,3,8,0,0)
\]

if \texttt{{[}PAD{]}} has index \(0\) and \(L=5\).

Padding is another representation choice. Models must be told which
positions are real and which are padding. This is done using an
attention mask or sequence mask.

\section{Special tokens}\label{special-tokens}

Many NLP systems add special tokens to the vocabulary. Common examples
include:

{\def\LTcaptype{none} % do not increment counter
\begin{longtable}[]{@{}ll@{}}
\toprule\noalign{}
Token & Meaning \\
\midrule\noalign{}
\endhead
\bottomrule\noalign{}
\endlastfoot
\texttt{{[}PAD{]}} & padding token \\
\texttt{{[}UNK{]}} & unknown token \\
\texttt{{[}BOS{]}} & beginning of sequence \\
\texttt{{[}EOS{]}} & end of sequence \\
\texttt{{[}MASK{]}} & masked token for prediction tasks \\
\texttt{{[}CLS{]}} & classification token in some encoder models \\
\texttt{{[}SEP{]}} & separator token between segments \\
\end{longtable}
}

Mathematically, special tokens are ordinary vocabulary elements with
assigned indices. Their meaning comes from how the model and loss
function use them.

\section{Tokenization for AI systems}\label{tokenization-for-ai-systems}

Tokenization is especially important in modern AI systems because
language models operate on token sequences, not directly on words. A
prompt, a retrieved document, or a generated answer is represented
internally as tokens.

This affects several practical AI questions.

\subsection{Context length}\label{context-length}

If a model has a maximum context length \(L\), then the tokenized input
must satisfy

\[
m \le L.
\]

Here \(m\) is the number of tokens after tokenization, not the number of
words. A text that looks short to a human may be long in tokens if it
contains code, formulas, tables, or uncommon terms.

\subsection{Retrieval-augmented
generation}\label{retrieval-augmented-generation}

In retrieval-augmented generation, documents are often split into chunks
before being embedded and retrieved. A chunking rule is partly a
tokenization problem. It determines the units that become searchable
objects.

A simple chunking map has the form

\[
\gamma:A^*\to (A^*)^q,
\]

where a document is divided into \(q\) text chunks. Each chunk is then
tokenized and embedded.

Bad chunking may separate definitions from examples, theorem statements
from assumptions, or student questions from answers. Good chunking
respects mathematical and semantic structure.

\subsection{Prompt design}\label{prompt-design}

A prompt is also a token sequence. Repetition, formatting, examples, and
equations all affect the tokenized representation. Therefore prompt
design has a mathematical side: it changes the conditioning information
available to a next-token probability model.

\section{AI-assisted learning
component}\label{ai-assisted-learning-component}

Use an AI assistant as a mathematical checker, not as a black box. For
this chapter, useful AI-assisted tasks include:

\begin{tcolorbox}[enhanced jigsaw, arc=.35mm, bottomrule=.15mm, bottomtitle=1mm, breakable, colback=white, colbacktitle=quarto-callout-tip-color!10!white, colframe=quarto-callout-tip-color-frame, coltitle=black, left=2mm, leftrule=.75mm, opacityback=0, opacitybacktitle=0.6, rightrule=.15mm, title=\textcolor{quarto-callout-tip-color}{\faLightbulb}\hspace{0.5em}{AI activity 1: tokenizer audit}, titlerule=0mm, toprule=.15mm, toptitle=1mm]

Give an AI assistant the following instruction:

\begin{quote}
Tokenize the sentence below in three ways: character-level, word-level,
and a plausible subword-level segmentation. Then explain what
information is preserved and what information is lost by each
representation. Sentence:
\texttt{Mathematical\ tokenization\ changes\ the\ sample\ space.}
\end{quote}

Then check whether the assistant clearly identifies the domain,
codomain, and representation loss for each tokenizer.

\end{tcolorbox}

\begin{tcolorbox}[enhanced jigsaw, arc=.35mm, bottomrule=.15mm, bottomtitle=1mm, breakable, colback=white, colbacktitle=quarto-callout-tip-color!10!white, colframe=quarto-callout-tip-color-frame, coltitle=black, left=2mm, leftrule=.75mm, opacityback=0, opacitybacktitle=0.6, rightrule=.15mm, title=\textcolor{quarto-callout-tip-color}{\faLightbulb}\hspace{0.5em}{AI activity 2: mathematical vocabulary design}, titlerule=0mm, toprule=.15mm, toptitle=1mm]

Ask an AI assistant:

\begin{quote}
I am building a tokenizer for graduate mathematics text. Give examples
of words or expressions that should not be carelessly split. Organize
the answer by mathematical area: linear algebra, probability,
statistics, optimization, and NLP.
\end{quote}

Use the answer to think about domain-specific vocabulary design.

\end{tcolorbox}

\begin{tcolorbox}[enhanced jigsaw, arc=.35mm, bottomrule=.15mm, bottomtitle=1mm, breakable, colback=white, colbacktitle=quarto-callout-warning-color!10!white, colframe=quarto-callout-warning-color-frame, coltitle=black, left=2mm, leftrule=.75mm, opacityback=0, opacitybacktitle=0.6, rightrule=.15mm, title=\textcolor{quarto-callout-warning-color}{\faExclamationTriangle}\hspace{0.5em}{AI caution}, titlerule=0mm, toprule=.15mm, toptitle=1mm]

Do not paste private student data, confidential emails, unpublished
manuscripts, or sensitive institutional documents into a public AI tool.
For tokenizer experiments, use synthetic examples or public text.

\end{tcolorbox}

\section{Mathematical summary}\label{mathematical-summary}

The main objects of this chapter are summarized below.

{\def\LTcaptype{none} % do not increment counter
\begin{longtable}[]{@{}
  >{\raggedright\arraybackslash}p{(\linewidth - 4\tabcolsep) * \real{0.3333}}
  >{\raggedright\arraybackslash}p{(\linewidth - 4\tabcolsep) * \real{0.3333}}
  >{\raggedright\arraybackslash}p{(\linewidth - 4\tabcolsep) * \real{0.3333}}@{}}
\toprule\noalign{}
\begin{minipage}[b]{\linewidth}\raggedright
Object
\end{minipage} & \begin{minipage}[b]{\linewidth}\raggedright
Mathematical form
\end{minipage} & \begin{minipage}[b]{\linewidth}\raggedright
Interpretation
\end{minipage} \\
\midrule\noalign{}
\endhead
\bottomrule\noalign{}
\endlastfoot
Alphabet & \(A\) & primitive symbols \\
String space & \(A^*\) & all finite raw strings \\
Corpus & \(D=\{x_1,\ldots,x_n\}\) & observed documents \\
Normalizer & \(N:A^*\to A^*\) & preprocessing map \\
Vocabulary & \(V=\{v_1,\ldots,v_d\}\) & finite token set \\
Tokenizer & \(\tau:A^*\to V^*\) & string-to-token sequence map \\
Index map & \(\iota:V\to\{1,\ldots,d\}\) & token-to-ID map \\
One-hot map & \(\phi:V\to\mathbb{R}^d\) & token-to-basis-vector map \\
Count vector & \(c(x)\in\mathbb{N}^d\) & bag-of-words representation \\
Document-term matrix & \(X\in\mathbb{R}^{n\times d}\) & corpus as
matrix \\
Empirical distribution & \(\hat p\in\Delta^{d-1}\) & token frequency
distribution \\
\end{longtable}
}

\section{Common misconceptions}\label{common-misconceptions}

\subsection{Misconception 1: tokenization is just
cleaning}\label{misconception-1-tokenization-is-just-cleaning}

Tokenization is not just cleaning. It defines the symbols used by the
model. It changes the statistical problem.

\subsection{Misconception 2: words are the natural
tokens}\label{misconception-2-words-are-the-natural-tokens}

Words are intuitive for humans, but they are not always the best
modeling units. Subwords often provide better coverage and a more stable
vocabulary.

\subsection{Misconception 3: larger vocabulary is always
better}\label{misconception-3-larger-vocabulary-is-always-better}

A larger vocabulary can represent more distinctions, but it also creates
a higher-dimensional and sparser estimation problem.

\subsection{Misconception 4: one-hot vectors encode
meaning}\label{misconception-4-one-hot-vectors-encode-meaning}

One-hot vectors encode identity, not meaning. Meaning enters later
through co-occurrence, embeddings, attention, and learned model
parameters.

\section{Worked exercise: comparing
tokenizers}\label{worked-exercise-comparing-tokenizers}

Consider the sentence

\begin{Shaded}
\begin{Highlighting}[]
\NormalTok{Transformer{-}based models tokenize mathematical notation differently.}
\end{Highlighting}
\end{Shaded}

\subsection{Character-level
representation}\label{character-level-representation}

A character tokenizer might produce every character, including
punctuation and spaces. This has strong coverage but long length.

\subsection{Word-level representation}\label{word-level-representation}

A word tokenizer might produce

\[
(\text{transformer},\text{based},\text{models},\text{tokenize},\text{mathematical},\text{notation},\text{differently}).
\]

This is interpretable but loses the hyphen structure unless the
tokenizer preserves it.

\subsection{Subword-level
representation}\label{subword-level-representation}

A subword tokenizer might produce pieces such as

\[
(\text{transform},\text{er},\text{-},\text{based},\text{models},\text{token},\text{ize},\ldots).
\]

This preserves more detail than a crude word tokenizer and avoids
unknown words more effectively.

\section{Exercises}\label{exercises-1}

\begin{enumerate}
\def\labelenumi{\arabic{enumi}.}
\item
  Let \(A=\{a,b\}\). List all elements of \(A^*\) of length at most
  \(3\). How many are there?
\item
  Let \(V=\{\text{data},\text{model},\text{token}\}\). Define a
  tokenizer for the sentence \texttt{data\ data\ model\ token}. Write
  the corresponding count vector.
\item
  Construct the document-term matrix for the corpus

\begin{Shaded}
\begin{Highlighting}[]
\NormalTok{data models data}
\NormalTok{language models language}
\NormalTok{data language token}
\end{Highlighting}
\end{Shaded}

  using the vocabulary ordered alphabetically.
\item
  Show that the bag-of-words map is not injective by giving two
  different sentences with the same count vector.
\item
  Suppose a vocabulary has size \(d=50,000\). What is the dimension of a
  one-hot token vector? How many nonzero entries does it have?
\item
  Let \(f(v)\) be the frequency of token \(v\). For the rule
  \(V_r=\{v:f(v)\ge r\}\), explain what happens to vocabulary size as
  \(r\) increases.
\item
  Consider a corpus with total token count \(M=1000\). A token appears
  twice. What is its empirical probability? Why might its estimate be
  unstable?
\item
  Explain why replacing all unseen words by \texttt{{[}UNK{]}} is a
  many-to-one map.
\item
  Implement a simple word tokenizer in Python that lowercases text and
  removes punctuation. Test it on a sentence containing mathematical
  notation such as \texttt{P(X\ \textgreater{}\ 2)\ =\ 0.13}.
\item
  Design a domain-specific tokenization rule for mathematical text. What
  should happen to symbols such as \texttt{\$x\_i\$},
  \texttt{\$\textbackslash{}nabla\ f\$}, or \texttt{\$A\^{}T\ A\$}?
\end{enumerate}

\section{Mini-project: tokenizer design for mathematical
documents}\label{mini-project-tokenizer-design-for-mathematical-documents}

Choose a short paragraph from a mathematical textbook or lecture note.
Compare three tokenization strategies:

\begin{enumerate}
\def\labelenumi{\arabic{enumi}.}
\tightlist
\item
  character-level tokenization;
\item
  word-level tokenization;
\item
  a manually designed math-aware tokenization.
\end{enumerate}

For each strategy, report:

\begin{itemize}
\tightlist
\item
  number of tokens;
\item
  vocabulary size;
\item
  examples of ambiguous tokens;
\item
  information preserved;
\item
  information lost;
\item
  whether the representation is suitable for search, classification, or
  language modeling.
\end{itemize}

\section{Lab connection}\label{lab-connection}

The associated independent-study lab should be saved as:

\begin{Shaded}
\begin{Highlighting}[]
\NormalTok{labs/chapter{-}02{-}text{-}tokens{-}vocabularies.ipynb}
\end{Highlighting}
\end{Shaded}

Recommended lab sections:

\begin{enumerate}
\def\labelenumi{\arabic{enumi}.}
\tightlist
\item
  strings, normalization, and regular expressions;
\item
  word tokenization from scratch;
\item
  vocabulary construction and integer encoding;
\item
  document-term matrices;
\item
  empirical token distributions;
\item
  a tiny BPE implementation;
\item
  comparison of character, word, and subword tokenization;
\item
  AI-assisted tokenizer audit.
\end{enumerate}

\section{Key takeaways}\label{key-takeaways}

\begin{itemize}
\tightlist
\item
  Text must be converted into finite mathematical objects before
  modeling.
\item
  A tokenizer is a map \(\tau:A^*\to V^*\).
\item
  A vocabulary determines the dimension of one-hot vectors and count
  vectors.
\item
  Count vectors and document-term matrices are the first bridge from NLP
  to linear algebra.
\item
  Word tokenization is interpretable but has serious unknown-word
  problems.
\item
  Subword tokenization balances vocabulary size, sequence length, and
  coverage.
\item
  Tokenization changes the probability space of language modeling.
\item
  In AI systems, tokenization affects context length, retrieval,
  prompting, and computational cost.
\end{itemize}

\section{Notes and references}\label{notes-and-references-1}

Classical treatments of text representations, tokenization, and
document-term matrices appear in information retrieval and statistical
NLP sources such as Manning et al. (2008), Manning and Schuetze (1999),
and Jurafsky and Martin (2024). Subword tokenization and byte-pair
encoding connect classical compression ideas with modern neural language
modeling.

\chapter{Documents as Vectors}\label{documents-as-vectors}

Term-document matrices, TF-IDF, and cosine geometry

\hfill\break

After text has been tokenized, the next classical step in NLP is to
represent each document by a vector. This chapter develops document
vectors from a mathematical point of view. We build count vectors,
term-document matrices, normalized frequency vectors, TF-IDF vectors,
and cosine similarity. The goal is not only to compute these objects,
but to understand what assumptions they encode and why they still matter
in modern AI systems such as search engines and retrieval-augmented
generation.

\section{Learning goals}\label{learning-goals-2}

After studying this chapter, you should be able to:

\begin{enumerate}
\def\labelenumi{\arabic{enumi}.}
\tightlist
\item
  define a document-vector representation as a map from token sequences
  to \(R^d\);
\item
  construct count vectors and document-term matrices from a corpus;
\item
  distinguish counts, binary indicators, relative frequencies, and
  normalized vectors;
\item
  explain TF-IDF as a mathematical weighting scheme for term
  specificity;
\item
  compute cosine similarity and interpret it geometrically;
\item
  rank documents in response to a query using sparse vector operations;
\item
  describe the statistical assumptions behind bag-of-words models;
\item
  implement document-vector methods in Python without relying on
  black-box NLP libraries;
\item
  use interactive visualizations to study how TF-IDF and normalization
  change retrieval;
\item
  use AI tools critically to test, explain, and improve vector-space
  representations.
\end{enumerate}

\section{Why documents become
vectors}\label{why-documents-become-vectors}

A document is naturally a sequence of tokens. If the vocabulary is

\[
V = \{v_1,v_2,\ldots,v_d\},
\]

then a tokenized document has the form

\[
x=(t_1,t_2,\ldots,t_m), \qquad t_k \in V.
\]

Many mathematical and statistical methods, however, operate on vectors
and matrices. The basic idea of this chapter is to replace a token
sequence by a vector

\[
\phi(x) \in R^d,
\]

where the \(j\)th coordinate measures how strongly token \(v_j\) appears
in the document.

\textbf{Central idea.} A document-vector model is a representation map

\[
\phi:V^* \to R^d.
\]

Different choices of \(\phi\) produce different mathematical models.
Count vectors, TF-IDF vectors, binary vectors, probability vectors, and
neural embeddings are all representation maps.

The simplest representation ignores word order and records only token
counts. Although this loses information, it produces a useful
finite-dimensional object. In classical NLP and information retrieval,
this is called the \textbf{vector space model}.

\section{Corpus, vocabulary, and
notation}\label{corpus-vocabulary-and-notation}

Let

\[
D = (x_1,x_2,\ldots,x_n)
\]

be a corpus of \(n\) tokenized documents. Each document \(x_i\) is a
finite sequence over the vocabulary \(V\).

We use the following notation throughout the chapter:

{\def\LTcaptype{none} % do not increment counter
\begin{longtable}[]{@{}ll@{}}
\toprule\noalign{}
Symbol & Meaning \\
\midrule\noalign{}
\endhead
\bottomrule\noalign{}
\endlastfoot
\(n\) & number of documents \\
\(d\) & vocabulary size \\
\(x_i\) & document \(i\) \\
\(v_j\) & vocabulary token \(j\) \\
\(c_{ij}\) & count of token \(v_j\) in document \(x_i\) \\
\(X\) & document-term matrix \\
\(q\) & query vector \\
\(s_i\) & similarity score between document \(i\) and query \\
\end{longtable}
}

The central matrix is

\[
X = (c_{ij}) \in R^{n \times d},
\]

where

\[
c_{ij}=\#\{k : t_{ik}=v_j\}.
\]

Here \(t_{ik}\) is the \(k\)th token in document \(x_i\).

\subsection{Row convention}\label{row-convention}

In this book we usually use the row convention:

\[
X_{i\cdot} = \text{vector representation of document } x_i.
\]

Thus each row is a document and each column is a term. Some books use
the transpose convention, where columns are documents. Both are valid,
but dimensions must be tracked carefully.

\textbf{Dimension check.} If \(X \in R^{n \times d}\) and a query vector
is \(q \in R^d\), then raw document-query scores have the form

\[
Xq \in R^n.
\]

The output contains one score for each document.

\section{Count vectors}\label{count-vectors}

The count vector of a document \(x\) is

\[
c(x)=(c_1(x),c_2(x),\ldots,c_d(x)) \in N^d,
\]

where

\[
c_j(x)=\#\{k:t_k=v_j\}.
\]

For example, suppose

\[
V=\{\text{data},\text{language},\text{learning},\text{models},\text{token}\}.
\]

If

\begin{Shaded}
\begin{Highlighting}[]
\NormalTok{data data learning models}
\end{Highlighting}
\end{Shaded}

is the document, then

\[
c(x)=(2,0,1,1,0).
\]

The count vector answers the question: \textbf{How many times does each
vocabulary term appear?}

\subsection{What counts preserve and
lose}\label{what-counts-preserve-and-lose}

The count vector preserves token frequency but discards order. For
example, the two documents

\begin{Shaded}
\begin{Highlighting}[]
\NormalTok{models predict language}
\end{Highlighting}
\end{Shaded}

and

\begin{Shaded}
\begin{Highlighting}[]
\NormalTok{language predict models}
\end{Highlighting}
\end{Shaded}

have the same count vector. Therefore, the count-vector map is not
injective.

\textbf{Bag-of-words assumption.} A count vector treats a document as a
multiset of tokens rather than as an ordered sequence. Mathematically,
this means the representation is invariant under permutations of the
token positions.

If \(\pi\) is a permutation of \(\{1,2,\ldots,m\}\) and

\[
x=(t_1,t_2,\ldots,t_m),
\]

then the permuted document

\[
\pi x=(t_{\pi(1)},t_{\pi(2)},\ldots,t_{\pi(m)})
\]

satisfies

\[
c(\pi x)=c(x).
\]

This invariance is both a strength and a weakness. It makes the model
simple and robust, but it cannot distinguish syntax, negation, or word
order.

\section{Binary vectors}\label{binary-vectors}

Sometimes we care only whether a token appears, not how often. The
binary representation is

\[
b_j(x)=
\begin{cases}
1, & c_j(x)>0,\\
0, & c_j(x)=0.
\end{cases}
\]

Thus

\[
b(x) \in \{0,1\}^d.
\]

Binary vectors are useful when repeated occurrences should not dominate.
For example, if one document repeats the word ``model'' many times, a
raw count vector may make it appear overly similar to any query
containing ``model.'' Binary vectors reduce this effect.

\section{Relative-frequency vectors}\label{relative-frequency-vectors}

Raw counts depend strongly on document length. If one document is twice
as long as another, its count vector may have a larger norm simply
because it contains more words. A simple correction is to divide by the
document length.

Let

\[
|x| = \sum_{j=1}^d c_j(x)
\]

be the total number of vocabulary tokens in document \(x\). The
relative-frequency vector is

\[
p_j(x)=\frac{c_j(x)}{|x|}, \qquad j=1,\ldots,d.
\]

Then

\[
p(x) \in R^d, \qquad p_j(x)\ge 0, \qquad \sum_{j=1}^d p_j(x)=1.
\]

Thus \(p(x)\) is an empirical probability distribution over the
vocabulary.

\textbf{Statistical viewpoint.} If a document is modeled as a sample
from a multinomial distribution over \(V\), then the relative-frequency
vector is the maximum likelihood estimate of the token probability
vector for that document.

This interpretation connects document vectors to statistics. A document
is not merely a string; it can be viewed as data generated by an unknown
distribution on words.

\section{The document-term matrix}\label{the-document-term-matrix}

Given \(n\) documents and \(d\) vocabulary terms, the count vectors form
a matrix

\[
X =
\begin{bmatrix}
--- & c(x_1)^T & ---\\
--- & c(x_2)^T & ---\\
& \vdots & \\
--- & c(x_n)^T & ---
\end{bmatrix}
\in R^{n\times d}.
\]

Equivalently,

\[
X_{ij}=c_j(x_i).
\]

This matrix is usually sparse because most documents contain only a
small fraction of all vocabulary terms.

\subsection{Example}\label{example}

Consider the corpus

\begin{enumerate}
\def\labelenumi{\arabic{enumi}.}
\tightlist
\item
  \texttt{machine\ learning\ models\ learn\ from\ data}
\item
  \texttt{language\ models\ predict\ the\ next\ token}
\item
  \texttt{statistics\ studies\ data\ uncertainty\ and\ models}
\end{enumerate}

Using the vocabulary

\[
V=\{\text{data},\text{language},\text{learning},\text{models},\text{predict},\text{statistics},\text{token}\},
\]

the document-term matrix is

\[
X=
\begin{bmatrix}
1&0&1&1&0&0&0\\
0&1&0&1&1&0&1\\
1&0&0&1&0&1&0
\end{bmatrix}.
\]

The first row says that the first document contains \texttt{data},
\texttt{learning}, and \texttt{models} once each.

\section{Python example: building a document-term
matrix}\label{python-example-building-a-document-term-matrix}

The following code builds a vocabulary and document-term matrix from
scratch. This is useful because it exposes the mathematical object
directly.

\protect\phantomsection\label{ch03-document-term-matrix}
\begin{Shaded}
\begin{Highlighting}[]
\ImportTok{import}\NormalTok{ re}
\ImportTok{import}\NormalTok{ numpy }\ImportTok{as}\NormalTok{ np}
\ImportTok{import}\NormalTok{ pandas }\ImportTok{as}\NormalTok{ pd}

\NormalTok{corpus }\OperatorTok{=}\NormalTok{ [}
    \StringTok{"machine learning models learn from data"}\NormalTok{,}
    \StringTok{"language models predict the next token"}\NormalTok{,}
    \StringTok{"statistics studies data uncertainty and models"}\NormalTok{,}
    \StringTok{"learning from data helps statistical models"}\NormalTok{,}
\NormalTok{]}

\KeywordTok{def}\NormalTok{ tokenize(text):}
\NormalTok{    text }\OperatorTok{=}\NormalTok{ text.lower()}
    \ControlFlowTok{return}\NormalTok{ re.findall(}\VerbatimStringTok{r"}\PreprocessorTok{[a{-}z]}\OperatorTok{+}\VerbatimStringTok{"}\NormalTok{, text)}

\NormalTok{tokenized }\OperatorTok{=}\NormalTok{ [tokenize(doc) }\ControlFlowTok{for}\NormalTok{ doc }\KeywordTok{in}\NormalTok{ corpus]}

\CommentTok{\# Build vocabulary in alphabetical order for reproducibility.}
\NormalTok{vocab }\OperatorTok{=} \BuiltInTok{sorted}\NormalTok{(}\BuiltInTok{set}\NormalTok{(token }\ControlFlowTok{for}\NormalTok{ doc }\KeywordTok{in}\NormalTok{ tokenized }\ControlFlowTok{for}\NormalTok{ token }\KeywordTok{in}\NormalTok{ doc))}
\NormalTok{token\_to\_id }\OperatorTok{=}\NormalTok{ \{token: j }\ControlFlowTok{for}\NormalTok{ j, token }\KeywordTok{in} \BuiltInTok{enumerate}\NormalTok{(vocab)\}}

\NormalTok{X }\OperatorTok{=}\NormalTok{ np.zeros((}\BuiltInTok{len}\NormalTok{(corpus), }\BuiltInTok{len}\NormalTok{(vocab)), dtype}\OperatorTok{=}\BuiltInTok{int}\NormalTok{)}
\ControlFlowTok{for}\NormalTok{ i, doc }\KeywordTok{in} \BuiltInTok{enumerate}\NormalTok{(tokenized):}
    \ControlFlowTok{for}\NormalTok{ token }\KeywordTok{in}\NormalTok{ doc:}
\NormalTok{        X[i, token\_to\_id[token]] }\OperatorTok{+=} \DecValTok{1}

\NormalTok{X\_df }\OperatorTok{=}\NormalTok{ pd.DataFrame(X, columns}\OperatorTok{=}\NormalTok{vocab, index}\OperatorTok{=}\NormalTok{[}\SpecialStringTok{f"doc\_}\SpecialCharTok{\{}\NormalTok{i}\OperatorTok{+}\DecValTok{1}\SpecialCharTok{\}}\SpecialStringTok{"} \ControlFlowTok{for}\NormalTok{ i }\KeywordTok{in} \BuiltInTok{range}\NormalTok{(}\BuiltInTok{len}\NormalTok{(corpus))])}
\BuiltInTok{print}\NormalTok{(X\_df)}
\end{Highlighting}
\end{Shaded}

The output matrix is the basic object used by many classical NLP
algorithms. Notice that the matrix construction requires only
tokenization, a vocabulary map, and counting.

\section{Linear algebra viewpoint}\label{linear-algebra-viewpoint}

Once documents are represented as vectors, we can use linear algebra.

Let

\[
x,y \in R^d
\]

be two document vectors. We may compute:

\begin{itemize}
\tightlist
\item
  their dot product \(x^Ty\);
\item
  their norms \(\|x\|_1\), \(\|x\|_2\), or \(\|x\|_\infty\);
\item
  their angle through cosine similarity;
\item
  their distance, such as \(\|x-y\|_2\);
\item
  their projection onto a low-dimensional subspace.
\end{itemize}

The dot product is

\[
x^Ty=\sum_{j=1}^d x_jy_j.
\]

For count vectors, this quantity is large when the two documents share
many terms with high counts.

However, dot products are sensitive to length. A longer document tends
to have a larger norm and may receive a large dot product even when it
is not especially similar in content.

\section{Cosine similarity}\label{cosine-similarity}

Cosine similarity normalizes by vector length:

\[
\operatorname{cos}(x,y)=\frac{x^Ty}{\|x\|_2\|y\|_2},
\]

provided \(x\ne 0\) and \(y\ne 0\).

Geometrically, this is the cosine of the angle between \(x\) and \(y\).
It measures direction rather than magnitude. If \(x\) and \(y\) point in
the same direction, then their cosine similarity is \(1\). If they are
orthogonal, the similarity is \(0\).

For nonnegative document vectors, cosine similarity lies between \(0\)
and \(1\).

\textbf{Scale invariance.} For any \(a>0\) and \(b>0\),

\[
\operatorname{cos}(ax,by)=\operatorname{cos}(x,y).
\]

Therefore, cosine similarity ignores positive rescaling of documents.
This is useful when document lengths differ.

\subsection{Matrix form of cosine
retrieval}\label{matrix-form-of-cosine-retrieval}

Suppose \(X\in R^{n\times d}\) contains document vectors and
\(q\in R^d\) is a query vector. The raw dot-product scores are

\[
r=Xq \in R^n.
\]

The cosine scores are

\[
s_i=\frac{X_{i\cdot}q}{\|X_{i\cdot}\|_2\|q\|_2}, \qquad i=1,\ldots,n.
\]

If we define the row-normalized matrix

\[
\widetilde X_{i\cdot}=\frac{X_{i\cdot}}{\|X_{i\cdot}\|_2}
\]

and the normalized query

\[
\widetilde q=\frac{q}{\|q\|_2},
\]

then

\[
s=\widetilde X\widetilde q.
\]

Thus retrieval by cosine similarity is matrix-vector multiplication
after normalization.

\section{Python example: cosine
search}\label{python-example-cosine-search}

\protect\phantomsection\label{ch03-cosine-search}
\begin{Shaded}
\begin{Highlighting}[]
\ImportTok{import}\NormalTok{ numpy }\ImportTok{as}\NormalTok{ np}
\ImportTok{import}\NormalTok{ pandas }\ImportTok{as}\NormalTok{ pd}
\ImportTok{import}\NormalTok{ re}

\NormalTok{corpus }\OperatorTok{=}\NormalTok{ [}
    \StringTok{"machine learning models learn from data"}\NormalTok{,}
    \StringTok{"language models predict the next token"}\NormalTok{,}
    \StringTok{"statistics studies data uncertainty and models"}\NormalTok{,}
    \StringTok{"learning from data helps statistical models"}\NormalTok{,}
\NormalTok{]}

\KeywordTok{def}\NormalTok{ tokenize(text):}
    \ControlFlowTok{return}\NormalTok{ re.findall(}\VerbatimStringTok{r"}\PreprocessorTok{[a{-}z]}\OperatorTok{+}\VerbatimStringTok{"}\NormalTok{, text.lower())}

\NormalTok{tokenized }\OperatorTok{=}\NormalTok{ [tokenize(doc) }\ControlFlowTok{for}\NormalTok{ doc }\KeywordTok{in}\NormalTok{ corpus]}
\NormalTok{vocab }\OperatorTok{=} \BuiltInTok{sorted}\NormalTok{(}\BuiltInTok{set}\NormalTok{(token }\ControlFlowTok{for}\NormalTok{ doc }\KeywordTok{in}\NormalTok{ tokenized }\ControlFlowTok{for}\NormalTok{ token }\KeywordTok{in}\NormalTok{ doc))}
\NormalTok{token\_to\_id }\OperatorTok{=}\NormalTok{ \{token: j }\ControlFlowTok{for}\NormalTok{ j, token }\KeywordTok{in} \BuiltInTok{enumerate}\NormalTok{(vocab)\}}

\KeywordTok{def}\NormalTok{ vectorize(tokens):}
\NormalTok{    vec }\OperatorTok{=}\NormalTok{ np.zeros(}\BuiltInTok{len}\NormalTok{(vocab))}
    \ControlFlowTok{for}\NormalTok{ token }\KeywordTok{in}\NormalTok{ tokens:}
        \ControlFlowTok{if}\NormalTok{ token }\KeywordTok{in}\NormalTok{ token\_to\_id:}
\NormalTok{            vec[token\_to\_id[token]] }\OperatorTok{+=} \DecValTok{1}
    \ControlFlowTok{return}\NormalTok{ vec}

\NormalTok{X }\OperatorTok{=}\NormalTok{ np.vstack([vectorize(doc) }\ControlFlowTok{for}\NormalTok{ doc }\KeywordTok{in}\NormalTok{ tokenized])}

\NormalTok{query }\OperatorTok{=} \StringTok{"data learning models"}
\NormalTok{q }\OperatorTok{=}\NormalTok{ vectorize(tokenize(query))}

\KeywordTok{def}\NormalTok{ cosine\_similarity\_matrix(X, q):}
\NormalTok{    q\_norm }\OperatorTok{=}\NormalTok{ np.linalg.norm(q)}
\NormalTok{    X\_norms }\OperatorTok{=}\NormalTok{ np.linalg.norm(X, axis}\OperatorTok{=}\DecValTok{1}\NormalTok{)}
\NormalTok{    scores }\OperatorTok{=}\NormalTok{ X }\OperatorTok{@}\NormalTok{ q}
\NormalTok{    scores }\OperatorTok{=}\NormalTok{ scores }\OperatorTok{/}\NormalTok{ (X\_norms }\OperatorTok{*}\NormalTok{ q\_norm)}
    \ControlFlowTok{return}\NormalTok{ scores}

\NormalTok{scores }\OperatorTok{=}\NormalTok{ cosine\_similarity\_matrix(X, q)}
\NormalTok{ranking }\OperatorTok{=}\NormalTok{ np.argsort(}\OperatorTok{{-}}\NormalTok{scores)}

\ControlFlowTok{for}\NormalTok{ rank, i }\KeywordTok{in} \BuiltInTok{enumerate}\NormalTok{(ranking, start}\OperatorTok{=}\DecValTok{1}\NormalTok{):}
    \BuiltInTok{print}\NormalTok{(}\SpecialStringTok{f"}\SpecialCharTok{\{}\NormalTok{rank}\SpecialCharTok{\}}\SpecialStringTok{. score=}\SpecialCharTok{\{}\NormalTok{scores[i]}\SpecialCharTok{:.3f\}}\SpecialStringTok{ | }\SpecialCharTok{\{}\NormalTok{corpus[i]}\SpecialCharTok{\}}\SpecialStringTok{"}\NormalTok{)}
\end{Highlighting}
\end{Shaded}

This code shows that document retrieval can be written as a small linear
algebra problem.

\section{Term frequency}\label{term-frequency}

Raw counts can be transformed before similarity is computed. The first
transformation is \textbf{term frequency}.

A simple term frequency is

\[
\operatorname{tf}_{ij}=c_{ij}.
\]

A length-normalized term frequency is

\[
\operatorname{tf}_{ij}=\frac{c_{ij}}{\sum_{k=1}^d c_{ik}}.
\]

A log-scaled term frequency is

\[
\operatorname{tf}_{ij}=
\begin{cases}
1+\log(c_{ij}), & c_{ij}>0,\\
0, & c_{ij}=0.
\end{cases}
\]

The log transform reduces the effect of repeated occurrences. For
example, a term appearing ten times should not always be ten times as
important as a term appearing once.

\section{Document frequency and inverse document
frequency}\label{document-frequency-and-inverse-document-frequency}

Some words appear in many documents. These common terms often have low
discriminative value. For example, words such as \texttt{the},
\texttt{and}, \texttt{model}, or \texttt{data} may appear frequently in
a technical corpus.

The document frequency of token \(v_j\) is

\[
\operatorname{df}_j=\#\{i:c_{ij}>0\}.
\]

It counts how many documents contain the term at least once.

The inverse document frequency is commonly defined as

\[
\operatorname{idf}_j=\log\left(\frac{n+1}{\operatorname{df}_j+1}\right)+1.
\]

The added \(1\) in numerator and denominator is smoothing. It avoids
division by zero and keeps the weight finite. Other conventions are also
used.

\textbf{Interpretation.} If \(\operatorname{df}_j/n\) is the empirical
probability that a random document contains token \(v_j\), then

\[
\log\left(\frac{n}{\operatorname{df}_j}\right)
\]

is the negative logarithm of this empirical probability. Rare terms
receive larger weights because they are more informative for
distinguishing documents.

\section{TF-IDF vectors}\label{tf-idf-vectors}

The TF-IDF weight of term \(v_j\) in document \(x_i\) is

\[
w_{ij}=\operatorname{tf}_{ij}\operatorname{idf}_j.
\]

In matrix form, if \(T=(\operatorname{tf}_{ij})\) and

\[
D_{\operatorname{idf}}=\operatorname{diag}(\operatorname{idf}_1,\ldots,\operatorname{idf}_d),
\]

then the TF-IDF matrix is

\[
W=TD_{\operatorname{idf}}.
\]

This is a column weighting of the term-frequency matrix. Terms that
appear in many documents get smaller weights. Terms that appear in fewer
documents get larger weights.

\subsection{Why TF-IDF is useful}\label{why-tf-idf-is-useful}

TF-IDF balances two effects:

\begin{enumerate}
\def\labelenumi{\arabic{enumi}.}
\tightlist
\item
  \textbf{Within-document importance.} A term is important if it appears
  often in a particular document.
\item
  \textbf{Across-document specificity.} A term is more informative if it
  does not appear in almost every document.
\end{enumerate}

Thus TF-IDF is a simple way to say:

\begin{quote}
A good keyword should be frequent in this document but not too frequent
across the whole corpus.
\end{quote}

\section{Python example: TF-IDF from
scratch}\label{python-example-tf-idf-from-scratch}

\protect\phantomsection\label{ch03-tfidf-from-scratch}
\begin{Shaded}
\begin{Highlighting}[]
\ImportTok{import}\NormalTok{ numpy }\ImportTok{as}\NormalTok{ np}
\ImportTok{import}\NormalTok{ pandas }\ImportTok{as}\NormalTok{ pd}

\CommentTok{\# X is the document{-}term count matrix from the previous example.}
\NormalTok{n, d }\OperatorTok{=}\NormalTok{ X.shape}

\CommentTok{\# Term frequency: length{-}normalized counts.}
\NormalTok{doc\_lengths }\OperatorTok{=}\NormalTok{ X.}\BuiltInTok{sum}\NormalTok{(axis}\OperatorTok{=}\DecValTok{1}\NormalTok{, keepdims}\OperatorTok{=}\VariableTok{True}\NormalTok{)}
\NormalTok{TF }\OperatorTok{=}\NormalTok{ X }\OperatorTok{/}\NormalTok{ doc\_lengths}

\CommentTok{\# Document frequency.}
\NormalTok{df }\OperatorTok{=}\NormalTok{ (X }\OperatorTok{\textgreater{}} \DecValTok{0}\NormalTok{).}\BuiltInTok{sum}\NormalTok{(axis}\OperatorTok{=}\DecValTok{0}\NormalTok{)}

\CommentTok{\# Smoothed inverse document frequency.}
\NormalTok{idf }\OperatorTok{=}\NormalTok{ np.log((n }\OperatorTok{+} \DecValTok{1}\NormalTok{) }\OperatorTok{/}\NormalTok{ (df }\OperatorTok{+} \DecValTok{1}\NormalTok{)) }\OperatorTok{+} \DecValTok{1}

\CommentTok{\# TF{-}IDF matrix.}
\NormalTok{W }\OperatorTok{=}\NormalTok{ TF }\OperatorTok{*}\NormalTok{ idf}

\NormalTok{W\_df }\OperatorTok{=}\NormalTok{ pd.DataFrame(W, columns}\OperatorTok{=}\NormalTok{vocab, index}\OperatorTok{=}\NormalTok{[}\SpecialStringTok{f"doc\_}\SpecialCharTok{\{}\NormalTok{i}\OperatorTok{+}\DecValTok{1}\SpecialCharTok{\}}\SpecialStringTok{"} \ControlFlowTok{for}\NormalTok{ i }\KeywordTok{in} \BuiltInTok{range}\NormalTok{(n)])}
\BuiltInTok{print}\NormalTok{(W\_df.}\BuiltInTok{round}\NormalTok{(}\DecValTok{3}\NormalTok{))}
\end{Highlighting}
\end{Shaded}

The resulting matrix has the same shape as the count matrix, but its
entries have changed. Common tokens receive smaller column weights,
while more specific tokens receive larger column weights.

\section{Comparing count vectors and TF-IDF
vectors}\label{comparing-count-vectors-and-tf-idf-vectors}

Count vectors often reward documents that share common terms with the
query. TF-IDF vectors reward documents that share more specific terms.

Suppose the query is

\begin{Shaded}
\begin{Highlighting}[]
\NormalTok{language token prediction}
\end{Highlighting}
\end{Shaded}

A count-vector search might rank a document highly because it contains
\texttt{language} and \texttt{token}. TF-IDF may increase the importance
of \texttt{token} if it appears in only a few documents.

This distinction matters in academic settings. For example, in a corpus
of course descriptions, words such as \texttt{course},
\texttt{students}, and \texttt{methods} may appear everywhere. Words
such as \texttt{Bayesian}, \texttt{topology}, or \texttt{transformer}
may be more specific.

\section{Interactive HTML: TF-IDF and cosine retrieval
explorer}\label{interactive-html-tf-idf-and-cosine-retrieval-explorer}

The following interactive component lets you adjust the IDF strength and
choose a query. It uses a tiny corpus so that the mathematics remains
visible. Move the slider to see how document rankings change.

The HTML version of this chapter contains an interactive Plotly.js
component for exploring TF-IDF weighting and cosine retrieval. In PDF
output, the same ideas can be studied by running the Python code in the
associated lab notebook.

\section{The geometry of sparse document
vectors}\label{the-geometry-of-sparse-document-vectors}

Document vectors often live in very high dimensions. A vocabulary may
contain thousands or millions of tokens, while a single document uses
only a small subset. Therefore, most entries of a document vector are
zero.

This creates a sparse geometry:

\[
\#\{j:x_j\ne 0\} \ll d.
\]

Sparse vectors are computationally useful because we do not need to
store or multiply by all zero entries. They are also geometrically
interesting. Two randomly chosen short documents may have no terms in
common, which makes their dot product zero.

\subsection{Orthogonality in sparse
spaces}\label{orthogonality-in-sparse-spaces}

If two count vectors have disjoint supports, then

\[
x^Ty=0.
\]

For nonnegative vectors, this means their cosine similarity is zero. In
a large vocabulary, many pairs of documents may be nearly orthogonal
even if they are topically related but use different words.

For example, the documents

\begin{Shaded}
\begin{Highlighting}[]
\NormalTok{neural network training}
\end{Highlighting}
\end{Shaded}

and

\begin{Shaded}
\begin{Highlighting}[]
\NormalTok{deep learning optimization}
\end{Highlighting}
\end{Shaded}

are semantically related, but a simple bag-of-words model may give low
similarity if the exact words do not overlap. This limitation motivates
later chapters on latent semantic analysis, word embeddings, and
transformer representations.

\section{Stop words and projection}\label{stop-words-and-projection}

A common preprocessing step removes very frequent words such as
\texttt{the}, \texttt{and}, \texttt{of}, and \texttt{to}.
Mathematically, stop-word removal is a coordinate projection.

Let \(S\subseteq \{1,2,\ldots,d\}\) be the set of coordinates
corresponding to terms we keep. The projection map is

\[
P_S:R^d\to R^{|S|},
\]

where

\[
P_S(x)=(x_j)_{j\in S}.
\]

This map discards coordinates. It can improve retrieval by removing
noise, but it can also remove meaningful words in some contexts. For
example, \texttt{not} is often a stop word in generic lists, but it is
crucial for sentiment analysis.

\textbf{Modeling warning.} Stop-word removal is not universally safe. It
is a projection that may remove coordinates relevant to the task.

\section{Statistical model behind
bag-of-words}\label{statistical-model-behind-bag-of-words}

The bag-of-words representation can be derived from a multinomial model.

Suppose document \(x_i\) has length \(m_i\) and is generated by
independent draws from a categorical distribution

\[
\theta_i=(\theta_{i1},\ldots,\theta_{id}),
\]

where

\[
\theta_{ij}\ge 0, \qquad \sum_{j=1}^d \theta_{ij}=1.
\]

The probability of observing counts

\[
c_i=(c_{i1},\ldots,c_{id})
\]

is

\[
P(c_i\mid \theta_i)=\frac{m_i!}{c_{i1}!c_{i2}!\cdots c_{id}!}\prod_{j=1}^d \theta_{ij}^{c_{ij}}.
\]

The log-likelihood, ignoring the multinomial coefficient, is

\[
\ell(\theta_i)=\sum_{j=1}^d c_{ij}\log \theta_{ij}.
\]

Maximizing this subject to \(\sum_j \theta_{ij}=1\) gives

\[
\widehat\theta_{ij}=\frac{c_{ij}}{m_i}.
\]

Thus the relative-frequency vector is the maximum likelihood estimate of
the document's token distribution under this simplified model.

\subsection{What the model assumes}\label{what-the-model-assumes}

The multinomial bag-of-words model assumes:

\begin{enumerate}
\def\labelenumi{\arabic{enumi}.}
\tightlist
\item
  token order is irrelevant;
\item
  tokens are conditionally independent given \(\theta_i\);
\item
  each document has its own word distribution;
\item
  dependence between neighboring words is ignored.
\end{enumerate}

These assumptions are false for real language, but they give a
mathematically transparent baseline.

\section{Connection to
classification}\label{connection-to-classification}

Document vectors are not only used for retrieval. They can also be used
for classification.

Suppose each document has a label

\[
y_i\in \{1,2,\ldots,K\}.
\]

A linear classifier has the form

\[
f(x_i)=\arg\max_{k\in\{1,\ldots,K\}} \left(w_k^T\phi(x_i)+b_k\right).
\]

Here \(\phi(x_i)\) may be a count vector, TF-IDF vector, or another
document representation. The same vectorization step supports logistic
regression, support vector machines, naive Bayes, and neural networks.

For binary classification, a common model is

\[
P(y_i=1\mid x_i)=\sigma(w^T\phi(x_i)+b),
\]

where

\[
\sigma(z)=\frac{1}{1+e^{-z}}.
\]

This shows how vector-space NLP connects directly to statistics and
machine learning.

\section{Low-dimensional visualization by
SVD}\label{low-dimensional-visualization-by-svd}

Because document vectors are high-dimensional, we often visualize them
by projecting to two dimensions. A simple method is singular value
decomposition.

Let \(W\in R^{n\times d}\) be a centered or weighted document-term
matrix. Its singular value decomposition is

\[
W=U\Sigma V^T.
\]

The best rank-\(k\) approximation in Frobenius norm is

\[
W_k=U_k\Sigma_kV_k^T.
\]

For visualization, we often use the rows of

\[
U_2\Sigma_2\in R^{n\times 2}
\]

as two-dimensional document coordinates.

\textbf{Preview.} This idea leads to latent semantic analysis, which
studies low-rank structure in document-term matrices. Chapter 7 returns
to this topic in detail.

\section{Python example: two-dimensional document
map}\label{python-example-two-dimensional-document-map}

\protect\phantomsection\label{ch03-svd-visualization}
\begin{Shaded}
\begin{Highlighting}[]
\ImportTok{import}\NormalTok{ numpy }\ImportTok{as}\NormalTok{ np}
\ImportTok{import}\NormalTok{ matplotlib.pyplot }\ImportTok{as}\NormalTok{ plt}

\CommentTok{\# Use the TF{-}IDF matrix W from the earlier example.}
\CommentTok{\# Center columns so the SVD highlights differences among documents.}
\NormalTok{W\_centered }\OperatorTok{=}\NormalTok{ W }\OperatorTok{{-}}\NormalTok{ W.mean(axis}\OperatorTok{=}\DecValTok{0}\NormalTok{, keepdims}\OperatorTok{=}\VariableTok{True}\NormalTok{)}

\NormalTok{U, S, Vt }\OperatorTok{=}\NormalTok{ np.linalg.svd(W\_centered, full\_matrices}\OperatorTok{=}\VariableTok{False}\NormalTok{)}
\NormalTok{coords }\OperatorTok{=}\NormalTok{ U[:, :}\DecValTok{2}\NormalTok{] }\OperatorTok{*}\NormalTok{ S[:}\DecValTok{2}\NormalTok{]}

\NormalTok{plt.figure(figsize}\OperatorTok{=}\NormalTok{(}\DecValTok{6}\NormalTok{, }\DecValTok{4}\NormalTok{))}
\NormalTok{plt.scatter(coords[:, }\DecValTok{0}\NormalTok{], coords[:, }\DecValTok{1}\NormalTok{])}
\ControlFlowTok{for}\NormalTok{ i, text }\KeywordTok{in} \BuiltInTok{enumerate}\NormalTok{(corpus):}
\NormalTok{    plt.annotate(}\SpecialStringTok{f"doc\_}\SpecialCharTok{\{}\NormalTok{i}\OperatorTok{+}\DecValTok{1}\SpecialCharTok{\}}\SpecialStringTok{"}\NormalTok{, (coords[i, }\DecValTok{0}\NormalTok{], coords[i, }\DecValTok{1}\NormalTok{]))}
\NormalTok{plt.axhline(}\DecValTok{0}\NormalTok{, linewidth}\OperatorTok{=}\FloatTok{0.8}\NormalTok{)}
\NormalTok{plt.axvline(}\DecValTok{0}\NormalTok{, linewidth}\OperatorTok{=}\FloatTok{0.8}\NormalTok{)}
\NormalTok{plt.xlabel(}\StringTok{"first SVD coordinate"}\NormalTok{)}
\NormalTok{plt.ylabel(}\StringTok{"second SVD coordinate"}\NormalTok{)}
\NormalTok{plt.title(}\StringTok{"Documents projected to two dimensions"}\NormalTok{)}
\NormalTok{plt.show()}
\end{Highlighting}
\end{Shaded}

This visualization is not a replacement for the full high-dimensional
representation. It is a diagnostic tool that helps us inspect whether
documents cluster by topic or vocabulary.

\section{Evaluation of retrieval}\label{evaluation-of-retrieval}

A document-vector method is often used for retrieval: given a query,
return a ranked list of documents. To evaluate retrieval, we need
relevance judgments.

For a query \(q\), let

\[
R(q)\subseteq \{1,2,\ldots,n\}
\]

be the set of relevant documents. If a system returns the top \(k\)
documents

\[
A_k(q),
\]

then precision at \(k\) is

\[
\operatorname{Precision@}k=\frac{|A_k(q)\cap R(q)|}{k}.
\]

Recall at \(k\) is

\[
\operatorname{Recall@}k=\frac{|A_k(q)\cap R(q)|}{|R(q)|}.
\]

Precision measures how many retrieved documents are relevant. Recall
measures how many relevant documents were retrieved.

For course search, advising chatbots, and RAG systems, retrieval
evaluation is essential. A language model can only answer from useful
evidence if the retrieval step finds useful evidence.

\section{AI component: document vectors in RAG
systems}\label{ai-component-document-vectors-in-rag-systems}

Modern retrieval-augmented generation systems often use neural
embeddings rather than TF-IDF. However, the mathematical pipeline is
similar:

\[
\text{document} \xrightarrow{\phi} R^p,
\qquad
\text{query} \xrightarrow{\phi} R^p,
\qquad
\text{rank by similarity}.
\]

In classical vector search, \(\phi\) may be TF-IDF. In neural vector
search, \(\phi\) may be a transformer embedding model. The retrieval
step still depends on geometry: dot products, cosine similarity, nearest
neighbors, and ranking.

\textbf{AI-assisted learning activity.} Ask an AI system:

\begin{quote}
Give three short documents for which TF-IDF search succeeds and three
short documents for which TF-IDF search fails. For each failure, explain
the mathematical reason in terms of vocabulary overlap, synonymy,
negation, or word order.
\end{quote}

Then verify the examples by computing count vectors and cosine
similarities yourself.

\section{AI component: auditing a
representation}\label{ai-component-auditing-a-representation}

When using AI tools to help build NLP systems, ask the tool to make the
representation assumptions explicit. A useful prompt is:

\begin{quote}
I am representing documents using TF-IDF vectors. List the mathematical
assumptions of this representation. Which assumptions are problematic
for sentiment analysis, course advising, and mathematical text search?
\end{quote}

A good answer should mention at least:

\begin{itemize}
\tightlist
\item
  loss of word order;
\item
  dependence on exact vocabulary overlap;
\item
  sensitivity to tokenization;
\item
  inability to understand synonyms unless they share terms;
\item
  possible removal of important stop words;
\item
  limited handling of negation;
\item
  dependence on the corpus used to compute IDF.
\end{itemize}

Use AI tools as mathematical critics, not as replacements for
mathematical verification.

\section{Worked example: hand
calculation}\label{worked-example-hand-calculation}

Consider three documents with vocabulary

\[
V=\{a,b,c,d\}.
\]

Suppose the count matrix is

\[
X=
\begin{bmatrix}
2&1&0&0\\
1&0&2&0\\
0&1&0&3
\end{bmatrix}.
\]

Let the query be

\[
q=(1,1,0,0)^T.
\]

The raw dot-product scores are

\[
Xq=
\begin{bmatrix}
3\\
1\\
1
\end{bmatrix}.
\]

The Euclidean norms are

\[
\|q\|_2=\sqrt{2},
\]

\[
\|X_{1\cdot}\|_2=\sqrt{5}, \qquad
\|X_{2\cdot}\|_2=\sqrt{5}, \qquad
\|X_{3\cdot}\|_2=\sqrt{10}.
\]

Thus the cosine scores are

\[
s_1=\frac{3}{\sqrt{5}\sqrt{2}}=\frac{3}{\sqrt{10}},
\]

\[
s_2=\frac{1}{\sqrt{5}\sqrt{2}}=\frac{1}{\sqrt{10}},
\]

and

\[
s_3=\frac{1}{\sqrt{10}\sqrt{2}}=\frac{1}{\sqrt{20}}.
\]

Therefore, document 1 is most similar to the query.

This small example illustrates the full retrieval calculation:
vectorize, multiply, normalize, rank.

\section{Common pitfalls}\label{common-pitfalls}

\subsection{Pitfall 1: confusing rows and
columns}\label{pitfall-1-confusing-rows-and-columns}

Always state whether documents are rows or columns. In this book,
documents are rows.

\subsection{Pitfall 2: forgetting zero
vectors}\label{pitfall-2-forgetting-zero-vectors}

Cosine similarity is undefined when one vector has zero norm. This can
happen if a query contains no vocabulary terms.

\subsection{Pitfall 3: assuming TF-IDF understands
meaning}\label{pitfall-3-assuming-tf-idf-understands-meaning}

TF-IDF is a lexical method. It does not know that \texttt{doctor} and
\texttt{physician} are related unless the corpus and vocabulary create
indirect associations.

\subsection{Pitfall 4: ignoring corpus
dependence}\label{pitfall-4-ignoring-corpus-dependence}

IDF weights depend on the corpus. If the corpus changes, the IDF weights
change.

\subsection{Pitfall 5: treating preprocessing as
neutral}\label{pitfall-5-treating-preprocessing-as-neutral}

Lowercasing, punctuation removal, stop-word removal, and stemming all
change the mathematical representation.

\section{Chapter summary}\label{chapter-summary-1}

The main construction in this chapter is the document-vector map

\[
\phi:V^*\to R^d.
\]

The simplest choice is the count vector. From count vectors we build the
document-term matrix

\[
X\in R^{n\times d}.
\]

We then transform this matrix using normalization and TF-IDF weighting.
Similarity between documents or between a query and documents is often
measured by cosine similarity:

\[
\operatorname{cos}(x,y)=\frac{x^Ty}{\|x\|_2\|y\|_2}.
\]

These ideas form the classical foundation of search and retrieval. They
also prepare us for modern NLP, where documents and queries are
represented by learned dense vectors instead of sparse TF-IDF vectors.

\section{Exercises}\label{exercises-2}

\subsection{Conceptual exercises}\label{conceptual-exercises-1}

\begin{enumerate}
\def\labelenumi{\arabic{enumi}.}
\tightlist
\item
  Explain why the count-vector map \(c:V^*\to N^d\) is not injective.
\item
  Give two documents with the same count vector but different meanings.
\item
  Explain why cosine similarity is often preferable to raw dot product
  for document retrieval.
\item
  Give an example where stop-word removal changes the meaning of a
  document.
\item
  Explain why IDF weights depend on the corpus, not only on the
  document.
\end{enumerate}

\subsection{Computational exercises}\label{computational-exercises}

\begin{enumerate}
\def\labelenumi{\arabic{enumi}.}
\setcounter{enumi}{5}
\tightlist
\item
  For the vocabulary \(V=\{a,b,c,d\}\) and document \(x=(a,b,a,d,d,d)\),
  compute the count vector, binary vector, and relative-frequency
  vector.
\item
  Given
\end{enumerate}

\[
X=
\begin{bmatrix}
1&2&0\\
0&1&3\\
1&0&1
\end{bmatrix},
\qquad
q=(1,1,0)^T,
\]

compute the raw scores \(Xq\) and the cosine similarity scores.

\begin{enumerate}
\def\labelenumi{\arabic{enumi}.}
\setcounter{enumi}{7}
\tightlist
\item
  For a corpus with \(n=10\) documents, suppose a term appears in \(2\)
  documents. Compute the smoothed IDF
\end{enumerate}

\[
\log\left(\frac{n+1}{\operatorname{df}+1}\right)+1.
\]

\begin{enumerate}
\def\labelenumi{\arabic{enumi}.}
\setcounter{enumi}{8}
\tightlist
\item
  Build a document-term matrix for five short course descriptions of
  your choice.
\item
  Compare rankings from raw counts and TF-IDF for the same query.
\end{enumerate}

\subsection{Proof-based exercises}\label{proof-based-exercises}

\begin{enumerate}
\def\labelenumi{\arabic{enumi}.}
\setcounter{enumi}{10}
\tightlist
\item
  Prove that cosine similarity is invariant under positive rescaling of
  both vectors.
\item
  Prove that if two nonnegative vectors have disjoint supports, then
  their cosine similarity is \(0\).
\item
  Show that the relative-frequency vector is the maximum likelihood
  estimate for a multinomial token distribution.
\item
  Let \(D_{\operatorname{idf}}\) be a positive diagonal matrix. Explain
  geometrically how multiplying \(X\) by \(D_{\operatorname{idf}}\)
  changes the coordinate system.
\end{enumerate}

\subsection{AI-assisted exercises}\label{ai-assisted-exercises}

\begin{enumerate}
\def\labelenumi{\arabic{enumi}.}
\setcounter{enumi}{14}
\tightlist
\item
  Ask an AI system to generate a corpus of six documents about
  mathematics, statistics, and machine learning. Build a TF-IDF matrix
  for the corpus and identify which terms receive the largest IDF
  weights.
\item
  Ask an AI system for an example where TF-IDF fails because two
  documents use synonyms. Verify the failure numerically using cosine
  similarity.
\item
  Ask an AI system to explain why \texttt{not} should not always be
  removed as a stop word. Translate the explanation into a statement
  about coordinate projection.
\end{enumerate}

\section{Mini-project: a small mathematical search
engine}\label{mini-project-a-small-mathematical-search-engine}

Build a small search engine for a topic of your choice, such as graduate
course descriptions, research abstracts, or textbook chapter summaries.

Your project should include:

\begin{enumerate}
\def\labelenumi{\arabic{enumi}.}
\tightlist
\item
  a corpus of at least \(20\) short documents;
\item
  a documented tokenizer;
\item
  a vocabulary list;
\item
  a count matrix;
\item
  a TF-IDF matrix;
\item
  cosine-similarity retrieval for at least five queries;
\item
  a short analysis comparing successful and unsuccessful searches;
\item
  one AI-assisted critique of the representation.
\end{enumerate}

The final report should explain the mathematical representation before
showing code.

\section{Python lab connection}\label{python-lab-connection-1}

The associated lab for this chapter is planned as:

\begin{Shaded}
\begin{Highlighting}[]
\NormalTok{labs/chapter{-}03{-}documents{-}as{-}vectors.ipynb}
\end{Highlighting}
\end{Shaded}

The lab should be written for independent study and should include:

\begin{itemize}
\tightlist
\item
  background on sparse vectors and document-term matrices;
\item
  count, binary, relative-frequency, and TF-IDF representations;
\item
  cosine similarity from scratch;
\item
  document retrieval experiments;
\item
  visualization using SVD;
\item
  comparison with a modern embedding-based retrieval method as an
  optional extension.
\end{itemize}

\section{Notes and references}\label{notes-and-references-2}

The vector space model and TF-IDF are classical foundations of
information retrieval. For more detail, see Manning et al. (2008) and
Manning and Schuetze (1999). Modern NLP still uses the same underlying
linear algebra ideas, even when the representation map is learned by a
neural network rather than manually constructed. Later chapters connect
these ideas to latent semantic analysis, word embeddings, and
transformer-based retrieval.

\part{Part II. Probability Models for Language}

\chapter{Probability and Language}\label{probability-and-language}

Random tokens, conditional distributions, likelihood, entropy, and
perplexity

\hfill\break

Natural language has uncertainty. A sentence may continue in many
possible ways, a word may have several meanings, and the same document
may be assigned different labels depending on context. This chapter
develops the probabilistic viewpoint needed for modern NLP. We treat
tokens, sentences, labels, and model outputs as random objects. The goal
is to understand language modeling as conditional probability estimation
and supervised NLP as learning probability distributions from data.

\section{Learning goals}\label{learning-goals-3}

After studying this chapter, you should be able to:

\begin{enumerate}
\def\labelenumi{\arabic{enumi}.}
\tightlist
\item
  describe a token sequence as a realization of random variables over a
  finite vocabulary;
\item
  distinguish empirical token frequencies, population probabilities, and
  model probabilities;
\item
  use the chain rule of probability to factor the probability of a
  sentence;
\item
  explain why language modeling is conditional distribution estimation;
\item
  derive maximum likelihood estimation for a unigram model;
\item
  compute negative log-likelihood, cross-entropy, entropy, KL
  divergence, and perplexity;
\item
  explain the difference between generative and discriminative NLP
  models;
\item
  implement simple probability models for text using Python;
\item
  use an interactive Plotly.js visualization to explore entropy,
  temperature, and next-token probabilities;
\item
  use AI tools to critique probabilistic assumptions without treating
  the AI response as a proof.
\end{enumerate}

\section{Why probability enters NLP}\label{why-probability-enters-nlp}

In the previous chapters, we treated text as data: strings, token
sequences, count vectors, and document-term matrices. These
representations are useful, but they do not yet answer questions such
as:

\begin{itemize}
\tightlist
\item
  How likely is a sentence?
\item
  Which word is most likely to come next?
\item
  How uncertain is a model's prediction?
\item
  How should a classifier combine evidence from many words?
\item
  How should we compare two language models statistically?
\end{itemize}

Probability provides the mathematical language for these questions.

Suppose the vocabulary is

\[
V=\{v_1,v_2,\ldots,v_d\}.
\]

A single token can be modeled as a random variable

\[
W:\Omega \to V,
\]

where \(\Omega\) is an underlying sample space. For a sentence or
document of length \(T\), we use a sequence of random variables

\[
W_1,W_2,\ldots,W_T,
\]

with each \(W_t\) taking values in \(V\). A particular observed sentence
is a realization

\[
w_{1:T}=(w_1,w_2,\ldots,w_T).
\]

\textbf{Central idea.} A probabilistic language model assigns
probabilities to token sequences:

\[
p(w_{1:T})=P(W_1=w_1,\ldots,W_T=w_T).
\]

Modern large language models usually estimate conditional probabilities
of the form

\[
p(w_t \mid w_1,\ldots,w_{t-1}).
\]

A probability model does not claim that language is random in a
philosophical sense. Rather, it gives a mathematical way to represent
uncertainty from the viewpoint of an observer who has partial
information.

\section{Probability distributions on a finite
vocabulary}\label{probability-distributions-on-a-finite-vocabulary}

For a fixed vocabulary \(V=\{v_1,\ldots,v_d\}\), a probability
distribution over tokens is a vector

\[
p=(p_1,p_2,\ldots,p_d) \in R^d
\]

such that

\[
p_j \geq 0, \qquad \sum_{j=1}^d p_j = 1.
\]

The set of all such vectors is the probability simplex

\[
\Delta^{d-1}
=\left\{p\in R^d: p_j\geq 0,\ \sum_{j=1}^d p_j=1\right\}.
\]

If \(p_j=P(W=v_j)\), then \(p\) is the distribution of the random token
\(W\).

\subsection{Example}\label{example-1}

Let

\[
V=\{\text{data},\text{language},\text{model},\text{learns}\}.
\]

One possible token distribution is

\[
p=(0.35,0.25,0.30,0.10).
\]

This means

\[
P(W=\text{data})=0.35, \qquad P(W=\text{learns})=0.10.
\]

A model's prediction at a particular position in a sentence is also a
probability vector in a simplex. If a neural language model reads the
context

\begin{Shaded}
\begin{Highlighting}[]
\NormalTok{language models}
\end{Highlighting}
\end{Shaded}

then it may output a conditional distribution

\[
p_\theta(\cdot \mid \text{language models}) \in \Delta^{d-1}.
\]

The subscript \(\theta\) denotes trainable model parameters.

\section{Empirical distributions}\label{empirical-distributions}

In practice we do not know the true population distribution of language.
We observe a corpus and estimate probabilities from counts.

Let \(c_j\) be the number of times token \(v_j\) appears in a corpus and
let

\[
N=\sum_{j=1}^d c_j
\]

be the total number of vocabulary tokens. The empirical token
distribution is

\[
\widehat p_j=\frac{c_j}{N}, \qquad j=1,\ldots,d.
\]

This is a probability vector because \(\widehat p_j\geq 0\) and

\[
\sum_{j=1}^d \widehat p_j
=
\sum_{j=1}^d \frac{c_j}{N}
=
\frac{1}{N}\sum_{j=1}^d c_j
=1.
\]

\textbf{Empirical distribution.} A corpus turns counts into
probabilities by normalization:

\[
\widehat p_j=\frac{\text{count of }v_j}{\text{total token count}}.
\]

This is the simplest statistical estimator in language modeling.

\subsection{Python example: empirical token
probabilities}\label{python-example-empirical-token-probabilities}

\begin{Shaded}
\begin{Highlighting}[]
\ImportTok{from}\NormalTok{ collections }\ImportTok{import}\NormalTok{ Counter}
\ImportTok{import}\NormalTok{ math}

\NormalTok{corpus }\OperatorTok{=}\NormalTok{ [}
    \StringTok{"language models learn from data"}\NormalTok{,}
    \StringTok{"data models predict tokens"}\NormalTok{,}
    \StringTok{"language data has structure"}\NormalTok{,}
    \StringTok{"models learn probability"}
\NormalTok{]}

\KeywordTok{def}\NormalTok{ tokenize(text):}
    \ControlFlowTok{return}\NormalTok{ text.lower().split()}

\NormalTok{tokens }\OperatorTok{=}\NormalTok{ []}
\ControlFlowTok{for}\NormalTok{ doc }\KeywordTok{in}\NormalTok{ corpus:}
\NormalTok{    tokens.extend(tokenize(doc))}

\NormalTok{counts }\OperatorTok{=}\NormalTok{ Counter(tokens)}
\NormalTok{N }\OperatorTok{=} \BuiltInTok{sum}\NormalTok{(counts.values())}
\NormalTok{empirical\_p }\OperatorTok{=}\NormalTok{ \{token: count }\OperatorTok{/}\NormalTok{ N }\ControlFlowTok{for}\NormalTok{ token, count }\KeywordTok{in}\NormalTok{ counts.items()\}}

\BuiltInTok{print}\NormalTok{(}\StringTok{"Token counts:"}\NormalTok{)}
\BuiltInTok{print}\NormalTok{(counts)}
\BuiltInTok{print}\NormalTok{(}\StringTok{"}\CharTok{\textbackslash{}n}\StringTok{Empirical probabilities:"}\NormalTok{)}
\ControlFlowTok{for}\NormalTok{ token, prob }\KeywordTok{in} \BuiltInTok{sorted}\NormalTok{(empirical\_p.items()):}
    \BuiltInTok{print}\NormalTok{(}\SpecialStringTok{f"}\SpecialCharTok{\{}\NormalTok{token}\SpecialCharTok{:12s\}}\SpecialStringTok{ }\SpecialCharTok{\{}\NormalTok{prob}\SpecialCharTok{:.3f\}}\SpecialStringTok{"}\NormalTok{)}
\end{Highlighting}
\end{Shaded}

\begin{verbatim}
Token counts:
Counter({'models': 3, 'data': 3, 'language': 2, 'learn': 2, 'from': 1, 'predict': 1, 'tokens': 1, 'has': 1, 'structure': 1, 'probability': 1})

Empirical probabilities:
data         0.188
from         0.062
has          0.062
language     0.125
learn        0.125
models       0.188
predict      0.062
probability  0.062
structure    0.062
tokens       0.062
\end{verbatim}

The output is an empirical distribution on the vocabulary built from the
corpus. This distribution is not a deep model, but it is the
mathematical beginning of language modeling.

\section{Population distribution, data distribution, and model
distribution}\label{population-distribution-data-distribution-and-model-distribution}

It is useful to distinguish three distributions.

{\def\LTcaptype{none} % do not increment counter
\begin{longtable}[]{@{}
  >{\raggedright\arraybackslash}p{(\linewidth - 4\tabcolsep) * \real{0.3333}}
  >{\raggedright\arraybackslash}p{(\linewidth - 4\tabcolsep) * \real{0.3333}}
  >{\raggedright\arraybackslash}p{(\linewidth - 4\tabcolsep) * \real{0.3333}}@{}}
\toprule\noalign{}
\begin{minipage}[b]{\linewidth}\raggedright
Distribution
\end{minipage} & \begin{minipage}[b]{\linewidth}\raggedright
Meaning
\end{minipage} & \begin{minipage}[b]{\linewidth}\raggedright
Usually known?
\end{minipage} \\
\midrule\noalign{}
\endhead
\bottomrule\noalign{}
\endlastfoot
\(p^*\) & true data-generating distribution of language in the target
domain & no \\
\(\widehat p\) & empirical distribution from the observed corpus &
yes \\
\(p_\theta\) & model distribution with parameters \(\theta\) & yes after
training \\
\end{longtable}
}

The true distribution \(p^*\) is an idealized object. It represents the
probability law of text in the population or domain we care about. We
never observe \(p^*\) directly. Instead, we observe samples and build an
empirical approximation \(\widehat p\). Training tries to choose
parameters \(\theta\) so that \(p_\theta\) is close to \(p^*\), using
the corpus as evidence.

A clean statistical formulation is:

\[
\text{Choose }\theta\text{ so that }p_\theta\text{ assigns high probability to observed text.}
\]

This principle leads to maximum likelihood estimation.

\section{Maximum likelihood for a unigram
model}\label{maximum-likelihood-for-a-unigram-model}

The simplest language model assumes that tokens are independent and
identically distributed from a vocabulary distribution \(p\).

Let the corpus tokens be

\[
w_1,w_2,\ldots,w_N.
\]

The unigram model says

\[
P(W_1=w_1,\ldots,W_N=w_N)
=
\prod_{i=1}^N p(w_i).
\]

If token \(v_j\) appears \(c_j\) times, the likelihood is

\[
L(p)
=
\prod_{j=1}^d p_j^{c_j}.
\]

The log-likelihood is

\[
\ell(p)
=
\log L(p)
=
\sum_{j=1}^d c_j\log p_j.
\]

We maximize \(\ell(p)\) subject to

\[
p_j\geq 0,\qquad \sum_{j=1}^d p_j=1.
\]

Ignoring zero-count boundary details, use a Lagrange multiplier:

\[
\mathcal L(p,\lambda)
=
\sum_{j=1}^d c_j\log p_j
+
\lambda\left(1-\sum_{j=1}^d p_j\right).
\]

For \(c_j>0\),

\[
\frac{\partial \mathcal L}{\partial p_j}
=
\frac{c_j}{p_j}-\lambda=0.
\]

Therefore

\[
p_j=\frac{c_j}{\lambda}.
\]

Using \(\sum_j p_j=1\) gives

\[
1=\sum_{j=1}^d \frac{c_j}{\lambda}
=\frac{N}{\lambda},
\]

so \(\lambda=N\). Hence the maximum likelihood estimator is

\[
\widehat p_j=\frac{c_j}{N}.
\]

\textbf{MLE for the unigram model.} Under the independent-token model,
the maximum likelihood estimate of the token probability is the
empirical frequency:

\[
\widehat p_j=\frac{c_j}{N}.
\]

This derivation is important because it shows how a simple counting rule
is actually a likelihood-based statistical estimator.

\section{Sentences and the chain
rule}\label{sentences-and-the-chain-rule}

A sentence is not just a bag of independent tokens. Word order and
context matter. The full probability of a sequence is

\[
P(W_1=w_1,\ldots,W_T=w_T).
\]

The chain rule of probability gives the exact factorization

\[
P(w_1,\ldots,w_T)
=
P(w_1)P(w_2\mid w_1)P(w_3\mid w_1,w_2)\cdots P(w_T\mid w_1,
\ldots,w_{T-1}).
\]

Equivalently,

\[
P(w_{1:T})
=
\prod_{t=1}^T P(w_t\mid w_{1:t-1}),
\]

where \(w_{1:t-1}=(w_1,\ldots,w_{t-1})\).

\textbf{Language modeling problem.} Estimate the conditional
distribution

\[
P(W_t=\cdot \mid W_1=w_1,\ldots,W_{t-1}=w_{t-1}).
\]

A language model is a machine for producing this distribution.

This formula is exact. The difficulty is statistical and computational:
there are many possible histories \(w_{1:t-1}\). N-gram models solve
this by truncating the history. Transformers solve it by learning a
flexible representation of the history through attention.

\section{Conditional probability and
prediction}\label{conditional-probability-and-prediction}

The basic object in next-token prediction is a conditional probability.
Given a context \(C\), a language model outputs

\[
p_\theta(v_j\mid C),\qquad j=1,\ldots,d.
\]

This is a probability vector satisfying

\[
p_\theta(v_j\mid C)\geq 0,
\qquad
\sum_{j=1}^d p_\theta(v_j\mid C)=1.
\]

The most likely next token is

\[
\widehat w
=
\arg\max_{v_j\in V} p_\theta(v_j\mid C).
\]

But the full probability vector is more informative than only the
maximum. It tells us how uncertain the model is. For example, after the
context

\begin{Shaded}
\begin{Highlighting}[]
\NormalTok{machine learning models can}
\end{Highlighting}
\end{Shaded}

the model may assign high probability to several plausible next tokens
such as ``predict,'' ``learn,'' ``classify,'' or ``generate.'' A
deterministic answer hides that uncertainty.

\section{Negative log-likelihood}\label{negative-log-likelihood}

Suppose a model assigns probability \(p_\theta(w_i)\) to each observed
token \(w_i\). The likelihood of the observed corpus is

\[
L(\theta)=\prod_{i=1}^N p_\theta(w_i).
\]

The log-likelihood is

\[
\ell(\theta)=\sum_{i=1}^N \log p_\theta(w_i).
\]

Instead of maximizing log-likelihood, machine learning usually minimizes
negative log-likelihood:

\[
\operatorname{NLL}(\theta)
=
-\sum_{i=1}^N \log p_\theta(w_i).
\]

For a sequence model, this becomes

\[
\operatorname{NLL}(\theta)
=
-\sum_{t=1}^T \log p_\theta(w_t\mid w_{1:t-1}).
\]

The negative sign appears because optimization algorithms are often
written as minimization algorithms. Minimizing negative log-likelihood
is equivalent to maximizing likelihood.

\subsection{Python example: negative
log-likelihood}\label{python-example-negative-log-likelihood}

\begin{Shaded}
\begin{Highlighting}[]
\NormalTok{observed\_tokens }\OperatorTok{=}\NormalTok{ [}\StringTok{"data"}\NormalTok{, }\StringTok{"models"}\NormalTok{, }\StringTok{"learn"}\NormalTok{, }\StringTok{"data"}\NormalTok{]}
\NormalTok{model\_p }\OperatorTok{=}\NormalTok{ \{}
    \StringTok{"data"}\NormalTok{: }\FloatTok{0.40}\NormalTok{,}
    \StringTok{"models"}\NormalTok{: }\FloatTok{0.25}\NormalTok{,}
    \StringTok{"learn"}\NormalTok{: }\FloatTok{0.20}\NormalTok{,}
    \StringTok{"language"}\NormalTok{: }\FloatTok{0.15}
\NormalTok{\}}

\NormalTok{nll }\OperatorTok{=} \OperatorTok{{-}}\BuiltInTok{sum}\NormalTok{(math.log(model\_p[token]) }\ControlFlowTok{for}\NormalTok{ token }\KeywordTok{in}\NormalTok{ observed\_tokens)}
\NormalTok{avg\_nll }\OperatorTok{=}\NormalTok{ nll }\OperatorTok{/} \BuiltInTok{len}\NormalTok{(observed\_tokens)}

\BuiltInTok{print}\NormalTok{(}\StringTok{"Negative log{-}likelihood:"}\NormalTok{, }\BuiltInTok{round}\NormalTok{(nll, }\DecValTok{4}\NormalTok{))}
\BuiltInTok{print}\NormalTok{(}\StringTok{"Average negative log{-}likelihood:"}\NormalTok{, }\BuiltInTok{round}\NormalTok{(avg\_nll, }\DecValTok{4}\NormalTok{))}
\BuiltInTok{print}\NormalTok{(}\StringTok{"Perplexity:"}\NormalTok{, }\BuiltInTok{round}\NormalTok{(math.exp(avg\_nll), }\DecValTok{4}\NormalTok{))}
\end{Highlighting}
\end{Shaded}

\begin{verbatim}
Negative log-likelihood: 4.8283
Average negative log-likelihood: 1.2071
Perplexity: 3.3437
\end{verbatim}

The average negative log-likelihood is also the empirical cross-entropy
of the observed data relative to the model.

\section{Entropy}\label{entropy}

Entropy measures uncertainty in a probability distribution. For a
distribution \(p\) on \(V\), the entropy is

\[
H(p)=-\sum_{j=1}^d p_j\log p_j.
\]

We use the convention \(0\log 0=0\), interpreted as a limit.

If \(p\) is concentrated on one token, entropy is small. If \(p\) is
spread evenly across many tokens, entropy is large.

\subsection{Example}\label{example-2}

For a four-token vocabulary, compare

\[
p=(1,0,0,0)
\]

and

\[
q=\left(\frac14,\frac14,\frac14,\frac14\right).
\]

Then

\[
H(p)=0,
\]

while

\[
H(q)=-4\left(\frac14\log\frac14\right)=\log 4.
\]

The uniform distribution has the largest entropy among distributions on
four outcomes.

\subsection{Python example: entropy of token
distributions}\label{python-example-entropy-of-token-distributions}

\begin{Shaded}
\begin{Highlighting}[]
\KeywordTok{def}\NormalTok{ entropy(probs):}
    \ControlFlowTok{return} \OperatorTok{{-}}\BuiltInTok{sum}\NormalTok{(p }\OperatorTok{*}\NormalTok{ math.log(p) }\ControlFlowTok{for}\NormalTok{ p }\KeywordTok{in}\NormalTok{ probs }\ControlFlowTok{if}\NormalTok{ p }\OperatorTok{\textgreater{}} \DecValTok{0}\NormalTok{)}

\NormalTok{sharp }\OperatorTok{=}\NormalTok{ [}\FloatTok{0.90}\NormalTok{, }\FloatTok{0.05}\NormalTok{, }\FloatTok{0.03}\NormalTok{, }\FloatTok{0.02}\NormalTok{]}
\NormalTok{spread }\OperatorTok{=}\NormalTok{ [}\FloatTok{0.25}\NormalTok{, }\FloatTok{0.25}\NormalTok{, }\FloatTok{0.25}\NormalTok{, }\FloatTok{0.25}\NormalTok{]}

\BuiltInTok{print}\NormalTok{(}\StringTok{"Entropy of sharp distribution:"}\NormalTok{, }\BuiltInTok{round}\NormalTok{(entropy(sharp), }\DecValTok{4}\NormalTok{))}
\BuiltInTok{print}\NormalTok{(}\StringTok{"Entropy of spread distribution:"}\NormalTok{, }\BuiltInTok{round}\NormalTok{(entropy(spread), }\DecValTok{4}\NormalTok{))}
\end{Highlighting}
\end{Shaded}

\begin{verbatim}
Entropy of sharp distribution: 0.428
Entropy of spread distribution: 1.3863
\end{verbatim}

In NLP, entropy helps quantify uncertainty in token prediction. A model
with high next-token entropy is unsure among many possible
continuations. A model with low next-token entropy strongly prefers a
small number of continuations.

\section{Cross-entropy}\label{cross-entropy}

Suppose \(p\) is the true distribution and \(q\) is a model
distribution. The cross-entropy from \(p\) to \(q\) is

\[
H(p,q)=-\sum_{j=1}^d p_j\log q_j.
\]

This measures the expected negative log-probability assigned by model
\(q\) when the data are drawn from \(p\).

If \(p\) is replaced by an empirical distribution \(\widehat p\), then

\[
H(\widehat p,q)
=
-\sum_{j=1}^d \widehat p_j\log q_j.
\]

Using \(\widehat p_j=c_j/N\), we get

\[
H(\widehat p,q)
=
-\sum_{j=1}^d \frac{c_j}{N}\log q_j
=
-\frac1N \sum_{j=1}^d c_j\log q_j.
\]

This is average negative log-likelihood.

\textbf{Cross-entropy and training.} When we train a language model with
next-token prediction, the usual loss is empirical cross-entropy:

\[
-\frac1N \sum_{i=1}^N \log p_\theta(w_i\mid \text{context}_i).
\]

Thus cross-entropy is not just an evaluation metric. It is the main
training objective for many language models.

\section{KL divergence}\label{kl-divergence}

The Kullback-Leibler divergence from \(p\) to \(q\) is

\[
D_{\mathrm{KL}}(p\|q)
=
\sum_{j=1}^d p_j\log\frac{p_j}{q_j}.
\]

It satisfies

\[
D_{\mathrm{KL}}(p\|q)\geq 0,
\]

with equality if \(p=q\) on the support of \(p\).

The relationship between entropy, cross-entropy, and KL divergence is

\[
H(p,q)=H(p)+D_{\mathrm{KL}}(p\|q).
\]

Since \(H(p)\) does not depend on the model \(q\), minimizing
cross-entropy is equivalent to minimizing KL divergence from the data
distribution to the model distribution.

This identity gives a clean mathematical interpretation of
language-model training:

\[
\text{training reduces the discrepancy between the data distribution and the model distribution.}
\]

\section{Perplexity}\label{perplexity}

Perplexity is a common metric for language models. If the average
negative log-likelihood per token is

\[
\overline L
=
-\frac1N\sum_{i=1}^N \log p_\theta(w_i\mid \text{context}_i),
\]

then perplexity is

\[
\operatorname{PPL}=\exp(\overline L).
\]

If logarithms base \(2\) are used, then

\[
\operatorname{PPL}=2^{\overline L_2}.
\]

Perplexity can be interpreted as an effective number of equally likely
choices. For example, a perplexity of \(10\) roughly means the model is
as uncertain as if it were choosing among \(10\) equally likely tokens
at each step.

\textbf{Important caution.} Lower perplexity usually means better
next-token prediction on similar data, but it does not automatically
mean better reasoning, truthfulness, robustness, or usefulness.

\section{Interactive HTML: temperature, entropy, and next-token
uncertainty}\label{interactive-html-temperature-entropy-and-next-token-uncertainty}

The softmax function converts scores, or logits, into probabilities.
Given logits \(z_1,\ldots,z_d\) and temperature \(T>0\), define

\[
p_j(T)=\frac{\exp(z_j/T)}{\sum_{k=1}^d \exp(z_k/T)}.
\]

Small temperature makes the distribution sharper. Large temperature
makes it flatter. This affects entropy and perplexity.

The HTML version of this chapter contains an interactive Plotly.js
visualization. It lets the reader change the softmax temperature \(T\)
and observe how the next-token distribution, entropy, and effective
perplexity change.

\subsection{Mathematical lesson from the
interactive}\label{mathematical-lesson-from-the-interactive}

The temperature transformation

\[
p_j(T)=\frac{\exp(z_j/T)}{\sum_{k=1}^d \exp(z_k/T)}
\]

is not a new training procedure. It is a transformation of the output
distribution. When \(T<1\), the largest logits become more dominant.
When \(T>1\), the probabilities become more balanced. Thus temperature
controls randomness in generation by changing the entropy of the
next-token distribution.

\section{Conditional language models as statistical
estimators}\label{conditional-language-models-as-statistical-estimators}

Let \(C_t=(W_1,\ldots,W_{t-1})\) denote the context before position
\(t\). A conditional language model estimates

\[
p^*(w_t\mid C_t).
\]

The model is a parameterized function

\[
f_\theta:C_t \mapsto p_\theta(\cdot\mid C_t)\in \Delta^{d-1}.
\]

Training data consist of context-target pairs

\[
(C_i,y_i),\qquad i=1,\ldots,N,
\]

where \(y_i\) is the observed next token. The empirical cross-entropy
loss is

\[
\widehat R(\theta)
=
-\frac1N\sum_{i=1}^N \log p_\theta(y_i\mid C_i).
\]

This is the same mathematical structure whether \(f_\theta\) is an
n-gram table, a recurrent neural network, or a transformer.

\textbf{Unified view.} Classical and neural language models differ in
how they represent the context \(C\) and parameterize
\(p_\theta(\cdot\mid C)\), but they share the same probabilistic target:
estimate a conditional distribution over the next token.

\section{From counts to conditional
probabilities}\label{from-counts-to-conditional-probabilities}

A simple conditional model estimates probabilities from co-occurrence
counts. Let \(a\) be a previous token and \(b\) be a next token. Define

\[
c(a,b)=\#\{t: w_{t-1}=a,\ w_t=b\}
\]

and

\[
c(a)=\sum_{b\in V} c(a,b).
\]

The empirical conditional probability is

\[
\widehat P(W_t=b\mid W_{t-1}=a)
=
\frac{c(a,b)}{c(a)}.
\]

This is the bigram maximum likelihood estimator. Chapter 5 develops this
idea systematically.

\subsection{Python example: conditional next-token
probabilities}\label{python-example-conditional-next-token-probabilities}

\begin{Shaded}
\begin{Highlighting}[]
\ImportTok{from}\NormalTok{ collections }\ImportTok{import}\NormalTok{ defaultdict, Counter}

\NormalTok{sentences }\OperatorTok{=}\NormalTok{ [}
    \StringTok{"language models learn from data"}\NormalTok{,}
    \StringTok{"language models predict tokens"}\NormalTok{,}
    \StringTok{"models learn probability from data"}\NormalTok{,}
    \StringTok{"data models learn patterns"}
\NormalTok{]}

\NormalTok{bigram\_counts }\OperatorTok{=}\NormalTok{ defaultdict(Counter)}

\ControlFlowTok{for}\NormalTok{ sentence }\KeywordTok{in}\NormalTok{ sentences:}
\NormalTok{    toks }\OperatorTok{=}\NormalTok{ tokenize(sentence)}
    \ControlFlowTok{for}\NormalTok{ a, b }\KeywordTok{in} \BuiltInTok{zip}\NormalTok{(toks[:}\OperatorTok{{-}}\DecValTok{1}\NormalTok{], toks[}\DecValTok{1}\NormalTok{:]):}
\NormalTok{        bigram\_counts[a][b] }\OperatorTok{+=} \DecValTok{1}

\NormalTok{previous\_token }\OperatorTok{=} \StringTok{"models"}
\NormalTok{total\_after\_previous }\OperatorTok{=} \BuiltInTok{sum}\NormalTok{(bigram\_counts[previous\_token].values())}
\NormalTok{conditional\_probs }\OperatorTok{=}\NormalTok{ \{}
\NormalTok{    token: count }\OperatorTok{/}\NormalTok{ total\_after\_previous}
    \ControlFlowTok{for}\NormalTok{ token, count }\KeywordTok{in}\NormalTok{ bigram\_counts[previous\_token].items()}
\NormalTok{\}}

\BuiltInTok{print}\NormalTok{(}\SpecialStringTok{f"Next{-}token distribution after \textquotesingle{}}\SpecialCharTok{\{}\NormalTok{previous\_token}\SpecialCharTok{\}}\SpecialStringTok{\textquotesingle{}:"}\NormalTok{)}
\ControlFlowTok{for}\NormalTok{ token, prob }\KeywordTok{in} \BuiltInTok{sorted}\NormalTok{(conditional\_probs.items()):}
    \BuiltInTok{print}\NormalTok{(}\SpecialStringTok{f"P(}\SpecialCharTok{\{}\NormalTok{token}\SpecialCharTok{\}}\SpecialStringTok{ | }\SpecialCharTok{\{}\NormalTok{previous\_token}\SpecialCharTok{\}}\SpecialStringTok{) = }\SpecialCharTok{\{}\NormalTok{prob}\SpecialCharTok{:.3f\}}\SpecialStringTok{"}\NormalTok{)}
\end{Highlighting}
\end{Shaded}

\begin{verbatim}
Next-token distribution after 'models':
P(learn | models) = 0.750
P(predict | models) = 0.250
\end{verbatim}

This model uses only the previous token as context. It is simple enough
to inspect by hand, but it already illustrates the central
conditional-probability structure of language modeling.

\section{Zero probabilities and
smoothing}\label{zero-probabilities-and-smoothing}

A serious problem occurs when a token sequence has not appeared in the
training corpus. If \(c(a,b)=0\), then the maximum likelihood bigram
estimate gives

\[
\widehat P(b\mid a)=0.
\]

This means the model says the event is impossible. But in language,
unseen does not mean impossible. A small training corpus cannot contain
every valid phrase.

A common remedy is smoothing. For example, add-alpha smoothing estimates

\[
\widehat P_\alpha(b\mid a)
=
\frac{c(a,b)+\alpha}{c(a)+\alpha d},
\]

where \(\alpha>0\) and \(d=|V|\).

When \(\alpha=1\), this is called add-one smoothing or Laplace
smoothing.

\textbf{Smoothing principle.} Smoothing moves probability mass away from
observed events and assigns some mass to unseen events. It is a
statistical correction for finite data.

Chapter 5 studies smoothing in more detail for n-gram language models.

\section{Generative and discriminative probabilistic
models}\label{generative-and-discriminative-probabilistic-models}

Many NLP tasks involve a label \(Y\) and an input text \(X\). For
example:

\begin{itemize}
\tightlist
\item
  \(X\) is a movie review and \(Y\) is its sentiment label;
\item
  \(X\) is an email and \(Y\) indicates spam or not spam;
\item
  \(X\) is a student question and \(Y\) is the correct resource
  category.
\end{itemize}

There are two common probabilistic viewpoints.

\subsection{Generative modeling}\label{generative-modeling}

A generative model describes the joint distribution

\[
P(X,Y).
\]

Using the product rule,

\[
P(X,Y)=P(Y)P(X\mid Y).
\]

For classification, Bayes' rule gives

\[
P(Y\mid X)=\frac{P(X\mid Y)P(Y)}{P(X)}.
\]

Naive Bayes, developed in Chapter 6, is a generative text classifier. It
models how each class generates words.

\subsection{Discriminative modeling}\label{discriminative-modeling}

A discriminative model directly estimates

\[
P(Y\mid X).
\]

Logistic regression, neural classifiers, and transformer classifiers are
usually discriminative models. They do not need to model the probability
of the text itself; they only model the label distribution given the
text.

{\def\LTcaptype{none} % do not increment counter
\begin{longtable}[]{@{}
  >{\raggedright\arraybackslash}p{(\linewidth - 4\tabcolsep) * \real{0.3333}}
  >{\raggedright\arraybackslash}p{(\linewidth - 4\tabcolsep) * \real{0.3333}}
  >{\raggedright\arraybackslash}p{(\linewidth - 4\tabcolsep) * \real{0.3333}}@{}}
\toprule\noalign{}
\begin{minipage}[b]{\linewidth}\raggedright
Type
\end{minipage} & \begin{minipage}[b]{\linewidth}\raggedright
Models
\end{minipage} & \begin{minipage}[b]{\linewidth}\raggedright
Main object
\end{minipage} \\
\midrule\noalign{}
\endhead
\bottomrule\noalign{}
\endlastfoot
Generative & Naive Bayes, language models & \(P(X,Y)\) or \(P(X)\) \\
Discriminative & Logistic regression, neural classifiers &
\(P(Y\mid X)\) \\
\end{longtable}
}

Both viewpoints are useful. Generative models are often interpretable
and useful with limited data. Discriminative models are often more
accurate when sufficient labeled data are available.

\section{Bayesian updating for text
classification}\label{bayesian-updating-for-text-classification}

Bayes' rule is central to probabilistic NLP classification. Suppose
\(Y\in\{0,1\}\) is a class label and \(X\) is a text. Then

\[
P(Y=1\mid X)
=
\frac{P(X\mid Y=1)P(Y=1)}{P(X)}.
\]

Since

\[
P(X)=P(X\mid Y=0)P(Y=0)+P(X\mid Y=1)P(Y=1),
\]

we can compute the posterior probability of the class once we have class
priors and class-conditional text probabilities.

This equation is the mathematical foundation of Naive Bayes classifiers.
Its practical success depends on how we model \(P(X\mid Y)\).

\section{Worked example: a tiny probabilistic
classifier}\label{worked-example-a-tiny-probabilistic-classifier}

Suppose we classify short texts into two labels:

\[
Y\in\{\text{positive},\text{negative}\}.
\]

Assume equal class priors:

\[
P(Y=\text{positive})=P(Y=\text{negative})=\frac12.
\]

Let the text be represented by tokens. A simple generative model
estimates class-conditional token probabilities. Suppose

\[
P(\text{good}\mid \text{positive})=0.30,
\qquad
P(\text{good}\mid \text{negative})=0.05,
\]

and

\[
P(\text{clear}\mid \text{positive})=0.20,
\qquad
P(\text{clear}\mid \text{negative})=0.10.
\]

For the two-token document

\begin{Shaded}
\begin{Highlighting}[]
\NormalTok{good clear}
\end{Highlighting}
\end{Shaded}

and under the naive conditional independence approximation,

\[
P(X\mid Y=\text{positive})
=0.30\cdot 0.20=0.060,
\]

while

\[
P(X\mid Y=\text{negative})
=0.05\cdot 0.10=0.005.
\]

Multiplying by equal priors gives scores

\[
0.5\cdot 0.060=0.030,
\qquad
0.5\cdot 0.005=0.0025.
\]

The posterior probability of the positive class is

\[
\frac{0.030}{0.030+0.0025}\approx 0.923.
\]

Therefore the model predicts positive sentiment.

\subsection{Python example: tiny Bayesian
classifier}\label{python-example-tiny-bayesian-classifier}

\begin{Shaded}
\begin{Highlighting}[]
\NormalTok{priors }\OperatorTok{=}\NormalTok{ \{}\StringTok{"positive"}\NormalTok{: }\FloatTok{0.5}\NormalTok{, }\StringTok{"negative"}\NormalTok{: }\FloatTok{0.5}\NormalTok{\}}
\NormalTok{class\_token\_probs }\OperatorTok{=}\NormalTok{ \{}
    \StringTok{"positive"}\NormalTok{: \{}\StringTok{"good"}\NormalTok{: }\FloatTok{0.30}\NormalTok{, }\StringTok{"clear"}\NormalTok{: }\FloatTok{0.20}\NormalTok{\},}
    \StringTok{"negative"}\NormalTok{: \{}\StringTok{"good"}\NormalTok{: }\FloatTok{0.05}\NormalTok{, }\StringTok{"clear"}\NormalTok{: }\FloatTok{0.10}\NormalTok{\}}
\NormalTok{\}}

\NormalTok{document }\OperatorTok{=}\NormalTok{ [}\StringTok{"good"}\NormalTok{, }\StringTok{"clear"}\NormalTok{]}

\NormalTok{unnormalized }\OperatorTok{=}\NormalTok{ \{\}}
\ControlFlowTok{for}\NormalTok{ label }\KeywordTok{in}\NormalTok{ priors:}
\NormalTok{    score }\OperatorTok{=}\NormalTok{ priors[label]}
    \ControlFlowTok{for}\NormalTok{ token }\KeywordTok{in}\NormalTok{ document:}
\NormalTok{        score }\OperatorTok{*=}\NormalTok{ class\_token\_probs[label][token]}
\NormalTok{    unnormalized[label] }\OperatorTok{=}\NormalTok{ score}

\NormalTok{normalizing\_constant }\OperatorTok{=} \BuiltInTok{sum}\NormalTok{(unnormalized.values())}
\NormalTok{posterior }\OperatorTok{=}\NormalTok{ \{}
\NormalTok{    label: score }\OperatorTok{/}\NormalTok{ normalizing\_constant}
    \ControlFlowTok{for}\NormalTok{ label, score }\KeywordTok{in}\NormalTok{ unnormalized.items()}
\NormalTok{\}}

\BuiltInTok{print}\NormalTok{(}\StringTok{"Unnormalized scores:"}\NormalTok{, unnormalized)}
\BuiltInTok{print}\NormalTok{(}\StringTok{"Posterior probabilities:"}\NormalTok{, posterior)}
\BuiltInTok{print}\NormalTok{(}\StringTok{"Prediction:"}\NormalTok{, }\BuiltInTok{max}\NormalTok{(posterior, key}\OperatorTok{=}\NormalTok{posterior.get))}
\end{Highlighting}
\end{Shaded}

\begin{verbatim}
Unnormalized scores: {'positive': 0.03, 'negative': 0.0025000000000000005}
Posterior probabilities: {'positive': 0.923076923076923, 'negative': 0.07692307692307694}
Prediction: positive
\end{verbatim}

This small calculation previews Chapter 6. The important point here is
not the particular classifier, but the probabilistic structure: prior,
likelihood, posterior, and normalization.

\section{Log probabilities and numerical
stability}\label{log-probabilities-and-numerical-stability}

Products of many probabilities become extremely small. For a sentence of
length \(T\), the probability

\[
\prod_{t=1}^T p_\theta(w_t\mid w_{1:t-1})
\]

may underflow on a computer. Therefore NLP systems usually work with log
probabilities:

\[
\log \prod_{t=1}^T p_\theta(w_t\mid w_{1:t-1})
=
\sum_{t=1}^T \log p_\theta(w_t\mid w_{1:t-1}).
\]

Log probabilities turn products into sums, making computation more
stable and easier to optimize.

\subsection{Python example: product versus log
probability}\label{python-example-product-versus-log-probability}

\begin{Shaded}
\begin{Highlighting}[]
\NormalTok{probs }\OperatorTok{=}\NormalTok{ [}\FloatTok{0.02}\NormalTok{] }\OperatorTok{*} \DecValTok{100}

\NormalTok{product\_prob }\OperatorTok{=} \FloatTok{1.0}
\ControlFlowTok{for}\NormalTok{ p }\KeywordTok{in}\NormalTok{ probs:}
\NormalTok{    product\_prob }\OperatorTok{*=}\NormalTok{ p}

\NormalTok{log\_prob }\OperatorTok{=} \BuiltInTok{sum}\NormalTok{(math.log(p) }\ControlFlowTok{for}\NormalTok{ p }\KeywordTok{in}\NormalTok{ probs)}

\BuiltInTok{print}\NormalTok{(}\StringTok{"Product probability:"}\NormalTok{, product\_prob)}
\BuiltInTok{print}\NormalTok{(}\StringTok{"Log probability:"}\NormalTok{, log\_prob)}
\end{Highlighting}
\end{Shaded}

\begin{verbatim}
Product probability: 1.2676506002282316e-170
Log probability: -391.20230054281404
\end{verbatim}

For long sequences, the raw probability may be too small to interpret
directly. The log probability remains manageable.

\section{Statistical assumptions in probabilistic
NLP}\label{statistical-assumptions-in-probabilistic-nlp}

A mathematical NLP model is built from assumptions. Some assumptions are
explicit; others are hidden in the architecture or training data.

Important assumptions include:

\begin{enumerate}
\def\labelenumi{\arabic{enumi}.}
\tightlist
\item
  \textbf{Sampling assumption.} The training corpus is representative of
  the target distribution.
\item
  \textbf{Tokenization assumption.} The chosen tokens preserve the
  relevant information.
\item
  \textbf{Context assumption.} The model uses enough context to predict
  the target.
\item
  \textbf{Stationarity assumption.} The relationship between context and
  token does not change too much between training and deployment.
\item
  \textbf{Loss assumption.} Cross-entropy is an appropriate objective
  for the downstream goal.
\item
  \textbf{Evaluation assumption.} Test data represent the kinds of
  examples where the model will be used.
\end{enumerate}

These assumptions are not merely philosophical. When they fail, a model
may have low training loss but poor real-world behavior.

\section{AI component: using AI to audit probabilistic
assumptions}\label{ai-component-using-ai-to-audit-probabilistic-assumptions}

AI tools can help test understanding, generate examples, and critique
assumptions. However, AI explanations should be checked mathematically.

\subsection{Prompt 1: Explain the model
object}\label{prompt-1-explain-the-model-object}

Ask an AI assistant:

\begin{Shaded}
\begin{Highlighting}[]
\NormalTok{In mathematical terms, explain why a next{-}token language model estimates a conditional probability distribution. Use the notation p\_theta(w\_t | w\_1, ..., w\_\{t{-}1\}). Then give one example where the model may be uncertain.}
\end{Highlighting}
\end{Shaded}

Check the answer for the following:

\begin{itemize}
\tightlist
\item
  Does it distinguish probability distribution from a single prediction?
\item
  Does it correctly state that probabilities over the vocabulary sum to
  one?
\item
  Does it confuse the model distribution with the true data
  distribution?
\end{itemize}

\subsection{Prompt 2: Derive the unigram
MLE}\label{prompt-2-derive-the-unigram-mle}

Ask:

\begin{Shaded}
\begin{Highlighting}[]
\NormalTok{Derive the maximum likelihood estimator for a unigram language model over a finite vocabulary. Use a Lagrange multiplier and show why the answer is c\_j / N.}
\end{Highlighting}
\end{Shaded}

Then verify each step:

\begin{itemize}
\tightlist
\item
  Is the likelihood written as \(\prod_j p_j^{c_j}\)?
\item
  Is the constraint \(\sum_j p_j=1\) included?
\item
  Is the solution normalized correctly?
\end{itemize}

\subsection{Prompt 3: Critique
perplexity}\label{prompt-3-critique-perplexity}

Ask:

\begin{Shaded}
\begin{Highlighting}[]
\NormalTok{Give three reasons why a lower perplexity language model may not always be better for a real NLP application.}
\end{Highlighting}
\end{Shaded}

A good answer should mention domain shift, factuality, alignment with
task goals, robustness, or evaluation mismatch.

\textbf{AI study rule.} Use AI to generate explanations and
counterexamples, but always reduce probabilistic claims to definitions:
sample space, random variables, distributions, conditioning, likelihood,
and loss.

\section{Common mistakes}\label{common-mistakes}

\subsection{Mistake 1: Treating the most likely token as the whole
model}\label{mistake-1-treating-the-most-likely-token-as-the-whole-model}

A language model returns a full probability distribution. The argmax
token is only one summary of that distribution.

\subsection{Mistake 2: Confusing frequency with
probability}\label{mistake-2-confusing-frequency-with-probability}

Empirical frequency estimates probability from data, but it is not the
same as the unknown population probability.

\subsection{Mistake 3: Ignoring the role of the
context}\label{mistake-3-ignoring-the-role-of-the-context}

The probability of a token depends on context. For example,
\(P(\text{bank})\) and \(P(\text{bank}\mid \text{river})\) represent
different quantities.

\subsection{Mistake 4: Comparing raw likelihoods across corpora of
different
lengths}\label{mistake-4-comparing-raw-likelihoods-across-corpora-of-different-lengths}

Longer corpora usually have smaller raw likelihoods. Use average
log-likelihood, cross-entropy, or perplexity to compare across datasets.

\subsection{Mistake 5: Assuming low perplexity implies high-quality
generation}\label{mistake-5-assuming-low-perplexity-implies-high-quality-generation}

Perplexity measures next-token prediction on a distribution. It does not
directly measure truthfulness, reasoning, safety, or usefulness.

\section{Chapter summary}\label{chapter-summary-2}

This chapter introduced the probabilistic foundation of NLP.

\begin{enumerate}
\def\labelenumi{\arabic{enumi}.}
\tightlist
\item
  A token is modeled as a random variable taking values in a finite
  vocabulary.
\item
  A sentence is a sequence of random variables.
\item
  A language model assigns probabilities to token sequences or
  conditional probabilities to next tokens.
\item
  Empirical frequencies estimate probability distributions from corpora.
\item
  The unigram maximum likelihood estimator is the empirical token
  distribution.
\item
  The chain rule factors a sentence probability into conditional
  next-token probabilities.
\item
  Cross-entropy and negative log-likelihood are the central training
  losses for language models.
\item
  Entropy measures uncertainty; perplexity exponentiates average
  negative log-likelihood.
\item
  Generative models estimate \(P(X,Y)\) or \(P(X)\); discriminative
  models estimate \(P(Y\mid X)\).
\item
  Probabilistic models are useful only when their assumptions match the
  intended data and task.
\end{enumerate}

\section{Exercises}\label{exercises-3}

\subsection{Conceptual exercises}\label{conceptual-exercises-2}

\begin{enumerate}
\def\labelenumi{\arabic{enumi}.}
\item
  Let \(V=\{a,b,c\}\) and \(p=(0.2,0.5,0.3)\). Verify that
  \(p\in \Delta^2\) and compute \(H(p)\).
\item
  Explain the difference between \(p^*\), \(\widehat p\), and
  \(p_\theta\) in your own words.
\item
  Give an example of a context \(C\) where the next-token distribution
  should have high entropy. Give another where it should have low
  entropy.
\item
  Why is the chain rule factorization of \(P(w_{1:T})\) exact, while an
  n-gram approximation is not exact?
\item
  Explain why raw likelihood is difficult to compare across test sets of
  different lengths.
\end{enumerate}

\subsection{Mathematical exercises}\label{mathematical-exercises-1}

\begin{enumerate}
\def\labelenumi{\arabic{enumi}.}
\setcounter{enumi}{5}
\item
  Let a corpus over \(V=\{v_1,v_2,v_3\}\) have counts \((10,5,5)\).
  Compute the unigram MLE.
\item
  For the distribution \(p=(1/2,1/4,1/4)\), compute entropy using
  natural logarithms.
\item
  Let \(p=(0.5,0.5)\) and \(q=(0.9,0.1)\). Compute \(H(p,q)\) and
  \(D_{\mathrm{KL}}(p\|q)\).
\item
  Suppose a model assigns probabilities \(0.5\), \(0.25\), and \(0.1\)
  to three observed next tokens. Compute the total negative
  log-likelihood, average negative log-likelihood, and perplexity.
\item
  Derive the MLE for a Bernoulli token model with vocabulary
  \(V=\{0,1\}\) and show that it equals the sample proportion.
\end{enumerate}

\subsection{Python exercises}\label{python-exercises}

\begin{enumerate}
\def\labelenumi{\arabic{enumi}.}
\setcounter{enumi}{10}
\item
  Write a Python function that takes a list of tokens and returns the
  empirical probability distribution.
\item
  Write a Python function that computes entropy for any probability
  vector.
\item
  Write a Python function that computes cross-entropy between an
  empirical distribution and a model distribution.
\item
  Build a bigram count table from a small corpus and compute
  \(\widehat P(b\mid a)\) for several choices of \(a\).
\item
  Modify the bigram probability code to use add-alpha smoothing. Compare
  \(\alpha=0\), \(\alpha=0.1\), and \(\alpha=1\).
\end{enumerate}

\section{Mini-project: probabilistic language
audit}\label{mini-project-probabilistic-language-audit}

Choose a small corpus of 20 to 50 short documents, such as course
descriptions, news headlines, or abstract titles. Build a probability
audit with the following parts:

\begin{enumerate}
\def\labelenumi{\arabic{enumi}.}
\tightlist
\item
  tokenize the corpus and construct the vocabulary;
\item
  compute empirical unigram probabilities;
\item
  compute entropy of the unigram distribution;
\item
  build a bigram table and inspect next-token distributions after five
  selected tokens;
\item
  compute smoothed and unsmoothed probabilities for unseen bigrams;
\item
  write a short interpretation of what the probability model captures
  and what it misses.
\end{enumerate}

The final report should include at least one table of probabilities, one
plot, and one paragraph explaining the assumptions behind the model.

\section{Lab connection}\label{lab-connection-1}

The corresponding Python lab should be designed for independent study
and should include:

\begin{itemize}
\tightlist
\item
  review of finite probability spaces;
\item
  empirical token probability estimation;
\item
  likelihood and log-likelihood calculations;
\item
  entropy, cross-entropy, KL divergence, and perplexity;
\item
  bigram conditional probability estimation;
\item
  smoothing experiments;
\item
  a short AI-assisted model-audit activity.
\end{itemize}

\chapter{N-Gram Language Models}\label{n-gram-language-models}

Markov assumptions, conditional probability tables, smoothing,
generation, and perplexity

\hfill\break

An \(n\)-gram language model is one of the simplest mathematical models
of language. It predicts the next token using only the previous \(n-1\)
tokens. Although modern transformer language models are much more
powerful, \(n\)-gram models remain an essential starting point because
their assumptions, estimators, failures, and evaluation metrics can be
written down exactly. In this chapter we study \(n\)-gram models as
finite-state conditional probability models estimated from counts.

\section{Learning goals}\label{learning-goals-4}

After studying this chapter, you should be able to:

\begin{enumerate}
\def\labelenumi{\arabic{enumi}.}
\tightlist
\item
  define unigrams, bigrams, trigrams, and general \(n\)-grams;
\item
  explain the Markov approximation behind an \(n\)-gram language model;
\item
  estimate \(n\)-gram probabilities by maximum likelihood from corpus
  counts;
\item
  represent an \(n\)-gram model as a conditional probability table over
  contexts;
\item
  explain why zero counts create zero probabilities for unseen
  sentences;
\item
  derive and implement add-\(\alpha\) smoothing;
\item
  compute sentence log-probability, average negative log-likelihood,
  cross-entropy, and perplexity;
\item
  sample text from a fitted \(n\)-gram model;
\item
  compare unigram, bigram, and trigram models mathematically and
  computationally;
\item
  use AI tools to critique a language model's assumptions without
  outsourcing the reasoning.
\end{enumerate}

\section{\texorpdfstring{Why study \(n\)-gram
models?}{Why study n-gram models?}}\label{why-study-n-gram-models}

An \(n\)-gram model is not a modern state-of-the-art NLP model. However,
it is mathematically transparent. Nearly every important idea in
language modeling appears in a simple form:

\begin{itemize}
\tightlist
\item
  language as a probability distribution on token sequences;
\item
  conditioning on previous context;
\item
  maximum likelihood estimation;
\item
  sparse data and overfitting;
\item
  smoothing and regularization;
\item
  evaluation by cross-entropy and perplexity;
\item
  generation by sequential sampling.
\end{itemize}

Modern neural language models replace explicit count tables with
high-dimensional nonlinear functions, but they still model conditional
probabilities of the form

\[
p_\theta(w_t \mid w_1,\ldots,w_{t-1}).
\]

The \(n\)-gram model is the finite-table version of the same idea.

\textbf{Central idea.} An \(n\)-gram model approximates

\[
P(W_t=w_t \mid W_1=w_1,\ldots,W_{t-1}=w_{t-1})
\]

by conditioning only on the most recent \(n-1\) tokens:

\[
P(W_t=w_t \mid W_{t-n+1}=w_{t-n+1},\ldots,W_{t-1}=w_{t-1}).
\]

This is a mathematical simplification: the entire past is compressed
into a short context.

\section{\texorpdfstring{Tokens, contexts, and
\(n\)-grams}{Tokens, contexts, and n-grams}}\label{tokens-contexts-and-n-grams}

Let

\[
V=\{v_1,v_2,\ldots,v_d\}
\]

be a finite vocabulary. A token sequence is

\[
w_{1:T}=(w_1,w_2,\ldots,w_T), \qquad w_t\in V.
\]

An \(n\)-gram is a consecutive block of \(n\) tokens:

\[
(w_t,w_{t+1},\ldots,w_{t+n-1}).
\]

Common cases are:

{\def\LTcaptype{none} % do not increment counter
\begin{longtable}[]{@{}
  >{\raggedright\arraybackslash}p{(\linewidth - 4\tabcolsep) * \real{0.3333}}
  >{\raggedright\arraybackslash}p{(\linewidth - 4\tabcolsep) * \real{0.3333}}
  >{\raggedright\arraybackslash}p{(\linewidth - 4\tabcolsep) * \real{0.3333}}@{}}
\toprule\noalign{}
\begin{minipage}[b]{\linewidth}\raggedright
Model
\end{minipage} & \begin{minipage}[b]{\linewidth}\raggedright
Unit
\end{minipage} & \begin{minipage}[b]{\linewidth}\raggedright
Example from \texttt{language\ models\ learn\ probability}
\end{minipage} \\
\midrule\noalign{}
\endhead
\bottomrule\noalign{}
\endlastfoot
unigram & one token & \texttt{language}, \texttt{models},
\texttt{learn}, \texttt{probability} \\
bigram & two tokens & \texttt{language\ models}, \texttt{models\ learn},
\texttt{learn\ probability} \\
trigram & three tokens & \texttt{language\ models\ learn},
\texttt{models\ learn\ probability} \\
\end{longtable}
}

For prediction, it is useful to split an \(n\)-gram into a context and a
next token. For an \(n\)-gram model, the context length is \(n-1\):

\[
h_t=(w_{t-n+1},\ldots,w_{t-1}), \qquad \text{next token } w_t.
\]

The model estimates

\[
p(w_t\mid h_t).
\]

For a bigram model, the context is one token. For a trigram model, the
context is two tokens.

\section{The chain rule and the Markov
approximation}\label{the-chain-rule-and-the-markov-approximation}

From Chapter 4, the chain rule gives the exact factorization

\[
P(W_{1:T}=w_{1:T})
=
\prod_{t=1}^{T}P(W_t=w_t\mid W_1=w_1,\ldots,W_{t-1}=w_{t-1}).
\]

This formula is exact but not directly useful for estimation when \(T\)
is large. The number of possible histories grows explosively. The
\(n\)-gram approximation assumes that the conditional distribution
depends only on the most recent \(n-1\) tokens:

\[
P(W_t=w_t\mid W_1=w_1,\ldots,W_{t-1}=w_{t-1})
\approx
P(W_t=w_t\mid W_{t-n+1}=w_{t-n+1},\ldots,W_{t-1}=w_{t-1}).
\]

Therefore,

\[
P(w_{1:T})
\approx
\prod_{t=1}^{T}P(w_t\mid w_{t-n+1:t-1}).
\]

Here \(w_{a:b}\) denotes the subsequence \((w_a,w_{a+1},\ldots,w_b)\).

\textbf{Markov viewpoint.} A bigram model is a first-order Markov chain
on the vocabulary:

\[
P(W_t=w_t\mid W_1,\ldots,W_{t-1})
\approx
P(W_t=w_t\mid W_{t-1}).
\]

A trigram model is a second-order Markov model because it conditions on
the previous two tokens.

This does not mean language truly has such short memory. It means we
choose a tractable approximation.

\section{Boundary tokens}\label{boundary-tokens}

A practical language model needs to represent the beginning and end of a
sentence. We usually add special tokens such as

\[
\langle s\rangle, \qquad \langle/ s\rangle.
\]

For a bigram model, the sentence

\begin{Shaded}
\begin{Highlighting}[]
\NormalTok{language models learn}
\end{Highlighting}
\end{Shaded}

is transformed into

\begin{Shaded}
\begin{Highlighting}[]
\NormalTok{\textless{}s\textgreater{} language models learn \textless{}/s\textgreater{}}
\end{Highlighting}
\end{Shaded}

and contributes the bigrams

\begin{Shaded}
\begin{Highlighting}[]
\NormalTok{(\textless{}s\textgreater{}, language), (language, models), (models, learn), (learn, \textless{}/s\textgreater{})}
\end{Highlighting}
\end{Shaded}

For a trigram model, we often add two start tokens:

\begin{Shaded}
\begin{Highlighting}[]
\NormalTok{\textless{}s\textgreater{} \textless{}s\textgreater{} language models learn \textless{}/s\textgreater{}}
\end{Highlighting}
\end{Shaded}

so that the first real word can be predicted from a two-token context.

Mathematically, boundary tokens allow every position to have a context
of the required length.

\section{Maximum likelihood estimation for
bigrams}\label{maximum-likelihood-estimation-for-bigrams}

Let \(C(u,v)\) denote the number of times token \(u\) is followed by
token \(v\) in the training corpus. Let

\[
C(u)=\sum_{v\in V} C(u,v)
\]

be the number of times \(u\) appears as a context token. The maximum
likelihood estimator for a bigram model is

\[
\widehat p(v\mid u)=\frac{C(u,v)}{C(u)}.
\]

This is simply a row-normalized count table.

\subsection{Derivation}\label{derivation}

For a fixed context \(u\), suppose we observe counts

\[
C(u,v_1),C(u,v_2),\ldots,C(u,v_d).
\]

We want to estimate a conditional probability vector

\[
p(\cdot\mid u)=(p(v_1\mid u),\ldots,p(v_d\mid u))\in \Delta^{d-1}.
\]

The conditional likelihood for this row is

\[
L_u(p)=\prod_{j=1}^{d} p(v_j\mid u)^{C(u,v_j)}.
\]

The log-likelihood is

\[
\ell_u(p)=\sum_{j=1}^{d} C(u,v_j)\log p(v_j\mid u).
\]

We maximize \(\ell_u(p)\) subject to

\[
\sum_{j=1}^{d}p(v_j\mid u)=1.
\]

Using a Lagrange multiplier \(\lambda\),

\[
\mathcal{L}(p,\lambda)
=
\sum_{j=1}^{d}C(u,v_j)\log p(v_j\mid u)
+
\lambda\left(\sum_{j=1}^{d}p(v_j\mid u)-1\right).
\]

Differentiating with respect to \(p(v_j\mid u)\) gives

\[
\frac{C(u,v_j)}{p(v_j\mid u)}+\lambda=0.
\]

Thus

\[
p(v_j\mid u)=-\frac{C(u,v_j)}{\lambda}.
\]

Summing over \(j\) and using the constraint gives

\[
1=\sum_{j=1}^{d}p(v_j\mid u)
=-\frac{1}{\lambda}\sum_{j=1}^{d}C(u,v_j)
=-\frac{C(u)}{\lambda}.
\]

So \(\lambda=-C(u)\) and therefore

\[
\widehat p(v_j\mid u)=\frac{C(u,v_j)}{C(u)}.
\]

This derivation is important because it shows that the count ratio is
not merely intuitive. It is the maximum likelihood estimate under a
multinomial model.

\section{\texorpdfstring{General \(n\)-gram maximum likelihood
estimator}{General n-gram maximum likelihood estimator}}\label{general-n-gram-maximum-likelihood-estimator}

Let \(h\) be a context of length \(n-1\). For example, in a trigram
model,

\[
h=(u_1,u_2).
\]

Let \(C(h,v)\) be the number of times context \(h\) is followed by token
\(v\), and let

\[
C(h)=\sum_{v\in V}C(h,v).
\]

The maximum likelihood estimator is

\[
\widehat p(v\mid h)=\frac{C(h,v)}{C(h)}.
\]

The same Lagrange multiplier argument applies independently to each
context row.

\textbf{Conditional probability table.} An \(n\)-gram model is a
collection of probability vectors

\[
\{p(\cdot\mid h): h\in V^{n-1}\}.
\]

Each row belongs to the probability simplex \(\Delta^{d-1}\).

For bigrams, this table is a \(d\times d\) matrix. For trigrams, it is a
\(d^2\times d\) table. For general \(n\), the number of possible
contexts is \(d^{n-1}\), which grows quickly.

\section{\texorpdfstring{Python example: building an \(n\)-gram model
from
scratch}{Python example: building an n-gram model from scratch}}\label{python-example-building-an-n-gram-model-from-scratch}

The following code builds an \(n\)-gram count model, estimates
probabilities, and displays the next-token distribution for a chosen
context.

\begin{Shaded}
\begin{Highlighting}[]
\ImportTok{from}\NormalTok{ collections }\ImportTok{import}\NormalTok{ Counter, defaultdict}
\ImportTok{import}\NormalTok{ math}
\ImportTok{import}\NormalTok{ random}

\NormalTok{corpus }\OperatorTok{=}\NormalTok{ [}
    \StringTok{"language models learn from data"}\NormalTok{,}
    \StringTok{"language models predict tokens"}\NormalTok{,}
    \StringTok{"statistical models learn probability"}\NormalTok{,}
    \StringTok{"neural models learn representations"}\NormalTok{,}
    \StringTok{"data science uses probability models"}\NormalTok{,}
    \StringTok{"models predict language tokens"}\NormalTok{,}
\NormalTok{]}

\KeywordTok{def}\NormalTok{ tokenize(sentence):}
    \ControlFlowTok{return}\NormalTok{ sentence.lower().split()}

\KeywordTok{def}\NormalTok{ add\_boundaries(tokens, n):}
    \ControlFlowTok{return}\NormalTok{ [}\StringTok{"\textless{}s\textgreater{}"}\NormalTok{] }\OperatorTok{*}\NormalTok{ (n }\OperatorTok{{-}} \DecValTok{1}\NormalTok{) }\OperatorTok{+}\NormalTok{ tokens }\OperatorTok{+}\NormalTok{ [}\StringTok{"\textless{}/s\textgreater{}"}\NormalTok{]}

\KeywordTok{def}\NormalTok{ fit\_ngram(corpus, n}\OperatorTok{=}\DecValTok{2}\NormalTok{):}
\NormalTok{    context\_counts }\OperatorTok{=}\NormalTok{ Counter()}
\NormalTok{    next\_counts }\OperatorTok{=}\NormalTok{ defaultdict(Counter)}
\NormalTok{    vocab }\OperatorTok{=} \BuiltInTok{set}\NormalTok{([}\StringTok{"\textless{}s\textgreater{}"}\NormalTok{, }\StringTok{"\textless{}/s\textgreater{}"}\NormalTok{])}

    \ControlFlowTok{for}\NormalTok{ sentence }\KeywordTok{in}\NormalTok{ corpus:}
\NormalTok{        tokens }\OperatorTok{=}\NormalTok{ add\_boundaries(tokenize(sentence), n)}
\NormalTok{        vocab.update(tokens)}
        \ControlFlowTok{for}\NormalTok{ i }\KeywordTok{in} \BuiltInTok{range}\NormalTok{(n }\OperatorTok{{-}} \DecValTok{1}\NormalTok{, }\BuiltInTok{len}\NormalTok{(tokens)):}
\NormalTok{            context }\OperatorTok{=} \BuiltInTok{tuple}\NormalTok{(tokens[i }\OperatorTok{{-}}\NormalTok{ n }\OperatorTok{+} \DecValTok{1}\NormalTok{:i])}
\NormalTok{            word }\OperatorTok{=}\NormalTok{ tokens[i]}
\NormalTok{            context\_counts[context] }\OperatorTok{+=} \DecValTok{1}
\NormalTok{            next\_counts[context][word] }\OperatorTok{+=} \DecValTok{1}

\NormalTok{    vocab }\OperatorTok{=} \BuiltInTok{sorted}\NormalTok{(vocab)}
    \ControlFlowTok{return}\NormalTok{ vocab, context\_counts, next\_counts}

\NormalTok{vocab, context\_counts, next\_counts }\OperatorTok{=}\NormalTok{ fit\_ngram(corpus, n}\OperatorTok{=}\DecValTok{2}\NormalTok{)}
\NormalTok{context }\OperatorTok{=}\NormalTok{ (}\StringTok{"models"}\NormalTok{,)}

\BuiltInTok{print}\NormalTok{(}\StringTok{"Vocabulary:"}\NormalTok{, vocab)}
\BuiltInTok{print}\NormalTok{(}\StringTok{"Context count:"}\NormalTok{, context\_counts[context])}
\BuiltInTok{print}\NormalTok{(}\StringTok{"Next{-}token counts after \textquotesingle{}models\textquotesingle{}:"}\NormalTok{, next\_counts[context])}

\BuiltInTok{print}\NormalTok{(}\StringTok{"}\CharTok{\textbackslash{}n}\StringTok{MLE next{-}token probabilities:"}\NormalTok{)}
\ControlFlowTok{for}\NormalTok{ word, count }\KeywordTok{in}\NormalTok{ next\_counts[context].most\_common():}
\NormalTok{    prob }\OperatorTok{=}\NormalTok{ count }\OperatorTok{/}\NormalTok{ context\_counts[context]}
    \BuiltInTok{print}\NormalTok{(}\SpecialStringTok{f"P(}\SpecialCharTok{\{}\NormalTok{word}\SpecialCharTok{:16s\}}\SpecialStringTok{ | models) = }\SpecialCharTok{\{}\NormalTok{prob}\SpecialCharTok{:.3f\}}\SpecialStringTok{"}\NormalTok{)}
\end{Highlighting}
\end{Shaded}

\begin{verbatim}
Vocabulary: ['</s>', '<s>', 'data', 'from', 'language', 'learn', 'models', 'neural', 'predict', 'probability', 'representations', 'science', 'statistical', 'tokens', 'uses']
Context count: 6
Next-token counts after 'models': Counter({'learn': 3, 'predict': 2, '</s>': 1})

MLE next-token probabilities:
P(learn            | models) = 0.500
P(predict          | models) = 0.333
P(</s>             | models) = 0.167
\end{verbatim}

The model has learned, from this tiny corpus, which words tend to follow
\texttt{models}.

\section{The zero-probability
problem}\label{the-zero-probability-problem}

Maximum likelihood \(n\)-gram models assign probability zero to any
next-token event not observed in the training corpus:

\[
C(h,v)=0
\quad\Longrightarrow\quad
\widehat p(v\mid h)=0.
\]

This is a serious issue. If a single bigram in a test sentence has
probability zero, then the entire sentence probability becomes zero:

\[
\prod_{t=1}^{T}\widehat p(w_t\mid h_t)=0.
\]

The log-probability becomes undefined because \(\log 0=-\infty\).

\subsection{Example}\label{example-3}

Suppose the training corpus contains

\begin{Shaded}
\begin{Highlighting}[]
\NormalTok{language models learn}
\end{Highlighting}
\end{Shaded}

but never contains

\begin{Shaded}
\begin{Highlighting}[]
\NormalTok{language models explain}
\end{Highlighting}
\end{Shaded}

A bigram model may estimate

\[
\widehat p(\text{explain}\mid \text{models})=0,
\]

even though the phrase is perfectly reasonable. The zero probability
reflects limited data, not linguistic impossibility.

\section{\texorpdfstring{Add-\(\alpha\)
smoothing}{Add-\textbackslash alpha smoothing}}\label{add-alpha-smoothing}

A simple solution is add-\(\alpha\) smoothing. For \(\alpha>0\), define

\[
\widehat p_\alpha(v\mid h)
=
\frac{C(h,v)+\alpha}{C(h)+\alpha |V|}.
\]

The numerator adds \(\alpha\) pseudo-counts to every possible next
token. The denominator adds the total pseudo-count mass \(\alpha |V|\).

When \(\alpha=1\), this is called add-one or Laplace smoothing. Smaller
values such as \(\alpha=0.1\) often behave better in language modeling
because the vocabulary may be large.

\textbf{Smoothing as regularization.} Add-\(\alpha\) smoothing moves
each empirical row toward the uniform distribution:

\[
\widehat p_\alpha(\cdot\mid h)
=
\frac{C(h)}{C(h)+\alpha |V|}\widehat p_{\text{MLE}}(\cdot\mid h)
+
\frac{\alpha |V|}{C(h)+\alpha |V|}\,u,
\]

where \(u(v)=1/|V|\).

This formula shows that smoothing is a convex combination of the
empirical MLE distribution and the uniform distribution.

\subsection{\texorpdfstring{Python example: add-\(\alpha\)
probabilities}{Python example: add-\textbackslash alpha probabilities}}\label{python-example-add-alpha-probabilities}

\begin{Shaded}
\begin{Highlighting}[]
\KeywordTok{def}\NormalTok{ next\_token\_distribution(context, vocab, context\_counts, next\_counts, alpha}\OperatorTok{=}\FloatTok{0.5}\NormalTok{):}
\NormalTok{    d }\OperatorTok{=} \BuiltInTok{len}\NormalTok{(vocab)}
\NormalTok{    total }\OperatorTok{=}\NormalTok{ context\_counts[context]}
    \ControlFlowTok{return}\NormalTok{ \{}
\NormalTok{        word: (next\_counts[context][word] }\OperatorTok{+}\NormalTok{ alpha) }\OperatorTok{/}\NormalTok{ (total }\OperatorTok{+}\NormalTok{ alpha }\OperatorTok{*}\NormalTok{ d)}
        \ControlFlowTok{for}\NormalTok{ word }\KeywordTok{in}\NormalTok{ vocab}
        \ControlFlowTok{if}\NormalTok{ word }\OperatorTok{!=} \StringTok{"\textless{}s\textgreater{}"}
\NormalTok{    \}}

\NormalTok{context }\OperatorTok{=}\NormalTok{ (}\StringTok{"models"}\NormalTok{,)}
\ControlFlowTok{for}\NormalTok{ alpha }\KeywordTok{in}\NormalTok{ [}\FloatTok{0.0}\NormalTok{, }\FloatTok{0.1}\NormalTok{, }\FloatTok{1.0}\NormalTok{]:}
    \BuiltInTok{print}\NormalTok{(}\SpecialStringTok{f"}\CharTok{\textbackslash{}n}\SpecialStringTok{alpha = }\SpecialCharTok{\{}\NormalTok{alpha}\SpecialCharTok{\}}\SpecialStringTok{"}\NormalTok{)}
    \ControlFlowTok{if}\NormalTok{ alpha }\OperatorTok{==} \FloatTok{0.0}\NormalTok{:}
        \CommentTok{\# MLE: show only observed words to avoid many zeros}
\NormalTok{        dist }\OperatorTok{=}\NormalTok{ \{}
\NormalTok{            word: next\_counts[context][word] }\OperatorTok{/}\NormalTok{ context\_counts[context]}
            \ControlFlowTok{for}\NormalTok{ word }\KeywordTok{in}\NormalTok{ vocab}
            \ControlFlowTok{if}\NormalTok{ context\_counts[context] }\OperatorTok{\textgreater{}} \DecValTok{0} \KeywordTok{and}\NormalTok{ next\_counts[context][word] }\OperatorTok{\textgreater{}} \DecValTok{0}
\NormalTok{        \}}
    \ControlFlowTok{else}\NormalTok{:}
\NormalTok{        dist }\OperatorTok{=}\NormalTok{ next\_token\_distribution(context, vocab, context\_counts, next\_counts, alpha)}
    \ControlFlowTok{for}\NormalTok{ word, prob }\KeywordTok{in} \BuiltInTok{sorted}\NormalTok{(dist.items(), key}\OperatorTok{=}\KeywordTok{lambda}\NormalTok{ x: }\OperatorTok{{-}}\NormalTok{x[}\DecValTok{1}\NormalTok{])[:}\DecValTok{8}\NormalTok{]:}
        \BuiltInTok{print}\NormalTok{(}\SpecialStringTok{f"P(}\SpecialCharTok{\{}\NormalTok{word}\SpecialCharTok{:16s\}}\SpecialStringTok{ | models) = }\SpecialCharTok{\{}\NormalTok{prob}\SpecialCharTok{:.3f\}}\SpecialStringTok{"}\NormalTok{)}
\end{Highlighting}
\end{Shaded}

\begin{verbatim}

alpha = 0.0
P(learn            | models) = 0.500
P(predict          | models) = 0.333
P(</s>             | models) = 0.167

alpha = 0.1
P(learn            | models) = 0.413
P(predict          | models) = 0.280
P(</s>             | models) = 0.147
P(data             | models) = 0.013
P(from             | models) = 0.013
P(language         | models) = 0.013
P(models           | models) = 0.013
P(neural           | models) = 0.013

alpha = 1.0
P(learn            | models) = 0.190
P(predict          | models) = 0.143
P(</s>             | models) = 0.095
P(data             | models) = 0.048
P(from             | models) = 0.048
P(language         | models) = 0.048
P(models           | models) = 0.048
P(neural           | models) = 0.048
\end{verbatim}

As \(\alpha\) increases, probability mass spreads more evenly across the
vocabulary.

\section{Sentence probability and
log-probability}\label{sentence-probability-and-log-probability}

Let \(w_{1:T}\) be a sentence after adding boundary tokens. Under an
\(n\)-gram model,

\[
p(w_{1:T})
=
\prod_{t=n}^{T}p(w_t\mid w_{t-n+1:t-1}).
\]

In computation, we usually work with log-probabilities:

\[
\log p(w_{1:T})
=
\sum_{t=n}^{T}\log p(w_t\mid w_{t-n+1:t-1}).
\]

This is numerically more stable because products of many small
probabilities can underflow to zero.

\subsection{Python example: sentence
log-probability}\label{python-example-sentence-log-probability}

\begin{Shaded}
\begin{Highlighting}[]
\KeywordTok{def}\NormalTok{ sentence\_logprob(sentence, n, vocab, context\_counts, next\_counts, alpha}\OperatorTok{=}\FloatTok{0.5}\NormalTok{):}
\NormalTok{    tokens }\OperatorTok{=}\NormalTok{ add\_boundaries(tokenize(sentence), n)}
\NormalTok{    d }\OperatorTok{=} \BuiltInTok{len}\NormalTok{(vocab)}
\NormalTok{    logp }\OperatorTok{=} \FloatTok{0.0}
\NormalTok{    details }\OperatorTok{=}\NormalTok{ []}

    \ControlFlowTok{for}\NormalTok{ i }\KeywordTok{in} \BuiltInTok{range}\NormalTok{(n }\OperatorTok{{-}} \DecValTok{1}\NormalTok{, }\BuiltInTok{len}\NormalTok{(tokens)):}
\NormalTok{        context }\OperatorTok{=} \BuiltInTok{tuple}\NormalTok{(tokens[i }\OperatorTok{{-}}\NormalTok{ n }\OperatorTok{+} \DecValTok{1}\NormalTok{:i])}
\NormalTok{        word }\OperatorTok{=}\NormalTok{ tokens[i]}
\NormalTok{        count\_context }\OperatorTok{=}\NormalTok{ context\_counts[context]}
\NormalTok{        count\_pair }\OperatorTok{=}\NormalTok{ next\_counts[context][word]}
\NormalTok{        prob }\OperatorTok{=}\NormalTok{ (count\_pair }\OperatorTok{+}\NormalTok{ alpha) }\OperatorTok{/}\NormalTok{ (count\_context }\OperatorTok{+}\NormalTok{ alpha }\OperatorTok{*}\NormalTok{ d)}
\NormalTok{        logp }\OperatorTok{+=}\NormalTok{ math.log(prob)}
\NormalTok{        details.append((context, word, prob))}

    \ControlFlowTok{return}\NormalTok{ logp, details}

\NormalTok{sentence }\OperatorTok{=} \StringTok{"language models learn probability"}
\NormalTok{logp, details }\OperatorTok{=}\NormalTok{ sentence\_logprob(sentence, }\DecValTok{2}\NormalTok{, vocab, context\_counts, next\_counts, alpha}\OperatorTok{=}\FloatTok{0.5}\NormalTok{)}

\BuiltInTok{print}\NormalTok{(}\StringTok{"Sentence:"}\NormalTok{, sentence)}
\BuiltInTok{print}\NormalTok{(}\StringTok{"Log{-}probability:"}\NormalTok{, }\BuiltInTok{round}\NormalTok{(logp, }\DecValTok{4}\NormalTok{))}
\BuiltInTok{print}\NormalTok{(}\StringTok{"Probability:"}\NormalTok{, }\BuiltInTok{round}\NormalTok{(math.exp(logp), }\DecValTok{8}\NormalTok{))}
\BuiltInTok{print}\NormalTok{(}\StringTok{"}\CharTok{\textbackslash{}n}\StringTok{Factors:"}\NormalTok{)}
\ControlFlowTok{for}\NormalTok{ context, word, prob }\KeywordTok{in}\NormalTok{ details:}
    \BuiltInTok{print}\NormalTok{(}\SpecialStringTok{f"P(}\SpecialCharTok{\{}\NormalTok{word}\SpecialCharTok{:16s\}}\SpecialStringTok{ | }\SpecialCharTok{\{}\NormalTok{context}\SpecialCharTok{\}}\SpecialStringTok{) = }\SpecialCharTok{\{}\NormalTok{prob}\SpecialCharTok{:.4f\}}\SpecialStringTok{"}\NormalTok{)}
\end{Highlighting}
\end{Shaded}

\begin{verbatim}
Sentence: language models learn probability
Log-probability: -8.2631
Probability: 0.00025785

Factors:
P(language         | ('<s>',)) = 0.1852
P(models           | ('language',)) = 0.2381
P(learn            | ('models',)) = 0.2593
P(probability      | ('learn',)) = 0.1429
P(</s>             | ('probability',)) = 0.1579
\end{verbatim}

The sentence probability is a product of local next-token probabilities.

\section{Cross-entropy and
perplexity}\label{cross-entropy-and-perplexity}

Let a test corpus contain \(N\) predicted tokens. Suppose the model
assigns probability \(p_i\) to the correct token at position \(i\). The
average negative log-likelihood is

\[
\operatorname{NLL}_{\text{avg}}
=
-\frac{1}{N}\sum_{i=1}^{N}\log p_i.
\]

When natural logarithms are used, this quantity is measured in nats per
token. The perplexity is

\[
\operatorname{Perplexity}
=
\exp\left(\operatorname{NLL}_{\text{avg}}\right).
\]

Equivalently,

\[
\operatorname{Perplexity}
=
\left(\prod_{i=1}^{N}\frac{1}{p_i}\right)^{1/N}.
\]

So perplexity is the geometric mean of the inverse probabilities
assigned to the correct next tokens.

\textbf{Interpretation.} Lower perplexity is better. A perplexity of
\(10\) roughly means that, on average, the model is as uncertain as
choosing uniformly among \(10\) plausible next tokens.

This interpretation is approximate, but it is useful for intuition.

\subsection{Python example: evaluating
perplexity}\label{python-example-evaluating-perplexity}

\begin{Shaded}
\begin{Highlighting}[]
\KeywordTok{def}\NormalTok{ evaluate\_perplexity(test\_corpus, n, vocab, context\_counts, next\_counts, alpha}\OperatorTok{=}\FloatTok{0.5}\NormalTok{):}
\NormalTok{    total\_logp }\OperatorTok{=} \FloatTok{0.0}
\NormalTok{    total\_tokens }\OperatorTok{=} \DecValTok{0}

    \ControlFlowTok{for}\NormalTok{ sentence }\KeywordTok{in}\NormalTok{ test\_corpus:}
\NormalTok{        tokens }\OperatorTok{=}\NormalTok{ add\_boundaries(tokenize(sentence), n)}
\NormalTok{        d }\OperatorTok{=} \BuiltInTok{len}\NormalTok{(vocab)}
        \ControlFlowTok{for}\NormalTok{ i }\KeywordTok{in} \BuiltInTok{range}\NormalTok{(n }\OperatorTok{{-}} \DecValTok{1}\NormalTok{, }\BuiltInTok{len}\NormalTok{(tokens)):}
\NormalTok{            context }\OperatorTok{=} \BuiltInTok{tuple}\NormalTok{(tokens[i }\OperatorTok{{-}}\NormalTok{ n }\OperatorTok{+} \DecValTok{1}\NormalTok{:i])}
\NormalTok{            word }\OperatorTok{=}\NormalTok{ tokens[i]}
\NormalTok{            prob }\OperatorTok{=}\NormalTok{ (next\_counts[context][word] }\OperatorTok{+}\NormalTok{ alpha) }\OperatorTok{/}\NormalTok{ (context\_counts[context] }\OperatorTok{+}\NormalTok{ alpha }\OperatorTok{*}\NormalTok{ d)}
\NormalTok{            total\_logp }\OperatorTok{+=}\NormalTok{ math.log(prob)}
\NormalTok{            total\_tokens }\OperatorTok{+=} \DecValTok{1}

\NormalTok{    avg\_nll }\OperatorTok{=} \OperatorTok{{-}}\NormalTok{total\_logp }\OperatorTok{/}\NormalTok{ total\_tokens}
    \ControlFlowTok{return}\NormalTok{ math.exp(avg\_nll), avg\_nll}

\NormalTok{test\_corpus }\OperatorTok{=}\NormalTok{ [}
    \StringTok{"language models predict probability"}\NormalTok{,}
    \StringTok{"neural models learn language"}\NormalTok{,}
\NormalTok{]}

\ControlFlowTok{for}\NormalTok{ alpha }\KeywordTok{in}\NormalTok{ [}\FloatTok{0.1}\NormalTok{, }\FloatTok{0.5}\NormalTok{, }\FloatTok{1.0}\NormalTok{]:}
\NormalTok{    ppl, avg\_nll }\OperatorTok{=}\NormalTok{ evaluate\_perplexity(test\_corpus, }\DecValTok{2}\NormalTok{, vocab, context\_counts, next\_counts, alpha)}
    \BuiltInTok{print}\NormalTok{(}\SpecialStringTok{f"alpha=}\SpecialCharTok{\{}\NormalTok{alpha}\SpecialCharTok{:3.1f\}}\SpecialStringTok{  avg NLL=}\SpecialCharTok{\{}\NormalTok{avg\_nll}\SpecialCharTok{:.3f\}}\SpecialStringTok{  perplexity=}\SpecialCharTok{\{}\NormalTok{ppl}\SpecialCharTok{:.3f\}}\SpecialStringTok{"}\NormalTok{)}
\end{Highlighting}
\end{Shaded}

\begin{verbatim}
alpha=0.1  avg NLL=1.926  perplexity=6.861
alpha=0.5  avg NLL=2.097  perplexity=8.141
alpha=1.0  avg NLL=2.253  perplexity=9.513
\end{verbatim}

Perplexity depends on the test set, the vocabulary, the tokenization
scheme, the order \(n\), and the smoothing method. Comparing
perplexities is meaningful only when these conditions are comparable.

\section{Interactive HTML: smoothing and next-token
uncertainty}\label{interactive-html-smoothing-and-next-token-uncertainty}

The interactive below uses a small bigram corpus. Move the smoothing
slider to see how add-\(\alpha\) smoothing changes the conditional
distribution after the context word \texttt{models}. The bar chart shows
the next-token probabilities. The table also reports the entropy of the
distribution.

The key mathematical pattern is shrinkage. When \(\alpha\) is small, the
distribution is close to the empirical count distribution. When
\(\alpha\) is large, the distribution moves toward uniform.

\section{\texorpdfstring{Sampling text from an \(n\)-gram
model}{Sampling text from an n-gram model}}\label{sampling-text-from-an-n-gram-model}

A fitted \(n\)-gram model is generative. Starting with boundary tokens,
repeatedly sample the next token from

\[
p(\cdot\mid h_t).
\]

Then update the context and continue until the end token appears.

\subsection{Python example: bigram
generation}\label{python-example-bigram-generation}

\begin{Shaded}
\begin{Highlighting}[]
\KeywordTok{def}\NormalTok{ sample\_next(distribution):}
\NormalTok{    words }\OperatorTok{=} \BuiltInTok{list}\NormalTok{(distribution.keys())}
\NormalTok{    probs }\OperatorTok{=} \BuiltInTok{list}\NormalTok{(distribution.values())}
    \ControlFlowTok{return}\NormalTok{ random.choices(words, weights}\OperatorTok{=}\NormalTok{probs, k}\OperatorTok{=}\DecValTok{1}\NormalTok{)[}\DecValTok{0}\NormalTok{]}

\KeywordTok{def}\NormalTok{ generate\_ngram(n, vocab, context\_counts, next\_counts, alpha}\OperatorTok{=}\FloatTok{0.5}\NormalTok{, max\_steps}\OperatorTok{=}\DecValTok{20}\NormalTok{, seed}\OperatorTok{=}\DecValTok{7}\NormalTok{):}
\NormalTok{    random.seed(seed)}
\NormalTok{    context }\OperatorTok{=} \BuiltInTok{tuple}\NormalTok{([}\StringTok{"\textless{}s\textgreater{}"}\NormalTok{] }\OperatorTok{*}\NormalTok{ (n }\OperatorTok{{-}} \DecValTok{1}\NormalTok{))}
\NormalTok{    generated }\OperatorTok{=}\NormalTok{ []}

    \ControlFlowTok{for}\NormalTok{ \_ }\KeywordTok{in} \BuiltInTok{range}\NormalTok{(max\_steps):}
\NormalTok{        dist }\OperatorTok{=}\NormalTok{ next\_token\_distribution(context, vocab, context\_counts, next\_counts, alpha}\OperatorTok{=}\NormalTok{alpha)}
\NormalTok{        word }\OperatorTok{=}\NormalTok{ sample\_next(dist)}
        \ControlFlowTok{if}\NormalTok{ word }\OperatorTok{==} \StringTok{"\textless{}/s\textgreater{}"}\NormalTok{:}
            \ControlFlowTok{break}
\NormalTok{        generated.append(word)}
        \ControlFlowTok{if}\NormalTok{ n }\OperatorTok{\textgreater{}} \DecValTok{1}\NormalTok{:}
\NormalTok{            context }\OperatorTok{=} \BuiltInTok{tuple}\NormalTok{((}\BuiltInTok{list}\NormalTok{(context) }\OperatorTok{+}\NormalTok{ [word])[}\OperatorTok{{-}}\NormalTok{(n }\OperatorTok{{-}} \DecValTok{1}\NormalTok{):])}

    \ControlFlowTok{return} \StringTok{" "}\NormalTok{.join(generated)}

\ControlFlowTok{for}\NormalTok{ seed }\KeywordTok{in} \BuiltInTok{range}\NormalTok{(}\DecValTok{5}\NormalTok{):}
    \BuiltInTok{print}\NormalTok{(generate\_ngram(}\DecValTok{2}\NormalTok{, vocab, context\_counts, next\_counts, alpha}\OperatorTok{=}\FloatTok{0.3}\NormalTok{, seed}\OperatorTok{=}\NormalTok{seed))}
\end{Highlighting}
\end{Shaded}

\begin{verbatim}
statistical representations language models learn neural science models learn probability tokens language models predict representations data tokens uses science uses
data science uses learn probability models predict tokens
tokens tokens
language neural models predict representations
language from language language data learn tokens representations representations
\end{verbatim}

The generated text may be grammatical in small fragments, but it often
lacks long-range coherence. This is a mathematical consequence of the
short-context assumption.

\section{Comparing unigram, bigram, and trigram
models}\label{comparing-unigram-bigram-and-trigram-models}

A unigram model assumes tokens are independent:

\[
p(w_{1:T})\approx \prod_{t=1}^{T}p(w_t).
\]

A bigram model keeps one previous token:

\[
p(w_{1:T})\approx \prod_{t=1}^{T}p(w_t\mid w_{t-1}).
\]

A trigram model keeps two previous tokens:

\[
p(w_{1:T})\approx \prod_{t=1}^{T}p(w_t\mid w_{t-2},w_{t-1}).
\]

As \(n\) increases, the model can represent richer local structure. But
the number of possible contexts grows as

\[
|V|^{n-1}.
\]

This creates a bias-variance tradeoff.

{\def\LTcaptype{none} % do not increment counter
\begin{longtable}[]{@{}lrll@{}}
\toprule\noalign{}
Model & Context length & Strength & Weakness \\
\midrule\noalign{}
\endhead
\bottomrule\noalign{}
\endlastfoot
unigram & 0 & stable estimates & ignores order \\
bigram & 1 & captures local transitions & limited memory \\
trigram & 2 & captures short phrases & sparse counts \\
larger \(n\) & \(n-1\) & more context & severe sparsity \\
\end{longtable}
}

\textbf{Bias-variance viewpoint.} Small \(n\) gives high bias because
the model is too simple. Large \(n\) can give high variance because many
contexts have few or no observations. Smoothing partially controls
variance but cannot fully solve the curse of dimensionality.

\section{Backoff and interpolation}\label{backoff-and-interpolation}

Add-\(\alpha\) smoothing is simple but not always strong enough. Two
common ideas are backoff and interpolation.

\subsection{Backoff}\label{backoff}

If a trigram context is rare or unseen, use a lower-order model:

\[
p(w_t\mid w_{t-2},w_{t-1})
\quad \longrightarrow \quad
p(w_t\mid w_{t-1})
\quad \longrightarrow \quad
p(w_t).
\]

The model backs off from specific context to more general context.

\subsection{Interpolation}\label{interpolation}

Interpolation combines several orders:

\[
p(w_t\mid w_{t-2},w_{t-1})
=
\lambda_3 p_3(w_t\mid w_{t-2},w_{t-1})
+
\lambda_2 p_2(w_t\mid w_{t-1})
+
\lambda_1 p_1(w_t),
\]

where

\[
\lambda_1,\lambda_2,\lambda_3\geq 0,
\qquad
\lambda_1+\lambda_2+\lambda_3=1.
\]

This is a convex combination of probability distributions. It is
mathematically similar to ensemble methods and shrinkage estimators.

\section{Matrix viewpoint for bigram
models}\label{matrix-viewpoint-for-bigram-models}

For a vocabulary of size \(d\), a bigram model can be represented by a
transition matrix

\[
P\in R^{d\times d},
\]

where

\[
P_{ij}=P(W_t=v_j\mid W_{t-1}=v_i).
\]

Each row is a probability vector:

\[
P_{ij}\geq 0,
\qquad
\sum_{j=1}^{d}P_{ij}=1.
\]

Thus \(P\) is a row-stochastic matrix. A distribution \(\pi_t\) over the
current token evolves by

\[
\pi_{t+1}=\pi_t P.
\]

This connects bigram language modeling to finite Markov chains.

\subsection{Python example: transition
matrix}\label{python-example-transition-matrix}

\begin{Shaded}
\begin{Highlighting}[]
\CommentTok{\# Build a smoothed bigram transition matrix over the full vocabulary.}
\CommentTok{\# Using the full vocabulary makes each row a probability vector.}
\NormalTok{selected }\OperatorTok{=}\NormalTok{ vocab}
\NormalTok{index }\OperatorTok{=}\NormalTok{ \{word: i }\ControlFlowTok{for}\NormalTok{ i, word }\KeywordTok{in} \BuiltInTok{enumerate}\NormalTok{(selected)\}}
\NormalTok{alpha }\OperatorTok{=} \FloatTok{0.2}

\NormalTok{P }\OperatorTok{=}\NormalTok{ []}
\ControlFlowTok{for}\NormalTok{ u }\KeywordTok{in}\NormalTok{ selected:}
\NormalTok{    row }\OperatorTok{=}\NormalTok{ []}
\NormalTok{    context }\OperatorTok{=}\NormalTok{ (u,)}
\NormalTok{    denom }\OperatorTok{=}\NormalTok{ context\_counts[context] }\OperatorTok{+}\NormalTok{ alpha }\OperatorTok{*} \BuiltInTok{len}\NormalTok{(selected)}
    \ControlFlowTok{for}\NormalTok{ v }\KeywordTok{in}\NormalTok{ selected:}
\NormalTok{        row.append((next\_counts[context][v] }\OperatorTok{+}\NormalTok{ alpha) }\OperatorTok{/}\NormalTok{ denom)}
\NormalTok{    P.append(row)}

\BuiltInTok{print}\NormalTok{(}\StringTok{"Rows approximately sum to 1:"}\NormalTok{)}
\ControlFlowTok{for}\NormalTok{ word, row }\KeywordTok{in} \BuiltInTok{list}\NormalTok{(}\BuiltInTok{zip}\NormalTok{(selected, P))[:}\DecValTok{8}\NormalTok{]:}
    \BuiltInTok{print}\NormalTok{(}\SpecialStringTok{f"}\SpecialCharTok{\{}\NormalTok{word}\SpecialCharTok{:12s\}}\SpecialStringTok{ row sum = }\SpecialCharTok{\{}\BuiltInTok{sum}\NormalTok{(row)}\SpecialCharTok{:.6f\}}\SpecialStringTok{"}\NormalTok{)}

\BuiltInTok{print}\NormalTok{(}\StringTok{"}\CharTok{\textbackslash{}n}\StringTok{Transition probability P(models {-}\textgreater{} learn):"}\NormalTok{)}
\BuiltInTok{print}\NormalTok{(P[index[}\StringTok{"models"}\NormalTok{]][index[}\StringTok{"learn"}\NormalTok{]])}
\end{Highlighting}
\end{Shaded}

\begin{verbatim}
Rows approximately sum to 1:
</s>         row sum = 1.000000
<s>          row sum = 1.000000
data         row sum = 1.000000
from         row sum = 1.000000
language     row sum = 1.000000
learn        row sum = 1.000000
models       row sum = 1.000000
neural       row sum = 1.000000

Transition probability P(models -> learn):
0.35555555555555557
\end{verbatim}

The transition matrix view is useful because it exposes the connection
between NLP and stochastic processes.

\section{Worked hand calculation}\label{worked-hand-calculation}

Consider the corpus consisting of two sentences:

\begin{Shaded}
\begin{Highlighting}[]
\NormalTok{\textless{}s\textgreater{} a b \textless{}/s\textgreater{}}
\NormalTok{\textless{}s\textgreater{} a c \textless{}/s\textgreater{}}
\end{Highlighting}
\end{Shaded}

The bigram counts after context \texttt{a} are

\[
C(a,b)=1,
\qquad
C(a,c)=1,
\qquad
C(a)=2.
\]

With no smoothing,

\[
\widehat p(b\mid a)=\frac{1}{2},
\qquad
\widehat p(c\mid a)=\frac{1}{2}.
\]

If the vocabulary is

\[
V=\{a,b,c,d,\langle/ s\rangle\},
\]

then \(|V|=5\). With add-\(\alpha\) smoothing and \(\alpha=1\),

\[
\widehat p_1(b\mid a)=\frac{1+1}{2+5}=\frac{2}{7},
\]

\[
\widehat p_1(c\mid a)=\frac{2}{7},
\]

and the unseen word \(d\) receives

\[
\widehat p_1(d\mid a)=\frac{0+1}{2+5}=\frac{1}{7}.
\]

This example shows both the benefit and cost of smoothing: unseen events
get nonzero probability, but observed events lose some probability mass.

\section{AI-assisted learning
component}\label{ai-assisted-learning-component-1}

AI systems can be helpful for studying \(n\)-gram models, but they
should be used as mathematical assistants rather than answer engines.

\subsection{Prompt 1: audit an
assumption}\label{prompt-1-audit-an-assumption}

\begin{Shaded}
\begin{Highlighting}[]
\NormalTok{I am studying n{-}gram language models. Explain the Markov assumption in a bigram model. Give one example where it works reasonably well and one example where it fails because of long{-}range dependence. Do not use code.}
\end{Highlighting}
\end{Shaded}

After reading the response, check whether the explanation distinguishes
clearly between an exact identity and an approximation.

\subsection{Prompt 2: verify a
derivation}\label{prompt-2-verify-a-derivation}

\begin{Shaded}
\begin{Highlighting}[]
\NormalTok{Derive the maximum likelihood estimator for P(v | h) in an n{-}gram model using a Lagrange multiplier. Use counts C(h,v). Then explain why the estimator is undefined when C(h)=0.}
\end{Highlighting}
\end{Shaded}

You should verify that the final answer is

\[
\widehat p(v\mid h)=\frac{C(h,v)}{C(h)}.
\]

\subsection{Prompt 3: critique
smoothing}\label{prompt-3-critique-smoothing}

\begin{Shaded}
\begin{Highlighting}[]
\NormalTok{In an n{-}gram model, add{-}alpha smoothing gives P\_alpha(v | h) = (C(h,v)+alpha)/(C(h)+alpha |V|). Explain this formula as shrinkage toward the uniform distribution. What problem does it solve, and what new problem can it introduce?}
\end{Highlighting}
\end{Shaded}

A strong answer should mention zero probabilities, probability mass
redistribution, and the fact that large vocabularies can make add-one
smoothing too strong.

\subsection{Prompt 4: connect to modern language
models}\label{prompt-4-connect-to-modern-language-models}

\begin{Shaded}
\begin{Highlighting}[]
\NormalTok{Compare an n{-}gram language model and a transformer language model from a mathematical point of view. Focus on how each estimates P(w\_t | context), how the context is represented, and why transformers can use longer{-}range dependencies.}
\end{Highlighting}
\end{Shaded}

This helps students see why Chapter 5 belongs in a modern NLP book
rather than only in a historical survey.

\section{Mathematical summary}\label{mathematical-summary-1}

The main formulas of this chapter are:

\[
P(w_{1:T})
=
\prod_{t=1}^{T}P(w_t\mid w_{1:t-1})
\]

for the exact chain rule,

\[
P(w_t\mid w_{1:t-1})
\approx
P(w_t\mid w_{t-n+1:t-1})
\]

for the \(n\)-gram Markov approximation,

\[
\widehat p(v\mid h)=\frac{C(h,v)}{C(h)}
\]

for maximum likelihood estimation,

\[
\widehat p_\alpha(v\mid h)=\frac{C(h,v)+\alpha}{C(h)+\alpha |V|}
\]

for add-\(\alpha\) smoothing, and

\[
\operatorname{Perplexity}
=
\exp\left(-\frac{1}{N}\sum_{i=1}^{N}\log p_i\right)
\]

for evaluation.

\section{Exercises}\label{exercises-4}

\subsection{Conceptual exercises}\label{conceptual-exercises-3}

\begin{enumerate}
\def\labelenumi{\arabic{enumi}.}
\tightlist
\item
  Explain in your own words why the exact chain rule is not enough to
  build a practical language model.
\item
  What is the difference between a bigram model and a trigram model?
\item
  Why does increasing \(n\) reduce bias but increase variance?
\item
  Why can an unseen bigram force an entire sentence to have probability
  zero?
\item
  Explain add-\(\alpha\) smoothing as a convex combination of the
  empirical distribution and the uniform distribution.
\item
  Why is perplexity related to average uncertainty?
\item
  Give an example of a sentence where a bigram model misses important
  long-range information.
\end{enumerate}

\subsection{Mathematical exercises}\label{mathematical-exercises-2}

\begin{enumerate}
\def\labelenumi{\arabic{enumi}.}
\tightlist
\item
  Let \(C(h,v_1)=3\), \(C(h,v_2)=2\), and \(C(h,v_3)=0\). Compute the
  MLE probabilities.
\item
  For the same counts, compute add-one smoothed probabilities.
\item
  For the same counts, compute add-\(0.5\) smoothed probabilities.
\item
  Prove that \(\widehat p_\alpha(\cdot\mid h)\) sums to \(1\).
\item
  Derive the MLE estimator for \(p(v\mid h)\) using a Lagrange
  multiplier.
\item
  Suppose a model assigns probabilities \(0.5,0.25,0.1,0.4\) to four
  correct next tokens. Compute the average negative log-likelihood and
  perplexity using natural logarithms.
\item
  Show that a bigram transition matrix is row-stochastic.
\item
  Suppose \(\pi_t\) is a row vector distribution over tokens and \(P\)
  is a bigram transition matrix. Explain why \(\pi_{t+1}=\pi_tP\) is
  again a probability distribution.
\end{enumerate}

\subsection{Python exercises}\label{python-exercises-1}

\begin{enumerate}
\def\labelenumi{\arabic{enumi}.}
\tightlist
\item
  Modify the code above to build a trigram model.
\item
  Compare bigram and trigram perplexity on a small train/test split.
\item
  Implement a function that returns the top \(k\) most likely next
  tokens for a given context.
\item
  Generate ten sentences using different values of \(\alpha\). Describe
  how the output changes.
\item
  Implement interpolation between unigram and bigram probabilities.
\item
  Build a heatmap of the bigram transition matrix for a selected
  vocabulary.
\item
  Write a function that identifies contexts with the highest entropy.
\end{enumerate}

\section{Mini-project: small-domain language
model}\label{mini-project-small-domain-language-model}

Choose a small text domain, such as course descriptions, research
abstracts, textbook section titles, or advising questions. Build
unigram, bigram, and trigram models.

Your report should include:

\begin{enumerate}
\def\labelenumi{\arabic{enumi}.}
\tightlist
\item
  the tokenization rule;
\item
  the vocabulary size;
\item
  the number of observed contexts for each model order;
\item
  several next-token predictions for selected contexts;
\item
  generated text samples;
\item
  perplexity on a held-out set;
\item
  a discussion of where the Markov assumption succeeds and fails.
\end{enumerate}

For a mathematically stronger project, include the transition matrix for
the bigram model and analyze row entropies.

\section{Lab connection}\label{lab-connection-2}

The corresponding computer lab should guide students through:

\begin{itemize}
\tightlist
\item
  building \(n\)-gram counts from scratch;
\item
  deriving and implementing MLE probabilities;
\item
  implementing add-\(\alpha\) smoothing;
\item
  computing sentence log-probability and perplexity;
\item
  generating text by sequential sampling;
\item
  visualizing smoothing effects with Plotly;
\item
  comparing unigram, bigram, and trigram models.
\end{itemize}

This lab is an important bridge between probability theory and modern
neural language modeling. The next chapters will use the same
probabilistic ideas in text classification and representation learning.

\chapter{Naive Bayes and Text
Classification}\label{naive-bayes-and-text-classification}

Bayes decision rules, multinomial likelihoods, smoothing, and linear
log-odds

\hfill\break

Text classification is the supervised learning problem of assigning a
label to a document. A document may be an email, review, abstract, news
article, student question, or clinical note. In this chapter we develop
one of the most important classical models for this task: \textbf{Naive
Bayes}. The model is simple, fast, interpretable, and mathematically
transparent. Its main lesson is that a probabilistic model for documents
can lead to a linear decision rule in word-count features.

\section{Learning goals}\label{learning-goals-5}

After studying this chapter, you should be able to:

\begin{enumerate}
\def\labelenumi{\arabic{enumi}.}
\tightlist
\item
  formulate text classification as a statistical decision problem;
\item
  derive the Bayes classifier under zero-one loss;
\item
  define the multinomial Naive Bayes model for count vectors;
\item
  estimate class priors and class-conditional token probabilities from
  data;
\item
  explain the conditional independence assumption and why it is called
  ``naive'';
\item
  derive the log-score form \[
  s_c(x)=\log \widehat \pi_c+\sum_{j=1}^d x_j\log \widehat \theta_{cj};
  \]
\item
  show that Naive Bayes is a linear classifier in count-vector space;
\item
  implement Naive Bayes from scratch in Python;
\item
  use add-\(\alpha\) smoothing to avoid zero probabilities;
\item
  evaluate a text classifier with confusion matrices, precision, recall,
  and \(F_1\) score;
\item
  use AI tools to critique assumptions, detect data leakage, and explain
  classifier behavior.
\end{enumerate}

\section{Why Naive Bayes matters}\label{why-naive-bayes-matters}

Naive Bayes is not a modern large language model. However, it remains a
useful mathematical baseline for NLP because it makes the modeling
assumptions explicit.

The model answers a clean question:

\begin{quote}
If documents in class \(c\) generate tokens according to a
class-specific probability distribution, which class most likely
generated a new document?
\end{quote}

This question leads naturally to:

\begin{itemize}
\tightlist
\item
  Bayes' theorem;
\item
  maximum likelihood estimation;
\item
  multinomial distributions;
\item
  Dirichlet/additive smoothing;
\item
  log-likelihood ratios;
\item
  linear decision boundaries;
\item
  interpretable feature weights.
\end{itemize}

Modern neural classifiers often achieve higher accuracy, but Naive Bayes
gives a clear first model against which more complicated systems should
be compared.

\textbf{Central idea.} A count vector
\(x=(x_1,\ldots,x_d)\in \mathbb{N}^d\) records how often each vocabulary
token appears in a document. Multinomial Naive Bayes estimates

\[
P(Y=c\mid X=x)
\]

by combining a class prior \(P(Y=c)\) with a class-conditional token
model \(P(X=x\mid Y=c)\).

\section{Text classification as a supervised learning
problem}\label{text-classification-as-a-supervised-learning-problem}

Let

\[
\mathcal{Y}=\{1,2,\ldots,K\}
\]

be a finite label set. Examples include:

{\def\LTcaptype{none} % do not increment counter
\begin{longtable}[]{@{}
  >{\raggedright\arraybackslash}p{(\linewidth - 4\tabcolsep) * \real{0.3333}}
  >{\raggedright\arraybackslash}p{(\linewidth - 4\tabcolsep) * \real{0.3333}}
  >{\raggedright\arraybackslash}p{(\linewidth - 4\tabcolsep) * \real{0.3333}}@{}}
\toprule\noalign{}
\begin{minipage}[b]{\linewidth}\raggedright
Task
\end{minipage} & \begin{minipage}[b]{\linewidth}\raggedright
Input
\end{minipage} & \begin{minipage}[b]{\linewidth}\raggedright
Label set
\end{minipage} \\
\midrule\noalign{}
\endhead
\bottomrule\noalign{}
\endlastfoot
sentiment analysis & movie review &
\(\{\text{negative},\text{positive}\}\) \\
spam detection & email & \(\{\text{ham},\text{spam}\}\) \\
topic classification & news article &
\(\{\text{sports},\text{business},\text{science},\ldots\}\) \\
intent classification & student question &
\(\{\text{transfer credit},\text{course registration},\text{visa},\ldots\}\) \\
\end{longtable}
}

A training dataset is

\[
D_n=\{(d_i,y_i)\}_{i=1}^n,
\]

where \(d_i\) is a document and \(y_i\in \mathcal{Y}\) is its label.
After choosing a representation map

\[
\phi: \text{documents}\longrightarrow \mathbb{R}^d,
\]

we obtain feature vectors

\[
x_i=\phi(d_i)\in \mathbb{R}^d.
\]

For Naive Bayes, the most common representation is a count vector:

\[
x_i=(x_{i1},x_{i2},\ldots,x_{id})\in \mathbb{N}^d,
\]

where \(x_{ij}\) is the number of times vocabulary token \(v_j\) appears
in document \(d_i\).

The goal is to learn a classifier

\[
f:\mathbb{N}^d\longrightarrow \mathcal{Y}.
\]

\section{The Bayes classifier}\label{the-bayes-classifier}

Assume \((X,Y)\) is a random pair with \(X\in \mathbb{N}^d\) and
\(Y\in \mathcal{Y}\). Under zero-one loss,

\[
L(\widehat y,y)=
\begin{cases}
0, & \widehat y=y,\\
1, & \widehat y\neq y,
\end{cases}
\]

the best possible prediction rule, if the true distribution were known,
is the \textbf{Bayes classifier}:

\[
f^*(x)=\arg\max_{c\in \mathcal{Y}} P(Y=c\mid X=x).
\]

This rule chooses the class with the largest posterior probability.

By Bayes' theorem,

\[
P(Y=c\mid X=x)
=
\frac{P(X=x\mid Y=c)P(Y=c)}{P(X=x)}.
\]

For classification, the denominator \(P(X=x)\) does not depend on \(c\).
Therefore,

\[
f^*(x)=\arg\max_{c\in \mathcal{Y}} P(X=x\mid Y=c)P(Y=c).
\]

This is the maximum a posteriori, or MAP, decision rule.

\textbf{Bayes decision rule.} To classify \(x\), compare the posterior
scores

\[
P(X=x\mid Y=c)P(Y=c), \qquad c=1,\ldots,K.
\]

Naive Bayes becomes a usable classifier after we specify and estimate
the class-conditional distribution \(P(X=x\mid Y=c)\).

\section{The multinomial event model}\label{the-multinomial-event-model}

Let

\[
V=\{v_1,v_2,\ldots,v_d\}
\]

be the vocabulary. In the multinomial Naive Bayes model, each class
\(c\) has a token probability vector

\[
\theta_c=(\theta_{c1},\theta_{c2},\ldots,\theta_{cd})\in \Delta^{d-1},
\]

where

\[
\theta_{cj}\ge 0, \qquad \sum_{j=1}^d \theta_{cj}=1.
\]

The interpretation is:

\[
\theta_{cj}=P(\text{a randomly drawn token is }v_j\mid Y=c).
\]

Now suppose a document has length

\[
m=\sum_{j=1}^d x_j.
\]

The multinomial model assumes that, conditional on the class, the
document's tokens are generated by repeated draws from \(\theta_c\). The
probability of observing the count vector \(x\) is

\[
P(X=x\mid Y=c)
=
\frac{m!}{x_1!x_2!\cdots x_d!}
\prod_{j=1}^d \theta_{cj}^{x_j}.
\]

The multinomial coefficient counts how many token sequences have the
same count vector.

For classification, if the document \(x\) is fixed, the coefficient

\[
\frac{m!}{x_1!x_2!\cdots x_d!}
\]

does not depend on the class \(c\). Therefore the MAP classifier can
ignore it:

\[
\widehat y
=
\arg\max_{c\in\mathcal{Y}}
\pi_c\prod_{j=1}^d \theta_{cj}^{x_j},
\]

where \(\pi_c=P(Y=c)\) is the class prior.

\section{The conditional independence
assumption}\label{the-conditional-independence-assumption}

The word ``naive'' refers to the conditional independence assumption. In
sequence form, if a document has tokens

\[
w_1,w_2,\ldots,w_m,
\]

the model assumes

\[
P(w_1,\ldots,w_m\mid Y=c)
\approx
\prod_{t=1}^m P(w_t\mid Y=c).
\]

That is, after conditioning on the class label, the tokens are treated
as independent draws from a class-specific distribution.

This is a strong assumption. Natural language has grammar, word order,
collocations, phrases, and long-range dependence. For example, in real
text,

\begin{Shaded}
\begin{Highlighting}[]
\NormalTok{not good}
\end{Highlighting}
\end{Shaded}

is not the same as the two independent words \texttt{not} and
\texttt{good}. Yet Naive Bayes often works surprisingly well for topic
and sentiment tasks because the presence of highly class-indicative
words can dominate the classification score.

\begin{tcolorbox}[enhanced jigsaw, arc=.35mm, bottomrule=.15mm, bottomtitle=1mm, breakable, colback=white, colbacktitle=quarto-callout-note-color!10!white, colframe=quarto-callout-note-color-frame, coltitle=black, left=2mm, leftrule=.75mm, opacityback=0, opacitybacktitle=0.6, rightrule=.15mm, title=\textcolor{quarto-callout-note-color}{\faInfo}\hspace{0.5em}{Note}, titlerule=0mm, toprule=.15mm, toptitle=1mm]

The Naive Bayes assumption is usually false as a generative model of
language. The important question is not whether the model is literally
true, but whether the resulting decision rule is useful and
interpretable for a particular task.

\end{tcolorbox}

\section{Estimating the parameters}\label{estimating-the-parameters}

Suppose we have training examples

\[
(x_i,y_i),\qquad i=1,\ldots,n.
\]

\subsection{Class priors}\label{class-priors}

The maximum likelihood estimate of the class prior is

\[
\widehat \pi_c
=
\frac{n_c}{n},
\]

where

\[
n_c=\sum_{i=1}^n \mathbf{1}\{y_i=c\}
\]

is the number of training documents in class \(c\).

\subsection{Token probabilities}\label{token-probabilities}

Let

\[
N_{cj}=\sum_{i:y_i=c} x_{ij}
\]

be the total number of times token \(v_j\) appears in documents with
label \(c\). Let

\[
N_c=\sum_{j=1}^d N_{cj}
\]

be the total number of tokens in class \(c\). The maximum likelihood
estimate is

\[
\widehat \theta_{cj}
=
\frac{N_{cj}}{N_c}.
\]

\subsection{Derivation by constrained
optimization}\label{derivation-by-constrained-optimization}

For a fixed class \(c\), the likelihood for \(\theta_c\) is proportional
to

\[
L_c(\theta_c)
=
\prod_{j=1}^d \theta_{cj}^{N_{cj}}.
\]

The log-likelihood is

\[
\ell_c(\theta_c)
=
\sum_{j=1}^d N_{cj}\log \theta_{cj},
\]

subject to

\[
\sum_{j=1}^d \theta_{cj}=1.
\]

Using a Lagrange multiplier \(\lambda\), define

\[
\mathcal{L}(\theta_c,\lambda)
=
\sum_{j=1}^d N_{cj}\log \theta_{cj}
+\lambda\left(1-\sum_{j=1}^d \theta_{cj}\right).
\]

The stationarity condition gives

\[
\frac{\partial \mathcal{L}}{\partial \theta_{cj}}
=
\frac{N_{cj}}{\theta_{cj}}-\lambda=0.
\]

Hence

\[
\theta_{cj}=\frac{N_{cj}}{\lambda}.
\]

Summing over \(j\) gives

\[
1=\sum_{j=1}^d \theta_{cj}=\frac{1}{\lambda}\sum_{j=1}^d N_{cj}=\frac{N_c}{\lambda},
\]

so \(\lambda=N_c\). Therefore

\[
\widehat \theta_{cj}=\frac{N_{cj}}{N_c}.
\]

\section{\texorpdfstring{Zero probabilities and add-\(\alpha\)
smoothing}{Zero probabilities and add-\textbackslash alpha smoothing}}\label{zero-probabilities-and-add-alpha-smoothing}

The maximum likelihood estimator assigns probability zero to any token
that never appears in class \(c\):

\[
N_{cj}=0 \quad \Longrightarrow \quad \widehat \theta_{cj}=0.
\]

This creates a serious problem. If a test document contains even one
such token, then

\[
\prod_{j=1}^d \widehat \theta_{cj}^{x_j}=0
\]

for that class. The model may rule out a class entirely because of one
unseen word.

To avoid this, we use add-\(\alpha\) smoothing:

\[
\widehat \theta_{cj}^{(\alpha)}
=
\frac{N_{cj}+\alpha}{N_c+\alpha d},
\qquad \alpha>0.
\]

Common choices include \(\alpha=1\) for Laplace smoothing and smaller
values such as \(\alpha=0.1\) or \(\alpha=0.01\) for less aggressive
smoothing.

\subsection{Bayesian interpretation}\label{bayesian-interpretation}

Add-\(\alpha\) smoothing can be interpreted through a symmetric
Dirichlet prior on each class token vector:

\[
\theta_c\sim \operatorname{Dirichlet}(\alpha,\ldots,\alpha).
\]

The smoothed estimate

\[
\frac{N_{cj}+\alpha}{N_c+\alpha d}
\]

is the posterior predictive probability of the next token being \(v_j\)
after observing the class-specific token counts.

\textbf{Smoothing as regularization.} Add-\(\alpha\) smoothing shrinks
the estimated token distribution away from the boundary of the simplex.
No token receives probability zero, so every class remains possible for
every document.

\section{Log-scores and numerical
stability}\label{log-scores-and-numerical-stability}

The product

\[
\pi_c\prod_{j=1}^d \theta_{cj}^{x_j}
\]

can be extremely small. In computation, multiplying many small
probabilities may underflow to zero. Taking logarithms avoids this
problem and converts products into sums:

\[
s_c(x)
=
\log \pi_c+
\sum_{j=1}^d x_j\log \theta_{cj}.
\]

The classifier is

\[
\widehat y
=
\arg\max_{c\in\mathcal{Y}} s_c(x).
\]

Because the logarithm is increasing, maximizing the log-score is
equivalent to maximizing the original posterior numerator.

\section{Naive Bayes as a linear
classifier}\label{naive-bayes-as-a-linear-classifier}

The score for class \(c\) can be written as

\[
s_c(x)=b_c+w_c^\top x,
\]

where

\[
b_c=\log \pi_c,
\qquad
w_c=(\log \theta_{c1},\ldots,\log \theta_{cd})^\top.
\]

Thus multinomial Naive Bayes is a linear classifier in count-vector
space.

For binary classification with classes \(1\) and \(0\), compare
log-scores:

\[
s_1(x)-s_0(x)
=
\log\frac{\pi_1}{\pi_0}
+
\sum_{j=1}^d x_j\log\frac{\theta_{1j}}{\theta_{0j}}.
\]

The decision rule is

\[
\widehat y=1
\quad \Longleftrightarrow \quad
\log\frac{\pi_1}{\pi_0}
+
\sum_{j=1}^d x_j\log\frac{\theta_{1j}}{\theta_{0j}}>0.
\]

The decision boundary is the hyperplane

\[
\log\frac{\pi_1}{\pi_0}
+
\sum_{j=1}^d x_j\log\frac{\theta_{1j}}{\theta_{0j}}=0.
\]

\begin{tcolorbox}[enhanced jigsaw, arc=.35mm, bottomrule=.15mm, bottomtitle=1mm, breakable, colback=white, colbacktitle=quarto-callout-tip-color!10!white, colframe=quarto-callout-tip-color-frame, coltitle=black, left=2mm, leftrule=.75mm, opacityback=0, opacitybacktitle=0.6, rightrule=.15mm, title=\textcolor{quarto-callout-tip-color}{\faLightbulb}\hspace{0.5em}{Tip}, titlerule=0mm, toprule=.15mm, toptitle=1mm]

The weights

\[
\log\frac{\theta_{1j}}{\theta_{0j}}
\]

are interpretable. A large positive value means token \(v_j\) strongly
favors class \(1\) over class \(0\). A large negative value means the
token strongly favors class \(0\).

\end{tcolorbox}

\section{Relation to logistic
regression}\label{relation-to-logistic-regression}

Both Naive Bayes and logistic regression can produce linear decision
boundaries, but they estimate different objects.

Naive Bayes is a \textbf{generative} model. It estimates

\[
P(Y=c),
\qquad
P(X=x\mid Y=c),
\]

and then uses Bayes' theorem to obtain a classifier.

Logistic regression is a \textbf{discriminative} model. It estimates the
conditional probability directly:

\[
P(Y=c\mid X=x).
\]

For binary classification, logistic regression models

\[
P(Y=1\mid X=x)
=
\sigma(\beta_0+\beta^\top x),
\]

where

\[
\sigma(z)=\frac{1}{1+e^{-z}}.
\]

Naive Bayes may work well with small data because it imposes strong
probabilistic structure. Logistic regression may work better with enough
labeled data because it directly optimizes conditional prediction.

\section{Bernoulli Naive Bayes}\label{bernoulli-naive-bayes}

The multinomial model uses token counts. Another common model uses
binary features:

\[
z_j=\mathbf{1}\{\text{token }v_j\text{ appears in the document}\}.
\]

For each class \(c\), Bernoulli Naive Bayes estimates

\[
\rho_{cj}=P(Z_j=1\mid Y=c).
\]

The class-conditional probability is

\[
P(Z=z\mid Y=c)
=
\prod_{j=1}^d \rho_{cj}^{z_j}(1-\rho_{cj})^{1-z_j}.
\]

The log-score becomes

\[
s_c(z)
=
\log \pi_c+
\sum_{j=1}^d
\left[z_j\log \rho_{cj}+(1-z_j)\log(1-\rho_{cj})\right].
\]

Bernoulli Naive Bayes models word presence or absence; multinomial Naive
Bayes models word counts. For longer documents, multinomial Naive Bayes
is often more natural. For short texts, binary presence features can be
useful.

\section{Worked example by hand}\label{worked-example-by-hand}

Suppose the vocabulary is

\[
V=\{\texttt{win},\texttt{meeting},\texttt{free}\}.
\]

There are two classes:

\[
\mathcal{Y}=\{\text{spam},\text{ham}\}.
\]

Assume the estimated parameters are:

\[
\widehat \pi_{\text{spam}}=0.4,
\qquad
\widehat \pi_{\text{ham}}=0.6,
\]

and

{\def\LTcaptype{none} % do not increment counter
\begin{longtable}[]{@{}
  >{\raggedright\arraybackslash}p{(\linewidth - 4\tabcolsep) * \real{0.2727}}
  >{\raggedleft\arraybackslash}p{(\linewidth - 4\tabcolsep) * \real{0.3636}}
  >{\raggedleft\arraybackslash}p{(\linewidth - 4\tabcolsep) * \real{0.3636}}@{}}
\toprule\noalign{}
\begin{minipage}[b]{\linewidth}\raggedright
token
\end{minipage} & \begin{minipage}[b]{\linewidth}\raggedleft
\(\widehat\theta_{\text{spam},j}\)
\end{minipage} & \begin{minipage}[b]{\linewidth}\raggedleft
\(\widehat\theta_{\text{ham},j}\)
\end{minipage} \\
\midrule\noalign{}
\endhead
\bottomrule\noalign{}
\endlastfoot
\texttt{win} & 0.50 & 0.05 \\
\texttt{meeting} & 0.05 & 0.70 \\
\texttt{free} & 0.45 & 0.25 \\
\end{longtable}
}

Consider the document

\begin{Shaded}
\begin{Highlighting}[]
\NormalTok{free win free}
\end{Highlighting}
\end{Shaded}

Its count vector is

\[
x=(1,0,2).
\]

The spam log-score is

\[
\log(0.4)+1\log(0.50)+0\log(0.05)+2\log(0.45).
\]

The ham log-score is

\[
\log(0.6)+1\log(0.05)+0\log(0.70)+2\log(0.25).
\]

Since \texttt{win} and \texttt{free} are much more likely under the spam
token distribution, the spam score will be larger.

\section{Python example: training from
scratch}\label{python-example-training-from-scratch}

The following code implements multinomial Naive Bayes without using a
machine learning library. The point is to see the mathematical
quantities directly.

\begin{Shaded}
\begin{Highlighting}[]
\ImportTok{import}\NormalTok{ re}
\ImportTok{from}\NormalTok{ collections }\ImportTok{import}\NormalTok{ Counter, defaultdict}
\ImportTok{import}\NormalTok{ numpy }\ImportTok{as}\NormalTok{ np}
\ImportTok{import}\NormalTok{ pandas }\ImportTok{as}\NormalTok{ pd}

\NormalTok{train\_docs }\OperatorTok{=}\NormalTok{ [}
\NormalTok{    (}\StringTok{"free money win now"}\NormalTok{, }\StringTok{"spam"}\NormalTok{),}
\NormalTok{    (}\StringTok{"win a free prize"}\NormalTok{, }\StringTok{"spam"}\NormalTok{),}
\NormalTok{    (}\StringTok{"cheap offer free entry"}\NormalTok{, }\StringTok{"spam"}\NormalTok{),}
\NormalTok{    (}\StringTok{"meeting schedule for project"}\NormalTok{, }\StringTok{"ham"}\NormalTok{),}
\NormalTok{    (}\StringTok{"project meeting tomorrow"}\NormalTok{, }\StringTok{"ham"}\NormalTok{),}
\NormalTok{    (}\StringTok{"please review the project schedule"}\NormalTok{, }\StringTok{"ham"}\NormalTok{),}
\NormalTok{]}

\NormalTok{test\_docs }\OperatorTok{=}\NormalTok{ [}
    \StringTok{"free project prize"}\NormalTok{,}
    \StringTok{"meeting tomorrow schedule"}\NormalTok{,}
    \StringTok{"win free offer"}\NormalTok{,}
    \StringTok{"please review money"}\NormalTok{,}
\NormalTok{]}

\KeywordTok{def}\NormalTok{ tokenize(text):}
    \ControlFlowTok{return}\NormalTok{ re.findall(}\VerbatimStringTok{r"}\PreprocessorTok{[a{-}z]}\OperatorTok{+}\VerbatimStringTok{"}\NormalTok{, text.lower())}

\CommentTok{\# Build vocabulary from training documents.}
\NormalTok{vocab }\OperatorTok{=} \BuiltInTok{sorted}\NormalTok{(\{tok }\ControlFlowTok{for}\NormalTok{ doc, label }\KeywordTok{in}\NormalTok{ train\_docs }\ControlFlowTok{for}\NormalTok{ tok }\KeywordTok{in}\NormalTok{ tokenize(doc)\})}
\NormalTok{token\_to\_index }\OperatorTok{=}\NormalTok{ \{tok: j }\ControlFlowTok{for}\NormalTok{ j, tok }\KeywordTok{in} \BuiltInTok{enumerate}\NormalTok{(vocab)\}}
\NormalTok{classes }\OperatorTok{=} \BuiltInTok{sorted}\NormalTok{(\{label }\ControlFlowTok{for}\NormalTok{ \_, label }\KeywordTok{in}\NormalTok{ train\_docs\})}

\BuiltInTok{print}\NormalTok{(}\StringTok{"Vocabulary:"}\NormalTok{, vocab)}
\BuiltInTok{print}\NormalTok{(}\StringTok{"Classes:"}\NormalTok{, classes)}
\end{Highlighting}
\end{Shaded}

\begin{verbatim}
Vocabulary: ['a', 'cheap', 'entry', 'for', 'free', 'meeting', 'money', 'now', 'offer', 'please', 'prize', 'project', 'review', 'schedule', 'the', 'tomorrow', 'win']
Classes: ['ham', 'spam']
\end{verbatim}

Now estimate the class priors and class-conditional token distributions.

\begin{Shaded}
\begin{Highlighting}[]
\NormalTok{alpha }\OperatorTok{=} \FloatTok{1.0}
\NormalTok{K }\OperatorTok{=} \BuiltInTok{len}\NormalTok{(classes)}
\NormalTok{d }\OperatorTok{=} \BuiltInTok{len}\NormalTok{(vocab)}

\NormalTok{class\_doc\_counts }\OperatorTok{=}\NormalTok{ Counter(label }\ControlFlowTok{for}\NormalTok{ \_, label }\KeywordTok{in}\NormalTok{ train\_docs)}
\NormalTok{class\_token\_counts }\OperatorTok{=}\NormalTok{ \{c: np.zeros(d, dtype}\OperatorTok{=}\BuiltInTok{float}\NormalTok{) }\ControlFlowTok{for}\NormalTok{ c }\KeywordTok{in}\NormalTok{ classes\}}
\NormalTok{class\_total\_tokens }\OperatorTok{=}\NormalTok{ \{c: }\FloatTok{0.0} \ControlFlowTok{for}\NormalTok{ c }\KeywordTok{in}\NormalTok{ classes\}}

\ControlFlowTok{for}\NormalTok{ doc, label }\KeywordTok{in}\NormalTok{ train\_docs:}
    \ControlFlowTok{for}\NormalTok{ tok }\KeywordTok{in}\NormalTok{ tokenize(doc):}
\NormalTok{        j }\OperatorTok{=}\NormalTok{ token\_to\_index[tok]}
\NormalTok{        class\_token\_counts[label][j] }\OperatorTok{+=} \DecValTok{1}
\NormalTok{        class\_total\_tokens[label] }\OperatorTok{+=} \DecValTok{1}

\NormalTok{priors }\OperatorTok{=}\NormalTok{ \{c: class\_doc\_counts[c] }\OperatorTok{/} \BuiltInTok{len}\NormalTok{(train\_docs) }\ControlFlowTok{for}\NormalTok{ c }\KeywordTok{in}\NormalTok{ classes\}}
\NormalTok{theta }\OperatorTok{=}\NormalTok{ \{\}}
\ControlFlowTok{for}\NormalTok{ c }\KeywordTok{in}\NormalTok{ classes:}
\NormalTok{    theta[c] }\OperatorTok{=}\NormalTok{ (class\_token\_counts[c] }\OperatorTok{+}\NormalTok{ alpha) }\OperatorTok{/}\NormalTok{ (class\_total\_tokens[c] }\OperatorTok{+}\NormalTok{ alpha }\OperatorTok{*}\NormalTok{ d)}

\NormalTok{prior\_table }\OperatorTok{=}\NormalTok{ pd.DataFrame(\{}\StringTok{"class"}\NormalTok{: classes, }\StringTok{"prior"}\NormalTok{: [priors[c] }\ControlFlowTok{for}\NormalTok{ c }\KeywordTok{in}\NormalTok{ classes]\})}
\BuiltInTok{print}\NormalTok{(prior\_table)}

\NormalTok{theta\_table }\OperatorTok{=}\NormalTok{ pd.DataFrame(theta, index}\OperatorTok{=}\NormalTok{vocab)}
\BuiltInTok{print}\NormalTok{(theta\_table.}\BuiltInTok{round}\NormalTok{(}\DecValTok{3}\NormalTok{))}
\end{Highlighting}
\end{Shaded}

\begin{verbatim}
  class  prior
0   ham    0.5
1  spam    0.5
            ham   spam
a         0.034  0.069
cheap     0.034  0.069
entry     0.034  0.069
for       0.069  0.034
free      0.034  0.138
meeting   0.103  0.034
money     0.034  0.069
now       0.034  0.069
offer     0.034  0.069
please    0.069  0.034
prize     0.034  0.069
project   0.138  0.034
review    0.069  0.034
schedule  0.103  0.034
the       0.069  0.034
tomorrow  0.069  0.034
win       0.034  0.103
\end{verbatim}

Define prediction using log-scores.

\begin{Shaded}
\begin{Highlighting}[]
\KeywordTok{def}\NormalTok{ vectorize(doc):}
\NormalTok{    x }\OperatorTok{=}\NormalTok{ np.zeros(d, dtype}\OperatorTok{=}\BuiltInTok{float}\NormalTok{)}
    \ControlFlowTok{for}\NormalTok{ tok }\KeywordTok{in}\NormalTok{ tokenize(doc):}
        \ControlFlowTok{if}\NormalTok{ tok }\KeywordTok{in}\NormalTok{ token\_to\_index:}
\NormalTok{            x[token\_to\_index[tok]] }\OperatorTok{+=} \DecValTok{1}
    \ControlFlowTok{return}\NormalTok{ x}

\KeywordTok{def}\NormalTok{ log\_scores(doc):}
\NormalTok{    x }\OperatorTok{=}\NormalTok{ vectorize(doc)}
\NormalTok{    scores }\OperatorTok{=}\NormalTok{ \{\}}
    \ControlFlowTok{for}\NormalTok{ c }\KeywordTok{in}\NormalTok{ classes:}
\NormalTok{        scores[c] }\OperatorTok{=}\NormalTok{ np.log(priors[c]) }\OperatorTok{+}\NormalTok{ np.}\BuiltInTok{sum}\NormalTok{(x }\OperatorTok{*}\NormalTok{ np.log(theta[c]))}
    \ControlFlowTok{return}\NormalTok{ scores}

\KeywordTok{def}\NormalTok{ predict(doc):}
\NormalTok{    scores }\OperatorTok{=}\NormalTok{ log\_scores(doc)}
    \ControlFlowTok{return} \BuiltInTok{max}\NormalTok{(scores, key}\OperatorTok{=}\NormalTok{scores.get), scores}

\ControlFlowTok{for}\NormalTok{ doc }\KeywordTok{in}\NormalTok{ test\_docs:}
\NormalTok{    pred, scores }\OperatorTok{=}\NormalTok{ predict(doc)}
    \BuiltInTok{print}\NormalTok{(}\SpecialStringTok{f"Document: }\SpecialCharTok{\{}\NormalTok{doc}\SpecialCharTok{!r\}}\SpecialStringTok{"}\NormalTok{)}
    \BuiltInTok{print}\NormalTok{(}\StringTok{"Prediction:"}\NormalTok{, pred)}
    \BuiltInTok{print}\NormalTok{(}\StringTok{"Scores:"}\NormalTok{, \{c: }\BuiltInTok{round}\NormalTok{(scores[c], }\DecValTok{3}\NormalTok{) }\ControlFlowTok{for}\NormalTok{ c }\KeywordTok{in}\NormalTok{ classes\})}
    \BuiltInTok{print}\NormalTok{()}
\end{Highlighting}
\end{Shaded}

\begin{verbatim}
Document: 'free project prize'
Prediction: spam
Scores: {'ham': -9.409, 'spam': -8.716}

Document: 'meeting tomorrow schedule'
Prediction: ham
Scores: {'ham': -7.905, 'spam': -10.795}

Document: 'win free offer'
Prediction: spam
Scores: {'ham': -10.795, 'spam': -7.617}

Document: 'please review money'
Prediction: ham
Scores: {'ham': -9.409, 'spam': -10.102}
\end{verbatim}

\section{Python example: interpretable log-odds
weights}\label{python-example-interpretable-log-odds-weights}

For binary classification, the log-odds weight for token \(v_j\) is

\[
\log\frac{\widehat\theta_{\text{spam},j}}{\widehat\theta_{\text{ham},j}}.
\]

\begin{Shaded}
\begin{Highlighting}[]
\NormalTok{positive\_class }\OperatorTok{=} \StringTok{"spam"}
\NormalTok{negative\_class }\OperatorTok{=} \StringTok{"ham"}

\NormalTok{log\_odds }\OperatorTok{=}\NormalTok{ np.log(theta[positive\_class] }\OperatorTok{/}\NormalTok{ theta[negative\_class])}
\NormalTok{weights }\OperatorTok{=}\NormalTok{ pd.DataFrame(\{}
    \StringTok{"token"}\NormalTok{: vocab,}
    \StringTok{"log\_odds\_spam\_vs\_ham"}\NormalTok{: log\_odds}
\NormalTok{\}).sort\_values(}\StringTok{"log\_odds\_spam\_vs\_ham"}\NormalTok{, ascending}\OperatorTok{=}\VariableTok{False}\NormalTok{)}

\BuiltInTok{print}\NormalTok{(weights.}\BuiltInTok{round}\NormalTok{(}\DecValTok{3}\NormalTok{))}
\end{Highlighting}
\end{Shaded}

\begin{verbatim}
       token  log_odds_spam_vs_ham
4       free                 1.386
16       win                 1.099
6      money                 0.693
10     prize                 0.693
1      cheap                 0.693
7        now                 0.693
0          a                 0.693
2      entry                 0.693
8      offer                 0.693
3        for                -0.693
9     please                -0.693
12    review                -0.693
14       the                -0.693
15  tomorrow                -0.693
5    meeting                -1.099
13  schedule                -1.099
11   project                -1.386
\end{verbatim}

Tokens with positive log-odds push the classifier toward \texttt{spam}.
Tokens with negative log-odds push it toward \texttt{ham}.

\section{Python example: accuracy and confusion
matrix}\label{python-example-accuracy-and-confusion-matrix}

Here is a tiny evaluation example. In a real project, train and test
sets should be separated carefully.

\begin{Shaded}
\begin{Highlighting}[]
\NormalTok{labeled\_test }\OperatorTok{=}\NormalTok{ [}
\NormalTok{    (}\StringTok{"free project prize"}\NormalTok{, }\StringTok{"spam"}\NormalTok{),}
\NormalTok{    (}\StringTok{"meeting tomorrow schedule"}\NormalTok{, }\StringTok{"ham"}\NormalTok{),}
\NormalTok{    (}\StringTok{"win free offer"}\NormalTok{, }\StringTok{"spam"}\NormalTok{),}
\NormalTok{    (}\StringTok{"please review money"}\NormalTok{, }\StringTok{"ham"}\NormalTok{),}
\NormalTok{]}

\NormalTok{confusion }\OperatorTok{=}\NormalTok{ pd.DataFrame(}\DecValTok{0}\NormalTok{, index}\OperatorTok{=}\NormalTok{classes, columns}\OperatorTok{=}\NormalTok{classes)}
\NormalTok{correct }\OperatorTok{=} \DecValTok{0}

\ControlFlowTok{for}\NormalTok{ doc, true\_label }\KeywordTok{in}\NormalTok{ labeled\_test:}
\NormalTok{    pred\_label, \_ }\OperatorTok{=}\NormalTok{ predict(doc)}
\NormalTok{    confusion.loc[true\_label, pred\_label] }\OperatorTok{+=} \DecValTok{1}
\NormalTok{    correct }\OperatorTok{+=} \BuiltInTok{int}\NormalTok{(pred\_label }\OperatorTok{==}\NormalTok{ true\_label)}

\NormalTok{accuracy }\OperatorTok{=}\NormalTok{ correct }\OperatorTok{/} \BuiltInTok{len}\NormalTok{(labeled\_test)}
\BuiltInTok{print}\NormalTok{(}\StringTok{"Accuracy:"}\NormalTok{, accuracy)}
\BuiltInTok{print}\NormalTok{(confusion)}
\end{Highlighting}
\end{Shaded}

\begin{verbatim}
Accuracy: 1.0
      ham  spam
ham     2     0
spam    0     2
\end{verbatim}

For binary classification with positive class \texttt{spam}, define:

\[
\text{precision}=\frac{TP}{TP+FP},
\qquad
\text{recall}=\frac{TP}{TP+FN},
\]

and

\[
F_1=\frac{2\cdot \text{precision}\cdot \text{recall}}{\text{precision}+\text{recall}}.
\]

\begin{Shaded}
\begin{Highlighting}[]
\NormalTok{pos }\OperatorTok{=} \StringTok{"spam"}
\NormalTok{neg }\OperatorTok{=} \StringTok{"ham"}
\NormalTok{TP }\OperatorTok{=}\NormalTok{ confusion.loc[pos, pos]}
\NormalTok{FP }\OperatorTok{=}\NormalTok{ confusion.loc[neg, pos]}
\NormalTok{FN }\OperatorTok{=}\NormalTok{ confusion.loc[pos, neg]}

\NormalTok{precision }\OperatorTok{=}\NormalTok{ TP }\OperatorTok{/}\NormalTok{ (TP }\OperatorTok{+}\NormalTok{ FP) }\ControlFlowTok{if}\NormalTok{ TP }\OperatorTok{+}\NormalTok{ FP }\OperatorTok{\textgreater{}} \DecValTok{0} \ControlFlowTok{else} \FloatTok{0.0}
\NormalTok{recall }\OperatorTok{=}\NormalTok{ TP }\OperatorTok{/}\NormalTok{ (TP }\OperatorTok{+}\NormalTok{ FN) }\ControlFlowTok{if}\NormalTok{ TP }\OperatorTok{+}\NormalTok{ FN }\OperatorTok{\textgreater{}} \DecValTok{0} \ControlFlowTok{else} \FloatTok{0.0}
\NormalTok{f1 }\OperatorTok{=} \DecValTok{2} \OperatorTok{*}\NormalTok{ precision }\OperatorTok{*}\NormalTok{ recall }\OperatorTok{/}\NormalTok{ (precision }\OperatorTok{+}\NormalTok{ recall) }\ControlFlowTok{if}\NormalTok{ precision }\OperatorTok{+}\NormalTok{ recall }\OperatorTok{\textgreater{}} \DecValTok{0} \ControlFlowTok{else} \FloatTok{0.0}

\BuiltInTok{print}\NormalTok{(\{}
    \StringTok{"precision"}\NormalTok{: }\BuiltInTok{round}\NormalTok{(}\BuiltInTok{float}\NormalTok{(precision), }\DecValTok{3}\NormalTok{),}
    \StringTok{"recall"}\NormalTok{: }\BuiltInTok{round}\NormalTok{(}\BuiltInTok{float}\NormalTok{(recall), }\DecValTok{3}\NormalTok{),}
    \StringTok{"f1"}\NormalTok{: }\BuiltInTok{round}\NormalTok{(}\BuiltInTok{float}\NormalTok{(f1), }\DecValTok{3}\NormalTok{),}
\NormalTok{\})}
\end{Highlighting}
\end{Shaded}

\begin{verbatim}
{'precision': 1.0, 'recall': 1.0, 'f1': 1.0}
\end{verbatim}

\section{Interactive HTML: smoothing and
log-odds}\label{interactive-html-smoothing-and-log-odds}

The next interactive component shows how add-\(\alpha\) smoothing
changes token probabilities and log-odds weights. The example uses two
classes and four tokens. Move the smoothing slider and token-count
sliders to see how the decision score changes.

The HTML version of this chapter contains an interactive Plotly.js
component for exploring how add-\(\alpha\) smoothing changes Naive Bayes
log-odds weights and predictions. In PDF output, the same ideas can be
studied by running the Python code examples above.

\section{Statistical interpretation of
smoothing}\label{statistical-interpretation-of-smoothing}

Smoothing changes a classifier in two ways.

First, it prevents impossible events. Every token receives positive
probability in every class:

\[
\widehat \theta_{cj}^{(\alpha)}>0.
\]

Second, it reduces the magnitude of extreme log-odds weights. If a token
appears many times in one class and zero times in another, the
unsmoothed log-ratio would be infinite. Smoothing makes the ratio
finite:

\[
\log\frac{N_{1j}+\alpha}{N_{0j}+\alpha}
+
\log\frac{N_0+\alpha d}{N_1+\alpha d}.
\]

Thus smoothing acts like a regularizer on the classifier's feature
weights.

\section{Model assumptions and failure
modes}\label{model-assumptions-and-failure-modes}

Naive Bayes works best when class-indicative tokens are strong and word
order is not essential. It can fail when meaning depends on syntax,
negation, or context.

\subsection{Failure mode 1: negation}\label{failure-mode-1-negation}

The phrases

\begin{Shaded}
\begin{Highlighting}[]
\NormalTok{good}
\end{Highlighting}
\end{Shaded}

and

\begin{Shaded}
\begin{Highlighting}[]
\NormalTok{not good}
\end{Highlighting}
\end{Shaded}

have different meanings. A unigram Naive Bayes model may treat
\texttt{good} as positive even when preceded by \texttt{not}.

Possible remedy: add bigram features such as \texttt{not\_good}.

\subsection{Failure mode 2: repeated correlated
words}\label{failure-mode-2-repeated-correlated-words}

Because Naive Bayes multiplies token probabilities, repeated correlated
words may make the model overconfident. The model behaves as though each
occurrence contributes independent evidence.

Possible remedy: use binary features, stronger smoothing, feature
selection, or a discriminative model.

\subsection{Failure mode 3: dataset
leakage}\label{failure-mode-3-dataset-leakage}

If the training data contains label artifacts, the classifier may learn
shortcuts. For example, if all spam examples contain a footer phrase not
present at deployment time, the model may appear accurate in validation
but fail in practice.

Possible remedy: inspect feature weights and validate on realistic
held-out data.

\section{AI-assisted learning
component}\label{ai-assisted-learning-component-2}

Use AI tools as a mathematical tutor and model auditor, not as a
replacement for your own derivations.

\subsection{Prompt 1: explain the
assumption}\label{prompt-1-explain-the-assumption}

\begin{Shaded}
\begin{Highlighting}[]
\NormalTok{I am studying multinomial Naive Bayes for text classification.}
\NormalTok{Explain the conditional independence assumption using probability notation.}
\NormalTok{Then give two examples where this assumption is false in natural language.}
\NormalTok{Do not use code.}
\end{Highlighting}
\end{Shaded}

After receiving an answer, check whether it correctly distinguishes

\[
P(w_1,\ldots,w_m\mid Y=c)
\]

from

\[
\prod_{t=1}^m P(w_t\mid Y=c).
\]

\subsection{Prompt 2: audit feature
weights}\label{prompt-2-audit-feature-weights}

\begin{Shaded}
\begin{Highlighting}[]
\NormalTok{Here are the largest positive and negative log{-}odds features from a Naive Bayes text classifier.}
\NormalTok{Identify which features seem semantically meaningful and which may indicate dataset leakage.}
\NormalTok{Explain your reasoning in a table.}
\end{Highlighting}
\end{Shaded}

This is useful because Naive Bayes is interpretable. Large log-odds
weights should be inspected before trusting a model.

\subsection{Prompt 3: compare models}\label{prompt-3-compare-models}

\begin{Shaded}
\begin{Highlighting}[]
\NormalTok{Compare multinomial Naive Bayes and logistic regression for count{-}vector text classification.}
\NormalTok{Use the language of generative and discriminative models.}
\NormalTok{Explain when Naive Bayes might perform well despite false assumptions.}
\end{Highlighting}
\end{Shaded}

A strong answer should mention that Naive Bayes estimates a joint model
while logistic regression estimates a conditional model.

\subsection{Prompt 4: generate
counterexamples}\label{prompt-4-generate-counterexamples}

\begin{Shaded}
\begin{Highlighting}[]
\NormalTok{Generate five short text examples where a unigram Naive Bayes sentiment classifier may fail.}
\NormalTok{For each example, identify which mathematical assumption is responsible for the failure.}
\end{Highlighting}
\end{Shaded}

This helps connect model assumptions to concrete language behavior.

\begin{tcolorbox}[enhanced jigsaw, arc=.35mm, bottomrule=.15mm, bottomtitle=1mm, breakable, colback=white, colbacktitle=quarto-callout-warning-color!10!white, colframe=quarto-callout-warning-color-frame, coltitle=black, left=2mm, leftrule=.75mm, opacityback=0, opacitybacktitle=0.6, rightrule=.15mm, title=\textcolor{quarto-callout-warning-color}{\faExclamationTriangle}\hspace{0.5em}{Warning}, titlerule=0mm, toprule=.15mm, toptitle=1mm]

When using AI tools, never paste private emails, student records, or
sensitive documents into an external model. Use synthetic or anonymized
text for examples.

\end{tcolorbox}

\section{A complete mathematical
summary}\label{a-complete-mathematical-summary}

The multinomial Naive Bayes pipeline can be summarized as follows.

\begin{enumerate}
\def\labelenumi{\arabic{enumi}.}
\tightlist
\item
  Choose a vocabulary
\end{enumerate}

\[
V=\{v_1,\ldots,v_d\}.
\]

\begin{enumerate}
\def\labelenumi{\arabic{enumi}.}
\setcounter{enumi}{1}
\tightlist
\item
  Represent each document as a count vector
\end{enumerate}

\[
x\in \mathbb{N}^d.
\]

\begin{enumerate}
\def\labelenumi{\arabic{enumi}.}
\setcounter{enumi}{2}
\tightlist
\item
  Estimate class priors
\end{enumerate}

\[
\widehat \pi_c=\frac{n_c}{n}.
\]

\begin{enumerate}
\def\labelenumi{\arabic{enumi}.}
\setcounter{enumi}{3}
\tightlist
\item
  Estimate smoothed token probabilities
\end{enumerate}

\[
\widehat \theta_{cj}
=
\frac{N_{cj}+\alpha}{N_c+\alpha d}.
\]

\begin{enumerate}
\def\labelenumi{\arabic{enumi}.}
\setcounter{enumi}{4}
\tightlist
\item
  Compute class log-scores
\end{enumerate}

\[
s_c(x)=\log \widehat \pi_c+
\sum_{j=1}^d x_j\log \widehat \theta_{cj}.
\]

\begin{enumerate}
\def\labelenumi{\arabic{enumi}.}
\setcounter{enumi}{5}
\tightlist
\item
  Predict
\end{enumerate}

\[
\widehat y=\arg\max_{c\in\mathcal{Y}}s_c(x).
\]

The classifier is probabilistic in origin but linear in its final
scoring form.

\section{Exercises}\label{exercises-5}

\begin{enumerate}
\def\labelenumi{\arabic{enumi}.}
\item
  \textbf{Bayes rule.} Starting from Bayes' theorem, derive the MAP
  classification rule \[
  \widehat y=\arg\max_c P(X=x\mid Y=c)P(Y=c).
  \]
\item
  \textbf{Multinomial likelihood.} Let \(x\in\mathbb{N}^d\) with
  \(m=\sum_j x_j\). Write the multinomial probability \(P(X=x\mid Y=c)\)
  and identify which factor can be ignored for classification.
\item
  \textbf{MLE derivation.} Use Lagrange multipliers to derive \[
  \widehat \theta_{cj}=\frac{N_{cj}}{N_c}.
  \]
\item
  \textbf{Smoothing.} Suppose \(d=1000\), \(N_c=5000\), \(N_{cj}=0\),
  and \(\alpha=1\). Compute the smoothed probability of token \(v_j\) in
  class \(c\).
\item
  \textbf{Log-odds.} In a binary classifier, show that the Naive Bayes
  decision boundary is a hyperplane in \(\mathbb{R}^d\).
\item
  \textbf{Interpretability.} Suppose a token has log-odds weight \(3.2\)
  for spam versus ham. Interpret this value as a ratio of
  class-conditional token probabilities.
\item
  \textbf{Bernoulli versus multinomial.} Give an example of a task where
  binary word-presence features may be better than count features.
\item
  \textbf{Negation.} Construct a small example showing why unigram Naive
  Bayes can mishandle negation.
\item
  \textbf{Implementation.} Modify the Python code above so that unknown
  test tokens are mapped to a special
  \texttt{\textless{}unk\textgreater{}} token.
\item
  \textbf{AI critique.} Ask an AI tool to explain Naive Bayes. Then
  identify one place where the answer is incomplete, imprecise, or
  missing a mathematical assumption.
\end{enumerate}

\section{Mini-project: a transparent text
classifier}\label{mini-project-a-transparent-text-classifier}

Build a Naive Bayes classifier for a small text classification task.
Possible datasets include synthetic course-question categories,
movie-review sentiment, or news-topic labels.

Your report should include:

\begin{enumerate}
\def\labelenumi{\arabic{enumi}.}
\tightlist
\item
  a description of the label set;
\item
  the vocabulary construction procedure;
\item
  the mathematical form of the model;
\item
  the smoothing value \(\alpha\) and why you chose it;
\item
  a confusion matrix;
\item
  precision, recall, and \(F_1\) score for at least one class;
\item
  the top positive and negative log-odds tokens;
\item
  an error analysis of at least five misclassified examples;
\item
  a discussion of one assumption that fails in your data.
\end{enumerate}

\section{Python lab connection}\label{python-lab-connection-2}

The associated lab for this chapter is planned as:

\begin{Shaded}
\begin{Highlighting}[]
\NormalTok{labs/chapter{-}06{-}naive{-}bayes{-}text{-}classification.ipynb}
\end{Highlighting}
\end{Shaded}

The lab should be designed for independent study and should include:

\begin{itemize}
\tightlist
\item
  a self-contained review of Bayes' theorem;
\item
  count-vector construction;
\item
  multinomial Naive Bayes from scratch;
\item
  smoothing experiments;
\item
  feature-weight interpretation;
\item
  comparison with logistic regression;
\item
  a short AI-assisted model-auditing exercise.
\end{itemize}

\section{What to remember}\label{what-to-remember}

\begin{itemize}
\tightlist
\item
  Naive Bayes is a generative classifier: it models \(P(Y)\) and
  \(P(X\mid Y)\).
\item
  The multinomial model treats a document as a bag of token counts.
\item
  The conditional independence assumption is mathematically convenient
  but linguistically false.
\item
  Add-\(\alpha\) smoothing prevents zero probabilities and regularizes
  log-odds weights.
\item
  The MAP rule becomes a linear classifier in count-vector space.
\item
  The feature weights are interpretable through log-ratios of
  class-conditional token probabilities.
\item
  Naive Bayes is a useful baseline before using larger neural NLP
  models.
\end{itemize}

\section{Notes and references}\label{notes-and-references-3}

Naive Bayes is a classical method in statistical NLP and information
retrieval. For broader context, see Manning and Schuetze (1999), Manning
et al. (2008), and Jurafsky and Martin (2024). The relationship between
generative and discriminative classifiers is also connected to standard
material in statistical learning theory.

\part{Part III. Linear Algebra of Language}

\chapter{Vector Space Models}\label{vector-space-models}

Document-term matrices, cosine retrieval, SVD, and latent semantic
structure

\hfill\break

In the previous chapters, documents were represented by count vectors,
TF-IDF vectors, and probabilistic models. In this chapter we study these
representations from a systematic linear-algebraic point of view. A
corpus becomes a matrix, documents become rows, vocabulary terms become
columns, similarity becomes an inner product, and latent semantic
structure is extracted by projection onto low-dimensional subspaces.

The central mathematical tool is the singular value decomposition. It
turns a large sparse document-term matrix into a sequence of best
low-rank approximations. This gives the classical method known as
\textbf{latent semantic analysis} and also prepares us for modern
embedding models.

\section{Learning goals}\label{learning-goals-6}

After studying this chapter, you should be able to:

\begin{enumerate}
\def\labelenumi{\arabic{enumi}.}
\tightlist
\item
  represent a corpus as a document-term matrix \(X\in R^{n\times d}\);
\item
  interpret rows, columns, row space, column space, rank, and Gram
  matrices in NLP terms;
\item
  compute document similarity using inner products and cosine
  similarity;
\item
  express retrieval scores in matrix form;
\item
  explain why normalization changes the geometry of document comparison;
\item
  state the singular value decomposition and connect it to eigenvalue
  problems;
\item
  use the Eckart--Young--Mirsky theorem to justify low-rank
  approximation;
\item
  construct latent semantic analysis from a truncated SVD;
\item
  compare raw vector-space retrieval with low-rank latent retrieval;
\item
  implement the main ideas in Python using NumPy;
\item
  use Plotly.js to visualize latent document geometry;
\item
  use AI tools to critique representation choices, identify hidden
  assumptions, and explain model behavior.
\end{enumerate}

Throughout this chapter, \(R\) denotes the real numbers. In code we use
ordinary NumPy arrays, but the mathematical object is a matrix
representation of a finite corpus.

\section{The main object: a document-term
matrix}\label{the-main-object-a-document-term-matrix}

Let

\[
D=\{d_1,d_2,\ldots,d_n\}
\]

be a corpus of \(n\) documents, and let

\[
V=\{v_1,v_2,\ldots,v_d\}
\]

be a fixed vocabulary of \(d\) tokens. A vector-space model chooses a
representation map

\[
\phi:D\longrightarrow R^d.
\]

The \(i\)th document is represented by

\[
x_i=\phi(d_i)=(x_{i1},x_{i2},\ldots,x_{id})\in R^d.
\]

The full corpus is represented by the matrix

\[
X=
\begin{bmatrix}
--- & x_1^T & ---\\
--- & x_2^T & ---\\
& \vdots & \\
--- & x_n^T & ---
\end{bmatrix}
\in R^{n\times d}.
\]

The entry \(X_{ij}\) is the weight of term \(v_j\) in document \(d_i\).
Depending on the representation, \(X_{ij}\) may be a count, a binary
indicator, a relative frequency, or a TF-IDF weight.

\textbf{Document-term matrix.} A corpus is converted into a matrix

\[
X\in R^{n\times d},
\]

where rows index documents and columns index vocabulary terms. Most
classical NLP methods can be described as operations on this matrix.

This viewpoint is powerful because it reduces many NLP operations to
familiar linear algebra:

{\def\LTcaptype{none} % do not increment counter
\begin{longtable}[]{@{}ll@{}}
\toprule\noalign{}
NLP operation & Linear algebra operation \\
\midrule\noalign{}
\endhead
\bottomrule\noalign{}
\endlastfoot
represent a document & choose a row vector \(x_i\in R^d\) \\
compare two documents & compute \(x_i^T x_j\) or cosine similarity \\
compare all documents & compute the Gram matrix \(XX^T\) \\
compare all terms & compute the Gram matrix \(X^TX\) \\
remove unimportant dimensions & project onto a subspace \\
find latent topics & compute a low-rank approximation \\
search with a query & compute matrix-vector scores \(Xq\) \\
\end{longtable}
}

\section{Rows, columns, and geometry}\label{rows-columns-and-geometry}

The row vector \(x_i\in R^d\) describes document \(d_i\) as a point in
vocabulary space. The \(j\)th column

\[
X_{:j}=\begin{bmatrix}X_{1j}\\X_{2j}\\\vdots\\X_{nj}\end{bmatrix}\in R^n
\]

describes how term \(v_j\) is distributed across documents.

These two perspectives are dual:

\begin{itemize}
\tightlist
\item
  rows describe \textbf{documents by terms};
\item
  columns describe \textbf{terms by documents}.
\end{itemize}

The document Gram matrix is

\[
G_D=XX^T\in R^{n\times n},
\]

with entries

\[
(G_D)_{ij}=x_i^Tx_j.
\]

The term Gram matrix is

\[
G_T=X^TX\in R^{d\times d},
\]

with entries

\[
(G_T)_{jk}=X_{:j}^T X_{:k}.
\]

Thus \(G_D\) measures document-document similarity, while \(G_T\)
measures term-term co-occurrence similarity.

\textbf{Two Gram matrices.}

\[
XX^T \quad \text{compares documents},
\]

while

\[
X^TX \quad \text{compares terms}.
\]

The singular value decomposition links these two matrices through the
same singular values.

\section{Inner product retrieval}\label{inner-product-retrieval}

Suppose a query is represented by a vector

\[
q\in R^d.
\]

A simple retrieval score for document \(d_i\) is

\[
s_i=x_i^Tq.
\]

For the whole corpus, the score vector is

\[
s=Xq\in R^n.
\]

Documents with larger scores are considered more relevant. This is the
foundation of classical vector-space search.

The inner product has a clear meaning for nonnegative text vectors:

\[
x_i^Tq=\sum_{j=1}^d x_{ij}q_j.
\]

Only terms that appear in both the document and the query contribute
positively. Large weights produce larger scores.

However, raw inner products favor long documents because long documents
tend to have larger norms. To reduce this length bias, we often use
cosine similarity.

\section{Cosine similarity and normalized
geometry}\label{cosine-similarity-and-normalized-geometry}

For nonzero vectors \(x,y\in R^d\), the cosine similarity is

\[
\operatorname{cos}(x,y)=\frac{x^Ty}{\lVert x\rVert_2\lVert y\rVert_2}.
\]

This is the inner product after projecting both vectors to the unit
sphere:

\[
\widehat x=\frac{x}{\lVert x\rVert_2},
\qquad
\widehat y=\frac{y}{\lVert y\rVert_2}.
\]

Then

\[
\operatorname{cos}(x,y)=\widehat x^T\widehat y.
\]

Cosine similarity compares direction rather than magnitude. In document
analysis, direction often represents relative term usage, while
magnitude is partly influenced by document length.

If \(\widetilde X\) is the row-normalized document matrix, then all
pairwise cosine similarities are contained in

\[
C=\widetilde X\widetilde X^T.
\]

Similarly, query scores are

\[
s=\widetilde X\widehat q.
\]

\textbf{Normalization changes the model.} Raw inner products ask: ``How
much weighted vocabulary overlap is there?'' Cosine similarity asks:
``How similar are the directions of the two term-weight vectors?'' These
are different statistical assumptions about what should count as
similarity.

\section{Python example: building a document-term
matrix}\label{python-example-building-a-document-term-matrix-1}

The following example uses a small corpus. The goal is not to build a
production search engine, but to expose every mathematical step.

\begin{Shaded}
\begin{Highlighting}[]
\ImportTok{import}\NormalTok{ re}
\ImportTok{import}\NormalTok{ numpy }\ImportTok{as}\NormalTok{ np}
\ImportTok{import}\NormalTok{ pandas }\ImportTok{as}\NormalTok{ pd}
\ImportTok{from}\NormalTok{ collections }\ImportTok{import}\NormalTok{ Counter}

\NormalTok{corpus }\OperatorTok{=}\NormalTok{ [}
    \StringTok{"linear algebra gives vector models for text data"}\NormalTok{,}
    \StringTok{"probability models describe uncertainty in language"}\NormalTok{,}
    \StringTok{"singular value decomposition finds latent structure"}\NormalTok{,}
    \StringTok{"vector spaces and matrices support document search"}\NormalTok{,}
    \StringTok{"language models use probability and optimization"}\NormalTok{,}
    \StringTok{"latent semantic analysis uses matrix factorization"}\NormalTok{,}
    \StringTok{"text retrieval compares query vectors with documents"}\NormalTok{,}
    \StringTok{"optimization trains modern neural language models"}
\NormalTok{]}

\NormalTok{labels }\OperatorTok{=}\NormalTok{ [}
    \StringTok{"linear algebra"}\NormalTok{,}
    \StringTok{"probability"}\NormalTok{,}
    \StringTok{"svd"}\NormalTok{,}
    \StringTok{"search"}\NormalTok{,}
    \StringTok{"probability"}\NormalTok{,}
    \StringTok{"svd"}\NormalTok{,}
    \StringTok{"search"}\NormalTok{,}
    \StringTok{"neural"}
\NormalTok{]}

\KeywordTok{def}\NormalTok{ tokenize(text):}
    \ControlFlowTok{return}\NormalTok{ re.findall(}\VerbatimStringTok{r"}\PreprocessorTok{[a{-}z]}\OperatorTok{+}\VerbatimStringTok{"}\NormalTok{, text.lower())}

\NormalTok{tokenized }\OperatorTok{=}\NormalTok{ [tokenize(doc) }\ControlFlowTok{for}\NormalTok{ doc }\KeywordTok{in}\NormalTok{ corpus]}
\NormalTok{vocab }\OperatorTok{=} \BuiltInTok{sorted}\NormalTok{(\{token }\ControlFlowTok{for}\NormalTok{ doc }\KeywordTok{in}\NormalTok{ tokenized }\ControlFlowTok{for}\NormalTok{ token }\KeywordTok{in}\NormalTok{ doc\})}
\NormalTok{token\_to\_col }\OperatorTok{=}\NormalTok{ \{token: j }\ControlFlowTok{for}\NormalTok{ j, token }\KeywordTok{in} \BuiltInTok{enumerate}\NormalTok{(vocab)\}}

\NormalTok{X\_counts }\OperatorTok{=}\NormalTok{ np.zeros((}\BuiltInTok{len}\NormalTok{(corpus), }\BuiltInTok{len}\NormalTok{(vocab)), dtype}\OperatorTok{=}\BuiltInTok{float}\NormalTok{)}
\ControlFlowTok{for}\NormalTok{ i, doc }\KeywordTok{in} \BuiltInTok{enumerate}\NormalTok{(tokenized):}
\NormalTok{    counts }\OperatorTok{=}\NormalTok{ Counter(doc)}
    \ControlFlowTok{for}\NormalTok{ token, count }\KeywordTok{in}\NormalTok{ counts.items():}
\NormalTok{        X\_counts[i, token\_to\_col[token]] }\OperatorTok{=}\NormalTok{ count}

\NormalTok{pd.DataFrame(X\_counts, columns}\OperatorTok{=}\NormalTok{vocab, index}\OperatorTok{=}\NormalTok{[}\SpecialStringTok{f"d}\SpecialCharTok{\{}\NormalTok{i}\OperatorTok{+}\DecValTok{1}\SpecialCharTok{\}}\SpecialStringTok{"} \ControlFlowTok{for}\NormalTok{ i }\KeywordTok{in} \BuiltInTok{range}\NormalTok{(}\BuiltInTok{len}\NormalTok{(corpus))])}
\end{Highlighting}
\end{Shaded}

{\def\LTcaptype{none} % do not increment counter
\begin{longtable}[]{@{}llllllllllllllllllllll@{}}
\toprule\noalign{}
& algebra & analysis & and & compares & data & decomposition & describe
& document & documents & factorization & ... & support & text & trains &
uncertainty & use & uses & value & vector & vectors & with \\
\midrule\noalign{}
\endhead
\bottomrule\noalign{}
\endlastfoot
d1 & 1.0 & 0.0 & 0.0 & 0.0 & 1.0 & 0.0 & 0.0 & 0.0 & 0.0 & 0.0 & ... &
0.0 & 1.0 & 0.0 & 0.0 & 0.0 & 0.0 & 0.0 & 1.0 & 0.0 & 0.0 \\
d2 & 0.0 & 0.0 & 0.0 & 0.0 & 0.0 & 0.0 & 1.0 & 0.0 & 0.0 & 0.0 & ... &
0.0 & 0.0 & 0.0 & 1.0 & 0.0 & 0.0 & 0.0 & 0.0 & 0.0 & 0.0 \\
d3 & 0.0 & 0.0 & 0.0 & 0.0 & 0.0 & 1.0 & 0.0 & 0.0 & 0.0 & 0.0 & ... &
0.0 & 0.0 & 0.0 & 0.0 & 0.0 & 0.0 & 1.0 & 0.0 & 0.0 & 0.0 \\
d4 & 0.0 & 0.0 & 1.0 & 0.0 & 0.0 & 0.0 & 0.0 & 1.0 & 0.0 & 0.0 & ... &
1.0 & 0.0 & 0.0 & 0.0 & 0.0 & 0.0 & 0.0 & 1.0 & 0.0 & 0.0 \\
d5 & 0.0 & 0.0 & 1.0 & 0.0 & 0.0 & 0.0 & 0.0 & 0.0 & 0.0 & 0.0 & ... &
0.0 & 0.0 & 0.0 & 0.0 & 1.0 & 0.0 & 0.0 & 0.0 & 0.0 & 0.0 \\
d6 & 0.0 & 1.0 & 0.0 & 0.0 & 0.0 & 0.0 & 0.0 & 0.0 & 0.0 & 1.0 & ... &
0.0 & 0.0 & 0.0 & 0.0 & 0.0 & 1.0 & 0.0 & 0.0 & 0.0 & 0.0 \\
d7 & 0.0 & 0.0 & 0.0 & 1.0 & 0.0 & 0.0 & 0.0 & 0.0 & 1.0 & 0.0 & ... &
0.0 & 1.0 & 0.0 & 0.0 & 0.0 & 0.0 & 0.0 & 0.0 & 1.0 & 1.0 \\
d8 & 0.0 & 0.0 & 0.0 & 0.0 & 0.0 & 0.0 & 0.0 & 0.0 & 0.0 & 0.0 & ... &
0.0 & 0.0 & 1.0 & 0.0 & 0.0 & 0.0 & 0.0 & 0.0 & 0.0 & 0.0 \\
\end{longtable}
}

The matrix is sparse: most entries are zero. Real NLP matrices may have
millions of documents and hundreds of thousands of vocabulary terms, so
sparse storage is essential in practice. Still, the dense representation
is useful pedagogically.

\section{Python example: cosine search from
scratch}\label{python-example-cosine-search-from-scratch}

Now we represent a query and rank the documents by cosine similarity.

\begin{Shaded}
\begin{Highlighting}[]
\KeywordTok{def}\NormalTok{ vectorize\_query(query, vocab\_index):}
\NormalTok{    q }\OperatorTok{=}\NormalTok{ np.zeros(}\BuiltInTok{len}\NormalTok{(vocab\_index), dtype}\OperatorTok{=}\BuiltInTok{float}\NormalTok{)}
    \ControlFlowTok{for}\NormalTok{ token }\KeywordTok{in}\NormalTok{ tokenize(query):}
        \ControlFlowTok{if}\NormalTok{ token }\KeywordTok{in}\NormalTok{ vocab\_index:}
\NormalTok{            q[vocab\_index[token]] }\OperatorTok{+=} \FloatTok{1.0}
    \ControlFlowTok{return}\NormalTok{ q}

\KeywordTok{def}\NormalTok{ row\_normalize(A, eps}\OperatorTok{=}\FloatTok{1e{-}12}\NormalTok{):}
\NormalTok{    norms }\OperatorTok{=}\NormalTok{ np.linalg.norm(A, axis}\OperatorTok{=}\DecValTok{1}\NormalTok{, keepdims}\OperatorTok{=}\VariableTok{True}\NormalTok{)}
    \ControlFlowTok{return}\NormalTok{ A }\OperatorTok{/}\NormalTok{ np.maximum(norms, eps)}

\NormalTok{query }\OperatorTok{=} \StringTok{"matrix decomposition for latent semantic search"}
\NormalTok{q\_counts }\OperatorTok{=}\NormalTok{ vectorize\_query(query, token\_to\_col)}

\NormalTok{Xn }\OperatorTok{=}\NormalTok{ row\_normalize(X\_counts)}
\NormalTok{qn }\OperatorTok{=}\NormalTok{ q\_counts }\OperatorTok{/} \BuiltInTok{max}\NormalTok{(np.linalg.norm(q\_counts), }\FloatTok{1e{-}12}\NormalTok{)}
\NormalTok{cos\_scores }\OperatorTok{=}\NormalTok{ Xn }\OperatorTok{@}\NormalTok{ qn}

\NormalTok{ranking }\OperatorTok{=}\NormalTok{ np.argsort(}\OperatorTok{{-}}\NormalTok{cos\_scores)}
\ControlFlowTok{for}\NormalTok{ rank, i }\KeywordTok{in} \BuiltInTok{enumerate}\NormalTok{(ranking, start}\OperatorTok{=}\DecValTok{1}\NormalTok{):}
    \BuiltInTok{print}\NormalTok{(}\SpecialStringTok{f"}\SpecialCharTok{\{}\NormalTok{rank}\SpecialCharTok{\}}\SpecialStringTok{. score=}\SpecialCharTok{\{}\NormalTok{cos\_scores[i]}\SpecialCharTok{:.3f\}}\SpecialStringTok{ | label=}\SpecialCharTok{\{}\NormalTok{labels[i]}\SpecialCharTok{:15s\}}\SpecialStringTok{ | }\SpecialCharTok{\{}\NormalTok{corpus[i]}\SpecialCharTok{\}}\SpecialStringTok{"}\NormalTok{)}
\end{Highlighting}
\end{Shaded}

\begin{verbatim}
1. score=0.500 | label=svd             | latent semantic analysis uses matrix factorization
2. score=0.333 | label=svd             | singular value decomposition finds latent structure
3. score=0.154 | label=search          | vector spaces and matrices support document search
4. score=0.144 | label=linear algebra  | linear algebra gives vector models for text data
5. score=0.000 | label=probability     | probability models describe uncertainty in language
6. score=0.000 | label=probability     | language models use probability and optimization
7. score=0.000 | label=search          | text retrieval compares query vectors with documents
8. score=0.000 | label=neural          | optimization trains modern neural language models
\end{verbatim}

The formula used here is exactly

\[
s=\widetilde X\widehat q.
\]

This is a complete vector-space retrieval system in a few lines of
NumPy.

\section{Term weighting as diagonal
scaling}\label{term-weighting-as-diagonal-scaling}

Raw count matrices are often not the best representation. Common words
may dominate similarity. TF-IDF corrects this by increasing the weight
of terms that are frequent in a document but rare across the corpus.

Let

\[
\operatorname{tf}_{ij}=X_{ij}
\]

or a transformed frequency such as \(\log(1+X_{ij})\). The document
frequency of term \(j\) is

\[
\operatorname{df}_j=\#\{i:X_{ij}>0\}.
\]

A smoothed inverse document frequency is

\[
\operatorname{idf}_j=\log\left(\frac{1+n}{1+\operatorname{df}_j}\right)+1.
\]

Then

\[
X^{\operatorname{tfidf}}_{ij}=\operatorname{tf}_{ij}\operatorname{idf}_j.
\]

In matrix form, this is column scaling:

\[
X^{\operatorname{tfidf}}=XW,
\]

where

\[
W=\operatorname{diag}(\operatorname{idf}_1,\ldots,\operatorname{idf}_d).
\]

Thus TF-IDF is not just a heuristic. It is a diagonal linear
transformation of the feature space.

\begin{Shaded}
\begin{Highlighting}[]
\KeywordTok{def}\NormalTok{ tfidf\_transform(X):}
\NormalTok{    n }\OperatorTok{=}\NormalTok{ X.shape[}\DecValTok{0}\NormalTok{]}
\NormalTok{    df }\OperatorTok{=}\NormalTok{ np.}\BuiltInTok{sum}\NormalTok{(X }\OperatorTok{\textgreater{}} \DecValTok{0}\NormalTok{, axis}\OperatorTok{=}\DecValTok{0}\NormalTok{)}
\NormalTok{    idf }\OperatorTok{=}\NormalTok{ np.log((}\DecValTok{1} \OperatorTok{+}\NormalTok{ n) }\OperatorTok{/}\NormalTok{ (}\DecValTok{1} \OperatorTok{+}\NormalTok{ df)) }\OperatorTok{+} \DecValTok{1}
    \ControlFlowTok{return}\NormalTok{ X }\OperatorTok{*}\NormalTok{ idf, idf}

\NormalTok{X\_tfidf, idf }\OperatorTok{=}\NormalTok{ tfidf\_transform(X\_counts)}

\NormalTok{idf\_table }\OperatorTok{=}\NormalTok{ pd.DataFrame(\{}\StringTok{"term"}\NormalTok{: vocab, }\StringTok{"idf"}\NormalTok{: idf\}).sort\_values(}\StringTok{"idf"}\NormalTok{, ascending}\OperatorTok{=}\VariableTok{False}\NormalTok{)}
\NormalTok{idf\_table.head(}\DecValTok{12}\NormalTok{)}
\end{Highlighting}
\end{Shaded}

{\def\LTcaptype{none} % do not increment counter
\begin{longtable}[]{@{}lll@{}}
\toprule\noalign{}
& term & idf \\
\midrule\noalign{}
\endhead
\bottomrule\noalign{}
\endlastfoot
0 & algebra & 2.504077 \\
18 & matrix & 2.504077 \\
21 & neural & 2.504077 \\
24 & query & 2.504077 \\
25 & retrieval & 2.504077 \\
26 & search & 2.504077 \\
27 & semantic & 2.504077 \\
28 & singular & 2.504077 \\
29 & spaces & 2.504077 \\
30 & structure & 2.504077 \\
31 & support & 2.504077 \\
33 & trains & 2.504077 \\
\end{longtable}
}

Now compare retrieval with TF-IDF weights.

\begin{Shaded}
\begin{Highlighting}[]
\NormalTok{q\_tfidf }\OperatorTok{=}\NormalTok{ q\_counts }\OperatorTok{*}\NormalTok{ idf}
\NormalTok{Xtn }\OperatorTok{=}\NormalTok{ row\_normalize(X\_tfidf)}
\NormalTok{qtn }\OperatorTok{=}\NormalTok{ q\_tfidf }\OperatorTok{/} \BuiltInTok{max}\NormalTok{(np.linalg.norm(q\_tfidf), }\FloatTok{1e{-}12}\NormalTok{)}
\NormalTok{tfidf\_scores }\OperatorTok{=}\NormalTok{ Xtn }\OperatorTok{@}\NormalTok{ qtn}

\NormalTok{ranking }\OperatorTok{=}\NormalTok{ np.argsort(}\OperatorTok{{-}}\NormalTok{tfidf\_scores)}
\ControlFlowTok{for}\NormalTok{ rank, i }\KeywordTok{in} \BuiltInTok{enumerate}\NormalTok{(ranking, start}\OperatorTok{=}\DecValTok{1}\NormalTok{):}
    \BuiltInTok{print}\NormalTok{(}\SpecialStringTok{f"}\SpecialCharTok{\{}\NormalTok{rank}\SpecialCharTok{\}}\SpecialStringTok{. score=}\SpecialCharTok{\{}\NormalTok{tfidf\_scores[i]}\SpecialCharTok{:.3f\}}\SpecialStringTok{ | label=}\SpecialCharTok{\{}\NormalTok{labels[i]}\SpecialCharTok{:15s\}}\SpecialStringTok{ | }\SpecialCharTok{\{}\NormalTok{corpus[i]}\SpecialCharTok{\}}\SpecialStringTok{"}\NormalTok{)}
\end{Highlighting}
\end{Shaded}

\begin{verbatim}
1. score=0.474 | label=svd             | latent semantic analysis uses matrix factorization
2. score=0.299 | label=svd             | singular value decomposition finds latent structure
3. score=0.165 | label=search          | vector spaces and matrices support document search
4. score=0.161 | label=linear algebra  | linear algebra gives vector models for text data
5. score=0.000 | label=probability     | probability models describe uncertainty in language
6. score=0.000 | label=probability     | language models use probability and optimization
7. score=0.000 | label=search          | text retrieval compares query vectors with documents
8. score=0.000 | label=neural          | optimization trains modern neural language models
\end{verbatim}

The only change is the geometry of the feature space. The similarity
calculation is still cosine similarity, but the axes have been rescaled.

\section{Rank and latent structure}\label{rank-and-latent-structure}

The rank of \(X\) is

\[
\operatorname{rank}(X)=\dim(\operatorname{Row}(X))=\dim(\operatorname{Col}(X)).
\]

In NLP, a high-dimensional vocabulary space often contains redundancy:

\begin{itemize}
\tightlist
\item
  different words may carry related meanings;
\item
  topics create correlations among terms;
\item
  document collections often lie near lower-dimensional structures;
\item
  noise and rare words inflate the apparent dimension.
\end{itemize}

A low-rank approximation tries to replace \(X\) by a simpler matrix

\[
X_k\in R^{n\times d},
\qquad
\operatorname{rank}(X_k)\le k,
\]

where \(k\ll \min(n,d)\).

The goal is to keep the most important geometric structure while
discarding high-dimensional noise.

\section{Singular value
decomposition}\label{singular-value-decomposition-1}

Every matrix \(X\in R^{n\times d}\) has a singular value decomposition

\[
X=U\Sigma V^T,
\]

where:

\begin{itemize}
\tightlist
\item
  \(U\in R^{n\times r}\) has orthonormal columns;
\item
  \(V\in R^{d\times r}\) has orthonormal columns;
\item
  \(\Sigma=\operatorname{diag}(\sigma_1,\ldots,\sigma_r)\);
\item
  \(r=\operatorname{rank}(X)\);
\item
  \(\sigma_1\ge \sigma_2\ge \cdots\ge \sigma_r>0\).
\end{itemize}

The columns of \(U\) are left singular vectors. The columns of \(V\) are
right singular vectors. The singular values measure how much variation
lies in each singular direction.

The SVD diagonalizes the two Gram matrices:

\[
XX^T=U\Sigma^2U^T,
\]

and

\[
X^TX=V\Sigma^2V^T.
\]

Thus left singular vectors describe principal directions among
documents, while right singular vectors describe principal directions
among terms.

\section{Best low-rank approximation}\label{best-low-rank-approximation}

For \(k<r\), define

\[
U_k=[u_1,\ldots,u_k],
\qquad
V_k=[v_1,\ldots,v_k],
\qquad
\Sigma_k=\operatorname{diag}(\sigma_1,\ldots,\sigma_k).
\]

The rank-\(k\) truncated SVD approximation is

\[
X_k=U_k\Sigma_kV_k^T.
\]

The Eckart--Young--Mirsky theorem says that \(X_k\) is the best
rank-\(k\) approximation to \(X\) in Frobenius norm:

\[
X_k=\arg\min_{\operatorname{rank}(A)\le k}\lVert X-A\rVert_F.
\]

Moreover,

\[
\lVert X-X_k\rVert_F^2=\sum_{j=k+1}^r \sigma_j^2.
\]

This theorem is the mathematical foundation of latent semantic analysis.

\textbf{SVD interpretation for NLP.} The truncated SVD replaces sparse
surface-level term matching by a lower-dimensional coordinate system.
Documents that do not share many exact words may become close if they
have similar coordinates in the latent space.

\section{Latent semantic analysis}\label{latent-semantic-analysis}

Latent semantic analysis, or LSA, applies truncated SVD to a
term-weighted document matrix such as \(X^{\operatorname{tfidf}}\).

Let

\[
X\approx X_k=U_k\Sigma_kV_k^T.
\]

The row vector

\[
z_i=(U_k\Sigma_k)_{i:}\in R^k
\]

is a \(k\)-dimensional latent representation of document \(d_i\). The
row vector

\[
t_j=(V_k\Sigma_k)_{j:}\in R^k
\]

may be interpreted as a latent representation of term \(v_j\).

A query vector \(q\in R^d\) can be projected into the same latent
coordinate system by

\[
z_q=q^TV_k\in R^k.
\]

Then retrieval can be performed by comparing \(z_q\) with the document
coordinates \(z_i\).

There are several conventions for scaling latent coordinates. The
important point is that documents and queries are compared after
projection onto the subspace spanned by the top right singular vectors.

\section{Python example: SVD and low-rank
coordinates}\label{python-example-svd-and-low-rank-coordinates}

\begin{Shaded}
\begin{Highlighting}[]
\CommentTok{\# Work with the TF{-}IDF matrix for SVD.}
\NormalTok{X }\OperatorTok{=}\NormalTok{ X\_tfidf.copy()}

\NormalTok{U, singular\_values, Vt }\OperatorTok{=}\NormalTok{ np.linalg.svd(X, full\_matrices}\OperatorTok{=}\VariableTok{False}\NormalTok{)}

\NormalTok{sv\_table }\OperatorTok{=}\NormalTok{ pd.DataFrame(\{}
    \StringTok{"component"}\NormalTok{: np.arange(}\DecValTok{1}\NormalTok{, }\BuiltInTok{len}\NormalTok{(singular\_values) }\OperatorTok{+} \DecValTok{1}\NormalTok{),}
    \StringTok{"singular\_value"}\NormalTok{: singular\_values,}
    \StringTok{"energy\_fraction"}\NormalTok{: singular\_values}\OperatorTok{**}\DecValTok{2} \OperatorTok{/}\NormalTok{ np.}\BuiltInTok{sum}\NormalTok{(singular\_values}\OperatorTok{**}\DecValTok{2}\NormalTok{),}
    \StringTok{"cumulative\_energy"}\NormalTok{: np.cumsum(singular\_values}\OperatorTok{**}\DecValTok{2}\NormalTok{) }\OperatorTok{/}\NormalTok{ np.}\BuiltInTok{sum}\NormalTok{(singular\_values}\OperatorTok{**}\DecValTok{2}\NormalTok{)}
\NormalTok{\})}
\NormalTok{sv\_table}
\end{Highlighting}
\end{Shaded}

{\def\LTcaptype{none} % do not increment counter
\begin{longtable}[]{@{}lllll@{}}
\toprule\noalign{}
& component & singular\_value & energy\_fraction & cumulative\_energy \\
\midrule\noalign{}
\endhead
\bottomrule\noalign{}
\endlastfoot
0 & 1 & 7.180600 & 0.184347 & 0.184347 \\
1 & 2 & 6.606961 & 0.156069 & 0.340416 \\
2 & 3 & 6.337221 & 0.143586 & 0.484001 \\
3 & 4 & 6.319106 & 0.142766 & 0.626768 \\
4 & 5 & 5.956870 & 0.126867 & 0.753635 \\
5 & 6 & 5.599287 & 0.112093 & 0.865728 \\
6 & 7 & 4.818235 & 0.083002 & 0.948730 \\
7 & 8 & 3.786813 & 0.051270 & 1.000000 \\
\end{longtable}
}

The squared singular values describe how much Frobenius-norm energy is
captured by each component.

\begin{Shaded}
\begin{Highlighting}[]
\KeywordTok{def}\NormalTok{ truncated\_svd\_coordinates(X, k):}
\NormalTok{    U, s, Vt }\OperatorTok{=}\NormalTok{ np.linalg.svd(X, full\_matrices}\OperatorTok{=}\VariableTok{False}\NormalTok{)}
\NormalTok{    Uk }\OperatorTok{=}\NormalTok{ U[:, :k]}
\NormalTok{    sk }\OperatorTok{=}\NormalTok{ s[:k]}
\NormalTok{    Vtk }\OperatorTok{=}\NormalTok{ Vt[:k, :]}
\NormalTok{    Z\_docs }\OperatorTok{=}\NormalTok{ Uk }\OperatorTok{*}\NormalTok{ sk}
    \ControlFlowTok{return}\NormalTok{ Z\_docs, Vtk.T, sk}

\NormalTok{Z2, V2, s2 }\OperatorTok{=}\NormalTok{ truncated\_svd\_coordinates(X, k}\OperatorTok{=}\DecValTok{2}\NormalTok{)}

\NormalTok{coords }\OperatorTok{=}\NormalTok{ pd.DataFrame(\{}
    \StringTok{"doc"}\NormalTok{: [}\SpecialStringTok{f"d}\SpecialCharTok{\{}\NormalTok{i}\OperatorTok{+}\DecValTok{1}\SpecialCharTok{\}}\SpecialStringTok{"} \ControlFlowTok{for}\NormalTok{ i }\KeywordTok{in} \BuiltInTok{range}\NormalTok{(}\BuiltInTok{len}\NormalTok{(corpus))],}
    \StringTok{"x"}\NormalTok{: Z2[:, }\DecValTok{0}\NormalTok{],}
    \StringTok{"y"}\NormalTok{: Z2[:, }\DecValTok{1}\NormalTok{],}
    \StringTok{"label"}\NormalTok{: labels,}
    \StringTok{"text"}\NormalTok{: corpus}
\NormalTok{\})}
\NormalTok{coords}
\end{Highlighting}
\end{Shaded}

{\def\LTcaptype{none} % do not increment counter
\begin{longtable}[]{@{}llllll@{}}
\toprule\noalign{}
& doc & x & y & label & text \\
\midrule\noalign{}
\endhead
\bottomrule\noalign{}
\endlastfoot
0 & d1 & -4.549237e+00 & 1.828238e+00 & linear algebra & linear algebra
gives vector models for text data \\
1 & d2 & -2.351264e+00 & -2.206304e+00 & probability & probability
models describe uncertainty in lan... \\
2 & d3 & -3.986034e-16 & -5.134640e-15 & svd & singular value
decomposition finds latent stru... \\
3 & d4 & -2.813666e+00 & -6.731915e-01 & search & vector spaces and
matrices support document se... \\
4 & d5 & -2.734224e+00 & -2.361934e+00 & probability & language models
use probability and optimization \\
5 & d6 & 7.041865e-18 & -2.499359e-17 & svd & latent semantic analysis
uses matrix factoriza... \\
6 & d7 & -2.101400e+00 & 4.953988e+00 & search & text retrieval compares
query vectors with doc... \\
7 & d8 & -2.351264e+00 & -2.206304e+00 & neural & optimization trains
modern neural language models \\
\end{longtable}
}

\section{Reconstruction error}\label{reconstruction-error}

A rank-\(k\) approximation is useful only if it captures enough
structure. The following code computes reconstruction error for each
possible \(k\).

\begin{Shaded}
\begin{Highlighting}[]
\KeywordTok{def}\NormalTok{ svd\_reconstruction(X, k):}
\NormalTok{    U, s, Vt }\OperatorTok{=}\NormalTok{ np.linalg.svd(X, full\_matrices}\OperatorTok{=}\VariableTok{False}\NormalTok{)}
    \ControlFlowTok{return}\NormalTok{ (U[:, :k] }\OperatorTok{*}\NormalTok{ s[:k]) }\OperatorTok{@}\NormalTok{ Vt[:k, :]}

\NormalTok{rows }\OperatorTok{=}\NormalTok{ []}
\ControlFlowTok{for}\NormalTok{ k }\KeywordTok{in} \BuiltInTok{range}\NormalTok{(}\DecValTok{1}\NormalTok{, }\BuiltInTok{min}\NormalTok{(X.shape) }\OperatorTok{+} \DecValTok{1}\NormalTok{):}
\NormalTok{    Xk }\OperatorTok{=}\NormalTok{ svd\_reconstruction(X, k)}
\NormalTok{    error }\OperatorTok{=}\NormalTok{ np.linalg.norm(X }\OperatorTok{{-}}\NormalTok{ Xk, }\BuiltInTok{ord}\OperatorTok{=}\StringTok{"fro"}\NormalTok{)}
\NormalTok{    relative\_error }\OperatorTok{=}\NormalTok{ error }\OperatorTok{/}\NormalTok{ np.linalg.norm(X, }\BuiltInTok{ord}\OperatorTok{=}\StringTok{"fro"}\NormalTok{)}
\NormalTok{    rows.append((k, error, relative\_error))}

\NormalTok{pd.DataFrame(rows, columns}\OperatorTok{=}\NormalTok{[}\StringTok{"k"}\NormalTok{, }\StringTok{"frobenius\_error"}\NormalTok{, }\StringTok{"relative\_error"}\NormalTok{])}
\end{Highlighting}
\end{Shaded}

{\def\LTcaptype{none} % do not increment counter
\begin{longtable}[]{@{}llll@{}}
\toprule\noalign{}
& k & frobenius\_error & relative\_error \\
\midrule\noalign{}
\endhead
\bottomrule\noalign{}
\endlastfoot
0 & 1 & 1.510414e+01 & 9.031353e-01 \\
1 & 2 & 1.358246e+01 & 8.121479e-01 \\
2 & 3 & 1.201344e+01 & 7.183304e-01 \\
3 & 4 & 1.021722e+01 & 6.109275e-01 \\
4 & 5 & 8.301045e+00 & 4.963518e-01 \\
5 & 6 & 6.128241e+00 & 3.664314e-01 \\
6 & 7 & 3.786813e+00 & 2.264283e-01 \\
7 & 8 & 5.821355e-14 & 3.480815e-15 \\
\end{longtable}
}

As \(k\) increases, the approximation error decreases. The mathematical
reason is that more singular directions are retained.

\section{Plotly visualization from
Python}\label{plotly-visualization-from-python}

Quarto can render interactive Plotly figures from Python. The next code
chunk produces a two-dimensional latent semantic map.

\begin{Shaded}
\begin{Highlighting}[]
\ControlFlowTok{try}\NormalTok{:}
    \ImportTok{import}\NormalTok{ plotly.express }\ImportTok{as}\NormalTok{ px}

\NormalTok{    fig }\OperatorTok{=}\NormalTok{ px.scatter(}
\NormalTok{        coords,}
\NormalTok{        x}\OperatorTok{=}\StringTok{"x"}\NormalTok{,}
\NormalTok{        y}\OperatorTok{=}\StringTok{"y"}\NormalTok{,}
\NormalTok{        color}\OperatorTok{=}\StringTok{"label"}\NormalTok{,}
\NormalTok{        hover\_data}\OperatorTok{=}\NormalTok{[}\StringTok{"doc"}\NormalTok{, }\StringTok{"text"}\NormalTok{],}
\NormalTok{        title}\OperatorTok{=}\StringTok{"Latent semantic map from rank{-}2 SVD"}
\NormalTok{    )}
\NormalTok{    fig.update\_layout(xaxis\_title}\OperatorTok{=}\StringTok{"latent dimension 1"}\NormalTok{, yaxis\_title}\OperatorTok{=}\StringTok{"latent dimension 2"}\NormalTok{)}
\NormalTok{    fig.show()}
\ControlFlowTok{except} \PreprocessorTok{Exception} \ImportTok{as}\NormalTok{ e:}
    \BuiltInTok{print}\NormalTok{(}\StringTok{"Plotly figure skipped:"}\NormalTok{, e)}
\end{Highlighting}
\end{Shaded}

\begin{verbatim}
Unable to display output for mime type(s): text/html
\end{verbatim}

Two-dimensional latent semantic representation from truncated SVD.

\begin{verbatim}
Unable to display output for mime type(s): text/html
\end{verbatim}

If Plotly is not available during rendering, the rest of the chapter
still works. The embedded JavaScript interactive below does not require
Python Plotly.

\section{Interactive HTML: latent semantic
projection}\label{interactive-html-latent-semantic-projection}

The interactive below uses a small fixed document-term matrix. Move the
slider to choose the rank \(k\) used by a precomputed low-rank
representation. The plot shows how documents move as the latent
dimension changes. The displayed quantity is the cumulative
singular-value energy

\[
\frac{\sum_{j=1}^k\sigma_j^2}{\sum_{j=1}^r\sigma_j^2}.
\]

This example is intentionally small. In real corpora, the visual
separation of topics may be much more subtle, and the first two
dimensions may not capture everything important.

\section{Low-rank retrieval}\label{low-rank-retrieval}

In raw vector-space retrieval, the query score is computed in the
original vocabulary space:

\[
s=Xq.
\]

In LSA retrieval, we first project the query into the latent
right-singular-vector space:

\[
z_q=q^TV_k.
\]

The document coordinates are

\[
Z=XV_k=U_k\Sigma_k.
\]

Then the latent retrieval score is

\[
s_i=z_i^Tz_q.
\]

In matrix form,

\[
s=Zz_q^T.
\]

If cosine similarity is desired, normalize the latent document
coordinates and the latent query coordinate before scoring.

\begin{Shaded}
\begin{Highlighting}[]
\KeywordTok{def}\NormalTok{ lsa\_search(X, q, k):}
\NormalTok{    U, s, Vt }\OperatorTok{=}\NormalTok{ np.linalg.svd(X, full\_matrices}\OperatorTok{=}\VariableTok{False}\NormalTok{)}
\NormalTok{    Vk }\OperatorTok{=}\NormalTok{ Vt[:k, :].T}
\NormalTok{    Z\_docs }\OperatorTok{=}\NormalTok{ X }\OperatorTok{@}\NormalTok{ Vk}
\NormalTok{    z\_query }\OperatorTok{=}\NormalTok{ q }\OperatorTok{@}\NormalTok{ Vk}
\NormalTok{    Zn }\OperatorTok{=}\NormalTok{ row\_normalize(Z\_docs)}
\NormalTok{    zqn }\OperatorTok{=}\NormalTok{ z\_query }\OperatorTok{/} \BuiltInTok{max}\NormalTok{(np.linalg.norm(z\_query), }\FloatTok{1e{-}12}\NormalTok{)}
    \ControlFlowTok{return}\NormalTok{ Zn }\OperatorTok{@}\NormalTok{ zqn}

\ControlFlowTok{for}\NormalTok{ k }\KeywordTok{in}\NormalTok{ [}\DecValTok{1}\NormalTok{, }\DecValTok{2}\NormalTok{, }\DecValTok{3}\NormalTok{, }\DecValTok{4}\NormalTok{]:}
\NormalTok{    scores }\OperatorTok{=}\NormalTok{ lsa\_search(X\_tfidf, q\_tfidf, k}\OperatorTok{=}\NormalTok{k)}
    \BuiltInTok{print}\NormalTok{(}\SpecialStringTok{f"}\CharTok{\textbackslash{}n}\SpecialStringTok{Rank k=}\SpecialCharTok{\{}\NormalTok{k}\SpecialCharTok{\}}\SpecialStringTok{"}\NormalTok{)}
    \ControlFlowTok{for}\NormalTok{ rank, i }\KeywordTok{in} \BuiltInTok{enumerate}\NormalTok{(np.argsort(}\OperatorTok{{-}}\NormalTok{scores), start}\OperatorTok{=}\DecValTok{1}\NormalTok{):}
        \BuiltInTok{print}\NormalTok{(}\SpecialStringTok{f"}\SpecialCharTok{\{}\NormalTok{rank}\SpecialCharTok{\}}\SpecialStringTok{. score=}\SpecialCharTok{\{}\NormalTok{scores[i]}\SpecialCharTok{: .3f\}}\SpecialStringTok{ | }\SpecialCharTok{\{}\NormalTok{corpus[i]}\SpecialCharTok{\}}\SpecialStringTok{"}\NormalTok{)}
\end{Highlighting}
\end{Shaded}

\begin{verbatim}

Rank k=1
1. score= 1.000 | linear algebra gives vector models for text data
2. score= 1.000 | probability models describe uncertainty in language
3. score= 1.000 | vector spaces and matrices support document search
4. score= 1.000 | language models use probability and optimization
5. score= 1.000 | text retrieval compares query vectors with documents
6. score= 1.000 | optimization trains modern neural language models
7. score= 0.001 | singular value decomposition finds latent structure
8. score= 0.000 | latent semantic analysis uses matrix factorization

Rank k=2
1. score= 0.980 | linear algebra gives vector models for text data
2. score= 0.914 | vector spaces and matrices support document search
3. score= 0.625 | language models use probability and optimization
4. score= 0.592 | probability models describe uncertainty in language
5. score= 0.592 | optimization trains modern neural language models
6. score= 0.552 | text retrieval compares query vectors with documents
7. score= 0.000 | latent semantic analysis uses matrix factorization
8. score=-0.000 | singular value decomposition finds latent structure

Rank k=3
1. score= 0.959 | singular value decomposition finds latent structure
2. score= 0.959 | latent semantic analysis uses matrix factorization
3. score= 0.278 | linear algebra gives vector models for text data
4. score= 0.259 | vector spaces and matrices support document search
5. score= 0.177 | language models use probability and optimization
6. score= 0.168 | optimization trains modern neural language models
7. score= 0.168 | probability models describe uncertainty in language
8. score= 0.156 | text retrieval compares query vectors with documents

Rank k=4
1. score= 0.922 | singular value decomposition finds latent structure
2. score= 0.922 | latent semantic analysis uses matrix factorization
3. score= 0.364 | vector spaces and matrices support document search
4. score= 0.320 | linear algebra gives vector models for text data
5. score= 0.065 | language models use probability and optimization
6. score= 0.031 | text retrieval compares query vectors with documents
7. score=-0.016 | optimization trains modern neural language models
8. score=-0.016 | probability models describe uncertainty in language
\end{verbatim}

Low-rank retrieval may retrieve documents that do not share many exact
query words but lie in a related latent direction. This is the main
advantage of LSA over purely surface-level matching.

\section{How LSA handles synonymy and
polysemy}\label{how-lsa-handles-synonymy-and-polysemy}

Two common language difficulties are \textbf{synonymy} and
\textbf{polysemy}.

Synonymy means that different words have related meanings. For example,
``car'' and ``automobile'' may not match exactly in a count-vector
model. If they occur in similar document contexts, a low-rank model may
place them near each other.

Polysemy means that one word has multiple meanings. For example,
``bank'' may refer to a financial institution or a river bank. A simple
vector-space model assigns one column to the token ``bank,'' so it
cannot fully separate meanings. LSA may partially separate contexts
through document-level directions, but it still lacks the contextual
flexibility of modern neural embeddings.

\textbf{What LSA can and cannot do.} LSA can discover linear correlation
structure in a corpus. It cannot truly understand syntax, word order, or
context-dependent meaning. Its strength is low-rank linear geometry; its
limitation is that language is not purely linear bag-of-words geometry.

\section{Centering, scaling, and preprocessing as linear
transformations}\label{centering-scaling-and-preprocessing-as-linear-transformations}

Many preprocessing steps can be described as transformations of \(X\).

\subsection{Row normalization}\label{row-normalization}

Row normalization replaces each document vector by

\[
\widehat x_i=\frac{x_i}{\lVert x_i\rVert_2}.
\]

This maps documents to the unit sphere in \(R^d\).

\subsection{Column scaling}\label{column-scaling}

TF-IDF is column scaling:

\[
X\mapsto XW.
\]

This changes the metric by giving different importance to different
vocabulary dimensions.

\subsection{Stop-word removal}\label{stop-word-removal}

Removing a set of stop-word columns is projection onto a coordinate
subspace. If \(S\) is the set of retained columns, then

\[
X\mapsto XP_S,
\]

where \(P_S\) selects the retained vocabulary coordinates.

\subsection{Centering}\label{centering}

Centering subtracts the column mean:

\[
X_c=X-\mathbf{1}\mu^T,
\]

where

\[
\mu=\frac{1}{n}X^T\mathbf{1}.
\]

Centering is standard in PCA, but it must be used carefully in sparse
nonnegative text matrices because it destroys sparsity and creates
negative entries.

\section{PCA and LSA: same theorem, different data
convention}\label{pca-and-lsa-same-theorem-different-data-convention}

PCA is usually applied to centered data. LSA is usually applied to
uncentered TF-IDF data. Both use the SVD.

For centered data \(X_c\), PCA finds directions maximizing variance:

\[
v_1=\arg\max_{\lVert v\rVert=1}\lVert X_cv\rVert_2^2.
\]

The solution is the top right singular vector of \(X_c\).

LSA instead applies SVD to a sparse term-weight matrix, often without
centering:

\[
X^{\operatorname{tfidf}}=U\Sigma V^T.
\]

The purpose is not variance analysis in the classical statistical sense,
but low-rank semantic approximation for retrieval and representation.

\section{A statistical
interpretation}\label{a-statistical-interpretation}

A document-term matrix can be viewed as data sampled from an underlying
distribution. Suppose each document is generated from a mixture of
latent topics. Then terms that often occur together across documents
produce correlated columns in \(X\).

A low-rank model assumes that much of this correlation can be explained
by a small number of latent factors. Informally,

\[
X\approx AB^T,
\]

where

\[
A\in R^{n\times k}
\]

contains document factors and

\[
B\in R^{d\times k}
\]

contains term factors.

The SVD gives the best unconstrained linear factorization under squared
error. Later models, such as nonnegative matrix factorization and topic
models, impose additional constraints or probabilistic assumptions.

\section{Worked hand calculation}\label{worked-hand-calculation-1}

Consider three documents and four terms:

\[
X=
\begin{bmatrix}
1&1&0&0\\
1&0&1&0\\
0&0&1&1
\end{bmatrix}.
\]

The document Gram matrix is

\[
XX^T=
\begin{bmatrix}
2&1&0\\
1&2&1\\
0&1&2
\end{bmatrix}.
\]

Thus documents 1 and 2 overlap in one term, documents 2 and 3 overlap in
one term, and documents 1 and 3 do not overlap. The cosine similarity
matrix is obtained by dividing each row and column by document norms.
Since every row has norm \(\sqrt{2}\),

\[
C=\frac{1}{2}XX^T
=
\begin{bmatrix}
1&1/2&0\\
1/2&1&1/2\\
0&1/2&1
\end{bmatrix}.
\]

This example shows why vector-space models can be understood entirely
through matrix products.

\section{AI-assisted learning
activities}\label{ai-assisted-learning-activities}

AI tools can help students learn vector-space models, but they should be
used as mathematical assistants rather than answer machines.

\subsection{Activity 1: representation
audit}\label{activity-1-representation-audit}

Give an AI tool the following prompt:

\begin{quote}
I represented documents using TF-IDF vectors and cosine similarity. List
the mathematical assumptions behind this representation. Which
assumptions might fail for short student advising questions?
\end{quote}

Then revise the output by checking whether it mentions:

\begin{itemize}
\tightlist
\item
  bag-of-words independence from word order;
\item
  sparse high-dimensional geometry;
\item
  length normalization;
\item
  rare-term upweighting;
\item
  vocabulary mismatch;
\item
  synonymy and polysemy.
\end{itemize}

\subsection{Activity 2: explain the SVD in NLP
terms}\label{activity-2-explain-the-svd-in-nlp-terms}

Ask:

\begin{quote}
Explain \(X\approx U_k\Sigma_kV_k^T\) to a student who knows linear
algebra but not NLP. Interpret \(U_k\), \(\Sigma_k\), and \(V_k\) for a
document-term matrix.
\end{quote}

Then identify whether the explanation correctly distinguishes document
coordinates from term directions.

\subsection{Activity 3: adversarial query
design}\label{activity-3-adversarial-query-design}

Create two queries:

\begin{enumerate}
\def\labelenumi{\arabic{enumi}.}
\tightlist
\item
  one that should work well with TF-IDF cosine search;
\item
  one that should fail because it uses synonyms not present in the
  corpus.
\end{enumerate}

Ask an AI tool to predict which query will fail and why. Then test the
prediction using Python.

\section{Common mistakes}\label{common-mistakes-1}

\begin{enumerate}
\def\labelenumi{\arabic{enumi}.}
\tightlist
\item
  \textbf{Confusing documents and terms.} If \(X\in R^{n\times d}\),
  rows are documents and columns are terms. Some books use the transpose
  convention, so always check the orientation.
\item
  \textbf{Using raw inner products without thinking about length.} Long
  documents can dominate raw scores.
\item
  \textbf{Calling SVD topics probabilities.} SVD factors are real-valued
  and may be negative. They are not probability distributions.
\item
  \textbf{Assuming low-rank always improves retrieval.} Low-rank
  projection can remove noise, but it can also remove important rare
  information.
\item
  \textbf{Overinterpreting two-dimensional plots.} A 2D map may be
  useful for visualization but may not preserve all relevant distances.
\end{enumerate}

\section{Exercises}\label{exercises-6}

\subsection{Conceptual exercises}\label{conceptual-exercises-4}

\begin{enumerate}
\def\labelenumi{\arabic{enumi}.}
\tightlist
\item
  Let \(X\in R^{n\times d}\) be a document-term matrix. Explain the NLP
  meaning of \(XX^T\) and \(X^TX\).
\item
  Why does cosine similarity reduce document-length bias compared with a
  raw inner product?
\item
  Explain TF-IDF as a diagonal column-scaling operation.
\item
  What does it mean for a corpus matrix to be approximately low rank?
\item
  Why can LSA retrieve related documents that do not share exact query
  words?
\item
  Give one limitation of LSA caused by the bag-of-words assumption.
\end{enumerate}

\subsection{Mathematical exercises}\label{mathematical-exercises-3}

\begin{enumerate}
\def\labelenumi{\arabic{enumi}.}
\tightlist
\item
  Prove that cosine similarity is the inner product of normalized
  vectors.
\item
  Let \(X=U\Sigma V^T\). Show that \(XX^T=U\Sigma^2U^T\) and
  \(X^TX=V\Sigma^2V^T\).
\item
  For the matrix \[
  X=
  \begin{bmatrix}
  1&1&0\\
  1&0&1\\
  0&1&1
  \end{bmatrix},
  \] compute \(XX^T\) and the cosine similarity matrix.
\item
  Suppose \(X_k=U_k\Sigma_kV_k^T\). Show that \(X_kV_k=U_k\Sigma_k\).
\item
  Let \(q\in R^d\) be a query vector. Derive the latent query coordinate
  \(z_q=q^TV_k\) and the latent score \(s_i=z_i^Tz_q\).
\item
  If \(\sigma_1,\ldots,\sigma_r\) are the singular values of \(X\),
  compute the fraction of Frobenius energy captured by the first \(k\)
  components.
\end{enumerate}

\subsection{Computational exercises}\label{computational-exercises-1}

\begin{enumerate}
\def\labelenumi{\arabic{enumi}.}
\tightlist
\item
  Add three new documents to the corpus and recompute the TF-IDF matrix.
\item
  Compare search rankings using counts, TF-IDF, and LSA with \(k=2\).
\item
  Plot the singular values of the document-term matrix.
\item
  Remove common words such as ``models'' and ``language'' and observe
  how the retrieval rankings change.
\item
  Replace cosine similarity with Euclidean distance after row
  normalization. What changes?
\item
  Try \(k=1,2,3,4\) for LSA retrieval. Which value gives the most
  sensible ranking for the sample query?
\end{enumerate}

\section{Mini-project: a small mathematical search
engine}\label{mini-project-a-small-mathematical-search-engine-1}

Build a small search engine for a focused corpus, such as course
descriptions, research abstracts, advising FAQs, or textbook chapter
summaries.

Your report should include:

\begin{enumerate}
\def\labelenumi{\arabic{enumi}.}
\tightlist
\item
  the tokenization rule;
\item
  the vocabulary size;
\item
  the definition of the document-term matrix;
\item
  the weighting scheme;
\item
  the similarity measure;
\item
  a comparison between TF-IDF search and LSA search;
\item
  at least one failure case;
\item
  a mathematical explanation of the failure case.
\end{enumerate}

A strong project should include both formulas and code. The goal is not
only to produce rankings, but to explain the geometry behind the
rankings.

\section{Lab connection}\label{lab-connection-3}

The corresponding Python lab should guide students through:

\begin{itemize}
\tightlist
\item
  constructing document-term matrices from scratch;
\item
  implementing cosine similarity;
\item
  building a TF-IDF search engine;
\item
  computing the SVD of a weighted corpus matrix;
\item
  visualizing document embeddings in two dimensions;
\item
  comparing raw and low-rank retrieval;
\item
  writing a short mathematical interpretation of the results.
\end{itemize}

This chapter is the bridge from count-based NLP to embedding-based NLP.
In the next chapter, we move from document vectors to word embeddings,
where the goal is to learn dense vector representations of individual
tokens from co-occurrence structure.

\chapter{Word Embeddings}\label{word-embeddings}

Distributional semantics, matrix factorization, skip-gram models, and
geometry of meaning

\hfill\break

In Chapter 7, a document was represented as a vector in vocabulary
space. The coordinates of that vector were counts, TF-IDF weights, or
low-rank latent coordinates derived from a document-term matrix. In this
chapter we reverse the viewpoint: instead of representing documents by
words, we represent \textbf{words themselves} as vectors.

A word embedding is a map from a discrete vocabulary into a continuous
vector space. This is one of the central mathematical ideas in modern
NLP. It converts symbolic data into geometric data: words have lengths,
angles, neighborhoods, projections, and directions. Once words become
vectors, we can use linear algebra, probability, optimization, and
geometry to study meaning.

\section{Learning goals}\label{learning-goals-7}

After studying this chapter, you should be able to:

\begin{enumerate}
\def\labelenumi{\arabic{enumi}.}
\tightlist
\item
  define a word embedding as a function \(e:V\to R^m\);
\item
  explain how one-hot vectors differ from dense embeddings;
\item
  construct co-occurrence matrices from a corpus;
\item
  interpret the distributional hypothesis mathematically;
\item
  compute empirical probabilities for word-context pairs;
\item
  define pointwise mutual information and positive PMI;
\item
  connect matrix factorization to word embeddings;
\item
  derive the skip-gram softmax objective;
\item
  explain negative sampling as a computational approximation;
\item
  state the GloVe weighted least-squares objective;
\item
  analyze embedding geometry using norms, cosine similarity, and
  analogies;
\item
  implement toy embeddings from scratch in Python;
\item
  visualize embedding spaces with Plotly and Plotly.js;
\item
  use AI tools to critique embedding assumptions and detect possible
  bias.
\end{enumerate}

Throughout this chapter, \(V=\{v_1,\ldots,v_d\}\) denotes a finite
vocabulary, \(d=|V|\) is the vocabulary size, and \(m\) is the embedding
dimension, usually with \(m\ll d\).

\section{From symbols to vectors}\label{from-symbols-to-vectors}

A token is a discrete object. If the vocabulary is

\[
V=\{v_1,v_2,\ldots,v_d\},
\]

then a token \(v_j\) can be represented by the one-hot vector

\[
\delta_j=(0,\ldots,0,1,0,\ldots,0)^T\in R^d,
\]

where the \(1\) occurs in position \(j\). This representation is exact
but not semantic. For \(i\neq j\),

\[
\delta_i^T\delta_j=0.
\]

Thus every pair of distinct words is equally dissimilar. The words
``cat'' and ``dog'' are no closer than ``cat'' and ``gradient''.

A word embedding replaces the one-hot representation with a dense
vector:

\[
e:V\longrightarrow R^m,
\qquad
v_j\longmapsto e_j.
\]

Equivalently, the embedding is stored in a matrix

\[
E=\begin{bmatrix}
--- & e_1^T & ---\\
--- & e_2^T & ---\\
& \vdots & \\
--- & e_d^T & ---
\end{bmatrix}
\in R^{d\times m}.
\]

The embedding of token \(v_j\) is the \(j\)th row of \(E\):

\[
e_j=E_{j:}\in R^m.
\]

If a token is encoded as a one-hot vector \(\delta_j\), then the lookup
operation can be written as matrix multiplication:

\[
\delta_j^T E=e_j^T.
\]

\textbf{Embedding lookup.} A word embedding matrix
\(E\in R^{d\times m}\) converts a one-hot token vector into a dense
vector:

\[
\delta_j^TE=e_j^T.
\]

This is why an embedding layer in a neural network is mathematically a
trainable matrix lookup.

\section{Why dense embeddings are
useful}\label{why-dense-embeddings-are-useful}

One-hot vectors live in a high-dimensional coordinate system where each
coordinate is an arbitrary vocabulary index. Dense embeddings live in a
lower-dimensional representation space in which geometry is learned from
data.

The central hope is that semantic or syntactic relationships become
geometric relationships:

\begin{itemize}
\tightlist
\item
  similar words have small distance or large cosine similarity;
\item
  related categories form clusters;
\item
  some relationships correspond approximately to directions;
\item
  analogy tasks can be expressed using vector arithmetic;
\item
  downstream models can use continuous features rather than sparse
  symbolic indicators.
\end{itemize}

This does not mean that embeddings contain ``true meaning.'' It means
that the training objective forces words with similar statistical
behavior to receive related vectors.

\section{The distributional
hypothesis}\label{the-distributional-hypothesis}

The guiding principle behind many word embedding methods is the
\textbf{distributional hypothesis}:

\begin{quote}
Words that occur in similar contexts tend to have related meanings.
\end{quote}

Mathematically, each word \(w\in V\) can be described by a context
distribution

\[
P(C=c\mid W=w),
\]

where \(C\) is a random context word. If two words \(w\) and \(w'\) have
similar conditional context distributions, then they should have nearby
representations.

One way to state this is:

\[
P(C=\cdot\mid W=w)\approx P(C=\cdot\mid W=w')
\quad\Longrightarrow\quad
e(w)\approx e(w').
\]

This is a statistical assumption, not a theorem. It is often useful, but
it can fail when words have multiple meanings, when corpora are biased,
or when context alone is not enough.

\textbf{Distributional semantics.} A word embedding is a
finite-dimensional summary of a word's empirical context distribution.
The geometry of the embedding space reflects statistical regularities in
the corpus.

\section{Co-occurrence matrices}\label{co-occurrence-matrices}

Let a corpus be represented as a token sequence

\[
w_1,w_2,\ldots,w_T,
\qquad w_t\in V.
\]

Choose a context window size \(h\geq 1\). For each position \(t\), the
context of \(w_t\) is the set of nearby positions

\[
\{s:1\leq s\leq T,\ 0<|s-t|\leq h\}.
\]

The word-context co-occurrence matrix is

\[
C\in R_{\geq 0}^{d\times d},
\]

where

\[
C_{ij}=\#\{(t,s):w_t=v_i,\ w_s=v_j,\ 0<|s-t|\leq h\}.
\]

Rows correspond to center words, and columns correspond to context
words. The row

\[
C_{i:}
\]

is a count-based description of the contexts in which word \(v_i\)
appears.

The simplest possible word representation is to use the row vector
\(C_{i:}\) itself. This is usually high-dimensional and sparse.
Embedding methods compress or transform this information into
lower-dimensional vectors.

\section{Empirical probabilities}\label{empirical-probabilities}

Let

\[
N=\sum_{i=1}^d\sum_{j=1}^d C_{ij}
\]

be the total number of observed word-context pairs. The empirical joint
distribution is

\[
\widehat P(W=v_i,C=v_j)=\frac{C_{ij}}{N}.
\]

The marginal distributions are

\[
\widehat P(W=v_i)=\frac{\sum_j C_{ij}}{N},
\qquad
\widehat P(C=v_j)=\frac{\sum_i C_{ij}}{N}.
\]

The empirical conditional distribution of contexts given a word is

\[
\widehat P(C=v_j\mid W=v_i)
=
\frac{C_{ij}}{\sum_{k=1}^d C_{ik}}.
\]

This conditional distribution is the probabilistic object that the
embedding should summarize.

\section{PMI and PPMI}\label{pmi-and-ppmi}

A raw co-occurrence count is not enough. Common words such as ``the'' or
``and'' co-occur with many words simply because they are frequent.
Pointwise mutual information corrects for marginal frequency.

The empirical pointwise mutual information of a word-context pair is

\[
\operatorname{PMI}(v_i,v_j)
=
\log\frac{\widehat P(W=v_i,C=v_j)}
{\widehat P(W=v_i)\widehat P(C=v_j)}.
\]

Equivalently, using counts,

\[
\operatorname{PMI}(v_i,v_j)
=
\log\frac{C_{ij}N}{\left(\sum_k C_{ik}\right)\left(\sum_k C_{kj}\right)},
\]

when \(C_{ij}>0\).

Positive PMI keeps only positive association:

\[
\operatorname{PPMI}(v_i,v_j)=\max\{\operatorname{PMI}(v_i,v_j),0\}.
\]

A positive PMI value means that the pair occurs more often than expected
under independence. A negative PMI value means that the pair occurs less
often than expected under independence.

\textbf{PMI compares joint occurrence with independence.}

\[
\operatorname{PMI}(w,c)=\log\frac{P(w,c)}{P(w)P(c)}.
\]

If \(W\) and \(C\) were independent, then \(P(w,c)=P(w)P(c)\) and the
PMI would be \(0\).

\section{Matrix factorization
viewpoint}\label{matrix-factorization-viewpoint}

Let \(M\in R^{d\times d}\) be a transformed co-occurrence matrix, such
as a PPMI matrix. A low-rank embedding seeks matrices

\[
E\in R^{d\times m},
\qquad
F\in R^{d\times m},
\]

such that

\[
M\approx EF^T.
\]

Here \(E\) contains word embeddings and \(F\) contains context
embeddings. The approximation says

\[
M_{ij}\approx e_i^Tf_j.
\]

Thus the inner product between a word vector and a context vector
approximates a transformed co-occurrence statistic.

A classical way to obtain such embeddings is the singular value
decomposition. Suppose

\[
M=U\Sigma V^T.
\]

The rank-\(m\) approximation is

\[
M_m=U_m\Sigma_m V_m^T.
\]

A common embedding choice is

\[
E=U_m\Sigma_m^{1/2},
\qquad
F=V_m\Sigma_m^{1/2}.
\]

Then

\[
EF^T=U_m\Sigma_mV_m^T=M_m.
\]

This directly connects word embeddings to the SVD theory from Chapter 7.

\section{Prediction viewpoint}\label{prediction-viewpoint}

Matrix factorization is one route to embeddings. Neural prediction is
another. In the prediction viewpoint, embeddings are learned because
they help predict nearby words.

Let \(w\) be a center word and \(c\) be a context word. The skip-gram
model assigns

\[
P_\theta(c\mid w)=
\frac{\exp(u_c^Tv_w)}
{\sum_{c'\in V}\exp(u_{c'}^Tv_w)},
\]

where

\begin{itemize}
\tightlist
\item
  \(v_w\in R^m\) is the input embedding of the center word;
\item
  \(u_c\in R^m\) is the output embedding of the context word;
\item
  \(\theta\) denotes all embedding parameters.
\end{itemize}

Given observed word-context pairs \(D=\{(w_\ell,c_\ell)\}_{\ell=1}^N\),
the negative log-likelihood is

\[
L(\theta)
=-\sum_{\ell=1}^N \log P_\theta(c_\ell\mid w_\ell).
\]

Equivalently,

\[
L(\theta)
=\sum_{\ell=1}^N
\left[
- u_{c_\ell}^Tv_{w_\ell}
+\log\sum_{c'\in V}\exp(u_{c'}^Tv_{w_\ell})
\right].
\]

This is a multiclass cross-entropy objective.

\section{Gradient of the skip-gram softmax
loss}\label{gradient-of-the-skip-gram-softmax-loss}

For one observed pair \((w,c)\), define scores

\[
z_k=u_k^Tv_w,
\qquad k=1,\ldots,d,
\]

and probabilities

\[
p_k=\frac{\exp(z_k)}{\sum_{r=1}^d\exp(z_r)}.
\]

The single-pair loss is

\[
\ell(w,c)=-\log p_c.
\]

The gradient with respect to a score is

\[
\frac{\partial \ell}{\partial z_k}=p_k-1_{\{k=c\}}.
\]

Therefore,

\[
\frac{\partial \ell}{\partial v_w}
=\sum_{k=1}^d (p_k-1_{\{k=c\}})u_k,
\]

and

\[
\frac{\partial \ell}{\partial u_k}
=(p_k-1_{\{k=c\}})v_w.
\]

This formula shows the learning mechanism:

\begin{itemize}
\tightlist
\item
  the true context vector \(u_c\) is pulled closer to \(v_w\);
\item
  all other context vectors are pushed away in proportion to their
  predicted probabilities;
\item
  the update is global because the softmax denominator sums over the
  whole vocabulary.
\end{itemize}

\section{Negative sampling}\label{negative-sampling}

The full softmax is expensive when \(d\) is large. Negative sampling
replaces the multiclass softmax problem by several binary classification
problems.

For an observed positive pair \((w,c)\) and negative samples
\(n_1,\ldots,n_K\), the negative-sampling objective for one pair is
often written as

\[
\ell_{\text{NS}}(w,c)
= -\log\sigma(u_c^Tv_w)
-\sum_{r=1}^K\log\sigma(-u_{n_r}^Tv_w),
\]

where

\[
\sigma(x)=\frac{1}{1+e^{-x}}.
\]

The positive pair is encouraged to have a large inner product, while
sampled negative pairs are encouraged to have small inner products.

Negative sampling changes the training objective. It is not merely a
computational trick for the same loss. It learns vectors by asking
whether a word-context pair came from the corpus or from a noise
distribution.

\section{GloVe as weighted least
squares}\label{glove-as-weighted-least-squares}

GloVe is another influential embedding method. It begins with a
co-occurrence matrix \(C\) and learns word vectors \(v_i\), context
vectors \(u_j\), and biases \(a_i,b_j\) by minimizing

\[
J=
\sum_{i=1}^d\sum_{j=1}^d
f(C_{ij})
\left(v_i^Tu_j+a_i+b_j-\log C_{ij}\right)^2.
\]

The weight function \(f(C_{ij})\) reduces the influence of rare and
extremely frequent pairs. The model tries to make inner products
approximate log co-occurrence counts, up to additive word and context
biases.

The mathematical form is a weighted matrix approximation problem. Unlike
ordinary SVD, it does not assign equal weight to every matrix entry.

\section{Embedding geometry}\label{embedding-geometry}

Once words have vectors, we can analyze them geometrically. The most
common similarity measure is cosine similarity:

\[
\operatorname{cos}(e_i,e_j)
=
\frac{e_i^Te_j}{\lVert e_i\rVert_2\lVert e_j\rVert_2}.
\]

Nearest-neighbor search asks for

\[
\operatorname{NN}(w)=\arg\max_{v\in V,\ v\neq w}\operatorname{cos}(e(w),e(v)).
\]

Analogy tasks use vector arithmetic. Given words \(a,b,c\), the
classical query

\[
a:b\approx c:?
\]

is approximated by

\[
\hat d
=
\arg\max_{w\in V\setminus\{a,b,c\}}
\operatorname{cos}\left(e(w),e(b)-e(a)+e(c)\right).
\]

A famous informal example is

\[
e(\text{king})-e(\text{man})+e(\text{woman})\approx e(\text{queen}).
\]

This behavior is not guaranteed by the definition of an embedding. It
emerges only when the training data and objective organize certain
relationships approximately as linear directions.

\section{What embeddings do not
solve}\label{what-embeddings-do-not-solve}

Word embeddings are powerful, but they have limitations.

First, a single vector per word cannot represent different meanings of
the same word. The word ``bank'' may refer to a financial institution or
the side of a river. A static embedding averages across contexts.

Second, embeddings can encode social and historical biases present in
the training corpus. Similarity in the embedding space reflects patterns
in data, not moral or factual correctness.

Third, nearest neighbors can be unstable in high dimensions. Very
frequent words, anisotropy, and hubness can distort neighborhood
structure.

Fourth, word embeddings alone do not model word order. They represent
tokens, not full contextual syntax.

These limitations motivate contextual embeddings and transformers in
later chapters.

\chapter{Python examples}\label{python-examples}

We now build a small embedding pipeline from scratch. The examples are
intentionally small so that every mathematical object is visible.

\section{A tiny corpus}\label{a-tiny-corpus}

\begin{Shaded}
\begin{Highlighting}[]
\ImportTok{import}\NormalTok{ re}
\ImportTok{from}\NormalTok{ collections }\ImportTok{import}\NormalTok{ Counter, defaultdict}

\ImportTok{import}\NormalTok{ numpy }\ImportTok{as}\NormalTok{ np}
\ImportTok{import}\NormalTok{ pandas }\ImportTok{as}\NormalTok{ pd}

\NormalTok{corpus }\OperatorTok{=}\NormalTok{ [}
    \StringTok{"linear algebra studies vectors matrices and spaces"}\NormalTok{,}
    \StringTok{"singular value decomposition finds low dimensional structure"}\NormalTok{,}
    \StringTok{"probability models uncertainty in language data"}\NormalTok{,}
    \StringTok{"language models predict words from context"}\NormalTok{,}
    \StringTok{"word embeddings represent words as vectors"}\NormalTok{,}
    \StringTok{"vectors and matrices are central to machine learning"}\NormalTok{,}
    \StringTok{"context helps determine the meaning of words"}\NormalTok{,}
    \StringTok{"statistics and probability help evaluate language models"}\NormalTok{,}
    \StringTok{"machine learning learns representations from data"}\NormalTok{,}
    \StringTok{"embeddings place related words near each other"}
\NormalTok{]}

\KeywordTok{def}\NormalTok{ tokenize(text):}
    \ControlFlowTok{return}\NormalTok{ re.findall(}\VerbatimStringTok{r"}\PreprocessorTok{[a{-}z]}\OperatorTok{+}\VerbatimStringTok{"}\NormalTok{, text.lower())}

\NormalTok{tokenized }\OperatorTok{=}\NormalTok{ [tokenize(doc) }\ControlFlowTok{for}\NormalTok{ doc }\KeywordTok{in}\NormalTok{ corpus]}
\NormalTok{tokenized[:}\DecValTok{2}\NormalTok{]}
\end{Highlighting}
\end{Shaded}

\begin{verbatim}
[['linear', 'algebra', 'studies', 'vectors', 'matrices', 'and', 'spaces'],
 ['singular',
  'value',
  'decomposition',
  'finds',
  'low',
  'dimensional',
  'structure']]
\end{verbatim}

\section{Vocabulary construction}\label{vocabulary-construction}

\begin{Shaded}
\begin{Highlighting}[]
\NormalTok{counts }\OperatorTok{=}\NormalTok{ Counter(token }\ControlFlowTok{for}\NormalTok{ doc }\KeywordTok{in}\NormalTok{ tokenized }\ControlFlowTok{for}\NormalTok{ token }\KeywordTok{in}\NormalTok{ doc)}
\NormalTok{vocab }\OperatorTok{=} \BuiltInTok{sorted}\NormalTok{(counts)}
\NormalTok{word\_to\_id }\OperatorTok{=}\NormalTok{ \{word: i }\ControlFlowTok{for}\NormalTok{ i, word }\KeywordTok{in} \BuiltInTok{enumerate}\NormalTok{(vocab)\}}
\NormalTok{id\_to\_word }\OperatorTok{=}\NormalTok{ \{i: word }\ControlFlowTok{for}\NormalTok{ word, i }\KeywordTok{in}\NormalTok{ word\_to\_id.items()\}}

\BuiltInTok{print}\NormalTok{(}\StringTok{"vocabulary size:"}\NormalTok{, }\BuiltInTok{len}\NormalTok{(vocab))}
\BuiltInTok{print}\NormalTok{(vocab)}
\end{Highlighting}
\end{Shaded}

\begin{verbatim}
vocabulary size: 48
['algebra', 'and', 'are', 'as', 'central', 'context', 'data', 'decomposition', 'determine', 'dimensional', 'each', 'embeddings', 'evaluate', 'finds', 'from', 'help', 'helps', 'in', 'language', 'learning', 'learns', 'linear', 'low', 'machine', 'matrices', 'meaning', 'models', 'near', 'of', 'other', 'place', 'predict', 'probability', 'related', 'represent', 'representations', 'singular', 'spaces', 'statistics', 'structure', 'studies', 'the', 'to', 'uncertainty', 'value', 'vectors', 'word', 'words']
\end{verbatim}

\section{Word-context pairs}\label{word-context-pairs}

The following function constructs center-context pairs using a symmetric
window of radius \(h\).

\begin{Shaded}
\begin{Highlighting}[]
\KeywordTok{def}\NormalTok{ make\_pairs(tokenized\_docs, window}\OperatorTok{=}\DecValTok{2}\NormalTok{):}
\NormalTok{    pairs }\OperatorTok{=}\NormalTok{ []}
    \ControlFlowTok{for}\NormalTok{ doc }\KeywordTok{in}\NormalTok{ tokenized\_docs:}
\NormalTok{        ids }\OperatorTok{=}\NormalTok{ [word\_to\_id[w] }\ControlFlowTok{for}\NormalTok{ w }\KeywordTok{in}\NormalTok{ doc]}
        \ControlFlowTok{for}\NormalTok{ t, center }\KeywordTok{in} \BuiltInTok{enumerate}\NormalTok{(ids):}
\NormalTok{            left }\OperatorTok{=} \BuiltInTok{max}\NormalTok{(}\DecValTok{0}\NormalTok{, t }\OperatorTok{{-}}\NormalTok{ window)}
\NormalTok{            right }\OperatorTok{=} \BuiltInTok{min}\NormalTok{(}\BuiltInTok{len}\NormalTok{(ids), t }\OperatorTok{+}\NormalTok{ window }\OperatorTok{+} \DecValTok{1}\NormalTok{)}
            \ControlFlowTok{for}\NormalTok{ s }\KeywordTok{in} \BuiltInTok{range}\NormalTok{(left, right):}
                \ControlFlowTok{if}\NormalTok{ s }\OperatorTok{!=}\NormalTok{ t:}
\NormalTok{                    pairs.append((center, ids[s]))}
    \ControlFlowTok{return}\NormalTok{ pairs}

\NormalTok{pairs }\OperatorTok{=}\NormalTok{ make\_pairs(tokenized, window}\OperatorTok{=}\DecValTok{2}\NormalTok{)}
\BuiltInTok{print}\NormalTok{(}\StringTok{"number of word{-}context pairs:"}\NormalTok{, }\BuiltInTok{len}\NormalTok{(pairs))}
\BuiltInTok{print}\NormalTok{([(id\_to\_word[i], id\_to\_word[j]) }\ControlFlowTok{for}\NormalTok{ i, j }\KeywordTok{in}\NormalTok{ pairs[:}\DecValTok{12}\NormalTok{]])}
\end{Highlighting}
\end{Shaded}

\begin{verbatim}
number of word-context pairs: 208
[('linear', 'algebra'), ('linear', 'studies'), ('algebra', 'linear'), ('algebra', 'studies'), ('algebra', 'vectors'), ('studies', 'linear'), ('studies', 'algebra'), ('studies', 'vectors'), ('studies', 'matrices'), ('vectors', 'algebra'), ('vectors', 'studies'), ('vectors', 'matrices')]
\end{verbatim}

\section{Co-occurrence matrix}\label{co-occurrence-matrix}

\begin{Shaded}
\begin{Highlighting}[]
\KeywordTok{def}\NormalTok{ cooccurrence\_matrix(pairs, d):}
\NormalTok{    C }\OperatorTok{=}\NormalTok{ np.zeros((d, d), dtype}\OperatorTok{=}\BuiltInTok{float}\NormalTok{)}
    \ControlFlowTok{for}\NormalTok{ i, j }\KeywordTok{in}\NormalTok{ pairs:}
\NormalTok{        C[i, j] }\OperatorTok{+=} \FloatTok{1.0}
    \ControlFlowTok{return}\NormalTok{ C}

\NormalTok{C }\OperatorTok{=}\NormalTok{ cooccurrence\_matrix(pairs, }\BuiltInTok{len}\NormalTok{(vocab))}
\NormalTok{C\_df }\OperatorTok{=}\NormalTok{ pd.DataFrame(C, index}\OperatorTok{=}\NormalTok{vocab, columns}\OperatorTok{=}\NormalTok{vocab)}
\NormalTok{C\_df.iloc[:}\DecValTok{8}\NormalTok{, :}\DecValTok{8}\NormalTok{]}
\end{Highlighting}
\end{Shaded}

{\def\LTcaptype{none} % do not increment counter
\begin{longtable}[]{@{}lllllllll@{}}
\toprule\noalign{}
& algebra & and & are & as & central & context & data & decomposition \\
\midrule\noalign{}
\endhead
\bottomrule\noalign{}
\endlastfoot
algebra & 0.0 & 0.0 & 0.0 & 0.0 & 0.0 & 0.0 & 0.0 & 0.0 \\
and & 0.0 & 0.0 & 1.0 & 0.0 & 0.0 & 0.0 & 0.0 & 0.0 \\
are & 0.0 & 1.0 & 0.0 & 0.0 & 1.0 & 0.0 & 0.0 & 0.0 \\
as & 0.0 & 0.0 & 0.0 & 0.0 & 0.0 & 0.0 & 0.0 & 0.0 \\
central & 0.0 & 0.0 & 1.0 & 0.0 & 0.0 & 0.0 & 0.0 & 0.0 \\
context & 0.0 & 0.0 & 0.0 & 0.0 & 0.0 & 0.0 & 0.0 & 0.0 \\
data & 0.0 & 0.0 & 0.0 & 0.0 & 0.0 & 0.0 & 0.0 & 0.0 \\
decomposition & 0.0 & 0.0 & 0.0 & 0.0 & 0.0 & 0.0 & 0.0 & 0.0 \\
\end{longtable}
}

\section{PPMI matrix}\label{ppmi-matrix}

\begin{Shaded}
\begin{Highlighting}[]
\KeywordTok{def}\NormalTok{ ppmi\_matrix(C, eps}\OperatorTok{=}\FloatTok{1e{-}12}\NormalTok{):}
\NormalTok{    total }\OperatorTok{=}\NormalTok{ C.}\BuiltInTok{sum}\NormalTok{()}
\NormalTok{    row\_sum }\OperatorTok{=}\NormalTok{ C.}\BuiltInTok{sum}\NormalTok{(axis}\OperatorTok{=}\DecValTok{1}\NormalTok{, keepdims}\OperatorTok{=}\VariableTok{True}\NormalTok{)}
\NormalTok{    col\_sum }\OperatorTok{=}\NormalTok{ C.}\BuiltInTok{sum}\NormalTok{(axis}\OperatorTok{=}\DecValTok{0}\NormalTok{, keepdims}\OperatorTok{=}\VariableTok{True}\NormalTok{)}
\NormalTok{    expected }\OperatorTok{=}\NormalTok{ row\_sum }\OperatorTok{@}\NormalTok{ col\_sum }\OperatorTok{/} \BuiltInTok{max}\NormalTok{(total, eps)}
    \ControlFlowTok{with}\NormalTok{ np.errstate(divide}\OperatorTok{=}\StringTok{"ignore"}\NormalTok{, invalid}\OperatorTok{=}\StringTok{"ignore"}\NormalTok{):}
\NormalTok{        pmi }\OperatorTok{=}\NormalTok{ np.log((C }\OperatorTok{*}\NormalTok{ total }\OperatorTok{+}\NormalTok{ eps) }\OperatorTok{/}\NormalTok{ (row\_sum }\OperatorTok{@}\NormalTok{ col\_sum }\OperatorTok{+}\NormalTok{ eps))}
\NormalTok{    pmi[C }\OperatorTok{==} \DecValTok{0}\NormalTok{] }\OperatorTok{=} \FloatTok{0.0}
    \ControlFlowTok{return}\NormalTok{ np.maximum(pmi, }\FloatTok{0.0}\NormalTok{)}

\NormalTok{M }\OperatorTok{=}\NormalTok{ ppmi\_matrix(C)}
\NormalTok{M\_df }\OperatorTok{=}\NormalTok{ pd.DataFrame(M, index}\OperatorTok{=}\NormalTok{vocab, columns}\OperatorTok{=}\NormalTok{vocab)}
\NormalTok{M\_df.}\BuiltInTok{round}\NormalTok{(}\DecValTok{2}\NormalTok{).iloc[:}\DecValTok{8}\NormalTok{, :}\DecValTok{8}\NormalTok{]}
\end{Highlighting}
\end{Shaded}

{\def\LTcaptype{none} % do not increment counter
\begin{longtable}[]{@{}lllllllll@{}}
\toprule\noalign{}
& algebra & and & are & as & central & context & data & decomposition \\
\midrule\noalign{}
\endhead
\bottomrule\noalign{}
\endlastfoot
algebra & 0.0 & 0.00 & 0.00 & 0.0 & 0.00 & 0.0 & 0.0 & 0.0 \\
and & 0.0 & 0.00 & 1.75 & 0.0 & 0.00 & 0.0 & 0.0 & 0.0 \\
are & 0.0 & 1.75 & 0.00 & 0.0 & 2.56 & 0.0 & 0.0 & 0.0 \\
as & 0.0 & 0.00 & 0.00 & 0.0 & 0.00 & 0.0 & 0.0 & 0.0 \\
central & 0.0 & 0.00 & 2.56 & 0.0 & 0.00 & 0.0 & 0.0 & 0.0 \\
context & 0.0 & 0.00 & 0.00 & 0.0 & 0.00 & 0.0 & 0.0 & 0.0 \\
data & 0.0 & 0.00 & 0.00 & 0.0 & 0.00 & 0.0 & 0.0 & 0.0 \\
decomposition & 0.0 & 0.00 & 0.00 & 0.0 & 0.00 & 0.0 & 0.0 & 0.0 \\
\end{longtable}
}

\section{SVD embeddings from PPMI}\label{svd-embeddings-from-ppmi}

We factor the PPMI matrix and keep the first \(m\) singular directions.

\begin{Shaded}
\begin{Highlighting}[]
\KeywordTok{def}\NormalTok{ svd\_embeddings(M, m}\OperatorTok{=}\DecValTok{2}\NormalTok{):}
\NormalTok{    U, s, Vt }\OperatorTok{=}\NormalTok{ np.linalg.svd(M, full\_matrices}\OperatorTok{=}\VariableTok{False}\NormalTok{)}
\NormalTok{    E }\OperatorTok{=}\NormalTok{ U[:, :m] }\OperatorTok{*}\NormalTok{ np.sqrt(s[:m])}
\NormalTok{    F }\OperatorTok{=}\NormalTok{ Vt[:m, :].T }\OperatorTok{*}\NormalTok{ np.sqrt(s[:m])}
    \ControlFlowTok{return}\NormalTok{ E, F, s}

\NormalTok{E2, F2, singular\_values }\OperatorTok{=}\NormalTok{ svd\_embeddings(M, m}\OperatorTok{=}\DecValTok{2}\NormalTok{)}
\NormalTok{emb\_df }\OperatorTok{=}\NormalTok{ pd.DataFrame(E2, index}\OperatorTok{=}\NormalTok{vocab, columns}\OperatorTok{=}\NormalTok{[}\StringTok{"z1"}\NormalTok{, }\StringTok{"z2"}\NormalTok{])}
\NormalTok{emb\_df.head()}
\end{Highlighting}
\end{Shaded}

{\def\LTcaptype{none} % do not increment counter
\begin{longtable}[]{@{}lll@{}}
\toprule\noalign{}
& z1 & z2 \\
\midrule\noalign{}
\endhead
\bottomrule\noalign{}
\endlastfoot
algebra & -0.244986 & 2.992904e-16 \\
and & -0.554060 & -1.067563e-15 \\
are & -0.407793 & 3.111442e-17 \\
as & -0.339528 & -5.242732e-16 \\
central & -0.417045 & -7.673284e-16 \\
\end{longtable}
}

\section{Nearest neighbors}\label{nearest-neighbors}

\begin{Shaded}
\begin{Highlighting}[]
\KeywordTok{def}\NormalTok{ row\_normalize(X, eps}\OperatorTok{=}\FloatTok{1e{-}12}\NormalTok{):}
\NormalTok{    norms }\OperatorTok{=}\NormalTok{ np.linalg.norm(X, axis}\OperatorTok{=}\DecValTok{1}\NormalTok{, keepdims}\OperatorTok{=}\VariableTok{True}\NormalTok{)}
    \ControlFlowTok{return}\NormalTok{ X }\OperatorTok{/}\NormalTok{ np.maximum(norms, eps)}

\KeywordTok{def}\NormalTok{ nearest\_neighbors(word, E, k}\OperatorTok{=}\DecValTok{5}\NormalTok{):}
\NormalTok{    idx }\OperatorTok{=}\NormalTok{ word\_to\_id[word]}
\NormalTok{    En }\OperatorTok{=}\NormalTok{ row\_normalize(E)}
\NormalTok{    scores }\OperatorTok{=}\NormalTok{ En }\OperatorTok{@}\NormalTok{ En[idx]}
\NormalTok{    order }\OperatorTok{=}\NormalTok{ np.argsort(}\OperatorTok{{-}}\NormalTok{scores)}
\NormalTok{    rows }\OperatorTok{=}\NormalTok{ []}
    \ControlFlowTok{for}\NormalTok{ j }\KeywordTok{in}\NormalTok{ order:}
        \ControlFlowTok{if}\NormalTok{ j }\OperatorTok{!=}\NormalTok{ idx:}
\NormalTok{            rows.append((id\_to\_word[j], }\BuiltInTok{float}\NormalTok{(scores[j])))}
        \ControlFlowTok{if} \BuiltInTok{len}\NormalTok{(rows) }\OperatorTok{==}\NormalTok{ k:}
            \ControlFlowTok{break}
    \ControlFlowTok{return}\NormalTok{ pd.DataFrame(rows, columns}\OperatorTok{=}\NormalTok{[}\StringTok{"word"}\NormalTok{, }\StringTok{"cosine\_similarity"}\NormalTok{])}

\NormalTok{nearest\_neighbors(}\StringTok{"language"}\NormalTok{, E2, k}\OperatorTok{=}\DecValTok{6}\NormalTok{)}
\end{Highlighting}
\end{Shaded}

{\def\LTcaptype{none} % do not increment counter
\begin{longtable}[]{@{}lll@{}}
\toprule\noalign{}
& word & cosine\_similarity \\
\midrule\noalign{}
\endhead
\bottomrule\noalign{}
\endlastfoot
0 & algebra & 1.0 \\
1 & meaning & 1.0 \\
2 & models & 1.0 \\
3 & near & 1.0 \\
4 & of & 1.0 \\
5 & other & 1.0 \\
\end{longtable}
}

Because the corpus is tiny, the neighbors are not stable. That is part
of the lesson: embeddings need enough data to estimate meaningful
context distributions.

\section{Compare count-context and PPMI-SVD
neighbors}\label{compare-count-context-and-ppmi-svd-neighbors}

\begin{Shaded}
\begin{Highlighting}[]
\BuiltInTok{print}\NormalTok{(}\StringTok{"Neighbors from raw co{-}occurrence rows"}\NormalTok{)}
\NormalTok{display(nearest\_neighbors(}\StringTok{"language"}\NormalTok{, C, k}\OperatorTok{=}\DecValTok{5}\NormalTok{))}

\BuiltInTok{print}\NormalTok{(}\StringTok{"Neighbors from rank{-}2 PPMI{-}SVD embeddings"}\NormalTok{)}
\NormalTok{display(nearest\_neighbors(}\StringTok{"language"}\NormalTok{, E2, k}\OperatorTok{=}\DecValTok{5}\NormalTok{))}
\end{Highlighting}
\end{Shaded}

\begin{verbatim}
Neighbors from raw co-occurrence rows
\end{verbatim}

{\def\LTcaptype{none} % do not increment counter
\begin{longtable}[]{@{}lll@{}}
\toprule\noalign{}
& word & cosine\_similarity \\
\midrule\noalign{}
\endhead
\bottomrule\noalign{}
\endlastfoot
0 & probability & 0.645497 \\
1 & in & 0.632456 \\
2 & uncertainty & 0.474342 \\
3 & evaluate & 0.474342 \\
4 & models & 0.400000 \\
\end{longtable}
}

\begin{verbatim}
Neighbors from rank-2 PPMI-SVD embeddings
\end{verbatim}

{\def\LTcaptype{none} % do not increment counter
\begin{longtable}[]{@{}lll@{}}
\toprule\noalign{}
& word & cosine\_similarity \\
\midrule\noalign{}
\endhead
\bottomrule\noalign{}
\endlastfoot
0 & algebra & 1.0 \\
1 & meaning & 1.0 \\
2 & models & 1.0 \\
3 & near & 1.0 \\
4 & of & 1.0 \\
\end{longtable}
}

Raw co-occurrence and PPMI-SVD embeddings can give different neighbors
because they encode different geometry. Raw counts emphasize frequency;
PPMI emphasizes association beyond chance; SVD compresses the result
into a low-dimensional subspace.

\section{A tiny skip-gram softmax
calculation}\label{a-tiny-skip-gram-softmax-calculation}

The next code implements the skip-gram probability for one center word.

\begin{Shaded}
\begin{Highlighting}[]
\NormalTok{rng }\OperatorTok{=}\NormalTok{ np.random.default\_rng(}\DecValTok{7243}\NormalTok{)}
\NormalTok{d }\OperatorTok{=} \BuiltInTok{len}\NormalTok{(vocab)}
\NormalTok{m }\OperatorTok{=} \DecValTok{4}

\NormalTok{V\_in }\OperatorTok{=}\NormalTok{ rng.normal(}\DecValTok{0}\NormalTok{, }\FloatTok{0.1}\NormalTok{, size}\OperatorTok{=}\NormalTok{(d, m))}
\NormalTok{U\_out }\OperatorTok{=}\NormalTok{ rng.normal(}\DecValTok{0}\NormalTok{, }\FloatTok{0.1}\NormalTok{, size}\OperatorTok{=}\NormalTok{(d, m))}

\KeywordTok{def}\NormalTok{ softmax(z):}
\NormalTok{    z }\OperatorTok{=}\NormalTok{ z }\OperatorTok{{-}}\NormalTok{ np.}\BuiltInTok{max}\NormalTok{(z)}
\NormalTok{    ez }\OperatorTok{=}\NormalTok{ np.exp(z)}
    \ControlFlowTok{return}\NormalTok{ ez }\OperatorTok{/}\NormalTok{ ez.}\BuiltInTok{sum}\NormalTok{()}

\KeywordTok{def}\NormalTok{ skipgram\_probs(center\_word):}
\NormalTok{    i }\OperatorTok{=}\NormalTok{ word\_to\_id[center\_word]}
\NormalTok{    scores }\OperatorTok{=}\NormalTok{ U\_out }\OperatorTok{@}\NormalTok{ V\_in[i]}
\NormalTok{    probs }\OperatorTok{=}\NormalTok{ softmax(scores)}
    \ControlFlowTok{return}\NormalTok{ pd.DataFrame(\{}\StringTok{"word"}\NormalTok{: vocab, }\StringTok{"probability"}\NormalTok{: probs\}).sort\_values(}
        \StringTok{"probability"}\NormalTok{, ascending}\OperatorTok{=}\VariableTok{False}
\NormalTok{    )}

\NormalTok{skipgram\_probs(}\StringTok{"language"}\NormalTok{).head(}\DecValTok{8}\NormalTok{)}
\end{Highlighting}
\end{Shaded}

{\def\LTcaptype{none} % do not increment counter
\begin{longtable}[]{@{}lll@{}}
\toprule\noalign{}
& word & probability \\
\midrule\noalign{}
\endhead
\bottomrule\noalign{}
\endlastfoot
39 & structure & 0.021541 \\
33 & related & 0.021359 \\
44 & value & 0.021199 \\
32 & probability & 0.021174 \\
14 & from & 0.021171 \\
47 & words & 0.021142 \\
8 & determine & 0.021086 \\
2 & are & 0.021082 \\
\end{longtable}
}

At initialization, the probabilities are almost uniform because the
random inner products are small. Training changes the vectors so
observed context words receive higher probability.

\section{One negative-sampling
update}\label{one-negative-sampling-update}

\begin{Shaded}
\begin{Highlighting}[]
\KeywordTok{def}\NormalTok{ sigmoid(x):}
    \ControlFlowTok{return} \FloatTok{1.0} \OperatorTok{/}\NormalTok{ (}\FloatTok{1.0} \OperatorTok{+}\NormalTok{ np.exp(}\OperatorTok{{-}}\NormalTok{x))}

\KeywordTok{def}\NormalTok{ negative\_sampling\_step(V\_in, U\_out, center, positive, negatives, lr}\OperatorTok{=}\FloatTok{0.05}\NormalTok{):}
\NormalTok{    i }\OperatorTok{=}\NormalTok{ word\_to\_id[center]}
\NormalTok{    j }\OperatorTok{=}\NormalTok{ word\_to\_id[positive]}
\NormalTok{    neg\_ids }\OperatorTok{=}\NormalTok{ [word\_to\_id[w] }\ControlFlowTok{for}\NormalTok{ w }\KeywordTok{in}\NormalTok{ negatives]}

\NormalTok{    v }\OperatorTok{=}\NormalTok{ V\_in[i].copy()}

    \CommentTok{\# Positive pair gradient for {-}log sigmoid(u\_j\^{}T v)}
\NormalTok{    score\_pos }\OperatorTok{=}\NormalTok{ U\_out[j] }\OperatorTok{@}\NormalTok{ v}
\NormalTok{    grad\_v }\OperatorTok{=}\NormalTok{ (sigmoid(score\_pos) }\OperatorTok{{-}} \FloatTok{1.0}\NormalTok{) }\OperatorTok{*}\NormalTok{ U\_out[j]}
\NormalTok{    grad\_u\_pos }\OperatorTok{=}\NormalTok{ (sigmoid(score\_pos) }\OperatorTok{{-}} \FloatTok{1.0}\NormalTok{) }\OperatorTok{*}\NormalTok{ v}

    \CommentTok{\# Negative pair gradients for {-}log sigmoid({-}u\_n\^{}T v)}
\NormalTok{    grad\_u\_negs }\OperatorTok{=}\NormalTok{ []}
    \ControlFlowTok{for}\NormalTok{ n }\KeywordTok{in}\NormalTok{ neg\_ids:}
\NormalTok{        score\_neg }\OperatorTok{=}\NormalTok{ U\_out[n] }\OperatorTok{@}\NormalTok{ v}
\NormalTok{        grad\_v }\OperatorTok{+=}\NormalTok{ sigmoid(score\_neg) }\OperatorTok{*}\NormalTok{ U\_out[n]}
\NormalTok{        grad\_u\_negs.append((n, sigmoid(score\_neg) }\OperatorTok{*}\NormalTok{ v))}

\NormalTok{    V\_in[i] }\OperatorTok{{-}=}\NormalTok{ lr }\OperatorTok{*}\NormalTok{ grad\_v}
\NormalTok{    U\_out[j] }\OperatorTok{{-}=}\NormalTok{ lr }\OperatorTok{*}\NormalTok{ grad\_u\_pos}
    \ControlFlowTok{for}\NormalTok{ n, grad\_u }\KeywordTok{in}\NormalTok{ grad\_u\_negs:}
\NormalTok{        U\_out[n] }\OperatorTok{{-}=}\NormalTok{ lr }\OperatorTok{*}\NormalTok{ grad\_u}

    \ControlFlowTok{return}\NormalTok{ V\_in, U\_out}

\NormalTok{center }\OperatorTok{=} \StringTok{"language"}
\NormalTok{positive }\OperatorTok{=} \StringTok{"models"}
\NormalTok{negatives }\OperatorTok{=}\NormalTok{ [}\StringTok{"algebra"}\NormalTok{, }\StringTok{"matrices"}\NormalTok{, }\StringTok{"singular"}\NormalTok{]}

\NormalTok{before }\OperatorTok{=}\NormalTok{ U\_out[word\_to\_id[positive]] }\OperatorTok{@}\NormalTok{ V\_in[word\_to\_id[center]]}
\NormalTok{V\_in, U\_out }\OperatorTok{=}\NormalTok{ negative\_sampling\_step(V\_in, U\_out, center, positive, negatives)}
\NormalTok{after }\OperatorTok{=}\NormalTok{ U\_out[word\_to\_id[positive]] }\OperatorTok{@}\NormalTok{ V\_in[word\_to\_id[center]]}

\BuiltInTok{print}\NormalTok{(}\StringTok{"positive inner product before:"}\NormalTok{, }\BuiltInTok{round}\NormalTok{(}\BuiltInTok{float}\NormalTok{(before), }\DecValTok{4}\NormalTok{))}
\BuiltInTok{print}\NormalTok{(}\StringTok{"positive inner product after: "}\NormalTok{, }\BuiltInTok{round}\NormalTok{(}\BuiltInTok{float}\NormalTok{(after), }\DecValTok{4}\NormalTok{))}
\end{Highlighting}
\end{Shaded}

\begin{verbatim}
positive inner product before: -0.0051
positive inner product after:  -0.0043
\end{verbatim}

The update tends to increase the inner product of the observed pair and
decrease the inner products of sampled negative pairs.

\section{Plotly visualization from
Python}\label{plotly-visualization-from-python-1}

\begin{Shaded}
\begin{Highlighting}[]
\ControlFlowTok{try}\NormalTok{:}
    \ImportTok{import}\NormalTok{ plotly.express }\ImportTok{as}\NormalTok{ px}

\NormalTok{    plot\_df }\OperatorTok{=}\NormalTok{ emb\_df.reset\_index().rename(columns}\OperatorTok{=}\NormalTok{\{}\StringTok{"index"}\NormalTok{: }\StringTok{"word"}\NormalTok{\})}
\NormalTok{    fig }\OperatorTok{=}\NormalTok{ px.scatter(}
\NormalTok{        plot\_df,}
\NormalTok{        x}\OperatorTok{=}\StringTok{"z1"}\NormalTok{,}
\NormalTok{        y}\OperatorTok{=}\StringTok{"z2"}\NormalTok{,}
\NormalTok{        text}\OperatorTok{=}\StringTok{"word"}\NormalTok{,}
\NormalTok{        hover\_name}\OperatorTok{=}\StringTok{"word"}\NormalTok{,}
\NormalTok{        title}\OperatorTok{=}\StringTok{"Toy word embedding map from PPMI{-}SVD"}
\NormalTok{    )}
\NormalTok{    fig.update\_traces(textposition}\OperatorTok{=}\StringTok{"top center"}\NormalTok{)}
\NormalTok{    fig.update\_layout(xaxis\_title}\OperatorTok{=}\StringTok{"embedding coordinate 1"}\NormalTok{, yaxis\_title}\OperatorTok{=}\StringTok{"embedding coordinate 2"}\NormalTok{)}
\NormalTok{    fig.show()}
\ControlFlowTok{except} \PreprocessorTok{Exception} \ImportTok{as}\NormalTok{ e:}
    \BuiltInTok{print}\NormalTok{(}\StringTok{"Plotly figure skipped:"}\NormalTok{, e)}
\end{Highlighting}
\end{Shaded}

\begin{verbatim}
Unable to display output for mime type(s): text/html
\end{verbatim}

Two-dimensional word embeddings obtained from PPMI and SVD.

\begin{verbatim}
Unable to display output for mime type(s): text/html
\end{verbatim}

\chapter{Interactive HTML: embedding analogy
explorer}\label{interactive-html-embedding-analogy-explorer}

Now define a target vector depending on a scalar \(\alpha\):

\[
z(\alpha)=e(\text{king})-\alpha e(\text{man})+\alpha e(\text{woman}).
\]

When \(\alpha=1\), this is the classical analogy vector

\[
e(\text{king})-e(\text{man})+e(\text{woman}).
\]

This interactive is deliberately idealized. Real embeddings are
high-dimensional, noisy, and trained from data. Nevertheless, the same
vector operation and cosine nearest-neighbor rule are used.

\chapter{Worked example: PPMI by
hand}\label{worked-example-ppmi-by-hand}

Suppose a tiny corpus produces the following co-occurrence counts:

\[
C=
\begin{bmatrix}
10 & 2\\
1 & 7
\end{bmatrix}.
\]

The total count is

\[
N=10+2+1+7=20.
\]

For entry \((1,1)\),

\[
P(W=1,C=1)=\frac{10}{20}=0.5,
\]

while

\[
P(W=1)=\frac{12}{20}=0.6,
\qquad
P(C=1)=\frac{11}{20}=0.55.
\]

Thus

\[
\operatorname{PMI}(1,1)
=\log\frac{0.5}{0.6\cdot 0.55}
=\log(1.515\ldots)>0.
\]

The pair \((1,1)\) occurs more often than expected under independence.

For entry \((1,2)\),

\[
P(W=1,C=2)=\frac{2}{20}=0.1,
\]

while

\[
P(W=1)P(C=2)=\frac{12}{20}\cdot\frac{9}{20}=0.27.
\]

Therefore

\[
\operatorname{PMI}(1,2)=\log\frac{0.1}{0.27}<0,
\]

and the PPMI value is set to zero.

\chapter{AI-assisted learning
components}\label{ai-assisted-learning-components}

AI tools can help students explore word embeddings, but they should be
used as mathematical assistants rather than authorities.

\section{Prompt 1: explain an embedding method
mathematically}\label{prompt-1-explain-an-embedding-method-mathematically}

Use an AI assistant to answer:

\begin{quote}
Explain the skip-gram model as a maximum likelihood problem. Define the
vocabulary, the input embedding matrix, the output embedding matrix, the
softmax probability, and the negative log-likelihood. Avoid
implementation details until the mathematical objective is clear.
\end{quote}

Then check whether the answer includes:

\begin{itemize}
\tightlist
\item
  a conditional probability \(P(c\mid w)\);
\item
  inner products between word and context vectors;
\item
  a softmax denominator over the vocabulary;
\item
  a sum over observed word-context pairs;
\item
  a clear distinction between input and output embeddings.
\end{itemize}

\section{Prompt 2: audit a nearest-neighbor
result}\label{prompt-2-audit-a-nearest-neighbor-result}

Ask:

\begin{quote}
I trained word embeddings and found the nearest neighbors of the word
``model''. How can I tell whether the neighbors reflect semantics,
syntax, frequency effects, or corpus bias? Give a mathematical
checklist.
\end{quote}

A good answer should mention cosine similarity, vector norms, word
frequency, training corpus, context window, and the difference between
association and truth.

\section{Prompt 3: derive the
gradient}\label{prompt-3-derive-the-gradient}

Ask:

\begin{quote}
For the skip-gram softmax loss \(\ell(w,c)=-\log P(c\mid w)\), derive
\(\partial \ell/\partial v_w\) and explain the sign of the update.
\end{quote}

Then compare the result with the gradient derived earlier in this
chapter.

\section{Prompt 4: detect hidden
assumptions}\label{prompt-4-detect-hidden-assumptions}

Ask:

\begin{quote}
What assumptions are hidden in the distributional hypothesis? Give
examples where words appear in similar contexts but should not be
considered semantically equivalent.
\end{quote}

This encourages students to connect mathematical modeling assumptions
with linguistic and social limitations.

\chapter{Exercises}\label{exercises-7}

\section{Conceptual exercises}\label{conceptual-exercises-5}

\begin{enumerate}
\def\labelenumi{\arabic{enumi}.}
\tightlist
\item
  Explain why one-hot vectors do not encode semantic similarity.
\item
  What information is contained in a row of a co-occurrence matrix?
\item
  Why does PMI correct for word frequency better than raw co-occurrence
  counts?
\item
  What is the difference between a word embedding matrix and a context
  embedding matrix?
\item
  In what sense is GloVe a matrix approximation method?
\item
  Why can static word embeddings struggle with polysemy?
\item
  Why might cosine similarity be preferred over Euclidean distance for
  word embeddings?
\item
  What does it mean for an analogy to be represented by a linear
  direction?
\end{enumerate}

\section{Mathematical exercises}\label{mathematical-exercises-4}

\begin{enumerate}
\def\labelenumi{\arabic{enumi}.}
\tightlist
\item
  Let \(E\in R^{d\times m}\) and \(\delta_j\in R^d\) be the \(j\)th
  one-hot vector. Prove that \(\delta_j^TE\) returns the \(j\)th row of
  \(E\).
\item
  Let \(C\) be a co-occurrence matrix. Derive the empirical joint and
  marginal distributions from \(C\).
\item
  Show that \(\operatorname{PMI}(w,c)=0\) when \(W\) and \(C\) are
  independent.
\item
  Let \(M=U\Sigma V^T\) and define \(E=U_m\Sigma_m^{1/2}\) and
  \(F=V_m\Sigma_m^{1/2}\). Show that \(EF^T\) equals the rank-\(m\)
  truncated SVD approximation.
\item
  Derive the gradient of the skip-gram softmax loss with respect to the
  center-word embedding.
\item
  For the negative-sampling loss, compute the derivative of
  \(-\log\sigma(u_c^Tv_w)\) with respect to \(v_w\).
\item
  Let \(x,y\in R^m\) be nonzero. Prove that cosine similarity is
  invariant under positive rescaling of either vector.
\item
  Suppose \(e(b)-e(a)\approx e(d)-e(c)\). Explain why this supports the
  analogy \(a:b\approx c:d\).
\end{enumerate}

\section{Computational exercises}\label{computational-exercises-2}

\begin{enumerate}
\def\labelenumi{\arabic{enumi}.}
\tightlist
\item
  Change the context window size in the Python code from \(h=2\) to
  \(h=1\) and \(h=3\). Compare the co-occurrence matrices.
\item
  Compute PPMI embeddings with dimensions \(m=2,3,5\). How do nearest
  neighbors change?
\item
  Replace PPMI with raw co-occurrence counts before SVD. Compare the
  geometry.
\item
  Implement a function that returns the top \(k\) word pairs with
  largest PMI values.
\item
  Train the toy negative-sampling model for several epochs over all
  word-context pairs. Visualize the resulting embeddings.
\item
  Use a larger text sample and compare embeddings trained from different
  corpora.
\item
  Investigate whether frequent words have larger embedding norms in your
  experiment.
\item
  Create a Plotly visualization that lets the user choose a word and
  displays its nearest neighbors.
\end{enumerate}

\chapter{Mini-project: build and audit toy word
embeddings}\label{mini-project-build-and-audit-toy-word-embeddings}

Choose a small corpus related to one domain, such as mathematics course
descriptions, university advising pages, research abstracts, or textbook
chapter summaries. Build word embeddings from scratch using the
following pipeline:

\begin{enumerate}
\def\labelenumi{\arabic{enumi}.}
\tightlist
\item
  tokenize the corpus;
\item
  build a vocabulary;
\item
  construct a word-context co-occurrence matrix;
\item
  compute a PPMI matrix;
\item
  compute SVD embeddings;
\item
  find nearest neighbors for selected words;
\item
  visualize embeddings in two dimensions;
\item
  write a short mathematical audit of the results.
\end{enumerate}

Your audit should answer:

\begin{itemize}
\tightlist
\item
  What does a row of the co-occurrence matrix mean?
\item
  What does PMI emphasize?
\item
  How much information is lost when projecting to two dimensions?
\item
  Which nearest neighbors are sensible?
\item
  Which nearest neighbors are artifacts of the corpus?
\item
  What would change if the corpus were larger?
\end{itemize}

\chapter{Lab connection}\label{lab-connection-4}

The computer lab for this chapter should guide students through a
complete word-embedding workflow:

\begin{enumerate}
\def\labelenumi{\arabic{enumi}.}
\tightlist
\item
  build co-occurrence matrices;
\item
  compute empirical probabilities;
\item
  compute PMI and PPMI;
\item
  obtain SVD embeddings;
\item
  implement nearest-neighbor search;
\item
  test analogy arithmetic;
\item
  implement a small skip-gram negative-sampling model;
\item
  use Plotly to visualize embeddings;
\item
  use AI to critique mathematical assumptions and possible bias.
\end{enumerate}

\chapter{Summary}\label{summary}

Word embeddings transform discrete tokens into continuous vectors.
Mathematically, an embedding is a map

\[
e:V\to R^m.
\]

Count-based methods begin with co-occurrence matrices and transform them
using PMI, PPMI, and matrix factorization. Prediction-based methods
learn embeddings by optimizing objectives such as skip-gram softmax or
negative sampling. In both cases, the resulting vectors encode
statistical regularities from the corpus.

The central message is:

\begin{quote}
Word meaning is not inserted into the vectors by hand. It is induced
from patterns of co-occurrence, prediction, and optimization.
\end{quote}

This chapter completes the bridge from classical vector-space models to
modern neural NLP. The next chapter studies the geometry of embedding
spaces more systematically, including clustering, visualization,
anisotropy, and bias.

\chapter{Geometry of Embedding
Spaces}\label{geometry-of-embedding-spaces}

Similarity, dimension reduction, clustering, and high-dimensional
caution

\hfill\break

Word embeddings turn tokens into points in a vector space. Once language
objects become vectors, NLP inherits the language of \textbf{geometry}:
distances, angles, projections, subspaces, neighborhoods, clusters, and
directions. This chapter develops the geometry of embedding spaces
carefully, with emphasis on the mathematical assumptions behind
similarity search and visualization.

\section{Learning goals}\label{learning-goals-8}

After studying this chapter, you should be able to:

\begin{itemize}
\tightlist
\item
  Treat a collection of word or sentence embeddings as a point cloud
  \(X \in R^{n \times m}\).
\item
  Compare inner-product similarity, cosine similarity, and Euclidean
  distance.
\item
  Explain why nearest-neighbor lists depend on normalization and metric
  choice.
\item
  Derive PCA from the singular value decomposition of a centered
  embedding matrix.
\item
  Interpret low-dimensional embedding visualizations responsibly.
\item
  State the \(k\)-means objective and understand clustering as an
  optimization problem.
\item
  Describe anisotropy, hubness, and distance concentration in
  high-dimensional spaces.
\item
  Use linear directions to study interpretable attributes and possible
  bias in embeddings.
\end{itemize}

\section{9.1 Why geometry matters in
NLP}\label{why-geometry-matters-in-nlp}

In earlier chapters, a word embedding was introduced as a map

\[
E:V \longrightarrow R^m,
\]

where \(V\) is a vocabulary and \(m\) is the embedding dimension. For a
token \(w \in V\), the vector \(E(w)\) is often written as
\(e_w \in R^m\).

If \(V = \{w_1,\dots,w_n\}\), the embedding model produces a matrix

\[
X=
\begin{bmatrix}
- & e_{w_1}^T & -\\
- & e_{w_2}^T & -\\
& \vdots & \\
- & e_{w_n}^T & -
\end{bmatrix}
\in R^{n\times m}.
\]

The rows of \(X\) form a finite point cloud in \(R^m\). Many NLP
operations can then be written geometrically.

{\def\LTcaptype{none} % do not increment counter
\begin{longtable}[]{@{}
  >{\raggedright\arraybackslash}p{(\linewidth - 2\tabcolsep) * \real{0.5000}}
  >{\raggedright\arraybackslash}p{(\linewidth - 2\tabcolsep) * \real{0.5000}}@{}}
\toprule\noalign{}
\begin{minipage}[b]{\linewidth}\raggedright
NLP question
\end{minipage} & \begin{minipage}[b]{\linewidth}\raggedright
Geometric formulation
\end{minipage} \\
\midrule\noalign{}
\endhead
\bottomrule\noalign{}
\endlastfoot
Which words are similar to \texttt{king}? & Find nearest neighbors of
\(e_{king}\). \\
Which documents match a query? & Find vectors with large inner product
or cosine similarity to the query vector. \\
Are topics separated? & Look for clusters or low-dimensional
structure. \\
What direction represents an attribute? & Estimate a vector direction
\(b\) and project embeddings onto \(b\). \\
Can we visualize embeddings? & Project \(R^m\) to \(R^2\) using PCA,
SVD, or another map. \\
\end{longtable}
}

The important point is that geometry is not merely a visualization tool.
It is part of the model. A statement such as ``these words are close''
is incomplete until we specify the representation, the metric, and any
preprocessing such as centering or normalization.

\textbf{Core viewpoint.} An embedding space is a finite metric or
inner-product structure induced by a representation map. The geometry we
observe is not intrinsic to language alone; it depends on the
construction of \(E\), the training data, the loss function, and the
similarity measure.

\section{9.2 Similarity measures}\label{similarity-measures}

Let \(x,y \in R^m\) be two embedding vectors.

\subsection{Inner product}\label{inner-product}

The inner product is

\[
\langle x,y\rangle = x^T y = \sum_{j=1}^m x_j y_j.
\]

Large inner products favor vectors that point in similar directions and
also have large norms. Therefore inner-product retrieval is sensitive to
vector length.

\subsection{Euclidean distance}\label{euclidean-distance}

The Euclidean distance is

\[
d_2(x,y)=\|x-y\|_2=
\left(\sum_{j=1}^m (x_j-y_j)^2\right)^{1/2}.
\]

Small Euclidean distance means that the vectors are close as points in
\(R^m\). It uses both direction and length.

\subsection{Cosine similarity}\label{cosine-similarity-1}

The cosine similarity is

\[
\cos(x,y)=\frac{x^T y}{\|x\|_2\|y\|_2},
\qquad x\ne 0,\ y\ne 0.
\]

Cosine similarity depends only on the angle between \(x\) and \(y\). It
is invariant under positive rescaling:

\[
\cos(cx,dy)=\cos(x,y),
\qquad c>0,\ d>0.
\]

This is useful when vector length is less meaningful than direction. In
many word and document embedding applications, cosine similarity is
preferred for nearest-neighbor search.

\subsection{Relationship after
normalization}\label{relationship-after-normalization}

Define the normalized vectors

\[
\hat{x}=\frac{x}{\|x\|_2},
\qquad
\hat{y}=\frac{y}{\|y\|_2}.
\]

Then

\[
\|\hat{x}-\hat{y}\|_2^2
=
\|\hat{x}\|_2^2+
\|\hat{y}\|_2^2-2\hat{x}^T\hat{y}
=2-2\cos(x,y).
\]

Thus, for normalized embeddings, minimizing Euclidean distance is
equivalent to maximizing cosine similarity.

\textbf{Metric choice matters.} The nearest neighbors of a word can
change when we switch from inner product to cosine similarity or
Euclidean distance. This is not a numerical accident. It reflects
different mathematical assumptions about whether vector length carries
semantic information.

\section{9.3 Nearest neighbors as an optimization
problem}\label{nearest-neighbors-as-an-optimization-problem}

Given embeddings \(e_1,\dots,e_n \in R^m\) and a query vector
\(q \in R^m\), nearest-neighbor search solves

\[
\operatorname{NN}_s(q)
=
\arg\max_{1\le i\le n} s(q,e_i),
\]

where \(s\) is a similarity function such as inner product or cosine
similarity. Equivalently, using a distance \(d\),

\[
\operatorname{NN}_d(q)
=
\arg\min_{1\le i\le n} d(q,e_i).
\]

For top-\(k\) retrieval, we return the \(k\) largest scores or the \(k\)
smallest distances.

In matrix form, if \(X\in R^{n\times m}\) stores the candidate
embeddings as rows, inner-product scores are

\[
s = Xq \in R^n.
\]

For cosine similarity, define a row-normalized matrix \(\widehat X\) and
normalized query \(\hat q\). Then

\[
s = \widehat X \hat q.
\]

This simple matrix-vector multiplication is one reason embedding
retrieval is computationally attractive.

\section{9.4 A small geometric example}\label{a-small-geometric-example}

Suppose an embedding model gives the following two-dimensional vectors:

{\def\LTcaptype{none} % do not increment counter
\begin{longtable}[]{@{}lrr@{}}
\toprule\noalign{}
token & \(x_1\) & \(x_2\) \\
\midrule\noalign{}
\endhead
\bottomrule\noalign{}
\endlastfoot
\texttt{cat} & 1.0 & 0.9 \\
\texttt{dog} & 1.1 & 0.8 \\
\texttt{lion} & 1.4 & 1.1 \\
\texttt{car} & -1.0 & -0.8 \\
\texttt{truck} & -1.2 & -0.9 \\
\texttt{bus} & -0.9 & -1.1 \\
\texttt{king} & 0.2 & 1.8 \\
\texttt{queen} & 0.1 & 1.7 \\
\end{longtable}
}

This toy space has three visible regions: animals, vehicles, and
royalty. But even in this simple example, a nearest-neighbor query
depends on whether we use cosine similarity or Euclidean distance. If
two points have similar directions but different magnitudes, cosine and
Euclidean rankings can disagree.

The next interactive example lets you explore this dependence directly.

\section{9.5 Interactive: embedding geometry
explorer}\label{interactive-embedding-geometry-explorer}

The figure below uses a small hand-built embedding space. Choose a query
token and a similarity measure. The plot highlights the nearest
neighbors and shows how the ranking changes under cosine similarity,
inner product, and Euclidean distance.

\begin{tcolorbox}[enhanced jigsaw, arc=.35mm, bottomrule=.15mm, bottomtitle=1mm, breakable, colback=white, colbacktitle=quarto-callout-note-color!10!white, colframe=quarto-callout-note-color-frame, coltitle=black, left=2mm, leftrule=.75mm, opacityback=0, opacitybacktitle=0.6, rightrule=.15mm, title=\textcolor{quarto-callout-note-color}{\faInfo}\hspace{0.5em}{Note}, titlerule=0mm, toprule=.15mm, toptitle=1mm]

For PDF output, the interactive figure is represented by the
mathematical discussion above. In HTML output, the controls allow
students to compare nearest-neighbor rankings under different similarity
measures.

\end{tcolorbox}

\section{9.6 Python example: nearest neighbors from
scratch}\label{python-example-nearest-neighbors-from-scratch}

The following example implements nearest-neighbor search without a
machine learning library. This is useful because it makes the metric
choice explicit.

\protect\phantomsection\label{py-ch09-nearest-neighbors}
\begin{Shaded}
\begin{Highlighting}[]
\ImportTok{import}\NormalTok{ numpy }\ImportTok{as}\NormalTok{ np}

\NormalTok{words }\OperatorTok{=}\NormalTok{ [}\StringTok{"cat"}\NormalTok{, }\StringTok{"dog"}\NormalTok{, }\StringTok{"lion"}\NormalTok{, }\StringTok{"car"}\NormalTok{, }\StringTok{"truck"}\NormalTok{, }\StringTok{"bus"}\NormalTok{, }\StringTok{"king"}\NormalTok{, }\StringTok{"queen"}\NormalTok{]}
\NormalTok{X }\OperatorTok{=}\NormalTok{ np.array([}
\NormalTok{    [}\FloatTok{1.00}\NormalTok{,  }\FloatTok{0.90}\NormalTok{],}
\NormalTok{    [}\FloatTok{1.10}\NormalTok{,  }\FloatTok{0.80}\NormalTok{],}
\NormalTok{    [}\FloatTok{1.45}\NormalTok{,  }\FloatTok{1.10}\NormalTok{],}
\NormalTok{    [}\OperatorTok{{-}}\FloatTok{1.00}\NormalTok{, }\OperatorTok{{-}}\FloatTok{0.80}\NormalTok{],}
\NormalTok{    [}\OperatorTok{{-}}\FloatTok{1.25}\NormalTok{, }\OperatorTok{{-}}\FloatTok{0.90}\NormalTok{],}
\NormalTok{    [}\OperatorTok{{-}}\FloatTok{0.90}\NormalTok{, }\OperatorTok{{-}}\FloatTok{1.15}\NormalTok{],}
\NormalTok{    [}\FloatTok{0.20}\NormalTok{,  }\FloatTok{1.85}\NormalTok{],}
\NormalTok{    [}\FloatTok{0.05}\NormalTok{,  }\FloatTok{1.70}\NormalTok{],}
\NormalTok{])}

\KeywordTok{def}\NormalTok{ cosine\_scores(X, q):}
\NormalTok{    X\_norm }\OperatorTok{=}\NormalTok{ X }\OperatorTok{/}\NormalTok{ np.linalg.norm(X, axis}\OperatorTok{=}\DecValTok{1}\NormalTok{, keepdims}\OperatorTok{=}\VariableTok{True}\NormalTok{)}
\NormalTok{    q\_norm }\OperatorTok{=}\NormalTok{ q }\OperatorTok{/}\NormalTok{ np.linalg.norm(q)}
    \ControlFlowTok{return}\NormalTok{ X\_norm }\OperatorTok{@}\NormalTok{ q\_norm}

\KeywordTok{def}\NormalTok{ euclidean\_distances(X, q):}
    \ControlFlowTok{return}\NormalTok{ np.linalg.norm(X }\OperatorTok{{-}}\NormalTok{ q, axis}\OperatorTok{=}\DecValTok{1}\NormalTok{)}

\KeywordTok{def}\NormalTok{ nearest\_neighbors(query\_word, metric}\OperatorTok{=}\StringTok{"cosine"}\NormalTok{, k}\OperatorTok{=}\DecValTok{3}\NormalTok{):}
\NormalTok{    idx }\OperatorTok{=}\NormalTok{ words.index(query\_word)}
\NormalTok{    q }\OperatorTok{=}\NormalTok{ X[idx]}

    \ControlFlowTok{if}\NormalTok{ metric }\OperatorTok{==} \StringTok{"cosine"}\NormalTok{:}
\NormalTok{        values }\OperatorTok{=}\NormalTok{ cosine\_scores(X, q)}
\NormalTok{        order }\OperatorTok{=}\NormalTok{ np.argsort(}\OperatorTok{{-}}\NormalTok{values)}
    \ControlFlowTok{elif}\NormalTok{ metric }\OperatorTok{==} \StringTok{"inner"}\NormalTok{:}
\NormalTok{        values }\OperatorTok{=}\NormalTok{ X }\OperatorTok{@}\NormalTok{ q}
\NormalTok{        order }\OperatorTok{=}\NormalTok{ np.argsort(}\OperatorTok{{-}}\NormalTok{values)}
    \ControlFlowTok{elif}\NormalTok{ metric }\OperatorTok{==} \StringTok{"euclidean"}\NormalTok{:}
\NormalTok{        values }\OperatorTok{=}\NormalTok{ euclidean\_distances(X, q)}
\NormalTok{        order }\OperatorTok{=}\NormalTok{ np.argsort(values)}
    \ControlFlowTok{else}\NormalTok{:}
        \ControlFlowTok{raise} \PreprocessorTok{ValueError}\NormalTok{(}\StringTok{"metric must be cosine, inner, or euclidean"}\NormalTok{)}

\NormalTok{    result }\OperatorTok{=}\NormalTok{ []}
    \ControlFlowTok{for}\NormalTok{ j }\KeywordTok{in}\NormalTok{ order:}
        \ControlFlowTok{if}\NormalTok{ j }\OperatorTok{!=}\NormalTok{ idx:}
\NormalTok{            result.append((words[j], }\BuiltInTok{float}\NormalTok{(values[j])))}
        \ControlFlowTok{if} \BuiltInTok{len}\NormalTok{(result) }\OperatorTok{==}\NormalTok{ k:}
            \ControlFlowTok{break}
    \ControlFlowTok{return}\NormalTok{ result}

\BuiltInTok{print}\NormalTok{(}\StringTok{"cosine:"}\NormalTok{, nearest\_neighbors(}\StringTok{"cat"}\NormalTok{, metric}\OperatorTok{=}\StringTok{"cosine"}\NormalTok{))}
\BuiltInTok{print}\NormalTok{(}\StringTok{"inner product:"}\NormalTok{, nearest\_neighbors(}\StringTok{"cat"}\NormalTok{, metric}\OperatorTok{=}\StringTok{"inner"}\NormalTok{))}
\BuiltInTok{print}\NormalTok{(}\StringTok{"euclidean:"}\NormalTok{, nearest\_neighbors(}\StringTok{"cat"}\NormalTok{, metric}\OperatorTok{=}\StringTok{"euclidean"}\NormalTok{))}
\end{Highlighting}
\end{Shaded}

\begin{verbatim}
cosine: [('lion', 0.9964891983367161), ('dog', 0.9945949192045226), ('king', 0.7449800754669957)]
inner product: [('lion', 2.44), ('king', 1.8650000000000002), ('dog', 1.82)]
euclidean: [('dog', 0.14142135623730953), ('lion', 0.49244289008980524), ('king', 1.241974234837422)]
\end{verbatim}

The same vectors can produce different neighbor rankings. In a real
embedding model, such differences can affect search results, clustering
results, and qualitative interpretations.

\section{9.7 Centering, normalization, and
anisotropy}\label{centering-normalization-and-anisotropy}

Embedding spaces often have strong global directions. For example, many
vectors may share a large common component. This can make most pairwise
cosine similarities positive and reduce the contrast between meaningful
neighbors and unrelated words.

Let the rows of \(X\in R^{n\times m}\) be embeddings. The sample mean
vector is

\[
\bar{x}=\frac{1}{n}\sum_{i=1}^n x_i.
\]

The centered embedding matrix is

\[
X_c = X - {\bf 1}\bar{x}^T,
\]

where \({\bf 1}\in R^n\) is the vector of all ones. Centering removes
the global mean direction.

Row normalization gives

\[
\tilde{x}_i=\frac{x_i}{\|x_i\|_2}.
\]

Depending on the task, one may use raw vectors, centered vectors,
normalized vectors, or centered-and-normalized vectors. These choices
change the geometry.

\textbf{Anisotropy.} A point cloud is anisotropic when its variance is
much larger in some directions than in others. In embedding spaces,
anisotropy can cause a few dominant directions to affect many similarity
computations.

A simple diagnostic is to inspect the singular values of the centered
matrix. If the first few singular values dominate, then much of the
variation lies in a low-dimensional subspace.

\section{9.8 PCA for embedding spaces}\label{pca-for-embedding-spaces}

PCA is a standard method for projecting high-dimensional embeddings to a
lower-dimensional space. It should be understood as a
variance-preserving linear projection, not as a perfect map of semantic
meaning.

Let \(X\in R^{n\times m}\) be the embedding matrix, and let

\[
X_c = X - {\bf 1}\bar{x}^T
\]

be the centered matrix. The sample covariance matrix is

\[
S=\frac{1}{n-1}X_c^T X_c \in R^{m\times m}.
\]

PCA finds orthonormal directions \(v_1,\dots,v_r\) solving

\[
v_1 = \arg\max_{\|v\|_2=1} v^T S v,
\]

and, for \(j>1\),

\[
v_j = \arg\max_{\|v\|_2=1,\ v\perp v_1,\dots,v_{j-1}} v^T S v.
\]

The solutions are eigenvectors of \(S\):

\[
S v_j = \lambda_j v_j,
\qquad
\lambda_1\ge \lambda_2\ge \cdots \ge 0.
\]

Equivalently, use the singular value decomposition

\[
X_c = U\Sigma V^T.
\]

The columns of \(V\) are the principal directions, and the
low-dimensional PCA coordinates are

\[
Z_k = X_c V_k \in R^{n\times k},
\]

where \(V_k=[v_1,\dots,v_k]\).

The fraction of variance explained by the first \(k\) components is

\[
\frac{\lambda_1+\cdots+\lambda_k}{\lambda_1+\cdots+\lambda_m}
=
\frac{\sigma_1^2+\cdots+\sigma_k^2}{\sigma_1^2+\cdots+\sigma_r^2}.
\]

\subsection{PCA is useful but limited}\label{pca-is-useful-but-limited}

PCA preserves large-variance linear structure. It does not necessarily
preserve all nearest-neighbor relationships, clusters, or nonlinear
geometry. A two-dimensional PCA plot is a useful diagnostic, but it
should not be treated as the embedding space itself.

\section{9.9 Python example: PCA by
SVD}\label{python-example-pca-by-svd}

\protect\phantomsection\label{py-ch09-pca-svd}
\begin{Shaded}
\begin{Highlighting}[]
\ImportTok{import}\NormalTok{ numpy }\ImportTok{as}\NormalTok{ np}

\CommentTok{\# Center the toy embedding matrix.}
\NormalTok{X\_centered }\OperatorTok{=}\NormalTok{ X }\OperatorTok{{-}}\NormalTok{ X.mean(axis}\OperatorTok{=}\DecValTok{0}\NormalTok{, keepdims}\OperatorTok{=}\VariableTok{True}\NormalTok{)}

\CommentTok{\# Singular value decomposition.}
\NormalTok{U, singular\_values, Vt }\OperatorTok{=}\NormalTok{ np.linalg.svd(X\_centered, full\_matrices}\OperatorTok{=}\VariableTok{False}\NormalTok{)}

\CommentTok{\# Two{-}dimensional PCA coordinates.}
\NormalTok{Z }\OperatorTok{=}\NormalTok{ X\_centered }\OperatorTok{@}\NormalTok{ Vt[:}\DecValTok{2}\NormalTok{].T}

\BuiltInTok{print}\NormalTok{(}\StringTok{"Singular values:"}\NormalTok{, np.}\BuiltInTok{round}\NormalTok{(singular\_values, }\DecValTok{3}\NormalTok{))}
\BuiltInTok{print}\NormalTok{(}\StringTok{"PCA coordinates:"}\NormalTok{)}
\ControlFlowTok{for}\NormalTok{ word, z }\KeywordTok{in} \BuiltInTok{zip}\NormalTok{(words, Z):}
    \BuiltInTok{print}\NormalTok{(}\SpecialStringTok{f"}\SpecialCharTok{\{}\NormalTok{word}\SpecialCharTok{:\textgreater{}6s\}}\SpecialStringTok{: (}\SpecialCharTok{\{}\NormalTok{z[}\DecValTok{0}\NormalTok{]}\SpecialCharTok{: .3f\}}\SpecialStringTok{, }\SpecialCharTok{\{}\NormalTok{z[}\DecValTok{1}\NormalTok{]}\SpecialCharTok{: .3f\}}\SpecialStringTok{)"}\NormalTok{)}
\end{Highlighting}
\end{Shaded}

\begin{verbatim}
Singular values: [3.946 1.522]
PCA coordinates:
   cat: (-0.941, -0.415)
   dog: (-0.928, -0.556)
  lion: (-1.382, -0.635)
   car: ( 1.643,  0.048)
 truck: ( 1.879,  0.177)
   bus: ( 1.849, -0.252)
  king: (-1.166,  0.806)
 queen: (-0.955,  0.827)
\end{verbatim}

For this toy example, the original vectors are already two-dimensional.
In a real embedding model, \(m\) might be 100, 300, 768, or larger. PCA
gives a low-dimensional view of this higher-dimensional structure.

\section{9.10 Clustering embeddings}\label{clustering-embeddings}

Clustering asks whether embedding vectors can be partitioned into groups
with similar geometry. A common model is \(k\)-means clustering.

Given embeddings \(x_1,\dots,x_n\in R^m\), the \(k\)-means objective is

\[
\min_{C_1,\dots,C_k,\ \mu_1,\dots,\mu_k}
\sum_{r=1}^k \sum_{i\in C_r} \|x_i-\mu_r\|_2^2,
\]

where \(C_1,\dots,C_k\) partition \(\{1,
\dots,n\}\) and \(\mu_r\in R^m\) is the centroid of cluster \(r\).

For fixed clusters, the optimal centroid is the sample mean:

\[
\mu_r = \frac{1}{|C_r|}\sum_{i\in C_r} x_i.
\]

The standard \(k\)-means algorithm alternates between two steps:

\begin{enumerate}
\def\labelenumi{\arabic{enumi}.}
\tightlist
\item
  Assign each point to the nearest centroid.
\item
  Replace each centroid by the mean of the points assigned to it.
\end{enumerate}

This decreases the objective at every iteration, but it can converge to
a local minimum.

\subsection{Spherical clustering}\label{spherical-clustering}

If cosine similarity is the desired geometry, it is common to normalize
vectors before clustering. Applying \(k\)-means to normalized vectors is
related to spherical clustering, where direction matters more than
length.

\section{\texorpdfstring{9.11 Python example: small \(k\)-means
implementation}{9.11 Python example: small k-means implementation}}\label{python-example-small-k-means-implementation}

\protect\phantomsection\label{py-ch09-kmeans-from-scratch}
\begin{Shaded}
\begin{Highlighting}[]
\ImportTok{import}\NormalTok{ numpy }\ImportTok{as}\NormalTok{ np}

\NormalTok{rng }\OperatorTok{=}\NormalTok{ np.random.default\_rng(}\DecValTok{7243}\NormalTok{)}

\KeywordTok{def}\NormalTok{ kmeans(X, k, n\_iter}\OperatorTok{=}\DecValTok{10}\NormalTok{):}
    \CommentTok{\# Choose initial centroids from the data.}
\NormalTok{    initial }\OperatorTok{=}\NormalTok{ rng.choice(X.shape[}\DecValTok{0}\NormalTok{], size}\OperatorTok{=}\NormalTok{k, replace}\OperatorTok{=}\VariableTok{False}\NormalTok{)}
\NormalTok{    centroids }\OperatorTok{=}\NormalTok{ X[initial].copy()}

    \ControlFlowTok{for}\NormalTok{ \_ }\KeywordTok{in} \BuiltInTok{range}\NormalTok{(n\_iter):}
\NormalTok{        distances }\OperatorTok{=}\NormalTok{ np.linalg.norm(X[:, }\VariableTok{None}\NormalTok{, :] }\OperatorTok{{-}}\NormalTok{ centroids[}\VariableTok{None}\NormalTok{, :, :], axis}\OperatorTok{=}\DecValTok{2}\NormalTok{)}
\NormalTok{        labels }\OperatorTok{=}\NormalTok{ np.argmin(distances, axis}\OperatorTok{=}\DecValTok{1}\NormalTok{)}

\NormalTok{        new\_centroids }\OperatorTok{=}\NormalTok{ centroids.copy()}
        \ControlFlowTok{for}\NormalTok{ r }\KeywordTok{in} \BuiltInTok{range}\NormalTok{(k):}
            \ControlFlowTok{if}\NormalTok{ np.}\BuiltInTok{any}\NormalTok{(labels }\OperatorTok{==}\NormalTok{ r):}
\NormalTok{                new\_centroids[r] }\OperatorTok{=}\NormalTok{ X[labels }\OperatorTok{==}\NormalTok{ r].mean(axis}\OperatorTok{=}\DecValTok{0}\NormalTok{)}
\NormalTok{        centroids }\OperatorTok{=}\NormalTok{ new\_centroids}

    \ControlFlowTok{return}\NormalTok{ labels, centroids}

\NormalTok{labels, centroids }\OperatorTok{=}\NormalTok{ kmeans(X, k}\OperatorTok{=}\DecValTok{3}\NormalTok{, n\_iter}\OperatorTok{=}\DecValTok{20}\NormalTok{)}

\ControlFlowTok{for}\NormalTok{ r }\KeywordTok{in} \BuiltInTok{range}\NormalTok{(}\DecValTok{3}\NormalTok{):}
\NormalTok{    cluster\_words }\OperatorTok{=}\NormalTok{ [w }\ControlFlowTok{for}\NormalTok{ w, lab }\KeywordTok{in} \BuiltInTok{zip}\NormalTok{(words, labels) }\ControlFlowTok{if}\NormalTok{ lab }\OperatorTok{==}\NormalTok{ r]}
    \BuiltInTok{print}\NormalTok{(}\SpecialStringTok{f"cluster }\SpecialCharTok{\{}\NormalTok{r}\SpecialCharTok{\}}\SpecialStringTok{:"}\NormalTok{, cluster\_words)}
\end{Highlighting}
\end{Shaded}

\begin{verbatim}
cluster 0: ['cat', 'dog', 'lion', 'king', 'queen']
cluster 1: ['car', 'truck']
cluster 2: ['bus']
\end{verbatim}

This example is intentionally small. The goal is not to obtain a perfect
clustering, but to see the mathematical objective behind the algorithm.

\section{9.12 High-dimensional caution: distance
concentration}\label{high-dimensional-caution-distance-concentration}

Embedding spaces used in NLP are often high-dimensional.
High-dimensional geometry can differ sharply from low-dimensional
intuition.

Suppose \(x,y\in R^m\) have independent standard normal coordinates.
Then

\[
\|x-y\|_2^2 = \sum_{j=1}^m (x_j-y_j)^2.
\]

Since \(x_j-y_j\sim N(0,2)\), the expected squared distance is

\[
E\|x-y\|_2^2 = 2m.
\]

As \(m\) grows, distances between random points become increasingly
concentrated around their mean. This can make nearest-neighbor structure
unstable unless the embedding space has meaningful low-dimensional or
clustered structure.

Another related issue is \textbf{hubness}: some points appear as nearest
neighbors of many other points. In embedding spaces, hubs may arise
because of frequency effects, anisotropy, or dominant common directions.

\section{9.13 Python example: distance
concentration}\label{python-example-distance-concentration}

\protect\phantomsection\label{py-ch09-distance-concentration}
\begin{Shaded}
\begin{Highlighting}[]
\ImportTok{import}\NormalTok{ numpy }\ImportTok{as}\NormalTok{ np}

\NormalTok{rng }\OperatorTok{=}\NormalTok{ np.random.default\_rng(}\DecValTok{9}\NormalTok{)}

\KeywordTok{def}\NormalTok{ distance\_summary(dim, n\_pairs}\OperatorTok{=}\DecValTok{5000}\NormalTok{):}
\NormalTok{    A }\OperatorTok{=}\NormalTok{ rng.normal(size}\OperatorTok{=}\NormalTok{(n\_pairs, dim))}
\NormalTok{    B }\OperatorTok{=}\NormalTok{ rng.normal(size}\OperatorTok{=}\NormalTok{(n\_pairs, dim))}
\NormalTok{    distances }\OperatorTok{=}\NormalTok{ np.linalg.norm(A }\OperatorTok{{-}}\NormalTok{ B, axis}\OperatorTok{=}\DecValTok{1}\NormalTok{)}
    \ControlFlowTok{return}\NormalTok{ distances.mean(), distances.std(), distances.std() }\OperatorTok{/}\NormalTok{ distances.mean()}

\ControlFlowTok{for}\NormalTok{ dim }\KeywordTok{in}\NormalTok{ [}\DecValTok{2}\NormalTok{, }\DecValTok{10}\NormalTok{, }\DecValTok{50}\NormalTok{, }\DecValTok{100}\NormalTok{, }\DecValTok{300}\NormalTok{, }\DecValTok{768}\NormalTok{]:}
\NormalTok{    mean\_d, std\_d, cv }\OperatorTok{=}\NormalTok{ distance\_summary(dim)}
    \BuiltInTok{print}\NormalTok{(}\SpecialStringTok{f"dim=}\SpecialCharTok{\{}\NormalTok{dim}\SpecialCharTok{:\textgreater{}3d\}}\SpecialStringTok{: mean=}\SpecialCharTok{\{}\NormalTok{mean\_d}\SpecialCharTok{:.3f\}}\SpecialStringTok{, std=}\SpecialCharTok{\{}\NormalTok{std\_d}\SpecialCharTok{:.3f\}}\SpecialStringTok{, coefficient of variation=}\SpecialCharTok{\{}\NormalTok{cv}\SpecialCharTok{:.3f\}}\SpecialStringTok{"}\NormalTok{)}
\end{Highlighting}
\end{Shaded}

\begin{verbatim}
dim=  2: mean=1.788, std=0.928, coefficient of variation=0.519
dim= 10: mean=4.367, std=0.987, coefficient of variation=0.226
dim= 50: mean=9.941, std=1.005, coefficient of variation=0.101
dim=100: mean=14.122, std=0.989, coefficient of variation=0.070
dim=300: mean=24.472, std=0.998, coefficient of variation=0.041
dim=768: mean=39.178, std=0.999, coefficient of variation=0.025
\end{verbatim}

The coefficient of variation decreases as dimension increases. This does
not mean all high-dimensional embeddings are useless. It means that
meaningful structure must come from learned correlations, training
objectives, and data geometry rather than from random high-dimensional
distance alone.

\section{9.14 Semantic directions and analogy
geometry}\label{semantic-directions-and-analogy-geometry}

Embedding spaces often encode attributes as approximate directions. For
example, the analogy pattern

\[
e_{king}-e_{man}+e_{woman}\approx e_{queen}
\]

suggests that the difference vector

\[
b = e_{woman}-e_{man}
\]

may represent a gender-related direction in that embedding space.

More generally, if we have paired examples \((a_i,b_i)\) representing an
attribute change, we can estimate a direction

\[
d = \frac{1}{r}\sum_{i=1}^r (e_{b_i}-e_{a_i}).
\]

Then the scalar projection of a word vector \(x\) onto this direction is

\[
\operatorname{score}_d(x)=\frac{x^T d}{\|d\|_2}.
\]

This score measures how far \(x\) lies along the direction \(d\). Such
directions can be useful for interpretation, but they must be used
carefully: a direction learned from a small set of examples may reflect
social bias, corpus artifacts, or confounding variables.

\textbf{Bias as geometry.} If an embedding space encodes an undesirable
association as a direction or subspace, then downstream models that use
the embeddings may inherit that association. Geometric analysis can
reveal such structure, but it does not automatically solve the ethical
or statistical problem.

\section{9.15 Python example: a toy attribute
direction}\label{python-example-a-toy-attribute-direction}

The example below constructs a toy \textbf{singular-to-plural} direction
and projects each word onto it. This avoids making a social claim and
lets us focus on the linear algebra.

\protect\phantomsection\label{py-ch09-attribute-direction}
\begin{Shaded}
\begin{Highlighting}[]
\ImportTok{import}\NormalTok{ numpy }\ImportTok{as}\NormalTok{ np}

\NormalTok{morph\_words }\OperatorTok{=}\NormalTok{ [}\StringTok{"cat"}\NormalTok{, }\StringTok{"cats"}\NormalTok{, }\StringTok{"dog"}\NormalTok{, }\StringTok{"dogs"}\NormalTok{, }\StringTok{"car"}\NormalTok{, }\StringTok{"cars"}\NormalTok{, }\StringTok{"book"}\NormalTok{, }\StringTok{"books"}\NormalTok{, }\StringTok{"city"}\NormalTok{, }\StringTok{"cities"}\NormalTok{]}
\NormalTok{Y }\OperatorTok{=}\NormalTok{ np.array([}
\NormalTok{    [}\FloatTok{1.00}\NormalTok{, }\FloatTok{0.20}\NormalTok{, }\FloatTok{0.00}\NormalTok{],   }\CommentTok{\# cat}
\NormalTok{    [}\FloatTok{1.05}\NormalTok{, }\FloatTok{0.20}\NormalTok{, }\FloatTok{0.90}\NormalTok{],   }\CommentTok{\# cats}
\NormalTok{    [}\FloatTok{0.90}\NormalTok{, }\FloatTok{0.30}\NormalTok{, }\FloatTok{0.05}\NormalTok{],   }\CommentTok{\# dog}
\NormalTok{    [}\FloatTok{0.95}\NormalTok{, }\FloatTok{0.30}\NormalTok{, }\FloatTok{0.95}\NormalTok{],   }\CommentTok{\# dogs}
\NormalTok{    [}\OperatorTok{{-}}\FloatTok{0.80}\NormalTok{, }\FloatTok{0.10}\NormalTok{, }\FloatTok{0.00}\NormalTok{],  }\CommentTok{\# car}
\NormalTok{    [}\OperatorTok{{-}}\FloatTok{0.75}\NormalTok{, }\FloatTok{0.10}\NormalTok{, }\FloatTok{0.85}\NormalTok{],  }\CommentTok{\# cars}
\NormalTok{    [}\FloatTok{0.10}\NormalTok{, }\OperatorTok{{-}}\FloatTok{0.90}\NormalTok{, }\FloatTok{0.10}\NormalTok{],  }\CommentTok{\# book}
\NormalTok{    [}\FloatTok{0.15}\NormalTok{, }\OperatorTok{{-}}\FloatTok{0.90}\NormalTok{, }\FloatTok{0.95}\NormalTok{],  }\CommentTok{\# books}
\NormalTok{    [}\OperatorTok{{-}}\FloatTok{0.20}\NormalTok{, }\OperatorTok{{-}}\FloatTok{0.70}\NormalTok{, }\FloatTok{0.05}\NormalTok{], }\CommentTok{\# city}
\NormalTok{    [}\OperatorTok{{-}}\FloatTok{0.15}\NormalTok{, }\OperatorTok{{-}}\FloatTok{0.70}\NormalTok{, }\FloatTok{0.90}\NormalTok{], }\CommentTok{\# cities}
\NormalTok{])}

\NormalTok{pairs }\OperatorTok{=}\NormalTok{ [(}\StringTok{"cat"}\NormalTok{, }\StringTok{"cats"}\NormalTok{), (}\StringTok{"dog"}\NormalTok{, }\StringTok{"dogs"}\NormalTok{), (}\StringTok{"car"}\NormalTok{, }\StringTok{"cars"}\NormalTok{), (}\StringTok{"book"}\NormalTok{, }\StringTok{"books"}\NormalTok{)]}
\NormalTok{plural\_direction }\OperatorTok{=}\NormalTok{ np.zeros(Y.shape[}\DecValTok{1}\NormalTok{])}
\ControlFlowTok{for}\NormalTok{ singular, plural }\KeywordTok{in}\NormalTok{ pairs:}
\NormalTok{    plural\_direction }\OperatorTok{+=}\NormalTok{ Y[morph\_words.index(plural)] }\OperatorTok{{-}}\NormalTok{ Y[morph\_words.index(singular)]}
\NormalTok{plural\_direction }\OperatorTok{/=} \BuiltInTok{len}\NormalTok{(pairs)}
\NormalTok{plural\_direction }\OperatorTok{/=}\NormalTok{ np.linalg.norm(plural\_direction)}

\NormalTok{scores }\OperatorTok{=}\NormalTok{ Y }\OperatorTok{@}\NormalTok{ plural\_direction}
\NormalTok{order }\OperatorTok{=}\NormalTok{ np.argsort(}\OperatorTok{{-}}\NormalTok{scores)}

\BuiltInTok{print}\NormalTok{(}\StringTok{"Projection onto the toy singular{-}to{-}plural direction:"}\NormalTok{)}
\ControlFlowTok{for}\NormalTok{ j }\KeywordTok{in}\NormalTok{ order:}
    \BuiltInTok{print}\NormalTok{(}\SpecialStringTok{f"}\SpecialCharTok{\{}\NormalTok{morph\_words[j]}\SpecialCharTok{:\textgreater{}6s\}}\SpecialStringTok{: }\SpecialCharTok{\{}\NormalTok{scores[j]}\SpecialCharTok{: .3f\}}\SpecialStringTok{"}\NormalTok{)}
\end{Highlighting}
\end{Shaded}

\begin{verbatim}
Projection onto the toy singular-to-plural direction:
  dogs:  1.003
  cats:  0.958
 books:  0.957
cities:  0.890
  cars:  0.806
  book:  0.106
   dog:  0.101
   cat:  0.057
  city:  0.039
   car: -0.046
\end{verbatim}

In real applications, one should estimate directions from multiple
pairs, test sensitivity to the choice of pairs, and evaluate whether the
direction is stable across corpora and model families.

\section{9.16 Debiasing as projection: a mathematical
sketch}\label{debiasing-as-projection-a-mathematical-sketch}

A simple geometric debiasing operation removes the component of a vector
along a direction \(b\ne 0\).

The projection of \(x\) onto the span of \(b\) is

\[
\operatorname{proj}_b(x)=\frac{x^T b}{b^T b}b.
\]

The component orthogonal to \(b\) is

\[
x_{\perp}=x-\operatorname{proj}_b(x).
\]

This operation guarantees

\[
x_{\perp}^T b = 0.
\]

However, this does not guarantee that all bias has been removed. Bias
may live in a higher-dimensional subspace, in nonlinear structure, or in
downstream training data. Projection is a useful mathematical tool, but
it is not a complete social or statistical solution.

\section{9.17 Static versus contextual embedding
geometry}\label{static-versus-contextual-embedding-geometry}

The previous chapter focused mostly on static embeddings, where each
word type has one vector:

\[
E:V\to R^m.
\]

Modern transformer models produce contextual embeddings. A token's
vector depends on the full input sequence:

\[
H_t = F_t(w_1,\dots,w_T) \in R^m.
\]

Thus the same word can have different vectors in different contexts. For
example, \texttt{bank} in ``river bank'' and \texttt{bank} in ``savings
bank'' may have different contextual embeddings.

The geometry of contextual embeddings is therefore a geometry of token
occurrences rather than only word types. The data matrix may contain
many rows for the same word in different contexts:

\[
X =
\begin{bmatrix}
- & H_{t_1}^{(1)T} & -\\
- & H_{t_2}^{(2)T} & -\\
& \vdots & \\
- & H_{t_n}^{(n)T} & -
\end{bmatrix}.
\]

This shift is important. Static embedding geometry asks: ``Where is the
word type?'' Contextual embedding geometry asks: ``Where is this
occurrence of the word in this sentence?''

\section{9.18 What can go wrong in embedding
visualization?}\label{what-can-go-wrong-in-embedding-visualization}

Low-dimensional plots are useful, but they can be misleading. Common
issues include:

\begin{itemize}
\tightlist
\item
  \textbf{Projection distortion:} points close in \(R^2\) may not be
  close in \(R^m\).
\item
  \textbf{Lost variance:} a two-dimensional PCA plot may explain only a
  small fraction of total variance.
\item
  \textbf{Metric mismatch:} visualization may use Euclidean distance
  while retrieval uses cosine similarity.
\item
  \textbf{Crowding:} many points in high dimension must be packed into a
  small two-dimensional region.
\item
  \textbf{Overinterpretation:} visible clusters may be artifacts of
  projection, preprocessing, or selected examples.
\end{itemize}

A responsible visualization should report the preprocessing, projection
method, metric, and explained variance when applicable.

\section{9.19 AI-assisted learning
component}\label{ai-assisted-learning-component-3}

Use an AI assistant as a mathematical critic, not as an authority. The
goal is to force precise statements.

\subsection{Prompt 1: metric audit}\label{prompt-1-metric-audit}

\begin{quote}
I have a word embedding matrix \(X\in R^{n\times m}\). I used cosine
similarity to find nearest neighbors. Ask me five mathematical questions
that check whether this choice is appropriate. Focus on normalization,
vector norms, anisotropy, and downstream task goals.
\end{quote}

\subsection{Prompt 2: PCA interpretation
critique}\label{prompt-2-pca-interpretation-critique}

\begin{quote}
I plotted word embeddings using the first two principal components. Give
a careful explanation of what this plot preserves and what it may
distort. Include the formulas for centering, covariance, and explained
variance.
\end{quote}

\subsection{Prompt 3: bias direction
analysis}\label{prompt-3-bias-direction-analysis}

\begin{quote}
Suppose I define a direction \(d\) by averaging difference vectors
\(e_{b_i}-e_{a_i}\). Explain how to test whether this direction is
stable and how projection onto this direction should be interpreted
statistically.
\end{quote}

\subsection{Prompt 4: nearest-neighbor
debugging}\label{prompt-4-nearest-neighbor-debugging}

\begin{quote}
My embedding nearest-neighbor results are dominated by frequent or
generic words. Give possible geometric explanations and suggest
mathematical diagnostics.
\end{quote}

After using an AI assistant, verify all formulas and test claims with a
small numerical example.

\section{9.20 Worked calculation}\label{worked-calculation}

Let

\[
x=(2,1),
\qquad
 y=(4,2),
\qquad
 z=(1,2).
\]

Then

\[
\cos(x,y)=\frac{2\cdot 4+1\cdot 2}{\sqrt{2^2+1^2}\sqrt{4^2+2^2}}
=\frac{10}{\sqrt{5}\sqrt{20}}
=1.
\]

So \(x\) and \(y\) have exactly the same direction. Their Euclidean
distance is

\[
\|x-y\|_2=\sqrt{(2-4)^2+(1-2)^2}=\sqrt{5}.
\]

For \(z\),

\[
\cos(x,z)=\frac{2\cdot 1+1\cdot 2}{\sqrt{5}\sqrt{5}}=\frac{4}{5}=0.8,
\]

while

\[
\|x-z\|_2=\sqrt{(2-1)^2+(1-2)^2}=\sqrt{2}.
\]

Thus \(y\) is more similar to \(x\) by cosine similarity, but \(z\) is
closer to \(x\) by Euclidean distance. This illustrates why metric
choice changes nearest-neighbor rankings.

\section{9.21 Mini-project: metric sensitivity of
embeddings}\label{mini-project-metric-sensitivity-of-embeddings}

Choose a small embedding set. You may use word embeddings, sentence
embeddings, or vectors constructed from TF-IDF and SVD.

\begin{enumerate}
\def\labelenumi{\arabic{enumi}.}
\tightlist
\item
  Select 20 query words or sentences.
\item
  Compute top-5 neighbors under cosine similarity, inner product, and
  Euclidean distance.
\item
  Compare the overlap between the three neighbor lists.
\item
  Compute vector norms and examine whether high-norm vectors dominate
  inner-product retrieval.
\item
  Center the embeddings and repeat the analysis.
\item
  Write a short mathematical report explaining how preprocessing changed
  the geometry.
\end{enumerate}

A useful overlap statistic for two top-\(k\) neighbor sets \(A\) and
\(B\) is the Jaccard index

\[
J(A,B)=\frac{|A\cap B|}{|A\cup B|}.
\]

\section{9.22 Exercises}\label{exercises-8}

\begin{enumerate}
\def\labelenumi{\arabic{enumi}.}
\item
  \textbf{Cosine and Euclidean distance.} Prove that for nonzero
  \(x,y\in R^m\), \[
  \left\|\frac{x}{\|x\|_2}-\frac{y}{\|y\|_2}\right\|_2^2
  =2-2\cos(x,y).
  \]
\item
  \textbf{Nearest-neighbor matrix form.} Let \(X\in R^{n\times m}\)
  contain row-normalized embeddings. Show that the vector \(Xq\) gives
  cosine similarities to a normalized query \(q\).
\item
  \textbf{PCA derivation.} Let \(S\) be a symmetric positive
  semidefinite covariance matrix. Show that maximizing \(v^T S v\)
  subject to \(\|v\|_2=1\) leads to the eigenvalue equation
  \(Sv=\lambda v\).
\item
  \textbf{Explained variance.} Given singular values
  \(\sigma_1\ge\cdots\ge\sigma_r>0\), write the fraction of variance
  explained by the first \(k\) principal components.
\item
  \textbf{\(k\)-means centroid.} For a fixed cluster \(C\), prove that
  \[
  \mu=\frac{1}{|C|}\sum_{i\in C}x_i
  \] minimizes \(\sum_{i\in C}\|x_i-\mu\|_2^2\).
\item
  \textbf{Projection removal.} Let \(b\ne 0\) and define \[
  x_{\perp}=x-\frac{x^T b}{b^T b}b.
  \] Prove that \(x_{\perp}^T b=0\).
\item
  \textbf{Distance concentration simulation.} Simulate random Gaussian
  points in dimensions \(m=2,10,100,1000\). Plot the distribution of
  pairwise distances. Describe what changes as \(m\) increases.
\item
  \textbf{Contextual embeddings.} Explain why the geometry of contextual
  embeddings is a geometry of token occurrences rather than word types.
\item
  \textbf{AI critique.} Ask an AI assistant to explain why a
  two-dimensional embedding plot can be misleading. Identify at least
  two statements that need mathematical qualification.
\end{enumerate}

\section{9.23 Python lab connection}\label{python-lab-connection-3}

The associated independent-study lab is planned as:

\begin{Shaded}
\begin{Highlighting}[]
\NormalTok{labs/chapter{-}09{-}geometry{-}of{-}embeddings.ipynb}
\end{Highlighting}
\end{Shaded}

The lab should include:

\begin{itemize}
\tightlist
\item
  nearest-neighbor search under several metrics;
\item
  row normalization and centering;
\item
  PCA by SVD;
\item
  \(k\)-means clustering from scratch;
\item
  distance concentration simulations;
\item
  a short audit of semantic directions and projection.
\end{itemize}

\section{9.24 What to remember}\label{what-to-remember-1}

\begin{itemize}
\tightlist
\item
  Embedding geometry depends on the representation, preprocessing, and
  metric.
\item
  Cosine similarity compares directions; Euclidean distance compares
  positions; inner products also reward large norms.
\item
  PCA is a linear variance-preserving projection, not a perfect picture
  of semantic space.
\item
  Clustering embeddings means optimizing a geometric objective, not
  discovering guaranteed natural categories.
\item
  High-dimensional spaces can exhibit distance concentration, hubness,
  and anisotropy.
\item
  Linear directions can be useful for interpretation, but they can also
  reveal or amplify bias.
\item
  Contextual embeddings move from word-type geometry to token-occurrence
  geometry.
\end{itemize}

\section{Notes and references}\label{notes-and-references-4}

This chapter builds on vector space models and word embeddings from
Chapters 7 and 8. Classical background on distributional representations
appears in Manning and Schuetze (1999) and Manning et al. (2008). Modern
word embedding models include Mikolov, Chen, et al. (2013) and
Pennington et al. (2014). The broader NLP context is developed in
Jurafsky and Martin (2024).

\part{Part IV. Neural Networks for Sequences}

\chapter{Neural Networks for Text}\label{neural-networks-for-text}

Embeddings, softmax classifiers, and trainable nonlinear maps

\hfill\break

Classical text models often use a fixed representation such as a count
vector, TF--IDF vector, or SVD coordinate. A \textbf{neural network for
text} replaces part of this representation pipeline by a trainable map.
Tokens become vectors, vectors are combined into a document or sequence
representation, and the resulting representation is sent through
nonlinear maps and a softmax output layer.

The goal of this chapter is not to treat neural networks as black boxes.
The goal is to write down the maps, dimensions, losses, and gradients
clearly enough that the model becomes a familiar object from linear
algebra, probability, and optimization.

\section{Learning goals}\label{learning-goals-9}

After studying this chapter, you should be able to:

\begin{enumerate}
\def\labelenumi{\arabic{enumi}.}
\tightlist
\item
  describe a neural text classifier as a composition of maps;
\item
  distinguish fixed text features from trainable embeddings;
\item
  compute softmax probabilities and cross-entropy loss;
\item
  derive the gradient of softmax cross-entropy with respect to logits;
\item
  explain how an embedding layer receives gradients from a
  sentence-level loss;
\item
  implement a small neural text classifier from scratch in Python;
\item
  interpret regularization, temperature, and calibration for neural text
  models.
\end{enumerate}

\section{10.1 The modeling pipeline}\label{the-modeling-pipeline}

A simple neural text classifier has four conceptual steps.

\[
\text{text}
\longrightarrow
\text{token indices}
\longrightarrow
\text{token embeddings}
\longrightarrow
\text{text representation}
\longrightarrow
\text{class probabilities}.
\]

Let \(V = \{1,\ldots,d\}\) be a finite vocabulary. A document or
sentence is represented as a token sequence

\[
x=(x_1,\ldots,x_T) \in V^T.
\]

A neural model is a parameterized function

\[
f_\theta: V^T \to \Delta^{K-1},
\]

where \(K\) is the number of classes and

\[
\Delta^{K-1}
=
\left\{p \in \mathbb R^K: p_k \ge 0,\ \sum_{k=1}^K p_k=1\right\}
\]

is the probability simplex. The output \(f_\theta(x)\) is interpreted as
a conditional distribution over labels:

\[
f_\theta(x)_k \approx P(Y=k\mid X=x).
\]

A minimal neural text classifier can be written as

\[
\begin{aligned}
e_i &= E_{x_i} \in \mathbb R^m, \qquad i=1,\ldots,T,\\
h(x) &= \frac{1}{T}\sum_{i=1}^T e_i \in \mathbb R^m,\\
r(x) &= \sigma(Ah(x)+a) \in \mathbb R^q,\\
z(x) &= Wr(x)+b \in \mathbb R^K,\\
f_\theta(x) &= \operatorname{softmax}(z(x)).
\end{aligned}
\]

Here:

\begin{itemize}
\tightlist
\item
  \(E\in \mathbb R^{d\times m}\) is an \textbf{embedding matrix};
\item
  \(E_j\) is the row vector representing token \(j\);
\item
  \(h(x)\) is an averaged sentence representation;
\item
  \(A\in \mathbb R^{q\times m}\) and \(a\in \mathbb R^q\) define a
  hidden layer;
\item
  \(\sigma\) is a nonlinear activation such as ReLU or tanh;
\item
  \(W\in \mathbb R^{K\times q}\) and \(b\in \mathbb R^K\) define the
  output logits;
\item
  \(\theta=(E,A,a,W,b)\) denotes all trainable parameters.
\end{itemize}

\textbf{Core idea.} A neural text classifier is a composition of a
discrete lookup map, vector averaging or sequence aggregation, nonlinear
transformations, and a probabilistic output map.

\section{10.2 From fixed features to trainable
features}\label{from-fixed-features-to-trainable-features}

Earlier chapters studied fixed representation maps such as

\[
\phi(x) \in \mathbb R^d,
\]

where \(\phi(x)\) may be a count vector, a TF--IDF vector, or an SVD
coordinate. A linear classifier then uses

\[
z(x)=W\phi(x)+b.
\]

In a neural model, part of \(\phi\) is learned from data. For example,
an embedding layer defines

\[
\phi_E(x)
=
\frac{1}{T}\sum_{i=1}^T E_{x_i}.
\]

The representation depends on trainable parameters \(E\). Therefore
training changes not only the classifier weights, but also the geometry
in which tokens are represented.

This is the main mathematical shift:

\[
\text{fixed representation } \phi
\quad \longrightarrow \quad
\text{trainable representation } \phi_\theta.
\]

The statistical question also changes. We are no longer only estimating
class probabilities after selecting features. We are jointly estimating
a feature map and a classifier.

\section{10.3 The embedding layer as a matrix
lookup}\label{the-embedding-layer-as-a-matrix-lookup}

Let \(E\in \mathbb R^{d\times m}\). If token \(j\in V\) is represented
by the one-hot vector \(u_j\in \mathbb R^d\), then its embedding is

\[
E^T u_j \in \mathbb R^m.
\]

Equivalently, if \(E_j\) denotes the \(j\)th row of \(E\), then

\[
E^T u_j = E_j^T.
\]

Thus an embedding layer is just a matrix multiplication by a one-hot
vector, implemented efficiently as a row lookup.

For a sequence \(x=(x_1,\ldots,x_T)\), the embedding layer produces a
matrix

\[
M(x)
=
\begin{bmatrix}
E_{x_1}\\
E_{x_2}\\
\vdots\\
E_{x_T}
\end{bmatrix}
\in \mathbb R^{T\times m}.
\]

A simple sentence representation is average pooling:

\[
h(x)=\frac{1}{T}\mathbf 1_T^T M(x).
\]

This representation is permutation-invariant. It ignores word order.
That is a serious limitation for some language tasks, but it is a useful
first model because its mathematics is transparent.

\textbf{Dimension check.} If \(E\in \mathbb R^{d\times m}\) and
\(M(x)\in \mathbb R^{T\times m}\), then average pooling gives
\(h(x)\in \mathbb R^m\). If the hidden layer has width \(q\) and there
are \(K\) classes, then \(A\in \mathbb R^{q\times m}\) and
\(W\in \mathbb R^{K\times q}\).

\section{10.4 Activations and nonlinear
maps}\label{activations-and-nonlinear-maps}

A neural network becomes more than a linear classifier because of
nonlinear activation functions. Common choices include

\[
\operatorname{ReLU}(t)=\max\{0,t\},
\]

\[
\tanh(t)=\frac{e^t-e^{-t}}{e^t+e^{-t}},
\]

and the logistic sigmoid

\[
\sigma(t)=\frac{1}{1+e^{-t}}.
\]

Applied coordinatewise, an activation maps vectors to vectors:

\[
\sigma(v)=(\sigma(v_1),\ldots,\sigma(v_q)).
\]

A one-hidden-layer classifier is

\[
z(x)=W\sigma(Ah(x)+a)+b.
\]

Without the activation, this collapses to another affine map:

\[
W(Ah+a)+b=(WA)h+(Wa+b).
\]

Therefore the activation is what allows the model to represent nonlinear
decision boundaries in the representation space.

\section{10.5 Softmax output and class
probabilities}\label{softmax-output-and-class-probabilities}

For \(K\) classes, the model produces logits

\[
z=(z_1,\ldots,z_K)\in \mathbb R^K.
\]

The softmax function maps logits to probabilities:

\[
\operatorname{softmax}(z)_k
=
\frac{e^{z_k}}{\sum_{\ell=1}^K e^{z_\ell}}.
\]

Softmax is invariant under adding a constant to all logits:

\[
\operatorname{softmax}(z+c\mathbf 1)=\operatorname{softmax}(z).
\]

This identity is important for numerical stability. In computation, we
usually evaluate

\[
\operatorname{softmax}(z)
=
\operatorname{softmax}(z-\max_j z_j\mathbf 1).
\]

For binary classification with classes \(0\) and \(1\), the two-class
softmax is equivalent to logistic regression. If

\[
z_1-z_0=s,
\]

then

\[
P(Y=1\mid x)=\frac{1}{1+e^{-s}}.
\]

Thus a neural binary classifier can be interpreted as logistic
regression applied after a learned nonlinear representation.

\section{10.6 Cross-entropy and maximum
likelihood}\label{cross-entropy-and-maximum-likelihood}

Suppose the training data are

\[
\mathcal D_n=\{(x_i,y_i)\}_{i=1}^n,
\]

where \(y_i\in\{1,\ldots,K\}\). Let \(p_\theta(x_i)=f_\theta(x_i)\). The
model assigns probability

\[
p_\theta(y_i\mid x_i)=p_\theta(x_i)_{y_i}.
\]

The conditional likelihood is

\[
L(\theta)
=
\prod_{i=1}^n p_\theta(y_i\mid x_i).
\]

Maximizing likelihood is equivalent to minimizing negative
log-likelihood:

\[
\ell(\theta)
=
-
\sum_{i=1}^n \log p_\theta(y_i\mid x_i).
\]

If \(e_{y_i}\in\mathbb R^K\) is the one-hot vector for the true label,
the per-example loss is

\[
\ell_i(\theta)
=
-
\sum_{k=1}^K (e_{y_i})_k \log p_\theta(x_i)_k.
\]

This is the cross-entropy between the empirical one-hot label
distribution and the model's predicted distribution.

For empirical risk minimization, we usually minimize the average loss:

\[
\widehat R_n(\theta)
=
\frac{1}{n}\sum_{i=1}^n \ell_i(\theta).
\]

\section{10.7 The key gradient: softmax plus
cross-entropy}\label{the-key-gradient-softmax-plus-cross-entropy}

The softmax cross-entropy gradient has a remarkably simple form. Let

\[
p=\operatorname{softmax}(z),
\qquad
\ell(z,y)=-\log p_y.
\]

Then

\[
\frac{\partial \ell}{\partial z_k}=p_k-\mathbf 1\{k=y\}.
\]

In vector notation,

\[
\nabla_z \ell=p-e_y.
\]

This formula is one reason softmax cross-entropy is convenient for
classification.

\subsection{Derivation}\label{derivation-1}

First note that

\[
\log p_y
=
\log\left(\frac{e^{z_y}}{\sum_{\ell=1}^K e^{z_\ell}}\right)
=
z_y-
\log\left(\sum_{\ell=1}^K e^{z_\ell}\right).
\]

Thus

\[
\ell(z,y)
=
-z_y+
\log\left(\sum_{\ell=1}^K e^{z_\ell}\right).
\]

Differentiating with respect to \(z_k\) gives

\[
\frac{\partial \ell}{\partial z_k}
=
-\mathbf 1\{k=y\}
+
\frac{e^{z_k}}{\sum_{\ell=1}^K e^{z_\ell}}
=
p_k-\mathbf 1\{k=y\}.
\]

\textbf{Statistical interpretation.} The vector \(p-e_y\) is the
difference between the predicted class distribution and the empirical
one-hot class distribution. Gradient descent pushes predicted
probabilities toward the observed labels.

\section{10.8 Backpropagation through an averaged embedding
model}\label{backpropagation-through-an-averaged-embedding-model}

Consider the simple model

\[
h=\frac{1}{T}\sum_{i=1}^T E_{x_i},
\qquad
z=Wh+b,
\qquad
p=\operatorname{softmax}(z),
\qquad
\ell=-\log p_y.
\]

Let

\[
g_z=\nabla_z\ell=p-e_y.
\]

Then the gradients are

\[
\nabla_W\ell=g_z h^T,
\]

\[
\nabla_b\ell=g_z,
\]

and

\[
g_h=\nabla_h\ell=W^Tg_z.
\]

Since

\[
h=\frac{1}{T}\sum_{i=1}^T E_{x_i},
\]

each token embedding in the document receives the gradient

\[
\nabla_{E_{x_i}}\ell = \frac{1}{T}g_h.
\]

If the same token appears multiple times, its row in \(E\) accumulates
multiple contributions. This is why embedding layers are trainable even
though token indices are discrete: the gradient flows into the
continuous row vectors selected by the indices.

\section{10.9 Optimization and
regularization}\label{optimization-and-regularization}

Neural text models are usually trained with stochastic optimization. For
a mini-batch \(B\subset\{1,\ldots,n\}\), define

\[
\widehat R_B(\theta)
=
\frac{1}{|B|}\sum_{i\in B}\ell_i(\theta).
\]

A stochastic gradient descent update is

\[
\theta_{t+1}=\theta_t-\eta_t\nabla_\theta \widehat R_B(\theta_t),
\]

where \(\eta_t>0\) is a learning rate.

A regularized objective may be written as

\[
J(\theta)=\widehat R_n(\theta)+\lambda\|\theta\|_2^2.
\]

In practice, one often regularizes only selected parameters, such as
\(A\) and \(W\), not necessarily the embedding matrix. Dropout, early
stopping, and weight decay are also commonly used.

From a statistical perspective, regularization is important because
neural models often have many more parameters than training examples.
Without regularization or inductive bias, empirical loss alone may not
control generalization.

\section{10.10 Temperature, confidence, and
calibration}\label{temperature-confidence-and-calibration}

For logits \(z\), define temperature-scaled softmax by

\[
p_k(\tau)
=
\frac{\exp(z_k/\tau)}{\sum_{\ell=1}^K\exp(z_\ell/\tau)},
\qquad \tau>0.
\]

Small \(\tau\) makes the distribution sharper. Large \(\tau\) makes it
flatter. Temperature does not change the ranking of logits, but it
changes the confidence assigned to predictions.

This distinction matters in NLP because a model can be accurate but
poorly calibrated. Calibration asks whether predicted probabilities
match empirical frequencies. For example, among all examples to which a
model assigns confidence near \(0.8\), approximately \(80\%\) should be
correct if the model is calibrated.

\section{10.11 Interactive: Softmax classifier and
temperature}\label{interactive-softmax-classifier-and-temperature}

The following interactive illustrates a binary softmax classifier in a
two-feature representation space. Think of the horizontal coordinate as
the strength of evidence for technical or mathematical language and the
vertical coordinate as the strength of evidence for casual or
nontechnical language. The sliders control the logit difference

\[
s(x)=w_1x_1+w_2x_2+b,
\]

and the model probability

\[
P(Y=1\mid x)=\frac{1}{1+\exp(-s(x)/\tau)}.
\]

The HTML version of this chapter contains an interactive Plotly.js
visualization of a two-feature softmax classifier. In PDF format, the
same mathematical object is

\[
P(Y=1\mid x)=\frac{1}{1+\exp(-(w_1x_1+w_2x_2+b)/\tau)}.
\]

The decision boundary is the line \(w_1x_1+w_2x_2+b=0\). Temperature
\(\tau\) changes confidence but not the boundary.

\section{10.12 Python example: data, tokenization, and
vocabulary}\label{python-example-data-tokenization-and-vocabulary}

The following example builds a tiny text classification dataset. The
goal is not high performance. The goal is to expose every mathematical
object in a neural text classifier.

\begin{Shaded}
\begin{Highlighting}[]
\ImportTok{import}\NormalTok{ re}
\ImportTok{from}\NormalTok{ collections }\ImportTok{import}\NormalTok{ Counter}

\NormalTok{texts }\OperatorTok{=}\NormalTok{ [}
    \StringTok{"linear algebra and probability are useful for machine learning"}\NormalTok{,}
    \StringTok{"statistics studies uncertainty and data"}\NormalTok{,}
    \StringTok{"neural networks learn representations from data"}\NormalTok{,}
    \StringTok{"the theorem gives a clean mathematical proof"}\NormalTok{,}
    \StringTok{"the movie was funny and enjoyable"}\NormalTok{,}
    \StringTok{"the dinner was tasty and relaxing"}\NormalTok{,}
    \StringTok{"music and stories made the evening fun"}\NormalTok{,}
    \StringTok{"the game was exciting and easy to watch"}\NormalTok{,}
\NormalTok{]}

\CommentTok{\# 1 = mathematical/technical, 0 = casual/nontechnical}
\NormalTok{labels }\OperatorTok{=}\NormalTok{ [}\DecValTok{1}\NormalTok{, }\DecValTok{1}\NormalTok{, }\DecValTok{1}\NormalTok{, }\DecValTok{1}\NormalTok{, }\DecValTok{0}\NormalTok{, }\DecValTok{0}\NormalTok{, }\DecValTok{0}\NormalTok{, }\DecValTok{0}\NormalTok{]}

\KeywordTok{def}\NormalTok{ tokenize(text):}
    \ControlFlowTok{return}\NormalTok{ re.findall(}\VerbatimStringTok{r"}\PreprocessorTok{[a{-}z]}\OperatorTok{+}\VerbatimStringTok{"}\NormalTok{, text.lower())}

\NormalTok{tokenized }\OperatorTok{=}\NormalTok{ [tokenize(t) }\ControlFlowTok{for}\NormalTok{ t }\KeywordTok{in}\NormalTok{ texts]}
\NormalTok{counts }\OperatorTok{=}\NormalTok{ Counter(tok }\ControlFlowTok{for}\NormalTok{ doc }\KeywordTok{in}\NormalTok{ tokenized }\ControlFlowTok{for}\NormalTok{ tok }\KeywordTok{in}\NormalTok{ doc)}
\NormalTok{vocab }\OperatorTok{=}\NormalTok{ \{tok: j }\ControlFlowTok{for}\NormalTok{ j, (tok, \_) }\KeywordTok{in} \BuiltInTok{enumerate}\NormalTok{(counts.most\_common())\}}

\BuiltInTok{print}\NormalTok{(}\StringTok{"Vocabulary size:"}\NormalTok{, }\BuiltInTok{len}\NormalTok{(vocab))}
\BuiltInTok{print}\NormalTok{(}\StringTok{"First 10 vocabulary items:"}\NormalTok{, }\BuiltInTok{list}\NormalTok{(vocab.items())[:}\DecValTok{10}\NormalTok{])}
\BuiltInTok{print}\NormalTok{(}\StringTok{"First document tokens:"}\NormalTok{, tokenized[}\DecValTok{0}\NormalTok{])}
\end{Highlighting}
\end{Shaded}

\begin{verbatim}
Vocabulary size: 42
First 10 vocabulary items: [('and', 0), ('the', 1), ('was', 2), ('data', 3), ('linear', 4), ('algebra', 5), ('probability', 6), ('are', 7), ('useful', 8), ('for', 9)]
First document tokens: ['linear', 'algebra', 'and', 'probability', 'are', 'useful', 'for', 'machine', 'learning']
\end{verbatim}

Convert each document to a list of token indices.

\begin{Shaded}
\begin{Highlighting}[]
\NormalTok{indexed }\OperatorTok{=}\NormalTok{ [[vocab[tok] }\ControlFlowTok{for}\NormalTok{ tok }\KeywordTok{in}\NormalTok{ doc] }\ControlFlowTok{for}\NormalTok{ doc }\KeywordTok{in}\NormalTok{ tokenized]}
\BuiltInTok{print}\NormalTok{(indexed[}\DecValTok{0}\NormalTok{])}
\end{Highlighting}
\end{Shaded}

\begin{verbatim}
[4, 5, 0, 6, 7, 8, 9, 10, 11]
\end{verbatim}

\section{10.13 Python example: softmax and
cross-entropy}\label{python-example-softmax-and-cross-entropy}

A numerically stable softmax subtracts the largest logit.

\begin{Shaded}
\begin{Highlighting}[]
\ImportTok{import}\NormalTok{ numpy }\ImportTok{as}\NormalTok{ np}

\NormalTok{rng }\OperatorTok{=}\NormalTok{ np.random.default\_rng(}\DecValTok{7243}\NormalTok{)}

\KeywordTok{def}\NormalTok{ softmax(Z):}
\NormalTok{    Z\_shift }\OperatorTok{=}\NormalTok{ Z }\OperatorTok{{-}}\NormalTok{ Z.}\BuiltInTok{max}\NormalTok{(axis}\OperatorTok{=}\DecValTok{1}\NormalTok{, keepdims}\OperatorTok{=}\VariableTok{True}\NormalTok{)}
\NormalTok{    expZ }\OperatorTok{=}\NormalTok{ np.exp(Z\_shift)}
    \ControlFlowTok{return}\NormalTok{ expZ }\OperatorTok{/}\NormalTok{ expZ.}\BuiltInTok{sum}\NormalTok{(axis}\OperatorTok{=}\DecValTok{1}\NormalTok{, keepdims}\OperatorTok{=}\VariableTok{True}\NormalTok{)}

\KeywordTok{def}\NormalTok{ cross\_entropy(P, y):}
\NormalTok{    n }\OperatorTok{=} \BuiltInTok{len}\NormalTok{(y)}
    \ControlFlowTok{return} \OperatorTok{{-}}\NormalTok{np.mean(np.log(P[np.arange(n), y] }\OperatorTok{+} \FloatTok{1e{-}12}\NormalTok{))}

\NormalTok{Z }\OperatorTok{=}\NormalTok{ np.array([[}\FloatTok{2.0}\NormalTok{, }\FloatTok{0.5}\NormalTok{], [}\FloatTok{0.2}\NormalTok{, }\FloatTok{1.7}\NormalTok{]])}
\NormalTok{P }\OperatorTok{=}\NormalTok{ softmax(Z)}
\BuiltInTok{print}\NormalTok{(P)}
\BuiltInTok{print}\NormalTok{(}\StringTok{"loss:"}\NormalTok{, cross\_entropy(P, np.array([}\DecValTok{0}\NormalTok{, }\DecValTok{1}\NormalTok{])))}
\end{Highlighting}
\end{Shaded}

\begin{verbatim}
[[0.81757448 0.18242552]
 [0.18242552 0.81757448]]
loss: 0.2014132779815293
\end{verbatim}

\section{10.14 Python example: averaged embedding classifier from
scratch}\label{python-example-averaged-embedding-classifier-from-scratch}

We now train the model

\[
h_i=\frac{1}{T_i}\sum_{t=1}^{T_i}E_{x_{it}},
\qquad
z_i=h_iW+b,
\qquad
p_i=\operatorname{softmax}(z_i).
\]

Here \(E\in\mathbb R^{d\times m}\) and \(W\in\mathbb R^{m\times K}\).

\begin{Shaded}
\begin{Highlighting}[]
\CommentTok{\# Model dimensions}
\NormalTok{D }\OperatorTok{=} \BuiltInTok{len}\NormalTok{(vocab)}
\NormalTok{m }\OperatorTok{=} \DecValTok{6}
\NormalTok{K }\OperatorTok{=} \DecValTok{2}
\NormalTok{n }\OperatorTok{=} \BuiltInTok{len}\NormalTok{(indexed)}

\NormalTok{y }\OperatorTok{=}\NormalTok{ np.array(labels, dtype}\OperatorTok{=}\BuiltInTok{int}\NormalTok{)}
\NormalTok{E }\OperatorTok{=} \FloatTok{0.05} \OperatorTok{*}\NormalTok{ rng.normal(size}\OperatorTok{=}\NormalTok{(D, m))}
\NormalTok{W }\OperatorTok{=} \FloatTok{0.05} \OperatorTok{*}\NormalTok{ rng.normal(size}\OperatorTok{=}\NormalTok{(m, K))}
\NormalTok{b }\OperatorTok{=}\NormalTok{ np.zeros(K)}

\KeywordTok{def}\NormalTok{ forward(indexed\_docs, E, W, b):}
\NormalTok{    H }\OperatorTok{=}\NormalTok{ np.zeros((}\BuiltInTok{len}\NormalTok{(indexed\_docs), E.shape[}\DecValTok{1}\NormalTok{]))}
    \ControlFlowTok{for}\NormalTok{ i, doc }\KeywordTok{in} \BuiltInTok{enumerate}\NormalTok{(indexed\_docs):}
\NormalTok{        H[i] }\OperatorTok{=}\NormalTok{ E[doc].mean(axis}\OperatorTok{=}\DecValTok{0}\NormalTok{)}
\NormalTok{    Z }\OperatorTok{=}\NormalTok{ H }\OperatorTok{@}\NormalTok{ W }\OperatorTok{+}\NormalTok{ b}
\NormalTok{    P }\OperatorTok{=}\NormalTok{ softmax(Z)}
    \ControlFlowTok{return}\NormalTok{ H, Z, P}

\NormalTok{H, Z, P }\OperatorTok{=}\NormalTok{ forward(indexed, E, W, b)}
\BuiltInTok{print}\NormalTok{(}\StringTok{"H shape:"}\NormalTok{, H.shape)}
\BuiltInTok{print}\NormalTok{(}\StringTok{"Z shape:"}\NormalTok{, Z.shape)}
\BuiltInTok{print}\NormalTok{(}\StringTok{"Initial loss:"}\NormalTok{, cross\_entropy(P, y))}
\end{Highlighting}
\end{Shaded}

\begin{verbatim}
H shape: (8, 6)
Z shape: (8, 2)
Initial loss: 0.6922852483744657
\end{verbatim}

Training by full-batch gradient descent.

\begin{Shaded}
\begin{Highlighting}[]
\NormalTok{lr }\OperatorTok{=} \FloatTok{1.0}
\NormalTok{history }\OperatorTok{=}\NormalTok{ []}

\ControlFlowTok{for}\NormalTok{ epoch }\KeywordTok{in} \BuiltInTok{range}\NormalTok{(}\DecValTok{300}\NormalTok{):}
\NormalTok{    H, Z, P }\OperatorTok{=}\NormalTok{ forward(indexed, E, W, b)}
\NormalTok{    loss }\OperatorTok{=}\NormalTok{ cross\_entropy(P, y)}
\NormalTok{    history.append(loss)}

    \CommentTok{\# Gradient of average cross{-}entropy with respect to logits}
\NormalTok{    Gz }\OperatorTok{=}\NormalTok{ P.copy()}
\NormalTok{    Gz[np.arange(n), y] }\OperatorTok{{-}=} \DecValTok{1}
\NormalTok{    Gz }\OperatorTok{/=}\NormalTok{ n}

    \CommentTok{\# Gradients for W and b}
\NormalTok{    grad\_W }\OperatorTok{=}\NormalTok{ H.T }\OperatorTok{@}\NormalTok{ Gz}
\NormalTok{    grad\_b }\OperatorTok{=}\NormalTok{ Gz.}\BuiltInTok{sum}\NormalTok{(axis}\OperatorTok{=}\DecValTok{0}\NormalTok{)}

    \CommentTok{\# Gradient with respect to averaged document representations}
\NormalTok{    grad\_H }\OperatorTok{=}\NormalTok{ Gz }\OperatorTok{@}\NormalTok{ W.T}

    \CommentTok{\# Gradient for embedding matrix}
\NormalTok{    grad\_E }\OperatorTok{=}\NormalTok{ np.zeros\_like(E)}
    \ControlFlowTok{for}\NormalTok{ i, doc }\KeywordTok{in} \BuiltInTok{enumerate}\NormalTok{(indexed):}
\NormalTok{        contribution }\OperatorTok{=}\NormalTok{ grad\_H[i] }\OperatorTok{/} \BuiltInTok{len}\NormalTok{(doc)}
        \ControlFlowTok{for}\NormalTok{ token\_id }\KeywordTok{in}\NormalTok{ doc:}
\NormalTok{            grad\_E[token\_id] }\OperatorTok{+=}\NormalTok{ contribution}

\NormalTok{    E }\OperatorTok{{-}=}\NormalTok{ lr }\OperatorTok{*}\NormalTok{ grad\_E}
\NormalTok{    W }\OperatorTok{{-}=}\NormalTok{ lr }\OperatorTok{*}\NormalTok{ grad\_W}
\NormalTok{    b }\OperatorTok{{-}=}\NormalTok{ lr }\OperatorTok{*}\NormalTok{ grad\_b}

\NormalTok{H, Z, P }\OperatorTok{=}\NormalTok{ forward(indexed, E, W, b)}
\NormalTok{pred }\OperatorTok{=}\NormalTok{ P.argmax(axis}\OperatorTok{=}\DecValTok{1}\NormalTok{)}
\BuiltInTok{print}\NormalTok{(}\StringTok{"Final loss:"}\NormalTok{, cross\_entropy(P, y))}
\BuiltInTok{print}\NormalTok{(}\StringTok{"Predictions:"}\NormalTok{, pred.tolist())}
\BuiltInTok{print}\NormalTok{(}\StringTok{"Training accuracy:"}\NormalTok{, }\BuiltInTok{float}\NormalTok{(np.mean(pred }\OperatorTok{==}\NormalTok{ y)))}
\BuiltInTok{print}\NormalTok{(}\StringTok{"First five losses:"}\NormalTok{, [}\BuiltInTok{float}\NormalTok{(np.}\BuiltInTok{round}\NormalTok{(v, }\DecValTok{4}\NormalTok{)) }\ControlFlowTok{for}\NormalTok{ v }\KeywordTok{in}\NormalTok{ history[:}\DecValTok{5}\NormalTok{]])}
\BuiltInTok{print}\NormalTok{(}\StringTok{"Last five losses:"}\NormalTok{, [}\BuiltInTok{float}\NormalTok{(np.}\BuiltInTok{round}\NormalTok{(v, }\DecValTok{4}\NormalTok{)) }\ControlFlowTok{for}\NormalTok{ v }\KeywordTok{in}\NormalTok{ history[}\OperatorTok{{-}}\DecValTok{5}\NormalTok{:]])}
\end{Highlighting}
\end{Shaded}

\begin{verbatim}
Final loss: 0.0019680287531346615
Predictions: [1, 1, 1, 1, 0, 0, 0, 0]
Training accuracy: 1.0
First five losses: [0.6923, 0.6919, 0.6915, 0.691, 0.6904]
Last five losses: [0.002, 0.002, 0.002, 0.002, 0.002]
\end{verbatim}

Plot the loss curve.

\begin{Shaded}
\begin{Highlighting}[]
\ImportTok{import}\NormalTok{ matplotlib.pyplot }\ImportTok{as}\NormalTok{ plt}

\NormalTok{plt.figure(figsize}\OperatorTok{=}\NormalTok{(}\DecValTok{7}\NormalTok{, }\DecValTok{4}\NormalTok{))}
\NormalTok{plt.plot(history)}
\NormalTok{plt.xlabel(}\StringTok{"epoch"}\NormalTok{)}
\NormalTok{plt.ylabel(}\StringTok{"cross{-}entropy loss"}\NormalTok{)}
\NormalTok{plt.title(}\StringTok{"Training loss for an averaged{-}embedding classifier"}\NormalTok{)}
\NormalTok{plt.show()}
\end{Highlighting}
\end{Shaded}

\pandocbounded{\includegraphics[keepaspectratio]{chapters/chapter-10-neural-networks-for-text_files/figure-pdf/cell-7-output-1.pdf}}

\section{10.15 Python example: inspecting learned token
directions}\label{python-example-inspecting-learned-token-directions}

For binary classification, the difference between the two output weight
columns gives a direction in embedding space:

\[
v=W_{:,1}-W_{:,0}.
\]

For an embedded token \(E_j\), the score

\[
E_j^T v
\]

measures how much that token pushes the classifier toward class \(1\)
rather than class \(0\), at least in the averaged linear model.

\begin{Shaded}
\begin{Highlighting}[]
\NormalTok{direction }\OperatorTok{=}\NormalTok{ W[:, }\DecValTok{1}\NormalTok{] }\OperatorTok{{-}}\NormalTok{ W[:, }\DecValTok{0}\NormalTok{]}
\NormalTok{scores }\OperatorTok{=}\NormalTok{ E }\OperatorTok{@}\NormalTok{ direction}
\NormalTok{id\_to\_token }\OperatorTok{=}\NormalTok{ \{j: tok }\ControlFlowTok{for}\NormalTok{ tok, j }\KeywordTok{in}\NormalTok{ vocab.items()\}}

\NormalTok{ranked }\OperatorTok{=} \BuiltInTok{sorted}\NormalTok{([(id\_to\_token[j], scores[j]) }\ControlFlowTok{for}\NormalTok{ j }\KeywordTok{in} \BuiltInTok{range}\NormalTok{(D)], key}\OperatorTok{=}\KeywordTok{lambda}\NormalTok{ t: t[}\DecValTok{1}\NormalTok{], reverse}\OperatorTok{=}\VariableTok{True}\NormalTok{)}
\BuiltInTok{print}\NormalTok{(}\StringTok{"Most class{-}1{-}associated tokens:"}\NormalTok{)}
\ControlFlowTok{for}\NormalTok{ tok, score }\KeywordTok{in}\NormalTok{ ranked[:}\DecValTok{8}\NormalTok{]:}
    \BuiltInTok{print}\NormalTok{(}\SpecialStringTok{f"}\SpecialCharTok{\{}\NormalTok{tok}\SpecialCharTok{:\textgreater{}15s\}}\SpecialStringTok{  }\SpecialCharTok{\{}\NormalTok{score}\SpecialCharTok{: .3f\}}\SpecialStringTok{"}\NormalTok{)}

\BuiltInTok{print}\NormalTok{(}\StringTok{"}\CharTok{\textbackslash{}n}\StringTok{Most class{-}0{-}associated tokens:"}\NormalTok{)}
\ControlFlowTok{for}\NormalTok{ tok, score }\KeywordTok{in}\NormalTok{ ranked[}\OperatorTok{{-}}\DecValTok{8}\NormalTok{:]:}
    \BuiltInTok{print}\NormalTok{(}\SpecialStringTok{f"}\SpecialCharTok{\{}\NormalTok{tok}\SpecialCharTok{:\textgreater{}15s\}}\SpecialStringTok{  }\SpecialCharTok{\{}\NormalTok{score}\SpecialCharTok{: .3f\}}\SpecialStringTok{"}\NormalTok{)}
\end{Highlighting}
\end{Shaded}

\begin{verbatim}
Most class-1-associated tokens:
           data   11.539
        theorem   8.061
   mathematical   7.781
          gives   7.360
          proof   7.333
    uncertainty   7.122
     statistics   6.976
        studies   6.951

Most class-0-associated tokens:
           made  -5.794
        stories  -6.040
          music  -6.161
        evening  -6.299
            fun  -6.635
            and  -7.697
            the  -13.089
            was  -13.491
\end{verbatim}

This kind of inspection is limited but useful. It helps students see
that a trainable embedding model still has interpretable linear algebra
inside it.

\section{10.16 Python example: adding one hidden ReLU
layer}\label{python-example-adding-one-hidden-relu-layer}

A one-hidden-layer model uses

\[
r_i=\operatorname{ReLU}(h_iA+a),
\qquad
z_i=r_iW+b.
\]

The additional nonlinear layer can represent more complicated decision
regions in \(h\)-space.

\begin{Shaded}
\begin{Highlighting}[]
\KeywordTok{def}\NormalTok{ relu(U):}
    \ControlFlowTok{return}\NormalTok{ np.maximum(U, }\FloatTok{0.0}\NormalTok{)}

\KeywordTok{def}\NormalTok{ relu\_grad(U):}
    \ControlFlowTok{return}\NormalTok{ (U }\OperatorTok{\textgreater{}} \DecValTok{0}\NormalTok{).astype(}\BuiltInTok{float}\NormalTok{)}

\NormalTok{q }\OperatorTok{=} \DecValTok{8}
\NormalTok{E2 }\OperatorTok{=} \FloatTok{0.05} \OperatorTok{*}\NormalTok{ rng.normal(size}\OperatorTok{=}\NormalTok{(D, m))}
\NormalTok{A }\OperatorTok{=} \FloatTok{0.05} \OperatorTok{*}\NormalTok{ rng.normal(size}\OperatorTok{=}\NormalTok{(m, q))}
\NormalTok{a }\OperatorTok{=}\NormalTok{ np.zeros(q)}
\NormalTok{W2 }\OperatorTok{=} \FloatTok{0.05} \OperatorTok{*}\NormalTok{ rng.normal(size}\OperatorTok{=}\NormalTok{(q, K))}
\NormalTok{b2 }\OperatorTok{=}\NormalTok{ np.zeros(K)}

\KeywordTok{def}\NormalTok{ forward\_relu(indexed\_docs, E, A, a, W, b):}
\NormalTok{    H }\OperatorTok{=}\NormalTok{ np.zeros((}\BuiltInTok{len}\NormalTok{(indexed\_docs), E.shape[}\DecValTok{1}\NormalTok{]))}
    \ControlFlowTok{for}\NormalTok{ i, doc }\KeywordTok{in} \BuiltInTok{enumerate}\NormalTok{(indexed\_docs):}
\NormalTok{        H[i] }\OperatorTok{=}\NormalTok{ E[doc].mean(axis}\OperatorTok{=}\DecValTok{0}\NormalTok{)}
\NormalTok{    U }\OperatorTok{=}\NormalTok{ H }\OperatorTok{@}\NormalTok{ A }\OperatorTok{+}\NormalTok{ a}
\NormalTok{    R }\OperatorTok{=}\NormalTok{ relu(U)}
\NormalTok{    Z }\OperatorTok{=}\NormalTok{ R }\OperatorTok{@}\NormalTok{ W }\OperatorTok{+}\NormalTok{ b}
\NormalTok{    P }\OperatorTok{=}\NormalTok{ softmax(Z)}
    \ControlFlowTok{return}\NormalTok{ H, U, R, Z, P}

\NormalTok{history\_relu }\OperatorTok{=}\NormalTok{ []}
\NormalTok{lr }\OperatorTok{=} \FloatTok{0.8}
\ControlFlowTok{for}\NormalTok{ epoch }\KeywordTok{in} \BuiltInTok{range}\NormalTok{(}\DecValTok{400}\NormalTok{):}
\NormalTok{    H, U, R, Z, P }\OperatorTok{=}\NormalTok{ forward\_relu(indexed, E2, A, a, W2, b2)}
\NormalTok{    loss }\OperatorTok{=}\NormalTok{ cross\_entropy(P, y)}
\NormalTok{    history\_relu.append(loss)}

\NormalTok{    Gz }\OperatorTok{=}\NormalTok{ P.copy()}
\NormalTok{    Gz[np.arange(n), y] }\OperatorTok{{-}=} \DecValTok{1}
\NormalTok{    Gz }\OperatorTok{/=}\NormalTok{ n}

\NormalTok{    grad\_W2 }\OperatorTok{=}\NormalTok{ R.T }\OperatorTok{@}\NormalTok{ Gz}
\NormalTok{    grad\_b2 }\OperatorTok{=}\NormalTok{ Gz.}\BuiltInTok{sum}\NormalTok{(axis}\OperatorTok{=}\DecValTok{0}\NormalTok{)}
\NormalTok{    grad\_R }\OperatorTok{=}\NormalTok{ Gz }\OperatorTok{@}\NormalTok{ W2.T}
\NormalTok{    grad\_U }\OperatorTok{=}\NormalTok{ grad\_R }\OperatorTok{*}\NormalTok{ relu\_grad(U)}
\NormalTok{    grad\_A }\OperatorTok{=}\NormalTok{ H.T }\OperatorTok{@}\NormalTok{ grad\_U}
\NormalTok{    grad\_a }\OperatorTok{=}\NormalTok{ grad\_U.}\BuiltInTok{sum}\NormalTok{(axis}\OperatorTok{=}\DecValTok{0}\NormalTok{)}
\NormalTok{    grad\_H }\OperatorTok{=}\NormalTok{ grad\_U }\OperatorTok{@}\NormalTok{ A.T}

\NormalTok{    grad\_E2 }\OperatorTok{=}\NormalTok{ np.zeros\_like(E2)}
    \ControlFlowTok{for}\NormalTok{ i, doc }\KeywordTok{in} \BuiltInTok{enumerate}\NormalTok{(indexed):}
\NormalTok{        contribution }\OperatorTok{=}\NormalTok{ grad\_H[i] }\OperatorTok{/} \BuiltInTok{len}\NormalTok{(doc)}
        \ControlFlowTok{for}\NormalTok{ token\_id }\KeywordTok{in}\NormalTok{ doc:}
\NormalTok{            grad\_E2[token\_id] }\OperatorTok{+=}\NormalTok{ contribution}

\NormalTok{    E2 }\OperatorTok{{-}=}\NormalTok{ lr }\OperatorTok{*}\NormalTok{ grad\_E2}
\NormalTok{    A }\OperatorTok{{-}=}\NormalTok{ lr }\OperatorTok{*}\NormalTok{ grad\_A}
\NormalTok{    a }\OperatorTok{{-}=}\NormalTok{ lr }\OperatorTok{*}\NormalTok{ grad\_a}
\NormalTok{    W2 }\OperatorTok{{-}=}\NormalTok{ lr }\OperatorTok{*}\NormalTok{ grad\_W2}
\NormalTok{    b2 }\OperatorTok{{-}=}\NormalTok{ lr }\OperatorTok{*}\NormalTok{ grad\_b2}

\NormalTok{H, U, R, Z, P }\OperatorTok{=}\NormalTok{ forward\_relu(indexed, E2, A, a, W2, b2)}
\BuiltInTok{print}\NormalTok{(}\StringTok{"Final ReLU{-}model loss:"}\NormalTok{, cross\_entropy(P, y))}
\BuiltInTok{print}\NormalTok{(}\StringTok{"Training accuracy:"}\NormalTok{, }\BuiltInTok{float}\NormalTok{(np.mean(P.argmax(axis}\OperatorTok{=}\DecValTok{1}\NormalTok{) }\OperatorTok{==}\NormalTok{ y)))}
\end{Highlighting}
\end{Shaded}

\begin{verbatim}
Final ReLU-model loss: 0.0005048229998677652
Training accuracy: 1.0
\end{verbatim}

\section{10.17 Worked calculation: one softmax
gradient}\label{worked-calculation-one-softmax-gradient}

Suppose a three-class model gives logits

\[
z=(2,1,0).
\]

Then

\[
p=\operatorname{softmax}(z)
=\frac{1}{e^2+e+1}(e^2,e,1).
\]

If the true class is \(1\) using zero-based indexing, then
\(e_y=(0,1,0)\) and

\[
\nabla_z\ell=p-e_y.
\]

Using approximate values,

\[
p\approx(0.665,0.245,0.090),
\]

so

\[
\nabla_z\ell\approx(0.665,-0.755,0.090).
\]

The negative entry for the true class means gradient descent will
increase the true-class logit. The positive entries for the other
classes mean gradient descent will decrease the competing logits.

\section{10.18 Neural networks versus classical text
classifiers}\label{neural-networks-versus-classical-text-classifiers}

The table summarizes important contrasts.

{\def\LTcaptype{none} % do not increment counter
\begin{longtable}[]{@{}
  >{\raggedright\arraybackslash}p{(\linewidth - 8\tabcolsep) * \real{0.2000}}
  >{\raggedright\arraybackslash}p{(\linewidth - 8\tabcolsep) * \real{0.2000}}
  >{\raggedright\arraybackslash}p{(\linewidth - 8\tabcolsep) * \real{0.2000}}
  >{\raggedright\arraybackslash}p{(\linewidth - 8\tabcolsep) * \real{0.2000}}
  >{\raggedright\arraybackslash}p{(\linewidth - 8\tabcolsep) * \real{0.2000}}@{}}
\toprule\noalign{}
\begin{minipage}[b]{\linewidth}\raggedright
Model
\end{minipage} & \begin{minipage}[b]{\linewidth}\raggedright
Representation
\end{minipage} & \begin{minipage}[b]{\linewidth}\raggedright
Decision function
\end{minipage} & \begin{minipage}[b]{\linewidth}\raggedright
Strength
\end{minipage} & \begin{minipage}[b]{\linewidth}\raggedright
Limitation
\end{minipage} \\
\midrule\noalign{}
\endhead
\bottomrule\noalign{}
\endlastfoot
Naive Bayes & counts & log posterior score & simple probabilistic model
& strong independence assumption \\
Logistic regression & fixed feature vector & linear logits &
interpretable and convex & depends on chosen features \\
Averaged embeddings & trainable token vectors & linear or nonlinear
logits & learns representation & weak treatment of word order \\
RNN/LSTM & sequence of embeddings & recurrent hidden state & uses order
& harder to parallelize \\
Transformer & sequence of contextual embeddings & attention-based maps &
long-range interactions & expensive and data-hungry \\
\end{longtable}
}

This chapter studies the third row. The next two chapters develop
recurrent models that use word order more explicitly.

\section{10.19 AI component: using AI to audit a neural text
model}\label{ai-component-using-ai-to-audit-a-neural-text-model}

Modern AI tools can help students debug, explain, and critique models,
but they should not replace mathematical verification. Use AI as a
structured reviewer.

\subsection{Prompt 1: dimension
checking}\label{prompt-1-dimension-checking}

\begin{quote}
I have a vocabulary of size \(d\), embedding dimension \(m\), hidden
width \(q\), and \(K\) classes. My model is \(E\in R^{d\times m}\),
\(h\) is the average of token embeddings,
\(r=\operatorname{ReLU}(Ah+a)\), and \(z=Wr+b\). Check all dimensions
and tell me which convention I am using for row versus column vectors.
\end{quote}

\subsection{Prompt 2: gradient
auditing}\label{prompt-2-gradient-auditing}

\begin{quote}
For the averaged embedding classifier \(h=T^{-1}\sum_i E_{x_i}\),
\(z=Wh+b\), and \(\ell=-\log\operatorname{softmax}(z)_y\), derive the
gradients with respect to \(W\), \(b\), \(h\), and each selected
embedding row. State every shape.
\end{quote}

\subsection{Prompt 3: modeling
critique}\label{prompt-3-modeling-critique}

\begin{quote}
I trained an averaged-embedding sentiment classifier. It performs well
on short reviews but fails on sentences with negation, such as ``not
good''. Explain mathematically why averaging embeddings may fail and
suggest a model class that can represent word order.
\end{quote}

\subsection{Prompt 4: statistical
caution}\label{prompt-4-statistical-caution}

\begin{quote}
My neural text classifier achieves perfect training accuracy on 40
examples. List possible statistical explanations and propose validation
diagnostics.
\end{quote}

The best use of AI is to ask it for checks that you can verify. In this
book, every AI-generated explanation should be reduced to mathematical
objects: maps, losses, gradients, distributions, and evaluation metrics.

\section{10.20 Common pitfalls}\label{common-pitfalls-1}

\begin{enumerate}
\def\labelenumi{\arabic{enumi}.}
\tightlist
\item
  \textbf{Confusing tokens and embeddings.} Token indices are discrete.
  Embeddings are continuous row vectors.
\item
  \textbf{Ignoring dimensions.} Many implementation errors are shape
  errors disguised as conceptual errors.
\item
  \textbf{Interpreting softmax scores as calibrated probabilities
  without checking calibration.}
\item
  \textbf{Believing that more nonlinear layers always improve
  generalization.} Larger models can overfit small text datasets.
\item
  \textbf{Forgetting word order.} Averaging embeddings is simple but
  cannot distinguish many different sequences.
\item
  \textbf{Reading too much into learned embedding coordinates.}
  Individual coordinates are usually not uniquely meaningful because
  rotations or reparameterizations can preserve predictions.
\end{enumerate}

\section{10.21 Exercises}\label{exercises-9}

\subsection{Conceptual exercises}\label{conceptual-exercises-6}

\begin{enumerate}
\def\labelenumi{\arabic{enumi}.}
\tightlist
\item
  Let \(E\in\mathbb R^{d\times m}\) and \(x=(x_1,\ldots,x_T)\). Write
  average pooling as a matrix operation using one-hot token vectors.
\item
  Prove that softmax is invariant under adding a constant multiple of
  \(\mathbf 1\) to the logits.
\item
  Show that two-class softmax is equivalent to logistic regression on
  the logit difference.
\item
  Explain why an averaged-embedding classifier is invariant under
  permutations of the input tokens.
\item
  Give an example of two sentences with different meanings but the same
  averaged embedding representation.
\end{enumerate}

\subsection{Computational exercises}\label{computational-exercises-3}

\begin{enumerate}
\def\labelenumi{\arabic{enumi}.}
\tightlist
\item
  Modify the from-scratch classifier so that it uses sum pooling instead
  of average pooling. Compare the behavior on short and long documents.
\item
  Add \(\ell_2\) regularization to the output weight matrix \(W\) and
  study how the learned token scores change.
\item
  Split the toy dataset into train and test sets. Repeat training under
  several random seeds.
\item
  Replace ReLU by tanh in the one-hidden-layer model. Compare the
  training loss curve.
\item
  Implement mini-batch stochastic gradient descent for the
  averaged-embedding classifier.
\end{enumerate}

\subsection{Mathematical exercises}\label{mathematical-exercises-5}

\begin{enumerate}
\def\labelenumi{\arabic{enumi}.}
\tightlist
\item
  Derive the gradient of the regularized loss \[
  \widehat R_n(W)+\lambda\|W\|_F^2
  \] with respect to \(W\).
\item
  For the model \(z=Wh+b\), derive the Hessian of the binary logistic
  loss with respect to \(h\) for one observation.
\item
  Let \(Q\in\mathbb R^{m\times m}\) be orthogonal. If embeddings are
  transformed by \(E_j\mapsto E_jQ\) and output weights by
  \(W\mapsto Q^TW\), show that logits are unchanged in the linear
  averaged-embedding model.
\item
  Suppose a token occurs \(c_j\) times in a document. Derive its
  contribution to the gradient of \(E_j\) under average pooling.
\item
  Explain why the empirical risk is generally nonconvex in \((E,W)\)
  even when there is no hidden layer.
\end{enumerate}

\section{10.22 Mini-project: from count vectors to trainable
embeddings}\label{mini-project-from-count-vectors-to-trainable-embeddings}

Build three classifiers on the same small text dataset:

\begin{enumerate}
\def\labelenumi{\arabic{enumi}.}
\tightlist
\item
  multinomial Naive Bayes;
\item
  logistic regression on TF--IDF features;
\item
  averaged trainable embeddings with a softmax output layer.
\end{enumerate}

Compare them using accuracy, confusion matrix, validation loss, and
interpretability. Your report should answer:

\begin{itemize}
\tightlist
\item
  Which model has the clearest probabilistic assumptions?
\item
  Which model has the most interpretable features?
\item
  Which model is most sensitive to random initialization?
\item
  Which model is most likely to overfit on a small dataset?
\item
  What mathematical object changes when embeddings are trainable?
\end{itemize}

\section{10.23 Python lab connection}\label{python-lab-connection-4}

The associated lab for this chapter is planned as:

\begin{Shaded}
\begin{Highlighting}[]
\NormalTok{labs/chapter{-}10{-}neural{-}networks{-}for{-}text.ipynb}
\end{Highlighting}
\end{Shaded}

The lab should be written for independent study and should include:

\begin{itemize}
\tightlist
\item
  a review of softmax and cross-entropy;
\item
  dimension exercises for embedding layers;
\item
  a from-scratch averaged-embedding classifier;
\item
  a one-hidden-layer ReLU classifier;
\item
  a comparison with TF--IDF logistic regression;
\item
  a short AI-assisted debugging reflection.
\end{itemize}

\section{10.24 Chapter summary}\label{chapter-summary-3}

A neural text classifier is a trainable composition of maps:

\[
V^T \to \mathbb R^{T\times m} \to \mathbb R^m \to \mathbb R^q \to \mathbb R^K \to \Delta^{K-1}.
\]

The embedding layer is a matrix lookup. Pooling turns token embeddings
into a document representation. Hidden layers apply nonlinear
transformations. Softmax turns logits into class probabilities.
Cross-entropy corresponds to negative conditional log-likelihood. The
gradient \(p-e_y\) links probabilistic prediction to optimization.

The model in this chapter is intentionally simple. It is powerful enough
to introduce trainable representations, but simple enough to analyze by
hand. The next chapters explain how recurrent models and gated models
add explicit sequence dependence.

\section{Notes and references}\label{notes-and-references-5}

This chapter connects classical neural network ideas with text
classification. Useful background sources include Goodfellow et al.
(2016a) for neural networks, Bishop (2006a) for probabilistic
classification and optimization, Jurafsky and Martin (2024) for NLP
context, and Manning et al. (2008) for text representations and
classification.

\chapter{Recurrent Neural Networks}\label{recurrent-neural-networks}

Dynamical systems, sequence likelihood, and backpropagation through time

\hfill\break

A bag-of-words classifier treats a sentence as an unordered collection
of tokens. A recurrent neural network treats a sentence as an ordered
sequence. At each time step, the model reads one token, updates an
internal state, and uses that state to make predictions.

Mathematically, an RNN is a parameterized discrete-time dynamical system
driven by a sequence of input vectors. This viewpoint makes the model
easier to understand: hidden states are states of a dynamical system,
recurrent weights control memory and stability, and backpropagation
through time is differentiation through repeated function composition.

\section{Learning goals}\label{learning-goals-10}

After studying this chapter, you should be able to:

\begin{enumerate}
\def\labelenumi{\arabic{enumi}.}
\tightlist
\item
  explain why order-sensitive language tasks require sequence models
  rather than pure bag-of-words models;
\item
  write a vanilla RNN as a nonlinear recurrence;
\item
  interpret hidden states as summaries of past tokens;
\item
  formulate RNNs for sequence classification and language modeling;
\item
  derive the negative log-likelihood for an RNN language model;
\item
  explain backpropagation through time using products of Jacobian
  matrices;
\item
  connect vanishing and exploding gradients to spectral radius and
  nonlinear saturation;
\item
  implement a small RNN forward pass and training loop in Python;
\item
  use an interactive visualization to study recurrent memory and
  stability.
\end{enumerate}

\section{11.1 Why recurrence?}\label{why-recurrence}

Consider the two short sentences:

\[
\text{not good}
\qquad \text{and} \qquad
\text{good not}.
\]

A bag-of-words representation records the same two tokens in both cases.
Therefore a bag-of-words model cannot distinguish them. More generally,
count vectors ignore the order of tokens, so they cannot directly
represent syntax, negation scope, local dependencies, or temporal
structure.

Let \(V\) be a vocabulary. A sentence of length \(T\) is an ordered
sequence

\[
x=(x_1,x_2,\ldots,x_T) \in V^T.
\]

A sequence model should process the tokens in order. Instead of mapping
the full sequence immediately to one fixed vector, it constructs a
sequence of hidden states

\[
h_1,h_2,\ldots,h_T,
\]

where \(h_t\) is intended to summarize the prefix

\[
(x_1,\ldots,x_t).
\]

The key modeling idea is recursive summarization:

\[
\text{new summary}
=\text{update}(\text{old summary},\text{new token}).
\]

This is exactly the mathematical structure of a discrete-time dynamical
system with external input.

\section{11.2 The vanilla RNN as a nonlinear
recurrence}\label{the-vanilla-rnn-as-a-nonlinear-recurrence}

Let each token \(x_t \in V\) be mapped to an embedding vector

\[
e_t = E_{x_t} \in \mathbb R^m,
\]

where \(E\in \mathbb R^{d\times m}\) is an embedding matrix and
\(d=|V|\).

A vanilla RNN defines hidden states \(h_t\in \mathbb R^r\) by

\[
h_t = \phi(W_h h_{t-1}+W_x e_t+b_h),
\qquad t=1,\ldots,T.
\]

Here:

\begin{itemize}
\tightlist
\item
  \(h_0\in \mathbb R^r\) is an initial state, often the zero vector;
\item
  \(W_h\in \mathbb R^{r\times r}\) is the recurrent weight matrix;
\item
  \(W_x\in \mathbb R^{r\times m}\) maps token embeddings into hidden
  space;
\item
  \(b_h\in \mathbb R^r\) is a bias vector;
\item
  \(\phi\) is a coordinatewise nonlinearity such as \(\tanh\) or ReLU.
\end{itemize}

A common choice is

\[
\phi(u)=\tanh(u),
\]

applied coordinatewise. Then every coordinate of \(h_t\) lies between
\(-1\) and \(1\).

\textbf{Dimension check.} If \(h_{t-1}\in \mathbb R^r\),
\(e_t\in \mathbb R^m\), \(W_h\in \mathbb R^{r\times r}\), and
\(W_x\in \mathbb R^{r\times m}\), then

\[
W_hh_{t-1}+W_xe_t+b_h\in \mathbb R^r.
\]

Therefore \(h_t\in \mathbb R^r\).

The same update rule is applied at every time step. This
\textbf{parameter sharing across time} is essential. It allows the model
to process sequences of different lengths and keeps the number of
parameters independent of \(T\).

\section{11.3 The state-space
viewpoint}\label{the-state-space-viewpoint}

The RNN recurrence can be written as

\[
h_t = F_\theta(h_{t-1},e_t),
\]

where

\[
F_\theta(h,e)=\phi(W_hh+W_xe+b_h).
\]

This is a controlled dynamical system. The input sequence
\(e_1,e_2,\ldots,e_T\) drives the state trajectory

\[
h_0 \to h_1 \to h_2 \to \cdots \to h_T.
\]

For language, the hidden state has an information-processing
interpretation:

\[
h_t \approx \text{a vector summary of the prefix } (x_1,\ldots,x_t).
\]

The hidden state is not guaranteed to contain all information about the
prefix. It is a compressed representation of a possibly long discrete
sequence. The quality of this compression depends on the hidden
dimension \(r\), the parameters, the nonlinearity, the training data,
and the task.

\textbf{RNN as learned sufficient statistic.} In ideal probabilistic
modeling, one might want a statistic \(S_t=S(x_1,\ldots,x_t)\) such that
the future depends on the past only through \(S_t\). An RNN attempts to
learn a finite-dimensional vector \(h_t\) that plays this role for a
prediction task.

\section{11.4 Output maps: many-to-one and many-to-many
models}\label{output-maps-many-to-one-and-many-to-many-models}

The hidden states can be used in different ways.

\subsection{Sequence classification}\label{sequence-classification}

For sentiment classification, topic classification, or spam detection,
the input is a sequence and the output is one label. A common model uses
only the final hidden state:

\[
z = W_o h_T+b_o,
\]

\[
\widehat p(y\mid x)=\operatorname{softmax}(z).
\]

This is called a \textbf{many-to-one} model because many input tokens
produce one output label.

The loss for one training example \((x,y)\) is cross-entropy:

\[
\ell(\theta;x,y)
= -\log \widehat p(y\mid x).
\]

\subsection{Sequence labeling}\label{sequence-labeling-1}

For part-of-speech tagging or named-entity recognition, the model
predicts a label at each time step:

\[
z_t = W_o h_t+b_o,
\]

\[
\widehat p(y_t\mid x_1,\ldots,x_t)=\operatorname{softmax}(z_t).
\]

This is a \textbf{many-to-many} model.

\subsection{Language modeling}\label{language-modeling-1}

For next-token prediction, the model uses the prefix hidden state to
predict the next token:

\[
\widehat p(x_{t+1}=j\mid x_1,\ldots,x_t)
= \operatorname{softmax}(W_o h_t+b_o)_j.
\]

The RNN language model estimates the conditional distribution

\[
P(X_{t+1}\mid X_1,\ldots,X_t).
\]

This is the recurrent analogue of the \(n\)-gram models from Chapter 5.
An \(n\)-gram model keeps a finite symbolic context. An RNN keeps a
finite-dimensional learned vector context.

\section{11.5 RNN language-model
likelihood}\label{rnn-language-model-likelihood}

Let a training sequence be

\[
x=(x_1,x_2,\ldots,x_T).
\]

A language model assigns probability using the chain rule:

\[
P_\theta(x_1,\ldots,x_T)
=\prod_{t=1}^T P_\theta(x_t\mid x_1,\ldots,x_{t-1}).
\]

In an RNN language model, the conditional probability is computed from
the previous hidden state:

\[
p_t=\operatorname{softmax}(W_oh_{t-1}+b_o),
\]

\[
P_\theta(x_t=j\mid x_1,\ldots,x_{t-1})=(p_t)_j.
\]

The hidden state is then updated using the observed token \(x_t\):

\[
h_t=\phi(W_hh_{t-1}+W_xE_{x_t}+b_h).
\]

The negative log-likelihood is

\[
L(\theta;x)
=
-\sum_{t=1}^T \log (p_t)_{x_t}.
\]

The average token-level cross-entropy is

\[
\operatorname{CE}(\theta;x)
=\frac{1}{T}L(\theta;x),
\]

and the perplexity is

\[
\operatorname{PPL}(\theta;x)
=\exp(\operatorname{CE}(\theta;x)).
\]

Thus training an RNN language model is maximum likelihood estimation, or
equivalently empirical risk minimization with cross-entropy loss.

\section{11.6 Unrolling the recurrence}\label{unrolling-the-recurrence}

For a fixed sequence length \(T\), the recurrent computation can be
unrolled into a deep feedforward network with shared weights.

For example,

\[
h_1=F_\theta(h_0,e_1),
\]

\[
h_2=F_\theta(h_1,e_2)=F_\theta(F_\theta(h_0,e_1),e_2),
\]

and so on. Therefore

\[
h_T=F_\theta(F_\theta(\cdots F_\theta(h_0,e_1)\cdots,e_{T-1}),e_T).
\]

The same parameters occur at every layer of this unrolled network. This
is why training RNNs is similar to training deep networks, but with
additional weight sharing.

For sequence classification, the loss may depend only on \(h_T\). For
language modeling, the loss depends on many hidden states:

\[
L=\sum_{t=1}^T \ell_t(h_{t-1}).
\]

In both cases, gradients must flow backward through the entire chain of
recurrent updates.

\section{11.7 Backpropagation through
time}\label{backpropagation-through-time}

Backpropagation through time, or \textbf{BPTT}, is ordinary
backpropagation applied to the unrolled recurrent computation graph.

To see the main issue, consider how the final state depends on an
earlier state. By the chain rule,

\[
\frac{\partial h_T}{\partial h_t}
=
\frac{\partial h_T}{\partial h_{T-1}}
\frac{\partial h_{T-1}}{\partial h_{T-2}}
\cdots
\frac{\partial h_{t+1}}{\partial h_t}.
\]

For the vanilla RNN,

\[
h_s=\phi(a_s),
\qquad
 a_s=W_hh_{s-1}+W_xe_s+b_h.
\]

The one-step Jacobian is

\[
J_s
=
\frac{\partial h_s}{\partial h_{s-1}}
=
D_s W_h,
\]

where

\[
D_s=\operatorname{diag}(\phi'(a_s)).
\]

Therefore

\[
\frac{\partial h_T}{\partial h_t}
=
J_TJ_{T-1}\cdots J_{t+1}.
\]

This product of many matrices is the central mathematical source of
vanishing and exploding gradients.

\textbf{Gradient stability.} If the norms of the Jacobians are typically
less than \(1\), products shrink rapidly. If they are typically greater
than \(1\), products may grow rapidly. Long-range learning in vanilla
RNNs is difficult because gradients involve long products of Jacobian
matrices.

\section{11.8 Vanishing and exploding
gradients}\label{vanishing-and-exploding-gradients}

Suppose, as a simplified model, that the nonlinear derivative is
approximately constant and

\[
J_s \approx J.
\]

Then

\[
J_TJ_{T-1}\cdots J_{t+1}\approx J^{T-t}.
\]

If \(J\) has eigenvalues with magnitudes smaller than \(1\), then
\(J^k\) tends to shrink directions as \(k\) increases. If \(J\) has
eigenvalues with magnitudes larger than \(1\), then \(J^k\) may grow
rapidly in some directions.

For a linear recurrence

\[
h_t=W_hh_{t-1},
\]

the state is

\[
h_T=W_h^T h_0.
\]

If \(W_h\) is diagonalizable as \(W_h=Q\Lambda Q^{-1}\), then

\[
W_h^T=Q\Lambda^TQ^{-1}.
\]

The powers of the eigenvalues determine the long-run behavior. This
connects RNN stability to familiar ideas from linear algebra and
dynamical systems.

For \$ anh\$ RNNs, saturation introduces another source of gradient
shrinkage. Since

\[
\frac{d}{du}\tanh(u)=1-\tanh^2(u),
\]

large values of \(|u|\) produce derivatives close to zero. Even if
\(W_h\) is not small, saturated nonlinearities can make \(D_s\) small.

\section{11.9 Truncated BPTT and gradient
clipping}\label{truncated-bptt-and-gradient-clipping}

For long sequences, full BPTT can be expensive and unstable. Two
practical techniques are common.

\subsection{Truncated backpropagation through
time}\label{truncated-backpropagation-through-time}

Instead of backpropagating through the entire history, one
backpropagates only through a window of length \(K\):

\[
\frac{\partial h_t}{\partial h_{t-K}}.
\]

This reduces computation and memory cost, but limits how directly the
model learns very long-range dependencies.

\subsection{Gradient clipping}\label{gradient-clipping}

When gradients explode, one can rescale the gradient vector \(g\) before
the optimizer update:

\[
g \leftarrow g\cdot \min\left\{1,\frac{c}{\|g\|}\right\}.
\]

Here \(c>0\) is a clipping threshold. Gradient clipping does not solve
all long-memory problems, but it prevents extremely large updates from
destabilizing training.

\section{\texorpdfstring{11.10 RNNs versus \(n\)-gram
models}{11.10 RNNs versus n-gram models}}\label{rnns-versus-n-gram-models}

An \(n\)-gram model approximates

\[
P(x_t\mid x_1,\ldots,x_{t-1})
\approx
P(x_t\mid x_{t-n+1},\ldots,x_{t-1}).
\]

It stores a discrete context of length \(n-1\). An RNN approximates

\[
P(x_t\mid x_1,\ldots,x_{t-1})
\approx
P(x_t\mid h_{t-1}),
\]

where \(h_{t-1}\) is a learned continuous summary of the prefix.

This has advantages and disadvantages.

{\def\LTcaptype{none} % do not increment counter
\begin{longtable}[]{@{}
  >{\raggedright\arraybackslash}p{(\linewidth - 6\tabcolsep) * \real{0.2500}}
  >{\raggedright\arraybackslash}p{(\linewidth - 6\tabcolsep) * \real{0.2500}}
  >{\raggedright\arraybackslash}p{(\linewidth - 6\tabcolsep) * \real{0.2500}}
  >{\raggedright\arraybackslash}p{(\linewidth - 6\tabcolsep) * \real{0.2500}}@{}}
\toprule\noalign{}
\begin{minipage}[b]{\linewidth}\raggedright
Model
\end{minipage} & \begin{minipage}[b]{\linewidth}\raggedright
Context representation
\end{minipage} & \begin{minipage}[b]{\linewidth}\raggedright
Strength
\end{minipage} & \begin{minipage}[b]{\linewidth}\raggedright
Limitation
\end{minipage} \\
\midrule\noalign{}
\endhead
\bottomrule\noalign{}
\endlastfoot
\(n\)-gram & last \(n-1\) tokens & transparent counts & sparse for large
\(n\) \\
RNN & hidden state \(h_t\in \mathbb R^r\) & shares statistical strength
through vectors & difficult long-range optimization \\
\end{longtable}
}

The RNN is more flexible because similar contexts can produce similar
hidden states. However, this flexibility comes from nonlinear
optimization, so the model is harder to analyze and train.

\section{11.11 Worked example: scalar linear
RNN}\label{worked-example-scalar-linear-rnn}

A scalar linear RNN is

\[
h_t=ah_{t-1}+bu_t,
\]

where \(u_t\) is a scalar input. Starting from \(h_0=0\), we obtain

\[
h_1=bu_1,
\]

\[
h_2=abu_1+bu_2,
\]

\[
h_3=a^2bu_1+abu_2+bu_3.
\]

By induction,

\[
h_T=\sum_{s=1}^T a^{T-s}b u_s.
\]

Thus the input from time \(s\) receives weight \(a^{T-s}\). If
\(|a|<1\), older inputs decay. If \(|a|>1\), older inputs may be
amplified. If \(a\approx 1\), the model can maintain information for a
longer time.

This simple formula is one of the cleanest mathematical explanations of
recurrent memory.

\section{11.12 Interactive: recurrent memory and
stability}\label{interactive-recurrent-memory-and-stability}

The following interactive studies the scalar recurrence

\[
h_t = \tanh(a h_{t-1}+b u_t).
\]

The parameter \(a\) plays the role of a recurrent weight. The parameter
\(b\) controls input strength. The input sequence has a pulse near the
beginning, followed by smaller oscillating inputs. Move the sliders to
see how memory, saturation, and stability change.

The HTML version of this chapter contains an interactive Plotly
visualization of the scalar recurrence \(h_t=\tanh(ah_{t-1}+bu_t)\).

\section{11.13 Python example: RNN forward pass from
scratch}\label{python-example-rnn-forward-pass-from-scratch}

The following code implements a small vanilla RNN forward pass for
sequence classification. It is intentionally written with NumPy rather
than a deep learning framework so the linear algebra is visible.

\protect\phantomsection\label{ch11-rnn-forward}
\begin{Shaded}
\begin{Highlighting}[]
\ImportTok{import}\NormalTok{ numpy }\ImportTok{as}\NormalTok{ np}

\NormalTok{rng }\OperatorTok{=}\NormalTok{ np.random.default\_rng(}\DecValTok{7243}\NormalTok{)}

\CommentTok{\# Tiny vocabulary and sequence}
\NormalTok{vocab }\OperatorTok{=}\NormalTok{ \{}\StringTok{"not"}\NormalTok{: }\DecValTok{0}\NormalTok{, }\StringTok{"very"}\NormalTok{: }\DecValTok{1}\NormalTok{, }\StringTok{"good"}\NormalTok{: }\DecValTok{2}\NormalTok{, }\StringTok{"bad"}\NormalTok{: }\DecValTok{3}\NormalTok{, }\StringTok{"movie"}\NormalTok{: }\DecValTok{4}\NormalTok{\}}
\NormalTok{sequence }\OperatorTok{=}\NormalTok{ [}\StringTok{"not"}\NormalTok{, }\StringTok{"very"}\NormalTok{, }\StringTok{"good"}\NormalTok{, }\StringTok{"movie"}\NormalTok{]}
\NormalTok{x }\OperatorTok{=}\NormalTok{ np.array([vocab[w] }\ControlFlowTok{for}\NormalTok{ w }\KeywordTok{in}\NormalTok{ sequence])}

\CommentTok{\# Dimensions}
\NormalTok{d }\OperatorTok{=} \BuiltInTok{len}\NormalTok{(vocab)   }\CommentTok{\# vocabulary size}
\NormalTok{m }\OperatorTok{=} \DecValTok{4}            \CommentTok{\# embedding dimension}
\NormalTok{r }\OperatorTok{=} \DecValTok{5}            \CommentTok{\# hidden dimension}
\NormalTok{K }\OperatorTok{=} \DecValTok{2}            \CommentTok{\# number of classes}

\CommentTok{\# Parameters}
\NormalTok{E }\OperatorTok{=}\NormalTok{ rng.normal(scale}\OperatorTok{=}\FloatTok{0.2}\NormalTok{, size}\OperatorTok{=}\NormalTok{(d, m))}
\NormalTok{W\_h }\OperatorTok{=}\NormalTok{ rng.normal(scale}\OperatorTok{=}\FloatTok{0.3}\NormalTok{, size}\OperatorTok{=}\NormalTok{(r, r))}
\NormalTok{W\_x }\OperatorTok{=}\NormalTok{ rng.normal(scale}\OperatorTok{=}\FloatTok{0.3}\NormalTok{, size}\OperatorTok{=}\NormalTok{(r, m))}
\NormalTok{b\_h }\OperatorTok{=}\NormalTok{ np.zeros(r)}
\NormalTok{W\_o }\OperatorTok{=}\NormalTok{ rng.normal(scale}\OperatorTok{=}\FloatTok{0.3}\NormalTok{, size}\OperatorTok{=}\NormalTok{(K, r))}
\NormalTok{b\_o }\OperatorTok{=}\NormalTok{ np.zeros(K)}

\KeywordTok{def}\NormalTok{ softmax(z):}
\NormalTok{    z }\OperatorTok{=}\NormalTok{ z }\OperatorTok{{-}}\NormalTok{ np.}\BuiltInTok{max}\NormalTok{(z)}
\NormalTok{    ez }\OperatorTok{=}\NormalTok{ np.exp(z)}
    \ControlFlowTok{return}\NormalTok{ ez }\OperatorTok{/}\NormalTok{ ez.}\BuiltInTok{sum}\NormalTok{()}

\KeywordTok{def}\NormalTok{ rnn\_forward(x):}
\NormalTok{    h }\OperatorTok{=}\NormalTok{ np.zeros(r)}
\NormalTok{    states }\OperatorTok{=}\NormalTok{ []}
    \ControlFlowTok{for}\NormalTok{ token\_id }\KeywordTok{in}\NormalTok{ x:}
\NormalTok{        e }\OperatorTok{=}\NormalTok{ E[token\_id]}
\NormalTok{        h }\OperatorTok{=}\NormalTok{ np.tanh(W\_h }\OperatorTok{@}\NormalTok{ h }\OperatorTok{+}\NormalTok{ W\_x }\OperatorTok{@}\NormalTok{ e }\OperatorTok{+}\NormalTok{ b\_h)}
\NormalTok{        states.append(h.copy())}
\NormalTok{    logits }\OperatorTok{=}\NormalTok{ W\_o }\OperatorTok{@}\NormalTok{ h }\OperatorTok{+}\NormalTok{ b\_o}
\NormalTok{    probs }\OperatorTok{=}\NormalTok{ softmax(logits)}
    \ControlFlowTok{return}\NormalTok{ probs, np.array(states)}

\NormalTok{probs, states }\OperatorTok{=}\NormalTok{ rnn\_forward(x)}
\BuiltInTok{print}\NormalTok{(}\StringTok{"Sequence:"}\NormalTok{, sequence)}
\BuiltInTok{print}\NormalTok{(}\StringTok{"Hidden{-}state matrix shape:"}\NormalTok{, states.shape)}
\BuiltInTok{print}\NormalTok{(}\StringTok{"Class probabilities:"}\NormalTok{, probs)}
\end{Highlighting}
\end{Shaded}

\begin{verbatim}
Sequence: ['not', 'very', 'good', 'movie']
Hidden-state matrix shape: (4, 5)
Class probabilities: [0.4960052 0.5039948]
\end{verbatim}

The hidden-state matrix has shape \(T\times r\). Each row is the hidden
state after reading one more token.

\section{11.14 Python example: visualizing hidden-state
trajectories}\label{python-example-visualizing-hidden-state-trajectories}

A hidden-state sequence is a trajectory in \(\mathbb R^r\). We can
project it to two dimensions with SVD to visualize how the RNN state
changes over time.

\begin{Shaded}
\begin{Highlighting}[]
\ImportTok{import}\NormalTok{ matplotlib.pyplot }\ImportTok{as}\NormalTok{ plt}

\CommentTok{\# Center the hidden states and compute a 2D SVD projection}
\NormalTok{H }\OperatorTok{=}\NormalTok{ states }\OperatorTok{{-}}\NormalTok{ states.mean(axis}\OperatorTok{=}\DecValTok{0}\NormalTok{, keepdims}\OperatorTok{=}\VariableTok{True}\NormalTok{)}
\NormalTok{U, S, Vt }\OperatorTok{=}\NormalTok{ np.linalg.svd(H, full\_matrices}\OperatorTok{=}\VariableTok{False}\NormalTok{)}
\NormalTok{coords }\OperatorTok{=}\NormalTok{ H }\OperatorTok{@}\NormalTok{ Vt[:}\DecValTok{2}\NormalTok{].T}

\NormalTok{plt.figure(figsize}\OperatorTok{=}\NormalTok{(}\DecValTok{6}\NormalTok{, }\DecValTok{4}\NormalTok{))}
\NormalTok{plt.plot(coords[:, }\DecValTok{0}\NormalTok{], coords[:, }\DecValTok{1}\NormalTok{], marker}\OperatorTok{=}\StringTok{"o"}\NormalTok{)}
\ControlFlowTok{for}\NormalTok{ i, token }\KeywordTok{in} \BuiltInTok{enumerate}\NormalTok{(sequence):}
\NormalTok{    plt.text(coords[i, }\DecValTok{0}\NormalTok{], coords[i, }\DecValTok{1}\NormalTok{], token)}
\NormalTok{plt.xlabel(}\StringTok{"first SVD direction"}\NormalTok{)}
\NormalTok{plt.ylabel(}\StringTok{"second SVD direction"}\NormalTok{)}
\NormalTok{plt.title(}\StringTok{"Hidden{-}state trajectory of a toy RNN"}\NormalTok{)}
\NormalTok{plt.grid(}\VariableTok{True}\NormalTok{, alpha}\OperatorTok{=}\FloatTok{0.3}\NormalTok{)}
\NormalTok{plt.show()}
\end{Highlighting}
\end{Shaded}

\pandocbounded{\includegraphics[keepaspectratio]{chapters/chapter-11-recurrent-neural-networks_files/figure-pdf/ch11-hidden-trajectory-output-1.pdf}}

The plot should not be overinterpreted. This is only a two-dimensional
projection of a higher-dimensional trajectory. Still, it provides useful
intuition: recurrent models convert a token sequence into a path through
representation space.

\section{11.15 Python example: RNN language-model
probabilities}\label{python-example-rnn-language-model-probabilities}

Now we implement the likelihood for a tiny RNN language model. The model
predicts the next token from the previous hidden state.

\protect\phantomsection\label{ch11-rnn-language-model}
\begin{Shaded}
\begin{Highlighting}[]
\NormalTok{rng }\OperatorTok{=}\NormalTok{ np.random.default\_rng(}\DecValTok{7339}\NormalTok{)}

\NormalTok{corpus }\OperatorTok{=}\NormalTok{ [}
    \StringTok{"math students learn probability"}\NormalTok{,}
    \StringTok{"statistics students learn inference"}\NormalTok{,}
    \StringTok{"language models learn probability"}\NormalTok{,}
    \StringTok{"students study language models"}
\NormalTok{]}

\NormalTok{tokens }\OperatorTok{=} \BuiltInTok{sorted}\NormalTok{(}\BuiltInTok{set}\NormalTok{(}\StringTok{" "}\NormalTok{.join(corpus).split()))}
\NormalTok{vocab }\OperatorTok{=}\NormalTok{ \{w: i }\ControlFlowTok{for}\NormalTok{ i, w }\KeywordTok{in} \BuiltInTok{enumerate}\NormalTok{(tokens)\}}
\NormalTok{id\_to\_word }\OperatorTok{=}\NormalTok{ \{i: w }\ControlFlowTok{for}\NormalTok{ w, i }\KeywordTok{in}\NormalTok{ vocab.items()\}}
\NormalTok{encoded }\OperatorTok{=}\NormalTok{ [[vocab[w] }\ControlFlowTok{for}\NormalTok{ w }\KeywordTok{in}\NormalTok{ sent.split()] }\ControlFlowTok{for}\NormalTok{ sent }\KeywordTok{in}\NormalTok{ corpus]}

\NormalTok{d }\OperatorTok{=} \BuiltInTok{len}\NormalTok{(vocab)}
\NormalTok{m }\OperatorTok{=} \DecValTok{6}
\NormalTok{r }\OperatorTok{=} \DecValTok{8}

\NormalTok{E\_lm }\OperatorTok{=}\NormalTok{ rng.normal(scale}\OperatorTok{=}\FloatTok{0.2}\NormalTok{, size}\OperatorTok{=}\NormalTok{(d, m))}
\NormalTok{W\_h\_lm }\OperatorTok{=}\NormalTok{ rng.normal(scale}\OperatorTok{=}\FloatTok{0.25}\NormalTok{, size}\OperatorTok{=}\NormalTok{(r, r))}
\NormalTok{W\_x\_lm }\OperatorTok{=}\NormalTok{ rng.normal(scale}\OperatorTok{=}\FloatTok{0.25}\NormalTok{, size}\OperatorTok{=}\NormalTok{(r, m))}
\NormalTok{b\_h\_lm }\OperatorTok{=}\NormalTok{ np.zeros(r)}
\NormalTok{W\_o\_lm }\OperatorTok{=}\NormalTok{ rng.normal(scale}\OperatorTok{=}\FloatTok{0.25}\NormalTok{, size}\OperatorTok{=}\NormalTok{(d, r))}
\NormalTok{b\_o\_lm }\OperatorTok{=}\NormalTok{ np.zeros(d)}

\KeywordTok{def}\NormalTok{ lm\_sequence\_nll(seq):}
\NormalTok{    h }\OperatorTok{=}\NormalTok{ np.zeros(r)}
\NormalTok{    nll }\OperatorTok{=} \FloatTok{0.0}
\NormalTok{    count }\OperatorTok{=} \DecValTok{0}
    \ControlFlowTok{for}\NormalTok{ token\_id }\KeywordTok{in}\NormalTok{ seq:}
\NormalTok{        logits }\OperatorTok{=}\NormalTok{ W\_o\_lm }\OperatorTok{@}\NormalTok{ h }\OperatorTok{+}\NormalTok{ b\_o\_lm}
\NormalTok{        p }\OperatorTok{=}\NormalTok{ softmax(logits)}
\NormalTok{        nll }\OperatorTok{{-}=}\NormalTok{ np.log(p[token\_id] }\OperatorTok{+} \FloatTok{1e{-}12}\NormalTok{)}
\NormalTok{        count }\OperatorTok{+=} \DecValTok{1}
\NormalTok{        e }\OperatorTok{=}\NormalTok{ E\_lm[token\_id]}
\NormalTok{        h }\OperatorTok{=}\NormalTok{ np.tanh(W\_h\_lm }\OperatorTok{@}\NormalTok{ h }\OperatorTok{+}\NormalTok{ W\_x\_lm }\OperatorTok{@}\NormalTok{ e }\OperatorTok{+}\NormalTok{ b\_h\_lm)}
    \ControlFlowTok{return}\NormalTok{ nll, count}

\NormalTok{total\_nll }\OperatorTok{=} \FloatTok{0.0}
\NormalTok{total\_tokens }\OperatorTok{=} \DecValTok{0}
\ControlFlowTok{for}\NormalTok{ seq }\KeywordTok{in}\NormalTok{ encoded:}
\NormalTok{    nll, count }\OperatorTok{=}\NormalTok{ lm\_sequence\_nll(seq)}
\NormalTok{    total\_nll }\OperatorTok{+=}\NormalTok{ nll}
\NormalTok{    total\_tokens }\OperatorTok{+=}\NormalTok{ count}

\NormalTok{cross\_entropy }\OperatorTok{=}\NormalTok{ total\_nll }\OperatorTok{/}\NormalTok{ total\_tokens}
\NormalTok{perplexity }\OperatorTok{=}\NormalTok{ np.exp(cross\_entropy)}
\BuiltInTok{print}\NormalTok{(}\StringTok{"Vocabulary:"}\NormalTok{, tokens)}
\BuiltInTok{print}\NormalTok{(}\StringTok{"Average cross{-}entropy:"}\NormalTok{, }\BuiltInTok{round}\NormalTok{(cross\_entropy, }\DecValTok{3}\NormalTok{))}
\BuiltInTok{print}\NormalTok{(}\StringTok{"Perplexity:"}\NormalTok{, }\BuiltInTok{round}\NormalTok{(perplexity, }\DecValTok{3}\NormalTok{))}
\end{Highlighting}
\end{Shaded}

\begin{verbatim}
Vocabulary: ['inference', 'language', 'learn', 'math', 'models', 'probability', 'statistics', 'students', 'study']
Average cross-entropy: 2.232
Perplexity: 9.315
\end{verbatim}

The parameters above are random, so the model is not expected to be
good. The purpose is to show the exact computation of likelihood,
cross-entropy, and perplexity.

\section{11.16 Python example: one training step by finite
differences}\label{python-example-one-training-step-by-finite-differences}

Full BPTT is lengthy to implement from scratch. To check our
understanding, we can examine a one-parameter finite-difference
derivative for the recurrent weight matrix.

\protect\phantomsection\label{ch11-finite-difference-gradient}
\begin{Shaded}
\begin{Highlighting}[]
\CommentTok{\# Use one sequence and one entry of W\_h\_lm.}
\NormalTok{seq }\OperatorTok{=}\NormalTok{ encoded[}\DecValTok{0}\NormalTok{]}

\KeywordTok{def}\NormalTok{ loss\_with\_current\_params():}
\NormalTok{    nll, \_ }\OperatorTok{=}\NormalTok{ lm\_sequence\_nll(seq)}
    \ControlFlowTok{return}\NormalTok{ nll}

\NormalTok{base\_loss }\OperatorTok{=}\NormalTok{ loss\_with\_current\_params()}

\NormalTok{i, j }\OperatorTok{=} \DecValTok{0}\NormalTok{, }\DecValTok{0}
\NormalTok{eps }\OperatorTok{=} \FloatTok{1e{-}5}
\NormalTok{old\_value }\OperatorTok{=}\NormalTok{ W\_h\_lm[i, j]}

\NormalTok{W\_h\_lm[i, j] }\OperatorTok{=}\NormalTok{ old\_value }\OperatorTok{+}\NormalTok{ eps}
\NormalTok{loss\_plus }\OperatorTok{=}\NormalTok{ loss\_with\_current\_params()}

\NormalTok{W\_h\_lm[i, j] }\OperatorTok{=}\NormalTok{ old\_value }\OperatorTok{{-}}\NormalTok{ eps}
\NormalTok{loss\_minus }\OperatorTok{=}\NormalTok{ loss\_with\_current\_params()}

\NormalTok{W\_h\_lm[i, j] }\OperatorTok{=}\NormalTok{ old\_value}
\NormalTok{finite\_difference }\OperatorTok{=}\NormalTok{ (loss\_plus }\OperatorTok{{-}}\NormalTok{ loss\_minus) }\OperatorTok{/}\NormalTok{ (}\DecValTok{2} \OperatorTok{*}\NormalTok{ eps)}

\BuiltInTok{print}\NormalTok{(}\StringTok{"Base loss:"}\NormalTok{, }\BuiltInTok{round}\NormalTok{(base\_loss, }\DecValTok{4}\NormalTok{))}
\BuiltInTok{print}\NormalTok{(}\StringTok{"Finite{-}difference derivative for W\_h[0,0]:"}\NormalTok{, }\BuiltInTok{round}\NormalTok{(finite\_difference, }\DecValTok{4}\NormalTok{))}
\end{Highlighting}
\end{Shaded}

\begin{verbatim}
Base loss: 8.8042
Finite-difference derivative for W_h[0,0]: 0.0139
\end{verbatim}

Finite differences are too slow for training real RNNs, but they are
useful for verifying gradient code. BPTT computes all such derivatives
efficiently by reusing intermediate Jacobian-vector products.

\section{11.17 Python example: gradient norm through
time}\label{python-example-gradient-norm-through-time}

The following simulation illustrates the product-of-Jacobians phenomenon
in a linear recurrence. We compute powers of a recurrent matrix and
observe how their norms behave.

\protect\phantomsection\label{ch11-gradient-norm-simulation}
\begin{Shaded}
\begin{Highlighting}[]
\NormalTok{rng }\OperatorTok{=}\NormalTok{ np.random.default\_rng(}\DecValTok{2026}\NormalTok{)}

\NormalTok{r }\OperatorTok{=} \DecValTok{6}
\NormalTok{A }\OperatorTok{=}\NormalTok{ rng.normal(size}\OperatorTok{=}\NormalTok{(r, r))}
\CommentTok{\# Normalize A and create versions with different spectral scales.}
\NormalTok{eig\_max }\OperatorTok{=} \BuiltInTok{max}\NormalTok{(}\BuiltInTok{abs}\NormalTok{(np.linalg.eigvals(A)))}
\NormalTok{A\_base }\OperatorTok{=}\NormalTok{ A }\OperatorTok{/}\NormalTok{ eig\_max}

\NormalTok{scales }\OperatorTok{=}\NormalTok{ [}\FloatTok{0.7}\NormalTok{, }\FloatTok{0.95}\NormalTok{, }\FloatTok{1.1}\NormalTok{]}
\NormalTok{steps }\OperatorTok{=}\NormalTok{ np.arange(}\DecValTok{1}\NormalTok{, }\DecValTok{31}\NormalTok{)}

\NormalTok{norms }\OperatorTok{=}\NormalTok{ \{\}}
\ControlFlowTok{for}\NormalTok{ scale }\KeywordTok{in}\NormalTok{ scales:}
\NormalTok{    M }\OperatorTok{=}\NormalTok{ scale }\OperatorTok{*}\NormalTok{ A\_base}
\NormalTok{    values }\OperatorTok{=}\NormalTok{ []}
\NormalTok{    P }\OperatorTok{=}\NormalTok{ np.eye(r)}
    \ControlFlowTok{for}\NormalTok{ \_ }\KeywordTok{in}\NormalTok{ steps:}
\NormalTok{        P }\OperatorTok{=}\NormalTok{ M }\OperatorTok{@}\NormalTok{ P}
\NormalTok{        values.append(np.linalg.norm(P, }\BuiltInTok{ord}\OperatorTok{=}\DecValTok{2}\NormalTok{))}
\NormalTok{    norms[scale] }\OperatorTok{=}\NormalTok{ values}

\ControlFlowTok{for}\NormalTok{ scale, values }\KeywordTok{in}\NormalTok{ norms.items():}
    \BuiltInTok{print}\NormalTok{(}\StringTok{"scale"}\NormalTok{, scale, }\StringTok{"first norm"}\NormalTok{, }\BuiltInTok{round}\NormalTok{(values[}\DecValTok{0}\NormalTok{], }\DecValTok{3}\NormalTok{), }\StringTok{"last norm"}\NormalTok{, }\BuiltInTok{round}\NormalTok{(values[}\OperatorTok{{-}}\DecValTok{1}\NormalTok{], }\DecValTok{3}\NormalTok{))}
\end{Highlighting}
\end{Shaded}

\begin{verbatim}
scale 0.7 first norm 1.118 last norm 0.0
scale 0.95 first norm 1.517 last norm 0.264
scale 1.1 first norm 1.757 last norm 21.454
\end{verbatim}

When the effective recurrent scale is below \(1\), powers tend to
shrink. When it is above \(1\), powers may grow. Nonlinearities and
data-dependent Jacobians complicate this picture, but the linear example
captures the core phenomenon.

\section{11.18 AI-assisted learning
component}\label{ai-assisted-learning-component-4}

Use an AI assistant as a mathematical critic rather than as a
replacement for derivation. Good prompts include:

\begin{enumerate}
\def\labelenumi{\arabic{enumi}.}
\tightlist
\item
  \textbf{Dimension audit.} ``Given the RNN equations
  \(h_t=\tanh(W_hh_{t-1}+W_xe_t+b_h)\) and \(z_t=W_oh_t+b_o\), check
  every matrix dimension and identify possible shape errors.''
\item
  \textbf{Gradient explanation.} ``Explain why products of Jacobian
  matrices can cause vanishing gradients in vanilla RNNs. Use a
  two-dimensional linear example.''
\item
  \textbf{Model comparison.} ``Compare an \(n\)-gram model and an RNN
  language model using conditional probability notation.''
\item
  \textbf{Code review.} ``Review this NumPy RNN forward pass and
  identify whether the hidden state is updated before or after
  next-token prediction.''
\item
  \textbf{Interpretability caution.} ``Give three reasons why a
  two-dimensional projection of hidden states should not be interpreted
  as the true geometry of the RNN.''
\end{enumerate}

The student should verify every answer by checking equations,
dimensions, and small numerical examples.

\section{11.19 Common pitfalls}\label{common-pitfalls-2}

\begin{enumerate}
\def\labelenumi{\arabic{enumi}.}
\tightlist
\item
  \textbf{Confusing token embeddings with hidden states.} The embedding
  \(e_t\) represents the current token. The hidden state \(h_t\)
  represents the prefix up to time \(t\).
\item
  \textbf{Forgetting parameter sharing.} The same recurrent weights are
  used at every time step.
\item
  \textbf{Updating in the wrong order for language modeling.} To predict
  \(x_t\), the model should use information from \(x_1,\ldots,x_{t-1}\),
  not from \(x_t\) itself.
\item
  \textbf{Ignoring gradient products.} Long-range learning depends on
  products of Jacobians.
\item
  \textbf{Assuming RNNs solve all sequence problems.} Vanilla RNNs often
  struggle with long dependencies. LSTM, GRU, and attention mechanisms
  were developed partly to address this problem.
\end{enumerate}

\section{11.20 Summary}\label{summary-1}

A recurrent neural network is a sequence model built from a repeated
update rule

\[
h_t=\phi(W_hh_{t-1}+W_xe_t+b_h).
\]

The hidden state is a learned continuous summary of the prefix. For
classification, the final state can feed a softmax classifier. For
language modeling, each hidden state predicts the next token
distribution.

The central mathematical challenge is that training requires
differentiating through many repeated compositions. Gradients contain
products of Jacobian matrices, which can vanish or explode. This
explains why vanilla RNNs are historically important but limited for
long-context language tasks. The next chapter studies gated recurrent
models, which modify the recurrence to control memory more explicitly.

\section{Exercises}\label{exercises-10}

\begin{enumerate}
\def\labelenumi{\arabic{enumi}.}
\item
  \textbf{Dimension check.} Let \(E\in \mathbb R^{10000\times 128}\),
  \(W_h\in \mathbb R^{256\times 256}\), and
  \(W_x\in \mathbb R^{256\times 128}\). What is the dimension of
  \(h_t\)? What is the dimension of \(W_xE_{x_t}\)?
\item
  \textbf{Bag-of-words failure.} Construct two sentences with the same
  bag-of-words vector but different meanings. Explain why a recurrent
  model can distinguish them in principle.
\item
  \textbf{Scalar recurrence.} For \(h_t=ah_{t-1}+bu_t\) with \(h_0=0\),
  prove that \[
  h_T=\sum_{s=1}^T a^{T-s}bu_s.
  \]
\item
  \textbf{Language-model likelihood.} Write the negative log-likelihood
  for the sequence \((x_1,x_2,x_3,x_4)\) under an RNN language model.
\item
  \textbf{Jacobian product.} For \(h_t=\tanh(W_hh_{t-1}+W_xe_t+b_h)\),
  derive the one-step Jacobian \(\partial h_t/\partial h_{t-1}\).
\item
  \textbf{Vanishing gradient.} Suppose \(J_s=0.8I\) for all \(s\).
  Compute \(J_TJ_{T-1}\cdots J_{t+1}\) and explain what happens as
  \(T-t\) grows.
\item
  \textbf{Exploding gradient.} Suppose \(J_s=1.2I\) for all \(s\).
  Repeat the previous exercise.
\item
  \textbf{RNN versus bigram model.} Compare the assumptions behind a
  bigram language model and an RNN language model.
\item
  \textbf{Hidden-state trajectory.} Modify the Python code to compare
  two sentences with the same words in different orders. Plot both
  hidden-state trajectories.
\item
  \textbf{Training experiment.} Implement a tiny RNN classifier in a
  deep learning framework and compare it with a bag-of-words classifier
  on a small sentiment dataset.
\end{enumerate}

\section{Mini-project: Order sensitivity in text
classification}\label{mini-project-order-sensitivity-in-text-classification}

Build a small dataset where word order matters. For example, generate
short phrases using templates such as:

\[
\text{not good},\quad \text{not bad},\quad \text{very good},\quad \text{very bad}.
\]

Compare three models:

\begin{enumerate}
\def\labelenumi{\arabic{enumi}.}
\tightlist
\item
  a bag-of-words logistic classifier;
\item
  an averaged-embedding neural classifier;
\item
  a vanilla RNN classifier.
\end{enumerate}

Report which models can distinguish negation and why. Your report should
include equations for each model, not only accuracy scores.

\section{Lab connection}\label{lab-connection-5}

The accompanying lab for this chapter should guide students through:

\begin{enumerate}
\def\labelenumi{\arabic{enumi}.}
\tightlist
\item
  implementing a vanilla RNN forward pass from scratch;
\item
  computing language-model negative log-likelihood;
\item
  visualizing hidden-state trajectories;
\item
  training a small RNN classifier using a modern deep learning library;
\item
  studying vanishing and exploding gradients numerically;
\item
  comparing RNNs with bag-of-words and \(n\)-gram baselines.
\end{enumerate}

\chapter{LSTM and GRU Models}\label{lstm-and-gru-models}

Gated recurrences, memory paths, and long-range dependence

\hfill\break

A vanilla recurrent neural network updates its hidden state by
repeatedly applying the same nonlinear map. This gives a clean
dynamical-system model for sequences, but it also creates a mathematical
difficulty: gradients must pass through long products of Jacobian
matrices. These products may shrink or explode.

Long short-term memory networks and gated recurrent units modify the
recurrence by adding \textbf{gates}. A gate is a vector with coordinates
between \(0\) and \(1\). It controls how much information is kept,
forgotten, written, or revealed. From a mathematical point of view,
gated recurrent models are controlled nonlinear dynamical systems with
coordinatewise interpolation. Their main purpose is to create more
stable paths for information and gradients across long sequences.

\section{Learning goals}\label{learning-goals-11}

After studying this chapter, you should be able to:

\begin{enumerate}
\def\labelenumi{\arabic{enumi}.}
\tightlist
\item
  explain why vanilla RNNs have difficulty learning long-range
  dependencies;
\item
  define sigmoid gates and interpret them as coordinatewise
  interpolation weights;
\item
  write the LSTM equations and check all dimensions;
\item
  interpret the LSTM cell state as a controlled memory variable;
\item
  explain how forget gates create a direct memory path through time;
\item
  write the GRU equations and compare them with LSTM equations;
\item
  connect gate values to gradient flow and effective memory length;
\item
  implement LSTM and GRU forward passes from scratch in Python;
\item
  use an interactive visualization to compare ungated and gated memory.
\end{enumerate}

\section{12.1 The limitation of a vanilla
RNN}\label{the-limitation-of-a-vanilla-rnn}

In Chapter 11, a vanilla RNN was written as

\[
h_t = \phi(W_hh_{t-1}+W_xe_t+b_h),
\]

where \(e_t\) is the embedding of token \(x_t\) and \(h_t\) is the
hidden state. This recurrence is mathematically natural, but it has a
difficult optimization structure.

To see the issue, consider a scalar linearized recurrence

\[
h_t = a h_{t-1} + b u_t.
\]

Starting from \(h_0=0\), we have

\[
h_T = \sum_{s=1}^T a^{T-s}b u_s.
\]

The contribution of an earlier input \(u_s\) to the final state \(h_T\)
is multiplied by \(a^{T-s}\). Therefore:

\begin{itemize}
\tightlist
\item
  if \(|a|<1\), old information decays exponentially;
\item
  if \(|a|>1\), old information may explode;
\item
  if \(|a|\approx 1\), old information may persist.
\end{itemize}

The same idea appears in nonlinear vector RNNs. Differentiating through
time gives products of Jacobians. A typical factor has the form

\[
J_t = \frac{\partial h_t}{\partial h_{t-1}}.
\]

The derivative from time \(s\) to time \(T\) contains a product

\[
\frac{\partial h_T}{\partial h_s}
= J_TJ_{T-1}\cdots J_{s+1}.
\]

Long products of matrices are numerically delicate. Their norms can
become very small or very large. This is the mathematical source of the
vanishing and exploding gradient problem.

\textbf{Core difficulty.} A vanilla RNN asks one nonlinear state vector
\(h_t\) to do two jobs at once:

\begin{enumerate}
\def\labelenumi{\arabic{enumi}.}
\tightlist
\item
  represent the current useful summary of the prefix;
\item
  carry information and gradients across many time steps.
\end{enumerate}

LSTM and GRU models separate these roles more carefully by using gates.

\section{12.2 Gates as coordinatewise control
variables}\label{gates-as-coordinatewise-control-variables}

The basic mathematical object in an LSTM or GRU is a \textbf{gate}. A
gate is a vector

\[
g_t \in (0,1)^r
\]

computed from the current input and previous hidden state:

\[
g_t = \sigma(W_ge_t+U_gh_{t-1}+b_g).
\]

Here:

\begin{itemize}
\tightlist
\item
  \(e_t\in R^m\) is the input embedding;
\item
  \(h_{t-1}\in R^r\) is the previous hidden state;
\item
  \(W_g\in R^{r\times m}\);
\item
  \(U_g\in R^{r\times r}\);
\item
  \(b_g\in R^r\);
\item
  \(\sigma\) is the logistic sigmoid function applied coordinatewise.
\end{itemize}

The sigmoid function is

\[
\sigma(z)=\frac{1}{1+e^{-z}}.
\]

Thus every coordinate satisfies

\[
0 < (g_t)_j < 1.
\]

This makes a gate useful for interpolation. For two vectors
\(a,b\in R^r\), the expression

\[
g_t\odot a + (1-g_t)\odot b
\]

is a coordinatewise convex combination of \(a\) and \(b\). The symbol
\(\odot\) denotes coordinatewise multiplication.

\textbf{Gate interpretation.} In coordinate \(j\),

\[
(g_t)_j a_j + (1-(g_t)_j)b_j
\]

lies between \(a_j\) and \(b_j\). If \((g_t)_j\approx 1\), the
coordinate mostly chooses \(a_j\). If \((g_t)_j\approx 0\), it mostly
chooses \(b_j\).

This simple interpolation idea is the mathematical heart of gated
recurrent networks.

\section{12.3 A minimal gated-memory
recurrence}\label{a-minimal-gated-memory-recurrence}

Before writing the full LSTM, consider a simplified memory recurrence:

\[
c_t = f_t\odot c_{t-1} + i_t\odot \widetilde c_t.
\]

Here:

\begin{itemize}
\tightlist
\item
  \(c_t\in R^r\) is a memory state;
\item
  \(f_t\in (0,1)^r\) is a forget gate;
\item
  \(i_t\in (0,1)^r\) is an input gate;
\item
  \(\widetilde c_t\in R^r\) is a candidate update.
\end{itemize}

The recurrence has two terms:

\[
\text{old memory kept} = f_t\odot c_{t-1},
\]

\[
\text{new information written} = i_t\odot \widetilde c_t.
\]

If a coordinate of \(f_t\) is close to \(1\), then the corresponding
coordinate of \(c_{t-1}\) is copied forward almost unchanged. If it is
close to \(0\), that coordinate is mostly erased.

Suppose for a moment that \(f_t=f\) is a constant scalar and the
candidate update is zero after time \(s\). Then the memory contribution
from time \(s\) to time \(T\) is multiplied by

\[
f^{T-s}.
\]

This resembles the scalar RNN factor \(a^{T-s}\), but with an important
difference: the model learns the value of \(f_t\) dynamically from the
data. It can choose long memory for some coordinates and short memory
for others.

More generally, the direct derivative of the cell state with respect to
the previous cell state is approximately

\[
\frac{\partial c_t}{\partial c_{t-1}} \approx \operatorname{diag}(f_t),
\]

when we focus on the direct path and ignore indirect dependence of the
gates on earlier states. Therefore the memory path from time \(s\) to
time \(T\) contains factors

\[
\operatorname{diag}(f_T)\operatorname{diag}(f_{T-1})\cdots \operatorname{diag}(f_{s+1}).
\]

Coordinatewise, this is a product of forget-gate values. This is much
easier to interpret than a product of arbitrary recurrent Jacobian
matrices.

\section{12.4 The LSTM equations}\label{the-lstm-equations}

An LSTM has two state vectors:

\[
h_t\in R^r,
\qquad
c_t\in R^r.
\]

The vector \(h_t\) is the hidden state exposed to the output layer. The
vector \(c_t\) is the memory cell state.

Given input embedding \(e_t\in R^m\), previous hidden state
\(h_{t-1}\in R^r\), and previous cell state \(c_{t-1}\in R^r\), the LSTM
computes

\[
f_t = \sigma(W_fe_t+U_fh_{t-1}+b_f),
\]

\[
i_t = \sigma(W_ie_t+U_ih_{t-1}+b_i),
\]

\[
o_t = \sigma(W_oe_t+U_oh_{t-1}+b_o),
\]

\[
\widetilde c_t = \tanh(W_ce_t+U_ch_{t-1}+b_c).
\]

Then it updates

\[
c_t = f_t\odot c_{t-1} + i_t\odot \widetilde c_t,
\]

and exposes

\[
h_t = o_t\odot \tanh(c_t).
\]

The three gates have standard names:

{\def\LTcaptype{none} % do not increment counter
\begin{longtable}[]{@{}lll@{}}
\toprule\noalign{}
Gate & Formula & Main role \\
\midrule\noalign{}
\endhead
\bottomrule\noalign{}
\endlastfoot
forget gate & \(f_t\) & how much previous memory to keep \\
input gate & \(i_t\) & how much new candidate memory to write \\
output gate & \(o_t\) & how much memory to reveal as hidden state \\
\end{longtable}
}

The candidate \(\widetilde c_t\) proposes new content. The input gate
decides how much of that content enters memory.

\section{12.5 Dimension check for LSTM}\label{dimension-check-for-lstm}

Let:

\[
e_t\in R^m,
\qquad
h_t,c_t\in R^r.
\]

For each gate \(g\in\{f,i,o\}\), the parameters satisfy

\[
W_g\in R^{r\times m},
\qquad
U_g\in R^{r\times r},
\qquad
b_g\in R^r.
\]

For the candidate memory,

\[
W_c\in R^{r\times m},
\qquad
U_c\in R^{r\times r},
\qquad
b_c\in R^r.
\]

Therefore

\[
W_ge_t+U_gh_{t-1}+b_g\in R^r,
\]

so each gate is in \(R^r\). The update

\[
c_t = f_t\odot c_{t-1}+i_t\odot \widetilde c_t
\]

is a sum of two vectors in \(R^r\). Hence \(c_t\in R^r\). Finally,

\[
h_t=o_t\odot \tanh(c_t)\in R^r.
\]

\textbf{Parameter count.} Each of the four affine maps has \(rm+r^2+r\)
parameters. Therefore one LSTM layer has

\[
4(rm+r^2+r)
\]

recurrent parameters, not counting the token embedding matrix and output
layer.

\section{12.6 Why the LSTM helps gradient
flow}\label{why-the-lstm-helps-gradient-flow}

The important mathematical feature is the additive cell-state update

\[
c_t = f_t\odot c_{t-1}+i_t\odot \widetilde c_t.
\]

This creates a direct path from \(c_{t-1}\) to \(c_t\). Along this path,
the derivative is governed by the forget gate:

\[
\frac{\partial c_t}{\partial c_{t-1}} \approx \operatorname{diag}(f_t).
\]

Therefore, for coordinate \(j\), the direct memory derivative from time
\(s\) to time \(T\) is approximately

\[
\prod_{t=s+1}^T (f_t)_j.
\]

This formula gives a clean interpretation of memory length.

If \((f_t)_j\approx 0.5\), then memory decays quickly:

\[
0.5^{20}\approx 9.5\times 10^{-7}.
\]

If \((f_t)_j\approx 0.95\), then memory decays slowly:

\[
0.95^{20}\approx 0.36.
\]

So a gate coordinate near \(1\) can preserve information over many
steps. A gate coordinate near \(0\) can deliberately forget.

\textbf{Effective memory length.} If the forget gate is approximately
constant with value \(f\in(0,1)\), then the contribution after \(k\)
steps is roughly \(f^k\). The half-life \(k_{1/2}\) solves

\[
f^{k_{1/2}}=\frac12.
\]

Thus

\[
k_{1/2}=\frac{\log(1/2)}{\log f}.
\]

For \(f=0.9\), the half-life is about \(6.58\) steps. For \(f=0.99\), it
is about \(68.97\) steps.

\section{12.7 Statistical view: adaptive state-space
modeling}\label{statistical-view-adaptive-state-space-modeling}

A gated recurrent model can be viewed as a parameterized state-space
model for sequences. At time \(t\), it computes latent states

\[
(h_t,c_t)=F_\theta(h_{t-1},c_{t-1},e_t).
\]

For classification, it defines

\[
\widehat p_\theta(y\mid x_1,\ldots,x_T)
=\operatorname{softmax}(W_yh_T+b_y).
\]

For language modeling, it defines

\[
\widehat p_\theta(x_{t+1}\mid x_1,
\ldots,x_t)
=\operatorname{softmax}(W_yh_t+b_y).
\]

Given a training set

\[
D=\{(x^{(n)},y^{(n)})\}_{n=1}^N,
\]

a sequence classifier is usually trained by minimizing empirical
cross-entropy:

\[
\widehat R(\theta)
=\frac1N\sum_{n=1}^N
\left[-\log \widehat p_\theta(y^{(n)}\mid x^{(n)})\right].
\]

For language modeling, the negative log-likelihood is

\[
L(\theta)
= -\sum_{n=1}^N\sum_{t=1}^{T_n}
\log \widehat p_\theta(x_t^{(n)}\mid x_1^{(n)},\ldots,x_{t-1}^{(n)}).
\]

Thus LSTM and GRU models are not mysterious from a statistical point of
view. They are flexible parameterized conditional probability models.
The gates are part of the parameterized representation map.

\section{12.8 The GRU equations}\label{the-gru-equations}

A gated recurrent unit is a simpler gated model. It uses one state
vector \(h_t\in R^r\) rather than separate hidden and cell states.

Given input embedding \(e_t\in R^m\) and previous state
\(h_{t-1}\in R^r\), a GRU computes

\[
z_t = \sigma(W_ze_t+U_zh_{t-1}+b_z),
\]

\[
r_t = \sigma(W_re_t+U_rh_{t-1}+b_r),
\]

\[
\widetilde h_t = \tanh(W_he_t+U_h(r_t\odot h_{t-1})+b_h).
\]

Then it updates

\[
h_t = (1-z_t)\odot h_{t-1}+z_t\odot \widetilde h_t.
\]

The two gates are:

{\def\LTcaptype{none} % do not increment counter
\begin{longtable}[]{@{}
  >{\raggedright\arraybackslash}p{(\linewidth - 4\tabcolsep) * \real{0.3333}}
  >{\raggedright\arraybackslash}p{(\linewidth - 4\tabcolsep) * \real{0.3333}}
  >{\raggedright\arraybackslash}p{(\linewidth - 4\tabcolsep) * \real{0.3333}}@{}}
\toprule\noalign{}
\begin{minipage}[b]{\linewidth}\raggedright
Gate
\end{minipage} & \begin{minipage}[b]{\linewidth}\raggedright
Formula
\end{minipage} & \begin{minipage}[b]{\linewidth}\raggedright
Main role
\end{minipage} \\
\midrule\noalign{}
\endhead
\bottomrule\noalign{}
\endlastfoot
update gate & \(z_t\) & controls interpolation between old state and
candidate state \\
reset gate & \(r_t\) & controls how much old state is used in the
candidate \\
\end{longtable}
}

The update gate \(z_t\) is especially important. If
\((z_t)_j\approx 0\), then

\[
(h_t)_j\approx (h_{t-1})_j.
\]

If \((z_t)_j\approx 1\), then

\[
(h_t)_j\approx (\widetilde h_t)_j.
\]

Thus the GRU also creates a direct copying path through time.

\section{12.9 LSTM versus GRU}\label{lstm-versus-gru}

The LSTM has a separate memory cell \(c_t\) and hidden state \(h_t\).
The GRU combines memory and hidden state into one vector \(h_t\).

{\def\LTcaptype{none} % do not increment counter
\begin{longtable}[]{@{}
  >{\raggedright\arraybackslash}p{(\linewidth - 6\tabcolsep) * \real{0.2143}}
  >{\raggedleft\arraybackslash}p{(\linewidth - 6\tabcolsep) * \real{0.2857}}
  >{\raggedleft\arraybackslash}p{(\linewidth - 6\tabcolsep) * \real{0.2857}}
  >{\raggedright\arraybackslash}p{(\linewidth - 6\tabcolsep) * \real{0.2143}}@{}}
\toprule\noalign{}
\begin{minipage}[b]{\linewidth}\raggedright
Model
\end{minipage} & \begin{minipage}[b]{\linewidth}\raggedleft
States
\end{minipage} & \begin{minipage}[b]{\linewidth}\raggedleft
Gates
\end{minipage} & \begin{minipage}[b]{\linewidth}\raggedright
Main memory path
\end{minipage} \\
\midrule\noalign{}
\endhead
\bottomrule\noalign{}
\endlastfoot
vanilla RNN & \(h_t\) & none & nonlinear recurrence \\
LSTM & \(h_t,c_t\) & forget, input, output & cell-state update
\(c_t=f_t\odot c_{t-1}+i_t\odot\widetilde c_t\) \\
GRU & \(h_t\) & update, reset & interpolation
\(h_t=(1-z_t)\odot h_{t-1}+z_t\odot\widetilde h_t\) \\
\end{longtable}
}

There is no universal rule that one is always better. The LSTM is more
expressive and exposes a separate memory state. The GRU has fewer gates
and often fewer parameters, which can make it easier to train on smaller
data. Both are mathematically designed to improve long-range credit
assignment relative to a vanilla RNN.

\section{12.10 Worked example: memory
half-life}\label{worked-example-memory-half-life}

Suppose a scalar LSTM memory coordinate follows

\[
c_t=f c_{t-1}
\]

with no new input after time \(0\). If \(c_0=1\), then

\[
c_t=f^t.
\]

The half-life is the number of steps until \(c_t=1/2\):

\[
f^t=\frac12.
\]

Taking logs gives

\[
t=\frac{\log(1/2)}{\log f}.
\]

For \(f=0.8\),

\[
t\approx \frac{-0.6931}{-0.2231}\approx 3.11.
\]

For \(f=0.95\),

\[
t\approx \frac{-0.6931}{-0.0513}\approx 13.51.
\]

For \(f=0.99\),

\[
t\approx \frac{-0.6931}{-0.01005}\approx 68.97.
\]

This calculation gives an intuitive numerical meaning to the forget
gate.

\section{12.11 Interactive: gated memory versus vanilla
recurrence}\label{interactive-gated-memory-versus-vanilla-recurrence}

The following interactive compares a scalar vanilla recurrence with a
simplified scalar LSTM memory recurrence.

Vanilla recurrence:

\[
h_t=\tanh(a h_{t-1}+b u_t).
\]

Gated memory recurrence:

\[
c_t=f c_{t-1}+i\tanh(bu_t),
\qquad
h_t=o\tanh(c_t).
\]

Move the sliders to see how the forget gate \(f\) controls memory
persistence.

The HTML version of this chapter contains an interactive Plotly
visualization comparing a scalar vanilla recurrence with a simplified
scalar LSTM memory cell.

\section{12.12 Python example: LSTM forward pass from
scratch}\label{python-example-lstm-forward-pass-from-scratch}

The following NumPy example implements one LSTM layer for a short token
sequence. The goal is not to build a production model. The goal is to
make every matrix and vector visible.

\protect\phantomsection\label{ch12-lstm-forward}
\begin{Shaded}
\begin{Highlighting}[]
\ImportTok{import}\NormalTok{ numpy }\ImportTok{as}\NormalTok{ np}

\NormalTok{rng }\OperatorTok{=}\NormalTok{ np.random.default\_rng(}\DecValTok{7243}\NormalTok{)}

\NormalTok{vocab }\OperatorTok{=}\NormalTok{ \{}\StringTok{"the"}\NormalTok{: }\DecValTok{0}\NormalTok{, }\StringTok{"movie"}\NormalTok{: }\DecValTok{1}\NormalTok{, }\StringTok{"was"}\NormalTok{: }\DecValTok{2}\NormalTok{, }\StringTok{"not"}\NormalTok{: }\DecValTok{3}\NormalTok{, }\StringTok{"very"}\NormalTok{: }\DecValTok{4}\NormalTok{, }\StringTok{"good"}\NormalTok{: }\DecValTok{5}\NormalTok{\}}
\NormalTok{sequence }\OperatorTok{=}\NormalTok{ [}\StringTok{"the"}\NormalTok{, }\StringTok{"movie"}\NormalTok{, }\StringTok{"was"}\NormalTok{, }\StringTok{"not"}\NormalTok{, }\StringTok{"very"}\NormalTok{, }\StringTok{"good"}\NormalTok{]}
\NormalTok{x }\OperatorTok{=}\NormalTok{ np.array([vocab[w] }\ControlFlowTok{for}\NormalTok{ w }\KeywordTok{in}\NormalTok{ sequence])}

\CommentTok{\# Dimensions}
\NormalTok{d }\OperatorTok{=} \BuiltInTok{len}\NormalTok{(vocab)   }\CommentTok{\# vocabulary size}
\NormalTok{m }\OperatorTok{=} \DecValTok{4}            \CommentTok{\# embedding dimension}
\NormalTok{r }\OperatorTok{=} \DecValTok{5}            \CommentTok{\# hidden/cell dimension}

\CommentTok{\# Embedding matrix}
\NormalTok{E }\OperatorTok{=}\NormalTok{ rng.normal(scale}\OperatorTok{=}\FloatTok{0.2}\NormalTok{, size}\OperatorTok{=}\NormalTok{(d, m))}

\CommentTok{\# LSTM parameters}
\NormalTok{W\_f }\OperatorTok{=}\NormalTok{ rng.normal(scale}\OperatorTok{=}\FloatTok{0.3}\NormalTok{, size}\OperatorTok{=}\NormalTok{(r, m))}
\NormalTok{U\_f }\OperatorTok{=}\NormalTok{ rng.normal(scale}\OperatorTok{=}\FloatTok{0.3}\NormalTok{, size}\OperatorTok{=}\NormalTok{(r, r))}
\NormalTok{b\_f }\OperatorTok{=}\NormalTok{ np.ones(r) }\OperatorTok{*} \FloatTok{0.5}       \CommentTok{\# positive forget bias encourages memory early in training}

\NormalTok{W\_i }\OperatorTok{=}\NormalTok{ rng.normal(scale}\OperatorTok{=}\FloatTok{0.3}\NormalTok{, size}\OperatorTok{=}\NormalTok{(r, m))}
\NormalTok{U\_i }\OperatorTok{=}\NormalTok{ rng.normal(scale}\OperatorTok{=}\FloatTok{0.3}\NormalTok{, size}\OperatorTok{=}\NormalTok{(r, r))}
\NormalTok{b\_i }\OperatorTok{=}\NormalTok{ np.zeros(r)}

\NormalTok{W\_o }\OperatorTok{=}\NormalTok{ rng.normal(scale}\OperatorTok{=}\FloatTok{0.3}\NormalTok{, size}\OperatorTok{=}\NormalTok{(r, m))}
\NormalTok{U\_o }\OperatorTok{=}\NormalTok{ rng.normal(scale}\OperatorTok{=}\FloatTok{0.3}\NormalTok{, size}\OperatorTok{=}\NormalTok{(r, r))}
\NormalTok{b\_o }\OperatorTok{=}\NormalTok{ np.zeros(r)}

\NormalTok{W\_c }\OperatorTok{=}\NormalTok{ rng.normal(scale}\OperatorTok{=}\FloatTok{0.3}\NormalTok{, size}\OperatorTok{=}\NormalTok{(r, m))}
\NormalTok{U\_c }\OperatorTok{=}\NormalTok{ rng.normal(scale}\OperatorTok{=}\FloatTok{0.3}\NormalTok{, size}\OperatorTok{=}\NormalTok{(r, r))}
\NormalTok{b\_c }\OperatorTok{=}\NormalTok{ np.zeros(r)}

\KeywordTok{def}\NormalTok{ sigmoid(z):}
    \ControlFlowTok{return} \DecValTok{1} \OperatorTok{/}\NormalTok{ (}\DecValTok{1} \OperatorTok{+}\NormalTok{ np.exp(}\OperatorTok{{-}}\NormalTok{z))}

\KeywordTok{def}\NormalTok{ lstm\_step(e\_t, h\_prev, c\_prev):}
\NormalTok{    f\_t }\OperatorTok{=}\NormalTok{ sigmoid(W\_f }\OperatorTok{@}\NormalTok{ e\_t }\OperatorTok{+}\NormalTok{ U\_f }\OperatorTok{@}\NormalTok{ h\_prev }\OperatorTok{+}\NormalTok{ b\_f)}
\NormalTok{    i\_t }\OperatorTok{=}\NormalTok{ sigmoid(W\_i }\OperatorTok{@}\NormalTok{ e\_t }\OperatorTok{+}\NormalTok{ U\_i }\OperatorTok{@}\NormalTok{ h\_prev }\OperatorTok{+}\NormalTok{ b\_i)}
\NormalTok{    o\_t }\OperatorTok{=}\NormalTok{ sigmoid(W\_o }\OperatorTok{@}\NormalTok{ e\_t }\OperatorTok{+}\NormalTok{ U\_o }\OperatorTok{@}\NormalTok{ h\_prev }\OperatorTok{+}\NormalTok{ b\_o)}
\NormalTok{    c\_tilde }\OperatorTok{=}\NormalTok{ np.tanh(W\_c }\OperatorTok{@}\NormalTok{ e\_t }\OperatorTok{+}\NormalTok{ U\_c }\OperatorTok{@}\NormalTok{ h\_prev }\OperatorTok{+}\NormalTok{ b\_c)}
\NormalTok{    c\_t }\OperatorTok{=}\NormalTok{ f\_t }\OperatorTok{*}\NormalTok{ c\_prev }\OperatorTok{+}\NormalTok{ i\_t }\OperatorTok{*}\NormalTok{ c\_tilde}
\NormalTok{    h\_t }\OperatorTok{=}\NormalTok{ o\_t }\OperatorTok{*}\NormalTok{ np.tanh(c\_t)}
\NormalTok{    cache }\OperatorTok{=}\NormalTok{ \{}\StringTok{"f"}\NormalTok{: f\_t, }\StringTok{"i"}\NormalTok{: i\_t, }\StringTok{"o"}\NormalTok{: o\_t, }\StringTok{"c\_tilde"}\NormalTok{: c\_tilde, }\StringTok{"c"}\NormalTok{: c\_t, }\StringTok{"h"}\NormalTok{: h\_t\}}
    \ControlFlowTok{return}\NormalTok{ h\_t, c\_t, cache}

\NormalTok{h }\OperatorTok{=}\NormalTok{ np.zeros(r)}
\NormalTok{c }\OperatorTok{=}\NormalTok{ np.zeros(r)}
\NormalTok{history }\OperatorTok{=}\NormalTok{ []}

\ControlFlowTok{for}\NormalTok{ token\_id }\KeywordTok{in}\NormalTok{ x:}
\NormalTok{    e\_t }\OperatorTok{=}\NormalTok{ E[token\_id]}
\NormalTok{    h, c, cache }\OperatorTok{=}\NormalTok{ lstm\_step(e\_t, h, c)}
\NormalTok{    history.append(cache)}

\BuiltInTok{print}\NormalTok{(}\StringTok{"Sequence:"}\NormalTok{, sequence)}
\BuiltInTok{print}\NormalTok{(}\StringTok{"Final hidden state shape:"}\NormalTok{, h.shape)}
\BuiltInTok{print}\NormalTok{(}\StringTok{"Final cell state shape:"}\NormalTok{, c.shape)}
\BuiltInTok{print}\NormalTok{(}\StringTok{"Final hidden state:"}\NormalTok{, np.}\BuiltInTok{round}\NormalTok{(h, }\DecValTok{3}\NormalTok{))}
\BuiltInTok{print}\NormalTok{(}\StringTok{"Final forget gate:"}\NormalTok{, np.}\BuiltInTok{round}\NormalTok{(history[}\OperatorTok{{-}}\DecValTok{1}\NormalTok{][}\StringTok{"f"}\NormalTok{], }\DecValTok{3}\NormalTok{))}

\CommentTok{\# Save the final LSTM states for later examples.}
\NormalTok{h\_lstm }\OperatorTok{=}\NormalTok{ h.copy()}
\NormalTok{c\_lstm }\OperatorTok{=}\NormalTok{ c.copy()}
\end{Highlighting}
\end{Shaded}

\begin{verbatim}
Sequence: ['the', 'movie', 'was', 'not', 'very', 'good']
Final hidden state shape: (5,)
Final cell state shape: (5,)
Final hidden state: [ 0.081  0.012  0.006  0.043 -0.037]
Final forget gate: [0.618 0.655 0.649 0.638 0.624]
\end{verbatim}

The forget bias was initialized to a positive value. This is common
because at the beginning of training it encourages the model not to
erase memory too aggressively.

\section{12.13 Python example: gate products and effective
memory}\label{python-example-gate-products-and-effective-memory}

The direct LSTM memory path is controlled by products of forget gates.
The following example simulates several constant forget-gate values and
plots the decay curves.

\begin{Shaded}
\begin{Highlighting}[]
\ImportTok{import}\NormalTok{ numpy }\ImportTok{as}\NormalTok{ np}
\ImportTok{import}\NormalTok{ matplotlib.pyplot }\ImportTok{as}\NormalTok{ plt}

\NormalTok{T }\OperatorTok{=} \DecValTok{80}
\NormalTok{steps }\OperatorTok{=}\NormalTok{ np.arange(T }\OperatorTok{+} \DecValTok{1}\NormalTok{)}
\NormalTok{forget\_values }\OperatorTok{=}\NormalTok{ [}\FloatTok{0.70}\NormalTok{, }\FloatTok{0.85}\NormalTok{, }\FloatTok{0.95}\NormalTok{, }\FloatTok{0.99}\NormalTok{]}

\NormalTok{plt.figure(figsize}\OperatorTok{=}\NormalTok{(}\DecValTok{7}\NormalTok{, }\DecValTok{4}\NormalTok{))}
\ControlFlowTok{for}\NormalTok{ f }\KeywordTok{in}\NormalTok{ forget\_values:}
\NormalTok{    weights }\OperatorTok{=}\NormalTok{ f }\OperatorTok{**}\NormalTok{ steps}
\NormalTok{    half\_life }\OperatorTok{=}\NormalTok{ np.log(}\FloatTok{0.5}\NormalTok{) }\OperatorTok{/}\NormalTok{ np.log(f)}
\NormalTok{    plt.plot(steps, weights, label}\OperatorTok{=}\SpecialStringTok{f"f=}\SpecialCharTok{\{}\NormalTok{f}\SpecialCharTok{\}}\SpecialStringTok{, half{-}life=}\SpecialCharTok{\{}\NormalTok{half\_life}\SpecialCharTok{:.1f\}}\SpecialStringTok{"}\NormalTok{)}

\NormalTok{plt.axhline(}\FloatTok{0.5}\NormalTok{, linestyle}\OperatorTok{=}\StringTok{"{-}{-}"}\NormalTok{, linewidth}\OperatorTok{=}\DecValTok{1}\NormalTok{)}
\NormalTok{plt.xlabel(}\StringTok{"number of steps k"}\NormalTok{)}
\NormalTok{plt.ylabel(}\StringTok{"memory weight f\^{}k"}\NormalTok{)}
\NormalTok{plt.title(}\StringTok{"Direct memory decay controlled by the forget gate"}\NormalTok{)}
\NormalTok{plt.legend()}
\NormalTok{plt.show()}
\end{Highlighting}
\end{Shaded}

\pandocbounded{\includegraphics[keepaspectratio]{chapters/chapter-12-lstm-gru-models_files/figure-pdf/ch12-forget-products-output-1.pdf}}

This plot gives a concrete mathematical interpretation of gate values. A
change from \(0.95\) to \(0.99\) has a large effect on long-range
memory.

\section{12.14 Python example: GRU forward pass from
scratch}\label{python-example-gru-forward-pass-from-scratch}

The GRU is similar in spirit but uses fewer state variables. The
following code implements a GRU forward pass.

\protect\phantomsection\label{ch12-gru-forward}
\begin{Shaded}
\begin{Highlighting}[]
\ImportTok{import}\NormalTok{ numpy }\ImportTok{as}\NormalTok{ np}

\NormalTok{rng }\OperatorTok{=}\NormalTok{ np.random.default\_rng(}\DecValTok{7339}\NormalTok{)}

\NormalTok{vocab }\OperatorTok{=}\NormalTok{ \{}\StringTok{"language"}\NormalTok{: }\DecValTok{0}\NormalTok{, }\StringTok{"models"}\NormalTok{: }\DecValTok{1}\NormalTok{, }\StringTok{"need"}\NormalTok{: }\DecValTok{2}\NormalTok{, }\StringTok{"long"}\NormalTok{: }\DecValTok{3}\NormalTok{, }\StringTok{"memory"}\NormalTok{: }\DecValTok{4}\NormalTok{\}}
\NormalTok{sequence }\OperatorTok{=}\NormalTok{ [}\StringTok{"language"}\NormalTok{, }\StringTok{"models"}\NormalTok{, }\StringTok{"need"}\NormalTok{, }\StringTok{"long"}\NormalTok{, }\StringTok{"memory"}\NormalTok{]}
\NormalTok{x }\OperatorTok{=}\NormalTok{ np.array([vocab[w] }\ControlFlowTok{for}\NormalTok{ w }\KeywordTok{in}\NormalTok{ sequence])}

\CommentTok{\# Dimensions}
\NormalTok{d }\OperatorTok{=} \BuiltInTok{len}\NormalTok{(vocab)}
\NormalTok{m }\OperatorTok{=} \DecValTok{4}
\NormalTok{r }\OperatorTok{=} \DecValTok{5}

\NormalTok{E }\OperatorTok{=}\NormalTok{ rng.normal(scale}\OperatorTok{=}\FloatTok{0.2}\NormalTok{, size}\OperatorTok{=}\NormalTok{(d, m))}

\NormalTok{W\_z }\OperatorTok{=}\NormalTok{ rng.normal(scale}\OperatorTok{=}\FloatTok{0.3}\NormalTok{, size}\OperatorTok{=}\NormalTok{(r, m))}
\NormalTok{U\_z }\OperatorTok{=}\NormalTok{ rng.normal(scale}\OperatorTok{=}\FloatTok{0.3}\NormalTok{, size}\OperatorTok{=}\NormalTok{(r, r))}
\NormalTok{b\_z }\OperatorTok{=}\NormalTok{ np.zeros(r)}

\NormalTok{W\_r }\OperatorTok{=}\NormalTok{ rng.normal(scale}\OperatorTok{=}\FloatTok{0.3}\NormalTok{, size}\OperatorTok{=}\NormalTok{(r, m))}
\NormalTok{U\_r }\OperatorTok{=}\NormalTok{ rng.normal(scale}\OperatorTok{=}\FloatTok{0.3}\NormalTok{, size}\OperatorTok{=}\NormalTok{(r, r))}
\NormalTok{b\_r }\OperatorTok{=}\NormalTok{ np.zeros(r)}

\NormalTok{W\_h }\OperatorTok{=}\NormalTok{ rng.normal(scale}\OperatorTok{=}\FloatTok{0.3}\NormalTok{, size}\OperatorTok{=}\NormalTok{(r, m))}
\NormalTok{U\_h }\OperatorTok{=}\NormalTok{ rng.normal(scale}\OperatorTok{=}\FloatTok{0.3}\NormalTok{, size}\OperatorTok{=}\NormalTok{(r, r))}
\NormalTok{b\_h }\OperatorTok{=}\NormalTok{ np.zeros(r)}

\KeywordTok{def}\NormalTok{ sigmoid(z):}
    \ControlFlowTok{return} \DecValTok{1} \OperatorTok{/}\NormalTok{ (}\DecValTok{1} \OperatorTok{+}\NormalTok{ np.exp(}\OperatorTok{{-}}\NormalTok{z))}

\KeywordTok{def}\NormalTok{ gru\_step(e\_t, h\_prev):}
\NormalTok{    z\_t }\OperatorTok{=}\NormalTok{ sigmoid(W\_z }\OperatorTok{@}\NormalTok{ e\_t }\OperatorTok{+}\NormalTok{ U\_z }\OperatorTok{@}\NormalTok{ h\_prev }\OperatorTok{+}\NormalTok{ b\_z)}
\NormalTok{    r\_t }\OperatorTok{=}\NormalTok{ sigmoid(W\_r }\OperatorTok{@}\NormalTok{ e\_t }\OperatorTok{+}\NormalTok{ U\_r }\OperatorTok{@}\NormalTok{ h\_prev }\OperatorTok{+}\NormalTok{ b\_r)}
\NormalTok{    h\_tilde }\OperatorTok{=}\NormalTok{ np.tanh(W\_h }\OperatorTok{@}\NormalTok{ e\_t }\OperatorTok{+}\NormalTok{ U\_h }\OperatorTok{@}\NormalTok{ (r\_t }\OperatorTok{*}\NormalTok{ h\_prev) }\OperatorTok{+}\NormalTok{ b\_h)}
\NormalTok{    h\_t }\OperatorTok{=}\NormalTok{ (}\DecValTok{1} \OperatorTok{{-}}\NormalTok{ z\_t) }\OperatorTok{*}\NormalTok{ h\_prev }\OperatorTok{+}\NormalTok{ z\_t }\OperatorTok{*}\NormalTok{ h\_tilde}
\NormalTok{    cache }\OperatorTok{=}\NormalTok{ \{}\StringTok{"z"}\NormalTok{: z\_t, }\StringTok{"r"}\NormalTok{: r\_t, }\StringTok{"h\_tilde"}\NormalTok{: h\_tilde, }\StringTok{"h"}\NormalTok{: h\_t\}}
    \ControlFlowTok{return}\NormalTok{ h\_t, cache}

\NormalTok{h }\OperatorTok{=}\NormalTok{ np.zeros(r)}
\NormalTok{states }\OperatorTok{=}\NormalTok{ []}
\NormalTok{gates }\OperatorTok{=}\NormalTok{ []}
\ControlFlowTok{for}\NormalTok{ token\_id }\KeywordTok{in}\NormalTok{ x:}
\NormalTok{    e\_t }\OperatorTok{=}\NormalTok{ E[token\_id]}
\NormalTok{    h, cache }\OperatorTok{=}\NormalTok{ gru\_step(e\_t, h)}
\NormalTok{    states.append(h.copy())}
\NormalTok{    gates.append(cache)}

\NormalTok{states }\OperatorTok{=}\NormalTok{ np.array(states)}
\BuiltInTok{print}\NormalTok{(}\StringTok{"Sequence:"}\NormalTok{, sequence)}
\BuiltInTok{print}\NormalTok{(}\StringTok{"Hidden{-}state trajectory shape:"}\NormalTok{, states.shape)}
\BuiltInTok{print}\NormalTok{(}\StringTok{"Final update gate:"}\NormalTok{, np.}\BuiltInTok{round}\NormalTok{(gates[}\OperatorTok{{-}}\DecValTok{1}\NormalTok{][}\StringTok{"z"}\NormalTok{], }\DecValTok{3}\NormalTok{))}
\BuiltInTok{print}\NormalTok{(}\StringTok{"Final reset gate:"}\NormalTok{, np.}\BuiltInTok{round}\NormalTok{(gates[}\OperatorTok{{-}}\DecValTok{1}\NormalTok{][}\StringTok{"r"}\NormalTok{], }\DecValTok{3}\NormalTok{))}
\end{Highlighting}
\end{Shaded}

\begin{verbatim}
Sequence: ['language', 'models', 'need', 'long', 'memory']
Hidden-state trajectory shape: (5, 5)
Final update gate: [0.489 0.517 0.502 0.449 0.454]
Final reset gate: [0.506 0.508 0.52  0.481 0.498]
\end{verbatim}

The update gate \(z_t\) controls how much the state changes. Small
update-gate values preserve the old hidden state, while large values
overwrite it with the candidate state.

\section{12.15 Python example: LSTM output layer for
classification}\label{python-example-lstm-output-layer-for-classification}

After reading a full sequence, an LSTM can feed \(h_T\) into a
classifier:

\[
z = W_yh_T+b_y,
\qquad
p=\operatorname{softmax}(z).
\]

The following code uses the final LSTM hidden state saved from the
earlier forward pass and computes class probabilities.

\protect\phantomsection\label{ch12-lstm-classifier-head}
\begin{Shaded}
\begin{Highlighting}[]
\ImportTok{import}\NormalTok{ numpy }\ImportTok{as}\NormalTok{ np}

\NormalTok{K }\OperatorTok{=} \DecValTok{2}
\NormalTok{hidden\_dim }\OperatorTok{=}\NormalTok{ h\_lstm.shape[}\DecValTok{0}\NormalTok{]}
\NormalTok{W\_y }\OperatorTok{=}\NormalTok{ rng.normal(scale}\OperatorTok{=}\FloatTok{0.3}\NormalTok{, size}\OperatorTok{=}\NormalTok{(K, hidden\_dim))}
\NormalTok{b\_y }\OperatorTok{=}\NormalTok{ np.zeros(K)}

\KeywordTok{def}\NormalTok{ softmax(z):}
\NormalTok{    z }\OperatorTok{=}\NormalTok{ z }\OperatorTok{{-}}\NormalTok{ np.}\BuiltInTok{max}\NormalTok{(z)}
\NormalTok{    exp\_z }\OperatorTok{=}\NormalTok{ np.exp(z)}
    \ControlFlowTok{return}\NormalTok{ exp\_z }\OperatorTok{/}\NormalTok{ exp\_z.}\BuiltInTok{sum}\NormalTok{()}

\NormalTok{logits }\OperatorTok{=}\NormalTok{ W\_y }\OperatorTok{@}\NormalTok{ h\_lstm }\OperatorTok{+}\NormalTok{ b\_y}
\NormalTok{probs }\OperatorTok{=}\NormalTok{ softmax(logits)}
\NormalTok{y\_true }\OperatorTok{=} \DecValTok{1}
\NormalTok{loss }\OperatorTok{=} \OperatorTok{{-}}\NormalTok{np.log(probs[y\_true])}

\BuiltInTok{print}\NormalTok{(}\StringTok{"Logits:"}\NormalTok{, np.}\BuiltInTok{round}\NormalTok{(logits, }\DecValTok{3}\NormalTok{))}
\BuiltInTok{print}\NormalTok{(}\StringTok{"Class probabilities:"}\NormalTok{, np.}\BuiltInTok{round}\NormalTok{(probs, }\DecValTok{3}\NormalTok{))}
\BuiltInTok{print}\NormalTok{(}\StringTok{"Cross{-}entropy loss:"}\NormalTok{, }\BuiltInTok{round}\NormalTok{(}\BuiltInTok{float}\NormalTok{(loss), }\DecValTok{3}\NormalTok{))}
\end{Highlighting}
\end{Shaded}

\begin{verbatim}
Logits: [-0.063 -0.025]
Class probabilities: [0.49 0.51]
Cross-entropy loss: 0.674
\end{verbatim}

The recurrent layer constructs the representation. The softmax layer
converts the representation into a conditional probability distribution.

\section{12.16 AI-assisted learning
component}\label{ai-assisted-learning-component-5}

Use an AI assistant as a mathematical critic, not as a replacement for
derivation. The following prompts are useful after you have written the
equations yourself.

\begin{tcolorbox}[enhanced jigsaw, arc=.35mm, bottomrule=.15mm, bottomtitle=1mm, breakable, colback=white, colbacktitle=quarto-callout-tip-color!10!white, colframe=quarto-callout-tip-color-frame, coltitle=black, left=2mm, leftrule=.75mm, opacityback=0, opacitybacktitle=0.6, rightrule=.15mm, title=\textcolor{quarto-callout-tip-color}{\faLightbulb}\hspace{0.5em}{Prompt 1: Dimension audit}, titlerule=0mm, toprule=.15mm, toptitle=1mm]

I wrote the following LSTM equations. Check every vector and matrix
dimension. Identify any inconsistent notation. Do not simplify the
model; only audit the mathematics.

\end{tcolorbox}

\begin{tcolorbox}[enhanced jigsaw, arc=.35mm, bottomrule=.15mm, bottomtitle=1mm, breakable, colback=white, colbacktitle=quarto-callout-tip-color!10!white, colframe=quarto-callout-tip-color-frame, coltitle=black, left=2mm, leftrule=.75mm, opacityback=0, opacitybacktitle=0.6, rightrule=.15mm, title=\textcolor{quarto-callout-tip-color}{\faLightbulb}\hspace{0.5em}{Prompt 2: Gate interpretation}, titlerule=0mm, toprule=.15mm, toptitle=1mm]

Explain the forget gate, input gate, and output gate of an LSTM using
coordinatewise convex combinations. Give a small scalar example and
connect it to memory half-life.

\end{tcolorbox}

\begin{tcolorbox}[enhanced jigsaw, arc=.35mm, bottomrule=.15mm, bottomtitle=1mm, breakable, colback=white, colbacktitle=quarto-callout-tip-color!10!white, colframe=quarto-callout-tip-color-frame, coltitle=black, left=2mm, leftrule=.75mm, opacityback=0, opacitybacktitle=0.6, rightrule=.15mm, title=\textcolor{quarto-callout-tip-color}{\faLightbulb}\hspace{0.5em}{Prompt 3: Compare models}, titlerule=0mm, toprule=.15mm, toptitle=1mm]

Compare a vanilla RNN, an LSTM, and a GRU from the point of view of
Jacobian products and gradient flow. Avoid vague statements such as
``LSTMs remember better'' unless they are connected to equations.

\end{tcolorbox}

\begin{tcolorbox}[enhanced jigsaw, arc=.35mm, bottomrule=.15mm, bottomtitle=1mm, breakable, colback=white, colbacktitle=quarto-callout-warning-color!10!white, colframe=quarto-callout-warning-color-frame, coltitle=black, left=2mm, leftrule=.75mm, opacityback=0, opacitybacktitle=0.6, rightrule=.15mm, title=\textcolor{quarto-callout-warning-color}{\faExclamationTriangle}\hspace{0.5em}{AI caution}, titlerule=0mm, toprule=.15mm, toptitle=1mm]

AI-generated explanations of LSTMs often describe gates only verbally.
For this course, a correct explanation must include the equations, the
dimensions, and the memory-path product controlled by forget or update
gates.

\end{tcolorbox}

\section{12.17 What to remember}\label{what-to-remember-2}

\begin{enumerate}
\def\labelenumi{\arabic{enumi}.}
\tightlist
\item
  A gate is a vector in \((0,1)^r\) computed by a sigmoid affine map.
\item
  Gates create coordinatewise interpolation between old information and
  new information.
\item
  The LSTM cell state follows
\end{enumerate}

\[
c_t=f_t\odot c_{t-1}+i_t\odot\widetilde c_t.
\]

\begin{enumerate}
\def\labelenumi{\arabic{enumi}.}
\setcounter{enumi}{3}
\tightlist
\item
  The direct LSTM memory path is governed by products of forget-gate
  coordinates.
\item
  A GRU uses an update gate to interpolate between the old hidden state
  and a candidate hidden state.
\item
  LSTM and GRU models are conditional probability models trained by
  cross-entropy or negative log-likelihood.
\item
  The mathematical advantage of gates is not magic memory; it is a more
  stable differentiable path through time.
\end{enumerate}

\section{12.18 Exercises}\label{exercises-11}

\begin{enumerate}
\def\labelenumi{\arabic{enumi}.}
\item
  \textbf{Gate range.} Prove that every coordinate of
  \(g_t=\sigma(W_ge_t+U_gh_{t-1}+b_g)\) lies strictly between \(0\) and
  \(1\).
\item
  \textbf{Dimension check.} Suppose \(e_t\in R^{20}\) and
  \(h_t,c_t\in R^{64}\). Give the dimensions of all LSTM matrices
  \(W_f,U_f,W_i,U_i,W_o,U_o,W_c,U_c\) and all biases.
\item
  \textbf{Parameter count.} For \(m=100\) and \(r=256\), compute the
  number of recurrent parameters in one LSTM layer. Do the same for one
  GRU layer.
\item
  \textbf{Memory half-life.} Compute the memory half-life for \(f=0.8\),
  \(0.9\), \(0.95\), and \(0.99\). Interpret the result.
\item
  \textbf{Scalar LSTM.} Let \(c_t=0.9c_{t-1}+0.5\tanh(u_t)\) and
  \(h_t=\tanh(c_t)\). Compute \(c_1,c_2,c_3\) for \(c_0=0\) and inputs
  \(u_1=1\), \(u_2=0\), \(u_3=0\).
\item
  \textbf{GRU interpolation.} For scalar
  \(h_t=(1-z)h_{t-1}+z\widetilde h_t\), show that \(h_t\) lies between
  \(h_{t-1}\) and \(\widetilde h_t\) whenever \(0\le z\le 1\).
\item
  \textbf{Gradient path.} For constant scalar forget gate \(f\), derive
  \(\partial c_T/\partial c_s=f^{T-s}\).
\item
  \textbf{Model comparison.} Explain why the LSTM has more parameters
  than the GRU for the same input and hidden dimensions.
\item
  \textbf{Language modeling.} Write the negative log-likelihood for an
  LSTM language model on a corpus of variable-length token sequences.
\item
  \textbf{Critical thinking.} Give one situation in which a gated RNN
  might still fail to capture long-range structure. Explain whether the
  issue is representation, optimization, data, or evaluation.
\end{enumerate}

\section{12.19 Mini-project}\label{mini-project}

Build a small NumPy experiment comparing a scalar vanilla recurrence and
a scalar gated memory recurrence on a delayed-signal task.

A possible task is:

\begin{enumerate}
\def\labelenumi{\arabic{enumi}.}
\tightlist
\item
  Generate sequences of length \(T\).
\item
  Put a binary signal at time \(1\).
\item
  Ask the model to predict the signal at time \(T\).
\item
  Compare how the memory weight changes with \(a\) in the vanilla
  recurrence and \(f\) in the gated recurrence.
\item
  Plot the final memory strength as a function of \(T\).
\end{enumerate}

The goal is not to train a large neural network. The goal is to
understand how recurrence parameters control memory mathematically.

\section{12.20 Python lab connection}\label{python-lab-connection-5}

The associated lab for this chapter is planned as:

\begin{Shaded}
\begin{Highlighting}[]
\NormalTok{labs/chapter{-}12{-}lstm{-}gru{-}models.ipynb}
\end{Highlighting}
\end{Shaded}

The lab should be written for independent study and should include:

\begin{itemize}
\tightlist
\item
  background on sigmoid, tanh, and coordinatewise multiplication;
\item
  a from-scratch LSTM forward pass;
\item
  a from-scratch GRU forward pass;
\item
  memory half-life experiments;
\item
  a small sequence classification example;
\item
  optional comparison with a deep learning library implementation.
\end{itemize}

\section{Notes and references}\label{notes-and-references-6}

The LSTM was introduced by Hochreiter and Schmidhuber, and the GRU was
introduced as a simpler gated recurrent architecture in neural machine
translation research. This chapter emphasizes the mathematical structure
of gated recurrence: coordinatewise gates, additive memory updates, and
gradient paths. For broader NLP context, see Jurafsky and Martin (2024)
and Goldberg (2017). For deep learning background, see Goodfellow et al.
(2016a).

\part{Part V. Attention and Transformers}

\chapter{Attention Mechanisms}\label{attention-mechanisms}

Queries, keys, values, and weighted information flow

\hfill\break

Attention is the mathematical operation that allows a model to decide,
at each token position, which other token positions are most relevant.
Instead of compressing a whole prefix into a single recurrent state,
attention forms many pairwise comparisons between tokens. These
comparisons become nonnegative weights, and the output at each position
is a weighted average of value vectors.

This chapter develops attention from first principles. The central
object is the matrix formula

\[
\operatorname{Attention}(Q,K,V)
=
\operatorname{softmax}_{\text{row}}\left(\frac{QK^T}{\sqrt{d_k}}\right)V.
\]

The goal is not to memorize the formula, but to understand what each
matrix means, why the scaling factor appears, why the rows of the
attention matrix are probability vectors, and how masking changes the
conditional structure of the model.

\section{Learning goals}\label{learning-goals-12}

After studying this chapter, you should be able to:

\begin{enumerate}
\def\labelenumi{\arabic{enumi}.}
\tightlist
\item
  explain the role of queries, keys, and values in an attention layer;
\item
  check the dimensions in the scaled dot-product attention formula;
\item
  interpret the attention matrix as a row-stochastic matrix;
\item
  explain why attention outputs are weighted averages of value vectors;
\item
  compare hard lookup, nearest-neighbor search, kernel smoothing, and
  attention;
\item
  implement attention from scratch in NumPy;
\item
  apply padding masks and causal masks correctly;
\item
  use attention weights to inspect token-token interactions;
\item
  explain why attention is useful for NLP sequence modeling.
\end{enumerate}

\section{13.1 From sequential memory to pairwise
interaction}\label{from-sequential-memory-to-pairwise-interaction}

In a recurrent neural network, information flows through a hidden-state
recurrence

\[
h_t = F(h_{t-1},x_t).
\]

The representation of token \(t\) depends on earlier tokens through
repeated application of \(F\). This is a natural model for ordered data,
but it creates a narrow information path. If token \(t\) needs
information from token \(s\) with \(s\ll t\), the information must pass
through all intermediate updates.

Attention uses a different idea. Given token representations

\[
x_1,x_2,\ldots,x_T \in R^{d_{\mathrm{model}}},
\]

attention lets every position compare itself directly with every other
position. Position \(i\) asks a question, position \(j\) offers a key,
and the similarity between the question and the key determines how much
information flows from \(j\) to \(i\).

In one sentence:

\textbf{Attention is adaptive weighted averaging.} For each output
position \(i\), attention computes weights \(a_{ij}\ge 0\) with
\(\sum_{j=1}^T a_{ij}=1\), then returns

\[
z_i = \sum_{j=1}^T a_{ij}v_j.
\]

The weights are not fixed in advance. They are computed from the current
sequence.

\section{13.2 Input matrix and token
representations}\label{input-matrix-and-token-representations}

Suppose a sequence has length \(T\). Let

\[
X =
\begin{bmatrix}
- & x_1^T & - \\
- & x_2^T & - \\
& \vdots & \\
- & x_T^T & -
\end{bmatrix}
\in R^{T\times d_{\mathrm{model}}}.
\]

The row \(x_i^T\) is the representation of token position \(i\). This
representation may combine token embedding, positional information, and
earlier neural transformations.

Attention does not use \(X\) directly. It first produces three
transformed views of the same sequence:

\[
Q = XW_Q,\qquad K = XW_K,\qquad V = XW_V.
\]

Here

\[
W_Q\in R^{d_{\mathrm{model}}\times d_k},\qquad
W_K\in R^{d_{\mathrm{model}}\times d_k},\qquad
W_V\in R^{d_{\mathrm{model}}\times d_v}.
\]

Therefore

\[
Q,K\in R^{T\times d_k},\qquad V\in R^{T\times d_v}.
\]

The rows are denoted

\[
q_i^T = \text{row } i \text{ of } Q,\qquad
k_j^T = \text{row } j \text{ of } K,\qquad
v_j^T = \text{row } j \text{ of } V.
\]

Interpretation:

\begin{itemize}
\tightlist
\item
  \(q_i\) is the query at output position \(i\);
\item
  \(k_j\) is the key advertised by input position \(j\);
\item
  \(v_j\) is the information carried by input position \(j\).
\end{itemize}

The query and key live in the same space \(R^{d_k}\) so that their dot
product is meaningful. The value vector may live in a different space
\(R^{d_v}\).

\section{13.3 Scores as pairwise
similarities}\label{scores-as-pairwise-similarities}

For each pair of positions \((i,j)\), define the raw attention score

\[
s_{ij}=q_i^T k_j.
\]

In matrix form,

\[
S = QK^T \in R^{T\times T}.
\]

The entry \(S_{ij}\) compares the query at position \(i\) with the key
at position \(j\).

A large positive score means that token \(i\) strongly attends to token
\(j\). A large negative score means that token \(j\) is unlikely to
receive much attention from token \(i\) after the softmax step.

\subsection{\texorpdfstring{Why divide by
\(\sqrt{d_k}\)?}{Why divide by \textbackslash sqrt\{d\_k\}?}}\label{why-divide-by-sqrtd_k}

Assume for intuition that the coordinates of \(q_i\) and \(k_j\) are
independent random variables with mean \(0\) and variance \(1\). Then

\[
q_i^T k_j = \sum_{r=1}^{d_k} q_{ir}k_{jr}.
\]

Each product has mean \(0\) and variance approximately \(1\), so the sum
has variance approximately \(d_k\). Thus the typical size of
\(q_i^Tk_j\) grows like \(\sqrt{d_k}\).

If \(d_k\) is large, the unscaled scores may be too large in magnitude.
The softmax then becomes nearly one-hot, which can produce unstable
gradients. The scaled score is

\[
\tilde{s}_{ij}=\frac{q_i^Tk_j}{\sqrt{d_k}}.
\]

The scaling keeps the variance of scores more stable as \(d_k\) changes.

\section{13.4 Row-wise softmax and the attention
matrix}\label{row-wise-softmax-and-the-attention-matrix}

Given scaled scores

\[
\tilde{S}=\frac{QK^T}{\sqrt{d_k}},
\]

define the attention weights row by row:

\[
a_{ij}
=
\frac{\exp(\tilde{s}_{ij})}{\sum_{\ell=1}^T \exp(\tilde{s}_{i\ell})}.
\]

Let

\[
A=\operatorname{softmax}_{\text{row}}(\tilde{S}).
\]

Then \(A\in R^{T\times T}\) and every row of \(A\) is a probability
vector:

\[
a_{ij}\ge 0,
\qquad
\sum_{j=1}^T a_{ij}=1.
\]

Thus attention turns arbitrary real-valued similarity scores into a
row-stochastic matrix.

\textbf{Attention as a Markov-style matrix.} The matrix \(A\) is
row-stochastic, like a transition matrix in a Markov chain. However,
\(A\) is not fixed. It is computed from \(X\) and therefore changes from
sequence to sequence.

\section{13.5 Weighted averages of
values}\label{weighted-averages-of-values}

The output of attention is

\[
Z=AV.
\]

Because \(A\in R^{T\times T}\) and \(V\in R^{T\times d_v}\), the output
satisfies

\[
Z\in R^{T\times d_v}.
\]

Row \(i\) is

\[
z_i = \sum_{j=1}^T a_{ij}v_j.
\]

Since \(a_{ij}\ge 0\) and \(\sum_j a_{ij}=1\), each output vector
\(z_i\) is a convex combination of the value vectors \(v_1,\ldots,v_T\).

This is the basic mechanism by which token representations exchange
information. Token \(i\) receives a weighted summary of all value
vectors, with weights determined by the similarity between \(q_i\) and
the keys.

\section{13.6 Attention as adaptive kernel
smoothing}\label{attention-as-adaptive-kernel-smoothing}

Students with a statistics background may find the following analogy
helpful. In kernel smoothing, an estimate at a point \(x\) often has the
form

\[
\hat{m}(x)=\sum_{j=1}^n w_j(x)y_j,
\]

where \(w_j(x)\ge 0\) and \(\sum_j w_j(x)=1\). The weights depend on the
similarity between \(x\) and training points \(x_j\).

Attention has the same weighted-average form:

\[
z_i = \sum_{j=1}^T a_{ij}v_j.
\]

The difference is that the similarity function is learned. Instead of
choosing a fixed kernel, the model learns \(W_Q\) and \(W_K\), and
therefore learns how to compare tokens.

This viewpoint is useful:

\begin{itemize}
\tightlist
\item
  \(q_i\) is the location where the model wants information;
\item
  \(k_j\) determines how relevant position \(j\) is to that location;
\item
  \(v_j\) is the response or information to be averaged;
\item
  \(a_{ij}\) is a data-dependent weight.
\end{itemize}

\section{13.7 Masking}\label{masking}

Not every token should always attend to every other token. A mask is a
matrix \(M\in R^{T\times T}\) added to the scaled scores before the
softmax:

\[
A=\operatorname{softmax}_{\text{row}}\left(\frac{QK^T}{\sqrt{d_k}}+M\right).
\]

Allowed positions have mask value \(0\). Disallowed positions receive a
very negative number, such as \(-10^9\):

\[
M_{ij}=\begin{cases}
0, & \text{if position } j \text{ is allowed for query } i,\\
-10^9, & \text{otherwise.}
\end{cases}
\]

After the softmax, masked positions receive attention weight
approximately \(0\).

\subsection{Padding mask}\label{padding-mask}

Sequences in a batch often have different lengths. Shorter sequences are
padded with artificial tokens. A padding mask prevents real tokens from
attending to padding positions.

\subsection{Causal mask}\label{causal-mask}

In autoregressive language modeling, token \(i\) should not attend to
future tokens \(j>i\). The causal mask is

\[
M_{ij}=\begin{cases}
0, & j\le i,\\
-10^9, & j>i.
\end{cases}
\]

This makes the attention matrix lower triangular after the softmax, up
to numerical precision.

\section{13.8 Numerical stability}\label{numerical-stability}

The softmax should be computed using the stable formula

\[
\operatorname{softmax}(r)_j
=
\frac{\exp(r_j-c)}{\sum_{\ell}\exp(r_{\ell}-c)},
\qquad
c=\max_{\ell} r_{\ell}.
\]

Subtracting \(c\) does not change the softmax because it multiplies
numerator and denominator by the same constant. It prevents overflow
when some scores are large.

For a score row \(r\in R^T\), the Jacobian of the softmax is

\[
\frac{\partial a_j}{\partial r_\ell}
= a_j(\delta_{j\ell}-a_\ell).
\]

This formula shows why extremely peaked attention can make learning
difficult: when one \(a_j\) is almost \(1\) and the rest are almost
\(0\), many derivatives are close to \(0\).

\section{13.9 Worked example}\label{worked-example}

Suppose a sequence has \(T=3\) positions and \(d_k=2\). Let

\[
Q=K=
\begin{bmatrix}
1 & 0\\
0 & 1\\
1 & 1
\end{bmatrix}.
\]

Then

\[
QK^T
=
\begin{bmatrix}
1 & 0 & 1\\
0 & 1 & 1\\
1 & 1 & 2
\end{bmatrix}.
\]

Since \(d_k=2\), the scaled score matrix is

\[
\tilde{S}=\frac{1}{\sqrt{2}}
\begin{bmatrix}
1 & 0 & 1\\
0 & 1 & 1\\
1 & 1 & 2
\end{bmatrix}.
\]

The first row assigns equal raw similarity to positions \(1\) and \(3\),
and lower similarity to position \(2\). After the row-wise softmax, the
first output vector becomes a weighted average of \(v_1\), \(v_2\), and
\(v_3\), with larger weights on \(v_1\) and \(v_3\).

This example illustrates a key point: attention is not simply copying
one token. Unless the softmax is very sharp, attention blends
information from multiple positions.

\section{13.10 Python: scaled dot-product attention from
scratch}\label{python-scaled-dot-product-attention-from-scratch}

\protect\phantomsection\label{py-ch13-scaled-dot-product-attention}
\begin{Shaded}
\begin{Highlighting}[]
\ImportTok{import}\NormalTok{ numpy }\ImportTok{as}\NormalTok{ np}
\ImportTok{import}\NormalTok{ pandas }\ImportTok{as}\NormalTok{ pd}

\NormalTok{np.set\_printoptions(precision}\OperatorTok{=}\DecValTok{3}\NormalTok{, suppress}\OperatorTok{=}\VariableTok{True}\NormalTok{)}

\KeywordTok{def}\NormalTok{ row\_softmax(S):}
    \CommentTok{"""Stable row{-}wise softmax."""}
\NormalTok{    S\_shifted }\OperatorTok{=}\NormalTok{ S }\OperatorTok{{-}}\NormalTok{ S.}\BuiltInTok{max}\NormalTok{(axis}\OperatorTok{=}\DecValTok{1}\NormalTok{, keepdims}\OperatorTok{=}\VariableTok{True}\NormalTok{)}
\NormalTok{    E }\OperatorTok{=}\NormalTok{ np.exp(S\_shifted)}
    \ControlFlowTok{return}\NormalTok{ E }\OperatorTok{/}\NormalTok{ E.}\BuiltInTok{sum}\NormalTok{(axis}\OperatorTok{=}\DecValTok{1}\NormalTok{, keepdims}\OperatorTok{=}\VariableTok{True}\NormalTok{)}

\KeywordTok{def}\NormalTok{ attention(Q, K, V, mask}\OperatorTok{=}\VariableTok{None}\NormalTok{):}
    \CommentTok{"""Scaled dot{-}product attention.}

\CommentTok{    Q: shape (T\_q, d\_k)}
\CommentTok{    K: shape (T\_k, d\_k)}
\CommentTok{    V: shape (T\_k, d\_v)}
\CommentTok{    mask: optional shape (T\_q, T\_k), added to scores before softmax}
\CommentTok{    """}
\NormalTok{    d\_k }\OperatorTok{=}\NormalTok{ Q.shape[}\DecValTok{1}\NormalTok{]}
\NormalTok{    scores }\OperatorTok{=}\NormalTok{ Q }\OperatorTok{@}\NormalTok{ K.T }\OperatorTok{/}\NormalTok{ np.sqrt(d\_k)}
    \ControlFlowTok{if}\NormalTok{ mask }\KeywordTok{is} \KeywordTok{not} \VariableTok{None}\NormalTok{:}
\NormalTok{        scores }\OperatorTok{=}\NormalTok{ scores }\OperatorTok{+}\NormalTok{ mask}
\NormalTok{    A }\OperatorTok{=}\NormalTok{ row\_softmax(scores)}
\NormalTok{    Z }\OperatorTok{=}\NormalTok{ A }\OperatorTok{@}\NormalTok{ V}
    \ControlFlowTok{return}\NormalTok{ Z, A, scores}

\NormalTok{Q }\OperatorTok{=}\NormalTok{ np.array([[}\FloatTok{1.0}\NormalTok{, }\FloatTok{0.0}\NormalTok{],}
\NormalTok{              [}\FloatTok{0.0}\NormalTok{, }\FloatTok{1.0}\NormalTok{],}
\NormalTok{              [}\FloatTok{1.0}\NormalTok{, }\FloatTok{1.0}\NormalTok{]])}
\NormalTok{K }\OperatorTok{=}\NormalTok{ Q.copy()}
\NormalTok{V }\OperatorTok{=}\NormalTok{ np.array([[}\FloatTok{1.0}\NormalTok{, }\FloatTok{0.0}\NormalTok{],}
\NormalTok{              [}\FloatTok{0.0}\NormalTok{, }\FloatTok{1.0}\NormalTok{],}
\NormalTok{              [}\FloatTok{1.0}\NormalTok{, }\FloatTok{1.0}\NormalTok{]])}

\NormalTok{Z, A, scores }\OperatorTok{=}\NormalTok{ attention(Q, K, V)}

\BuiltInTok{print}\NormalTok{(}\StringTok{"scaled scores"}\NormalTok{)}
\BuiltInTok{print}\NormalTok{(scores)}
\BuiltInTok{print}\NormalTok{(}\StringTok{"attention weights"}\NormalTok{)}
\BuiltInTok{print}\NormalTok{(A)}
\BuiltInTok{print}\NormalTok{(}\StringTok{"output"}\NormalTok{)}
\BuiltInTok{print}\NormalTok{(Z)}
\end{Highlighting}
\end{Shaded}

\begin{verbatim}
scaled scores
[[0.707 0.    0.707]
 [0.    0.707 0.707]
 [0.707 0.707 1.414]]
attention weights
[[0.401 0.198 0.401]
 [0.198 0.401 0.401]
 [0.248 0.248 0.503]]
output
[[0.802 0.599]
 [0.599 0.802]
 [0.752 0.752]]
\end{verbatim}

The rows of \texttt{A} should sum to \(1\).

\begin{Shaded}
\begin{Highlighting}[]
\NormalTok{A.}\BuiltInTok{sum}\NormalTok{(axis}\OperatorTok{=}\DecValTok{1}\NormalTok{)}
\end{Highlighting}
\end{Shaded}

\protect\phantomsection\label{py-ch13-row-sums}
\begin{verbatim}
array([1., 1., 1.])
\end{verbatim}

\section{13.11 Python: attention with learned
projections}\label{python-attention-with-learned-projections}

In a neural network, \(Q\), \(K\), and \(V\) are usually obtained from
an input matrix \(X\) by learned linear maps.

\protect\phantomsection\label{py-ch13-projections}
\begin{Shaded}
\begin{Highlighting}[]
\NormalTok{rng }\OperatorTok{=}\NormalTok{ np.random.default\_rng(}\DecValTok{7243}\NormalTok{)}

\NormalTok{T }\OperatorTok{=} \DecValTok{5}
\NormalTok{D\_model }\OperatorTok{=} \DecValTok{4}
\NormalTok{d\_k }\OperatorTok{=} \DecValTok{3}
\NormalTok{d\_v }\OperatorTok{=} \DecValTok{2}

\NormalTok{X }\OperatorTok{=}\NormalTok{ rng.normal(size}\OperatorTok{=}\NormalTok{(T, D\_model))}
\NormalTok{W\_Q }\OperatorTok{=}\NormalTok{ rng.normal(scale}\OperatorTok{=}\FloatTok{0.3}\NormalTok{, size}\OperatorTok{=}\NormalTok{(D\_model, d\_k))}
\NormalTok{W\_K }\OperatorTok{=}\NormalTok{ rng.normal(scale}\OperatorTok{=}\FloatTok{0.3}\NormalTok{, size}\OperatorTok{=}\NormalTok{(D\_model, d\_k))}
\NormalTok{W\_V }\OperatorTok{=}\NormalTok{ rng.normal(scale}\OperatorTok{=}\FloatTok{0.3}\NormalTok{, size}\OperatorTok{=}\NormalTok{(D\_model, d\_v))}

\NormalTok{Q }\OperatorTok{=}\NormalTok{ X }\OperatorTok{@}\NormalTok{ W\_Q}
\NormalTok{K }\OperatorTok{=}\NormalTok{ X }\OperatorTok{@}\NormalTok{ W\_K}
\NormalTok{V }\OperatorTok{=}\NormalTok{ X }\OperatorTok{@}\NormalTok{ W\_V}

\NormalTok{Z, A, scores }\OperatorTok{=}\NormalTok{ attention(Q, K, V)}

\BuiltInTok{print}\NormalTok{(}\StringTok{"X shape:"}\NormalTok{, X.shape)}
\BuiltInTok{print}\NormalTok{(}\StringTok{"Q shape:"}\NormalTok{, Q.shape)}
\BuiltInTok{print}\NormalTok{(}\StringTok{"K shape:"}\NormalTok{, K.shape)}
\BuiltInTok{print}\NormalTok{(}\StringTok{"V shape:"}\NormalTok{, V.shape)}
\BuiltInTok{print}\NormalTok{(}\StringTok{"A shape:"}\NormalTok{, A.shape)}
\BuiltInTok{print}\NormalTok{(}\StringTok{"Z shape:"}\NormalTok{, Z.shape)}
\end{Highlighting}
\end{Shaded}

\begin{verbatim}
X shape: (5, 4)
Q shape: (5, 3)
K shape: (5, 3)
V shape: (5, 2)
A shape: (5, 5)
Z shape: (5, 2)
\end{verbatim}

Dimension checking is not bookkeeping only. It is the fastest way to
find conceptual mistakes in an attention implementation.

\section{13.12 Python: causal mask}\label{python-causal-mask}

A causal mask prevents position \(i\) from using information at future
positions \(j>i\).

\begin{Shaded}
\begin{Highlighting}[]
\NormalTok{T }\OperatorTok{=} \DecValTok{5}
\NormalTok{causal\_mask }\OperatorTok{=}\NormalTok{ np.triu(np.ones((T, T)), k}\OperatorTok{=}\DecValTok{1}\NormalTok{)}
\NormalTok{causal\_mask }\OperatorTok{=}\NormalTok{ np.where(causal\_mask }\OperatorTok{==} \DecValTok{1}\NormalTok{, }\OperatorTok{{-}}\FloatTok{1e9}\NormalTok{, }\FloatTok{0.0}\NormalTok{)}

\NormalTok{Z\_causal, A\_causal, scores\_causal }\OperatorTok{=}\NormalTok{ attention(Q, K, V, mask}\OperatorTok{=}\NormalTok{causal\_mask)}

\NormalTok{pd.DataFrame(A\_causal, columns}\OperatorTok{=}\NormalTok{[}\SpecialStringTok{f"key }\SpecialCharTok{\{}\NormalTok{j}\SpecialCharTok{\}}\SpecialStringTok{"} \ControlFlowTok{for}\NormalTok{ j }\KeywordTok{in} \BuiltInTok{range}\NormalTok{(T)],}
\NormalTok{             index}\OperatorTok{=}\NormalTok{[}\SpecialStringTok{f"query }\SpecialCharTok{\{}\NormalTok{i}\SpecialCharTok{\}}\SpecialStringTok{"} \ControlFlowTok{for}\NormalTok{ i }\KeywordTok{in} \BuiltInTok{range}\NormalTok{(T)])}
\end{Highlighting}
\end{Shaded}

\protect\phantomsection\label{py-ch13-causal-mask}
{\def\LTcaptype{none} % do not increment counter
\begin{longtable}[]{@{}llllll@{}}
\toprule\noalign{}
& key 0 & key 1 & key 2 & key 3 & key 4 \\
\midrule\noalign{}
\endhead
\bottomrule\noalign{}
\endlastfoot
query 0 & 1.000000 & 0.000000 & 0.000000 & 0.000000 & 0.000000 \\
query 1 & 0.489487 & 0.510513 & 0.000000 & 0.000000 & 0.000000 \\
query 2 & 0.332598 & 0.322160 & 0.345243 & 0.000000 & 0.000000 \\
query 3 & 0.221961 & 0.237611 & 0.270132 & 0.270296 & 0.000000 \\
query 4 & 0.213651 & 0.216603 & 0.183614 & 0.166865 & 0.219267 \\
\end{longtable}
}

The upper triangular entries are approximately zero.

\begin{Shaded}
\begin{Highlighting}[]
\NormalTok{np.triu(A\_causal, k}\OperatorTok{=}\DecValTok{1}\NormalTok{).}\BuiltInTok{max}\NormalTok{()}
\end{Highlighting}
\end{Shaded}

\protect\phantomsection\label{py-ch13-causal-check}
\begin{verbatim}
0.0
\end{verbatim}

\section{13.13 Python: inspecting token-token
attention}\label{python-inspecting-token-token-attention}

The following example uses a tiny sequence. The numbers are random, so
they should not be interpreted semantically. The purpose is to inspect
the shape and row-stochastic structure of the attention matrix.

\begin{Shaded}
\begin{Highlighting}[]
\NormalTok{tokens }\OperatorTok{=}\NormalTok{ [}\StringTok{"the"}\NormalTok{, }\StringTok{"model"}\NormalTok{, }\StringTok{"attends"}\NormalTok{, }\StringTok{"to"}\NormalTok{, }\StringTok{"tokens"}\NormalTok{]}
\NormalTok{attention\_table }\OperatorTok{=}\NormalTok{ pd.DataFrame(A\_causal, columns}\OperatorTok{=}\NormalTok{tokens, index}\OperatorTok{=}\NormalTok{tokens)}
\NormalTok{attention\_table}
\end{Highlighting}
\end{Shaded}

\protect\phantomsection\label{py-ch13-token-attention-table}
{\def\LTcaptype{none} % do not increment counter
\begin{longtable}[]{@{}llllll@{}}
\toprule\noalign{}
& the & model & attends & to & tokens \\
\midrule\noalign{}
\endhead
\bottomrule\noalign{}
\endlastfoot
the & 1.000000 & 0.000000 & 0.000000 & 0.000000 & 0.000000 \\
model & 0.489487 & 0.510513 & 0.000000 & 0.000000 & 0.000000 \\
attends & 0.332598 & 0.322160 & 0.345243 & 0.000000 & 0.000000 \\
to & 0.221961 & 0.237611 & 0.270132 & 0.270296 & 0.000000 \\
tokens & 0.213651 & 0.216603 & 0.183614 & 0.166865 & 0.219267 \\
\end{longtable}
}

A table of attention weights can be useful, but it must be interpreted
carefully. Attention weights show one internal routing mechanism; they
are not automatically a complete explanation of the model's prediction.

\section{13.14 Interactive HTML: scaled attention
explorer}\label{interactive-html-scaled-attention-explorer}

The following interactive component lets you vary the scale of the
dot-product scores and choose whether to apply a causal mask. It uses a
local Plotly file from \texttt{assets/plotly.min.js}, so it does not
depend on a CDN.

The slider changes the sharpness of the attention distribution. Low
scale gives diffuse averaging. High scale gives attention weights closer
to hard selection.

\section{13.15 Multi-head attention
preview}\label{multi-head-attention-preview}

A single attention head compares tokens in one learned query-key space.
Multi-head attention uses several heads in parallel:

\[
\operatorname{head}_h
=
\operatorname{Attention}(XW_Q^{(h)},XW_K^{(h)},XW_V^{(h)}).
\]

The heads are concatenated and projected:

\[
\operatorname{MultiHead}(X)
=
\operatorname{Concat}(\operatorname{head}_1,\ldots,\operatorname{head}_H)W_O.
\]

Different heads can learn different comparison rules. One head may
attend to local syntax, another to long-range agreement, and another to
delimiter or position-like patterns. The next chapter studies
self-attention and multi-head attention in more detail.

\section{13.16 AI-assisted learning: attention
audit}\label{ai-assisted-learning-attention-audit}

Use an AI system as a mathematical critic, not as an authority.

\subsection{Prompt 1: dimension
checking}\label{prompt-1-dimension-checking-1}

Ask:

\begin{quote}
I have \(X\in R^{T\times d_{\mathrm{model}}}\),
\(W_Q,W_K\in R^{d_{\mathrm{model}}\times d_k}\), and
\(W_V\in R^{d_{\mathrm{model}}\times d_v}\). Check every dimension in
scaled dot-product attention and explain the shape of the final output.
\end{quote}

Then verify every matrix product yourself.

\subsection{Prompt 2: interpretation
critique}\label{prompt-2-interpretation-critique}

Ask:

\begin{quote}
Explain why attention weights should not automatically be interpreted as
a complete explanation of a transformer model's prediction.
\end{quote}

A good answer should distinguish internal routing from causal
explanation.

\subsection{Prompt 3: mask debugging}\label{prompt-3-mask-debugging}

Ask:

\begin{quote}
I implemented a causal mask for self-attention. What numerical checks
should I run to verify that future positions have zero attention weight?
\end{quote}

You should expect checks involving the upper triangular part of the
attention matrix.

\section{13.17 Common mistakes}\label{common-mistakes-2}

\begin{enumerate}
\def\labelenumi{\arabic{enumi}.}
\tightlist
\item
  \textbf{Confusing keys and values.} Keys determine weights; values
  carry the information being averaged.
\item
  \textbf{Applying softmax over the wrong axis.} In standard attention,
  softmax is applied row by row over keys for each query.
\item
  \textbf{Forgetting the scaling factor.} Without \(1/\sqrt{d_k}\),
  scores may become too large when \(d_k\) is large.
\item
  \textbf{Using a mask after softmax.} Masks should be applied to scores
  before softmax, followed by renormalization.
\item
  \textbf{Ignoring dimensions.} Most attention implementation errors are
  shape errors in disguise.
\end{enumerate}

\section{13.18 Exercises}\label{exercises-12}

\begin{enumerate}
\def\labelenumi{\arabic{enumi}.}
\item
  \textbf{Dimension check.} Suppose \(T=12\), \(d_{\mathrm{model}}=64\),
  \(d_k=16\), and \(d_v=20\). Write the shapes of \(X\), \(W_Q\),
  \(W_K\), \(W_V\), \(Q\), \(K\), \(V\), \(QK^T\), \(A\), and \(Z\).
\item
  \textbf{Row-stochastic proof.} Prove that each row of
  \(A=\operatorname{softmax}_{\text{row}}(S)\) sums to \(1\) and has
  nonnegative entries.
\item
  \textbf{Convex combination.} Prove that each attention output \(z_i\)
  lies in the convex hull of the value vectors \(v_1,\ldots,v_T\).
\item
  \textbf{Scaling argument.} Assume the coordinates of \(q\) and \(k\)
  are independent with mean \(0\) and variance \(1\). Compute the mean
  and variance of \(q^Tk\).
\item
  \textbf{Masking.} For \(T=4\), write the causal mask matrix using
  \(0\) and \(-\infty\) notation.
\item
  \textbf{Attention versus kernel smoothing.} Compare attention with
  Nadaraya-Watson kernel regression. What is fixed in kernel regression
  but learned in attention?
\item
  \textbf{Implementation.} Implement attention in NumPy. Verify that the
  attention matrix is row-stochastic before and after applying a causal
  mask.
\item
  \textbf{Interpretation.} Give an example where a token receives high
  attention weight but is not necessarily the sole explanation for the
  model's final prediction.
\end{enumerate}

\section{13.19 Mini-project}\label{mini-project-1}

Build a small attention visualizer for a short sentence. Use randomly
initialized query, key, and value matrices first. Then manually choose
query and key vectors that make one token attend strongly to another
token. Your report should include:

\begin{itemize}
\tightlist
\item
  the sentence and token list;
\item
  the matrices \(Q\), \(K\), and \(V\);
\item
  the scaled score matrix;
\item
  the attention matrix;
\item
  a heatmap visualization;
\item
  a short explanation of how changing the scale changes entropy;
\item
  a comparison of unmasked and causal attention.
\end{itemize}

\section{13.20 Python lab connection}\label{python-lab-connection-6}

The associated lab for this chapter is planned as:

\begin{Shaded}
\begin{Highlighting}[]
\NormalTok{labs/chapter{-}13{-}attention{-}mechanisms.ipynb}
\end{Highlighting}
\end{Shaded}

The lab should be written for independent study and should include:

\begin{itemize}
\tightlist
\item
  a review of matrix multiplication and row-wise softmax;
\item
  attention from scratch in NumPy;
\item
  masking examples;
\item
  numerical stability checks;
\item
  attention heatmaps;
\item
  a short bridge to multi-head self-attention.
\end{itemize}

\section{13.21 What to remember}\label{what-to-remember-3}

Attention is a learned, data-dependent weighted-average operator. The
mathematical structure is simple and powerful:

\[
Q=XW_Q,
\qquad
K=XW_K,
\qquad
V=XW_V,
\]

\[
A=\operatorname{softmax}_{\text{row}}\left(\frac{QK^T}{\sqrt{d_k}}+M\right),
\qquad
Z=AV.
\]

The matrix \(A\) determines how information moves among token positions.
Because \(A\) is computed from the sequence itself, attention can adapt
its information flow to each input.

\section{Notes and references}\label{notes-and-references-7}

Attention was introduced into neural sequence modeling before the
transformer architecture and became central in encoder-decoder models
for machine translation. Scaled dot-product attention and multi-head
attention are the core operations in the transformer architecture.
Useful references include Bahdanau et al. (2015), Vaswani et al. (2017),
Jurafsky and Martin (2024), and Goodfellow et al. (2016b).

\chapter{Self-Attention}\label{self-attention}

Contextual token representations, permutation equivariance, and
multi-head structure

\hfill\break

Self-attention is the operation that made the Transformer architecture
possible. In an RNN, information moves through a time-indexed hidden
state. In self-attention, all token positions in a sequence compare
themselves with all other token positions directly. The result is a new
representation for each token, called a \textbf{contextual embedding},
whose meaning depends on the entire sequence.

The mathematical object in this chapter is the self-attention map

\[
X \mapsto \operatorname{softmax}_{\text{row}}\left(\frac{XW_QW_K^TX^T}{\sqrt{d_k}}+M\right)XW_V.
\]

This formula is compact, but it hides several important ideas: learned
pairwise similarity, row-stochastic averaging, masking, positional
information, permutation equivariance, and parallel attention heads. We
study each idea carefully.

\section{Learning goals}\label{learning-goals-13}

After studying this chapter, you should be able to:

\begin{enumerate}
\def\labelenumi{\arabic{enumi}.}
\tightlist
\item
  define self-attention as attention with \(Q,K,V\) all computed from
  the same sequence matrix \(X\);
\item
  check every matrix dimension in a self-attention layer;
\item
  explain why self-attention produces contextual token embeddings;
\item
  prove the permutation equivariance of self-attention without
  positional encodings;
\item
  explain why positional information is necessary for language modeling;
\item
  construct padding masks and causal masks;
\item
  interpret multi-head attention as several learned comparison
  geometries in parallel;
\item
  implement single-head and multi-head self-attention from scratch in
  NumPy;
\item
  use attention heatmaps carefully without over-interpreting them;
\item
  explain the computational cost of self-attention.
\end{enumerate}

\section{14.1 From attention to
self-attention}\label{from-attention-to-self-attention}

In Chapter 13, attention was defined for three matrices

\[
Q\in R^{T_q\times d_k},\qquad K\in R^{T_k\times d_k},\qquad V\in R^{T_k\times d_v},
\]

by

\[
\operatorname{Attention}(Q,K,V)
=
\operatorname{softmax}_{\text{row}}\left(\frac{QK^T}{\sqrt{d_k}}\right)V.
\]

In \textbf{self-attention}, the queries, keys, and values are all
computed from the same sequence. Let

\[
X\in R^{T\times d_{\mathrm{model}}}
\]

be the matrix of token representations for a sequence of length \(T\).
The \(i\)th row \(x_i^T\) represents token position \(i\). A
self-attention head uses learned matrices

\[
W_Q\in R^{d_{\mathrm{model}}\times d_k},\qquad
W_K\in R^{d_{\mathrm{model}}\times d_k},\qquad
W_V\in R^{d_{\mathrm{model}}\times d_v}
\]

to form

\[
Q=XW_Q,
\qquad
K=XW_K,
\qquad
V=XW_V.
\]

Then

\[
Z
=
\operatorname{SelfAttn}(X)
=
\operatorname{softmax}_{\text{row}}\left(\frac{XW_QW_K^TX^T}{\sqrt{d_k}}\right)XW_V.
\]

The output \(Z\in R^{T\times d_v}\) has one row for each token position.
Each row is a weighted average of transformed value vectors from the
same sequence.

\textbf{One sentence summary.} Self-attention lets every token position
build a new representation by asking: which other token positions in
this same sequence are relevant to me?

\section{14.2 Token matrices and dimension
bookkeeping}\label{token-matrices-and-dimension-bookkeeping}

A sequence of tokens

\[
w_1,w_2,\ldots,w_T
\]

is first mapped to vectors. If \(E\in R^{|V|\times d_{\mathrm{model}}}\)
is a token embedding matrix and token \(w_i\) has integer index \(r_i\),
then

\[
e_i=E_{r_i,:}\in R^{d_{\mathrm{model}}}.
\]

Usually we also add positional information
\(p_i\in R^{d_{\mathrm{model}}}\):

\[
x_i=e_i+p_i.
\]

Stacking rows gives

\[
X=
\begin{bmatrix}
x_1^T\\
x_2^T\\
\vdots\\
x_T^T
\end{bmatrix}
\in R^{T\times d_{\mathrm{model}}}.
\]

For one attention head:

\[
XW_Q\in R^{T\times d_k},
\qquad
XW_K\in R^{T\times d_k},
\qquad
XW_V\in R^{T\times d_v}.
\]

The score matrix is

\[
S=\frac{QK^T}{\sqrt{d_k}}\in R^{T\times T}.
\]

The attention matrix is

\[
A=\operatorname{softmax}_{\text{row}}(S)\in R^{T\times T}.
\]

The output is

\[
Z=AV\in R^{T\times d_v}.
\]

The row \(z_i^T\) is the new representation of token position \(i\).

\section{14.3 Contextual embeddings}\label{contextual-embeddings}

A static word embedding assigns one vector to each token type. For
example, the word \texttt{bank} might have a single embedding vector.
But the meaning of \texttt{bank} changes in different contexts:

\begin{itemize}
\tightlist
\item
  \texttt{river\ bank}
\item
  \texttt{bank\ loan}
\item
  \texttt{central\ bank}
\end{itemize}

Self-attention produces \textbf{contextual embeddings}. The output
vector \(z_i\) at position \(i\) depends not only on \(x_i\), but also
on all rows of \(X\) through the attention weights.

Row \(i\) of the output is

\[
z_i
=
\sum_{j=1}^T a_{ij}v_j,
\]

where

\[
a_{ij}
=
\frac{\exp(q_i^Tk_j/\sqrt{d_k})}
{\sum_{\ell=1}^T \exp(q_i^Tk_\ell/\sqrt{d_k})}.
\]

Since \(q_i\), \(k_j\), and \(v_j\) are all functions of the input token
representations, the vector \(z_i\) changes when the surrounding
sentence changes.

\textbf{Contextualization.} A token representation is contextual if
changing another token in the same sequence can change the
representation at the original position.

\subsection{A small mathematical test for
contextuality}\label{a-small-mathematical-test-for-contextuality}

Let \(F:R^{T\times d}\to R^{T\times d_v}\) be a sequence model. The
representation at position \(i\) is contextual if there exists some
\(j\ne i\) such that

\[
\frac{\partial F_i(X)}{\partial x_j}\ne 0.
\]

For self-attention, this derivative is usually nonzero because the
scores \(q_i^Tk_j\) and the value \(v_j\) both depend on \(x_j\).

\section{14.4 Self-attention as a learned graph on
tokens}\label{self-attention-as-a-learned-graph-on-tokens}

For a fixed input sequence \(X\), the attention matrix

\[
A(X)\in R^{T\times T}
\]

can be viewed as the weighted adjacency matrix of a directed graph on
token positions. The edge weight from query position \(i\) to key/value
position \(j\) is \(a_{ij}\).

This graph is:

\begin{itemize}
\tightlist
\item
  \textbf{weighted}, because \(a_{ij}\in[0,1]\);
\item
  \textbf{directed}, because usually \(a_{ij}\ne a_{ji}\);
\item
  \textbf{row-stochastic}, because \(\sum_j a_{ij}=1\);
\item
  \textbf{data-dependent}, because \(A\) changes with \(X\).
\end{itemize}

Self-attention therefore constructs a sequence-specific graph and then
performs a weighted message-passing step:

\[
z_i=\sum_{j=1}^T a_{ij}v_j.
\]

This connects self-attention with graph neural networks. The difference
is that the graph is not given externally; it is inferred from learned
similarities inside the sequence.

\section{14.5 Permutation equivariance without positional
information}\label{permutation-equivariance-without-positional-information}

A subtle but important mathematical property of self-attention is
permutation equivariance. Without positional information or masks,
self-attention does not know the order of tokens. It only sees a set of
input rows.

Let \(P\in R^{T\times T}\) be a permutation matrix. The matrix \(PX\) is
obtained by permuting the rows of \(X\). Define one-head self-attention
by

\[
F(X)=\operatorname{softmax}_{\text{row}}\left(\frac{XW_QW_K^TX^T}{\sqrt{d_k}}\right)XW_V.
\]

Then

\[
F(PX)=PF(X).
\]

This property is called \textbf{permutation equivariance}.

\subsection{Proof}\label{proof}

Let

\[
Q=XW_Q,
\qquad
K=XW_K,
\qquad
V=XW_V.
\]

For the permuted input \(PX\),

\[
Q'=PXW_Q=PQ,
\qquad
K'=PXW_K=PK,
\qquad
V'=PXW_V=PV.
\]

The score matrix becomes

\[
S'
=
\frac{Q'(K')^T}{\sqrt{d_k}}
=
\frac{(PQ)(PK)^T}{\sqrt{d_k}}
=
\frac{P Q K^T P^T}{\sqrt{d_k}}
=
P S P^T.
\]

Row-wise softmax commutes with simultaneous row and column permutation,
so

\[
\operatorname{softmax}_{\text{row}}(S')
=
P\operatorname{softmax}_{\text{row}}(S)P^T
=
PAP^T.
\]

Therefore

\[
F(PX)
=
(PAP^T)(PV)
=PA(P^TP)V
=PAV
=PF(X).
\]

This proves the claim.

\subsection{Why this matters}\label{why-this-matters}

A language model must distinguish

\begin{Shaded}
\begin{Highlighting}[]
\NormalTok{student studies model}
\end{Highlighting}
\end{Shaded}

from

\begin{Shaded}
\begin{Highlighting}[]
\NormalTok{model studies student}
\end{Highlighting}
\end{Shaded}

because the order changes the meaning. Self-attention without positional
information treats these as row permutations of the same collection of
token vectors. The output rows are permuted in the same way, but the
model has no intrinsic sense of first, second, or third position.

Therefore Transformers add positional information.

\section{14.6 Positional information breaks pure permutation
symmetry}\label{positional-information-breaks-pure-permutation-symmetry}

A common input representation is

\[
x_i=e(w_i)+p_i,
\]

where \(e(w_i)\) is the token embedding and \(p_i\) is a position
vector. The sequence matrix is then

\[
X=
\begin{bmatrix}
(e(w_1)+p_1)^T\\
\cdots\\
(e(w_T)+p_T)^T
\end{bmatrix}.
\]

Now token identity and token position are mixed in each row. The model
can learn different behaviors for early positions, nearby positions, or
relative patterns.

There are several common approaches:

\begin{itemize}
\tightlist
\item
  learned absolute position embeddings;
\item
  sinusoidal position encodings;
\item
  relative position biases;
\item
  rotary position embeddings.
\end{itemize}

Chapter 16 studies these in more detail. For now, the key point is
simple: self-attention provides token-token interaction, but positional
information provides order.

\section{14.7 Masked self-attention}\label{masked-self-attention}

A mask is added to the score matrix before softmax:

\[
A
=
\operatorname{softmax}_{\text{row}}
\left(
\frac{QK^T}{\sqrt{d_k}}+M
\right).
\]

Allowed entries of \(M\) are \(0\). Disallowed entries are large
negative numbers.

\subsection{Padding mask}\label{padding-mask-1}

A padding mask prevents real tokens from attending to artificial padding
tokens. If token position \(j\) is padding, then set

\[
M_{ij}=-10^9
\]

for every query position \(i\).

\subsection{Causal mask}\label{causal-mask-1}

A causal mask prevents position \(i\) from attending to future positions
\(j>i\):

\[
M_{ij}=\begin{cases}
0, & j\le i,\\
-10^9, & j>i.
\end{cases}
\]

Causal masking is necessary for autoregressive language modeling. When
predicting token \(w_{t+1}\), the representation at position \(t\) must
not use information from tokens after \(t\).

The causal factorization is

\[
P(w_1,\ldots,w_T)
=
\prod_{t=1}^T P(w_t\mid w_1,\ldots,w_{t-1}).
\]

Causal self-attention enforces this conditional structure inside the
network.

\section{14.8 Multi-head
self-attention}\label{multi-head-self-attention}

A single attention head uses one learned query-key geometry. Multi-head
attention uses \(H\) such heads in parallel.

For head \(h\), define

\[
Q^{(h)}=XW_Q^{(h)},
\qquad
K^{(h)}=XW_K^{(h)},
\qquad
V^{(h)}=XW_V^{(h)}.
\]

The head output is

\[
Z^{(h)}
=
\operatorname{softmax}_{\text{row}}
\left(
\frac{Q^{(h)}(K^{(h)})^T}{\sqrt{d_k}}+M
\right)V^{(h)}.
\]

The head outputs are concatenated:

\[
Z_{\mathrm{cat}}
=
\operatorname{Concat}(Z^{(1)},\ldots,Z^{(H)})
\in R^{T\times Hd_v}.
\]

Then an output projection maps back to the model dimension:

\[
\operatorname{MHA}(X)
=
Z_{\mathrm{cat}}W_O,
\qquad
W_O\in R^{Hd_v\times d_{\mathrm{model}}}.
\]

If \(d_v=d_k=d_{\mathrm{model}}/H\), then the concatenated output has
dimension \(d_{\mathrm{model}}\), which is convenient for residual
connections.

\textbf{Multi-head viewpoint.} Each head learns a different similarity
geometry. One head may emphasize nearby tokens, another may emphasize
matching syntax, and another may emphasize semantic roles. The
architecture does not force these roles; they emerge, if useful, from
the training objective.

\section{14.9 Self-attention as matrix-valued feature
construction}\label{self-attention-as-matrix-valued-feature-construction}

From a statistical learning point of view, a self-attention layer is a
data-dependent feature map.

For a fixed parameter collection \(\theta=(W_Q,W_K,W_V)\),

\[
F_\theta(X)=A_\theta(X)XW_V,
\]

where

\[
A_\theta(X)
=
\operatorname{softmax}_{\text{row}}
\left(\frac{XW_QW_K^TX^T}{\sqrt{d_k}}\right).
\]

This is nonlinear in \(X\) because \(A_\theta(X)\) itself depends on
\(X\). However, conditional on the attention matrix \(A_\theta(X)\), the
output is linear in the values:

\[
Z=A_\theta(X)V.
\]

So self-attention has two layers of structure:

\begin{enumerate}
\def\labelenumi{\arabic{enumi}.}
\tightlist
\item
  construct a similarity matrix from \(X\);
\item
  use that matrix to average transformed vectors.
\end{enumerate}

This decomposition is often the cleanest way to understand the
operation.

\section{14.10 Attention entropy and
sharpness}\label{attention-entropy-and-sharpness}

For each query position \(i\), the attention row
\(a_i=(a_{i1},\ldots,a_{iT})\) is a probability vector. Its entropy is

\[
H(a_i)=-\sum_{j=1}^T a_{ij}\log a_{ij}.
\]

Low entropy means the query attends mostly to a few positions. High
entropy means the query averages broadly across the sequence.

If the scores are multiplied by a scale parameter \(\gamma>0\),

\[
a_{ij}(\gamma)
=
\frac{\exp(\gamma s_{ij})}{\sum_{\ell=1}^T \exp(\gamma s_{i\ell})},
\]

then increasing \(\gamma\) usually makes attention sharper. Decreasing
\(\gamma\) makes attention more uniform.

This is one reason the factor \(1/\sqrt{d_k}\) matters: it controls the
typical scale of the scores before the softmax.

\section{14.11 Computational complexity}\label{computational-complexity}

For sequence length \(T\) and key dimension \(d_k\), the score matrix

\[
QK^T
\]

requires roughly \(O(T^2d_k)\) arithmetic operations and stores
\(O(T^2)\) attention weights. The multiplication

\[
AV
\]

requires roughly \(O(T^2d_v)\) operations.

Thus standard full self-attention scales quadratically in sequence
length. This is powerful for moderate sequences but expensive for very
long documents. Many efficient Transformer variants modify the attention
pattern to reduce this cost.

For this book, the main mathematical point is:

The power of full self-attention comes from all-pairs token interaction.
The cost of full self-attention also comes from all-pairs token
interaction.

\section{14.12 Worked example: two query
positions}\label{worked-example-two-query-positions}

Suppose a sequence has three value vectors

\[
v_1=(1,0),\qquad v_2=(0,1),\qquad v_3=(1,1).
\]

Assume that after softmax the first two attention rows are

\[
a_1=(0.7,0.2,0.1),
\qquad
a_2=(0.1,0.3,0.6).
\]

Then

\[
z_1=0.7v_1+0.2v_2+0.1v_3
=(0.8,0.3),
\]

and

\[
z_2=0.1v_1+0.3v_2+0.6v_3
=(0.7,0.9).
\]

The same value vectors produce different contextual embeddings because
the attention rows are different.

\section{14.13 Python example: single-head self-attention from
scratch}\label{python-example-single-head-self-attention-from-scratch}

The next code block implements one self-attention head using NumPy.

\begin{Shaded}
\begin{Highlighting}[]
\ImportTok{import}\NormalTok{ numpy }\ImportTok{as}\NormalTok{ np}

\NormalTok{np.set\_printoptions(precision}\OperatorTok{=}\DecValTok{3}\NormalTok{, suppress}\OperatorTok{=}\VariableTok{True}\NormalTok{)}

\KeywordTok{def}\NormalTok{ row\_softmax(S):}
\NormalTok{    S\_shifted }\OperatorTok{=}\NormalTok{ S }\OperatorTok{{-}}\NormalTok{ S.}\BuiltInTok{max}\NormalTok{(axis}\OperatorTok{=}\DecValTok{1}\NormalTok{, keepdims}\OperatorTok{=}\VariableTok{True}\NormalTok{)}
\NormalTok{    E }\OperatorTok{=}\NormalTok{ np.exp(S\_shifted)}
    \ControlFlowTok{return}\NormalTok{ E }\OperatorTok{/}\NormalTok{ E.}\BuiltInTok{sum}\NormalTok{(axis}\OperatorTok{=}\DecValTok{1}\NormalTok{, keepdims}\OperatorTok{=}\VariableTok{True}\NormalTok{)}

\KeywordTok{def}\NormalTok{ self\_attention(X, W\_Q, W\_K, W\_V, mask}\OperatorTok{=}\VariableTok{None}\NormalTok{):}
\NormalTok{    Q }\OperatorTok{=}\NormalTok{ X }\OperatorTok{@}\NormalTok{ W\_Q}
\NormalTok{    K }\OperatorTok{=}\NormalTok{ X }\OperatorTok{@}\NormalTok{ W\_K}
\NormalTok{    V }\OperatorTok{=}\NormalTok{ X }\OperatorTok{@}\NormalTok{ W\_V}
\NormalTok{    scores }\OperatorTok{=}\NormalTok{ Q }\OperatorTok{@}\NormalTok{ K.T }\OperatorTok{/}\NormalTok{ np.sqrt(Q.shape[}\DecValTok{1}\NormalTok{])}
    \ControlFlowTok{if}\NormalTok{ mask }\KeywordTok{is} \KeywordTok{not} \VariableTok{None}\NormalTok{:}
\NormalTok{        scores }\OperatorTok{=}\NormalTok{ scores }\OperatorTok{+}\NormalTok{ mask}
\NormalTok{    A }\OperatorTok{=}\NormalTok{ row\_softmax(scores)}
\NormalTok{    Z }\OperatorTok{=}\NormalTok{ A }\OperatorTok{@}\NormalTok{ V}
    \ControlFlowTok{return}\NormalTok{ Z, A, scores}

\CommentTok{\# Toy token matrix: T = 4 tokens, d\_model = 3 features.}
\NormalTok{X }\OperatorTok{=}\NormalTok{ np.array([}
\NormalTok{    [}\FloatTok{1.0}\NormalTok{, }\FloatTok{0.0}\NormalTok{, }\FloatTok{0.2}\NormalTok{],}
\NormalTok{    [}\FloatTok{0.8}\NormalTok{, }\FloatTok{0.2}\NormalTok{, }\FloatTok{0.1}\NormalTok{],}
\NormalTok{    [}\FloatTok{0.1}\NormalTok{, }\FloatTok{1.0}\NormalTok{, }\FloatTok{0.3}\NormalTok{],}
\NormalTok{    [}\FloatTok{0.0}\NormalTok{, }\FloatTok{0.7}\NormalTok{, }\FloatTok{1.0}\NormalTok{],}
\NormalTok{])}

\NormalTok{W\_Q }\OperatorTok{=}\NormalTok{ np.array([}
\NormalTok{    [}\FloatTok{1.0}\NormalTok{, }\FloatTok{0.0}\NormalTok{],}
\NormalTok{    [}\FloatTok{0.2}\NormalTok{, }\FloatTok{0.8}\NormalTok{],}
\NormalTok{    [}\FloatTok{0.0}\NormalTok{, }\FloatTok{0.5}\NormalTok{],}
\NormalTok{])}
\NormalTok{W\_K }\OperatorTok{=}\NormalTok{ np.array([}
\NormalTok{    [}\FloatTok{0.9}\NormalTok{, }\FloatTok{0.1}\NormalTok{],}
\NormalTok{    [}\FloatTok{0.1}\NormalTok{, }\FloatTok{0.9}\NormalTok{],}
\NormalTok{    [}\FloatTok{0.2}\NormalTok{, }\FloatTok{0.4}\NormalTok{],}
\NormalTok{])}
\NormalTok{W\_V }\OperatorTok{=}\NormalTok{ np.array([}
\NormalTok{    [}\FloatTok{1.0}\NormalTok{, }\FloatTok{0.0}\NormalTok{],}
\NormalTok{    [}\FloatTok{0.0}\NormalTok{, }\FloatTok{1.0}\NormalTok{],}
\NormalTok{    [}\FloatTok{0.5}\NormalTok{, }\FloatTok{0.5}\NormalTok{],}
\NormalTok{])}

\NormalTok{Z, A, scores }\OperatorTok{=}\NormalTok{ self\_attention(X, W\_Q, W\_K, W\_V)}
\BuiltInTok{print}\NormalTok{(}\StringTok{"Attention matrix A:"}\NormalTok{)}
\BuiltInTok{print}\NormalTok{(A)}
\BuiltInTok{print}\NormalTok{(}\StringTok{"}\CharTok{\textbackslash{}n}\StringTok{Output Z:"}\NormalTok{)}
\BuiltInTok{print}\NormalTok{(Z)}
\BuiltInTok{print}\NormalTok{(}\StringTok{"}\CharTok{\textbackslash{}n}\StringTok{Row sums of A:"}\NormalTok{, A.}\BuiltInTok{sum}\NormalTok{(axis}\OperatorTok{=}\DecValTok{1}\NormalTok{))}
\end{Highlighting}
\end{Shaded}

\begin{verbatim}
Attention matrix A:
[[0.312 0.277 0.204 0.207]
 [0.292 0.267 0.22  0.222]
 [0.196 0.204 0.299 0.301]
 [0.178 0.192 0.315 0.315]]

Output Z:
[[0.734 0.583]
 [0.713 0.615]
 [0.614 0.776]
 [0.595 0.806]]

Row sums of A: [1. 1. 1. 1.]
\end{verbatim}

Each row of \(A\) sums to \(1\). The output row \(z_i\) is a weighted
average of the rows of \(V=XW_V\).

\section{14.14 Python example: causal
masking}\label{python-example-causal-masking}

The causal mask prevents future attention.

\begin{Shaded}
\begin{Highlighting}[]
\NormalTok{T }\OperatorTok{=}\NormalTok{ X.shape[}\DecValTok{0}\NormalTok{]}
\NormalTok{causal\_mask }\OperatorTok{=}\NormalTok{ np.triu(np.ones((T, T)), k}\OperatorTok{=}\DecValTok{1}\NormalTok{)}
\NormalTok{causal\_mask }\OperatorTok{=}\NormalTok{ np.where(causal\_mask }\OperatorTok{==} \DecValTok{1}\NormalTok{, }\OperatorTok{{-}}\FloatTok{1e9}\NormalTok{, }\FloatTok{0.0}\NormalTok{)}

\NormalTok{Z\_causal, A\_causal, scores\_causal }\OperatorTok{=}\NormalTok{ self\_attention(}
\NormalTok{    X, W\_Q, W\_K, W\_V, mask}\OperatorTok{=}\NormalTok{causal\_mask}
\NormalTok{)}

\BuiltInTok{print}\NormalTok{(}\StringTok{"Causal attention matrix:"}\NormalTok{)}
\BuiltInTok{print}\NormalTok{(A\_causal)}
\BuiltInTok{print}\NormalTok{(}\StringTok{"}\CharTok{\textbackslash{}n}\StringTok{Entries above the diagonal are approximately zero."}\NormalTok{)}
\end{Highlighting}
\end{Shaded}

\begin{verbatim}
Causal attention matrix:
[[1.    0.    0.    0.   ]
 [0.522 0.478 0.    0.   ]
 [0.28  0.292 0.428 0.   ]
 [0.178 0.192 0.315 0.315]]

Entries above the diagonal are approximately zero.
\end{verbatim}

The first token can attend only to itself. The second token can attend
to positions \(1\) and \(2\). The last token can attend to all previous
positions and itself.

\section{14.15 Python example: multi-head
self-attention}\label{python-example-multi-head-self-attention}

The following implementation uses a list of projection matrices, one set
per head.

\begin{Shaded}
\begin{Highlighting}[]
\KeywordTok{def}\NormalTok{ multi\_head\_self\_attention(X, heads, W\_O}\OperatorTok{=}\VariableTok{None}\NormalTok{, mask}\OperatorTok{=}\VariableTok{None}\NormalTok{):}
\NormalTok{    outputs }\OperatorTok{=}\NormalTok{ []}
\NormalTok{    attentions }\OperatorTok{=}\NormalTok{ []}
    \ControlFlowTok{for}\NormalTok{ W\_Q\_h, W\_K\_h, W\_V\_h }\KeywordTok{in}\NormalTok{ heads:}
\NormalTok{        Z\_h, A\_h, \_ }\OperatorTok{=}\NormalTok{ self\_attention(X, W\_Q\_h, W\_K\_h, W\_V\_h, mask}\OperatorTok{=}\NormalTok{mask)}
\NormalTok{        outputs.append(Z\_h)}
\NormalTok{        attentions.append(A\_h)}
\NormalTok{    Z\_cat }\OperatorTok{=}\NormalTok{ np.concatenate(outputs, axis}\OperatorTok{=}\DecValTok{1}\NormalTok{)}
    \ControlFlowTok{if}\NormalTok{ W\_O }\KeywordTok{is} \VariableTok{None}\NormalTok{:}
        \ControlFlowTok{return}\NormalTok{ Z\_cat, attentions}
    \ControlFlowTok{return}\NormalTok{ Z\_cat }\OperatorTok{@}\NormalTok{ W\_O, attentions}

\NormalTok{heads }\OperatorTok{=}\NormalTok{ [}
\NormalTok{    (W\_Q, W\_K, W\_V),}
\NormalTok{    (}
\NormalTok{        np.array([[}\FloatTok{0.2}\NormalTok{, }\FloatTok{0.9}\NormalTok{], [}\FloatTok{1.0}\NormalTok{, }\FloatTok{0.0}\NormalTok{], [}\FloatTok{0.3}\NormalTok{, }\FloatTok{0.4}\NormalTok{]]),}
\NormalTok{        np.array([[}\FloatTok{0.1}\NormalTok{, }\FloatTok{0.8}\NormalTok{], [}\FloatTok{0.8}\NormalTok{, }\FloatTok{0.2}\NormalTok{], [}\FloatTok{0.4}\NormalTok{, }\FloatTok{0.6}\NormalTok{]]),}
\NormalTok{        np.array([[}\FloatTok{0.5}\NormalTok{, }\FloatTok{0.5}\NormalTok{], [}\FloatTok{1.0}\NormalTok{, }\FloatTok{0.0}\NormalTok{], [}\FloatTok{0.0}\NormalTok{, }\FloatTok{1.0}\NormalTok{]]),}
\NormalTok{    ),}
\NormalTok{]}

\NormalTok{Z\_multi, A\_heads }\OperatorTok{=}\NormalTok{ multi\_head\_self\_attention(X, heads)}
\BuiltInTok{print}\NormalTok{(}\StringTok{"Multi{-}head output shape:"}\NormalTok{, Z\_multi.shape)}
\BuiltInTok{print}\NormalTok{(}\StringTok{"Head 1 attention:"}\NormalTok{)}
\BuiltInTok{print}\NormalTok{(A\_heads[}\DecValTok{0}\NormalTok{])}
\BuiltInTok{print}\NormalTok{(}\StringTok{"Head 2 attention:"}\NormalTok{)}
\BuiltInTok{print}\NormalTok{(A\_heads[}\DecValTok{1}\NormalTok{])}
\end{Highlighting}
\end{Shaded}

\begin{verbatim}
Multi-head output shape: (4, 4)
Head 1 attention:
[[0.312 0.277 0.204 0.207]
 [0.292 0.267 0.22  0.222]
 [0.196 0.204 0.299 0.301]
 [0.178 0.192 0.315 0.315]]
Head 2 attention:
[[0.266 0.24  0.222 0.272]
 [0.249 0.232 0.239 0.28 ]
 [0.181 0.19  0.304 0.325]
 [0.194 0.198 0.289 0.32 ]]
\end{verbatim}

The two heads may produce different attention patterns because their
projection matrices define different query-key geometries.

\section{14.16 Python example: attention
entropy}\label{python-example-attention-entropy}

Attention rows are probability vectors, so we can compute row entropies.

\begin{Shaded}
\begin{Highlighting}[]
\KeywordTok{def}\NormalTok{ entropy\_rows(A):}
\NormalTok{    A\_safe }\OperatorTok{=}\NormalTok{ np.clip(A, }\FloatTok{1e{-}12}\NormalTok{, }\FloatTok{1.0}\NormalTok{)}
    \ControlFlowTok{return} \OperatorTok{{-}}\NormalTok{(A\_safe }\OperatorTok{*}\NormalTok{ np.log(A\_safe)).}\BuiltInTok{sum}\NormalTok{(axis}\OperatorTok{=}\DecValTok{1}\NormalTok{)}

\BuiltInTok{print}\NormalTok{(}\StringTok{"Unmasked row entropies:"}\NormalTok{, entropy\_rows(A))}
\BuiltInTok{print}\NormalTok{(}\StringTok{"Causal row entropies:"}\NormalTok{, entropy\_rows(A\_causal))}
\end{Highlighting}
\end{Shaded}

\begin{verbatim}
Unmasked row entropies: [1.369 1.379 1.366 1.352]
Causal row entropies: [0.    0.692 1.079 1.352]
\end{verbatim}

Entropy is not automatically good or bad. It is a diagnostic. The
appropriate sharpness depends on the task, the layer, and the head.

\section{14.17 Interactive HTML: self-attention head
explorer}\label{interactive-html-self-attention-head-explorer}

The following interactive component compares three hand-designed
self-attention heads. It uses the local Plotly file
\texttt{assets/plotly.min.js}, so it should work without a network
connection after the book is rendered.

Use this interactive as a mathematical object, not merely a picture. Ask
how the attention matrix changes when scores are sharpened, when future
tokens are masked, and when the head geometry changes.

\section{14.18 AI-assisted learning: self-attention
audit}\label{ai-assisted-learning-self-attention-audit}

AI systems can help you check your reasoning, but they should not
replace it. Use them as mathematical critics.

\subsection{Prompt 1: dimension audit}\label{prompt-1-dimension-audit-1}

Ask:

\begin{quote}
Let \(X\in R^{T\times d_{\mathrm{model}}}\) and let
\(W_Q,W_K\in R^{d_{\mathrm{model}}\times d_k}\),
\(W_V\in R^{d_{\mathrm{model}}\times d_v}\). Check every matrix
dimension in one-head self-attention and explain why the output has
shape \(T\times d_v\).
\end{quote}

Then verify the answer line by line.

\subsection{Prompt 2: proof critique}\label{prompt-2-proof-critique}

Ask:

\begin{quote}
I claim that self-attention without positional encodings is permutation
equivariant. Give a proof using a permutation matrix \(P\), and identify
exactly where row-wise softmax is used.
\end{quote}

Check whether the proof distinguishes invariance from equivariance.

\subsection{Prompt 3: causal-mask
explanation}\label{prompt-3-causal-mask-explanation}

Ask:

\begin{quote}
Explain why causal self-attention is needed for autoregressive language
modeling. Connect the mask to the factorization
\(P(w_1,\ldots,w_T)=\prod_t P(w_t\mid w_1,\ldots,w_{t-1})\).
\end{quote}

Then write your own explanation in one paragraph.

\section{14.19 Common mistakes}\label{common-mistakes-3}

\begin{enumerate}
\def\labelenumi{\arabic{enumi}.}
\tightlist
\item
  \textbf{Confusing attention with explanation.} Attention weights are
  useful diagnostics, but they are not complete causal explanations of a
  model's decision.
\item
  \textbf{Forgetting row-wise softmax.} The softmax is applied
  separately to each query row.
\item
  \textbf{Ignoring positional information.} Self-attention alone is
  permutation equivariant; language requires order.
\item
  \textbf{Using the wrong mask orientation.} In causal attention, query
  position \(i\) may attend to key position \(j\) only when \(j\le i\).
\item
  \textbf{Confusing token type and token position.} The same word
  appearing twice can have two different contextual embeddings.
\item
  \textbf{Overlooking quadratic cost.} Full attention stores a
  \(T\times T\) matrix for each head.
\end{enumerate}

\section{14.20 Exercises}\label{exercises-13}

\begin{enumerate}
\def\labelenumi{\arabic{enumi}.}
\item
  \textbf{Dimension check.} Let \(T=12\), \(d_{\mathrm{model}}=64\),
  \(d_k=16\), and \(d_v=16\). Give the shapes of \(Q,K,V,QK^T,A\), and
  \(Z\) for one attention head.
\item
  \textbf{Contextuality.} Explain why changing token \(x_j\) can change
  output vector \(z_i\) even when \(i\ne j\).
\item
  \textbf{Permutation equivariance.} Prove that one-head self-attention
  without masks and positional embeddings satisfies \(F(PX)=PF(X)\) for
  every permutation matrix \(P\).
\item
  \textbf{Position necessity.} Give two three-token sentences with the
  same words but different order and different meaning. Explain why
  token embeddings alone cannot distinguish them under a
  permutation-equivariant model.
\item
  \textbf{Causal mask.} Write the \(5\times 5\) causal mask pattern
  showing allowed and disallowed entries.
\item
  \textbf{Entropy.} For \(a=(0.7,0.2,0.1)\) and \(b=(1/3,1/3,1/3)\),
  compute \(H(a)\) and \(H(b)\). Which distribution is sharper?
\item
  \textbf{Multi-head dimensions.} Suppose \(d_{\mathrm{model}}=128\),
  \(H=8\), and \(d_v=16\). What is the dimension of the concatenated
  head output? What shape should \(W_O\) have to return to
  \(d_{\mathrm{model}}\)?
\item
  \textbf{Implementation.} Modify the Python function
  \texttt{self\_attention} to return both the unnormalized score matrix
  and the attention matrix.
\item
  \textbf{Mask orientation test.} Construct a toy sequence of length
  \(4\) and verify that your causal mask makes the upper triangular
  entries of \(A\) approximately zero.
\item
  \textbf{Interpretation.} Give one reason an attention heatmap may be
  useful and one reason it may be misleading.
\end{enumerate}

\section{14.21 Mini-project: attention pattern
laboratory}\label{mini-project-attention-pattern-laboratory}

Build a small attention laboratory in Python.

\begin{enumerate}
\def\labelenumi{\arabic{enumi}.}
\tightlist
\item
  Choose a sentence of length \(5\) to \(8\).
\item
  Assign simple hand-designed token features.
\item
  Construct at least three different heads by choosing different
  projection matrices.
\item
  Plot attention heatmaps for each head.
\item
  Compare unmasked and causal attention.
\item
  Compute row entropies.
\item
  Write a short report explaining what each head appears to measure.
\end{enumerate}

The goal is not to train a model. The goal is to understand how
different query-key geometries produce different attention patterns.

\section{14.22 Python lab connection}\label{python-lab-connection-7}

The associated lab for this chapter is planned as:

\begin{Shaded}
\begin{Highlighting}[]
\NormalTok{labs/chapter{-}14{-}self{-}attention.ipynb}
\end{Highlighting}
\end{Shaded}

The lab should be written for independent study and should include:

\begin{itemize}
\tightlist
\item
  a review of matrix shapes;
\item
  one-head self-attention from scratch;
\item
  padding and causal masks;
\item
  multi-head self-attention;
\item
  attention heatmap visualization;
\item
  attention entropy diagnostics;
\item
  a short comparison with RNNs and n-gram models.
\end{itemize}

\section{14.23 What to remember}\label{what-to-remember-4}

\begin{itemize}
\tightlist
\item
  Self-attention computes \(Q,K,V\) from the same sequence matrix \(X\).
\item
  Each output token representation is a weighted average of value
  vectors from the sequence.
\item
  Attention weights form a row-stochastic matrix.
\item
  Self-attention produces contextual embeddings.
\item
  Without positional information and masks, self-attention is
  permutation equivariant.
\item
  Causal masks enforce the conditional structure of autoregressive
  language modeling.
\item
  Multi-head attention learns several comparison geometries in parallel.
\item
  Full self-attention has quadratic cost in sequence length.
\end{itemize}

\section{Notes and references}\label{notes-and-references-8}

Self-attention is the central operation in the Transformer architecture.
This chapter connects the attention formula from Chapter 13 to
contextual token representations, symmetry, masking, and multi-head
design. Useful background sources include Vaswani et al. (2017),
Jurafsky and Martin (2024), and Goodfellow et al. (2016b).

\chapter{The Transformer
Architecture}\label{the-transformer-architecture}

Residual nonlinear operators on token sequences

\hfill\break

The Transformer combines the self-attention mechanism from Chapter 14
with residual connections, normalization, and position-wise neural
networks. From a mathematical point of view, a Transformer layer is a
learned nonlinear map

\[
F:R^{T\times d}\to R^{T\times d}
\]

that transforms a matrix of token representations into a new matrix of
contextual token representations. The key architectural idea is not a
single formula, but a carefully organized composition of simple maps:

\[
\text{linear maps}\;\longrightarrow\;\text{softmax attention}\;\longrightarrow\;\text{residual addition}\;\longrightarrow\;\text{normalization}\;\longrightarrow\;\text{position-wise nonlinear maps}.
\]

This chapter studies the Transformer as an operator on sequence
matrices. We emphasize dimensions, probability models, optimization
objectives, and why the architecture is easier to train than a plain
deep stack of nonlinear maps.

\section{Learning goals}\label{learning-goals-14}

After studying this chapter, you should be able to:

\begin{enumerate}
\def\labelenumi{\arabic{enumi}.}
\tightlist
\item
  describe a Transformer layer as a map from \(R^{T\times d}\) to
  \(R^{T\times d}\);
\item
  distinguish encoder blocks, decoder blocks, and encoder-decoder
  blocks;
\item
  write the formulas for multi-head attention, residual connections,
  layer normalization, and feedforward sublayers;
\item
  explain why residual connections help gradient flow;
\item
  explain layer normalization as a row-wise standardization map;
\item
  compute the parameter count of a basic Transformer block;
\item
  connect decoder-only Transformers with autoregressive language
  modeling;
\item
  implement a small Transformer block from scratch in NumPy;
\item
  use a Plotly.js interactive to visualize how attention, residual
  mixing, and nonlinear feedforward maps change token representations;
\item
  describe the main mathematical difference between RNNs and
  Transformers.
\end{enumerate}

\section{15.1 The sequence-matrix
viewpoint}\label{the-sequence-matrix-viewpoint}

Let a token sequence have length \(T\). After token embedding and
positional encoding, the input to a Transformer layer is a matrix

\[
X\in R^{T\times d},
\]

where row \(x_i^T\in R^d\) represents token position \(i\). The
Transformer layer produces another matrix

\[
Y=F(X)\in R^{T\times d}.
\]

The row \(y_i^T\) is a contextual representation of token position
\(i\). The same layer is applied to every example in a batch, but for
mathematical clarity we first ignore the batch dimension.

A simple Transformer encoder layer has the form

\[
\begin{aligned}
H &= \operatorname{LayerNorm}\bigl(X + \operatorname{MHA}(X)\bigr),\\
Y &= \operatorname{LayerNorm}\bigl(H + \operatorname{FFN}(H)\bigr).
\end{aligned}
\]

This is called a \textbf{post-norm} block because layer normalization is
applied after the residual addition. Many modern implementations use a
\textbf{pre-norm} block:

\[
\begin{aligned}
H &= X + \operatorname{MHA}(\operatorname{LayerNorm}(X)),\\
Y &= H + \operatorname{FFN}(\operatorname{LayerNorm}(H)).
\end{aligned}
\]

The two designs use the same basic pieces, but place normalization in
different locations. Pre-norm blocks are often easier to train in very
deep networks because the residual path is more direct.

\textbf{Main mathematical abstraction.} A Transformer is a deep
composition

\[
F_L\circ F_{L-1}\circ\cdots\circ F_1
\]

of sequence-to-sequence maps, where each \(F_\ell\) mixes token
positions by attention and transforms feature coordinates by feedforward
nonlinear maps.

\section{15.2 Multi-head attention inside a Transformer
block}\label{multi-head-attention-inside-a-transformer-block}

A Transformer does not usually use a single attention head. It uses
\(h\) heads in parallel. Let

\[
X\in R^{T\times d}.
\]

For head \(r=1,\ldots,h\), define

\[
Q_r=XW_Q^{(r)},\qquad
K_r=XW_K^{(r)},\qquad
V_r=XW_V^{(r)},
\]

where

\[
W_Q^{(r)},W_K^{(r)}\in R^{d\times d_k},
\qquad
W_V^{(r)}\in R^{d\times d_v}.
\]

The \(r\)th head output is

\[
Z_r=\operatorname{softmax}_{\text{row}}\left(\frac{Q_rK_r^T+M}{\sqrt{d_k}}\right)V_r
\in R^{T\times d_v}.
\]

The matrix \(M\in R^{T\times T}\) is optional. For encoder
self-attention, \(M\) may be a padding mask. For decoder self-attention,
\(M\) also includes a causal mask.

The head outputs are concatenated:

\[
Z=\operatorname{concat}(Z_1,\ldots,Z_h)\in R^{T\times hd_v}.
\]

Then a final output projection gives

\[
\operatorname{MHA}(X)=ZW_O,
\qquad
W_O\in R^{hd_v\times d}.
\]

A common choice is

\[
d_k=d_v=d/h,
\]

so \(hd_v=d\) and the output has the same feature dimension as \(X\).

\subsection{Why multiple heads?}\label{why-multiple-heads}

A single attention head has one learned geometry for comparing tokens.
Multiple heads allow several comparison geometries in parallel. One head
might emphasize local syntactic relations, another might emphasize
long-range dependencies, and another might attend to special tokens or
punctuation. This interpretation should be used carefully: attention
heads are learned numerical objects, not guaranteed linguistic modules.

Mathematically, multi-head attention increases expressiveness by
applying several different bilinear score functions

\[
(x_i,x_j)\mapsto x_i^TW_Q^{(r)}(W_K^{(r)})^Tx_j
\]

and several different value transformations

\[
x_j\mapsto x_j^TW_V^{(r)}.
\]

Thus each head defines a different kernel-like similarity and a
different message representation.

\section{15.3 Residual connections}\label{residual-connections}

A residual connection adds the input of a sublayer to its output. If
\(G:R^{T\times d}\to R^{T\times d}\) is a sublayer, the residual map is

\[
R(X)=X+G(X).
\]

The derivative of this map has the form

\[
DR(X)=I+DG(X),
\]

where \(I\) denotes the identity map on \(R^{T\times d}\). This is the
main mathematical reason residual connections help optimization: even if
\(DG(X)\) has small singular values, the identity term provides a direct
path for gradients.

To see the idea in a scalar example, suppose

\[
y=x+g(x).
\]

Then

\[
\frac{dy}{dx}=1+g'(x).
\]

In a deep non-residual composition, derivatives multiply many factors.
In a residual network, each layer has an identity contribution, so
information and gradients can pass through the network more easily.

\textbf{Residual learning.} A residual block does not need to learn the
entire transformation. It only needs to learn a correction to the
identity map.

\section{15.4 Layer normalization}\label{layer-normalization}

Layer normalization standardizes each token vector across its feature
coordinates. For a row vector

\[
x=(x_1,\ldots,x_d)\in R^d,
\]

define

\[
\mu(x)=\frac{1}{d}\sum_{j=1}^d x_j,
\qquad
\sigma^2(x)=\frac{1}{d}\sum_{j=1}^d (x_j-\mu(x))^2.
\]

Layer normalization maps \(x\) to

\[
\operatorname{LN}(x)_j
=
\gamma_j\frac{x_j-\mu(x)}{\sqrt{\sigma^2(x)+\varepsilon}}+\beta_j,
\qquad j=1,\ldots,d,
\]

where \(\gamma,\beta\in R^d\) are trainable parameters and
\(\varepsilon>0\) is a small constant for numerical stability.

For a sequence matrix \(X\in R^{T\times d}\), layer normalization is
applied row by row:

\[
\operatorname{LN}(X)_{i,:}=\operatorname{LN}(X_{i,:}).
\]

Layer normalization is different from batch normalization. It does not
compute statistics across examples in a batch. Instead, it computes
statistics across the coordinates of one token representation.

\subsection{Geometric interpretation}\label{geometric-interpretation}

Ignoring \(\gamma\) and \(\beta\), layer normalization subtracts the
component in the direction of the all-ones vector and rescales to
approximately fixed length. If

\[
\mathbf{1}=(1,\ldots,1)^T\in R^d,
\]

then centering replaces \(x\) by

\[
x-\mu(x)\mathbf{1}.
\]

This vector lies in the hyperplane

\[
\{z\in R^d: z^T\mathbf{1}=0\}.
\]

Thus layer normalization controls scale and location of token features,
making the optimization problem more stable.

\section{15.5 Position-wise feedforward
networks}\label{position-wise-feedforward-networks}

After attention mixes information across token positions, a Transformer
block applies the same feedforward network to each token position
independently. For one row \(h_i\in R^d\), define

\[
\operatorname{FFN}(h_i)=\sigma(h_iW_1+b_1)W_2+b_2,
\]

where

\[
W_1\in R^{d\times d_{\mathrm{ff}}},
\qquad
W_2\in R^{d_{\mathrm{ff}}\times d}.
\]

The activation \(\sigma\) is applied coordinatewise. Common choices
include ReLU and GELU. In matrix form,

\[
\operatorname{FFN}(H)=\sigma(HW_1+\mathbf{1}b_1^T)W_2+\mathbf{1}b_2^T.
\]

This sublayer does not mix token positions. It transforms the feature
representation at each position. Attention mixes positions; the
feedforward network mixes coordinates.

\textbf{Two kinds of mixing.} In a Transformer block, attention mixes
rows of \(X\), while the feedforward network mixes columns of \(X\).

\section{15.6 Encoder blocks}\label{encoder-blocks}

A Transformer encoder is designed to read an entire input sequence and
produce contextual representations for all positions. An encoder block
uses bidirectional self-attention: every non-padding token may attend to
every other non-padding token.

For an encoder block, a typical post-norm formula is

\[
\begin{aligned}
H &= \operatorname{LN}(X+\operatorname{MHA}_{\text{enc}}(X)),\\
Y &= \operatorname{LN}(H+\operatorname{FFN}(H)).
\end{aligned}
\]

The attention mask usually prevents real tokens from attending to
padding positions. It does not prevent a token from attending to future
positions because the encoder is not generating left-to-right.

Encoder-only architectures are natural for:

\begin{itemize}
\tightlist
\item
  document classification;
\item
  token classification;
\item
  sentence similarity;
\item
  extractive question answering;
\item
  sequence representation.
\end{itemize}

Mathematically, an encoder is a map

\[
F_{\text{enc}}:R^{T\times d}\to R^{T\times d}
\]

that is allowed to use information from the full input sequence at every
output position.

\section{15.7 Decoder blocks}\label{decoder-blocks}

A Transformer decoder is designed for generation. In a decoder-only
language model, token position \(t\) should not use information from
positions \(t+1,t+2,\ldots,T\) when predicting the next token. This is
enforced by a causal mask.

For decoder self-attention, the mask satisfies

\[
M_{ij}=\begin{cases}
0, & j\le i,\\
-\infty, & j>i.
\end{cases}
\]

Then after applying row-wise softmax, all future-position weights are
zero:

\[
a_{ij}=0\qquad\text{for }j>i.
\]

A decoder-only block has the same two main sublayers as an encoder
block, but with causal self-attention:

\[
\begin{aligned}
H &= \operatorname{LN}(X+\operatorname{MHA}_{\text{causal}}(X)),\\
Y &= \operatorname{LN}(H+\operatorname{FFN}(H)).
\end{aligned}
\]

Decoder-only architectures are natural for autoregressive language
modeling:

\[
P(w_1,\ldots,w_T)=\prod_{t=1}^T P(w_t\mid w_1,\ldots,w_{t-1}).
\]

The model predicts a probability distribution for the next token at each
position.

\section{15.8 Encoder-decoder
Transformers}\label{encoder-decoder-transformers}

The original Transformer architecture for machine translation used both
an encoder and a decoder. The encoder reads the source sequence. The
decoder generates the target sequence. The decoder uses two kinds of
attention:

\begin{enumerate}
\def\labelenumi{\arabic{enumi}.}
\tightlist
\item
  \textbf{causal self-attention} over previously generated target
  tokens;
\item
  \textbf{cross-attention} from target positions to source positions.
\end{enumerate}

Let

\[
X_{\text{src}}\in R^{T_s\times d},
\qquad
X_{\text{tgt}}\in R^{T_t\times d}.
\]

The encoder produces

\[
E=F_{\text{enc}}(X_{\text{src}})\in R^{T_s\times d}.
\]

In cross-attention, queries come from the decoder state, while keys and
values come from the encoder output:

\[
Q=HW_Q,
\qquad
K=EW_K,
\qquad
V=EW_V.
\]

The cross-attention output is

\[
\operatorname{CrossAttn}(H,E)
=
\operatorname{softmax}_{\text{row}}\left(\frac{QK^T}{\sqrt{d_k}}\right)V.
\]

This lets each target position decide which source positions are
relevant.

\section{15.9 Output layer and language
modeling}\label{output-layer-and-language-modeling}

For a decoder-only language model, the final Transformer representation
at position \(t\) is a vector

\[
h_t\in R^d.
\]

A linear output layer maps \(h_t\) to vocabulary logits:

\[
z_t=h_tW_U+b_U\in R^{|V|},
\]

where

\[
W_U\in R^{d\times |V|},
\qquad
b_U\in R^{|V|}.
\]

The conditional distribution over the next token is

\[
P_\theta(w_{t+1}=v\mid w_1,\ldots,w_t)
=\frac{\exp(z_{t,v})}{\sum_{u\in V}\exp(z_{t,u})}.
\]

For a training sequence \(w_1,\ldots,w_T\), the negative log-likelihood
is

\[
\mathcal{L}(\theta)
= -\sum_{t=1}^{T-1}\log P_\theta(w_{t+1}\mid w_1,\ldots,w_t).
\]

Dividing by \(T-1\) gives the average cross-entropy loss. Exponentiating
average cross-entropy gives perplexity.

\section{15.10 Parameter count of a basic
block}\label{parameter-count-of-a-basic-block}

Assume a Transformer block has model dimension \(d\), number of heads
\(h\), and feedforward dimension \(d_{\mathrm{ff}}\). Suppose each head
uses

\[
d_k=d_v=d/h.
\]

The query, key, and value projections across all heads together have
approximately

\[
3d^2
\]

weights. The output projection has

\[
d^2
\]

weights. Therefore multi-head attention has about

\[
4d^2
\]

weights, ignoring biases.

The feedforward sublayer has

\[
dd_{\mathrm{ff}}+d_{\mathrm{ff}}d=2dd_{\mathrm{ff}}
\]

weights, ignoring biases. Layer normalization contributes \(2d\)
parameters per normalization layer: scale \(\gamma\) and shift
\(\beta\).

Thus a rough block parameter count is

\[
4d^2+2dd_{\mathrm{ff}}+4d,
\]

for a block with two layer normalization modules and no bias counting in
the linear layers.

A common choice is

\[
d_{\mathrm{ff}}=4d.
\]

Then the feedforward sublayer has about \(8d^2\) weights, which is often
more than the attention sublayer.

\section{15.11 Computational
complexity}\label{computational-complexity-1}

Self-attention forms the score matrix

\[
S=QK^T\in R^{T\times T}.
\]

Computing \(S\) costs on the order of

\[
O(T^2d)
\]

operations. Multiplying the attention matrix by the values also costs
about

\[
O(T^2d).
\]

The feedforward sublayer costs

\[
O(Tdd_{\mathrm{ff}}).
\]

Thus the cost of attention grows quadratically with sequence length
\(T\). This is the main computational bottleneck for long-context
Transformers.

RNNs process sequences sequentially with cost roughly linear in \(T\),
but they are hard to parallelize over time. Transformers use \(T^2\)
pairwise comparisons, but these comparisons can be computed in parallel
on modern hardware.

\textbf{Tradeoff.} RNNs have sequential recurrence. Transformers replace
recurrence with parallel pairwise interaction.

\section{15.12 Interactive: Transformer block geometry
explorer}\label{interactive-transformer-block-geometry-explorer}

The following interactive uses a two-dimensional toy sequence. It is not
a real Transformer model, but it visualizes three structural ideas:

\begin{itemize}
\tightlist
\item
  attention mixes token positions;
\item
  residual connections pull the transformed representation back toward
  the input;
\item
  the feedforward sublayer bends the representation coordinatewise.
\end{itemize}

Try increasing the attention mixing parameter and residual strength.
Notice that a Transformer layer is not simply a classifier or a lookup
table; it is a learned geometric transformation of a whole sequence
matrix.

\protect\phantomsection\label{ch15-transformer-status}
Loading local Plotly for Chapter 15\ldots{}

Attention mixing \(\alpha\) Residual strength \(\rho\) FFN nonlinearity
\(\eta\)

\protect\phantomsection\label{ch15-transformer-demo}

\protect\phantomsection\label{ch15-transformer-readout}

\section{15.13 Python example: layer
normalization}\label{python-example-layer-normalization}

The following implementation applies layer normalization row by row to a
token matrix.

\protect\phantomsection\label{py-ch15-layernorm}
\begin{Shaded}
\begin{Highlighting}[]
\ImportTok{import}\NormalTok{ numpy }\ImportTok{as}\NormalTok{ np}

\NormalTok{np.set\_printoptions(precision}\OperatorTok{=}\DecValTok{4}\NormalTok{, suppress}\OperatorTok{=}\VariableTok{True}\NormalTok{)}

\KeywordTok{def}\NormalTok{ layer\_norm(X, gamma}\OperatorTok{=}\VariableTok{None}\NormalTok{, beta}\OperatorTok{=}\VariableTok{None}\NormalTok{, eps}\OperatorTok{=}\FloatTok{1e{-}5}\NormalTok{):}
    \CommentTok{"""Layer normalization for X with shape (T, d)."""}
\NormalTok{    X }\OperatorTok{=}\NormalTok{ np.asarray(X, dtype}\OperatorTok{=}\BuiltInTok{float}\NormalTok{)}
\NormalTok{    mu }\OperatorTok{=}\NormalTok{ X.mean(axis}\OperatorTok{=}\DecValTok{1}\NormalTok{, keepdims}\OperatorTok{=}\VariableTok{True}\NormalTok{)}
\NormalTok{    var }\OperatorTok{=}\NormalTok{ ((X }\OperatorTok{{-}}\NormalTok{ mu) }\OperatorTok{**} \DecValTok{2}\NormalTok{).mean(axis}\OperatorTok{=}\DecValTok{1}\NormalTok{, keepdims}\OperatorTok{=}\VariableTok{True}\NormalTok{)}
\NormalTok{    Z }\OperatorTok{=}\NormalTok{ (X }\OperatorTok{{-}}\NormalTok{ mu) }\OperatorTok{/}\NormalTok{ np.sqrt(var }\OperatorTok{+}\NormalTok{ eps)}
    \ControlFlowTok{if}\NormalTok{ gamma }\KeywordTok{is} \VariableTok{None}\NormalTok{:}
\NormalTok{        gamma }\OperatorTok{=}\NormalTok{ np.ones(X.shape[}\DecValTok{1}\NormalTok{])}
    \ControlFlowTok{if}\NormalTok{ beta }\KeywordTok{is} \VariableTok{None}\NormalTok{:}
\NormalTok{        beta }\OperatorTok{=}\NormalTok{ np.zeros(X.shape[}\DecValTok{1}\NormalTok{])}
    \ControlFlowTok{return}\NormalTok{ Z }\OperatorTok{*}\NormalTok{ gamma }\OperatorTok{+}\NormalTok{ beta}

\NormalTok{X }\OperatorTok{=}\NormalTok{ np.array([}
\NormalTok{    [}\FloatTok{2.0}\NormalTok{, }\FloatTok{4.0}\NormalTok{, }\FloatTok{6.0}\NormalTok{, }\FloatTok{8.0}\NormalTok{],}
\NormalTok{    [}\FloatTok{1.0}\NormalTok{, }\FloatTok{1.0}\NormalTok{, }\FloatTok{3.0}\NormalTok{, }\FloatTok{7.0}\NormalTok{],}
\NormalTok{    [}\OperatorTok{{-}}\FloatTok{2.0}\NormalTok{, }\FloatTok{0.0}\NormalTok{, }\FloatTok{2.0}\NormalTok{, }\FloatTok{4.0}\NormalTok{]}
\NormalTok{])}

\NormalTok{Y }\OperatorTok{=}\NormalTok{ layer\_norm(X)}
\BuiltInTok{print}\NormalTok{(Y)}
\BuiltInTok{print}\NormalTok{(}\StringTok{"row means after LN:"}\NormalTok{, Y.mean(axis}\OperatorTok{=}\DecValTok{1}\NormalTok{))}
\BuiltInTok{print}\NormalTok{(}\StringTok{"row variances after LN:"}\NormalTok{, Y.var(axis}\OperatorTok{=}\DecValTok{1}\NormalTok{))}
\end{Highlighting}
\end{Shaded}

\begin{verbatim}
[[-1.3416 -0.4472  0.4472  1.3416]
 [-0.8165 -0.8165  0.      1.633 ]
 [-1.3416 -0.4472  0.4472  1.3416]]
row means after LN: [0. 0. 0.]
row variances after LN: [1. 1. 1.]
\end{verbatim}

Layer normalization changes the feature coordinates of each token
vector, but it does not mix token positions.

\section{15.14 Python example: multi-head attention from
scratch}\label{python-example-multi-head-attention-from-scratch}

The next example implements a small multi-head self-attention layer. The
purpose is dimension bookkeeping, not performance.

\protect\phantomsection\label{py-ch15-mha-from-scratch}
\begin{Shaded}
\begin{Highlighting}[]
\ImportTok{import}\NormalTok{ numpy }\ImportTok{as}\NormalTok{ np}

\NormalTok{rng }\OperatorTok{=}\NormalTok{ np.random.default\_rng(}\DecValTok{7243}\NormalTok{)}

\NormalTok{T }\OperatorTok{=} \DecValTok{5}
\NormalTok{model\_dim }\OperatorTok{=} \DecValTok{8}
\NormalTok{num\_heads }\OperatorTok{=} \DecValTok{2}
\NormalTok{head\_dim }\OperatorTok{=}\NormalTok{ model\_dim }\OperatorTok{//}\NormalTok{ num\_heads}

\NormalTok{X }\OperatorTok{=}\NormalTok{ rng.normal(size}\OperatorTok{=}\NormalTok{(T, model\_dim))}

\NormalTok{WQ }\OperatorTok{=}\NormalTok{ rng.normal(scale}\OperatorTok{=}\FloatTok{0.2}\NormalTok{, size}\OperatorTok{=}\NormalTok{(num\_heads, model\_dim, head\_dim))}
\NormalTok{WK }\OperatorTok{=}\NormalTok{ rng.normal(scale}\OperatorTok{=}\FloatTok{0.2}\NormalTok{, size}\OperatorTok{=}\NormalTok{(num\_heads, model\_dim, head\_dim))}
\NormalTok{WV }\OperatorTok{=}\NormalTok{ rng.normal(scale}\OperatorTok{=}\FloatTok{0.2}\NormalTok{, size}\OperatorTok{=}\NormalTok{(num\_heads, model\_dim, head\_dim))}
\NormalTok{WO }\OperatorTok{=}\NormalTok{ rng.normal(scale}\OperatorTok{=}\FloatTok{0.2}\NormalTok{, size}\OperatorTok{=}\NormalTok{(num\_heads }\OperatorTok{*}\NormalTok{ head\_dim, model\_dim))}

\KeywordTok{def}\NormalTok{ softmax\_rows(S):}
\NormalTok{    S\_shift }\OperatorTok{=}\NormalTok{ S }\OperatorTok{{-}}\NormalTok{ S.}\BuiltInTok{max}\NormalTok{(axis}\OperatorTok{=}\DecValTok{1}\NormalTok{, keepdims}\OperatorTok{=}\VariableTok{True}\NormalTok{)}
\NormalTok{    E }\OperatorTok{=}\NormalTok{ np.exp(S\_shift)}
    \ControlFlowTok{return}\NormalTok{ E }\OperatorTok{/}\NormalTok{ E.}\BuiltInTok{sum}\NormalTok{(axis}\OperatorTok{=}\DecValTok{1}\NormalTok{, keepdims}\OperatorTok{=}\VariableTok{True}\NormalTok{)}

\KeywordTok{def}\NormalTok{ multi\_head\_attention(X, WQ, WK, WV, WO, mask}\OperatorTok{=}\VariableTok{None}\NormalTok{):}
\NormalTok{    head\_outputs }\OperatorTok{=}\NormalTok{ []}
\NormalTok{    attention\_matrices }\OperatorTok{=}\NormalTok{ []}
    \ControlFlowTok{for}\NormalTok{ r }\KeywordTok{in} \BuiltInTok{range}\NormalTok{(WQ.shape[}\DecValTok{0}\NormalTok{]):}
\NormalTok{        Q }\OperatorTok{=}\NormalTok{ X }\OperatorTok{@}\NormalTok{ WQ[r]}
\NormalTok{        K }\OperatorTok{=}\NormalTok{ X }\OperatorTok{@}\NormalTok{ WK[r]}
\NormalTok{        V }\OperatorTok{=}\NormalTok{ X }\OperatorTok{@}\NormalTok{ WV[r]}
\NormalTok{        scores }\OperatorTok{=}\NormalTok{ (Q }\OperatorTok{@}\NormalTok{ K.T) }\OperatorTok{/}\NormalTok{ np.sqrt(Q.shape[}\DecValTok{1}\NormalTok{])}
        \ControlFlowTok{if}\NormalTok{ mask }\KeywordTok{is} \KeywordTok{not} \VariableTok{None}\NormalTok{:}
\NormalTok{            scores }\OperatorTok{=}\NormalTok{ scores }\OperatorTok{+}\NormalTok{ mask}
\NormalTok{        A }\OperatorTok{=}\NormalTok{ softmax\_rows(scores)}
\NormalTok{        Z }\OperatorTok{=}\NormalTok{ A }\OperatorTok{@}\NormalTok{ V}
\NormalTok{        head\_outputs.append(Z)}
\NormalTok{        attention\_matrices.append(A)}
\NormalTok{    Z\_all }\OperatorTok{=}\NormalTok{ np.concatenate(head\_outputs, axis}\OperatorTok{=}\DecValTok{1}\NormalTok{)}
    \ControlFlowTok{return}\NormalTok{ Z\_all }\OperatorTok{@}\NormalTok{ WO, attention\_matrices}

\NormalTok{Y, As }\OperatorTok{=}\NormalTok{ multi\_head\_attention(X, WQ, WK, WV, WO)}

\BuiltInTok{print}\NormalTok{(}\StringTok{"X shape:"}\NormalTok{, X.shape)}
\BuiltInTok{print}\NormalTok{(}\StringTok{"head attention shape:"}\NormalTok{, As[}\DecValTok{0}\NormalTok{].shape)}
\BuiltInTok{print}\NormalTok{(}\StringTok{"output shape:"}\NormalTok{, Y.shape)}
\BuiltInTok{print}\NormalTok{(}\StringTok{"row sums of first attention matrix:"}\NormalTok{, As[}\DecValTok{0}\NormalTok{].}\BuiltInTok{sum}\NormalTok{(axis}\OperatorTok{=}\DecValTok{1}\NormalTok{))}
\end{Highlighting}
\end{Shaded}

\begin{verbatim}
X shape: (5, 8)
head attention shape: (5, 5)
output shape: (5, 8)
row sums of first attention matrix: [1. 1. 1. 1. 1.]
\end{verbatim}

The row sums are approximately one because each attention matrix is
row-stochastic.

\section{15.15 Python example: a full Transformer encoder
block}\label{python-example-a-full-transformer-encoder-block}

Now we add residual connections, layer normalization, and a
position-wise feedforward network.

\protect\phantomsection\label{py-ch15-transformer-block}
\begin{Shaded}
\begin{Highlighting}[]
\ImportTok{import}\NormalTok{ numpy }\ImportTok{as}\NormalTok{ np}

\NormalTok{rng }\OperatorTok{=}\NormalTok{ np.random.default\_rng(}\DecValTok{2026}\NormalTok{)}

\NormalTok{T }\OperatorTok{=} \DecValTok{6}
\NormalTok{d }\OperatorTok{=} \DecValTok{8}
\NormalTok{h }\OperatorTok{=} \DecValTok{2}
\NormalTok{dh }\OperatorTok{=}\NormalTok{ d }\OperatorTok{//}\NormalTok{ h}
\NormalTok{dff }\OperatorTok{=} \DecValTok{4} \OperatorTok{*}\NormalTok{ d}

\NormalTok{X }\OperatorTok{=}\NormalTok{ rng.normal(size}\OperatorTok{=}\NormalTok{(T, d))}

\NormalTok{params }\OperatorTok{=}\NormalTok{ \{}
    \StringTok{"WQ"}\NormalTok{: rng.normal(scale}\OperatorTok{=}\FloatTok{0.2}\NormalTok{, size}\OperatorTok{=}\NormalTok{(h, d, dh)),}
    \StringTok{"WK"}\NormalTok{: rng.normal(scale}\OperatorTok{=}\FloatTok{0.2}\NormalTok{, size}\OperatorTok{=}\NormalTok{(h, d, dh)),}
    \StringTok{"WV"}\NormalTok{: rng.normal(scale}\OperatorTok{=}\FloatTok{0.2}\NormalTok{, size}\OperatorTok{=}\NormalTok{(h, d, dh)),}
    \StringTok{"WO"}\NormalTok{: rng.normal(scale}\OperatorTok{=}\FloatTok{0.2}\NormalTok{, size}\OperatorTok{=}\NormalTok{(h }\OperatorTok{*}\NormalTok{ dh, d)),}
    \StringTok{"W1"}\NormalTok{: rng.normal(scale}\OperatorTok{=}\FloatTok{0.2}\NormalTok{, size}\OperatorTok{=}\NormalTok{(d, dff)),}
    \StringTok{"b1"}\NormalTok{: np.zeros(dff),}
    \StringTok{"W2"}\NormalTok{: rng.normal(scale}\OperatorTok{=}\FloatTok{0.2}\NormalTok{, size}\OperatorTok{=}\NormalTok{(dff, d)),}
    \StringTok{"b2"}\NormalTok{: np.zeros(d),}
\NormalTok{\}}

\KeywordTok{def}\NormalTok{ relu(U):}
    \ControlFlowTok{return}\NormalTok{ np.maximum(U, }\FloatTok{0.0}\NormalTok{)}

\KeywordTok{def}\NormalTok{ ffn(H, params):}
    \ControlFlowTok{return}\NormalTok{ relu(H }\OperatorTok{@}\NormalTok{ params[}\StringTok{"W1"}\NormalTok{] }\OperatorTok{+}\NormalTok{ params[}\StringTok{"b1"}\NormalTok{]) }\OperatorTok{@}\NormalTok{ params[}\StringTok{"W2"}\NormalTok{] }\OperatorTok{+}\NormalTok{ params[}\StringTok{"b2"}\NormalTok{]}

\KeywordTok{def}\NormalTok{ transformer\_encoder\_block(X, params):}
\NormalTok{    attn\_out, attention\_mats }\OperatorTok{=}\NormalTok{ multi\_head\_attention(}
\NormalTok{        X,}
\NormalTok{        params[}\StringTok{"WQ"}\NormalTok{],}
\NormalTok{        params[}\StringTok{"WK"}\NormalTok{],}
\NormalTok{        params[}\StringTok{"WV"}\NormalTok{],}
\NormalTok{        params[}\StringTok{"WO"}\NormalTok{]}
\NormalTok{    )}
\NormalTok{    H }\OperatorTok{=}\NormalTok{ layer\_norm(X }\OperatorTok{+}\NormalTok{ attn\_out)}
\NormalTok{    Y }\OperatorTok{=}\NormalTok{ layer\_norm(H }\OperatorTok{+}\NormalTok{ ffn(H, params))}
    \ControlFlowTok{return}\NormalTok{ Y, attention\_mats}

\NormalTok{Y, attention\_mats }\OperatorTok{=}\NormalTok{ transformer\_encoder\_block(X, params)}

\BuiltInTok{print}\NormalTok{(}\StringTok{"input shape:"}\NormalTok{, X.shape)}
\BuiltInTok{print}\NormalTok{(}\StringTok{"output shape:"}\NormalTok{, Y.shape)}
\BuiltInTok{print}\NormalTok{(}\StringTok{"first token before:"}\NormalTok{, X[}\DecValTok{0}\NormalTok{])}
\BuiltInTok{print}\NormalTok{(}\StringTok{"first token after:"}\NormalTok{, Y[}\DecValTok{0}\NormalTok{])}
\end{Highlighting}
\end{Shaded}

\begin{verbatim}
input shape: (6, 8)
output shape: (6, 8)
first token before: [-0.7931  0.2406 -1.8963  1.3958  0.6383 -0.292  -0.3119  0.3038]
first token after: [-0.757   0.0496 -1.8095  0.6765  1.1463 -0.6849 -0.0028  1.3818]
\end{verbatim}

The block maps \(R^{T\times d}\) to \(R^{T\times d}\). The sequence
length and model dimension are preserved.

\section{15.16 Python example: causal decoder
mask}\label{python-example-causal-decoder-mask}

A decoder-only Transformer uses a causal mask so that position \(i\)
cannot attend to position \(j>i\).

\protect\phantomsection\label{py-ch15-causal-mask}
\begin{Shaded}
\begin{Highlighting}[]
\ImportTok{import}\NormalTok{ numpy }\ImportTok{as}\NormalTok{ np}

\NormalTok{T }\OperatorTok{=} \DecValTok{6}
\NormalTok{causal\_mask }\OperatorTok{=}\NormalTok{ np.triu(np.ones((T, T)), k}\OperatorTok{=}\DecValTok{1}\NormalTok{)}
\NormalTok{causal\_mask }\OperatorTok{=}\NormalTok{ np.where(causal\_mask }\OperatorTok{==} \DecValTok{1}\NormalTok{, }\OperatorTok{{-}}\FloatTok{1e9}\NormalTok{, }\FloatTok{0.0}\NormalTok{)}

\BuiltInTok{print}\NormalTok{(causal\_mask)}

\NormalTok{Y\_causal, causal\_As }\OperatorTok{=}\NormalTok{ multi\_head\_attention(}
\NormalTok{    X[:T],}
\NormalTok{    params[}\StringTok{"WQ"}\NormalTok{],}
\NormalTok{    params[}\StringTok{"WK"}\NormalTok{],}
\NormalTok{    params[}\StringTok{"WV"}\NormalTok{],}
\NormalTok{    params[}\StringTok{"WO"}\NormalTok{],}
\NormalTok{    mask}\OperatorTok{=}\NormalTok{causal\_mask}
\NormalTok{)}

\BuiltInTok{print}\NormalTok{(}\StringTok{"First head causal attention matrix:"}\NormalTok{)}
\BuiltInTok{print}\NormalTok{(causal\_As[}\DecValTok{0}\NormalTok{])}
\BuiltInTok{print}\NormalTok{(}\StringTok{"largest future attention weight:"}\NormalTok{, np.triu(causal\_As[}\DecValTok{0}\NormalTok{], k}\OperatorTok{=}\DecValTok{1}\NormalTok{).}\BuiltInTok{max}\NormalTok{())}
\end{Highlighting}
\end{Shaded}

\begin{verbatim}
[[ 0.e+00 -1.e+09 -1.e+09 -1.e+09 -1.e+09 -1.e+09]
 [ 0.e+00  0.e+00 -1.e+09 -1.e+09 -1.e+09 -1.e+09]
 [ 0.e+00  0.e+00  0.e+00 -1.e+09 -1.e+09 -1.e+09]
 [ 0.e+00  0.e+00  0.e+00  0.e+00 -1.e+09 -1.e+09]
 [ 0.e+00  0.e+00  0.e+00  0.e+00  0.e+00 -1.e+09]
 [ 0.e+00  0.e+00  0.e+00  0.e+00  0.e+00  0.e+00]]
First head causal attention matrix:
[[1.     0.     0.     0.     0.     0.    ]
 [0.5115 0.4885 0.     0.     0.     0.    ]
 [0.3315 0.3412 0.3272 0.     0.     0.    ]
 [0.3058 0.2589 0.2758 0.1595 0.     0.    ]
 [0.1792 0.1936 0.1667 0.2458 0.2147 0.    ]
 [0.1508 0.1865 0.1687 0.1265 0.1511 0.2164]]
largest future attention weight: 0.0
\end{verbatim}

The largest future attention weight should be numerically close to zero.

\section{15.17 Python example: output probabilities for language
modeling}\label{python-example-output-probabilities-for-language-modeling}

The final hidden state at each position is converted into vocabulary
logits and then probabilities.

\protect\phantomsection\label{py-ch15-output-probabilities}
\begin{Shaded}
\begin{Highlighting}[]
\ImportTok{import}\NormalTok{ numpy }\ImportTok{as}\NormalTok{ np}

\NormalTok{rng }\OperatorTok{=}\NormalTok{ np.random.default\_rng(}\DecValTok{15}\NormalTok{)}

\NormalTok{vocab }\OperatorTok{=}\NormalTok{ [}\StringTok{"\textless{}bos\textgreater{}"}\NormalTok{, }\StringTok{"the"}\NormalTok{, }\StringTok{"model"}\NormalTok{, }\StringTok{"learns"}\NormalTok{, }\StringTok{"language"}\NormalTok{, }\StringTok{"."}\NormalTok{]}
\NormalTok{V }\OperatorTok{=} \BuiltInTok{len}\NormalTok{(vocab)}
\NormalTok{T }\OperatorTok{=}\NormalTok{ Y\_causal.shape[}\DecValTok{0}\NormalTok{]}
\NormalTok{d }\OperatorTok{=}\NormalTok{ Y\_causal.shape[}\DecValTok{1}\NormalTok{]}

\NormalTok{WU }\OperatorTok{=}\NormalTok{ rng.normal(scale}\OperatorTok{=}\FloatTok{0.2}\NormalTok{, size}\OperatorTok{=}\NormalTok{(d, V))}
\NormalTok{bU }\OperatorTok{=}\NormalTok{ np.zeros(V)}

\NormalTok{logits }\OperatorTok{=}\NormalTok{ Y\_causal }\OperatorTok{@}\NormalTok{ WU }\OperatorTok{+}\NormalTok{ bU}
\NormalTok{probs }\OperatorTok{=}\NormalTok{ softmax\_rows(logits)}

\ControlFlowTok{for}\NormalTok{ t }\KeywordTok{in} \BuiltInTok{range}\NormalTok{(T):}
\NormalTok{    top }\OperatorTok{=} \BuiltInTok{int}\NormalTok{(np.argmax(probs[t]))}
    \BuiltInTok{print}\NormalTok{(}\SpecialStringTok{f"position }\SpecialCharTok{\{}\NormalTok{t}\SpecialCharTok{\}}\SpecialStringTok{: top next{-}token candidate = }\SpecialCharTok{\{}\NormalTok{vocab[top]}\SpecialCharTok{:\textgreater{}8s\}}\SpecialStringTok{, probability = }\SpecialCharTok{\{}\NormalTok{probs[t, top]}\SpecialCharTok{:.3f\}}\SpecialStringTok{"}\NormalTok{)}
\end{Highlighting}
\end{Shaded}

\begin{verbatim}
position 0: top next-token candidate =    <bos>, probability = 0.248
position 1: top next-token candidate =    <bos>, probability = 0.223
position 2: top next-token candidate =    <bos>, probability = 0.181
position 3: top next-token candidate =      the, probability = 0.183
position 4: top next-token candidate =      the, probability = 0.181
position 5: top next-token candidate =    <bos>, probability = 0.195
\end{verbatim}

During training, the true next token is known, so the model is trained
by minimizing cross-entropy across positions.

\section{15.18 Parameter-count function}\label{parameter-count-function}

The following function computes the rough parameter count from Section
15.10.

\protect\phantomsection\label{py-ch15-parameter-count}
\begin{Shaded}
\begin{Highlighting}[]
\KeywordTok{def}\NormalTok{ transformer\_block\_parameter\_count(d, dff, include\_bias}\OperatorTok{=}\VariableTok{False}\NormalTok{):}
\NormalTok{    attention }\OperatorTok{=} \DecValTok{4} \OperatorTok{*}\NormalTok{ d }\OperatorTok{*}\NormalTok{ d}
\NormalTok{    feedforward }\OperatorTok{=} \DecValTok{2} \OperatorTok{*}\NormalTok{ d }\OperatorTok{*}\NormalTok{ dff}
\NormalTok{    layernorm }\OperatorTok{=} \DecValTok{4} \OperatorTok{*}\NormalTok{ d}
\NormalTok{    total }\OperatorTok{=}\NormalTok{ attention }\OperatorTok{+}\NormalTok{ feedforward }\OperatorTok{+}\NormalTok{ layernorm}
    \ControlFlowTok{if}\NormalTok{ include\_bias:}
\NormalTok{        total }\OperatorTok{+=} \DecValTok{4} \OperatorTok{*}\NormalTok{ d      }\CommentTok{\# Q, K, V, O biases approximately}
\NormalTok{        total }\OperatorTok{+=}\NormalTok{ dff }\OperatorTok{+}\NormalTok{ d    }\CommentTok{\# FFN biases}
    \ControlFlowTok{return}\NormalTok{ \{}
        \StringTok{"attention\_weights"}\NormalTok{: attention,}
        \StringTok{"feedforward\_weights"}\NormalTok{: feedforward,}
        \StringTok{"layernorm\_parameters"}\NormalTok{: layernorm,}
        \StringTok{"total"}\NormalTok{: total,}
\NormalTok{    \}}

\ControlFlowTok{for}\NormalTok{ d }\KeywordTok{in}\NormalTok{ [}\DecValTok{128}\NormalTok{, }\DecValTok{256}\NormalTok{, }\DecValTok{512}\NormalTok{]:}
\NormalTok{    counts }\OperatorTok{=}\NormalTok{ transformer\_block\_parameter\_count(d}\OperatorTok{=}\NormalTok{d, dff}\OperatorTok{=}\DecValTok{4}\OperatorTok{*}\NormalTok{d)}
    \BuiltInTok{print}\NormalTok{(}\SpecialStringTok{f"d=}\SpecialCharTok{\{}\NormalTok{d}\SpecialCharTok{\}}\SpecialStringTok{, dff=}\SpecialCharTok{\{}\DecValTok{4}\OperatorTok{*}\NormalTok{d}\SpecialCharTok{\}}\SpecialStringTok{, total parameters per block ≈ }\SpecialCharTok{\{}\NormalTok{counts[}\StringTok{\textquotesingle{}total\textquotesingle{}}\NormalTok{]}\SpecialCharTok{:,\}}\SpecialStringTok{"}\NormalTok{)}
\end{Highlighting}
\end{Shaded}

\begin{verbatim}
d=128, dff=512, total parameters per block ≈ 197,120
d=256, dff=1024, total parameters per block ≈ 787,456
d=512, dff=2048, total parameters per block ≈ 3,147,776
\end{verbatim}

This calculation explains why feedforward layers are often a major part
of the parameter budget.

\section{15.19 Encoder-only, decoder-only, and encoder-decoder
models}\label{encoder-only-decoder-only-and-encoder-decoder-models}

The three main Transformer families differ in their attention pattern
and task objective.

{\def\LTcaptype{none} % do not increment counter
\begin{longtable}[]{@{}
  >{\raggedright\arraybackslash}p{(\linewidth - 6\tabcolsep) * \real{0.2500}}
  >{\raggedright\arraybackslash}p{(\linewidth - 6\tabcolsep) * \real{0.2500}}
  >{\raggedright\arraybackslash}p{(\linewidth - 6\tabcolsep) * \real{0.2500}}
  >{\raggedright\arraybackslash}p{(\linewidth - 6\tabcolsep) * \real{0.2500}}@{}}
\toprule\noalign{}
\begin{minipage}[b]{\linewidth}\raggedright
Architecture
\end{minipage} & \begin{minipage}[b]{\linewidth}\raggedright
Attention pattern
\end{minipage} & \begin{minipage}[b]{\linewidth}\raggedright
Main output
\end{minipage} & \begin{minipage}[b]{\linewidth}\raggedright
Typical tasks
\end{minipage} \\
\midrule\noalign{}
\endhead
\bottomrule\noalign{}
\endlastfoot
Encoder-only & bidirectional self-attention & contextual input
representation & classification, retrieval, tagging \\
Decoder-only & causal self-attention & next-token distribution &
generation, completion, chat \\
Encoder-decoder & source encoder plus target decoder cross-attention &
conditional target distribution & translation, summarization \\
\end{longtable}
}

Mathematically, encoder-only models estimate representations useful for
conditional predictions such as

\[
P(y\mid x_1,\ldots,x_T).
\]

Decoder-only models estimate left-to-right conditional distributions
such as

\[
P(w_{t+1}\mid w_1,\ldots,w_t).
\]

Encoder-decoder models estimate conditional sequence distributions such
as

\[
P(y_1,\ldots,y_m\mid x_1,\ldots,x_T)
=\prod_{t=1}^m P(y_t\mid y_1,\ldots,y_{t-1},x_1,\ldots,x_T).
\]

\section{15.20 Mathematical comparison with
RNNs}\label{mathematical-comparison-with-rnns}

An RNN computes hidden states recursively:

\[
h_t=f(h_{t-1},x_t).
\]

Thus information from token \(j\) reaches token \(i\) through a chain of
intermediate states. A Transformer computes direct token-token scores:

\[
s_{ij}=q_i^Tk_j/\sqrt{d_k}.
\]

Thus token \(i\) can attend directly to token \(j\) in one layer. This
shorter interaction path is one reason Transformers can model long-range
dependencies effectively.

However, the direct interaction comes with quadratic cost in \(T\). The
mathematical tradeoff is:

\[
\text{recurrence path length }O(T)\quad\text{versus}\quad\text{attention score matrix size }O(T^2).
\]

\section{15.21 AI-assisted learning
activities}\label{ai-assisted-learning-activities-1}

Use an AI assistant as a mathematical critic, not just as a code
generator.

\subsection{Prompt 1: dimension audit}\label{prompt-1-dimension-audit-2}

Give the assistant the formulas for multi-head attention and ask:

\begin{quote}
Check every matrix dimension in this Transformer block. Identify the
shape of each intermediate matrix and explain where shape errors are
most likely.
\end{quote}

Then verify the response by doing the calculation yourself.

\subsection{Prompt 2: residual-gradient
explanation}\label{prompt-2-residual-gradient-explanation}

Ask:

\begin{quote}
Explain why \(R(X)=X+G(X)\) helps gradient flow. Give both a scalar
explanation and a matrix derivative explanation.
\end{quote}

A good response should mention the identity term in the derivative.

\subsection{Prompt 3: architecture
classification}\label{prompt-3-architecture-classification}

Give the assistant several tasks such as sentiment classification,
next-token prediction, translation, semantic search, and summarization.
Ask it to classify which Transformer family is most natural and justify
the answer mathematically.

\subsection{Prompt 4: implementation
debugging}\label{prompt-4-implementation-debugging}

Paste your NumPy Transformer block and ask:

\begin{quote}
Find possible shape, masking, and numerical-stability bugs. Do not
rewrite the whole program. Explain each issue mathematically.
\end{quote}

\section{15.22 Common mistakes}\label{common-mistakes-4}

\begin{enumerate}
\def\labelenumi{\arabic{enumi}.}
\tightlist
\item
  \textbf{Confusing attention with the entire Transformer.} Attention is
  one sublayer. A Transformer block also uses residual connections,
  normalization, feedforward networks, and output projections.
\item
  \textbf{Forgetting the causal mask in decoder-only language models.}
  Without causal masking, training leaks future tokens.
\item
  \textbf{Applying softmax over the wrong axis.} Attention softmax
  should usually be row-wise over key positions for each query position.
\item
  \textbf{Thinking layer normalization mixes tokens.} Layer
  normalization standardizes feature coordinates within each token row.
\item
  \textbf{Assuming attention weights are explanations.} Attention
  weights are useful diagnostics, but they are not automatically causal
  explanations of model behavior.
\item
  \textbf{Ignoring sequence length cost.} The \(T\times T\) score matrix
  can dominate memory and computation for long sequences.
\end{enumerate}

\section{15.23 Exercises}\label{exercises-14}

\begin{enumerate}
\def\labelenumi{\arabic{enumi}.}
\item
  \textbf{Dimension check.} Let \(T=12\), \(d=64\), \(h=4\), and
  \(d_k=d_v=16\). Write the shapes of \(Q_r\), \(K_r\), \(V_r\),
  \(Q_rK_r^T\), \(A_r\), \(Z_r\),
  \(\operatorname{concat}(Z_1,\ldots,Z_h)\), and the final multi-head
  output.
\item
  \textbf{Residual derivative.} Suppose \(g:R\to R\) and
  \(R(x)=x+g(x)\). Show that \(R'(x)=1+g'(x)\). Explain why this is
  helpful in a deep composition.
\item
  \textbf{Layer normalization.} Compute layer normalization by hand for
  \(x=(1,3,5)\) with \(\gamma=(1,1,1)\), \(\beta=(0,0,0)\), and
  \(\varepsilon=0\).
\item
  \textbf{Parameter count.} For \(d=768\), \(h=12\), and
  \(d_{\mathrm{ff}}=3072\), compute the approximate number of parameters
  in one Transformer block, ignoring biases.
\item
  \textbf{Causal mask.} Write the causal mask matrix for \(T=5\).
  Explain why row \(3\) can attend to columns \(1,2,3\) but not columns
  \(4,5\) if indexing starts at \(1\).
\item
  \textbf{Encoder versus decoder.} Give one mathematical reason why
  bidirectional attention is inappropriate for next-token prediction
  training if no mask is used.
\item
  \textbf{Complexity.} Compare the attention cost for sequence lengths
  \(T=512\), \(T=2048\), and \(T=8192\), assuming \(d\) is fixed. How
  many times larger is \(T^2\) when \(T\) quadruples?
\item
  \textbf{Implementation.} Modify the Python multi-head attention
  example so that \(h=4\) and \(d=8\). What must be true about \(d\) and
  \(h\)?
\item
  \textbf{Interpretation.} In the interactive, set residual strength
  near \(1\). What happens to the transformed points? Explain this using
  the formula \(X+G(X)\).
\item
  \textbf{Model families.} Classify each task as encoder-only,
  decoder-only, or encoder-decoder: spam detection, machine translation,
  autocomplete, named-entity recognition, and abstractive summarization.
\end{enumerate}

\section{15.24 Mini-project: Build a tiny Transformer encoder
block}\label{mini-project-build-a-tiny-transformer-encoder-block}

Build a small Transformer encoder block from scratch in NumPy.

Your project should include:

\begin{enumerate}
\def\labelenumi{\arabic{enumi}.}
\tightlist
\item
  token embeddings for a toy vocabulary;
\item
  positional encodings or learned position vectors;
\item
  multi-head self-attention;
\item
  residual connections;
\item
  layer normalization;
\item
  a position-wise feedforward network;
\item
  a short report checking all dimensions.
\end{enumerate}

Then compare two inputs that contain the same words in different orders.
Explain how positional information changes the output.

\section{15.25 Lab connection}\label{lab-connection-6}

The Chapter 15 lab should ask students to implement a minimal
Transformer block step by step. Students should first verify shapes
using small matrices, then add masking, and finally connect the block to
a simple next-token prediction objective.

Suggested lab sections:

\begin{enumerate}
\def\labelenumi{\arabic{enumi}.}
\tightlist
\item
  matrix shapes and token embeddings;
\item
  row-wise softmax and attention;
\item
  multi-head attention;
\item
  residual connections and layer normalization;
\item
  feedforward network;
\item
  causal masking;
\item
  output logits and cross-entropy;
\item
  reflection questions comparing RNNs and Transformers.
\end{enumerate}

The goal is not to train a large model. The goal is to understand the
mathematical architecture well enough to debug real implementations.

\chapter{Positional Encoding and Sequence
Geometry}\label{positional-encoding-and-sequence-geometry}

How order enters attention-based language models

\hfill\break

Self-attention compares tokens by content. If the input to a Transformer
is only a matrix of token embeddings, then self-attention sees a set of
token vectors rather than an ordered sentence. This is a serious
mathematical problem: the sentences

\[
\text{the dog chased the cat}
\]

and

\[
\text{the cat chased the dog}
\]

contain nearly the same words, but they have different meanings because
their order is different. Positional encoding is the mechanism by which
a Transformer receives information about order.

This chapter studies position from a mathematical point of view. We
begin with permutation equivariance, then study absolute positional
embeddings, sinusoidal positional encodings, Fourier-feature geometry,
relative positions, rotary positional embeddings, and attention biases.
The goal is not to memorize architectural variations, but to understand
what mathematical information about order is injected into a sequence
model.

\section{Learning goals}\label{learning-goals-15}

After studying this chapter, you should be able to:

\begin{enumerate}
\def\labelenumi{\arabic{enumi}.}
\tightlist
\item
  explain why vanilla self-attention without position is permutation
  equivariant;
\item
  distinguish token embeddings, position embeddings, and contextual
  embeddings;
\item
  write the sinusoidal positional encoding formulas and check their
  dimensions;
\item
  explain sinusoidal encoding as a finite collection of Fourier
  features;
\item
  compute dot products between sinusoidal position vectors and interpret
  them as functions of relative distance;
\item
  compare learned absolute position embeddings with fixed sinusoidal
  encodings;
\item
  describe relative-position attention, rotary positional embeddings,
  and additive attention biases;
\item
  implement positional encodings in NumPy;
\item
  visualize position vectors and their similarity structure with
  Plotly.js;
\item
  identify common mathematical mistakes involving position in
  Transformers.
\end{enumerate}

\section{16.1 The order problem}\label{the-order-problem}

Let a sequence of length \(T\) be represented by a matrix

\[
X =
\begin{bmatrix}
- & x_1^T & - \\
- & x_2^T & - \\
  & \vdots &   \\
- & x_T^T & -
\end{bmatrix}
\in R^{T\times d},
\]

where \(x_t\in R^d\) is the input representation of the token at
position \(t\). A self-attention layer computes

\[
\operatorname{Attn}(X)
=
\operatorname{softmax}\left(
\frac{XW_Q(XW_K)^T}{\sqrt{d_k}}
\right)XW_V.
\]

This formula compares rows of \(X\) using dot products. If the rows are
permuted, the dot-product matrix is permuted in the same way. Thus
self-attention respects the symmetry of the input row order unless we
explicitly break that symmetry by adding position information.

\textbf{Key issue.} A sequence is not only a multiset of tokens. A
sequence has an order. A Transformer must therefore represent both
content and position.

For example, the two sequences

\[
(w_1,w_2,w_3)=(\text{students},\text{use},\text{AI})
\]

and

\[
(w_1,w_2,w_3)=(\text{AI},\text{use},\text{students})
\]

have the same bag of words but different syntax and meaning.
Bag-of-words models from Chapters 2--3 intentionally discard order.
Transformers should not.

\section{16.2 Permutation equivariance of
self-attention}\label{permutation-equivariance-of-self-attention}

Let \(P\in R^{T\times T}\) be a permutation matrix. Multiplying \(X\) on
the left by \(P\) permutes the rows of \(X\):

\[
X' = PX.
\]

For a single self-attention head without position, define

\[
Q=XW_Q,
\qquad
K=XW_K,
\qquad
V=XW_V.
\]

For the permuted input \(X'=PX\),

\[
Q'=PXW_Q=PQ,
\qquad
K'=PK,
\qquad
V'=PV.
\]

The score matrix becomes

\[
S'
=
\frac{Q'(K')^T}{\sqrt{d_k}}
=
\frac{(PQ)(PK)^T}{\sqrt{d_k}}
=
P\left(\frac{QK^T}{\sqrt{d_k}}\right)P^T.
\]

The row-wise softmax satisfies

\[
\operatorname{softmax}_{\text{row}}(PSP^T)
=
P\operatorname{softmax}_{\text{row}}(S)P^T.
\]

Therefore

\[
\begin{aligned}
\operatorname{Attn}(PX)
&=
\operatorname{softmax}_{\text{row}}(S')V' \\
&=
P\operatorname{softmax}_{\text{row}}(S)P^T PV \\
&=
P\operatorname{Attn}(X).
\end{aligned}
\]

So self-attention is \textbf{permutation equivariant}:

\[
\operatorname{Attn}(PX)=P\operatorname{Attn}(X).
\]

This is not wrong; it is often a useful symmetry. But for language, the
original row index matters. Positional encoding breaks pure permutation
symmetry by replacing the token-only representation with a
content-plus-position representation.

\section{16.3 Content vectors plus position
vectors}\label{content-vectors-plus-position-vectors}

Let \(e:V\to R^d\) be a token embedding map. If token \(w_t\) appears at
position \(t\), the token embedding is

\[
x_t=e(w_t).
\]

A simple positional strategy creates a vector \(p_t\in R^d\) for every
position \(t\), then uses

\[
z_t = x_t + p_t.
\]

The matrix input to the Transformer is then

\[
Z = X + P_{\text{pos}},
\]

where

\[
P_{\text{pos}}
=
\begin{bmatrix}
- & p_1^T & - \\
- & p_2^T & - \\
  & \vdots &   \\
- & p_T^T & -
\end{bmatrix}
\in R^{T\times d}.
\]

The addition is simple but mathematically powerful. Because the vector
at row \(t\) now contains both token identity and position identity,
attention scores can depend on position through inner products such as

\[
(z_i W_Q)(z_j W_K)^T.
\]

Expanding \(z_i=x_i+p_i\) and \(z_j=x_j+p_j\) shows the possible
interactions:

\[
(z_i W_Q)(z_j W_K)^T
=
(x_iW_Q)(x_jW_K)^T
+
(x_iW_Q)(p_jW_K)^T
+
(p_iW_Q)(x_jW_K)^T
+
(p_iW_Q)(p_jW_K)^T.
\]

These terms allow the model to learn rules involving content, source
position, target position, and relative placement.

\section{16.4 Learned absolute position
embeddings}\label{learned-absolute-position-embeddings}

The most direct approach is to learn a position-embedding matrix

\[
E_{\text{pos}}\in R^{T_{\max}\times d},
\]

where row \(t\) is the learned vector \(p_t^T\). For a sequence of
length \(T\le T_{\max}\), use the first \(T\) rows.

The trainable parameters are

\[
T_{\max}d.
\]

This method is flexible because the model learns the position vectors
that help the training objective. But it has two limitations.

First, the model is usually trained with a fixed maximum length
\(T_{\max}\). Extrapolating beyond this length is not natural unless
additional design choices are made.

Second, learned absolute embeddings do not automatically encode a smooth
geometry of distance. The model must learn from data that positions 10
and 11 are closer than positions 10 and 500.

\textbf{Absolute position viewpoint.} Learned absolute position
embeddings identify each index \(t\) with a trainable vector \(p_t\).
They are flexible labels for positions, but they do not automatically
impose translation or distance structure.

\section{16.5 Sinusoidal positional
encoding}\label{sinusoidal-positional-encoding}

The original Transformer introduced a fixed sinusoidal positional
encoding. Suppose \(d\) is even. For position \(t\) and coordinate pair
indexed by \(i=0,1,\ldots,d/2-1\), define

\[
PE(t,2i)=\sin\left(\frac{t}{10000^{2i/d}}\right),
\]

and

\[
PE(t,2i+1)=\cos\left(\frac{t}{10000^{2i/d}}\right).
\]

Thus each position \(t\) is mapped to a vector

\[
p_t = PE(t)\in R^d.
\]

The denominator

\[
10000^{2i/d}
\]

creates wavelengths on a logarithmic scale. Low-index coordinate pairs
vary rapidly with \(t\); high-index coordinate pairs vary slowly.

Equivalently, define frequencies

\[
\omega_i = 10000^{-2i/d}.
\]

Then the coordinate pair is

\[
(\sin(\omega_i t),\cos(\omega_i t)).
\]

This shows the Fourier-feature structure: sinusoidal encoding represents
each position through multiple periodic functions at different
frequencies.

\section{16.6 Geometry of sinusoidal
encodings}\label{geometry-of-sinusoidal-encodings}

For a single frequency \(\omega\), define

\[
u_t(\omega)=(\sin(\omega t),\cos(\omega t))\in R^2.
\]

Then

\[
\|u_t(\omega)\|_2^2
=
\sin^2(\omega t)+\cos^2(\omega t)=1.
\]

The dot product between two positions is

\[
\begin{aligned}
u_t(\omega)^T u_s(\omega)
&=\sin(\omega t)\sin(\omega s)+\cos(\omega t)\cos(\omega s)\\
&=\cos(\omega(t-s)).
\end{aligned}
\]

Therefore the similarity between two sinusoidal position vectors depends
on the relative displacement \(t-s\). With multiple frequencies,

\[
p_t^T p_s
=
\sum_i \cos(\omega_i(t-s)).
\]

This is a crucial mathematical property: although the encoding is
absolute because it assigns a vector to each \(t\), inner products
between position vectors naturally contain relative-distance
information.

\textbf{Relative-distance property.} For sinusoidal encodings, the dot
product \(p_t^T p_s\) is a function of \(t-s\). Thus the geometry of
position vectors contains relative offset information.

\section{16.7 Linear transformations of sinusoidal
encodings}\label{linear-transformations-of-sinusoidal-encodings}

There is another useful property. For a fixed shift \(\Delta\), the pair

\[
(\sin(\omega(t+\Delta)),\cos(\omega(t+\Delta)))
\]

can be obtained from

\[
(\sin(\omega t),\cos(\omega t))
\]

by a two-dimensional rotation. Using angle-addition formulas,

\[
\sin(\omega(t+\Delta))
=
\sin(\omega t)\cos(\omega\Delta)
+
\cos(\omega t)\sin(\omega\Delta),
\]

and

\[
\cos(\omega(t+\Delta))
=
\cos(\omega t)\cos(\omega\Delta)
-
\sin(\omega t)\sin(\omega\Delta).
\]

In matrix form,

\[
\begin{bmatrix}
\sin(\omega(t+\Delta))\\
\cos(\omega(t+\Delta))
\end{bmatrix}
=
\begin{bmatrix}
\cos(\omega\Delta) & \sin(\omega\Delta)\\
-\sin(\omega\Delta) & \cos(\omega\Delta)
\end{bmatrix}
\begin{bmatrix}
\sin(\omega t)\\
\cos(\omega t)
\end{bmatrix}.
\]

So a fixed relative shift is represented by a block-diagonal linear map
across the sinusoidal coordinate pairs. This helps explain why
sinusoidal encodings can support relative-position reasoning.

\section{16.8 Scaling token embeddings and position
embeddings}\label{scaling-token-embeddings-and-position-embeddings}

In many Transformer implementations, token embeddings are scaled before
adding positions:

\[
z_t=\sqrt{d}\,x_t+p_t.
\]

The reason is variance control. Suppose the coordinates of \(x_t\) have
typical scale around \(1/\sqrt{d}\). Then

\[
\|x_t\|_2
\]

may be order 1. Multiplying by \(\sqrt{d}\) changes the coordinate scale
and can balance the token content against the added positional signal,
depending on how embeddings are initialized.

From a mathematical perspective, the relative magnitude of \(x_t\) and
\(p_t\) matters because attention scores involve dot products. If
\(p_t\) is too small, the model may ignore order. If \(p_t\) is too
large, the model may overemphasize position and underemphasize content.

\section{16.9 Relative position
representations}\label{relative-position-representations}

Absolute position embeddings assign vectors to positions:

\[
t\mapsto p_t.
\]

Relative-position methods instead emphasize offsets:

\[
(i,j)\mapsto r_{i-j}.
\]

In attention, a relative-position term may be added to the score between
query position \(i\) and key position \(j\):

\[
S_{ij}
=
\frac{q_i^T k_j}{\sqrt{d_k}}
+
 b_{i-j}.
\]

Here \(b_{i-j}\) is a learned or prescribed scalar bias depending on
relative displacement. This gives the model direct access to information
such as ``nearby tokens should often matter more'' or ``previous tokens
matter in a causal language model.''

A vector-valued version uses relative embeddings \(r_{i-j}\in R^{d_k}\),
for example

\[
S_{ij}
=
\frac{q_i^T(k_j+r_{i-j})}{\sqrt{d_k}}.
\]

The main idea is that attention scores become functions of both token
similarity and relative location.

\section{16.10 Rotary positional
embeddings}\label{rotary-positional-embeddings}

Rotary positional embeddings, often abbreviated RoPE, encode position by
rotating query and key vectors. The simplest version applies a
two-dimensional rotation to each coordinate pair.

For a vector pair \((a,b)\) and angle \(\theta\), define

\[
R_\theta
\begin{bmatrix}a\\ b\end{bmatrix}
=
\begin{bmatrix}
\cos\theta & -\sin\theta\\
\sin\theta & \cos\theta
\end{bmatrix}
\begin{bmatrix}a\\ b\end{bmatrix}.
\]

At position \(t\), coordinate pair \(i\) is rotated by angle

\[
\theta_{t,i}=\omega_i t.
\]

Let \(R_t\) denote the resulting block-diagonal rotation matrix. RoPE
uses

\[
\tilde q_t = R_t q_t,
\qquad
\tilde k_s = R_s k_s.
\]

The attention score uses

\[
\tilde q_t^T\tilde k_s
=
q_t^T R_t^T R_s k_s.
\]

Because rotations satisfy

\[
R_t^T R_s = R_{s-t},
\]

the score depends on the relative offset \(s-t\) through the rotation
matrix. Thus RoPE builds relative position directly into dot-product
attention.

\textbf{RoPE viewpoint.} Instead of adding a position vector to the
token vector, RoPE rotates query and key coordinates. Dot products
between rotated vectors naturally encode relative displacement.

\section{16.11 Attention bias and distance
decay}\label{attention-bias-and-distance-decay}

Another common approach is to add a deterministic or learned distance
bias to attention scores:

\[
S_{ij}
=\frac{q_i^T k_j}{\sqrt{d_k}} + b(i,j).
\]

For a causal decoder, \(b(i,j)\) may include the mask

\[
b(i,j)=
\begin{cases}
0, & j\le i,\\
-\infty, & j>i.
\end{cases}
\]

A distance-decay bias may take the form

\[
b(i,j)=-\lambda |i-j|,
\]

or, in causal form,

\[
b(i,j)=-\lambda(i-j)\quad \text{for } j\le i.
\]

This bias makes attention to distant positions less likely unless the
content score is strong enough to overcome the penalty. From the
perspective of probability, the row-wise softmax turns scores into a
distribution over positions, so additive biases act like log-prior
preferences over where attention should be placed.

\section{16.12 Positional encoding as feature
design}\label{positional-encoding-as-feature-design}

Before deep learning, feature engineering was a central part of NLP.
Positional encoding is a modern version of feature design. It chooses a
map

\[
\psi:\{1,2,\ldots,T_{\max}\}\to R^d
\]

and combines it with the token embedding map

\[
e:V\to R^d.
\]

The model receives

\[
z_t=e(w_t)+\psi(t).
\]

Different choices of \(\psi\) encode different assumptions:

{\def\LTcaptype{none} % do not increment counter
\begin{longtable}[]{@{}
  >{\raggedright\arraybackslash}p{(\linewidth - 6\tabcolsep) * \real{0.2308}}
  >{\raggedleft\arraybackslash}p{(\linewidth - 6\tabcolsep) * \real{0.3077}}
  >{\raggedright\arraybackslash}p{(\linewidth - 6\tabcolsep) * \real{0.2308}}
  >{\raggedright\arraybackslash}p{(\linewidth - 6\tabcolsep) * \real{0.2308}}@{}}
\toprule\noalign{}
\begin{minipage}[b]{\linewidth}\raggedright
Method
\end{minipage} & \begin{minipage}[b]{\linewidth}\raggedleft
Mathematical object
\end{minipage} & \begin{minipage}[b]{\linewidth}\raggedright
Strength
\end{minipage} & \begin{minipage}[b]{\linewidth}\raggedright
Limitation
\end{minipage} \\
\midrule\noalign{}
\endhead
\bottomrule\noalign{}
\endlastfoot
Learned absolute & trainable table \(E_{\text{pos}}\) & flexible & weak
extrapolation \\
Sinusoidal & Fourier features & relative-offset geometry & fixed
frequencies \\
Relative bias & scalar \(b_{i-j}\) & direct distance modeling & less
content-position interaction \\
RoPE & block rotations \(R_t\) & relative offset in dot products &
requires paired dimensions \\
Linear bias & monotone score penalty & simple length behavior & less
expressive \\
\end{longtable}
}

The correct choice depends on the task, sequence length, training data,
and architecture.

\section{16.13 Interactive: sinusoidal position
geometry}\label{interactive-sinusoidal-position-geometry}

The following interactive figure visualizes sinusoidal positional
encoding. The left plot shows position vectors projected to the first
two sinusoidal coordinates. The right plot shows the normalized
similarity between position \(0\) and position \(t\).

\subsection{Sinusoidal positional encoding
explorer}\label{sinusoidal-positional-encoding-explorer}

Sequence length \(T\)\protect\phantomsection\label{ch16-T-val}{80}
Embedding dimension \(d\)\protect\phantomsection\label{ch16-d-val}{32}
Frequency scale\protect\phantomsection\label{ch16-scale-val}{10000}

\protect\phantomsection\label{ch16-positional-plot}

\protect\phantomsection\label{ch16-positional-status}
Loading local Plotly.js\ldots{}

\section{16.14 Python example: sinusoidal encodings from
scratch}\label{python-example-sinusoidal-encodings-from-scratch}

The following code constructs the sinusoidal matrix

\[
P\in R^{T\times d}
\]

and checks its shape.

\begin{Shaded}
\begin{Highlighting}[]
\ImportTok{import}\NormalTok{ numpy }\ImportTok{as}\NormalTok{ np}


\KeywordTok{def}\NormalTok{ sinusoidal\_encoding(T, d, scale}\OperatorTok{=}\FloatTok{10000.0}\NormalTok{):}
    \CommentTok{"""Return a T by d matrix of sinusoidal positional encodings."""}
    \ControlFlowTok{if}\NormalTok{ d }\OperatorTok{\%} \DecValTok{2} \OperatorTok{!=} \DecValTok{0}\NormalTok{:}
        \ControlFlowTok{raise} \PreprocessorTok{ValueError}\NormalTok{(}\StringTok{"d must be even for paired sine/cosine coordinates."}\NormalTok{)}
\NormalTok{    positions }\OperatorTok{=}\NormalTok{ np.arange(T)[:, }\VariableTok{None}\NormalTok{]}
\NormalTok{    pair\_ids }\OperatorTok{=}\NormalTok{ np.arange(d }\OperatorTok{//} \DecValTok{2}\NormalTok{)[}\VariableTok{None}\NormalTok{, :]}
\NormalTok{    denominators }\OperatorTok{=}\NormalTok{ scale }\OperatorTok{**}\NormalTok{ (}\DecValTok{2} \OperatorTok{*}\NormalTok{ pair\_ids }\OperatorTok{/}\NormalTok{ d)}
\NormalTok{    angles }\OperatorTok{=}\NormalTok{ positions }\OperatorTok{/}\NormalTok{ denominators}

\NormalTok{    P }\OperatorTok{=}\NormalTok{ np.zeros((T, d))}
\NormalTok{    P[:, }\DecValTok{0}\NormalTok{::}\DecValTok{2}\NormalTok{] }\OperatorTok{=}\NormalTok{ np.sin(angles)}
\NormalTok{    P[:, }\DecValTok{1}\NormalTok{::}\DecValTok{2}\NormalTok{] }\OperatorTok{=}\NormalTok{ np.cos(angles)}
    \ControlFlowTok{return}\NormalTok{ P}

\NormalTok{P }\OperatorTok{=}\NormalTok{ sinusoidal\_encoding(T}\OperatorTok{=}\DecValTok{8}\NormalTok{, d}\OperatorTok{=}\DecValTok{6}\NormalTok{)}
\BuiltInTok{print}\NormalTok{(P.shape)}
\BuiltInTok{print}\NormalTok{(np.}\BuiltInTok{round}\NormalTok{(P[:}\DecValTok{3}\NormalTok{], }\DecValTok{3}\NormalTok{))}
\end{Highlighting}
\end{Shaded}

\begin{verbatim}
(8, 6)
[[ 0.     1.     0.     1.     0.     1.   ]
 [ 0.841  0.54   0.046  0.999  0.002  1.   ]
 [ 0.909 -0.416  0.093  0.996  0.004  1.   ]]
\end{verbatim}

Each row is the position vector for one index. Each pair of columns is
one sine-cosine frequency pair.

\section{16.15 Python example: similarity depends on relative
offset}\label{python-example-similarity-depends-on-relative-offset}

For sinusoidal encodings, the dot product \(p_t^T p_s\) depends on
\(t-s\). The following experiment checks this numerically.

\begin{Shaded}
\begin{Highlighting}[]
\NormalTok{T }\OperatorTok{=} \DecValTok{20}
\NormalTok{d }\OperatorTok{=} \DecValTok{16}
\NormalTok{P }\OperatorTok{=}\NormalTok{ sinusoidal\_encoding(T, d)}

\NormalTok{G }\OperatorTok{=}\NormalTok{ P }\OperatorTok{@}\NormalTok{ P.T}

\CommentTok{\# Compare pairs with the same offset.}
\ControlFlowTok{for}\NormalTok{ offset }\KeywordTok{in}\NormalTok{ [}\DecValTok{1}\NormalTok{, }\DecValTok{2}\NormalTok{, }\DecValTok{5}\NormalTok{]:}
\NormalTok{    values }\OperatorTok{=}\NormalTok{ []}
    \ControlFlowTok{for}\NormalTok{ t }\KeywordTok{in} \BuiltInTok{range}\NormalTok{(T }\OperatorTok{{-}}\NormalTok{ offset):}
\NormalTok{        values.append(G[t, t }\OperatorTok{+}\NormalTok{ offset])}
    \BuiltInTok{print}\NormalTok{(}\StringTok{"offset"}\NormalTok{, offset, }\StringTok{"min"}\NormalTok{, }\BuiltInTok{round}\NormalTok{(}\BuiltInTok{min}\NormalTok{(values), }\DecValTok{6}\NormalTok{), }\StringTok{"max"}\NormalTok{, }\BuiltInTok{round}\NormalTok{(}\BuiltInTok{max}\NormalTok{(values), }\DecValTok{6}\NormalTok{))}
\end{Highlighting}
\end{Shaded}

\begin{verbatim}
offset 1 min 7.485166 max 7.485166
offset 2 min 6.368277 max 6.368277
offset 5 min 6.13704 max 6.13704
\end{verbatim}

The values are nearly identical for each fixed offset. Small numerical
differences may appear because of floating-point arithmetic, but the
symbolic formula says the values depend only on the offset.

\section{16.16 Python example: adding position to token
embeddings}\label{python-example-adding-position-to-token-embeddings}

The next example illustrates the standard input to a Transformer block:

\[
Z=X+P_{\text{pos}}.
\]

\begin{Shaded}
\begin{Highlighting}[]
\NormalTok{rng }\OperatorTok{=}\NormalTok{ np.random.default\_rng(}\DecValTok{7243}\NormalTok{)}

\NormalTok{T }\OperatorTok{=} \DecValTok{6}
\NormalTok{d }\OperatorTok{=} \DecValTok{8}
\NormalTok{vocab\_size }\OperatorTok{=} \DecValTok{12}

\CommentTok{\# A small token{-}embedding matrix.}
\NormalTok{E }\OperatorTok{=}\NormalTok{ rng.normal(scale}\OperatorTok{=}\FloatTok{0.2}\NormalTok{, size}\OperatorTok{=}\NormalTok{(vocab\_size, d))}

\CommentTok{\# A tokenized sequence represented by integer ids.}
\NormalTok{token\_ids }\OperatorTok{=}\NormalTok{ np.array([}\DecValTok{2}\NormalTok{, }\DecValTok{5}\NormalTok{, }\DecValTok{5}\NormalTok{, }\DecValTok{1}\NormalTok{, }\DecValTok{8}\NormalTok{, }\DecValTok{3}\NormalTok{])}

\NormalTok{X }\OperatorTok{=}\NormalTok{ E[token\_ids]}
\NormalTok{P\_pos }\OperatorTok{=}\NormalTok{ sinusoidal\_encoding(T, d)}
\NormalTok{Z }\OperatorTok{=}\NormalTok{ X }\OperatorTok{+}\NormalTok{ P\_pos}

\BuiltInTok{print}\NormalTok{(}\StringTok{"Token embedding matrix X:"}\NormalTok{, X.shape)}
\BuiltInTok{print}\NormalTok{(}\StringTok{"Position matrix P\_pos:"}\NormalTok{, P\_pos.shape)}
\BuiltInTok{print}\NormalTok{(}\StringTok{"Transformer input Z:"}\NormalTok{, Z.shape)}
\end{Highlighting}
\end{Shaded}

\begin{verbatim}
Token embedding matrix X: (6, 8)
Position matrix P_pos: (6, 8)
Transformer input Z: (6, 8)
\end{verbatim}

This example also shows why position and token embeddings must have the
same dimension \(d\) when they are combined by addition.

\section{16.17 Python example: permutation equivariance without
position}\label{python-example-permutation-equivariance-without-position}

The following simple self-attention function demonstrates the
permutation-equivariance property numerically.

\begin{Shaded}
\begin{Highlighting}[]
\KeywordTok{def}\NormalTok{ row\_softmax(S):}
\NormalTok{    S\_shifted }\OperatorTok{=}\NormalTok{ S }\OperatorTok{{-}}\NormalTok{ S.}\BuiltInTok{max}\NormalTok{(axis}\OperatorTok{=}\DecValTok{1}\NormalTok{, keepdims}\OperatorTok{=}\VariableTok{True}\NormalTok{)}
\NormalTok{    E }\OperatorTok{=}\NormalTok{ np.exp(S\_shifted)}
    \ControlFlowTok{return}\NormalTok{ E }\OperatorTok{/}\NormalTok{ E.}\BuiltInTok{sum}\NormalTok{(axis}\OperatorTok{=}\DecValTok{1}\NormalTok{, keepdims}\OperatorTok{=}\VariableTok{True}\NormalTok{)}


\KeywordTok{def}\NormalTok{ self\_attention\_no\_position(X, Wq, Wk, Wv):}
\NormalTok{    Q }\OperatorTok{=}\NormalTok{ X }\OperatorTok{@}\NormalTok{ Wq}
\NormalTok{    K }\OperatorTok{=}\NormalTok{ X }\OperatorTok{@}\NormalTok{ Wk}
\NormalTok{    V }\OperatorTok{=}\NormalTok{ X }\OperatorTok{@}\NormalTok{ Wv}
\NormalTok{    S }\OperatorTok{=}\NormalTok{ Q }\OperatorTok{@}\NormalTok{ K.T }\OperatorTok{/}\NormalTok{ np.sqrt(Q.shape[}\DecValTok{1}\NormalTok{])}
\NormalTok{    A }\OperatorTok{=}\NormalTok{ row\_softmax(S)}
    \ControlFlowTok{return}\NormalTok{ A }\OperatorTok{@}\NormalTok{ V}

\NormalTok{rng }\OperatorTok{=}\NormalTok{ np.random.default\_rng(}\DecValTok{16}\NormalTok{)}
\NormalTok{T, d, dk, dv }\OperatorTok{=} \DecValTok{5}\NormalTok{, }\DecValTok{4}\NormalTok{, }\DecValTok{3}\NormalTok{, }\DecValTok{3}
\NormalTok{X }\OperatorTok{=}\NormalTok{ rng.normal(size}\OperatorTok{=}\NormalTok{(T, d))}
\NormalTok{Wq }\OperatorTok{=}\NormalTok{ rng.normal(size}\OperatorTok{=}\NormalTok{(d, dk))}
\NormalTok{Wk }\OperatorTok{=}\NormalTok{ rng.normal(size}\OperatorTok{=}\NormalTok{(d, dk))}
\NormalTok{Wv }\OperatorTok{=}\NormalTok{ rng.normal(size}\OperatorTok{=}\NormalTok{(d, dv))}

\NormalTok{perm }\OperatorTok{=}\NormalTok{ np.array([}\DecValTok{2}\NormalTok{, }\DecValTok{0}\NormalTok{, }\DecValTok{4}\NormalTok{, }\DecValTok{1}\NormalTok{, }\DecValTok{3}\NormalTok{])}
\NormalTok{Pmat }\OperatorTok{=}\NormalTok{ np.eye(T)[perm]}

\NormalTok{Y }\OperatorTok{=}\NormalTok{ self\_attention\_no\_position(X, Wq, Wk, Wv)}
\NormalTok{Y\_perm }\OperatorTok{=}\NormalTok{ self\_attention\_no\_position(Pmat }\OperatorTok{@}\NormalTok{ X, Wq, Wk, Wv)}

\BuiltInTok{print}\NormalTok{(np.allclose(Y\_perm, Pmat }\OperatorTok{@}\NormalTok{ Y))}
\end{Highlighting}
\end{Shaded}

\begin{verbatim}
True
\end{verbatim}

The output \texttt{True} means that permuting the input rows only
permutes the output rows. Without position, the layer does not know
which row was originally first, second, or third.

\section{16.18 Python example: rotary position
operation}\label{python-example-rotary-position-operation}

Here is a simple implementation of paired coordinate rotations. For
clarity, this code uses the same frequency pattern as sinusoidal
encoding.

\begin{Shaded}
\begin{Highlighting}[]
\KeywordTok{def}\NormalTok{ apply\_rope(X, scale}\OperatorTok{=}\FloatTok{10000.0}\NormalTok{):}
    \CommentTok{"""Apply rotary position rotations to a T by d matrix X with even d."""}
\NormalTok{    T, d }\OperatorTok{=}\NormalTok{ X.shape}
    \ControlFlowTok{if}\NormalTok{ d }\OperatorTok{\%} \DecValTok{2} \OperatorTok{!=} \DecValTok{0}\NormalTok{:}
        \ControlFlowTok{raise} \PreprocessorTok{ValueError}\NormalTok{(}\StringTok{"d must be even for pairwise rotations."}\NormalTok{)}
\NormalTok{    positions }\OperatorTok{=}\NormalTok{ np.arange(T)[:, }\VariableTok{None}\NormalTok{]}
\NormalTok{    pair\_ids }\OperatorTok{=}\NormalTok{ np.arange(d }\OperatorTok{//} \DecValTok{2}\NormalTok{)[}\VariableTok{None}\NormalTok{, :]}
\NormalTok{    omega }\OperatorTok{=}\NormalTok{ scale }\OperatorTok{**}\NormalTok{ (}\OperatorTok{{-}}\DecValTok{2} \OperatorTok{*}\NormalTok{ pair\_ids }\OperatorTok{/}\NormalTok{ d)}
\NormalTok{    theta }\OperatorTok{=}\NormalTok{ positions }\OperatorTok{*}\NormalTok{ omega}

\NormalTok{    cos }\OperatorTok{=}\NormalTok{ np.cos(theta)}
\NormalTok{    sin }\OperatorTok{=}\NormalTok{ np.sin(theta)}

\NormalTok{    X\_even }\OperatorTok{=}\NormalTok{ X[:, }\DecValTok{0}\NormalTok{::}\DecValTok{2}\NormalTok{]}
\NormalTok{    X\_odd }\OperatorTok{=}\NormalTok{ X[:, }\DecValTok{1}\NormalTok{::}\DecValTok{2}\NormalTok{]}

\NormalTok{    Y }\OperatorTok{=}\NormalTok{ np.empty\_like(X)}
\NormalTok{    Y[:, }\DecValTok{0}\NormalTok{::}\DecValTok{2}\NormalTok{] }\OperatorTok{=}\NormalTok{ X\_even }\OperatorTok{*}\NormalTok{ cos }\OperatorTok{{-}}\NormalTok{ X\_odd }\OperatorTok{*}\NormalTok{ sin}
\NormalTok{    Y[:, }\DecValTok{1}\NormalTok{::}\DecValTok{2}\NormalTok{] }\OperatorTok{=}\NormalTok{ X\_even }\OperatorTok{*}\NormalTok{ sin }\OperatorTok{+}\NormalTok{ X\_odd }\OperatorTok{*}\NormalTok{ cos}
    \ControlFlowTok{return}\NormalTok{ Y}

\NormalTok{Q }\OperatorTok{=}\NormalTok{ rng.normal(size}\OperatorTok{=}\NormalTok{(}\DecValTok{6}\NormalTok{, }\DecValTok{8}\NormalTok{))}
\NormalTok{K }\OperatorTok{=}\NormalTok{ rng.normal(size}\OperatorTok{=}\NormalTok{(}\DecValTok{6}\NormalTok{, }\DecValTok{8}\NormalTok{))}

\NormalTok{Q\_rot }\OperatorTok{=}\NormalTok{ apply\_rope(Q)}
\NormalTok{K\_rot }\OperatorTok{=}\NormalTok{ apply\_rope(K)}
\NormalTok{S\_rot }\OperatorTok{=}\NormalTok{ Q\_rot }\OperatorTok{@}\NormalTok{ K\_rot.T }\OperatorTok{/}\NormalTok{ np.sqrt(Q.shape[}\DecValTok{1}\NormalTok{])}

\BuiltInTok{print}\NormalTok{(S\_rot.shape)}
\BuiltInTok{print}\NormalTok{(np.}\BuiltInTok{round}\NormalTok{(S\_rot[:}\DecValTok{2}\NormalTok{, :}\DecValTok{4}\NormalTok{], }\DecValTok{3}\NormalTok{))}
\end{Highlighting}
\end{Shaded}

\begin{verbatim}
(6, 6)
[[ 0.812 -0.286 -1.151  0.467]
 [-0.245 -0.604 -0.397 -0.75 ]]
\end{verbatim}

RoPE is not just a position vector added to \(X\). It changes the
geometry of dot products by rotating queries and keys according to their
positions.

\section{16.19 Python example: Plotly
visualization}\label{python-example-plotly-visualization}

The interactive above is written directly in JavaScript so it works in
the rendered HTML book. The same idea can be explored in Python with
Plotly.

\begin{Shaded}
\begin{Highlighting}[]
\ImportTok{import}\NormalTok{ plotly.graph\_objects }\ImportTok{as}\NormalTok{ go}
\ImportTok{from}\NormalTok{ plotly.subplots }\ImportTok{import}\NormalTok{ make\_subplots}

\NormalTok{T }\OperatorTok{=} \DecValTok{80}
\NormalTok{d }\OperatorTok{=} \DecValTok{32}
\NormalTok{P }\OperatorTok{=}\NormalTok{ sinusoidal\_encoding(T, d)}
\NormalTok{G0 }\OperatorTok{=}\NormalTok{ (P }\OperatorTok{@}\NormalTok{ P[}\DecValTok{0}\NormalTok{]) }\OperatorTok{/}\NormalTok{ (P[}\DecValTok{0}\NormalTok{] }\OperatorTok{@}\NormalTok{ P[}\DecValTok{0}\NormalTok{])}

\NormalTok{fig }\OperatorTok{=}\NormalTok{ make\_subplots(rows}\OperatorTok{=}\DecValTok{1}\NormalTok{, cols}\OperatorTok{=}\DecValTok{2}\NormalTok{, subplot\_titles}\OperatorTok{=}\NormalTok{(}\StringTok{"First sine{-}cosine pair"}\NormalTok{, }\StringTok{"Similarity to position 0"}\NormalTok{))}
\NormalTok{fig.add\_trace(}
\NormalTok{    go.Scatter(x}\OperatorTok{=}\NormalTok{P[:, }\DecValTok{0}\NormalTok{], y}\OperatorTok{=}\NormalTok{P[:, }\DecValTok{1}\NormalTok{], mode}\OperatorTok{=}\StringTok{"lines+markers"}\NormalTok{, name}\OperatorTok{=}\StringTok{"positions"}\NormalTok{),}
\NormalTok{    row}\OperatorTok{=}\DecValTok{1}\NormalTok{, col}\OperatorTok{=}\DecValTok{1}
\NormalTok{)}
\NormalTok{fig.add\_trace(}
\NormalTok{    go.Scatter(x}\OperatorTok{=}\NormalTok{np.arange(T), y}\OperatorTok{=}\NormalTok{G0, mode}\OperatorTok{=}\StringTok{"lines"}\NormalTok{, name}\OperatorTok{=}\StringTok{"similarity"}\NormalTok{),}
\NormalTok{    row}\OperatorTok{=}\DecValTok{1}\NormalTok{, col}\OperatorTok{=}\DecValTok{2}
\NormalTok{)}
\NormalTok{fig.update\_xaxes(title\_text}\OperatorTok{=}\StringTok{"sin coordinate"}\NormalTok{, row}\OperatorTok{=}\DecValTok{1}\NormalTok{, col}\OperatorTok{=}\DecValTok{1}\NormalTok{)}
\NormalTok{fig.update\_yaxes(title\_text}\OperatorTok{=}\StringTok{"cos coordinate"}\NormalTok{, row}\OperatorTok{=}\DecValTok{1}\NormalTok{, col}\OperatorTok{=}\DecValTok{1}\NormalTok{)}
\NormalTok{fig.update\_xaxes(title\_text}\OperatorTok{=}\StringTok{"offset from position 0"}\NormalTok{, row}\OperatorTok{=}\DecValTok{1}\NormalTok{, col}\OperatorTok{=}\DecValTok{2}\NormalTok{)}
\NormalTok{fig.update\_yaxes(title\_text}\OperatorTok{=}\StringTok{"normalized dot product"}\NormalTok{, row}\OperatorTok{=}\DecValTok{1}\NormalTok{, col}\OperatorTok{=}\DecValTok{2}\NormalTok{)}
\NormalTok{fig.update\_layout(height}\OperatorTok{=}\DecValTok{420}\NormalTok{, showlegend}\OperatorTok{=}\VariableTok{False}\NormalTok{)}
\NormalTok{fig.show()}
\end{Highlighting}
\end{Shaded}

\begin{verbatim}
Unable to display output for mime type(s): text/html
\end{verbatim}

\begin{verbatim}
Unable to display output for mime type(s): text/html
\end{verbatim}

\section{16.20 AI-assisted learning
components}\label{ai-assisted-learning-components-1}

Use an AI assistant as a mathematical critic, not only as a code
generator.

\subsection{Prompt 1: prove the symmetry
problem}\label{prompt-1-prove-the-symmetry-problem}

Ask:

\begin{quote}
Prove that self-attention without positional information is permutation
equivariant. Use matrices and explain every dimension.
\end{quote}

Then check whether the proof correctly uses

\[
P^TP=I
\]

and whether it handles the row-wise softmax carefully.

\subsection{Prompt 2: compare position
methods}\label{prompt-2-compare-position-methods}

Ask:

\begin{quote}
Compare learned absolute positional embeddings, sinusoidal encodings,
relative-position biases, and RoPE from the viewpoint of inner products
and extrapolation to longer sequences.
\end{quote}

A good answer should mention trainable tables, Fourier features,
additive score biases, and rotations of query/key coordinates.

\subsection{Prompt 3: audit an
implementation}\label{prompt-3-audit-an-implementation}

Give the assistant your positional-encoding code and ask:

\begin{quote}
Check the dimension conventions. Are positions rows or columns? Are sine
and cosine coordinates interleaved correctly? Does the code work for odd
embedding dimension?
\end{quote}

Then verify the answer by running tests yourself.

\section{16.21 Common mistakes}\label{common-mistakes-5}

\begin{enumerate}
\def\labelenumi{\arabic{enumi}.}
\tightlist
\item
  \textbf{Confusing token order with vocabulary order.} The vocabulary
  index of a token is not its sentence position.
\item
  \textbf{Forgetting that self-attention is permutation equivariant
  without position.} Attention alone does not create order information.
\item
  \textbf{Using mismatched dimensions.} If token and position vectors
  are added, both must lie in \(R^d\).
\item
  \textbf{Assuming sinusoidal encoding is random.} It is deterministic
  and structured.
\item
  \textbf{Assuming all positional methods are equivalent.} Absolute
  embeddings, relative biases, and RoPE affect the model in different
  mathematical ways.
\item
  \textbf{Ignoring sequence length.} A method that works for short
  sequences may behave differently for long-context modeling.
\end{enumerate}

\section{16.22 Worked calculation}\label{worked-calculation-1}

Let \(d=4\) and frequencies \(\omega_0=1\), \(\omega_1=0.1\). The
sinusoidal position vector is

\[
p_t=(\sin t,\cos t,\sin(0.1t),\cos(0.1t)).
\]

Then

\[
p_t^Tp_s
=
\cos(t-s)+\cos(0.1(t-s)).
\]

In particular,

\[
p_3^Tp_1
=\cos(2)+\cos(0.2).
\]

The result depends on \(3-1=2\), not on the absolute pair \((3,1)\)
separately. This illustrates how sinusoidal encodings provide
relative-offset information through inner products.

\section{16.23 Exercises}\label{exercises-15}

\begin{enumerate}
\def\labelenumi{\arabic{enumi}.}
\tightlist
\item
  Prove that if \(F(X)=XW\) is applied row-wise, then \(F(PX)=PF(X)\)
  for every permutation matrix \(P\).
\item
  Prove the permutation-equivariance result for a single self-attention
  head.
\item
  Show that the row-wise softmax satisfies
  \(\operatorname{softmax}_{\text{row}}(PSP^T)=P\operatorname{softmax}_{\text{row}}(S)P^T\).
\item
  For a sinusoidal encoding with one frequency \(\omega\), prove that
  \(u_t(\omega)^Tu_s(\omega)=\cos(\omega(t-s))\).
\item
  Derive the block rotation matrix that maps \(u_t(\omega)\) to
  \(u_{t+\Delta}(\omega)\).
\item
  Implement sinusoidal encoding for \(T=100\) and \(d=32\). Plot the
  similarity matrix \(PP^T\).
\item
  Compare learned absolute position embeddings and sinusoidal encodings
  in terms of parameter count.
\item
  Implement a relative-position bias matrix \(B\) with entries
  \(B_{ij}=-\lambda |i-j|\). How does increasing \(\lambda\) change
  attention weights?
\item
  Implement the causal relative bias \(B_{ij}=-\infty\) for \(j>i\) and
  \(B_{ij}=-\lambda(i-j)\) for \(j\le i\).
\item
  Use RoPE on random queries and keys. Compare attention score matrices
  before and after rotation.
\item
  Explain why RoPE requires an even number of rotated coordinates, or
  else a convention for the leftover coordinate.
\item
  Write a short mathematical paragraph explaining why position is not a
  small implementation detail but a structural part of the Transformer
  model.
\end{enumerate}

\section{16.24 Mini-project}\label{mini-project-2}

Choose one short corpus, such as course descriptions, news headlines, or
theorem statements. Build a tiny Transformer-style input representation
in three ways:

\begin{enumerate}
\def\labelenumi{\arabic{enumi}.}
\tightlist
\item
  token embeddings only;
\item
  token embeddings plus learned absolute position embeddings;
\item
  token embeddings plus sinusoidal position encodings.
\end{enumerate}

For each method, compute the self-attention score matrix before training
using the same random \(W_Q\) and \(W_K\). Compare how the score
structure changes. Your report should include:

\begin{itemize}
\tightlist
\item
  the mathematical formulas for the three input representations;
\item
  at least one heatmap for each score matrix;
\item
  a discussion of permutation equivariance;
\item
  a discussion of whether the model can distinguish word order before
  training.
\end{itemize}

\section{16.25 Lab connection}\label{lab-connection-7}

The accompanying lab for this chapter should ask students to build
positional encodings from scratch, test permutation equivariance,
visualize sinusoidal similarity matrices, implement RoPE, and compare
attention weights with and without distance biases.

A useful independent-study lab structure is:

\begin{enumerate}
\def\labelenumi{\arabic{enumi}.}
\tightlist
\item
  implement sinusoidal position encoding;
\item
  visualize the encoding matrix and similarity matrix;
\item
  prove and test permutation equivariance of attention;
\item
  add absolute position vectors and show that equivariance is broken in
  the desired way;
\item
  implement relative-position bias;
\item
  implement RoPE;
\item
  write a short reflection on how sequence geometry changes attention.
\end{enumerate}

\section{16.26 Summary}\label{summary-2}

Self-attention is powerful because every token can compare itself with
every other token. But without position information, self-attention is
permutation equivariant: it does not know the original order of the
rows. Positional encoding solves this problem by injecting order into
the input vectors or attention scores.

The simplest method adds position vectors to token embeddings. Learned
absolute embeddings are flexible but may not extrapolate naturally.
Sinusoidal encodings use Fourier features and have dot products that
depend on relative offsets. Relative-position biases add displacement
information directly to attention scores. RoPE rotates query and key
vectors so their dot products encode relative displacement.

From a mathematical perspective, positional encoding is the study of how
to represent the ordered set

\[
\{1,2,\ldots,T\}
\]

inside the geometry of a Transformer. It is where sequence structure
enters attention-based language modeling.

\part{Part VI. Large Language Models and Retrieval}

\chapter{Autoregressive Language
Models}\label{autoregressive-language-models}

Next-token prediction, likelihood, and mathematical generation

\hfill\break

This chapter begins the part of the book on large language models. The
goal is to explain an autoregressive language model as a clean
mathematical object: a conditional probability model for the next token.
Once this viewpoint is clear, training, perplexity, temperature,
top-\(k\) sampling, nucleus sampling, greedy decoding, and beam search
become consequences of probability and optimization rather than
mysterious engineering tricks.

\section{Learning goals}\label{learning-goals-16}

After studying this chapter, you should be able to:

\begin{itemize}
\tightlist
\item
  Write the probability of a token sequence using the autoregressive
  chain rule.
\item
  Explain next-token prediction as conditional distribution estimation.
\item
  Derive the negative log-likelihood and cross-entropy objective for
  language modeling.
\item
  Compute sentence log-probability and perplexity from next-token
  probabilities.
\item
  Compare greedy decoding, sampling, temperature scaling, top-\(k\)
  sampling, nucleus sampling, and beam search.
\item
  Explain why causal masking is required in decoder-only Transformer
  language models.
\item
  Implement a small autoregressive generator and decoding methods from
  scratch in Python.
\end{itemize}

\section{17.1 The basic object: a distribution on
sequences}\label{the-basic-object-a-distribution-on-sequences}

Let \(V\) be a finite vocabulary with \(|V|=d\). A token sequence of
length \(T\) is

\[
w_{1:T} = (w_1,w_2,\ldots,w_T) \in V^T.
\]

A language model assigns probabilities to sequences. For a fixed length
\(T\), this means a probability mass function

\[
P(w_{1:T}), \qquad w_{1:T} \in V^T,
\]

satisfying

\[
P(w_{1:T}) \geq 0, \qquad \sum_{w_{1:T} \in V^T} P(w_{1:T}) = 1.
\]

The space \(V^T\) has size \(d^T\), so directly specifying \(P\) is
impossible except for tiny vocabularies and tiny lengths. The central
mathematical idea is to factor the joint distribution into conditional
distributions.

By the chain rule of probability,

\[
P(w_{1:T})
=
P(w_1)P(w_2 \mid w_1)P(w_3 \mid w_{1:2})\cdots P(w_T \mid w_{1:T-1}).
\]

Equivalently,

\[
P(w_{1:T})
=
\prod_{t=1}^{T} P(w_t \mid w_{<t}),
\]

where \(w_{<t}=(w_1,\ldots,w_{t-1})\). This identity is exact. An
autoregressive language model is a parameterized approximation of these
conditional probabilities.

\begin{tcolorbox}[enhanced jigsaw, arc=.35mm, bottomrule=.15mm, bottomtitle=1mm, breakable, colback=white, colbacktitle=quarto-callout-note-color!10!white, colframe=quarto-callout-note-color-frame, coltitle=black, left=2mm, leftrule=.75mm, opacityback=0, opacitybacktitle=0.6, rightrule=.15mm, title=\textcolor{quarto-callout-note-color}{\faInfo}\hspace{0.5em}{Autoregressive viewpoint}, titlerule=0mm, toprule=.15mm, toptitle=1mm]

An \textbf{autoregressive language model} estimates

\[
p_\theta(w_t \mid w_{<t})
\]

for each position \(t\). The model then defines a sequence distribution
by

\[
p_\theta(w_{1:T})
=
\prod_{t=1}^{T} p_\theta(w_t \mid w_{<t}).
\]

\end{tcolorbox}

The word \emph{autoregressive} means that the model generates each new
token using previously generated tokens as input. In time series
notation, the future is predicted from the past. In NLP notation, the
next token is predicted from the prefix.

\section{17.2 Prefixes, contexts, and conditional
distributions}\label{prefixes-contexts-and-conditional-distributions}

For every prefix \(c=w_{<t}\), an autoregressive model outputs a
probability vector

\[
p_\theta(\cdot \mid c) \in \Delta^{d-1},
\]

where

\[
\Delta^{d-1}
=
\left\{p\in R^d: p_j\geq 0,\; \sum_{j=1}^d p_j=1\right\}
\]

is the probability simplex. If the vocabulary is ordered as

\[
V=\{v_1,\ldots,v_d\},
\]

then the model output is

\[
p_\theta(\cdot \mid c)
=
\left(p_\theta(v_1\mid c),\ldots,p_\theta(v_d\mid c)\right).
\]

A neural language model usually produces real-valued logits

\[
z_\theta(c) \in R^d,
\]

and converts them to probabilities by the softmax map

\[
p_\theta(v_j\mid c)
=
\frac{\exp(z_j(c))}{\sum_{\ell=1}^{d}\exp(z_\ell(c))}.
\]

The logits are scores. The softmax turns scores into a categorical
distribution.

\begin{tcolorbox}[enhanced jigsaw, arc=.35mm, bottomrule=.15mm, bottomtitle=1mm, breakable, colback=white, colbacktitle=quarto-callout-important-color!10!white, colframe=quarto-callout-important-color-frame, coltitle=black, left=2mm, leftrule=.75mm, opacityback=0, opacitybacktitle=0.6, rightrule=.15mm, title=\textcolor{quarto-callout-important-color}{\faExclamation}\hspace{0.5em}{Representation, model, and probability}, titlerule=0mm, toprule=.15mm, toptitle=1mm]

For a prefix \(c\), a modern neural language model usually has the form

\[
c \longmapsto h_\theta(c) \in R^m
\longmapsto z_\theta(c)=Wh_\theta(c)+b \in R^d
\longmapsto p_\theta(\cdot\mid c)=\operatorname{softmax}(z_\theta(c)).
\]

The representation \(h_\theta(c)\) is not itself a probability
distribution. The probability distribution appears after the output
logits are transformed by softmax.

\end{tcolorbox}

\section{17.3 Training data and maximum
likelihood}\label{training-data-and-maximum-likelihood}

Suppose the training corpus consists of token sequences

\[
D=\{w^{(i)}_{1:T_i}: i=1,\ldots,n\}.
\]

The likelihood of the data under an autoregressive model is

\[
L(\theta;D)
=
\prod_{i=1}^{n}\prod_{t=1}^{T_i}
p_\theta\left(w_t^{(i)} \mid w_{<t}^{(i)}\right).
\]

The log-likelihood is

\[
\ell(\theta;D)
=
\sum_{i=1}^{n}\sum_{t=1}^{T_i}
\log p_\theta\left(w_t^{(i)} \mid w_{<t}^{(i)}\right).
\]

Training usually minimizes the negative log-likelihood,

\[
\mathcal L(\theta)
=
-
\sum_{i=1}^{n}\sum_{t=1}^{T_i}
\log p_\theta\left(w_t^{(i)} \mid w_{<t}^{(i)}\right).
\]

This is also the sum of cross-entropy losses. At each position \(t\),
the target is the one-hot vector for the observed next token. If \(y_t\)
is the index of \(w_t\), then the one-position loss is

\[
\ell_t(\theta)=-\log p_\theta(y_t\mid w_{<t}).
\]

For logits \(z\in R^d\) and target class \(y\), the gradient of the
softmax cross-entropy is

\[
\frac{\partial}{\partial z_j}\left[-\log \operatorname{softmax}(z)_y\right]
=
\operatorname{softmax}(z)_j - 1\{j=y\}.
\]

This simple formula is one reason next-token prediction scales well
computationally.

\section{17.4 Teacher forcing}\label{teacher-forcing}

During training, the model is given the true prefix from the corpus. For
example, from the sentence

\begin{Shaded}
\begin{Highlighting}[]
\NormalTok{language models predict the next token}
\end{Highlighting}
\end{Shaded}

we form training pairs:

\begin{Shaded}
\begin{Highlighting}[]
\NormalTok{prefix: language                 target: models}
\NormalTok{prefix: language models          target: predict}
\NormalTok{prefix: language models predict  target: the}
\end{Highlighting}
\end{Shaded}

This training approach is called \textbf{teacher forcing}. The model
learns from correct prefixes even though, at generation time, it must
condition on tokens that it generated itself.

Mathematically, training minimizes empirical conditional cross-entropy
under the empirical distribution of true prefixes. Generation samples
from the model distribution over its own prefixes. This mismatch is one
source of error accumulation in generated text.

\section{17.5 Sentence probability, average loss, and
perplexity}\label{sentence-probability-average-loss-and-perplexity}

For one sequence \(w_{1:T}\), the model log-probability is

\[
\log p_\theta(w_{1:T})
=
\sum_{t=1}^{T}\log p_\theta(w_t\mid w_{<t}).
\]

The average negative log-likelihood per token is

\[
\bar{\mathcal L}
=
-
\frac{1}{T}
\sum_{t=1}^{T}\log p_\theta(w_t\mid w_{<t}).
\]

The perplexity is

\[
\operatorname{PPL}
=
\exp(\bar{\mathcal L}).
\]

Perplexity can be interpreted as an effective average number of
plausible next-token choices. Lower perplexity usually means the model
assigns higher probability to the observed text, but perplexity is not a
complete measure of usefulness or truthfulness.

\begin{tcolorbox}[enhanced jigsaw, arc=.35mm, bottomrule=.15mm, bottomtitle=1mm, breakable, colback=white, colbacktitle=quarto-callout-warning-color!10!white, colframe=quarto-callout-warning-color-frame, coltitle=black, left=2mm, leftrule=.75mm, opacityback=0, opacitybacktitle=0.6, rightrule=.15mm, title=\textcolor{quarto-callout-warning-color}{\faExclamationTriangle}\hspace{0.5em}{Perplexity is not understanding}, titlerule=0mm, toprule=.15mm, toptitle=1mm]

A model can have low perplexity on a dataset by learning strong
statistical regularities. This does not imply that it reasons correctly,
uses facts reliably, or understands meaning in a human sense. Perplexity
is a likelihood-based metric, not a complete evaluation of language
behavior.

\end{tcolorbox}

\section{17.6 Decoder-only Transformers as autoregressive
models}\label{decoder-only-transformers-as-autoregressive-models}

A decoder-only Transformer represents the prefix \(w_{<t}\) using a
stack of causal self-attention blocks. If the input sequence is

\[
w_{1:T},
\]

then each position \(t\) produces a hidden vector \(h_t\in R^m\). The
output logits for the next token are often computed by

\[
z_t = Wh_t + b \in R^d,
\]

and the next-token distribution is

\[
p_\theta(w_{t+1}=v_j\mid w_{1:t})
=
\frac{\exp((z_t)_j)}{\sum_{\ell=1}^{d}\exp((z_t)_\ell)}.
\]

The causal mask ensures that the hidden vector at position \(t\) cannot
use tokens from positions greater than \(t\). Without causal masking, a
language model trained to predict the next token could leak information
from the future.

For a sequence length \(T\), the causal attention mask has entries

\[
M_{ij}=
\begin{cases}
0, & j\leq i,\\
-\infty, & j>i.
\end{cases}
\]

The attention weights are computed as

\[
A=\operatorname{softmax}\left(\frac{QK^T}{\sqrt{d_k}}+M\right),
\]

where the softmax is applied row by row.

\section{17.7 Generation as repeated
sampling}\label{generation-as-repeated-sampling}

Once a model produces \(p_\theta(\cdot\mid c)\), generation repeats
three steps:

\begin{enumerate}
\def\labelenumi{\arabic{enumi}.}
\tightlist
\item
  Compute logits \(z_\theta(c)\) for the current context \(c\).
\item
  Convert logits to a probability distribution, possibly modified by a
  decoding rule.
\item
  Choose the next token and append it to the context.
\end{enumerate}

This gives a recursive stochastic process:

\[
W_t \sim p_\theta(\cdot\mid W_{<t}),
\qquad t=1,2,\ldots.
\]

Different decoding rules define different stochastic or deterministic
processes.

\section{17.8 Greedy decoding}\label{greedy-decoding}

Greedy decoding chooses the most likely next token:

\[
\hat w_t
=
\arg\max_{v\in V} p_\theta(v\mid \hat w_{<t}).
\]

Greedy decoding is deterministic. It is simple and often produces fluent
text, but it can become repetitive and may miss globally better
sequences. It is a local optimization rule, not a global
maximum-likelihood sequence search.

\section{17.9 Temperature scaling}\label{temperature-scaling}

Given logits \(z\in R^d\), temperature scaling defines

\[
p_j(\tau)
=
\frac{\exp(z_j/\tau)}{\sum_{\ell=1}^{d}\exp(z_\ell/\tau)},
\qquad \tau>0.
\]

Small temperature sharpens the distribution. Large temperature flattens
it. As \(\tau\to 0\), the distribution concentrates on the largest
logit. As \(\tau\to \infty\), it approaches the uniform distribution.

Temperature does not change the ordering of logits, but it changes
entropy. The entropy is

\[
H(p(\tau))=-\sum_{j=1}^{d}p_j(\tau)\log p_j(\tau).
\]

In applications, temperature is a diversity-control parameter.

\section{\texorpdfstring{17.10 Top-\(k\)
sampling}{17.10 Top-k sampling}}\label{top-k-sampling}

Top-\(k\) sampling keeps only the \(k\) highest-probability tokens. Let
\(S_k\) be the set of \(k\) token indices with largest probabilities.
The top-\(k\) distribution is

\[
q_j
=
\begin{cases}
\dfrac{p_j}{\sum_{\ell\in S_k}p_\ell}, & j\in S_k,\\
0, & j\notin S_k.
\end{cases}
\]

Then the next token is sampled from \(q\). Top-\(k\) sampling prevents
extremely unlikely tokens from being selected while still allowing
randomness among likely tokens.

\section{\texorpdfstring{17.11 Nucleus sampling, or top-\(p\)
sampling}{17.11 Nucleus sampling, or top-p sampling}}\label{nucleus-sampling-or-top-p-sampling}

Nucleus sampling chooses the smallest set \(S_p\) of high-probability
tokens whose cumulative probability is at least a threshold \(p_0\).
After sorting probabilities in descending order,

\[
p_{(1)}\geq p_{(2)}\geq \cdots \geq p_{(d)},
\]

choose the smallest \(m\) such that

\[
\sum_{r=1}^{m}p_{(r)}\geq p_0.
\]

The nucleus is

\[
S_p=\{(1),\ldots,(m)\}.
\]

Then probabilities are renormalized on \(S_p\). Unlike top-\(k\), the
number of allowed tokens changes with the confidence of the
distribution. If the distribution is sharp, the nucleus may be small. If
the distribution is flat, the nucleus may be larger.

\section{17.12 Beam search}\label{beam-search}

Beam search is an approximate search for a high-probability sequence. It
maintains a set of \(B\) candidate prefixes called the beam. At each
step, each candidate prefix is extended by possible next tokens, and
only the \(B\) highest-scoring extended prefixes are kept.

The score of a sequence is usually its log-probability:

\[
s(w_{1:T})
=
\sum_{t=1}^{T}\log p_\theta(w_t\mid w_{<t}).
\]

Because log-probabilities are usually negative, longer sequences tend to
have smaller scores. Some implementations use length normalization, for
example

\[
s_\alpha(w_{1:T})
=
\frac{1}{T^\alpha}
\sum_{t=1}^{T}\log p_\theta(w_t\mid w_{<t}),
\]

where \(\alpha\geq 0\).

Beam search is common in translation and structured generation, but for
open-ended generation it can produce dull or repetitive text because it
aggressively seeks high-probability continuations.

\section{17.13 Worked example: likelihood and
perplexity}\label{worked-example-likelihood-and-perplexity}

Suppose a vocabulary contains four tokens:

\[
V=\{\texttt{math},\texttt{models},\texttt{learn},\texttt{data}\}.
\]

A model assigns the following next-token probabilities to the observed
sequence

\begin{Shaded}
\begin{Highlighting}[]
\NormalTok{math models learn data}
\end{Highlighting}
\end{Shaded}

at the four positions:

\[
0.40,\quad 0.25,\quad 0.50,\quad 0.20.
\]

The sequence probability is

\[
0.40\cdot 0.25\cdot 0.50\cdot 0.20=0.01.
\]

The log-probability is

\[
\log(0.01)\approx -4.6052.
\]

The average negative log-likelihood is

\[
-\frac{1}{4}\log(0.01)\approx 1.1513.
\]

The perplexity is

\[
\exp(1.1513)\approx 3.1623.
\]

This means that, on average, the model behaves as if it were choosing
among about \(3.16\) equally likely next-token options at each step.

\section{17.14 Python example: softmax, likelihood, and
perplexity}\label{python-example-softmax-likelihood-and-perplexity}

\protect\phantomsection\label{ch17-softmax-perplexity}
\begin{Shaded}
\begin{Highlighting}[]
\ImportTok{import}\NormalTok{ numpy }\ImportTok{as}\NormalTok{ np}

\NormalTok{np.set\_printoptions(precision}\OperatorTok{=}\DecValTok{4}\NormalTok{, suppress}\OperatorTok{=}\VariableTok{True}\NormalTok{)}

\KeywordTok{def}\NormalTok{ softmax(logits, temperature}\OperatorTok{=}\FloatTok{1.0}\NormalTok{):}
\NormalTok{    logits }\OperatorTok{=}\NormalTok{ np.asarray(logits, dtype}\OperatorTok{=}\BuiltInTok{float}\NormalTok{) }\OperatorTok{/}\NormalTok{ temperature}
\NormalTok{    shifted }\OperatorTok{=}\NormalTok{ logits }\OperatorTok{{-}}\NormalTok{ np.}\BuiltInTok{max}\NormalTok{(logits)}
\NormalTok{    exp\_scores }\OperatorTok{=}\NormalTok{ np.exp(shifted)}
    \ControlFlowTok{return}\NormalTok{ exp\_scores }\OperatorTok{/}\NormalTok{ exp\_scores.}\BuiltInTok{sum}\NormalTok{()}

\KeywordTok{def}\NormalTok{ sequence\_log\_probability(next\_token\_probabilities):}
\NormalTok{    probs }\OperatorTok{=}\NormalTok{ np.asarray(next\_token\_probabilities, dtype}\OperatorTok{=}\BuiltInTok{float}\NormalTok{)}
    \ControlFlowTok{if}\NormalTok{ np.}\BuiltInTok{any}\NormalTok{(probs }\OperatorTok{\textless{}=} \DecValTok{0}\NormalTok{):}
        \ControlFlowTok{raise} \PreprocessorTok{ValueError}\NormalTok{(}\StringTok{"All observed{-}token probabilities must be positive."}\NormalTok{)}
    \ControlFlowTok{return}\NormalTok{ np.log(probs).}\BuiltInTok{sum}\NormalTok{()}

\KeywordTok{def}\NormalTok{ perplexity(next\_token\_probabilities):}
\NormalTok{    probs }\OperatorTok{=}\NormalTok{ np.asarray(next\_token\_probabilities, dtype}\OperatorTok{=}\BuiltInTok{float}\NormalTok{)}
\NormalTok{    avg\_nll }\OperatorTok{=} \OperatorTok{{-}}\NormalTok{np.log(probs).mean()}
    \ControlFlowTok{return} \BuiltInTok{float}\NormalTok{(np.exp(avg\_nll))}

\NormalTok{observed\_probs }\OperatorTok{=}\NormalTok{ [}\FloatTok{0.40}\NormalTok{, }\FloatTok{0.25}\NormalTok{, }\FloatTok{0.50}\NormalTok{, }\FloatTok{0.20}\NormalTok{]}
\NormalTok{logp }\OperatorTok{=}\NormalTok{ sequence\_log\_probability(observed\_probs)}
\NormalTok{ppl }\OperatorTok{=}\NormalTok{ perplexity(observed\_probs)}

\BuiltInTok{print}\NormalTok{(}\StringTok{"Sequence log{-}probability:"}\NormalTok{, logp)}
\BuiltInTok{print}\NormalTok{(}\StringTok{"Perplexity:"}\NormalTok{, ppl)}

\NormalTok{logits }\OperatorTok{=}\NormalTok{ np.array([}\FloatTok{2.0}\NormalTok{, }\FloatTok{1.2}\NormalTok{, }\FloatTok{0.5}\NormalTok{, }\OperatorTok{{-}}\FloatTok{0.4}\NormalTok{])}
\BuiltInTok{print}\NormalTok{(}\StringTok{"Softmax probabilities:"}\NormalTok{, softmax(logits))}
\BuiltInTok{print}\NormalTok{(}\StringTok{"Temperature 0.5:"}\NormalTok{, softmax(logits, temperature}\OperatorTok{=}\FloatTok{0.5}\NormalTok{))}
\BuiltInTok{print}\NormalTok{(}\StringTok{"Temperature 2.0:"}\NormalTok{, softmax(logits, temperature}\OperatorTok{=}\FloatTok{2.0}\NormalTok{))}
\end{Highlighting}
\end{Shaded}

\begin{verbatim}
Sequence log-probability: -4.605170185988091
Perplexity: 3.162277660168379
Softmax probabilities: [0.5672 0.2548 0.1266 0.0515]
Temperature 0.5: [0.7937 0.1602 0.0395 0.0065]
Temperature 2.0: [0.4092 0.2743 0.1933 0.1232]
\end{verbatim}

\section{17.15 Python example: causal
mask}\label{python-example-causal-mask}

\protect\phantomsection\label{ch17-causal-mask}
\begin{Shaded}
\begin{Highlighting}[]
\NormalTok{T }\OperatorTok{=} \DecValTok{6}
\NormalTok{causal\_mask }\OperatorTok{=}\NormalTok{ np.triu(np.ones((T, T)), k}\OperatorTok{=}\DecValTok{1}\NormalTok{)}
\NormalTok{causal\_mask }\OperatorTok{=}\NormalTok{ np.where(causal\_mask }\OperatorTok{==} \DecValTok{1}\NormalTok{, }\OperatorTok{{-}}\FloatTok{1e9}\NormalTok{, }\FloatTok{0.0}\NormalTok{)}

\BuiltInTok{print}\NormalTok{(causal\_mask)}
\end{Highlighting}
\end{Shaded}

\begin{verbatim}
[[ 0.e+00 -1.e+09 -1.e+09 -1.e+09 -1.e+09 -1.e+09]
 [ 0.e+00  0.e+00 -1.e+09 -1.e+09 -1.e+09 -1.e+09]
 [ 0.e+00  0.e+00  0.e+00 -1.e+09 -1.e+09 -1.e+09]
 [ 0.e+00  0.e+00  0.e+00  0.e+00 -1.e+09 -1.e+09]
 [ 0.e+00  0.e+00  0.e+00  0.e+00  0.e+00 -1.e+09]
 [ 0.e+00  0.e+00  0.e+00  0.e+00  0.e+00  0.e+00]]
\end{verbatim}

The large negative values play the role of \(-\infty\) numerically.
After applying softmax, future positions receive probability
approximately zero.

\section{17.16 Python example: decoding
rules}\label{python-example-decoding-rules}

\protect\phantomsection\label{ch17-decoding-rules}
\begin{Shaded}
\begin{Highlighting}[]
\NormalTok{rng }\OperatorTok{=}\NormalTok{ np.random.default\_rng(}\DecValTok{7243}\NormalTok{)}

\NormalTok{vocab }\OperatorTok{=}\NormalTok{ np.array([}\StringTok{"math"}\NormalTok{, }\StringTok{"models"}\NormalTok{, }\StringTok{"learn"}\NormalTok{, }\StringTok{"data"}\NormalTok{, }\StringTok{"language"}\NormalTok{, }\StringTok{"proof"}\NormalTok{])}
\NormalTok{base\_logits }\OperatorTok{=}\NormalTok{ np.array([}\FloatTok{2.5}\NormalTok{, }\FloatTok{2.0}\NormalTok{, }\FloatTok{1.2}\NormalTok{, }\FloatTok{0.7}\NormalTok{, }\FloatTok{0.3}\NormalTok{, }\OperatorTok{{-}}\FloatTok{0.8}\NormalTok{])}

\KeywordTok{def}\NormalTok{ top\_k\_filter(probs, k):}
\NormalTok{    probs }\OperatorTok{=}\NormalTok{ np.asarray(probs, dtype}\OperatorTok{=}\BuiltInTok{float}\NormalTok{)}
    \ControlFlowTok{if}\NormalTok{ k }\OperatorTok{\textgreater{}=} \BuiltInTok{len}\NormalTok{(probs):}
        \ControlFlowTok{return}\NormalTok{ probs }\OperatorTok{/}\NormalTok{ probs.}\BuiltInTok{sum}\NormalTok{()}
\NormalTok{    keep }\OperatorTok{=}\NormalTok{ np.argsort(probs)[::}\OperatorTok{{-}}\DecValTok{1}\NormalTok{][:k]}
\NormalTok{    q }\OperatorTok{=}\NormalTok{ np.zeros\_like(probs)}
\NormalTok{    q[keep] }\OperatorTok{=}\NormalTok{ probs[keep]}
    \ControlFlowTok{return}\NormalTok{ q }\OperatorTok{/}\NormalTok{ q.}\BuiltInTok{sum}\NormalTok{()}

\KeywordTok{def}\NormalTok{ top\_p\_filter(probs, p\_threshold):}
\NormalTok{    probs }\OperatorTok{=}\NormalTok{ np.asarray(probs, dtype}\OperatorTok{=}\BuiltInTok{float}\NormalTok{)}
\NormalTok{    order }\OperatorTok{=}\NormalTok{ np.argsort(probs)[::}\OperatorTok{{-}}\DecValTok{1}\NormalTok{]}
\NormalTok{    cumulative }\OperatorTok{=}\NormalTok{ np.cumsum(probs[order])}
\NormalTok{    m }\OperatorTok{=} \BuiltInTok{int}\NormalTok{(np.searchsorted(cumulative, p\_threshold, side}\OperatorTok{=}\StringTok{"left"}\NormalTok{)) }\OperatorTok{+} \DecValTok{1}
\NormalTok{    keep }\OperatorTok{=}\NormalTok{ order[:m]}
\NormalTok{    q }\OperatorTok{=}\NormalTok{ np.zeros\_like(probs)}
\NormalTok{    q[keep] }\OperatorTok{=}\NormalTok{ probs[keep]}
    \ControlFlowTok{return}\NormalTok{ q }\OperatorTok{/}\NormalTok{ q.}\BuiltInTok{sum}\NormalTok{()}

\NormalTok{p }\OperatorTok{=}\NormalTok{ softmax(base\_logits, temperature}\OperatorTok{=}\FloatTok{1.0}\NormalTok{)}
\BuiltInTok{print}\NormalTok{(}\StringTok{"Base distribution"}\NormalTok{)}
\ControlFlowTok{for}\NormalTok{ token, prob }\KeywordTok{in} \BuiltInTok{zip}\NormalTok{(vocab, p):}
    \BuiltInTok{print}\NormalTok{(}\SpecialStringTok{f"}\SpecialCharTok{\{}\NormalTok{token}\SpecialCharTok{:\textgreater{}9s\}}\SpecialStringTok{: }\SpecialCharTok{\{}\NormalTok{prob}\SpecialCharTok{:.4f\}}\SpecialStringTok{"}\NormalTok{)}

\BuiltInTok{print}\NormalTok{(}\StringTok{"}\CharTok{\textbackslash{}n}\StringTok{Top{-}k distribution, k=3"}\NormalTok{)}
\NormalTok{qk }\OperatorTok{=}\NormalTok{ top\_k\_filter(p, k}\OperatorTok{=}\DecValTok{3}\NormalTok{)}
\ControlFlowTok{for}\NormalTok{ token, prob }\KeywordTok{in} \BuiltInTok{zip}\NormalTok{(vocab, qk):}
    \BuiltInTok{print}\NormalTok{(}\SpecialStringTok{f"}\SpecialCharTok{\{}\NormalTok{token}\SpecialCharTok{:\textgreater{}9s\}}\SpecialStringTok{: }\SpecialCharTok{\{}\NormalTok{prob}\SpecialCharTok{:.4f\}}\SpecialStringTok{"}\NormalTok{)}

\BuiltInTok{print}\NormalTok{(}\StringTok{"}\CharTok{\textbackslash{}n}\StringTok{Top{-}p distribution, p=0.80"}\NormalTok{)}
\NormalTok{qp }\OperatorTok{=}\NormalTok{ top\_p\_filter(p, p\_threshold}\OperatorTok{=}\FloatTok{0.80}\NormalTok{)}
\ControlFlowTok{for}\NormalTok{ token, prob }\KeywordTok{in} \BuiltInTok{zip}\NormalTok{(vocab, qp):}
    \BuiltInTok{print}\NormalTok{(}\SpecialStringTok{f"}\SpecialCharTok{\{}\NormalTok{token}\SpecialCharTok{:\textgreater{}9s\}}\SpecialStringTok{: }\SpecialCharTok{\{}\NormalTok{prob}\SpecialCharTok{:.4f\}}\SpecialStringTok{"}\NormalTok{)}

\NormalTok{sampled\_token }\OperatorTok{=}\NormalTok{ rng.choice(vocab, p}\OperatorTok{=}\NormalTok{qp)}
\BuiltInTok{print}\NormalTok{(}\StringTok{"}\CharTok{\textbackslash{}n}\StringTok{Sample from top{-}p distribution:"}\NormalTok{, sampled\_token)}
\end{Highlighting}
\end{Shaded}

\begin{verbatim}
Base distribution
     math: 0.4562
   models: 0.2767
    learn: 0.1243
     data: 0.0754
 language: 0.0505
    proof: 0.0168

Top-k distribution, k=3
     math: 0.5322
   models: 0.3228
    learn: 0.1450
     data: 0.0000
 language: 0.0000
    proof: 0.0000

Top-p distribution, p=0.80
     math: 0.5322
   models: 0.3228
    learn: 0.1450
     data: 0.0000
 language: 0.0000
    proof: 0.0000

Sample from top-p distribution: models
\end{verbatim}

\section{17.17 Python example: a tiny autoregressive
generator}\label{python-example-a-tiny-autoregressive-generator}

The following example is intentionally simple. It uses a hand-built
bigram-like table so that the mathematics of generation is visible.

\protect\phantomsection\label{ch17-tiny-generator}
\begin{Shaded}
\begin{Highlighting}[]
\NormalTok{rng }\OperatorTok{=}\NormalTok{ np.random.default\_rng(}\DecValTok{2026}\NormalTok{)}

\NormalTok{vocab }\OperatorTok{=}\NormalTok{ [}\StringTok{"\textless{}bos\textgreater{}"}\NormalTok{, }\StringTok{"math"}\NormalTok{, }\StringTok{"models"}\NormalTok{, }\StringTok{"learn"}\NormalTok{, }\StringTok{"from"}\NormalTok{, }\StringTok{"data"}\NormalTok{, }\StringTok{"language"}\NormalTok{, }\StringTok{"\textless{}eos\textgreater{}"}\NormalTok{]}
\NormalTok{idx }\OperatorTok{=}\NormalTok{ \{w: i }\ControlFlowTok{for}\NormalTok{ i, w }\KeywordTok{in} \BuiltInTok{enumerate}\NormalTok{(vocab)\}}

\CommentTok{\# Rows are contexts, columns are next{-}token logits.}
\NormalTok{Z }\OperatorTok{=}\NormalTok{ np.full((}\BuiltInTok{len}\NormalTok{(vocab), }\BuiltInTok{len}\NormalTok{(vocab)), }\OperatorTok{{-}}\FloatTok{3.0}\NormalTok{)}

\KeywordTok{def}\NormalTok{ set\_logits(context, scores):}
\NormalTok{    row }\OperatorTok{=}\NormalTok{ idx[context]}
    \ControlFlowTok{for}\NormalTok{ token, score }\KeywordTok{in}\NormalTok{ scores.items():}
\NormalTok{        Z[row, idx[token]] }\OperatorTok{=}\NormalTok{ score}

\NormalTok{set\_logits(}\StringTok{"\textless{}bos\textgreater{}"}\NormalTok{, \{}\StringTok{"math"}\NormalTok{: }\FloatTok{2.0}\NormalTok{, }\StringTok{"language"}\NormalTok{: }\FloatTok{1.8}\NormalTok{, }\StringTok{"models"}\NormalTok{: }\FloatTok{0.5}\NormalTok{\})}
\NormalTok{set\_logits(}\StringTok{"math"}\NormalTok{, \{}\StringTok{"models"}\NormalTok{: }\FloatTok{2.5}\NormalTok{, }\StringTok{"learn"}\NormalTok{: }\FloatTok{1.0}\NormalTok{, }\StringTok{"data"}\NormalTok{: }\FloatTok{0.2}\NormalTok{\})}
\NormalTok{set\_logits(}\StringTok{"language"}\NormalTok{, \{}\StringTok{"models"}\NormalTok{: }\FloatTok{2.4}\NormalTok{, }\StringTok{"data"}\NormalTok{: }\FloatTok{0.4}\NormalTok{\})}
\NormalTok{set\_logits(}\StringTok{"models"}\NormalTok{, \{}\StringTok{"learn"}\NormalTok{: }\FloatTok{2.2}\NormalTok{, }\StringTok{"from"}\NormalTok{: }\FloatTok{1.4}\NormalTok{, }\StringTok{"\textless{}eos\textgreater{}"}\NormalTok{: }\FloatTok{0.1}\NormalTok{\})}
\NormalTok{set\_logits(}\StringTok{"learn"}\NormalTok{, \{}\StringTok{"from"}\NormalTok{: }\FloatTok{2.3}\NormalTok{, }\StringTok{"data"}\NormalTok{: }\FloatTok{1.0}\NormalTok{\})}
\NormalTok{set\_logits(}\StringTok{"from"}\NormalTok{, \{}\StringTok{"data"}\NormalTok{: }\FloatTok{2.6}\NormalTok{, }\StringTok{"language"}\NormalTok{: }\FloatTok{0.4}\NormalTok{\})}
\NormalTok{set\_logits(}\StringTok{"data"}\NormalTok{, \{}\StringTok{"\textless{}eos\textgreater{}"}\NormalTok{: }\FloatTok{2.2}\NormalTok{, }\StringTok{"models"}\NormalTok{: }\FloatTok{0.4}\NormalTok{\})}
\NormalTok{set\_logits(}\StringTok{"\textless{}eos\textgreater{}"}\NormalTok{, \{}\StringTok{"\textless{}eos\textgreater{}"}\NormalTok{: }\FloatTok{3.0}\NormalTok{\})}

\KeywordTok{def}\NormalTok{ next\_distribution(context\_token, temperature}\OperatorTok{=}\FloatTok{1.0}\NormalTok{, top\_k}\OperatorTok{=}\VariableTok{None}\NormalTok{, top\_p}\OperatorTok{=}\VariableTok{None}\NormalTok{):}
\NormalTok{    probs }\OperatorTok{=}\NormalTok{ softmax(Z[idx[context\_token]], temperature}\OperatorTok{=}\NormalTok{temperature)}
    \ControlFlowTok{if}\NormalTok{ top\_k }\KeywordTok{is} \KeywordTok{not} \VariableTok{None}\NormalTok{:}
\NormalTok{        probs }\OperatorTok{=}\NormalTok{ top\_k\_filter(probs, top\_k)}
    \ControlFlowTok{if}\NormalTok{ top\_p }\KeywordTok{is} \KeywordTok{not} \VariableTok{None}\NormalTok{:}
\NormalTok{        probs }\OperatorTok{=}\NormalTok{ top\_p\_filter(probs, top\_p)}
    \ControlFlowTok{return}\NormalTok{ probs}

\KeywordTok{def}\NormalTok{ generate(max\_steps}\OperatorTok{=}\DecValTok{10}\NormalTok{, temperature}\OperatorTok{=}\FloatTok{1.0}\NormalTok{, top\_k}\OperatorTok{=}\VariableTok{None}\NormalTok{, top\_p}\OperatorTok{=}\VariableTok{None}\NormalTok{):}
\NormalTok{    tokens }\OperatorTok{=}\NormalTok{ [}\StringTok{"\textless{}bos\textgreater{}"}\NormalTok{]}
    \ControlFlowTok{for}\NormalTok{ \_ }\KeywordTok{in} \BuiltInTok{range}\NormalTok{(max\_steps):}
\NormalTok{        probs }\OperatorTok{=}\NormalTok{ next\_distribution(tokens[}\OperatorTok{{-}}\DecValTok{1}\NormalTok{], temperature}\OperatorTok{=}\NormalTok{temperature, top\_k}\OperatorTok{=}\NormalTok{top\_k, top\_p}\OperatorTok{=}\NormalTok{top\_p)}
\NormalTok{        next\_token }\OperatorTok{=}\NormalTok{ rng.choice(vocab, p}\OperatorTok{=}\NormalTok{probs)}
\NormalTok{        tokens.append(next\_token)}
        \ControlFlowTok{if}\NormalTok{ next\_token }\OperatorTok{==} \StringTok{"\textless{}eos\textgreater{}"}\NormalTok{:}
            \ControlFlowTok{break}
    \ControlFlowTok{return}\NormalTok{ tokens}

\ControlFlowTok{for}\NormalTok{ temp }\KeywordTok{in}\NormalTok{ [}\FloatTok{0.5}\NormalTok{, }\FloatTok{1.0}\NormalTok{, }\FloatTok{1.8}\NormalTok{]:}
    \BuiltInTok{print}\NormalTok{(}\SpecialStringTok{f"temperature=}\SpecialCharTok{\{}\NormalTok{temp}\SpecialCharTok{\}}\SpecialStringTok{:"}\NormalTok{, }\StringTok{" "}\NormalTok{.join(generate(temperature}\OperatorTok{=}\NormalTok{temp, top\_p}\OperatorTok{=}\FloatTok{0.90}\NormalTok{)))}
\end{Highlighting}
\end{Shaded}

\begin{verbatim}
temperature=0.5: <bos> math models learn from data <eos>
temperature=1.0: <bos> language models learn from language data <eos>
temperature=1.8: <bos> language models language models learn from data <eos>
\end{verbatim}

\section{17.18 Python example: beam
search}\label{python-example-beam-search}

\protect\phantomsection\label{ch17-beam-search}
\begin{Shaded}
\begin{Highlighting}[]
\ImportTok{import}\NormalTok{ math}

\KeywordTok{def}\NormalTok{ beam\_search(beam\_width}\OperatorTok{=}\DecValTok{3}\NormalTok{, max\_steps}\OperatorTok{=}\DecValTok{6}\NormalTok{, alpha}\OperatorTok{=}\FloatTok{0.0}\NormalTok{):}
\NormalTok{    beam }\OperatorTok{=}\NormalTok{ [([}\StringTok{"\textless{}bos\textgreater{}"}\NormalTok{], }\FloatTok{0.0}\NormalTok{)]}
\NormalTok{    completed }\OperatorTok{=}\NormalTok{ []}

    \ControlFlowTok{for}\NormalTok{ \_ }\KeywordTok{in} \BuiltInTok{range}\NormalTok{(max\_steps):}
\NormalTok{        candidates }\OperatorTok{=}\NormalTok{ []}
        \ControlFlowTok{for}\NormalTok{ tokens, score }\KeywordTok{in}\NormalTok{ beam:}
            \ControlFlowTok{if}\NormalTok{ tokens[}\OperatorTok{{-}}\DecValTok{1}\NormalTok{] }\OperatorTok{==} \StringTok{"\textless{}eos\textgreater{}"}\NormalTok{:}
\NormalTok{                completed.append((tokens, score))}
\NormalTok{                candidates.append((tokens, score))}
                \ControlFlowTok{continue}
\NormalTok{            probs }\OperatorTok{=}\NormalTok{ next\_distribution(tokens[}\OperatorTok{{-}}\DecValTok{1}\NormalTok{], temperature}\OperatorTok{=}\FloatTok{1.0}\NormalTok{)}
            \ControlFlowTok{for}\NormalTok{ j, prob }\KeywordTok{in} \BuiltInTok{enumerate}\NormalTok{(probs):}
                \ControlFlowTok{if}\NormalTok{ prob }\OperatorTok{\textgreater{}} \DecValTok{0}\NormalTok{:}
\NormalTok{                    new\_tokens }\OperatorTok{=}\NormalTok{ tokens }\OperatorTok{+}\NormalTok{ [vocab[j]]}
\NormalTok{                    new\_score }\OperatorTok{=}\NormalTok{ score }\OperatorTok{+}\NormalTok{ math.log(prob)}
\NormalTok{                    candidates.append((new\_tokens, new\_score))}

        \KeywordTok{def}\NormalTok{ normalized\_score(item):}
\NormalTok{            tokens, score }\OperatorTok{=}\NormalTok{ item}
\NormalTok{            length }\OperatorTok{=} \BuiltInTok{max}\NormalTok{(}\BuiltInTok{len}\NormalTok{(tokens) }\OperatorTok{{-}} \DecValTok{1}\NormalTok{, }\DecValTok{1}\NormalTok{)}
            \ControlFlowTok{return}\NormalTok{ score }\OperatorTok{/}\NormalTok{ (length }\OperatorTok{**}\NormalTok{ alpha)}

\NormalTok{        candidates.sort(key}\OperatorTok{=}\NormalTok{normalized\_score, reverse}\OperatorTok{=}\VariableTok{True}\NormalTok{)}
\NormalTok{        beam }\OperatorTok{=}\NormalTok{ candidates[:beam\_width]}

\NormalTok{    completed.extend(beam)}
\NormalTok{    completed.sort(key}\OperatorTok{=}\KeywordTok{lambda}\NormalTok{ item: item[}\DecValTok{1}\NormalTok{] }\OperatorTok{/}\NormalTok{ (}\BuiltInTok{max}\NormalTok{(}\BuiltInTok{len}\NormalTok{(item[}\DecValTok{0}\NormalTok{]) }\OperatorTok{{-}} \DecValTok{1}\NormalTok{, }\DecValTok{1}\NormalTok{) }\OperatorTok{**}\NormalTok{ alpha), reverse}\OperatorTok{=}\VariableTok{True}\NormalTok{)}
    \ControlFlowTok{return}\NormalTok{ completed[:beam\_width]}

\ControlFlowTok{for}\NormalTok{ tokens, score }\KeywordTok{in}\NormalTok{ beam\_search(beam\_width}\OperatorTok{=}\DecValTok{3}\NormalTok{, max\_steps}\OperatorTok{=}\DecValTok{6}\NormalTok{, alpha}\OperatorTok{=}\FloatTok{0.6}\NormalTok{):}
    \BuiltInTok{print}\NormalTok{(}\SpecialStringTok{f"score=}\SpecialCharTok{\{}\NormalTok{score}\SpecialCharTok{: .4f\}}\SpecialStringTok{ |"}\NormalTok{, }\StringTok{" "}\NormalTok{.join(tokens))}
\end{Highlighting}
\end{Shaded}

\begin{verbatim}
score=-2.0655 | <bos> math models learn from data <eos>
score=-2.1205 | <bos> language models learn from data <eos>
score=-2.6012 | <bos> math models from data <eos>
\end{verbatim}

\section{17.19 Interactive: next-token decoding
explorer}\label{interactive-next-token-decoding-explorer}

This interactive illustrates how logits become a next-token distribution
and how decoding choices change that distribution. The example is
intentionally small so that the probability simplex is visible.

\section{17.20 Python Plotly version}\label{python-plotly-version}

The inline HTML interactive above is designed for the Quarto webpage.
The following Python version is useful in notebooks.

\protect\phantomsection\label{ch17-python-plotly-version}
\begin{Shaded}
\begin{Highlighting}[]
\ImportTok{import}\NormalTok{ plotly.graph\_objects }\ImportTok{as}\NormalTok{ go}

\NormalTok{interactive\_tokens }\OperatorTok{=}\NormalTok{ [}\StringTok{"math"}\NormalTok{, }\StringTok{"models"}\NormalTok{, }\StringTok{"learn"}\NormalTok{, }\StringTok{"data"}\NormalTok{, }\StringTok{"language"}\NormalTok{, }\StringTok{"proof"}\NormalTok{, }\StringTok{"random"}\NormalTok{, }\StringTok{"graph"}\NormalTok{]}
\NormalTok{interactive\_logits }\OperatorTok{=}\NormalTok{ np.array([}\FloatTok{2.7}\NormalTok{, }\FloatTok{2.2}\NormalTok{, }\FloatTok{1.4}\NormalTok{, }\FloatTok{0.9}\NormalTok{, }\FloatTok{0.2}\NormalTok{, }\OperatorTok{{-}}\FloatTok{0.2}\NormalTok{, }\OperatorTok{{-}}\FloatTok{0.8}\NormalTok{, }\OperatorTok{{-}}\FloatTok{1.4}\NormalTok{])}
\NormalTok{interactive\_probs }\OperatorTok{=}\NormalTok{ softmax(interactive\_logits, temperature}\OperatorTok{=}\FloatTok{1.0}\NormalTok{)}

\NormalTok{fig }\OperatorTok{=}\NormalTok{ go.Figure(}
\NormalTok{    data}\OperatorTok{=}\NormalTok{[go.Bar(x}\OperatorTok{=}\NormalTok{interactive\_tokens, y}\OperatorTok{=}\NormalTok{interactive\_probs)],}
\NormalTok{    layout}\OperatorTok{=}\NormalTok{go.Layout(}
\NormalTok{        title}\OperatorTok{=}\StringTok{"Next{-}token distribution from logits"}\NormalTok{,}
\NormalTok{        xaxis\_title}\OperatorTok{=}\StringTok{"token"}\NormalTok{,}
\NormalTok{        yaxis\_title}\OperatorTok{=}\StringTok{"probability"}
\NormalTok{    )}
\NormalTok{)}
\NormalTok{fig.show()}
\end{Highlighting}
\end{Shaded}

\protect\phantomsection\label{ch17-python-plotly-version-1}
\begin{verbatim}
Unable to display output for mime type(s): text/html
\end{verbatim}

\protect\phantomsection\label{ch17-python-plotly-version-2}
\begin{verbatim}
Unable to display output for mime type(s): text/html
\end{verbatim}

\section{17.21 AI-assisted learning
activities}\label{ai-assisted-learning-activities-2}

Use an AI assistant as a mathematical tutor, not as a black box. Good
prompts force the assistant to state the conditioning, the probability
distribution, and the loss function.

\begin{tcolorbox}[enhanced jigsaw, arc=.35mm, bottomrule=.15mm, bottomtitle=1mm, breakable, colback=white, colbacktitle=quarto-callout-tip-color!10!white, colframe=quarto-callout-tip-color-frame, coltitle=black, left=2mm, leftrule=.75mm, opacityback=0, opacitybacktitle=0.6, rightrule=.15mm, title=\textcolor{quarto-callout-tip-color}{\faLightbulb}\hspace{0.5em}{Prompt 1: identify the distribution}, titlerule=0mm, toprule=.15mm, toptitle=1mm]

Ask:

\begin{quote}
In the sentence \texttt{language\ models\ predict\ tokens}, list the
autoregressive training pairs. For each pair, identify the prefix, the
target token, and the conditional probability that appears in the
likelihood.
\end{quote}

Then check whether the answer correctly distinguishes
\(p(w_t\mid w_{<t})\) from \(p(w_{<t}\mid w_t)\).

\end{tcolorbox}

\begin{tcolorbox}[enhanced jigsaw, arc=.35mm, bottomrule=.15mm, bottomtitle=1mm, breakable, colback=white, colbacktitle=quarto-callout-tip-color!10!white, colframe=quarto-callout-tip-color-frame, coltitle=black, left=2mm, leftrule=.75mm, opacityback=0, opacitybacktitle=0.6, rightrule=.15mm, title=\textcolor{quarto-callout-tip-color}{\faLightbulb}\hspace{0.5em}{Prompt 2: decoding audit}, titlerule=0mm, toprule=.15mm, toptitle=1mm]

Ask:

\begin{quote}
Explain the difference between temperature, top-k sampling, and nucleus
sampling using probability notation. Give one example where top-k and
top-p keep different numbers of tokens.
\end{quote}

Then verify the explanation by computing the filtered distributions
yourself.

\end{tcolorbox}

\begin{tcolorbox}[enhanced jigsaw, arc=.35mm, bottomrule=.15mm, bottomtitle=1mm, breakable, colback=white, colbacktitle=quarto-callout-tip-color!10!white, colframe=quarto-callout-tip-color-frame, coltitle=black, left=2mm, leftrule=.75mm, opacityback=0, opacitybacktitle=0.6, rightrule=.15mm, title=\textcolor{quarto-callout-tip-color}{\faLightbulb}\hspace{0.5em}{Prompt 3: perplexity interpretation}, titlerule=0mm, toprule=.15mm, toptitle=1mm]

Ask:

\begin{quote}
A language model has average negative log-likelihood 1.7 nats per token.
What is its perplexity? What does this number mean, and what does it not
mean?
\end{quote}

A good answer should mention \(\exp(1.7)\) and should not claim that
perplexity measures factual correctness.

\end{tcolorbox}

\section{17.22 Common mistakes}\label{common-mistakes-6}

\begin{enumerate}
\def\labelenumi{\arabic{enumi}.}
\item
  \textbf{Confusing the chain rule with a Markov assumption.} The
  factorization \(P(w_{1:T})=\prod_t P(w_t\mid w_{<t})\) is exact. A
  Markov model makes the additional approximation
  \(P(w_t\mid w_{<t})\approx P(w_t\mid w_{t-r:t-1})\).
\item
  \textbf{Treating logits as probabilities.} Logits can be any real
  numbers. Probabilities appear after softmax.
\item
  \textbf{Ignoring length in sequence scores.} The raw log-probability
  of a sequence usually decreases as length increases. Comparing
  sequences of different lengths requires care.
\item
  \textbf{Believing sampling is the same as training.} Training usually
  uses teacher forcing. Generation conditions on previously generated
  tokens.
\item
  \textbf{Thinking low perplexity guarantees high-quality answers.}
  Perplexity is a likelihood metric. It does not by itself measure
  truthfulness, fairness, safety, or usefulness.
\end{enumerate}

\section{17.23 Exercises}\label{exercises-16}

\subsection{Conceptual exercises}\label{conceptual-exercises-7}

\begin{enumerate}
\def\labelenumi{\arabic{enumi}.}
\item
  Let \(V=\{a,b,c\}\). Write the full chain-rule factorization of
  \(P(w_1,w_2,w_3,w_4)\).
\item
  Explain why the autoregressive factorization is exact but an
  \(n\)-gram approximation is not exact in general.
\item
  Suppose two models assign the following probabilities to the observed
  next tokens:

  \[
  (0.5,0.4,0.25,0.5)
  \quad\text{and}\quad
  (0.8,0.1,0.8,0.1).
  \]

  Compute the average negative log-likelihood and perplexity for each.
  Which model has better likelihood?
\item
  Show that temperature scaling does not change the identity of the
  largest-probability token when there is a unique largest logit.
\item
  Give an example where top-\(k\) sampling with \(k=3\) keeps less total
  probability mass than top-\(p\) sampling with \(p_0=0.90\).
\end{enumerate}

\subsection{Mathematical exercises}\label{mathematical-exercises-6}

\begin{enumerate}
\def\labelenumi{\arabic{enumi}.}
\setcounter{enumi}{5}
\item
  Derive the gradient

  \[
  \nabla_z[-\log \operatorname{softmax}(z)_y]
  =
  \operatorname{softmax}(z)-e_y.
  \]
\item
  Let \(p(\tau)=\operatorname{softmax}(z/\tau)\). Compute
  \(\lim_{\tau\to\infty}p(\tau)\).
\item
  Let \(M\) be the causal mask with \(M_{ij}=-\infty\) for \(j>i\) and
  \(0\) otherwise. Prove that the row-wise softmax of \(S+M\) assigns
  zero weight to future positions.
\item
  Suppose a model generates tokens by
  \(W_t\sim p_\theta(\cdot\mid W_{<t})\). Explain why this defines a
  stochastic process on \(V^*\).
\item
  In beam search, explain why raw log-probability favors shorter
  sequences when end-of-sequence tokens are allowed.
\end{enumerate}

\subsection{Computational exercises}\label{computational-exercises-4}

\begin{enumerate}
\def\labelenumi{\arabic{enumi}.}
\setcounter{enumi}{10}
\item
  Implement top-\(k\) and top-\(p\) filtering for a probability vector.
  Test your code on at least three distributions: sharp, medium, and
  flat.
\item
  Write a function that takes a list of observed next-token
  probabilities and returns log-probability, average NLL, and
  perplexity.
\item
  Build a tiny bigram language model from a small corpus. Generate text
  using greedy decoding, temperature sampling, top-\(k\) sampling, and
  top-\(p\) sampling.
\item
  Implement beam search for your tiny bigram model. Compare results for
  beam widths \(1\), \(3\), and \(5\).
\item
  Simulate generation under different temperatures and estimate the
  empirical entropy of generated next tokens.
\end{enumerate}

\section{17.24 Mini-project: a mathematical decoding
report}\label{mini-project-a-mathematical-decoding-report}

Choose a small corpus related to a course topic, for example mathematics
course descriptions, research abstracts, or advising FAQ text. Build a
simple autoregressive model using bigram or trigram probabilities. Then
write a short mathematical report containing:

\begin{enumerate}
\def\labelenumi{\arabic{enumi}.}
\tightlist
\item
  the vocabulary construction,
\item
  the estimated conditional probabilities,
\item
  the training negative log-likelihood,
\item
  generated examples under at least four decoding rules,
\item
  a comparison of entropy, diversity, repetition, and coherence,
\item
  a discussion of how these ideas scale to Transformer language models.
\end{enumerate}

The goal is not to build a strong model. The goal is to make the
probability structure of generation completely explicit.

\section{17.25 Python lab connection}\label{python-lab-connection-8}

The associated lab for this chapter is planned as:

\begin{Shaded}
\begin{Highlighting}[]
\NormalTok{labs/chapter{-}17{-}autoregressive{-}language{-}models.ipynb}
\end{Highlighting}
\end{Shaded}

The lab should be written for independent study. It should include
background on categorical distributions, log-probabilities, softmax,
perplexity, decoding rules, and a small from-scratch text generator.

\section{What to remember}\label{what-to-remember-5}

\begin{itemize}
\tightlist
\item
  An autoregressive language model estimates
  \(p_\theta(w_t\mid w_{<t})\).
\item
  The sequence probability is the product of next-token conditional
  probabilities.
\item
  Training usually minimizes next-token negative log-likelihood.
\item
  Perplexity is the exponential of average negative log-likelihood.
\item
  Generation is a recursive process in which sampled tokens become
  future context.
\item
  Temperature, top-\(k\), and top-\(p\) modify the next-token
  distribution before sampling.
\item
  Causal masking prevents a decoder-only model from seeing future
  tokens.
\end{itemize}

\section{Notes and references}\label{notes-and-references-9}

This chapter connects the probabilistic language-modeling ideas from
Chapters 4 and 5 with the Transformer machinery developed in Chapters
13--16. Classical background includes statistical language modeling and
information theory. Modern autoregressive models build on neural
sequence modeling, attention, and large-scale maximum likelihood
training.

\chapter{Training Large Language
Models}\label{training-large-language-models}

Pretraining, fine-tuning, instruction tuning, and alignment as
optimization problems

\hfill\break

Chapter 17 described an autoregressive language model as a conditional
probability model for the next token. This chapter explains how such a
model is trained. The main mathematical idea is simple: training a large
language model is empirical risk minimization over token sequences.
Pretraining, supervised fine-tuning, instruction tuning, preference
learning, and alignment differ mainly in the data distribution, the loss
function, and the constraints placed on the updated model.

\section{Learning goals}\label{learning-goals-17}

After studying this chapter, you should be able to:

\begin{itemize}
\tightlist
\item
  Write language-model pretraining as maximum likelihood estimation.
\item
  Derive the token-level cross-entropy loss from the autoregressive
  factorization.
\item
  Explain teacher forcing, causal masking, and padding masks in the
  training objective.
\item
  Compute the gradient of the softmax cross-entropy loss with respect to
  logits.
\item
  Distinguish pretraining, supervised fine-tuning, instruction tuning,
  and preference tuning.
\item
  Explain KL-regularized fine-tuning as constrained distributional
  movement away from a reference model.
\item
  Describe pairwise preference modeling using the Bradley--Terry
  logistic model.
\item
  Interpret RLHF and DPO-style methods as optimization problems on
  conditional distributions.
\item
  Implement small pieces of the training objective from scratch in
  Python.
\end{itemize}

\section{18.1 From a language model to a statistical
estimator}\label{from-a-language-model-to-a-statistical-estimator}

Let \(V\) be a finite vocabulary and let a token sequence be

\[
w_{1:T}=(w_1,\ldots,w_T) \in V^T.
\]

An autoregressive language model with parameters \(\theta\) defines

\[
p_\theta(w_{1:T})
=
\prod_{t=1}^T p_\theta(w_t \mid w_{<t}).
\]

Training means choosing \(\theta\) using data. Suppose the training
corpus consists of \(N\) token sequences

\[
D=\{w^{(1)}_{1:T_1},\ldots,w^{(N)}_{1:T_N}\}.
\]

The maximum likelihood estimator is

\[
\hat \theta
=
\arg\max_\theta
\prod_{i=1}^N
p_\theta(w^{(i)}_{1:T_i}).
\]

Because products of many small probabilities are numerically unstable,
we maximize log-likelihood or equivalently minimize negative
log-likelihood:

\[
\hat\theta
=
\arg\min_\theta
\left[-\sum_{i=1}^N \log p_\theta(w^{(i)}_{1:T_i})\right].
\]

Using the autoregressive factorization,

\[
-\log p_\theta(w^{(i)}_{1:T_i})
=
-\sum_{t=1}^{T_i}
\log p_\theta(w_t^{(i)} \mid w_{<t}^{(i)}).
\]

Thus the training problem is a large supervised learning problem in
disguise. Each prefix \(w_{<t}\) is an input and the next token \(w_t\)
is the label.

\begin{tcolorbox}[enhanced jigsaw, arc=.35mm, bottomrule=.15mm, bottomtitle=1mm, breakable, colback=white, colbacktitle=quarto-callout-note-color!10!white, colframe=quarto-callout-note-color-frame, coltitle=black, left=2mm, leftrule=.75mm, opacityback=0, opacitybacktitle=0.6, rightrule=.15mm, title=\textcolor{quarto-callout-note-color}{\faInfo}\hspace{0.5em}{Pretraining objective}, titlerule=0mm, toprule=.15mm, toptitle=1mm]

The standard autoregressive pretraining loss is

\[
L_{\text{pre}}(\theta)
=
-\frac{1}{M}
\sum_{i=1}^N\sum_{t=1}^{T_i}
\log p_\theta(w_t^{(i)} \mid w_{<t}^{(i)}),
\]

where \(M=\sum_i T_i\) is the number of predicted tokens. This is the
empirical cross-entropy between the data distribution and the model
distribution.

\end{tcolorbox}

\section{18.2 Training data as samples from an unknown
distribution}\label{training-data-as-samples-from-an-unknown-distribution}

A mathematically clean way to view the corpus is as a sample from an
unknown population distribution \(P_{\text{data}}\) on token sequences.
The population risk for language modeling is

\[
R(\theta)
=
E_{W_{1:T}\sim P_{\text{data}}}
\left[
-\sum_{t=1}^T \log p_\theta(W_t \mid W_{<t})
\right].
\]

The empirical training loss replaces the expectation by an average over
the observed corpus. This gives the usual statistical learning picture:

\[
\text{population risk} \quad \longrightarrow \quad \text{empirical risk} \quad \longrightarrow \quad \text{optimization algorithm}.
\]

For MS Statistics students, this viewpoint is important because the
training corpus is not neutral. It induces an empirical distribution.
The model learns from that distribution and inherits its coverage,
imbalances, style, and noise.

\section{18.3 Tokens, contexts, and
labels}\label{tokens-contexts-and-labels}

Given a sequence

\[
w_{1:T}=(w_1,w_2,\ldots,w_T),
\]

training forms the pairs

\[
(w_{<1},w_1),\;(w_{<2},w_2),\;\ldots,\;(w_{<T},w_T).
\]

The first context \(w_{<1}\) is usually represented by a
beginning-of-sequence token. The target at position \(t\) is not the
entire future. It is only the next token \(w_t\). However, minimizing
next-token loss over many positions encourages the model to capture
long-range structure because the prefix may contain earlier semantic and
syntactic information.

For a vocabulary ordered as \(V=\{v_1,\ldots,v_d\}\), the model produces
logits

\[
z_\theta(w_{<t})\in R^d.
\]

The probability assigned to token \(v_j\) is

\[
p_\theta(v_j\mid w_{<t})
=
\frac{\exp(z_j)}{\sum_{k=1}^d \exp(z_k)}.
\]

If the true next token is \(v_y\), the token loss is

\[
\ell(z,y)
=
-\log \frac{\exp(z_y)}{\sum_{k=1}^d \exp(z_k)}
=
-z_y+\log\sum_{k=1}^d \exp(z_k).
\]

This is the same softmax cross-entropy loss used in multiclass
classification, applied at every token position.

\section{18.4 Gradient of token-level
cross-entropy}\label{gradient-of-token-level-cross-entropy}

Let

\[
p=\operatorname{softmax}(z).
\]

Let \(e_y\in R^d\) be the one-hot vector for the correct class. The
token loss is

\[
\ell(z,y)=-\log p_y.
\]

A direct calculation gives

\[
\frac{\partial \ell}{\partial z_j}=p_j-(e_y)_j.
\]

Therefore,

\[
\nabla_z \ell = p-e_y.
\]

This simple formula is central. The model is pushed to decrease
probabilities assigned to incorrect tokens and increase the probability
assigned to the observed token.

\textbf{Interpretation.} If the model assigns too little probability to
the observed next token, the residual \(p-e_y\) is large in the
direction that corrects the logits. Training a language model is
therefore repeated probability calibration over a huge number of
next-token classification problems.

\section{18.5 Teacher forcing}\label{teacher-forcing-1}

During training, the model is given the true previous tokens as context.
This is called \textbf{teacher forcing}. For example, if the training
sequence is

\begin{Shaded}
\begin{Highlighting}[]
\NormalTok{\textless{}bos\textgreater{} math models learn from data \textless{}eos\textgreater{}}
\end{Highlighting}
\end{Shaded}

then the model is trained on contexts such as

\begin{Shaded}
\begin{Highlighting}[]
\NormalTok{\textless{}bos\textgreater{} math models}
\end{Highlighting}
\end{Shaded}

with target

\begin{Shaded}
\begin{Highlighting}[]
\NormalTok{learn}
\end{Highlighting}
\end{Shaded}

The model does not train on its own generated mistakes at this stage.
Mathematically, the loss uses contexts from the empirical data
distribution, not contexts produced by the model. This difference is one
source of the gap between training and generation.

\section{18.6 Batching, padding, and
masks}\label{batching-padding-and-masks}

In practice, sequences of different lengths are placed into a
rectangular array. Suppose a batch has \(B\) sequences and maximum
length \(T_{\max}\). The token array has shape

\[
B\times T_{\max}.
\]

Shorter sequences are padded. Let

\[
m_{i,t}\in \{0,1\}
\]

be a mask indicating whether token position \((i,t)\) is real or
padding. The masked training loss is

\[
L(\theta)
=
-
\frac{1}{\sum_{i,t}m_{i,t}}
\sum_{i=1}^B\sum_{t=1}^{T_{\max}}
 m_{i,t}
\log p_\theta(w_{i,t}\mid w_{i,<t}).
\]

Padding positions must not contribute to the loss. In decoder-only
Transformers, causal masks are also used so that position \(t\) cannot
attend to positions greater than \(t\).

\section{18.7 Optimization: stochastic gradients and
AdamW}\label{optimization-stochastic-gradients-and-adamw}

The loss function for a large language model is nonconvex and
high-dimensional. Let \(\theta\in R^P\) denote all parameters. A generic
stochastic-gradient update is

\[
\theta_{s+1}
=
\theta_s-\eta_s \widehat{\nabla L}(\theta_s),
\]

where \(\widehat{\nabla L}\) is a gradient estimate computed from a
minibatch and \(\eta_s\) is a learning rate.

Modern training often uses adaptive optimizers. A simplified Adam-style
update maintains moving averages of gradients and squared gradients:

\[
m_s=\beta_1m_{s-1}+(1-\beta_1)g_s,
\]

\[
v_s=\beta_2v_{s-1}+(1-\beta_2)g_s\odot g_s.
\]

Then the parameter update has the coordinatewise form

\[
\theta_{s+1}
=
\theta_s
-
\eta_s
\frac{\widehat m_s}{\sqrt{\widehat v_s}+\epsilon}.
\]

AdamW separates weight decay from the gradient update. A schematic
update is

\[
\theta_{s+1}
=
(1-\eta_s\lambda)\theta_s
-
\eta_s
\frac{\widehat m_s}{\sqrt{\widehat v_s}+\epsilon}.
\]

Here \(\lambda\) controls weight decay. From an applied mathematics
viewpoint, AdamW is a stochastic, adaptive, preconditioned first-order
method with explicit shrinkage.

\section{18.8 Validation loss, overfitting, and early
stopping}\label{validation-loss-overfitting-and-early-stopping}

The training loss estimates performance on the training distribution. To
estimate generalization, we hold out validation data and compute

\[
L_{\text{val}}(\theta)
=
-
\frac{1}{M_{\text{val}}}
\sum_{(i,t)\in \text{validation}}
\log p_\theta(w_t^{(i)}\mid w_{<t}^{(i)}).
\]

A useful diagnostic is the training-validation gap:

\[
L_{\text{val}}(\theta)-L_{\text{train}}(\theta).
\]

A large positive gap suggests overfitting, distribution mismatch, or
both. In very large pretraining, the situation is more subtle because
data diversity, model capacity, optimization budget, and evaluation
distribution interact.

\section{18.9 Supervised fine-tuning}\label{supervised-fine-tuning}

Pretraining learns a broad next-token model. \textbf{Supervised
fine-tuning} updates the model on a more targeted dataset. For
instruction tuning, each training example often has a prompt and a
desired response:

\[
x_i=\text{prompt},\qquad y_i=\text{response}.
\]

The conditional fine-tuning objective is

\[
L_{\text{SFT}}(\theta)
=
-
\frac{1}{M}
\sum_{i=1}^N\sum_{t=1}^{T_i}
\log p_\theta(y_{i,t}\mid x_i,y_{i,<t}).
\]

The prompt tokens provide context, but the loss is often applied only to
the response tokens. This is an important masking detail. If prompt
tokens are included in the loss, the model is trained to reproduce
prompts rather than only to generate useful responses.

\begin{tcolorbox}[enhanced jigsaw, arc=.35mm, bottomrule=.15mm, bottomtitle=1mm, breakable, colback=white, colbacktitle=quarto-callout-warning-color!10!white, colframe=quarto-callout-warning-color-frame, coltitle=black, left=2mm, leftrule=.75mm, opacityback=0, opacitybacktitle=0.6, rightrule=.15mm, title=\textcolor{quarto-callout-warning-color}{\faExclamationTriangle}\hspace{0.5em}{Masking matters}, titlerule=0mm, toprule=.15mm, toptitle=1mm]

For instruction fine-tuning, one often wants

\[
\text{loss on response tokens only}.
\]

The mask determines which tokens count as labels. This small
implementation choice changes the statistical problem being solved.

\end{tcolorbox}

\section{18.10 Fine-tuning as distributional
movement}\label{fine-tuning-as-distributional-movement}

Let \(p_0\) be the pretrained reference model and let \(p_\theta\) be
the fine-tuned model. Fine-tuning changes conditional distributions:

\[
p_0(\cdot\mid c)
\quad \longrightarrow \quad
p_\theta(\cdot\mid c).
\]

If the fine-tuning dataset is small, moving too far from the reference
model can reduce general ability. A common mathematical safeguard is a
KL penalty:

\[
L_{\text{KL-SFT}}(\theta)
=
L_{\text{SFT}}(\theta)
+
\beta
E_c\left[
D_{\mathrm{KL}}\left(
 p_\theta(\cdot\mid c)\;\|\;p_0(\cdot\mid c)
\right)
\right].
\]

The KL divergence is

\[
D_{\mathrm{KL}}(p\|q)=\sum_{j=1}^d p_j\log\frac{p_j}{q_j}.
\]

The coefficient \(\beta\geq 0\) controls how strongly the fine-tuned
model is discouraged from moving away from the reference model.

\section{18.11 Preference data and reward
models}\label{preference-data-and-reward-models}

Some training data does not specify a single correct response. Instead,
it gives preferences. For a prompt \(x\), suppose humans or another
evaluator prefer response \(y^+\) over response \(y^-\). A common model
assumes there is a scalar reward function

\[
r_\phi(x,y)\in R
\]

and that preference probabilities follow a Bradley--Terry logistic
model:

\[
P_\phi(y^+ \succ y^-\mid x)
=
\sigma\left(r_\phi(x,y^+)-r_\phi(x,y^-)\right),
\]

where

\[
\sigma(a)=\frac{1}{1+\exp(-a)}.
\]

The negative log-likelihood for one preference pair is

\[
\ell_{\text{pref}}(\phi)
=
-
\log \sigma\left(r_\phi(x,y^+)-r_\phi(x,y^-)\right).
\]

This is logistic regression on pairwise differences of rewards.

\section{18.12 KL-regularized policy optimization
viewpoint}\label{kl-regularized-policy-optimization-viewpoint}

A language model can also be viewed as a policy. Given a prompt \(x\),
it samples a response \(Y\) from

\[
Y\sim p_\theta(\cdot\mid x).
\]

Suppose a reward model \(r_\phi(x,Y)\) assigns a scalar score. A
KL-regularized objective is

\[
J(\theta)
=
E_{Y\sim p_\theta(\cdot\mid x)}[r_\phi(x,Y)]
-
\beta
D_{\mathrm{KL}}\left(
 p_\theta(\cdot\mid x)\;\|\;p_0(\cdot\mid x)
\right).
\]

The first term encourages high-reward responses. The second term
prevents the model from drifting too far from a reference model. This is
the mathematical core behind many alignment-style training procedures.

The optimization is difficult because the response \(Y\) is a discrete
sequence. Gradient estimates can have high variance. Practical
algorithms introduce approximations, clipping, baselines, and minibatch
estimates. The mathematical principle, however, is the same: maximize
reward while controlling distributional change.

\section{18.13 Direct preference optimization
viewpoint}\label{direct-preference-optimization-viewpoint}

Preference tuning can sometimes be written directly in terms of
log-probabilities rather than explicitly training a separate reward
model. A common form compares the preferred response \(y^+\) and
rejected response \(y^-\) under the current model and a reference model.

Define

\[
\Delta_\theta
=
\log p_\theta(y^+\mid x)-\log p_\theta(y^-\mid x),
\]

and

\[
\Delta_0
=
\log p_0(y^+\mid x)-\log p_0(y^-\mid x).
\]

A direct preference loss has the logistic form

\[
L_{\text{DPO}}(\theta)
=
-
\log \sigma\left(\beta(\Delta_\theta-\Delta_0)\right).
\]

This encourages the updated model to increase the relative
log-probability of the preferred response compared with the reference
model. The reference model acts as a statistical anchor.

\section{18.14 Parameter-efficient
fine-tuning}\label{parameter-efficient-fine-tuning}

Fine-tuning every parameter can be expensive. A parameter-efficient
method freezes most of the pretrained model and trains only a small
number of new parameters.

A simple mathematical model is a low-rank update to a weight matrix.
Suppose a pretrained layer contains

\[
W_0\in R^{m\times n}.
\]

Instead of training all entries of \(W_0\), we train matrices

\[
A\in R^{r\times n},\qquad B\in R^{m\times r},
\]

where \(r\ll \min(m,n)\), and use

\[
W=W_0+BA.
\]

The update \(BA\) has rank at most \(r\). This turns a large update in
\(R^{m\times n}\) into a smaller optimization problem with \(r(m+n)\)
trainable parameters.

\textbf{Linear algebra viewpoint.} Parameter-efficient fine-tuning
restricts the update direction to a low-dimensional subspace or manifold
of the full parameter space.

\section{18.15 What changes across training
stages?}\label{what-changes-across-training-stages}

The following table summarizes the mathematical differences.

{\def\LTcaptype{none} % do not increment counter
\begin{longtable}[]{@{}
  >{\raggedright\arraybackslash}p{(\linewidth - 6\tabcolsep) * \real{0.2500}}
  >{\raggedright\arraybackslash}p{(\linewidth - 6\tabcolsep) * \real{0.2500}}
  >{\raggedright\arraybackslash}p{(\linewidth - 6\tabcolsep) * \real{0.2500}}
  >{\raggedright\arraybackslash}p{(\linewidth - 6\tabcolsep) * \real{0.2500}}@{}}
\toprule\noalign{}
\begin{minipage}[b]{\linewidth}\raggedright
Stage
\end{minipage} & \begin{minipage}[b]{\linewidth}\raggedright
Data
\end{minipage} & \begin{minipage}[b]{\linewidth}\raggedright
Objective
\end{minipage} & \begin{minipage}[b]{\linewidth}\raggedright
Main mathematical object
\end{minipage} \\
\midrule\noalign{}
\endhead
\bottomrule\noalign{}
\endlastfoot
Pretraining & large unlabeled text corpus & next-token negative
log-likelihood & conditional distribution
\(p_\theta(w_t\mid w_{<t})\) \\
Supervised fine-tuning & prompt-response pairs & response-token
cross-entropy & conditional response model \(p_\theta(y\mid x)\) \\
Instruction tuning & instruction-following examples & masked conditional
cross-entropy & prompt-conditioned generation \\
Reward modeling & pairwise preferences & logistic pairwise loss & scalar
reward \(r_\phi(x,y)\) \\
KL-regularized tuning & prompts and rewards/preferences & reward minus
KL penalty & constrained distributional shift \\
Parameter-efficient tuning & task data & any fine-tuning loss &
low-dimensional parameter update \\
\end{longtable}
}

\section{18.16 Python example: token-level
cross-entropy}\label{python-example-token-level-cross-entropy}

The next examples use small arrays so that every quantity is visible.

\protect\phantomsection\label{ch18-cross-entropy}
\begin{Shaded}
\begin{Highlighting}[]
\ImportTok{import}\NormalTok{ numpy }\ImportTok{as}\NormalTok{ np}

\NormalTok{np.set\_printoptions(precision}\OperatorTok{=}\DecValTok{4}\NormalTok{, suppress}\OperatorTok{=}\VariableTok{True}\NormalTok{)}

\KeywordTok{def}\NormalTok{ softmax(logits):}
\NormalTok{    logits }\OperatorTok{=}\NormalTok{ np.asarray(logits, dtype}\OperatorTok{=}\BuiltInTok{float}\NormalTok{)}
\NormalTok{    shifted }\OperatorTok{=}\NormalTok{ logits }\OperatorTok{{-}}\NormalTok{ logits.}\BuiltInTok{max}\NormalTok{(axis}\OperatorTok{={-}}\DecValTok{1}\NormalTok{, keepdims}\OperatorTok{=}\VariableTok{True}\NormalTok{)}
\NormalTok{    exp\_values }\OperatorTok{=}\NormalTok{ np.exp(shifted)}
    \ControlFlowTok{return}\NormalTok{ exp\_values }\OperatorTok{/}\NormalTok{ exp\_values.}\BuiltInTok{sum}\NormalTok{(axis}\OperatorTok{={-}}\DecValTok{1}\NormalTok{, keepdims}\OperatorTok{=}\VariableTok{True}\NormalTok{)}

\KeywordTok{def}\NormalTok{ cross\_entropy\_from\_logits(logits, targets):}
\NormalTok{    probs }\OperatorTok{=}\NormalTok{ softmax(logits)}
\NormalTok{    n }\OperatorTok{=} \BuiltInTok{len}\NormalTok{(targets)}
\NormalTok{    losses }\OperatorTok{=} \OperatorTok{{-}}\NormalTok{np.log(probs[np.arange(n), targets])}
    \ControlFlowTok{return}\NormalTok{ losses, losses.mean(), probs}

\NormalTok{vocab }\OperatorTok{=}\NormalTok{ [}\StringTok{"math"}\NormalTok{, }\StringTok{"models"}\NormalTok{, }\StringTok{"learn"}\NormalTok{, }\StringTok{"data"}\NormalTok{]}
\NormalTok{logits }\OperatorTok{=}\NormalTok{ np.array([}
\NormalTok{    [}\FloatTok{2.0}\NormalTok{, }\FloatTok{1.0}\NormalTok{, }\FloatTok{0.0}\NormalTok{, }\OperatorTok{{-}}\FloatTok{1.0}\NormalTok{],}
\NormalTok{    [}\FloatTok{0.2}\NormalTok{, }\FloatTok{1.8}\NormalTok{, }\FloatTok{1.1}\NormalTok{, }\OperatorTok{{-}}\FloatTok{0.5}\NormalTok{],}
\NormalTok{    [}\OperatorTok{{-}}\FloatTok{0.3}\NormalTok{, }\FloatTok{0.4}\NormalTok{, }\FloatTok{2.2}\NormalTok{, }\FloatTok{1.0}\NormalTok{]}
\NormalTok{])}
\NormalTok{targets }\OperatorTok{=}\NormalTok{ np.array([}\DecValTok{0}\NormalTok{, }\DecValTok{2}\NormalTok{, }\DecValTok{3}\NormalTok{])}

\NormalTok{losses, mean\_loss, probs }\OperatorTok{=}\NormalTok{ cross\_entropy\_from\_logits(logits, targets)}
\BuiltInTok{print}\NormalTok{(}\StringTok{"probabilities:}\CharTok{\textbackslash{}n}\StringTok{"}\NormalTok{, probs)}
\BuiltInTok{print}\NormalTok{(}\StringTok{"token losses:"}\NormalTok{, losses)}
\BuiltInTok{print}\NormalTok{(}\StringTok{"mean loss:"}\NormalTok{, mean\_loss)}
\BuiltInTok{print}\NormalTok{(}\StringTok{"perplexity:"}\NormalTok{, np.exp(mean\_loss))}
\end{Highlighting}
\end{Shaded}

\begin{verbatim}
probabilities:
 [[0.6439 0.2369 0.0871 0.0321]
 [0.1122 0.5559 0.2761 0.0557]
 [0.053  0.1067 0.6458 0.1945]]
token losses: [0.4402 1.2871 1.6373]
mean loss: 1.1215378808659175
perplexity: 3.0695712100966532
\end{verbatim}

\section{18.17 Python example: gradient with respect to
logits}\label{python-example-gradient-with-respect-to-logits}

\protect\phantomsection\label{ch18-logit-gradient}
\begin{Shaded}
\begin{Highlighting}[]
\NormalTok{probs }\OperatorTok{=}\NormalTok{ softmax(logits)}
\NormalTok{one\_hot }\OperatorTok{=}\NormalTok{ np.zeros\_like(probs)}
\NormalTok{one\_hot[np.arange(}\BuiltInTok{len}\NormalTok{(targets)), targets] }\OperatorTok{=} \FloatTok{1.0}

\NormalTok{grad\_logits }\OperatorTok{=}\NormalTok{ (probs }\OperatorTok{{-}}\NormalTok{ one\_hot) }\OperatorTok{/} \BuiltInTok{len}\NormalTok{(targets)}
\BuiltInTok{print}\NormalTok{(}\StringTok{"gradient of mean cross{-}entropy with respect to logits:}\CharTok{\textbackslash{}n}\StringTok{"}\NormalTok{, grad\_logits)}
\BuiltInTok{print}\NormalTok{(}\StringTok{"row sums should be zero:"}\NormalTok{, grad\_logits.}\BuiltInTok{sum}\NormalTok{(axis}\OperatorTok{=}\DecValTok{1}\NormalTok{))}
\end{Highlighting}
\end{Shaded}

\begin{verbatim}
gradient of mean cross-entropy with respect to logits:
 [[-0.1187  0.079   0.029   0.0107]
 [ 0.0374  0.1853 -0.2413  0.0186]
 [ 0.0177  0.0356  0.2153 -0.2685]]
row sums should be zero: [ 0. -0.  0.]
\end{verbatim}

The row sums are zero because adding the same constant to all logits
does not change the softmax distribution.

\section{18.18 Python example: masked response
loss}\label{python-example-masked-response-loss}

Instruction fine-tuning often uses a loss mask. In the following toy
example, prompt tokens have mask value \(0\) and response tokens have
mask value \(1\).

\protect\phantomsection\label{ch18-masked-loss}
\begin{Shaded}
\begin{Highlighting}[]
\CommentTok{\# Shape: batch size 2, length 5, vocabulary size 4.}
\NormalTok{rng }\OperatorTok{=}\NormalTok{ np.random.default\_rng(}\DecValTok{18}\NormalTok{)}
\NormalTok{logits\_batch }\OperatorTok{=}\NormalTok{ rng.normal(size}\OperatorTok{=}\NormalTok{(}\DecValTok{2}\NormalTok{, }\DecValTok{5}\NormalTok{, }\DecValTok{4}\NormalTok{))}
\NormalTok{target\_batch }\OperatorTok{=}\NormalTok{ np.array([}
\NormalTok{    [}\DecValTok{0}\NormalTok{, }\DecValTok{1}\NormalTok{, }\DecValTok{2}\NormalTok{, }\DecValTok{3}\NormalTok{, }\DecValTok{0}\NormalTok{],}
\NormalTok{    [}\DecValTok{1}\NormalTok{, }\DecValTok{2}\NormalTok{, }\DecValTok{0}\NormalTok{, }\DecValTok{0}\NormalTok{, }\DecValTok{3}\NormalTok{]}
\NormalTok{])}

\CommentTok{\# Only response tokens count.}
\NormalTok{loss\_mask }\OperatorTok{=}\NormalTok{ np.array([}
\NormalTok{    [}\DecValTok{0}\NormalTok{, }\DecValTok{0}\NormalTok{, }\DecValTok{1}\NormalTok{, }\DecValTok{1}\NormalTok{, }\DecValTok{1}\NormalTok{],}
\NormalTok{    [}\DecValTok{0}\NormalTok{, }\DecValTok{1}\NormalTok{, }\DecValTok{1}\NormalTok{, }\DecValTok{0}\NormalTok{, }\DecValTok{1}\NormalTok{]}
\NormalTok{], dtype}\OperatorTok{=}\BuiltInTok{float}\NormalTok{)}

\NormalTok{probs\_batch }\OperatorTok{=}\NormalTok{ softmax(logits\_batch)}
\NormalTok{selected }\OperatorTok{=}\NormalTok{ probs\_batch[}
\NormalTok{    np.arange(}\DecValTok{2}\NormalTok{)[:, }\VariableTok{None}\NormalTok{],}
\NormalTok{    np.arange(}\DecValTok{5}\NormalTok{)[}\VariableTok{None}\NormalTok{, :],}
\NormalTok{    target\_batch}
\NormalTok{]}

\NormalTok{token\_losses }\OperatorTok{=} \OperatorTok{{-}}\NormalTok{np.log(selected)}
\NormalTok{masked\_mean }\OperatorTok{=}\NormalTok{ (token\_losses }\OperatorTok{*}\NormalTok{ loss\_mask).}\BuiltInTok{sum}\NormalTok{() }\OperatorTok{/}\NormalTok{ loss\_mask.}\BuiltInTok{sum}\NormalTok{()}

\BuiltInTok{print}\NormalTok{(}\StringTok{"token losses:}\CharTok{\textbackslash{}n}\StringTok{"}\NormalTok{, token\_losses)}
\BuiltInTok{print}\NormalTok{(}\StringTok{"loss mask:}\CharTok{\textbackslash{}n}\StringTok{"}\NormalTok{, loss\_mask)}
\BuiltInTok{print}\NormalTok{(}\StringTok{"masked mean loss:"}\NormalTok{, masked\_mean)}
\end{Highlighting}
\end{Shaded}

\begin{verbatim}
token losses:
 [[1.6153 1.6    4.1549 1.4498 0.7166]
 [1.274  1.6516 2.1492 2.0518 3.3522]]
loss mask:
 [[0. 0. 1. 1. 1.]
 [0. 1. 1. 0. 1.]]
masked mean loss: 2.245731514471048
\end{verbatim}

\section{18.19 Python example: KL penalty between current and reference
models}\label{python-example-kl-penalty-between-current-and-reference-models}

\protect\phantomsection\label{ch18-kl-penalty}
\begin{Shaded}
\begin{Highlighting}[]
\KeywordTok{def}\NormalTok{ kl\_divergence(p, q):}
\NormalTok{    p }\OperatorTok{=}\NormalTok{ np.asarray(p, dtype}\OperatorTok{=}\BuiltInTok{float}\NormalTok{)}
\NormalTok{    q }\OperatorTok{=}\NormalTok{ np.asarray(q, dtype}\OperatorTok{=}\BuiltInTok{float}\NormalTok{)}
    \ControlFlowTok{return}\NormalTok{ np.}\BuiltInTok{sum}\NormalTok{(p }\OperatorTok{*}\NormalTok{ (np.log(p }\OperatorTok{+} \FloatTok{1e{-}12}\NormalTok{) }\OperatorTok{{-}}\NormalTok{ np.log(q }\OperatorTok{+} \FloatTok{1e{-}12}\NormalTok{)), axis}\OperatorTok{={-}}\DecValTok{1}\NormalTok{)}

\NormalTok{current\_logits }\OperatorTok{=}\NormalTok{ np.array([}
\NormalTok{    [}\FloatTok{2.2}\NormalTok{, }\FloatTok{1.4}\NormalTok{, }\FloatTok{0.1}\NormalTok{, }\OperatorTok{{-}}\FloatTok{0.5}\NormalTok{],}
\NormalTok{    [}\FloatTok{0.0}\NormalTok{, }\FloatTok{1.2}\NormalTok{, }\FloatTok{2.1}\NormalTok{, }\OperatorTok{{-}}\FloatTok{0.8}\NormalTok{]}
\NormalTok{])}
\NormalTok{reference\_logits }\OperatorTok{=}\NormalTok{ np.array([}
\NormalTok{    [}\FloatTok{1.8}\NormalTok{, }\FloatTok{1.3}\NormalTok{, }\FloatTok{0.4}\NormalTok{, }\OperatorTok{{-}}\FloatTok{0.1}\NormalTok{],}
\NormalTok{    [}\FloatTok{0.2}\NormalTok{, }\FloatTok{1.0}\NormalTok{, }\FloatTok{1.6}\NormalTok{, }\OperatorTok{{-}}\FloatTok{0.2}\NormalTok{]}
\NormalTok{])}

\NormalTok{p\_current }\OperatorTok{=}\NormalTok{ softmax(current\_logits)}
\NormalTok{p\_reference }\OperatorTok{=}\NormalTok{ softmax(reference\_logits)}
\NormalTok{kl\_values }\OperatorTok{=}\NormalTok{ kl\_divergence(p\_current, p\_reference)}

\BuiltInTok{print}\NormalTok{(}\StringTok{"current probabilities:}\CharTok{\textbackslash{}n}\StringTok{"}\NormalTok{, p\_current)}
\BuiltInTok{print}\NormalTok{(}\StringTok{"reference probabilities:}\CharTok{\textbackslash{}n}\StringTok{"}\NormalTok{, p\_reference)}
\BuiltInTok{print}\NormalTok{(}\StringTok{"KL(current || reference):"}\NormalTok{, kl\_values)}
\BuiltInTok{print}\NormalTok{(}\StringTok{"mean KL:"}\NormalTok{, kl\_values.mean())}
\end{Highlighting}
\end{Shaded}

\begin{verbatim}
current probabilities:
 [[0.6101 0.2741 0.0747 0.041 ]
 [0.0773 0.2567 0.6313 0.0347]]
reference probabilities:
 [[0.4993 0.3029 0.1231 0.0747]
 [0.1258 0.2799 0.51   0.0843]]
KL(current || reference): [0.0331 0.044 ]
mean KL: 0.0385315810437777
\end{verbatim}

\section{18.20 Python example: preference
loss}\label{python-example-preference-loss}

\protect\phantomsection\label{ch18-preference-loss}
\begin{Shaded}
\begin{Highlighting}[]
\KeywordTok{def}\NormalTok{ sigmoid(x):}
    \ControlFlowTok{return} \FloatTok{1.0} \OperatorTok{/}\NormalTok{ (}\FloatTok{1.0} \OperatorTok{+}\NormalTok{ np.exp(}\OperatorTok{{-}}\NormalTok{x))}

\KeywordTok{def}\NormalTok{ preference\_loss(reward\_chosen, reward\_rejected):}
\NormalTok{    margin }\OperatorTok{=}\NormalTok{ reward\_chosen }\OperatorTok{{-}}\NormalTok{ reward\_rejected}
    \ControlFlowTok{return} \OperatorTok{{-}}\NormalTok{np.log(sigmoid(margin) }\OperatorTok{+} \FloatTok{1e{-}12}\NormalTok{)}

\NormalTok{reward\_chosen }\OperatorTok{=}\NormalTok{ np.array([}\FloatTok{2.0}\NormalTok{, }\FloatTok{0.8}\NormalTok{, }\FloatTok{1.5}\NormalTok{, }\FloatTok{0.4}\NormalTok{])}
\NormalTok{reward\_rejected }\OperatorTok{=}\NormalTok{ np.array([}\FloatTok{0.5}\NormalTok{, }\FloatTok{1.1}\NormalTok{, }\FloatTok{0.0}\NormalTok{, }\FloatTok{0.2}\NormalTok{])}
\NormalTok{losses }\OperatorTok{=}\NormalTok{ preference\_loss(reward\_chosen, reward\_rejected)}

\BuiltInTok{print}\NormalTok{(}\StringTok{"reward margins:"}\NormalTok{, reward\_chosen }\OperatorTok{{-}}\NormalTok{ reward\_rejected)}
\BuiltInTok{print}\NormalTok{(}\StringTok{"preference probabilities:"}\NormalTok{, sigmoid(reward\_chosen }\OperatorTok{{-}}\NormalTok{ reward\_rejected))}
\BuiltInTok{print}\NormalTok{(}\StringTok{"preference losses:"}\NormalTok{, losses)}
\BuiltInTok{print}\NormalTok{(}\StringTok{"mean preference loss:"}\NormalTok{, losses.mean())}
\end{Highlighting}
\end{Shaded}

\begin{verbatim}
reward margins: [ 1.5 -0.3  1.5  0.2]
preference probabilities: [0.8176 0.4256 0.8176 0.5498]
preference losses: [0.2014 0.8544 0.2014 0.5981]
mean preference loss: 0.46383016745225225
\end{verbatim}

\section{18.21 Python example: low-rank
adaptation}\label{python-example-low-rank-adaptation}

\protect\phantomsection\label{ch18-low-rank-adaptation}
\begin{Shaded}
\begin{Highlighting}[]
\NormalTok{rng }\OperatorTok{=}\NormalTok{ np.random.default\_rng(}\DecValTok{1818}\NormalTok{)}

\NormalTok{m, n, r }\OperatorTok{=} \DecValTok{6}\NormalTok{, }\DecValTok{5}\NormalTok{, }\DecValTok{2}
\NormalTok{W0 }\OperatorTok{=}\NormalTok{ rng.normal(scale}\OperatorTok{=}\FloatTok{0.2}\NormalTok{, size}\OperatorTok{=}\NormalTok{(m, n))}
\NormalTok{A }\OperatorTok{=}\NormalTok{ rng.normal(scale}\OperatorTok{=}\FloatTok{0.05}\NormalTok{, size}\OperatorTok{=}\NormalTok{(r, n))}
\NormalTok{B }\OperatorTok{=}\NormalTok{ rng.normal(scale}\OperatorTok{=}\FloatTok{0.05}\NormalTok{, size}\OperatorTok{=}\NormalTok{(m, r))}

\NormalTok{Delta }\OperatorTok{=}\NormalTok{ B }\OperatorTok{@}\NormalTok{ A}
\NormalTok{W }\OperatorTok{=}\NormalTok{ W0 }\OperatorTok{+}\NormalTok{ Delta}

\BuiltInTok{print}\NormalTok{(}\StringTok{"W0 shape:"}\NormalTok{, W0.shape)}
\BuiltInTok{print}\NormalTok{(}\StringTok{"A shape:"}\NormalTok{, A.shape)}
\BuiltInTok{print}\NormalTok{(}\StringTok{"B shape:"}\NormalTok{, B.shape)}
\BuiltInTok{print}\NormalTok{(}\StringTok{"rank of update B A:"}\NormalTok{, np.linalg.matrix\_rank(Delta))}
\BuiltInTok{print}\NormalTok{(}\StringTok{"trainable full matrix parameters:"}\NormalTok{, m }\OperatorTok{*}\NormalTok{ n)}
\BuiltInTok{print}\NormalTok{(}\StringTok{"trainable low{-}rank parameters:"}\NormalTok{, r }\OperatorTok{*}\NormalTok{ (m }\OperatorTok{+}\NormalTok{ n))}
\end{Highlighting}
\end{Shaded}

\begin{verbatim}
W0 shape: (6, 5)
A shape: (2, 5)
B shape: (6, 2)
rank of update B A: 2
trainable full matrix parameters: 30
trainable low-rank parameters: 22
\end{verbatim}

\section{18.22 Interactive: training objective
explorer}\label{interactive-training-objective-explorer}

This interactive shows how three ingredients affect a stylized training
curve: learning rate, batch size, and KL regularization strength. It is
not a physical law of neural-network training; it is a visualization of
the qualitative roles of step size, gradient noise, and distributional
anchoring.

\section{18.23 AI-assisted learning
components}\label{ai-assisted-learning-components-2}

Use an AI assistant as a mathematical critic, not as an authority. The
following prompts are designed for students to check understanding.

\subsection{Prompt 1: identify the
objective}\label{prompt-1-identify-the-objective}

\begin{Shaded}
\begin{Highlighting}[]
\NormalTok{I am studying language model training.  Given a training setup, help me identify the data distribution, model distribution, loss function, mask, and optimization variable.  Do not give code first.  Start with mathematical notation.}
\end{Highlighting}
\end{Shaded}

\subsection{Prompt 2: audit a fine-tuning
dataset}\label{prompt-2-audit-a-fine-tuning-dataset}

\begin{Shaded}
\begin{Highlighting}[]
\NormalTok{I have a prompt{-}response fine{-}tuning dataset.  Help me decide which tokens should be included in the loss mask and explain how the choice changes the statistical objective.}
\end{Highlighting}
\end{Shaded}

\subsection{Prompt 3: compare alignment
objectives}\label{prompt-3-compare-alignment-objectives}

\begin{Shaded}
\begin{Highlighting}[]
\NormalTok{Compare supervised fine{-}tuning, reward{-}model training, KL{-}regularized policy optimization, and direct preference optimization.  For each one, state the objective function and the random variables.}
\end{Highlighting}
\end{Shaded}

\subsection{Prompt 4: detect
overclaiming}\label{prompt-4-detect-overclaiming}

\begin{Shaded}
\begin{Highlighting}[]
\NormalTok{A student says: "RLHF makes a model truthful."  Critique this statement mathematically.  Discuss what objective is optimized and why optimizing that objective does not automatically prove truthfulness.}
\end{Highlighting}
\end{Shaded}

\section{18.24 Common mistakes}\label{common-mistakes-7}

\begin{enumerate}
\def\labelenumi{\arabic{enumi}.}
\tightlist
\item
  \textbf{Confusing pretraining with memorization.} The objective
  estimates conditional token distributions. Memorization may occur, but
  it is not the mathematical objective.
\item
  \textbf{Ignoring masks.} The mask determines which tokens are treated
  as labels.
\item
  \textbf{Treating KL divergence as symmetric.} In general,
  \(D_{\mathrm{KL}}(p\|q)\neq D_{\mathrm{KL}}(q\|p)\).
\item
  \textbf{Assuming lower training loss always means better behavior.}
  Fine-tuning may improve one distribution while hurting another.
\item
  \textbf{Equating reward with truth.} A reward model estimates
  preferences or scores, not a theorem about correctness.
\item
  \textbf{Forgetting the reference model.} Many preference objectives
  are meaningful only relative to a reference distribution.
\end{enumerate}

\section{18.25 Worked calculation}\label{worked-calculation-2}

Suppose a vocabulary has three tokens and the model assigns
probabilities

\[
p_\theta(\cdot\mid c)=(0.70,0.20,0.10).
\]

If the observed token is token \(2\), the token loss is

\[
\ell=-\log(0.20)\approx 1.609.
\]

If a reference model has

\[
p_0(\cdot\mid c)=(0.60,0.30,0.10),
\]

then

\[
D_{\mathrm{KL}}(p_\theta\|p_0)
=0.70\log\frac{0.70}{0.60}
+0.20\log\frac{0.20}{0.30}
+0.10\log\frac{0.10}{0.10}.
\]

Numerically,

\[
D_{\mathrm{KL}}(p_\theta\|p_0)
\approx 0.70(0.154)-0.20(0.405)+0
\approx 0.027.
\]

With \(\beta=0.5\), the KL-regularized token objective is approximately

\[
1.609+0.5(0.027)=1.6225.
\]

The penalty is small here because the current model is close to the
reference distribution.

\section{18.26 Exercises}\label{exercises-17}

\begin{enumerate}
\def\labelenumi{\arabic{enumi}.}
\tightlist
\item
  \textbf{Pretraining objective.} Starting from the autoregressive
  factorization, derive the negative log-likelihood loss for a corpus of
  variable-length sequences.
\item
  \textbf{Softmax gradient.} Prove that \(\nabla_z\ell=p-e_y\) for the
  multiclass cross-entropy loss.
\item
  \textbf{Masking.} Construct a prompt-response example and write a
  binary mask that applies loss only to the response tokens.
\item
  \textbf{KL direction.} Compute both \(D_{\mathrm{KL}}(p\|q)\) and
  \(D_{\mathrm{KL}}(q\|p)\) for \(p=(0.8,0.2)\) and \(q=(0.5,0.5)\).
\item
  \textbf{Preference loss.} For rewards \(r^+=1.4\) and \(r^-=0.2\),
  compute the Bradley--Terry preference probability and loss.
\item
  \textbf{Low-rank update.} If \(W_0\in R^{4096\times 4096}\) and rank
  \(r=8\), compare the number of full fine-tuning parameters with the
  number of low-rank update parameters.
\item
  \textbf{Validation gap.} Suppose training loss decreases but
  validation loss increases. Give at least two mathematical
  explanations.
\item
  \textbf{Objective comparison.} Write one paragraph comparing
  supervised fine-tuning and preference tuning as empirical risk
  minimization problems.
\end{enumerate}

\section{18.27 Mini-project}\label{mini-project-3}

Choose a small text dataset and implement a miniature training pipeline.

\begin{enumerate}
\def\labelenumi{\arabic{enumi}.}
\tightlist
\item
  Build a vocabulary and encode token sequences.
\item
  Construct next-token input-target pairs.
\item
  Train a simple model such as a bigram table, linear classifier, or
  tiny neural model.
\item
  Track training loss and validation loss.
\item
  Add a mask that excludes selected tokens from the loss.
\item
  Add a KL penalty relative to a reference distribution.
\item
  Write a mathematical report explaining the objective, the masks, and
  the learned conditional distributions.
\end{enumerate}

\section{18.28 Python lab connection}\label{python-lab-connection-9}

The associated independent-study lab is planned as:

\begin{Shaded}
\begin{Highlighting}[]
\NormalTok{labs/chapter{-}18{-}training{-}large{-}language{-}models.ipynb}
\end{Highlighting}
\end{Shaded}

The lab should include:

\begin{itemize}
\tightlist
\item
  background on maximum likelihood and cross-entropy,
\item
  implementation of token-level losses,
\item
  masked instruction-tuning losses,
\item
  KL-regularized fine-tuning on toy distributions,
\item
  pairwise preference loss,
\item
  a small training loop,
\item
  interpretation questions that force students to separate data, model,
  objective, and optimizer.
\end{itemize}

\section{18.29 What to remember}\label{what-to-remember-6}

\begin{itemize}
\tightlist
\item
  Pretraining is maximum likelihood estimation for next-token
  prediction.
\item
  The token-level loss is multiclass cross-entropy.
\item
  Teacher forcing uses true previous tokens during training.
\item
  Masks define which tokens contribute to the empirical risk.
\item
  Fine-tuning changes a pretrained conditional distribution.
\item
  KL penalties control movement away from a reference model.
\item
  Preference data can be modeled with pairwise logistic losses.
\item
  Alignment-style objectives optimize mathematical proxies for desired
  behavior; they do not automatically prove correctness, truthfulness,
  or safety.
\end{itemize}

\section{Notes and references}\label{notes-and-references-10}

This chapter connects maximum likelihood estimation, stochastic
optimization, preference learning, and modern language-model training.
Useful background sources include Jurafsky and Martin (2024), Bishop
(2006b), Goodfellow et al. (2016b), Vaswani et al. (2017), {Brown et
al.} (2020), {Ouyang et al.} (2022a), and Rafailov et al. (2023a).

\chapter{Chapter 19. Embeddings, Retrieval, and
RAG}\label{chapter-19.-embeddings-retrieval-and-rag}

Semantic search, nearest neighbors, and retrieval-augmented generation

\hfill\break

\section{19.1 Why retrieval belongs in a mathematical NLP
book}\label{why-retrieval-belongs-in-a-mathematical-nlp-book}

Large language models estimate conditional distributions over tokens. In
an autoregressive model, the central object is

\[
p_\theta(w_t \mid w_{<t}).
\]

This distribution is learned from training data and depends on the
information contained in the prompt. A model may have learned general
linguistic and statistical patterns, but it may not know the contents of
a new course policy, a recent paper, a private document collection, or a
local advising rule. Retrieval-augmented generation, usually abbreviated
as \textbf{RAG}, addresses this limitation by adding an explicit search
step.

A RAG system separates two mathematical operations:

\begin{enumerate}
\def\labelenumi{\arabic{enumi}.}
\tightlist
\item
  \textbf{Retrieval:} given a query \(q\), choose documents or chunks
  \(c_1,\dots,c_k\) from a collection \(C\).
\item
  \textbf{Generation:} produce an answer from a conditional model using
  both the query and the retrieved context.
\end{enumerate}

A simplified RAG model has the form

\[
p_\theta(y \mid q)
\approx
p_\theta(y \mid q, r_k(q)),
\]

where

\[
r_k(q) = (c_{i_1},\dots,c_{i_k})
\]

is the top-\(k\) retrieved context. In a more probabilistic formulation,
the retrieved document is treated as a latent variable:

\[
p(y \mid q)
= \sum_{c \in C} p_\theta(y \mid q,c) p_\eta(c \mid q).
\]

This chapter studies the mathematics behind \(p_\eta(c \mid q)\):
embeddings, similarity functions, nearest-neighbor search, vector
indexes, ranking, and retrieval evaluation. The emphasis is not on
building a production search engine, but on understanding the objects
that make semantic retrieval possible.

\begin{tcolorbox}[enhanced jigsaw, arc=.35mm, bottomrule=.15mm, bottomtitle=1mm, breakable, colback=white, colbacktitle=quarto-callout-tip-color!10!white, colframe=quarto-callout-tip-color-frame, coltitle=black, left=2mm, leftrule=.75mm, opacityback=0, opacitybacktitle=0.6, rightrule=.15mm, title=\textcolor{quarto-callout-tip-color}{\faLightbulb}\hspace{0.5em}{Main mathematical theme}, titlerule=0mm, toprule=.15mm, toptitle=1mm]

RAG turns part of language modeling into a problem in finite-dimensional
geometry. Queries and documents are mapped into vectors; retrieval
becomes an optimization problem over a point cloud; generation becomes
conditional modeling with selected external evidence.

\end{tcolorbox}

\section{19.2 Learning goals}\label{learning-goals-18}

After studying this chapter, students should be able to:

\begin{itemize}
\tightlist
\item
  define a sentence or document embedding as a map into \(R^m\);
\item
  compare sparse lexical retrieval and dense semantic retrieval;
\item
  compute dot-product, cosine, and Euclidean nearest neighbors;
\item
  explain why normalization changes retrieval geometry;
\item
  formulate retrieval as a maximum-score optimization problem;
\item
  understand approximate nearest-neighbor search at a high level;
\item
  explain chunking, indexing, retrieval, reranking, and context
  construction;
\item
  evaluate retrieval using precision, recall, MRR, and nDCG;
\item
  describe failure modes of RAG using probability and geometry.
\end{itemize}

\section{19.3 From word embeddings to document
embeddings}\label{from-word-embeddings-to-document-embeddings}

In Chapter 8 we represented individual words using an embedding map

\[
e: V \to R^m.
\]

For retrieval, we usually need embeddings of longer text units: queries,
sentences, paragraphs, document chunks, or full documents. Let \(T\)
denote the set of all finite token sequences. A text embedding model is
a map

\[
f_\theta: T \to R^m.
\]

For a query \(q \in T\) and a document chunk \(c \in T\), we compute

\[
z_q = f_\theta(q),
\qquad
z_c = f_\theta(c).
\]

The retrieval score is then a function

\[
s(q,c) = S(z_q,z_c),
\]

where \(S:R^m \times R^m \to R\) is a similarity measure. Common choices
are

\[
S_{\text{dot}}(z_q,z_c)=z_q^T z_c,
\]

\[
S_{\text{cos}}(z_q,z_c)=\frac{z_q^T z_c}{\|z_q\|\|z_c\|},
\]

and

\[
S_{\text{neg-dist}}(z_q,z_c)=-\|z_q-z_c\|_2^2.
\]

The top-\(k\) retrieval problem is

\[
\operatorname{TopK}(q)
=
\operatorname{argtopk}_{c \in C} S(f_\theta(q),f_\theta(c)).
\]

This is a finite optimization problem over a collection of vectors.

\section{19.4 Sparse retrieval versus dense
retrieval}\label{sparse-retrieval-versus-dense-retrieval}

Classical information retrieval often represents a document by a sparse
vector. If the vocabulary has size \(d\), then a document \(c\) may be
represented as

\[
x_c \in R^d,
\]

where \(x_{cj}\) is a count or TF-IDF weight for token \(j\). A query
\(q\) is represented similarly by \(x_q\), and the score is often a dot
product or cosine similarity.

Dense retrieval instead represents the same text by

\[
z_c \in R^m,
\]

where \(m\) is usually much smaller than \(d\), and the coordinates are
learned. Dense embeddings are designed so that semantically related
texts can be close even when they do not share exact words.

{\def\LTcaptype{none} % do not increment counter
\begin{longtable}[]{@{}
  >{\raggedright\arraybackslash}p{(\linewidth - 8\tabcolsep) * \real{0.1875}}
  >{\raggedleft\arraybackslash}p{(\linewidth - 8\tabcolsep) * \real{0.2500}}
  >{\raggedright\arraybackslash}p{(\linewidth - 8\tabcolsep) * \real{0.1875}}
  >{\raggedright\arraybackslash}p{(\linewidth - 8\tabcolsep) * \real{0.1875}}
  >{\raggedright\arraybackslash}p{(\linewidth - 8\tabcolsep) * \real{0.1875}}@{}}
\toprule\noalign{}
\begin{minipage}[b]{\linewidth}\raggedright
Representation
\end{minipage} & \begin{minipage}[b]{\linewidth}\raggedleft
Vector dimension
\end{minipage} & \begin{minipage}[b]{\linewidth}\raggedright
Coordinates mean
\end{minipage} & \begin{minipage}[b]{\linewidth}\raggedright
Strength
\end{minipage} & \begin{minipage}[b]{\linewidth}\raggedright
Weakness
\end{minipage} \\
\midrule\noalign{}
\endhead
\bottomrule\noalign{}
\endlastfoot
Bag-of-words or TF-IDF & \(d=|V|\) & token weights & exact lexical
matching & weak semantic matching \\
Dense embedding & \(m\) & learned features & semantic matching & can
miss exact constraints \\
Hybrid retrieval & both & lexical and semantic & often robust & more
complex \\
\end{longtable}
}

The mathematical difference is important. Sparse retrieval mostly uses
explicit coordinates corresponding to words. Dense retrieval uses an
embedding map learned from data. Thus, dense retrieval transfers part of
the retrieval problem into representation learning.

\section{19.5 Similarity geometry}\label{similarity-geometry}

Suppose all embeddings are nonzero. Define normalized embeddings by

\[
\widehat z = \frac{z}{\|z\|_2}.
\]

For normalized vectors,

\[
\widehat z_q^T \widehat z_c
= \cos(\widehat z_q,\widehat z_c).
\]

Also,

\[
\|\widehat z_q - \widehat z_c\|_2^2
= \|\widehat z_q\|^2 + \|\widehat z_c\|^2 - 2\widehat z_q^T\widehat z_c
= 2 - 2\widehat z_q^T\widehat z_c.
\]

Therefore, for normalized vectors, maximizing cosine similarity is
equivalent to minimizing Euclidean distance:

\[
\operatorname{argmax}_c \widehat z_q^T \widehat z_c
=
\operatorname{argmin}_c \|\widehat z_q-\widehat z_c\|_2.
\]

This is why many vector databases support both inner-product search and
Euclidean nearest-neighbor search. With normalized embeddings, these
objectives are equivalent up to a monotone transformation.

\begin{tcolorbox}[enhanced jigsaw, arc=.35mm, bottomrule=.15mm, bottomtitle=1mm, breakable, colback=white, colbacktitle=quarto-callout-warning-color!10!white, colframe=quarto-callout-warning-color-frame, coltitle=black, left=2mm, leftrule=.75mm, opacityback=0, opacitybacktitle=0.6, rightrule=.15mm, title=\textcolor{quarto-callout-warning-color}{\faExclamationTriangle}\hspace{0.5em}{Normalization changes the retrieval problem}, titlerule=0mm, toprule=.15mm, toptitle=1mm]

Without normalization, the dot product depends on both direction and
vector length:

\[
z_q^T z_c = \|z_q\|\|z_c\|\cos(z_q,z_c).
\]

Thus, a long vector can score highly because of its norm, not only
because of semantic alignment. Cosine similarity removes this norm
effect.

\end{tcolorbox}

\section{19.6 Retrieval as matrix
multiplication}\label{retrieval-as-matrix-multiplication}

Let the collection contain \(n\) chunks \(c_1,\dots,c_n\). Stack their
embeddings as rows of a matrix

\[
Z =
\begin{bmatrix}
- z_{c_1}^T - \\
- z_{c_2}^T - \\
\vdots \\
- z_{c_n}^T -
\end{bmatrix}
\in R^{n \times m}.
\]

For a query embedding \(z_q \in R^m\), dot-product scores are

\[
s = Z z_q \in R^n.
\]

The \(i\)th component is

\[
s_i = z_{c_i}^T z_q.
\]

For cosine retrieval, first normalize each row of \(Z\) and the query
vector:

\[
\widehat Z_i = \frac{z_{c_i}}{\|z_{c_i}\|},
\qquad
\widehat z_q = \frac{z_q}{\|z_q\|}.
\]

Then

\[
s = \widehat Z \widehat z_q.
\]

So dense retrieval is mathematically simple: compute an embedding,
multiply by a matrix, sort the scores.

\section{19.7 A small retrieval
example}\label{a-small-retrieval-example}

Consider six short chunks:

\begin{enumerate}
\def\labelenumi{\arabic{enumi}.}
\tightlist
\item
  ``students may transfer graduate credits with approval''
\item
  ``the course covers stochastic processes and Markov chains''
\item
  ``the library has study rooms and research databases''
\item
  ``machine learning uses embeddings and optimization''
\item
  ``applied mathematics students study probability and linear algebra''
\item
  ``transfer credit forms require the program director signature''
\end{enumerate}

A query such as ``How do I transfer credits?'' should retrieve chunks 1
and 6. A query such as ``What math background is needed for machine
learning?'' should retrieve chunks 4 and 5.

A sparse lexical system will rely on words like ``transfer'' and
``credit.'' A dense embedding system may also learn that ``approval,''
``form,'' and ``signature'' are related to administrative transfer
rules. The central question is how well the embedding map preserves the
task-relevant notion of similarity.

\section{19.8 Python: dense retrieval from scratch with toy
embeddings}\label{python-dense-retrieval-from-scratch-with-toy-embeddings}

The following example builds a tiny semantic search engine using
manually designed low-dimensional embeddings. The goal is not to train a
real embedding model; the goal is to make the retrieval geometry
transparent.

\begin{Shaded}
\begin{Highlighting}[]
\ImportTok{import}\NormalTok{ numpy }\ImportTok{as}\NormalTok{ np}
\ImportTok{import}\NormalTok{ pandas }\ImportTok{as}\NormalTok{ pd}

\NormalTok{chunks }\OperatorTok{=}\NormalTok{ [}
    \StringTok{"students may transfer graduate credits with approval"}\NormalTok{,}
    \StringTok{"the course covers stochastic processes and Markov chains"}\NormalTok{,}
    \StringTok{"the library has study rooms and research databases"}\NormalTok{,}
    \StringTok{"machine learning uses embeddings and optimization"}\NormalTok{,}
    \StringTok{"applied mathematics students study probability and linear algebra"}\NormalTok{,}
    \StringTok{"transfer credit forms require the program director signature"}\NormalTok{,}
\NormalTok{]}

\CommentTok{\# Coordinates are interpretable toy semantic features:}
\CommentTok{\# [administration, probability/math, ML/AI, campus resources]}
\NormalTok{toy\_embeddings }\OperatorTok{=}\NormalTok{ np.array([}
\NormalTok{    [}\FloatTok{0.95}\NormalTok{, }\FloatTok{0.10}\NormalTok{, }\FloatTok{0.05}\NormalTok{, }\FloatTok{0.05}\NormalTok{],}
\NormalTok{    [}\FloatTok{0.05}\NormalTok{, }\FloatTok{0.95}\NormalTok{, }\FloatTok{0.05}\NormalTok{, }\FloatTok{0.02}\NormalTok{],}
\NormalTok{    [}\FloatTok{0.02}\NormalTok{, }\FloatTok{0.10}\NormalTok{, }\FloatTok{0.05}\NormalTok{, }\FloatTok{0.95}\NormalTok{],}
\NormalTok{    [}\FloatTok{0.05}\NormalTok{, }\FloatTok{0.25}\NormalTok{, }\FloatTok{0.95}\NormalTok{, }\FloatTok{0.05}\NormalTok{],}
\NormalTok{    [}\FloatTok{0.05}\NormalTok{, }\FloatTok{0.90}\NormalTok{, }\FloatTok{0.30}\NormalTok{, }\FloatTok{0.05}\NormalTok{],}
\NormalTok{    [}\FloatTok{0.98}\NormalTok{, }\FloatTok{0.05}\NormalTok{, }\FloatTok{0.05}\NormalTok{, }\FloatTok{0.02}\NormalTok{],}
\NormalTok{], dtype}\OperatorTok{=}\BuiltInTok{float}\NormalTok{)}

\NormalTok{queries }\OperatorTok{=}\NormalTok{ \{}
    \StringTok{"How do I transfer credits?"}\NormalTok{: np.array([}\FloatTok{1.0}\NormalTok{, }\FloatTok{0.05}\NormalTok{, }\FloatTok{0.05}\NormalTok{, }\FloatTok{0.02}\NormalTok{]),}
    \StringTok{"What math background helps ML?"}\NormalTok{: np.array([}\FloatTok{0.03}\NormalTok{, }\FloatTok{0.75}\NormalTok{, }\FloatTok{0.75}\NormalTok{, }\FloatTok{0.03}\NormalTok{]),}
    \StringTok{"Where can I study on campus?"}\NormalTok{: np.array([}\FloatTok{0.02}\NormalTok{, }\FloatTok{0.05}\NormalTok{, }\FloatTok{0.05}\NormalTok{, }\FloatTok{1.0}\NormalTok{]),}
\NormalTok{\}}

\KeywordTok{def}\NormalTok{ normalize\_rows(A):}
\NormalTok{    norms }\OperatorTok{=}\NormalTok{ np.linalg.norm(A, axis}\OperatorTok{=}\DecValTok{1}\NormalTok{, keepdims}\OperatorTok{=}\VariableTok{True}\NormalTok{)}
    \ControlFlowTok{return}\NormalTok{ A }\OperatorTok{/}\NormalTok{ np.maximum(norms, }\FloatTok{1e{-}12}\NormalTok{)}

\KeywordTok{def}\NormalTok{ cosine\_search(query\_vec, doc\_matrix, k}\OperatorTok{=}\DecValTok{3}\NormalTok{):}
\NormalTok{    Z }\OperatorTok{=}\NormalTok{ normalize\_rows(doc\_matrix)}
\NormalTok{    q }\OperatorTok{=}\NormalTok{ query\_vec }\OperatorTok{/} \BuiltInTok{max}\NormalTok{(np.linalg.norm(query\_vec), }\FloatTok{1e{-}12}\NormalTok{)}
\NormalTok{    scores }\OperatorTok{=}\NormalTok{ Z }\OperatorTok{@}\NormalTok{ q}
\NormalTok{    order }\OperatorTok{=}\NormalTok{ np.argsort(scores)[::}\OperatorTok{{-}}\DecValTok{1}\NormalTok{][:k]}
    \ControlFlowTok{return}\NormalTok{ pd.DataFrame(\{}
        \StringTok{"rank"}\NormalTok{: np.arange(}\DecValTok{1}\NormalTok{, }\BuiltInTok{len}\NormalTok{(order) }\OperatorTok{+} \DecValTok{1}\NormalTok{),}
        \StringTok{"chunk"}\NormalTok{: [chunks[i] }\ControlFlowTok{for}\NormalTok{ i }\KeywordTok{in}\NormalTok{ order],}
        \StringTok{"score"}\NormalTok{: scores[order].}\BuiltInTok{round}\NormalTok{(}\DecValTok{3}\NormalTok{),}
\NormalTok{    \})}

\NormalTok{cosine\_search(queries[}\StringTok{"How do I transfer credits?"}\NormalTok{], toy\_embeddings, k}\OperatorTok{=}\DecValTok{3}\NormalTok{)}
\end{Highlighting}
\end{Shaded}

\protect\phantomsection\label{ch19-toy-dense-retrieval}
{\def\LTcaptype{none} % do not increment counter
\begin{longtable}[]{@{}llll@{}}
\toprule\noalign{}
& rank & chunk & score \\
\midrule\noalign{}
\endhead
\bottomrule\noalign{}
\endlastfoot
0 & 1 & transfer credit forms require the program dire... & 1.000 \\
1 & 2 & students may transfer graduate credits with ap... & 0.998 \\
2 & 3 & applied mathematics students study probability... & 0.116 \\
\end{longtable}
}

The retrieval score is not a probability. It is a geometric score. Later
we may convert scores to a probability-like distribution using a
softmax, but this adds modeling assumptions.

\section{19.9 Sparse retrieval with
TF-IDF}\label{sparse-retrieval-with-tf-idf}

Dense retrieval is not always better than sparse retrieval. If a query
contains a precise course number, policy phrase, chemical name, theorem
name, or error message, exact lexical matching can be essential.

The following code builds a small TF-IDF retriever without relying on a
specialized library.

\begin{Shaded}
\begin{Highlighting}[]
\ImportTok{import}\NormalTok{ re}
\ImportTok{from}\NormalTok{ collections }\ImportTok{import}\NormalTok{ Counter}

\KeywordTok{def}\NormalTok{ tokenize(text):}
    \ControlFlowTok{return}\NormalTok{ re.findall(}\VerbatimStringTok{r"}\PreprocessorTok{[a{-}z]}\OperatorTok{+}\VerbatimStringTok{"}\NormalTok{, text.lower())}

\NormalTok{vocab }\OperatorTok{=} \BuiltInTok{sorted}\NormalTok{(}\BuiltInTok{set}\NormalTok{(token }\ControlFlowTok{for}\NormalTok{ c }\KeywordTok{in}\NormalTok{ chunks }\ControlFlowTok{for}\NormalTok{ token }\KeywordTok{in}\NormalTok{ tokenize(c)))}
\NormalTok{idx }\OperatorTok{=}\NormalTok{ \{w: j }\ControlFlowTok{for}\NormalTok{ j, w }\KeywordTok{in} \BuiltInTok{enumerate}\NormalTok{(vocab)\}}

\NormalTok{X }\OperatorTok{=}\NormalTok{ np.zeros((}\BuiltInTok{len}\NormalTok{(chunks), }\BuiltInTok{len}\NormalTok{(vocab)))}
\ControlFlowTok{for}\NormalTok{ i, c }\KeywordTok{in} \BuiltInTok{enumerate}\NormalTok{(chunks):}
\NormalTok{    counts }\OperatorTok{=}\NormalTok{ Counter(tokenize(c))}
    \ControlFlowTok{for}\NormalTok{ w, count }\KeywordTok{in}\NormalTok{ counts.items():}
\NormalTok{        X[i, idx[w]] }\OperatorTok{=}\NormalTok{ count}

\CommentTok{\# smoothed IDF}
\NormalTok{n\_docs }\OperatorTok{=}\NormalTok{ X.shape[}\DecValTok{0}\NormalTok{]}
\NormalTok{df }\OperatorTok{=}\NormalTok{ (X }\OperatorTok{\textgreater{}} \DecValTok{0}\NormalTok{).}\BuiltInTok{sum}\NormalTok{(axis}\OperatorTok{=}\DecValTok{0}\NormalTok{)}
\NormalTok{idf }\OperatorTok{=}\NormalTok{ np.log((}\DecValTok{1} \OperatorTok{+}\NormalTok{ n\_docs) }\OperatorTok{/}\NormalTok{ (}\DecValTok{1} \OperatorTok{+}\NormalTok{ df)) }\OperatorTok{+} \DecValTok{1}
\NormalTok{X\_tfidf }\OperatorTok{=}\NormalTok{ X }\OperatorTok{*}\NormalTok{ idf}

\KeywordTok{def}\NormalTok{ vectorize\_query(text):}
\NormalTok{    q }\OperatorTok{=}\NormalTok{ np.zeros(}\BuiltInTok{len}\NormalTok{(vocab))}
\NormalTok{    counts }\OperatorTok{=}\NormalTok{ Counter(tokenize(text))}
    \ControlFlowTok{for}\NormalTok{ w, count }\KeywordTok{in}\NormalTok{ counts.items():}
        \ControlFlowTok{if}\NormalTok{ w }\KeywordTok{in}\NormalTok{ idx:}
\NormalTok{            q[idx[w]] }\OperatorTok{=}\NormalTok{ count}
    \ControlFlowTok{return}\NormalTok{ q }\OperatorTok{*}\NormalTok{ idf}

\KeywordTok{def}\NormalTok{ sparse\_search(query, k}\OperatorTok{=}\DecValTok{3}\NormalTok{):}
\NormalTok{    q }\OperatorTok{=}\NormalTok{ vectorize\_query(query)}
\NormalTok{    Z }\OperatorTok{=}\NormalTok{ normalize\_rows(X\_tfidf)}
\NormalTok{    q }\OperatorTok{=}\NormalTok{ q }\OperatorTok{/} \BuiltInTok{max}\NormalTok{(np.linalg.norm(q), }\FloatTok{1e{-}12}\NormalTok{)}
\NormalTok{    scores }\OperatorTok{=}\NormalTok{ Z }\OperatorTok{@}\NormalTok{ q}
\NormalTok{    order }\OperatorTok{=}\NormalTok{ np.argsort(scores)[::}\OperatorTok{{-}}\DecValTok{1}\NormalTok{][:k]}
    \ControlFlowTok{return}\NormalTok{ pd.DataFrame(\{}
        \StringTok{"rank"}\NormalTok{: np.arange(}\DecValTok{1}\NormalTok{, }\BuiltInTok{len}\NormalTok{(order) }\OperatorTok{+} \DecValTok{1}\NormalTok{),}
        \StringTok{"chunk"}\NormalTok{: [chunks[i] }\ControlFlowTok{for}\NormalTok{ i }\KeywordTok{in}\NormalTok{ order],}
        \StringTok{"score"}\NormalTok{: scores[order].}\BuiltInTok{round}\NormalTok{(}\DecValTok{3}\NormalTok{),}
\NormalTok{    \})}

\NormalTok{sparse\_search(}\StringTok{"transfer credit approval form"}\NormalTok{, k}\OperatorTok{=}\DecValTok{3}\NormalTok{)}
\end{Highlighting}
\end{Shaded}

\protect\phantomsection\label{ch19-tfidf-retrieval}
{\def\LTcaptype{none} % do not increment counter
\begin{longtable}[]{@{}llll@{}}
\toprule\noalign{}
& rank & chunk & score \\
\midrule\noalign{}
\endhead
\bottomrule\noalign{}
\endlastfoot
0 & 1 & students may transfer graduate credits with ap... & 0.406 \\
1 & 2 & transfer credit forms require the program dire... & 0.383 \\
2 & 3 & applied mathematics students study probability... & 0.000 \\
\end{longtable}
}

TF-IDF can be understood as a diagonal rescaling of a document-term
matrix. If \(X \in R^{n \times d}\) is a count matrix and
\(D \in R^{d \times d}\) is a diagonal matrix of
inverse-document-frequency weights, then the TF-IDF matrix is

\[
X_{\text{tfidf}} = X D.
\]

Retrieval is then performed using row-normalized vectors from
\(X_{\text{tfidf}}\).

\section{19.10 Hybrid retrieval}\label{hybrid-retrieval}

In many applications, sparse and dense retrieval are complementary. One
simple hybrid score is

\[
s_{\text{hybrid}}(q,c)
= \lambda s_{\text{dense}}(q,c)
+ (1-\lambda)s_{\text{sparse}}(q,c),
\qquad 0 \leq \lambda \leq 1.
\]

Here \(\lambda\) controls the tradeoff between semantic matching and
exact-word matching. This formula is simple but useful. It also makes a
conceptual point: retrieval design is often a choice of scoring
function.

\begin{Shaded}
\begin{Highlighting}[]
\KeywordTok{def}\NormalTok{ dense\_scores(query\_vec):}
\NormalTok{    Z }\OperatorTok{=}\NormalTok{ normalize\_rows(toy\_embeddings)}
\NormalTok{    q }\OperatorTok{=}\NormalTok{ query\_vec }\OperatorTok{/} \BuiltInTok{max}\NormalTok{(np.linalg.norm(query\_vec), }\FloatTok{1e{-}12}\NormalTok{)}
    \ControlFlowTok{return}\NormalTok{ Z }\OperatorTok{@}\NormalTok{ q}

\KeywordTok{def}\NormalTok{ sparse\_scores(query\_text):}
\NormalTok{    q }\OperatorTok{=}\NormalTok{ vectorize\_query(query\_text)}
\NormalTok{    Z }\OperatorTok{=}\NormalTok{ normalize\_rows(X\_tfidf)}
\NormalTok{    q }\OperatorTok{=}\NormalTok{ q }\OperatorTok{/} \BuiltInTok{max}\NormalTok{(np.linalg.norm(q), }\FloatTok{1e{-}12}\NormalTok{)}
    \ControlFlowTok{return}\NormalTok{ Z }\OperatorTok{@}\NormalTok{ q}

\KeywordTok{def}\NormalTok{ hybrid\_search(query\_text, query\_vec, lam}\OperatorTok{=}\FloatTok{0.5}\NormalTok{, k}\OperatorTok{=}\DecValTok{3}\NormalTok{):}
\NormalTok{    d }\OperatorTok{=}\NormalTok{ dense\_scores(query\_vec)}
\NormalTok{    s }\OperatorTok{=}\NormalTok{ sparse\_scores(query\_text)}
\NormalTok{    score }\OperatorTok{=}\NormalTok{ lam }\OperatorTok{*}\NormalTok{ d }\OperatorTok{+}\NormalTok{ (}\DecValTok{1} \OperatorTok{{-}}\NormalTok{ lam) }\OperatorTok{*}\NormalTok{ s}
\NormalTok{    order }\OperatorTok{=}\NormalTok{ np.argsort(score)[::}\OperatorTok{{-}}\DecValTok{1}\NormalTok{][:k]}
    \ControlFlowTok{return}\NormalTok{ pd.DataFrame(\{}
        \StringTok{"rank"}\NormalTok{: np.arange(}\DecValTok{1}\NormalTok{, }\BuiltInTok{len}\NormalTok{(order) }\OperatorTok{+} \DecValTok{1}\NormalTok{),}
        \StringTok{"chunk"}\NormalTok{: [chunks[i] }\ControlFlowTok{for}\NormalTok{ i }\KeywordTok{in}\NormalTok{ order],}
        \StringTok{"dense"}\NormalTok{: d[order].}\BuiltInTok{round}\NormalTok{(}\DecValTok{3}\NormalTok{),}
        \StringTok{"sparse"}\NormalTok{: s[order].}\BuiltInTok{round}\NormalTok{(}\DecValTok{3}\NormalTok{),}
        \StringTok{"hybrid"}\NormalTok{: score[order].}\BuiltInTok{round}\NormalTok{(}\DecValTok{3}\NormalTok{),}
\NormalTok{    \})}

\NormalTok{hybrid\_search(}
    \StringTok{"transfer credit approval form"}\NormalTok{,}
\NormalTok{    queries[}\StringTok{"How do I transfer credits?"}\NormalTok{],}
\NormalTok{    lam}\OperatorTok{=}\FloatTok{0.5}\NormalTok{,}
\NormalTok{    k}\OperatorTok{=}\DecValTok{4}\NormalTok{,}
\NormalTok{)}
\end{Highlighting}
\end{Shaded}

\protect\phantomsection\label{ch19-hybrid-search}
{\def\LTcaptype{none} % do not increment counter
\begin{longtable}[]{@{}llllll@{}}
\toprule\noalign{}
& rank & chunk & dense & sparse & hybrid \\
\midrule\noalign{}
\endhead
\bottomrule\noalign{}
\endlastfoot
0 & 1 & students may transfer graduate credits with ap... & 0.998 &
0.406 & 0.702 \\
1 & 2 & transfer credit forms require the program dire... & 1.000 &
0.383 & 0.691 \\
2 & 3 & applied mathematics students study probability... & 0.116 &
0.000 & 0.058 \\
3 & 4 & machine learning uses embeddings and optimization & 0.112 &
0.000 & 0.056 \\
\end{longtable}
}

\section{19.11 Reranking}\label{reranking}

A common retrieval pipeline has two stages:

\begin{enumerate}
\def\labelenumi{\arabic{enumi}.}
\tightlist
\item
  \textbf{Retriever:} quickly returns a candidate set \(C_k(q)\).
\item
  \textbf{Reranker:} applies a more expensive scoring function to
  reorder the candidates.
\end{enumerate}

Mathematically, the first stage solves an approximate top-\(k\) problem:

\[
C_k(q) \approx \operatorname{argtopk}_{c \in C} S_1(q,c).
\]

The reranker then solves

\[
\operatorname{argtopr}_{c \in C_k(q)} S_2(q,c),
\qquad r \leq k.
\]

The first score \(S_1\) is usually cheap, such as vector similarity. The
second score \(S_2\) may use a cross-encoder or a larger language model
that jointly reads the query and candidate chunk. This is more expensive
because the representation of the pair \((q,c)\) must be computed
jointly for each candidate.

\section{19.12 Chunking as a mathematical design
choice}\label{chunking-as-a-mathematical-design-choice}

Before indexing a document collection, long documents are usually split
into chunks. Let a document be a sequence

\[
d = (w_1,\dots,w_T).
\]

A chunking function maps \(d\) to shorter subsequences:

\[
g(d) = (c_1,\dots,c_M).
\]

For example, fixed-length chunking with length \(L\) and stride \(s\)
creates chunks

\[
c_j = (w_{1+(j-1)s},\dots,w_{L+(j-1)s}).
\]

Chunking is not merely a technical preprocessing step. It affects
retrieval quality.

\begin{itemize}
\tightlist
\item
  If chunks are too short, they may lack context.
\item
  If chunks are too long, they may contain irrelevant material and waste
  the model context window.
\item
  If chunks overlap too much, the index becomes redundant.
\item
  If chunks do not respect document structure, important definitions or
  assumptions may be separated.
\end{itemize}

In mathematical terms, chunking determines the finite candidate set
\(C\) over which retrieval is performed. A poor candidate set cannot be
fixed completely by a good similarity function.

\section{19.13 RAG as conditional
generation}\label{rag-as-conditional-generation}

A simple RAG pipeline is:

\[
q \xrightarrow{\text{embed}} z_q
\xrightarrow{\text{retrieve}} c_{i_1},\dots,c_{i_k}
\xrightarrow{\text{prompt}} x
\xrightarrow{\text{LM}} y.
\]

The context-augmented prompt may be written as

\[
x = \operatorname{concat}(q,c_{i_1},\dots,c_{i_k},I),
\]

where \(I\) denotes the instruction template. The generator then samples
or chooses

\[
y \sim p_\theta(\cdot \mid x).
\]

This clarifies an important point: retrieval does not force the
generator to be correct. Retrieval changes the conditioning information.
The final answer still depends on the conditional distribution learned
by the language model and on the decoding procedure.

\section{19.14 A probabilistic latent-document
view}\label{a-probabilistic-latent-document-view}

A more formal RAG model treats retrieved documents as latent variables.
Let \(C\) be the collection. Then

\[
p(y \mid q)
= \sum_{c \in C} p_\theta(y \mid q,c) p_\eta(c \mid q).
\]

In practice, summing over all \(c \in C\) is expensive, so we
approximate by the top-\(k\) retrieved chunks:

\[
p(y \mid q)
\approx
\sum_{c \in C_k(q)} p_\theta(y \mid q,c) p_\eta(c \mid q).
\]

If the retriever assigns scores \(s(q,c)\), a soft retrieval
distribution can be defined by

\[
p_\eta(c \mid q)
= \frac{\exp(s(q,c)/\tau)}{\sum_{c' \in C} \exp(s(q,c')/\tau)}.
\]

The temperature \(\tau\) controls concentration. A small \(\tau\)
concentrates probability on the top chunks. A large \(\tau\) spreads
probability across many chunks.

This formulation makes retrieval uncertainty explicit. If several chunks
have similar scores, then the system should be less confident about
which evidence is most relevant.

\section{19.15 Evaluating retrieval}\label{evaluating-retrieval}

Suppose for a query \(q\) we know a set of relevant documents
\(R(q) \subseteq C\). If a retriever returns an ordered list

\[
L_k(q) = (c_{i_1},\dots,c_{i_k}),
\]

then precision at \(k\) is

\[
\operatorname{Precision@}k(q)
= \frac{|L_k(q) \cap R(q)|}{k}.
\]

Recall at \(k\) is

\[
\operatorname{Recall@}k(q)
= \frac{|L_k(q) \cap R(q)|}{|R(q)|}.
\]

The reciprocal rank is

\[
\operatorname{RR}(q)
= \frac{1}{\min\{j:c_{i_j}\in R(q)\}},
\]

if a relevant result appears, and \(0\) otherwise. Mean reciprocal rank
averages this quantity over queries.

For graded relevance labels \(rel_j\), discounted cumulative gain is

\[
\operatorname{DCG@}k
= \sum_{j=1}^k \frac{2^{rel_j}-1}{\log_2(j+1)}.
\]

Normalized DCG divides by the best possible DCG:

\[
\operatorname{nDCG@}k
= \frac{\operatorname{DCG@}k}{\operatorname{IDCG@}k}.
\]

These metrics measure retrieval, not necessarily final answer quality. A
RAG system can retrieve the right document and still generate a poor
answer; it can also retrieve incomplete evidence and produce an answer
that sounds plausible but is unsupported.

\section{19.16 Python: retrieval
metrics}\label{python-retrieval-metrics}

\protect\phantomsection\label{ch19-retrieval-metrics}
\begin{Shaded}
\begin{Highlighting}[]
\NormalTok{retrieved }\OperatorTok{=}\NormalTok{ [}\StringTok{"c6"}\NormalTok{, }\StringTok{"c1"}\NormalTok{, }\StringTok{"c4"}\NormalTok{, }\StringTok{"c2"}\NormalTok{, }\StringTok{"c3"}\NormalTok{]}
\NormalTok{relevant }\OperatorTok{=}\NormalTok{ \{}\StringTok{"c1"}\NormalTok{, }\StringTok{"c6"}\NormalTok{\}}

\KeywordTok{def}\NormalTok{ precision\_at\_k(retrieved, relevant, k):}
\NormalTok{    top }\OperatorTok{=}\NormalTok{ retrieved[:k]}
    \ControlFlowTok{return} \BuiltInTok{sum}\NormalTok{(item }\KeywordTok{in}\NormalTok{ relevant }\ControlFlowTok{for}\NormalTok{ item }\KeywordTok{in}\NormalTok{ top) }\OperatorTok{/}\NormalTok{ k}

\KeywordTok{def}\NormalTok{ recall\_at\_k(retrieved, relevant, k):}
\NormalTok{    top }\OperatorTok{=}\NormalTok{ retrieved[:k]}
    \ControlFlowTok{return} \BuiltInTok{sum}\NormalTok{(item }\KeywordTok{in}\NormalTok{ relevant }\ControlFlowTok{for}\NormalTok{ item }\KeywordTok{in}\NormalTok{ top) }\OperatorTok{/} \BuiltInTok{len}\NormalTok{(relevant)}

\KeywordTok{def}\NormalTok{ reciprocal\_rank(retrieved, relevant):}
    \ControlFlowTok{for}\NormalTok{ j, item }\KeywordTok{in} \BuiltInTok{enumerate}\NormalTok{(retrieved, start}\OperatorTok{=}\DecValTok{1}\NormalTok{):}
        \ControlFlowTok{if}\NormalTok{ item }\KeywordTok{in}\NormalTok{ relevant:}
            \ControlFlowTok{return} \DecValTok{1} \OperatorTok{/}\NormalTok{ j}
    \ControlFlowTok{return} \FloatTok{0.0}

\KeywordTok{def}\NormalTok{ dcg\_at\_k(relevances, k):}
\NormalTok{    relevances }\OperatorTok{=}\NormalTok{ np.asarray(relevances[:k], dtype}\OperatorTok{=}\BuiltInTok{float}\NormalTok{)}
\NormalTok{    discounts }\OperatorTok{=}\NormalTok{ np.log2(np.arange(}\DecValTok{2}\NormalTok{, k }\OperatorTok{+} \DecValTok{2}\NormalTok{))}
    \ControlFlowTok{return}\NormalTok{ np.}\BuiltInTok{sum}\NormalTok{((}\DecValTok{2} \OperatorTok{**}\NormalTok{ relevances }\OperatorTok{{-}} \DecValTok{1}\NormalTok{) }\OperatorTok{/}\NormalTok{ discounts)}

\BuiltInTok{print}\NormalTok{(}\StringTok{"Precision@3:"}\NormalTok{, precision\_at\_k(retrieved, relevant, }\DecValTok{3}\NormalTok{))}
\BuiltInTok{print}\NormalTok{(}\StringTok{"Recall@3:"}\NormalTok{, recall\_at\_k(retrieved, relevant, }\DecValTok{3}\NormalTok{))}
\BuiltInTok{print}\NormalTok{(}\StringTok{"Reciprocal rank:"}\NormalTok{, reciprocal\_rank(retrieved, relevant))}

\NormalTok{relevance\_scores }\OperatorTok{=}\NormalTok{ [}\DecValTok{2}\NormalTok{, }\DecValTok{2}\NormalTok{, }\DecValTok{0}\NormalTok{, }\DecValTok{0}\NormalTok{, }\DecValTok{0}\NormalTok{]}
\NormalTok{ideal\_scores }\OperatorTok{=} \BuiltInTok{sorted}\NormalTok{(relevance\_scores, reverse}\OperatorTok{=}\VariableTok{True}\NormalTok{)}
\NormalTok{ndcg3 }\OperatorTok{=}\NormalTok{ dcg\_at\_k(relevance\_scores, }\DecValTok{3}\NormalTok{) }\OperatorTok{/}\NormalTok{ dcg\_at\_k(ideal\_scores, }\DecValTok{3}\NormalTok{)}
\BuiltInTok{print}\NormalTok{(}\StringTok{"nDCG@3:"}\NormalTok{, }\BuiltInTok{round}\NormalTok{(ndcg3, }\DecValTok{3}\NormalTok{))}
\end{Highlighting}
\end{Shaded}

\begin{verbatim}
Precision@3: 0.6666666666666666
Recall@3: 1.0
Reciprocal rank: 1.0
nDCG@3: 1.0
\end{verbatim}

\section{19.17 Approximate nearest-neighbor
search}\label{approximate-nearest-neighbor-search}

Exact retrieval by matrix multiplication requires computing \(Zz_q\)
over all \(n\) stored vectors. This costs \(O(nm)\) per query. For small
\(n\), this is fine. For millions or billions of chunks, exact search
may be too expensive.

Approximate nearest-neighbor methods build an index that returns a
candidate set quickly. The mathematical tradeoff is between:

\begin{itemize}
\tightlist
\item
  \textbf{speed:} how fast queries are answered;
\item
  \textbf{memory:} how large the index is;
\item
  \textbf{recall:} how often the true nearest neighbors are returned;
\item
  \textbf{update cost:} how hard it is to add or remove vectors.
\end{itemize}

At a high level, many methods use one of the following ideas.

\subsection{Partitioning}\label{partitioning}

Partition the vector space into regions. Search only regions near the
query. This reduces the number of candidate vectors.

\subsection{Graph search}\label{graph-search}

Build a graph whose vertices are stored vectors and whose edges connect
nearby vectors. Starting from an entry point, greedily move through the
graph toward vectors closer to the query.

\subsection{Quantization}\label{quantization}

Approximate vectors by compressed codes. Search is faster and uses less
memory, but distances are approximate.

For this book, the key point is that vector indexing is an algorithmic
approximation to the mathematical top-\(k\) problem.

\section{19.18 Worked calculation: cosine
retrieval}\label{worked-calculation-cosine-retrieval}

Suppose a query and three chunk embeddings are

\[
z_q = (1,1),
\qquad
z_1=(2,2),
\qquad
z_2=(2,0),
\qquad
z_3=(0,3).
\]

Dot-product scores are

\[
z_1^Tz_q=4,
\qquad
z_2^Tz_q=2,
\qquad
z_3^Tz_q=3.
\]

Thus dot-product retrieval ranks \(z_1,z_3,z_2\). Cosine scores are

\[
\cos(z_q,z_1)=1,
\]

\[
\cos(z_q,z_2)=\frac{2}{\sqrt{2}\cdot 2}=\frac{1}{\sqrt 2},
\]

\[
\cos(z_q,z_3)=\frac{3}{\sqrt{2}\cdot 3}=\frac{1}{\sqrt 2}.
\]

Under cosine similarity, \(z_2\) and \(z_3\) tie. The difference occurs
because dot product rewards the longer norm of \(z_3\).

\section{19.19 Interactive: dense retrieval and hybrid
scoring}\label{interactive-dense-retrieval-and-hybrid-scoring}

The interactive below uses a two-dimensional toy embedding space. Each
point is a document chunk. The query direction can be changed, and the
slider \(\lambda\) blends dense semantic score with sparse lexical
score:

\[
s_{\text{hybrid}} = \lambda s_{\text{dense}} + (1-\lambda)s_{\text{sparse}}.
\]

\section{19.20 Python Plotly version}\label{python-plotly-version-1}

The HTML interactive above is useful in the rendered book. The following
Python version is useful in a notebook.

\protect\phantomsection\label{ch19-python-plotly-retrieval-map}
\begin{Shaded}
\begin{Highlighting}[]
\ImportTok{import}\NormalTok{ plotly.graph\_objects }\ImportTok{as}\NormalTok{ go}

\NormalTok{points }\OperatorTok{=}\NormalTok{ toy\_embeddings[:, [}\DecValTok{0}\NormalTok{, }\DecValTok{1}\NormalTok{]]}
\NormalTok{query\_vec }\OperatorTok{=}\NormalTok{ queries[}\StringTok{"How do I transfer credits?"}\NormalTok{][[}\DecValTok{0}\NormalTok{, }\DecValTok{1}\NormalTok{]]}

\NormalTok{fig }\OperatorTok{=}\NormalTok{ go.Figure()}
\NormalTok{fig.add\_trace(go.Scatter(}
\NormalTok{    x}\OperatorTok{=}\NormalTok{points[:, }\DecValTok{0}\NormalTok{], y}\OperatorTok{=}\NormalTok{points[:, }\DecValTok{1}\NormalTok{],}
\NormalTok{    mode}\OperatorTok{=}\StringTok{"markers+text"}\NormalTok{,}
\NormalTok{    text}\OperatorTok{=}\NormalTok{[}\SpecialStringTok{f"chunk }\SpecialCharTok{\{}\NormalTok{i}\OperatorTok{+}\DecValTok{1}\SpecialCharTok{\}}\SpecialStringTok{"} \ControlFlowTok{for}\NormalTok{ i }\KeywordTok{in} \BuiltInTok{range}\NormalTok{(}\BuiltInTok{len}\NormalTok{(chunks))],}
\NormalTok{    textposition}\OperatorTok{=}\StringTok{"top center"}\NormalTok{,}
\NormalTok{    name}\OperatorTok{=}\StringTok{"chunks"}
\NormalTok{))}
\NormalTok{fig.add\_trace(go.Scatter(}
\NormalTok{    x}\OperatorTok{=}\NormalTok{[}\DecValTok{0}\NormalTok{, query\_vec[}\DecValTok{0}\NormalTok{]], y}\OperatorTok{=}\NormalTok{[}\DecValTok{0}\NormalTok{, query\_vec[}\DecValTok{1}\NormalTok{]],}
\NormalTok{    mode}\OperatorTok{=}\StringTok{"lines+markers+text"}\NormalTok{,}
\NormalTok{    text}\OperatorTok{=}\NormalTok{[}\StringTok{""}\NormalTok{, }\StringTok{"query"}\NormalTok{],}
\NormalTok{    textposition}\OperatorTok{=}\StringTok{"bottom right"}\NormalTok{,}
\NormalTok{    name}\OperatorTok{=}\StringTok{"query"}
\NormalTok{))}
\NormalTok{fig.update\_layout(}
\NormalTok{    title}\OperatorTok{=}\StringTok{"Toy retrieval embedding space"}\NormalTok{,}
\NormalTok{    xaxis\_title}\OperatorTok{=}\StringTok{"administration coordinate"}\NormalTok{,}
\NormalTok{    yaxis\_title}\OperatorTok{=}\StringTok{"math/ML coordinate"}
\NormalTok{)}
\NormalTok{fig.show()}
\end{Highlighting}
\end{Shaded}

\section{19.21 Error modes in RAG}\label{error-modes-in-rag}

RAG systems can fail for mathematically distinct reasons.

\subsection{Retrieval failure}\label{retrieval-failure}

The relevant chunk is not retrieved:

\[
R(q) \cap L_k(q) = \varnothing.
\]

This may happen because the embedding model fails to place the query
near the right document, because the relevant document was not indexed,
or because the chunking strategy split the answer from its context.

\subsection{Ranking failure}\label{ranking-failure}

A relevant chunk is retrieved, but ranked too low to be included in the
final context. This is a recall and ranking problem.

\subsection{Context overload}\label{context-overload}

Too many chunks are provided. The answer-relevant evidence may be
present, but the generator may attend to the wrong material or combine
incompatible chunks.

\subsection{Generation failure}\label{generation-failure}

The correct evidence is present, but the language model misreads it,
ignores it, or produces unsupported text.

\subsection{Citation failure}\label{citation-failure}

The answer may be mostly correct but cite the wrong chunk. This is a
mismatch between answer generation and evidence attribution.

\begin{tcolorbox}[enhanced jigsaw, arc=.35mm, bottomrule=.15mm, bottomtitle=1mm, breakable, colback=white, colbacktitle=quarto-callout-note-color!10!white, colframe=quarto-callout-note-color-frame, coltitle=black, left=2mm, leftrule=.75mm, opacityback=0, opacitybacktitle=0.6, rightrule=.15mm, title=\textcolor{quarto-callout-note-color}{\faInfo}\hspace{0.5em}{A useful diagnostic separation}, titlerule=0mm, toprule=.15mm, toptitle=1mm]

When debugging a RAG system, separate three questions:

\begin{enumerate}
\def\labelenumi{\arabic{enumi}.}
\tightlist
\item
  Did retrieval find the right evidence?
\item
  Did the prompt place the evidence in a usable form?
\item
  Did the generator faithfully use the evidence?
\end{enumerate}

These are different mathematical and engineering problems.

\end{tcolorbox}

\section{19.22 RAG for academic advising and course
resources}\label{rag-for-academic-advising-and-course-resources}

A university resource chatbot is a natural RAG application. The document
collection may include:

\begin{itemize}
\tightlist
\item
  program requirements;
\item
  course descriptions;
\item
  transfer-credit policies;
\item
  advising FAQs;
\item
  academic calendar pages;
\item
  graduate catalog entries;
\item
  official forms and procedures.
\end{itemize}

For such a system, mathematical care matters. The answer should be
grounded in retrieved official sources, and the system should
communicate uncertainty when retrieval is weak.

A minimal retrieval-aware answer generator should keep track of:

\[
\text{query},\quad
\text{retrieved chunks},\quad
\text{similarity scores},\quad
\text{source metadata},\quad
\text{answer}.
\]

The source metadata are essential. Without them, the system may generate
fluent but unverifiable responses.

\section{19.23 AI-assisted learning
activities}\label{ai-assisted-learning-activities-3}

Use an AI assistant as a mathematical critic, not only as a code
generator.

\subsection{Activity 1: explain the scoring
geometry}\label{activity-1-explain-the-scoring-geometry}

Ask an AI assistant:

\begin{quote}
In dense retrieval, explain the difference between dot-product
similarity and cosine similarity. Give a two-dimensional example where
they rank documents differently.
\end{quote}

Then verify the example by hand.

\subsection{Activity 2: audit a RAG
pipeline}\label{activity-2-audit-a-rag-pipeline}

Give the assistant a proposed RAG design and ask:

\begin{quote}
Identify possible retrieval failure modes, ranking failure modes, and
generation failure modes. For each one, suggest a quantitative
diagnostic.
\end{quote}

Check whether the suggested metrics actually correspond to the failure
modes.

\subsection{Activity 3: create adversarial
queries}\label{activity-3-create-adversarial-queries}

Ask the assistant to produce queries that are likely to confuse a
retriever:

\begin{itemize}
\tightlist
\item
  queries with synonyms;
\item
  queries with exact codes or course numbers;
\item
  queries requiring negation;
\item
  queries requiring date-sensitive information;
\item
  queries requiring multiple documents.
\end{itemize}

Then test whether dense, sparse, or hybrid retrieval handles each case
better.

\section{19.24 Common mistakes}\label{common-mistakes-8}

\begin{enumerate}
\def\labelenumi{\arabic{enumi}.}
\tightlist
\item
  \textbf{Treating similarity scores as probabilities.} A cosine score
  is not automatically a calibrated probability of relevance.
\item
  \textbf{Ignoring chunking.} Retrieval quality depends on the candidate
  units.
\item
  \textbf{Using only dense retrieval.} Exact keywords, identifiers, and
  policy names often require lexical matching.
\item
  \textbf{Assuming retrieved means true.} Retrieved text may be
  outdated, irrelevant, or misinterpreted.
\item
  \textbf{Evaluating only final answers.} Retrieval metrics and
  generation metrics should be separated.
\item
  \textbf{Overloading the context window.} More retrieved text is not
  always better.
\end{enumerate}

\section{19.25 Exercises}\label{exercises-18}

\subsection{Exercise 1: normalized dot
product}\label{exercise-1-normalized-dot-product}

Prove that for unit vectors \(u\) and \(v\),

\[
\|u-v\|_2^2 = 2 - 2u^Tv.
\]

Explain why nearest-neighbor search by Euclidean distance is equivalent
to maximum cosine similarity on the unit sphere.

\subsection{Exercise 2: hybrid score}\label{exercise-2-hybrid-score}

Suppose two chunks have scores

\[
(s_{\text{dense}},s_{\text{sparse}})_1=(0.9,0.1),
\]

and

\[
(s_{\text{dense}},s_{\text{sparse}})_2=(0.6,0.8).
\]

For which values of \(\lambda\) does chunk 1 rank above chunk 2 under

\[
s_{\text{hybrid}}=\lambda s_{\text{dense}}+(1-\lambda)s_{\text{sparse}}?
\]

\subsection{Exercise 3: retrieval
metrics}\label{exercise-3-retrieval-metrics}

For a query, the returned list is

\[
(c_3,c_5,c_1,c_2,c_4),
\]

and the relevant set is

\[
\{c_1,c_4,c_6\}.
\]

Compute Precision@3, Recall@3, and reciprocal rank.

\subsection{Exercise 4: chunking
design}\label{exercise-4-chunking-design}

Consider a 4000-token document. You use chunk length \(L=500\) and
stride \(s=250\). How many chunks are produced if every chunk must fit
inside the document? What changes if the last partial chunk is kept?

\subsection{Exercise 5: latent-document
RAG}\label{exercise-5-latent-document-rag}

Starting from

\[
p(y \mid q)=\sum_{c \in C}p_\theta(y\mid q,c)p_\eta(c\mid q),
\]

explain why top-\(k\) truncation can introduce bias. When might the
approximation be reasonable?

\section{19.26 Mini-project: build a small mathematical RAG
retriever}\label{mini-project-build-a-small-mathematical-rag-retriever}

Build a small retriever for a collection of course or program documents.

Minimum requirements:

\begin{enumerate}
\def\labelenumi{\arabic{enumi}.}
\tightlist
\item
  Split documents into chunks.
\item
  Implement sparse TF-IDF retrieval.
\item
  Implement dense retrieval using either a toy embedding, a pretrained
  sentence embedding model, or an API embedding model.
\item
  Implement hybrid scoring.
\item
  Create at least 10 test queries with known relevant chunks.
\item
  Report Precision@3, Recall@3, and MRR.
\item
  Analyze at least three failure cases.
\end{enumerate}

A strong project should also include reranking, source metadata, and a
short discussion of how retrieval affects generated answers.

\section{19.27 Lab connection}\label{lab-connection-8}

The companion lab for this chapter should guide students through:

\begin{itemize}
\tightlist
\item
  constructing a document collection;
\item
  chunking text;
\item
  building TF-IDF vectors;
\item
  building dense vectors;
\item
  comparing cosine and dot-product retrieval;
\item
  implementing hybrid search;
\item
  computing retrieval metrics;
\item
  connecting retrieved chunks to a simple answer-generation prompt.
\end{itemize}

The lab should emphasize that RAG is not magic. It is a pipeline made of
mathematical components: representation maps, similarity functions,
ranking rules, and conditional generation.

\section{19.28 Summary}\label{summary-3}

Retrieval-augmented generation combines language modeling with search.
The retrieval part is fundamentally geometric: map queries and chunks to
vectors, compute similarities, and return top-ranked evidence. The
generation part is probabilistic: condition a language model on the
query and retrieved context.

The main mathematical objects in this chapter were:

\[
f_\theta:T \to R^m,
\]

\[
S(z_q,z_c),
\]

\[
\operatorname{TopK}(q)=\operatorname{argtopk}_{c\in C}S(f_\theta(q),f_\theta(c)),
\]

and

\[
p(y\mid q)\approx p_\theta(y\mid q,r_k(q)).
\]

Understanding these objects helps us design better retrieval systems,
evaluate them separately from generation, and diagnose failure modes in
applied NLP systems.

\section{19.29 Further reading}\label{further-reading}

The original retrieval-augmented generation framework is due to Lewis et
al. ({Lewis et al.} 2020). For approximate nearest-neighbor search and
vector indexing, see Johnson, Douze, and Jegou (Johnson et al. 2019).
For background on vector space models and embeddings, see earlier
chapters and the references on latent semantic analysis, Word2Vec,
GloVe, and Transformers.

\chapter{Chapter 20. Evaluation of Language
Models}\label{chapter-20.-evaluation-of-language-models}

Metrics, Uncertainty, Calibration, and Responsible Measurement

\hfill\break

\chapter{Chapter 20. Evaluation of Language
Models}\label{chapter-20.-evaluation-of-language-models-1}

Evaluation is the mathematical discipline of asking whether a language
system is doing what we want it to do. For a small classifier,
evaluation may mean accuracy on a labeled test set. For a language
model, evaluation may mean perplexity on held-out text. For a
summarizer, evaluation may compare generated summaries with reference
summaries. For a retrieval-augmented generation system, evaluation must
include both retrieval quality and answer quality. For a deployed
assistant, evaluation must also consider calibration, robustness,
factuality, bias, latency, privacy, and human usefulness.

The purpose of this chapter is not to collect many disconnected scores.
The goal is to show that most evaluation methods have a common
statistical structure:

\[
\text{model} \quad \longrightarrow \quad \text{predictions} \quad \longrightarrow \quad \text{loss or score} \quad \longrightarrow \quad \text{statistical estimate}.
\]

For MA Applied Mathematics and MS Statistics students, this chapter
should feel familiar. Evaluation is estimation. A reported score is
usually a sample average, a ratio of random counts, or a function of
empirical distributions. Therefore, every evaluation score has
assumptions, variance, failure modes, and a domain of validity.

\begin{tcolorbox}[enhanced jigsaw, arc=.35mm, bottomrule=.15mm, bottomtitle=1mm, breakable, colback=white, colbacktitle=quarto-callout-tip-color!10!white, colframe=quarto-callout-tip-color-frame, coltitle=black, left=2mm, leftrule=.75mm, opacityback=0, opacitybacktitle=0.6, rightrule=.15mm, title=\textcolor{quarto-callout-tip-color}{\faLightbulb}\hspace{0.5em}{Guiding principle}, titlerule=0mm, toprule=.15mm, toptitle=1mm]

A metric is not a property of a model alone. It is a property of a
model, a data distribution, a task definition, an annotation procedure,
and a scoring rule.

\end{tcolorbox}

\section{20.1 Learning goals}\label{learning-goals-19}

After studying this chapter, you should be able to:

\begin{enumerate}
\def\labelenumi{\arabic{enumi}.}
\tightlist
\item
  express NLP evaluation as estimation of expected loss or expected
  utility;
\item
  compute accuracy, precision, recall, \(F_1\), and confusion matrices
  from first principles;
\item
  derive cross-entropy and perplexity for language models;
\item
  explain calibration, Brier score, and expected calibration error;
\item
  distinguish token-level, sequence-level, retrieval-level, and
  human-centered evaluation;
\item
  implement overlap-based metrics such as ROUGE-style recall and
  BLEU-style precision;
\item
  evaluate retrieval and RAG pipelines using Precision@\(k\),
  Recall@\(k\), MRR, and nDCG;
\item
  quantify uncertainty in evaluation using standard errors and bootstrap
  confidence intervals;
\item
  critique metric choice using mathematical reasoning.
\end{enumerate}

\section{20.2 The mathematical form of
evaluation}\label{the-mathematical-form-of-evaluation}

Let \((X,Y)\) be a random input-output pair drawn from an unknown
population distribution \(P\). A model \(f_\theta\) maps an input \(x\)
to either a prediction, a score vector, a probability distribution, or a
generated sequence. A loss function compares the model output to the
target:

\[
L(f_\theta(X),Y).
\]

The population risk is

\[
R(\theta) = E_P[L(f_\theta(X),Y)].
\]

Since \(P\) is unknown, we estimate risk from a test set

\[
D_{\text{test}} = \{(x_i,y_i)\}_{i=1}^n.
\]

The empirical risk is

\[
\widehat R(\theta) = \frac{1}{n}\sum_{i=1}^n L(f_\theta(x_i),y_i).
\]

This formula is the mathematical backbone of evaluation. Different NLP
tasks use different forms of \(X\), \(Y\), \(f_\theta\), and \(L\).

{\def\LTcaptype{none} % do not increment counter
\begin{longtable}[]{@{}
  >{\raggedright\arraybackslash}p{(\linewidth - 8\tabcolsep) * \real{0.1579}}
  >{\raggedleft\arraybackslash}p{(\linewidth - 8\tabcolsep) * \real{0.2105}}
  >{\raggedleft\arraybackslash}p{(\linewidth - 8\tabcolsep) * \real{0.2105}}
  >{\raggedleft\arraybackslash}p{(\linewidth - 8\tabcolsep) * \real{0.2105}}
  >{\raggedleft\arraybackslash}p{(\linewidth - 8\tabcolsep) * \real{0.2105}}@{}}
\toprule\noalign{}
\begin{minipage}[b]{\linewidth}\raggedright
Task
\end{minipage} & \begin{minipage}[b]{\linewidth}\raggedleft
Input \(X\)
\end{minipage} & \begin{minipage}[b]{\linewidth}\raggedleft
Target \(Y\)
\end{minipage} & \begin{minipage}[b]{\linewidth}\raggedleft
Model output
\end{minipage} & \begin{minipage}[b]{\linewidth}\raggedleft
Example score
\end{minipage} \\
\midrule\noalign{}
\endhead
\bottomrule\noalign{}
\endlastfoot
Sentiment classification & document & class label & class probability &
accuracy, \(F_1\) \\
Language modeling & prefix & next token & categorical distribution &
cross-entropy, perplexity \\
Summarization & document & reference summary & generated summary &
ROUGE, human rating \\
Retrieval & query & relevant documents & ranked list & Recall@\(k\),
MRR \\
RAG answering & query plus corpus & answer & generated answer with
evidence & faithfulness, answer correctness \\
\end{longtable}
}

The important statistical question is:

\[
\widehat R(\theta) \approx R(\theta)?
\]

This approximation depends on sample size, test-set quality,
distribution shift, and whether the metric faithfully measures the
desired behavior.

\section{20.3 Data splits and the danger of
leakage}\label{data-splits-and-the-danger-of-leakage}

A standard supervised workflow uses three data sets:

\[
D_{\text{train}}, \qquad D_{\text{validation}}, \qquad D_{\text{test}}.
\]

The training set is used to fit parameters. The validation set is used
to choose hyperparameters, thresholds, prompts, decoding strategies,
retrieval settings, and checkpoints. The test set is used only once for
final evaluation.

If decisions are repeatedly made using the test set, then the test score
becomes optimistically biased. This is a form of selection bias. For
example, suppose we try many models \(\theta_1,\ldots,\theta_m\) and
report the best test score. Even if all models have the same population
risk, the minimum empirical risk

\[
\min_{1 \le j \le m} \widehat R(\theta_j)
\]

will usually be smaller than the true risk of the selected model.

\begin{tcolorbox}[enhanced jigsaw, arc=.35mm, bottomrule=.15mm, bottomtitle=1mm, breakable, colback=white, colbacktitle=quarto-callout-warning-color!10!white, colframe=quarto-callout-warning-color-frame, coltitle=black, left=2mm, leftrule=.75mm, opacityback=0, opacitybacktitle=0.6, rightrule=.15mm, title=\textcolor{quarto-callout-warning-color}{\faExclamationTriangle}\hspace{0.5em}{Test-set leakage}, titlerule=0mm, toprule=.15mm, toptitle=1mm]

A test set is not a data bank for repeated design decisions. It is an
estimator of future performance. Once it guides design choices, it
starts functioning like a validation set.

\end{tcolorbox}

\section{20.4 Classification metrics}\label{classification-metrics}

Consider a classification task with classes \(\{1,\ldots,K\}\). Let
\(y_i\) be the true class and \(\hat y_i\) be the predicted class. The
accuracy is

\[
\operatorname{Accuracy}=\frac{1}{n}\sum_{i=1}^n \mathbf 1\{\hat y_i=y_i\}.
\]

Accuracy is simple and useful when classes are balanced and errors have
similar cost. However, for imbalanced tasks, accuracy can be misleading.
A classifier that always predicts the majority class may have high
accuracy but low usefulness.

\subsection{Binary confusion matrix}\label{binary-confusion-matrix}

For binary classification, let the positive class be \(1\) and the
negative class be \(0\). The confusion counts are:

\[
TP = \sum_i \mathbf 1\{y_i=1,\hat y_i=1\},
\]

\[
FP = \sum_i \mathbf 1\{y_i=0,\hat y_i=1\},
\]

\[
FN = \sum_i \mathbf 1\{y_i=1,\hat y_i=0\},
\]

\[
TN = \sum_i \mathbf 1\{y_i=0,\hat y_i=0\}.
\]

Precision measures how reliable positive predictions are:

\[
\operatorname{Precision}=\frac{TP}{TP+FP}.
\]

Recall measures how many true positives are found:

\[
\operatorname{Recall}=\frac{TP}{TP+FN}.
\]

The \(F_1\) score is the harmonic mean:

\[
F_1 = \frac{2\cdot \operatorname{Precision}\cdot \operatorname{Recall}}{\operatorname{Precision}+\operatorname{Recall}}.
\]

The harmonic mean is small if either precision or recall is small. Thus
\(F_1\) rewards balance.

\subsection{Thresholds and decision
rules}\label{thresholds-and-decision-rules}

Many classifiers output a score \(s(x)\) or probability estimate
\(p_\theta(y=1\mid x)\). A threshold \(\tau\) produces a hard
prediction:

\[
\hat y_\tau(x)=\mathbf 1\{p_\theta(y=1\mid x)\ge \tau\}.
\]

Changing \(\tau\) changes \(TP\), \(FP\), \(FN\), and \(TN\). Therefore
precision, recall, and \(F_1\) are not properties of the probability
model alone. They are properties of the probability model plus the
threshold.

\begin{tcolorbox}[enhanced jigsaw, arc=.35mm, bottomrule=.15mm, bottomtitle=1mm, breakable, colback=white, colbacktitle=quarto-callout-note-color!10!white, colframe=quarto-callout-note-color-frame, coltitle=black, left=2mm, leftrule=.75mm, opacityback=0, opacitybacktitle=0.6, rightrule=.15mm, title=\textcolor{quarto-callout-note-color}{\faInfo}\hspace{0.5em}{Mathematical viewpoint}, titlerule=0mm, toprule=.15mm, toptitle=1mm]

A threshold is a decision rule. The same probabilistic model can produce
many classifiers depending on how probabilities are converted into
actions.

\end{tcolorbox}

\section{20.5 Python example: classification metrics from
scratch}\label{python-example-classification-metrics-from-scratch}

\begin{Shaded}
\begin{Highlighting}[]
\ImportTok{import}\NormalTok{ numpy }\ImportTok{as}\NormalTok{ np}

\CommentTok{\# A small binary text{-}classification evaluation example.}
\NormalTok{y\_true }\OperatorTok{=}\NormalTok{ np.array([}\DecValTok{1}\NormalTok{, }\DecValTok{1}\NormalTok{, }\DecValTok{0}\NormalTok{, }\DecValTok{1}\NormalTok{, }\DecValTok{0}\NormalTok{, }\DecValTok{0}\NormalTok{, }\DecValTok{1}\NormalTok{, }\DecValTok{0}\NormalTok{, }\DecValTok{1}\NormalTok{, }\DecValTok{0}\NormalTok{])}
\NormalTok{p\_score }\OperatorTok{=}\NormalTok{ np.array([}\FloatTok{0.92}\NormalTok{, }\FloatTok{0.75}\NormalTok{, }\FloatTok{0.62}\NormalTok{, }\FloatTok{0.55}\NormalTok{, }\FloatTok{0.48}\NormalTok{, }\FloatTok{0.40}\NormalTok{, }\FloatTok{0.36}\NormalTok{, }\FloatTok{0.30}\NormalTok{, }\FloatTok{0.22}\NormalTok{, }\FloatTok{0.10}\NormalTok{])}

\KeywordTok{def}\NormalTok{ binary\_metrics(y\_true, p\_score, threshold}\OperatorTok{=}\FloatTok{0.5}\NormalTok{):}
\NormalTok{    y\_pred }\OperatorTok{=}\NormalTok{ (p\_score }\OperatorTok{\textgreater{}=}\NormalTok{ threshold).astype(}\BuiltInTok{int}\NormalTok{)}
\NormalTok{    tp }\OperatorTok{=} \BuiltInTok{int}\NormalTok{(np.}\BuiltInTok{sum}\NormalTok{((y\_true }\OperatorTok{==} \DecValTok{1}\NormalTok{) }\OperatorTok{\&}\NormalTok{ (y\_pred }\OperatorTok{==} \DecValTok{1}\NormalTok{)))}
\NormalTok{    fp }\OperatorTok{=} \BuiltInTok{int}\NormalTok{(np.}\BuiltInTok{sum}\NormalTok{((y\_true }\OperatorTok{==} \DecValTok{0}\NormalTok{) }\OperatorTok{\&}\NormalTok{ (y\_pred }\OperatorTok{==} \DecValTok{1}\NormalTok{)))}
\NormalTok{    fn }\OperatorTok{=} \BuiltInTok{int}\NormalTok{(np.}\BuiltInTok{sum}\NormalTok{((y\_true }\OperatorTok{==} \DecValTok{1}\NormalTok{) }\OperatorTok{\&}\NormalTok{ (y\_pred }\OperatorTok{==} \DecValTok{0}\NormalTok{)))}
\NormalTok{    tn }\OperatorTok{=} \BuiltInTok{int}\NormalTok{(np.}\BuiltInTok{sum}\NormalTok{((y\_true }\OperatorTok{==} \DecValTok{0}\NormalTok{) }\OperatorTok{\&}\NormalTok{ (y\_pred }\OperatorTok{==} \DecValTok{0}\NormalTok{)))}
\NormalTok{    accuracy }\OperatorTok{=}\NormalTok{ (tp }\OperatorTok{+}\NormalTok{ tn) }\OperatorTok{/} \BuiltInTok{len}\NormalTok{(y\_true)}
\NormalTok{    precision }\OperatorTok{=}\NormalTok{ tp }\OperatorTok{/}\NormalTok{ (tp }\OperatorTok{+}\NormalTok{ fp) }\ControlFlowTok{if}\NormalTok{ tp }\OperatorTok{+}\NormalTok{ fp }\OperatorTok{\textgreater{}} \DecValTok{0} \ControlFlowTok{else} \FloatTok{0.0}
\NormalTok{    recall }\OperatorTok{=}\NormalTok{ tp }\OperatorTok{/}\NormalTok{ (tp }\OperatorTok{+}\NormalTok{ fn) }\ControlFlowTok{if}\NormalTok{ tp }\OperatorTok{+}\NormalTok{ fn }\OperatorTok{\textgreater{}} \DecValTok{0} \ControlFlowTok{else} \FloatTok{0.0}
\NormalTok{    f1 }\OperatorTok{=} \DecValTok{2} \OperatorTok{*}\NormalTok{ precision }\OperatorTok{*}\NormalTok{ recall }\OperatorTok{/}\NormalTok{ (precision }\OperatorTok{+}\NormalTok{ recall) }\ControlFlowTok{if}\NormalTok{ precision }\OperatorTok{+}\NormalTok{ recall }\OperatorTok{\textgreater{}} \DecValTok{0} \ControlFlowTok{else} \FloatTok{0.0}
    \ControlFlowTok{return}\NormalTok{ \{}
        \StringTok{"threshold"}\NormalTok{: threshold,}
        \StringTok{"TP"}\NormalTok{: tp,}
        \StringTok{"FP"}\NormalTok{: fp,}
        \StringTok{"FN"}\NormalTok{: fn,}
        \StringTok{"TN"}\NormalTok{: tn,}
        \StringTok{"accuracy"}\NormalTok{: accuracy,}
        \StringTok{"precision"}\NormalTok{: precision,}
        \StringTok{"recall"}\NormalTok{: recall,}
        \StringTok{"f1"}\NormalTok{: f1,}
\NormalTok{    \}}

\ControlFlowTok{for}\NormalTok{ tau }\KeywordTok{in}\NormalTok{ [}\FloatTok{0.3}\NormalTok{, }\FloatTok{0.5}\NormalTok{, }\FloatTok{0.7}\NormalTok{]:}
    \BuiltInTok{print}\NormalTok{(binary\_metrics(y\_true, p\_score, tau))}
\end{Highlighting}
\end{Shaded}

\begin{verbatim}
{'threshold': 0.3, 'TP': 4, 'FP': 4, 'FN': 1, 'TN': 1, 'accuracy': 0.5, 'precision': 0.5, 'recall': 0.8, 'f1': 0.6153846153846154}
{'threshold': 0.5, 'TP': 3, 'FP': 1, 'FN': 2, 'TN': 4, 'accuracy': 0.7, 'precision': 0.75, 'recall': 0.6, 'f1': 0.6666666666666665}
{'threshold': 0.7, 'TP': 2, 'FP': 0, 'FN': 3, 'TN': 5, 'accuracy': 0.7, 'precision': 1.0, 'recall': 0.4, 'f1': 0.5714285714285715}
\end{verbatim}

\section{20.6 Macro, micro, and weighted
averaging}\label{macro-micro-and-weighted-averaging}

For multiclass classification, we can compute precision, recall, and
\(F_1\) separately for each class. For class \(c\), treat \(c\) as
positive and all other classes as negative. This gives \(F_{1,c}\).

Macro averaging gives equal weight to each class:

\[
F_{1,\text{macro}} = \frac{1}{K}\sum_{c=1}^K F_{1,c}.
\]

Weighted averaging weights each class by its support \(n_c\):

\[
F_{1,\text{weighted}} = \sum_{c=1}^K \frac{n_c}{n}F_{1,c}.
\]

Micro averaging pools all true positives, false positives, and false
negatives across classes before computing the score. For single-label
multiclass classification, micro \(F_1\) equals accuracy.

The choice matters. Macro \(F_1\) emphasizes minority classes. Weighted
\(F_1\) is closer to the empirical class distribution. Micro metrics
emphasize overall instance-level performance.

\begin{Shaded}
\begin{Highlighting}[]
\NormalTok{y\_true\_mc }\OperatorTok{=}\NormalTok{ np.array([}\StringTok{"math"}\NormalTok{, }\StringTok{"math"}\NormalTok{, }\StringTok{"stats"}\NormalTok{, }\StringTok{"stats"}\NormalTok{, }\StringTok{"cs"}\NormalTok{, }\StringTok{"math"}\NormalTok{, }\StringTok{"cs"}\NormalTok{, }\StringTok{"stats"}\NormalTok{])}
\NormalTok{y\_pred\_mc }\OperatorTok{=}\NormalTok{ np.array([}\StringTok{"math"}\NormalTok{, }\StringTok{"stats"}\NormalTok{, }\StringTok{"stats"}\NormalTok{, }\StringTok{"stats"}\NormalTok{, }\StringTok{"math"}\NormalTok{, }\StringTok{"math"}\NormalTok{, }\StringTok{"cs"}\NormalTok{, }\StringTok{"cs"}\NormalTok{])}
\NormalTok{classes }\OperatorTok{=} \BuiltInTok{sorted}\NormalTok{(}\BuiltInTok{set}\NormalTok{(y\_true\_mc) }\OperatorTok{|} \BuiltInTok{set}\NormalTok{(y\_pred\_mc))}

\KeywordTok{def}\NormalTok{ per\_class\_f1(y\_true, y\_pred, classes):}
\NormalTok{    rows }\OperatorTok{=}\NormalTok{ []}
    \ControlFlowTok{for}\NormalTok{ c }\KeywordTok{in}\NormalTok{ classes:}
\NormalTok{        tp }\OperatorTok{=}\NormalTok{ np.}\BuiltInTok{sum}\NormalTok{((y\_true }\OperatorTok{==}\NormalTok{ c) }\OperatorTok{\&}\NormalTok{ (y\_pred }\OperatorTok{==}\NormalTok{ c))}
\NormalTok{        fp }\OperatorTok{=}\NormalTok{ np.}\BuiltInTok{sum}\NormalTok{((y\_true }\OperatorTok{!=}\NormalTok{ c) }\OperatorTok{\&}\NormalTok{ (y\_pred }\OperatorTok{==}\NormalTok{ c))}
\NormalTok{        fn }\OperatorTok{=}\NormalTok{ np.}\BuiltInTok{sum}\NormalTok{((y\_true }\OperatorTok{==}\NormalTok{ c) }\OperatorTok{\&}\NormalTok{ (y\_pred }\OperatorTok{!=}\NormalTok{ c))}
\NormalTok{        precision }\OperatorTok{=}\NormalTok{ tp }\OperatorTok{/}\NormalTok{ (tp }\OperatorTok{+}\NormalTok{ fp) }\ControlFlowTok{if}\NormalTok{ tp }\OperatorTok{+}\NormalTok{ fp }\OperatorTok{\textgreater{}} \DecValTok{0} \ControlFlowTok{else} \FloatTok{0.0}
\NormalTok{        recall }\OperatorTok{=}\NormalTok{ tp }\OperatorTok{/}\NormalTok{ (tp }\OperatorTok{+}\NormalTok{ fn) }\ControlFlowTok{if}\NormalTok{ tp }\OperatorTok{+}\NormalTok{ fn }\OperatorTok{\textgreater{}} \DecValTok{0} \ControlFlowTok{else} \FloatTok{0.0}
\NormalTok{        f1 }\OperatorTok{=} \DecValTok{2} \OperatorTok{*}\NormalTok{ precision }\OperatorTok{*}\NormalTok{ recall }\OperatorTok{/}\NormalTok{ (precision }\OperatorTok{+}\NormalTok{ recall) }\ControlFlowTok{if}\NormalTok{ precision }\OperatorTok{+}\NormalTok{ recall }\OperatorTok{\textgreater{}} \DecValTok{0} \ControlFlowTok{else} \FloatTok{0.0}
\NormalTok{        support }\OperatorTok{=} \BuiltInTok{int}\NormalTok{(np.}\BuiltInTok{sum}\NormalTok{(y\_true }\OperatorTok{==}\NormalTok{ c))}
\NormalTok{        rows.append((c, support, precision, recall, f1))}
    \ControlFlowTok{return}\NormalTok{ rows}

\NormalTok{rows }\OperatorTok{=}\NormalTok{ per\_class\_f1(y\_true\_mc, y\_pred\_mc, classes)}
\ControlFlowTok{for}\NormalTok{ row }\KeywordTok{in}\NormalTok{ rows:}
    \BuiltInTok{print}\NormalTok{(row)}

\NormalTok{macro\_f1 }\OperatorTok{=}\NormalTok{ np.mean([row[}\OperatorTok{{-}}\DecValTok{1}\NormalTok{] }\ControlFlowTok{for}\NormalTok{ row }\KeywordTok{in}\NormalTok{ rows])}
\NormalTok{weighted\_f1 }\OperatorTok{=} \BuiltInTok{sum}\NormalTok{(row[}\DecValTok{1}\NormalTok{] }\OperatorTok{*}\NormalTok{ row[}\OperatorTok{{-}}\DecValTok{1}\NormalTok{] }\ControlFlowTok{for}\NormalTok{ row }\KeywordTok{in}\NormalTok{ rows) }\OperatorTok{/} \BuiltInTok{len}\NormalTok{(y\_true\_mc)}
\NormalTok{accuracy }\OperatorTok{=}\NormalTok{ np.mean(y\_true\_mc }\OperatorTok{==}\NormalTok{ y\_pred\_mc)}
\BuiltInTok{print}\NormalTok{(}\StringTok{"macro F1:"}\NormalTok{, macro\_f1)}
\BuiltInTok{print}\NormalTok{(}\StringTok{"weighted F1:"}\NormalTok{, weighted\_f1)}
\BuiltInTok{print}\NormalTok{(}\StringTok{"accuracy:"}\NormalTok{, accuracy)}
\end{Highlighting}
\end{Shaded}

\begin{verbatim}
('cs', 2, 0.5, 0.5, 0.5)
('math', 3, 0.6666666666666666, 0.6666666666666666, 0.6666666666666666)
('stats', 3, 0.6666666666666666, 0.6666666666666666, 0.6666666666666666)
macro F1: 0.611111111111111
weighted F1: 0.625
accuracy: 0.625
\end{verbatim}

\section{20.7 Language-model evaluation: cross-entropy and
perplexity}\label{language-model-evaluation-cross-entropy-and-perplexity}

For a language model, the target is usually the next token. Given a
sequence \(w_{1:T}\), an autoregressive model assigns

\[
p_\theta(w_{1:T}) = \prod_{t=1}^T p_\theta(w_t \mid w_{<t}).
\]

The average negative log-likelihood per token is

\[
\widehat H = -\frac{1}{N}\sum_{i=1}^N \log p_\theta(w_i \mid \text{context}_i),
\]

where \(N\) is the number of evaluated target tokens. This quantity is
also called empirical cross-entropy when the tokens are sampled from a
data distribution.

Perplexity is

\[
\operatorname{PPL} = \exp(\widehat H)
\]

when natural logarithms are used. If logarithms base \(2\) are used,
then

\[
\operatorname{PPL}=2^{\widehat H_2}.
\]

Perplexity can be interpreted as an effective average branching factor.
A lower perplexity means the model assigns higher probability to the
observed next tokens.

\subsection{Important limitations of
perplexity}\label{important-limitations-of-perplexity}

Perplexity is useful for comparing models under the same tokenizer and
evaluation distribution. However, it can be misleading across
tokenizers. A model using shorter subword units may have a different
number of tokens than a model using longer tokens. Per-token perplexity
is therefore not automatically comparable across tokenization schemes.

Also, low perplexity does not guarantee usefulness, factuality, safety,
or reasoning quality. A model may predict common text well but fail on
specialized mathematical or institutional questions.

\begin{Shaded}
\begin{Highlighting}[]
\CommentTok{\# Token{-}level negative log{-}likelihood and perplexity.}
\CommentTok{\# Rows are time steps; columns are vocabulary items.}
\NormalTok{logits }\OperatorTok{=}\NormalTok{ np.array([}
\NormalTok{    [}\FloatTok{2.0}\NormalTok{, }\FloatTok{1.0}\NormalTok{, }\FloatTok{0.0}\NormalTok{],}
\NormalTok{    [}\FloatTok{0.2}\NormalTok{, }\FloatTok{1.8}\NormalTok{, }\FloatTok{0.1}\NormalTok{],}
\NormalTok{    [}\FloatTok{0.1}\NormalTok{, }\FloatTok{0.4}\NormalTok{, }\FloatTok{2.2}\NormalTok{],}
\NormalTok{    [}\FloatTok{1.4}\NormalTok{, }\FloatTok{0.8}\NormalTok{, }\FloatTok{0.2}\NormalTok{],}
\NormalTok{])}

\NormalTok{targets }\OperatorTok{=}\NormalTok{ np.array([}\DecValTok{0}\NormalTok{, }\DecValTok{1}\NormalTok{, }\DecValTok{2}\NormalTok{, }\DecValTok{1}\NormalTok{])}

\KeywordTok{def}\NormalTok{ log\_softmax(z):}
\NormalTok{    z }\OperatorTok{=}\NormalTok{ z }\OperatorTok{{-}}\NormalTok{ np.}\BuiltInTok{max}\NormalTok{(z, axis}\OperatorTok{={-}}\DecValTok{1}\NormalTok{, keepdims}\OperatorTok{=}\VariableTok{True}\NormalTok{)}
    \ControlFlowTok{return}\NormalTok{ z }\OperatorTok{{-}}\NormalTok{ np.log(np.}\BuiltInTok{sum}\NormalTok{(np.exp(z), axis}\OperatorTok{={-}}\DecValTok{1}\NormalTok{, keepdims}\OperatorTok{=}\VariableTok{True}\NormalTok{))}

\NormalTok{log\_probs }\OperatorTok{=}\NormalTok{ log\_softmax(logits)}
\NormalTok{nll }\OperatorTok{=} \OperatorTok{{-}}\NormalTok{log\_probs[np.arange(}\BuiltInTok{len}\NormalTok{(targets)), targets]}
\NormalTok{mean\_nll }\OperatorTok{=}\NormalTok{ np.mean(nll)}
\NormalTok{perplexity }\OperatorTok{=}\NormalTok{ np.exp(mean\_nll)}

\BuiltInTok{print}\NormalTok{(}\StringTok{"token NLL values:"}\NormalTok{, np.}\BuiltInTok{round}\NormalTok{(nll, }\DecValTok{4}\NormalTok{))}
\BuiltInTok{print}\NormalTok{(}\StringTok{"mean NLL:"}\NormalTok{, }\BuiltInTok{round}\NormalTok{(mean\_nll, }\DecValTok{4}\NormalTok{))}
\BuiltInTok{print}\NormalTok{(}\StringTok{"perplexity:"}\NormalTok{, }\BuiltInTok{round}\NormalTok{(perplexity, }\DecValTok{4}\NormalTok{))}
\end{Highlighting}
\end{Shaded}

\begin{verbatim}
token NLL values: [0.4076 0.3254 0.2529 1.2152]
mean NLL: 0.5503
perplexity: 1.7337
\end{verbatim}

\section{20.8 Calibration and probabilistic
honesty}\label{calibration-and-probabilistic-honesty}

Accuracy measures whether predictions are correct. Calibration measures
whether predicted probabilities mean what they say.

A binary classifier is calibrated if

\[
P(Y=1 \mid \hat p(X)=q)=q
\]

for all probability values \(q\). In practice, probabilities are binned.
For bin \(B_m\), define

\[
\operatorname{conf}(B_m)=\frac{1}{|B_m|}\sum_{i\in B_m}\hat p_i,
\]

and

\[
\operatorname{acc}(B_m)=\frac{1}{|B_m|}\sum_{i\in B_m}y_i.
\]

The expected calibration error is approximated by

\[
\operatorname{ECE}=\sum_{m=1}^M \frac{|B_m|}{n}\left|\operatorname{acc}(B_m)-\operatorname{conf}(B_m)\right|.
\]

A related proper scoring rule is the Brier score:

\[
\operatorname{Brier}=\frac{1}{n}\sum_{i=1}^n (\hat p_i-y_i)^2.
\]

A proper scoring rule rewards honest probability estimates.
Cross-entropy and Brier score are both proper scoring rules for
probabilistic prediction.

\begin{Shaded}
\begin{Highlighting}[]
\NormalTok{y }\OperatorTok{=}\NormalTok{ np.array([}\DecValTok{1}\NormalTok{, }\DecValTok{0}\NormalTok{, }\DecValTok{1}\NormalTok{, }\DecValTok{1}\NormalTok{, }\DecValTok{0}\NormalTok{, }\DecValTok{0}\NormalTok{, }\DecValTok{1}\NormalTok{, }\DecValTok{0}\NormalTok{, }\DecValTok{1}\NormalTok{, }\DecValTok{0}\NormalTok{, }\DecValTok{1}\NormalTok{, }\DecValTok{0}\NormalTok{])}
\NormalTok{p }\OperatorTok{=}\NormalTok{ np.array([}\FloatTok{0.95}\NormalTok{, }\FloatTok{0.88}\NormalTok{, }\FloatTok{0.82}\NormalTok{, }\FloatTok{0.76}\NormalTok{, }\FloatTok{0.69}\NormalTok{, }\FloatTok{0.55}\NormalTok{, }\FloatTok{0.52}\NormalTok{, }\FloatTok{0.45}\NormalTok{, }\FloatTok{0.37}\NormalTok{, }\FloatTok{0.28}\NormalTok{, }\FloatTok{0.21}\NormalTok{, }\FloatTok{0.08}\NormalTok{])}

\KeywordTok{def}\NormalTok{ calibration\_table(y, p, bins}\OperatorTok{=}\DecValTok{4}\NormalTok{):}
\NormalTok{    edges }\OperatorTok{=}\NormalTok{ np.linspace(}\DecValTok{0}\NormalTok{, }\DecValTok{1}\NormalTok{, bins }\OperatorTok{+} \DecValTok{1}\NormalTok{)}
\NormalTok{    rows }\OperatorTok{=}\NormalTok{ []}
\NormalTok{    ece }\OperatorTok{=} \FloatTok{0.0}
    \ControlFlowTok{for}\NormalTok{ j }\KeywordTok{in} \BuiltInTok{range}\NormalTok{(bins):}
\NormalTok{        lo, hi }\OperatorTok{=}\NormalTok{ edges[j], edges[j }\OperatorTok{+} \DecValTok{1}\NormalTok{]}
        \ControlFlowTok{if}\NormalTok{ j }\OperatorTok{==}\NormalTok{ bins }\OperatorTok{{-}} \DecValTok{1}\NormalTok{:}
\NormalTok{            mask }\OperatorTok{=}\NormalTok{ (p }\OperatorTok{\textgreater{}=}\NormalTok{ lo) }\OperatorTok{\&}\NormalTok{ (p }\OperatorTok{\textless{}=}\NormalTok{ hi)}
        \ControlFlowTok{else}\NormalTok{:}
\NormalTok{            mask }\OperatorTok{=}\NormalTok{ (p }\OperatorTok{\textgreater{}=}\NormalTok{ lo) }\OperatorTok{\&}\NormalTok{ (p }\OperatorTok{\textless{}}\NormalTok{ hi)}
        \ControlFlowTok{if}\NormalTok{ np.}\BuiltInTok{sum}\NormalTok{(mask) }\OperatorTok{==} \DecValTok{0}\NormalTok{:}
            \ControlFlowTok{continue}
\NormalTok{        conf }\OperatorTok{=}\NormalTok{ np.mean(p[mask])}
\NormalTok{        acc }\OperatorTok{=}\NormalTok{ np.mean(y[mask])}
\NormalTok{        weight }\OperatorTok{=}\NormalTok{ np.mean(mask)}
\NormalTok{        ece }\OperatorTok{+=}\NormalTok{ weight }\OperatorTok{*} \BuiltInTok{abs}\NormalTok{(acc }\OperatorTok{{-}}\NormalTok{ conf)}
\NormalTok{        rows.append((lo, hi, }\BuiltInTok{int}\NormalTok{(np.}\BuiltInTok{sum}\NormalTok{(mask)), conf, acc))}
    \ControlFlowTok{return}\NormalTok{ rows, ece}

\NormalTok{rows, ece }\OperatorTok{=}\NormalTok{ calibration\_table(y, p, bins}\OperatorTok{=}\DecValTok{4}\NormalTok{)}
\ControlFlowTok{for}\NormalTok{ row }\KeywordTok{in}\NormalTok{ rows:}
    \BuiltInTok{print}\NormalTok{(row)}
\BuiltInTok{print}\NormalTok{(}\StringTok{"ECE:"}\NormalTok{, }\BuiltInTok{round}\NormalTok{(ece, }\DecValTok{4}\NormalTok{))}
\BuiltInTok{print}\NormalTok{(}\StringTok{"Brier:"}\NormalTok{, }\BuiltInTok{round}\NormalTok{(np.mean((p }\OperatorTok{{-}}\NormalTok{ y) }\OperatorTok{**} \DecValTok{2}\NormalTok{), }\DecValTok{4}\NormalTok{))}
\end{Highlighting}
\end{Shaded}

\begin{verbatim}
(0.0, 0.25, 2, 0.145, 0.5)
(0.25, 0.5, 3, 0.3666666666666667, 0.3333333333333333)
(0.5, 0.75, 3, 0.5866666666666667, 0.3333333333333333)
(0.75, 1.0, 4, 0.8525, 0.75)
ECE: 0.165
Brier: 0.2654
\end{verbatim}

\section{20.9 Interactive: threshold, precision, recall, and
calibration}\label{interactive-threshold-precision-recall-and-calibration}

The following interactive example shows how a single set of probability
scores can lead to different classification metrics as the decision
threshold changes. It also shows how temperature scaling changes
probabilistic confidence.

\section{20.10 Overlap metrics for generated
text}\label{overlap-metrics-for-generated-text}

Generated text is difficult to evaluate because there may be many
correct answers. For example, two summaries can be semantically similar
while sharing few exact words. Nevertheless, overlap metrics remain
useful as quick baselines.

Let \(G\) be a generated sequence and \(R\) be a reference sequence. For
a token \(n\)-gram \(g\), let \(c_G(g)\) and \(c_R(g)\) be its counts. A
clipped precision score is

\[
P_n = \frac{\sum_g \min\{c_G(g),c_R(g)\}}{\sum_g c_G(g)}.
\]

BLEU-style evaluation combines clipped \(n\)-gram precisions with a
brevity penalty. A simplified version is

\[
\operatorname{BLEU}_N = BP \cdot \exp\left(\frac{1}{N}\sum_{n=1}^N \log P_n\right),
\]

where

\[
BP = \begin{cases}
1, & |G| > |R|,\\
\exp(1-|R|/|G|), & |G| \le |R|.
\end{cases}
\]

ROUGE-style evaluation often emphasizes recall. For unigrams, a simple
ROUGE-1 recall score is

\[
\operatorname{ROUGE\text{-}1\ recall} = \frac{\sum_g \min\{c_G(g),c_R(g)\}}{\sum_g c_R(g)}.
\]

The conceptual difference is important:

\begin{itemize}
\tightlist
\item
  BLEU-style precision asks: how much of the generated text is supported
  by the reference?
\item
  ROUGE-style recall asks: how much of the reference was covered by the
  generated text?
\end{itemize}

\begin{Shaded}
\begin{Highlighting}[]
\ImportTok{from}\NormalTok{ collections }\ImportTok{import}\NormalTok{ Counter}
\ImportTok{import}\NormalTok{ math}

\KeywordTok{def}\NormalTok{ tokenize\_simple(text):}
    \ControlFlowTok{return}\NormalTok{ text.lower().replace(}\StringTok{"."}\NormalTok{, }\StringTok{""}\NormalTok{).replace(}\StringTok{","}\NormalTok{, }\StringTok{""}\NormalTok{).split()}

\KeywordTok{def}\NormalTok{ ngrams(tokens, n):}
    \ControlFlowTok{return}\NormalTok{ [}\BuiltInTok{tuple}\NormalTok{(tokens[i:i}\OperatorTok{+}\NormalTok{n]) }\ControlFlowTok{for}\NormalTok{ i }\KeywordTok{in} \BuiltInTok{range}\NormalTok{(}\BuiltInTok{len}\NormalTok{(tokens) }\OperatorTok{{-}}\NormalTok{ n }\OperatorTok{+} \DecValTok{1}\NormalTok{)]}

\KeywordTok{def}\NormalTok{ clipped\_precision(generated, reference, n):}
\NormalTok{    g\_counts }\OperatorTok{=}\NormalTok{ Counter(ngrams(generated, n))}
\NormalTok{    r\_counts }\OperatorTok{=}\NormalTok{ Counter(ngrams(reference, n))}
\NormalTok{    overlap }\OperatorTok{=} \BuiltInTok{sum}\NormalTok{(}\BuiltInTok{min}\NormalTok{(count, r\_counts[gram]) }\ControlFlowTok{for}\NormalTok{ gram, count }\KeywordTok{in}\NormalTok{ g\_counts.items())}
\NormalTok{    total }\OperatorTok{=} \BuiltInTok{sum}\NormalTok{(g\_counts.values())}
    \ControlFlowTok{return}\NormalTok{ overlap }\OperatorTok{/}\NormalTok{ total }\ControlFlowTok{if}\NormalTok{ total }\OperatorTok{\textgreater{}} \DecValTok{0} \ControlFlowTok{else} \FloatTok{0.0}

\KeywordTok{def}\NormalTok{ rouge1\_recall(generated, reference):}
\NormalTok{    g\_counts }\OperatorTok{=}\NormalTok{ Counter(generated)}
\NormalTok{    r\_counts }\OperatorTok{=}\NormalTok{ Counter(reference)}
\NormalTok{    overlap }\OperatorTok{=} \BuiltInTok{sum}\NormalTok{(}\BuiltInTok{min}\NormalTok{(g\_counts[token], count) }\ControlFlowTok{for}\NormalTok{ token, count }\KeywordTok{in}\NormalTok{ r\_counts.items())}
\NormalTok{    total }\OperatorTok{=} \BuiltInTok{sum}\NormalTok{(r\_counts.values())}
    \ControlFlowTok{return}\NormalTok{ overlap }\OperatorTok{/}\NormalTok{ total }\ControlFlowTok{if}\NormalTok{ total }\OperatorTok{\textgreater{}} \DecValTok{0} \ControlFlowTok{else} \FloatTok{0.0}

\KeywordTok{def}\NormalTok{ bleu\_like(generated, reference, max\_n}\OperatorTok{=}\DecValTok{2}\NormalTok{, eps}\OperatorTok{=}\FloatTok{1e{-}12}\NormalTok{):}
\NormalTok{    precisions }\OperatorTok{=}\NormalTok{ [}\BuiltInTok{max}\NormalTok{(clipped\_precision(generated, reference, n), eps) }\ControlFlowTok{for}\NormalTok{ n }\KeywordTok{in} \BuiltInTok{range}\NormalTok{(}\DecValTok{1}\NormalTok{, max\_n }\OperatorTok{+} \DecValTok{1}\NormalTok{)]}
\NormalTok{    geo }\OperatorTok{=}\NormalTok{ math.exp(}\BuiltInTok{sum}\NormalTok{(math.log(p) }\ControlFlowTok{for}\NormalTok{ p }\KeywordTok{in}\NormalTok{ precisions) }\OperatorTok{/}\NormalTok{ max\_n)}
\NormalTok{    bp }\OperatorTok{=} \FloatTok{1.0} \ControlFlowTok{if} \BuiltInTok{len}\NormalTok{(generated) }\OperatorTok{\textgreater{}} \BuiltInTok{len}\NormalTok{(reference) }\ControlFlowTok{else}\NormalTok{ math.exp(}\DecValTok{1} \OperatorTok{{-}} \BuiltInTok{len}\NormalTok{(reference) }\OperatorTok{/} \BuiltInTok{max}\NormalTok{(}\BuiltInTok{len}\NormalTok{(generated), }\DecValTok{1}\NormalTok{))}
    \ControlFlowTok{return}\NormalTok{ bp }\OperatorTok{*}\NormalTok{ geo}

\NormalTok{reference }\OperatorTok{=}\NormalTok{ tokenize\_simple(}\StringTok{"the model retrieves relevant evidence before answering"}\NormalTok{)}
\NormalTok{generated }\OperatorTok{=}\NormalTok{ tokenize\_simple(}\StringTok{"the system retrieves evidence before generating an answer"}\NormalTok{)}

\BuiltInTok{print}\NormalTok{(}\StringTok{"clipped unigram precision:"}\NormalTok{, clipped\_precision(generated, reference, }\DecValTok{1}\NormalTok{))}
\BuiltInTok{print}\NormalTok{(}\StringTok{"clipped bigram precision:"}\NormalTok{, clipped\_precision(generated, reference, }\DecValTok{2}\NormalTok{))}
\BuiltInTok{print}\NormalTok{(}\StringTok{"ROUGE{-}1 recall:"}\NormalTok{, rouge1\_recall(generated, reference))}
\BuiltInTok{print}\NormalTok{(}\StringTok{"BLEU{-}like score:"}\NormalTok{, bleu\_like(generated, reference, max\_n}\OperatorTok{=}\DecValTok{2}\NormalTok{))}
\end{Highlighting}
\end{Shaded}

\begin{verbatim}
clipped unigram precision: 0.5
clipped bigram precision: 0.14285714285714285
ROUGE-1 recall: 0.5714285714285714
BLEU-like score: 0.26726124191242434
\end{verbatim}

\section{20.11 Semantic evaluation and embedding
similarity}\label{semantic-evaluation-and-embedding-similarity}

Exact token overlap fails when two answers are semantically similar but
lexically different. A common alternative is to embed generated and
reference text into vectors:

\[
z_G = f(G), \qquad z_R = f(R),
\]

and score similarity by cosine similarity:

\[
\operatorname{sim}(G,R)=\frac{z_G^T z_R}{\|z_G\|\|z_R\|}.
\]

This is mathematically clean but not automatically reliable. The
embedding model itself has limitations. It may miss precise mathematical
distinctions, factual errors, negations, or numerical mistakes.

For example, the sentences

\begin{Shaded}
\begin{Highlighting}[]
\NormalTok{The null hypothesis is rejected.}
\NormalTok{The null hypothesis is not rejected.}
\end{Highlighting}
\end{Shaded}

may be close in embedding space because most words are shared, but their
statistical meanings are different.

Thus semantic similarity is useful, but it should not replace
task-specific correctness checks.

\section{20.12 Retrieval and RAG
evaluation}\label{retrieval-and-rag-evaluation}

For retrieval, the model returns an ordered list of documents

\[
(d_1,d_2,\ldots,d_k).
\]

Let \(R(q)\) be the set of relevant documents for query \(q\).
Precision@\(k\) is

\[
\operatorname{Precision@}k = \frac{|\{d_1,\ldots,d_k\}\cap R(q)|}{k}.
\]

Recall@\(k\) is

\[
\operatorname{Recall@}k = \frac{|\{d_1,\ldots,d_k\}\cap R(q)|}{|R(q)|}.
\]

The reciprocal rank is

\[
RR(q)=\frac{1}{\operatorname{rank of first relevant document}}.
\]

The mean reciprocal rank is the average over queries:

\[
\operatorname{MRR}=\frac{1}{m}\sum_{j=1}^m RR(q_j).
\]

If relevance has graded scores, discounted cumulative gain is

\[
\operatorname{DCG@}k = \sum_{i=1}^k \frac{2^{rel_i}-1}{\log_2(i+1)}.
\]

Normalized DCG divides by the ideal DCG:

\[
\operatorname{nDCG@}k = \frac{\operatorname{DCG@}k}{\operatorname{IDCG@}k}.
\]

\subsection{RAG-specific
decomposition}\label{rag-specific-decomposition}

A RAG system should be evaluated in at least three pieces:

\[
\text{RAG quality} = \text{retrieval quality} + \text{grounding quality} + \text{answer quality}.
\]

Retrieval quality asks whether useful evidence was retrieved. Grounding
quality asks whether the answer is supported by the retrieved evidence.
Answer quality asks whether the final response is correct, complete, and
usable.

A failure may occur at any stage:

\begin{enumerate}
\def\labelenumi{\arabic{enumi}.}
\tightlist
\item
  the retriever misses the relevant document;
\item
  the retriever finds it but ranks it too low;
\item
  the generator ignores the evidence;
\item
  the generator cites evidence incorrectly;
\item
  the answer is correct but incomplete;
\item
  the answer is fluent but unsupported.
\end{enumerate}

\begin{Shaded}
\begin{Highlighting}[]
\KeywordTok{def}\NormalTok{ precision\_at\_k(ranked, relevant, k):}
\NormalTok{    top }\OperatorTok{=}\NormalTok{ ranked[:k]}
    \ControlFlowTok{return} \BuiltInTok{len}\NormalTok{(}\BuiltInTok{set}\NormalTok{(top) }\OperatorTok{\&} \BuiltInTok{set}\NormalTok{(relevant)) }\OperatorTok{/}\NormalTok{ k}

\KeywordTok{def}\NormalTok{ recall\_at\_k(ranked, relevant, k):}
\NormalTok{    top }\OperatorTok{=}\NormalTok{ ranked[:k]}
    \ControlFlowTok{return} \BuiltInTok{len}\NormalTok{(}\BuiltInTok{set}\NormalTok{(top) }\OperatorTok{\&} \BuiltInTok{set}\NormalTok{(relevant)) }\OperatorTok{/} \BuiltInTok{len}\NormalTok{(relevant) }\ControlFlowTok{if}\NormalTok{ relevant }\ControlFlowTok{else} \FloatTok{0.0}

\KeywordTok{def}\NormalTok{ reciprocal\_rank(ranked, relevant):}
\NormalTok{    relevant }\OperatorTok{=} \BuiltInTok{set}\NormalTok{(relevant)}
    \ControlFlowTok{for}\NormalTok{ i, doc }\KeywordTok{in} \BuiltInTok{enumerate}\NormalTok{(ranked, start}\OperatorTok{=}\DecValTok{1}\NormalTok{):}
        \ControlFlowTok{if}\NormalTok{ doc }\KeywordTok{in}\NormalTok{ relevant:}
            \ControlFlowTok{return} \DecValTok{1} \OperatorTok{/}\NormalTok{ i}
    \ControlFlowTok{return} \FloatTok{0.0}

\KeywordTok{def}\NormalTok{ dcg\_at\_k(relevances, k):}
    \ControlFlowTok{return} \BuiltInTok{sum}\NormalTok{((}\DecValTok{2} \OperatorTok{**}\NormalTok{ rel }\OperatorTok{{-}} \DecValTok{1}\NormalTok{) }\OperatorTok{/}\NormalTok{ math.log2(i }\OperatorTok{+} \DecValTok{2}\NormalTok{) }\ControlFlowTok{for}\NormalTok{ i, rel }\KeywordTok{in} \BuiltInTok{enumerate}\NormalTok{(relevances[:k]))}

\KeywordTok{def}\NormalTok{ ndcg\_at\_k(ranked\_relevances, k):}
\NormalTok{    dcg }\OperatorTok{=}\NormalTok{ dcg\_at\_k(ranked\_relevances, k)}
\NormalTok{    ideal }\OperatorTok{=}\NormalTok{ dcg\_at\_k(}\BuiltInTok{sorted}\NormalTok{(ranked\_relevances, reverse}\OperatorTok{=}\VariableTok{True}\NormalTok{), k)}
    \ControlFlowTok{return}\NormalTok{ dcg }\OperatorTok{/}\NormalTok{ ideal }\ControlFlowTok{if}\NormalTok{ ideal }\OperatorTok{\textgreater{}} \DecValTok{0} \ControlFlowTok{else} \FloatTok{0.0}

\NormalTok{ranked\_docs }\OperatorTok{=}\NormalTok{ [}\StringTok{"d7"}\NormalTok{, }\StringTok{"d2"}\NormalTok{, }\StringTok{"d9"}\NormalTok{, }\StringTok{"d1"}\NormalTok{, }\StringTok{"d4"}\NormalTok{]}
\NormalTok{relevant\_docs }\OperatorTok{=}\NormalTok{ \{}\StringTok{"d1"}\NormalTok{, }\StringTok{"d2"}\NormalTok{, }\StringTok{"d4"}\NormalTok{\}}
\NormalTok{ranked\_relevances }\OperatorTok{=}\NormalTok{ [}\DecValTok{0}\NormalTok{, }\DecValTok{2}\NormalTok{, }\DecValTok{0}\NormalTok{, }\DecValTok{3}\NormalTok{, }\DecValTok{1}\NormalTok{]}

\BuiltInTok{print}\NormalTok{(}\StringTok{"Precision@3:"}\NormalTok{, precision\_at\_k(ranked\_docs, relevant\_docs, }\DecValTok{3}\NormalTok{))}
\BuiltInTok{print}\NormalTok{(}\StringTok{"Recall@3:"}\NormalTok{, recall\_at\_k(ranked\_docs, relevant\_docs, }\DecValTok{3}\NormalTok{))}
\BuiltInTok{print}\NormalTok{(}\StringTok{"MRR contribution:"}\NormalTok{, reciprocal\_rank(ranked\_docs, relevant\_docs))}
\BuiltInTok{print}\NormalTok{(}\StringTok{"nDCG@5:"}\NormalTok{, ndcg\_at\_k(ranked\_relevances, }\DecValTok{5}\NormalTok{))}
\end{Highlighting}
\end{Shaded}

\begin{verbatim}
Precision@3: 0.3333333333333333
Recall@3: 0.3333333333333333
MRR contribution: 0.5
nDCG@5: 0.5636640860885236
\end{verbatim}

\section{20.13 Human evaluation}\label{human-evaluation}

For open-ended generation, human evaluation is often necessary. A
typical human evaluation protocol defines a rubric. For example, a
generated answer may be rated on:

\begin{itemize}
\tightlist
\item
  correctness;
\item
  completeness;
\item
  clarity;
\item
  faithfulness to evidence;
\item
  mathematical precision;
\item
  usefulness for the intended audience.
\end{itemize}

Suppose \(r_{ij}\) is the rating assigned by evaluator \(j\) to item
\(i\). A simple average score is

\[
\bar r = \frac{1}{mn}\sum_{i=1}^n\sum_{j=1}^m r_{ij}.
\]

However, human evaluation has its own statistical issues:

\begin{itemize}
\tightlist
\item
  annotators may disagree;
\item
  rubrics may be ambiguous;
\item
  ratings may be ordinal rather than interval-valued;
\item
  examples may not represent the target population;
\item
  evaluators may be influenced by fluency or length.
\end{itemize}

Inter-rater agreement should be checked when human scores are used as
evidence. A simple pairwise agreement rate is

\[
\operatorname{Agreement}=\frac{\text{number of matching rating decisions}}{\text{number of compared decisions}}.
\]

For categorical ratings, Cohen's kappa adjusts for chance agreement:

\[
\kappa = \frac{p_o-p_e}{1-p_e},
\]

where \(p_o\) is observed agreement and \(p_e\) is expected chance
agreement under the annotators' marginal distributions.

\section{20.14 Statistical uncertainty in
evaluation}\label{statistical-uncertainty-in-evaluation}

An evaluation score computed on a finite test set is a random variable.
For a loss sample

\[
\ell_i = L(f_\theta(x_i),y_i),
\]

the empirical mean is

\[
\bar \ell = \frac{1}{n}\sum_{i=1}^n \ell_i.
\]

The sample variance is

\[
s^2 = \frac{1}{n-1}\sum_{i=1}^n (\ell_i-\bar \ell)^2.
\]

An approximate standard error is

\[
SE(\bar \ell)=\frac{s}{\sqrt n}.
\]

A rough \(95\%\) confidence interval is

\[
\bar \ell \pm 1.96\frac{s}{\sqrt n}.
\]

For metrics such as \(F_1\), BLEU, or nDCG, the formula may be
complicated. A bootstrap interval is often easier. The nonparametric
bootstrap repeatedly samples test examples with replacement and
recomputes the metric.

\begin{Shaded}
\begin{Highlighting}[]
\NormalTok{rng }\OperatorTok{=}\NormalTok{ np.random.default\_rng(}\DecValTok{7243}\NormalTok{)}

\CommentTok{\# Simulated item{-}level correctness for an NLP system.}
\NormalTok{correct }\OperatorTok{=}\NormalTok{ rng.binomial(}\DecValTok{1}\NormalTok{, }\FloatTok{0.78}\NormalTok{, size}\OperatorTok{=}\DecValTok{200}\NormalTok{)}
\NormalTok{accuracy\_hat }\OperatorTok{=}\NormalTok{ np.mean(correct)}
\NormalTok{se }\OperatorTok{=}\NormalTok{ np.std(correct, ddof}\OperatorTok{=}\DecValTok{1}\NormalTok{) }\OperatorTok{/}\NormalTok{ np.sqrt(}\BuiltInTok{len}\NormalTok{(correct))}
\NormalTok{normal\_ci }\OperatorTok{=}\NormalTok{ (accuracy\_hat }\OperatorTok{{-}} \FloatTok{1.96} \OperatorTok{*}\NormalTok{ se, accuracy\_hat }\OperatorTok{+} \FloatTok{1.96} \OperatorTok{*}\NormalTok{ se)}

\CommentTok{\# Bootstrap confidence interval.}
\NormalTok{B }\OperatorTok{=} \DecValTok{2000}
\NormalTok{boot }\OperatorTok{=}\NormalTok{ []}
\ControlFlowTok{for}\NormalTok{ \_ }\KeywordTok{in} \BuiltInTok{range}\NormalTok{(B):}
\NormalTok{    sample }\OperatorTok{=}\NormalTok{ rng.choice(correct, size}\OperatorTok{=}\BuiltInTok{len}\NormalTok{(correct), replace}\OperatorTok{=}\VariableTok{True}\NormalTok{)}
\NormalTok{    boot.append(np.mean(sample))}
\NormalTok{boot\_ci }\OperatorTok{=}\NormalTok{ np.quantile(boot, [}\FloatTok{0.025}\NormalTok{, }\FloatTok{0.975}\NormalTok{])}

\BuiltInTok{print}\NormalTok{(}\StringTok{"accuracy:"}\NormalTok{, }\BuiltInTok{round}\NormalTok{(accuracy\_hat, }\DecValTok{4}\NormalTok{))}
\BuiltInTok{print}\NormalTok{(}\StringTok{"normal CI:"}\NormalTok{, }\BuiltInTok{tuple}\NormalTok{(np.}\BuiltInTok{round}\NormalTok{(normal\_ci, }\DecValTok{4}\NormalTok{)))}
\BuiltInTok{print}\NormalTok{(}\StringTok{"bootstrap CI:"}\NormalTok{, }\BuiltInTok{tuple}\NormalTok{(np.}\BuiltInTok{round}\NormalTok{(boot\_ci, }\DecValTok{4}\NormalTok{)))}
\end{Highlighting}
\end{Shaded}

\begin{verbatim}
accuracy: 0.83
normal CI: (0.7778, 0.8822)
bootstrap CI: (0.775, 0.88)
\end{verbatim}

\section{20.15 Comparing two models}\label{comparing-two-models}

Suppose two systems A and B are evaluated on the same \(n\) examples.
Let

\[
d_i = L_A(x_i,y_i)-L_B(x_i,y_i).
\]

If lower loss is better, then \(d_i>0\) means B is better on item \(i\).
The mean difference is

\[
\bar d = \frac{1}{n}\sum_{i=1}^n d_i.
\]

A paired comparison is usually better than comparing two independent
averages because both models are evaluated on the same items. The paired
standard error is

\[
SE(\bar d)=\frac{s_d}{\sqrt n}.
\]

where \(s_d\) is the sample standard deviation of the item-level
differences.

\begin{Shaded}
\begin{Highlighting}[]
\NormalTok{rng }\OperatorTok{=}\NormalTok{ np.random.default\_rng(}\DecValTok{7339}\NormalTok{)}
\NormalTok{loss\_A }\OperatorTok{=}\NormalTok{ rng.normal(loc}\OperatorTok{=}\FloatTok{0.42}\NormalTok{, scale}\OperatorTok{=}\FloatTok{0.12}\NormalTok{, size}\OperatorTok{=}\DecValTok{100}\NormalTok{)}
\NormalTok{loss\_B }\OperatorTok{=}\NormalTok{ loss\_A }\OperatorTok{{-}}\NormalTok{ rng.normal(loc}\OperatorTok{=}\FloatTok{0.03}\NormalTok{, scale}\OperatorTok{=}\FloatTok{0.08}\NormalTok{, size}\OperatorTok{=}\DecValTok{100}\NormalTok{)}
\NormalTok{d }\OperatorTok{=}\NormalTok{ loss\_A }\OperatorTok{{-}}\NormalTok{ loss\_B}
\NormalTok{mean\_diff }\OperatorTok{=}\NormalTok{ np.mean(d)}
\NormalTok{se\_diff }\OperatorTok{=}\NormalTok{ np.std(d, ddof}\OperatorTok{=}\DecValTok{1}\NormalTok{) }\OperatorTok{/}\NormalTok{ np.sqrt(}\BuiltInTok{len}\NormalTok{(d))}
\NormalTok{ci\_diff }\OperatorTok{=}\NormalTok{ (mean\_diff }\OperatorTok{{-}} \FloatTok{1.96} \OperatorTok{*}\NormalTok{ se\_diff, mean\_diff }\OperatorTok{+} \FloatTok{1.96} \OperatorTok{*}\NormalTok{ se\_diff)}

\BuiltInTok{print}\NormalTok{(}\StringTok{"mean paired loss reduction A {-} B:"}\NormalTok{, }\BuiltInTok{round}\NormalTok{(mean\_diff, }\DecValTok{4}\NormalTok{))}
\BuiltInTok{print}\NormalTok{(}\StringTok{"approximate 95\% CI:"}\NormalTok{, }\BuiltInTok{tuple}\NormalTok{(np.}\BuiltInTok{round}\NormalTok{(ci\_diff, }\DecValTok{4}\NormalTok{)))}
\end{Highlighting}
\end{Shaded}

\begin{verbatim}
mean paired loss reduction A - B: 0.0297
approximate 95% CI: (0.012, 0.0474)
\end{verbatim}

A confidence interval for the paired difference is more informative than
saying one score is simply larger than another. It tells us whether the
observed difference is large compared with test-set variability.

\section{20.16 Robustness and distribution
shift}\label{robustness-and-distribution-shift}

Let \(P_{\text{test}}\) be the distribution represented by the
evaluation set and \(P_{\text{deploy}}\) be the distribution where the
system is actually used. A model with low empirical risk on
\(P_{\text{test}}\) may fail when

\[
P_{\text{deploy}} \ne P_{\text{test}}.
\]

In NLP, distribution shift may come from:

\begin{itemize}
\tightlist
\item
  a different domain;
\item
  a different institution;
\item
  new terminology;
\item
  multilingual or code-switched text;
\item
  different prompt styles;
\item
  adversarial inputs;
\item
  longer contexts;
\item
  updated policies or facts.
\end{itemize}

Robust evaluation therefore should include subpopulation analysis. If
the test set is partitioned into groups \(G_1,\ldots,G_m\), compute

\[
\widehat R_j = \frac{1}{|G_j|}\sum_{i\in G_j} L(f_\theta(x_i),y_i).
\]

A single average can hide serious failures on smaller groups.

\section{20.17 Evaluation for mathematical NLP
tasks}\label{evaluation-for-mathematical-nlp-tasks}

For this book's audience, mathematical precision matters. A language
model can be fluent but mathematically wrong. For math-heavy NLP tasks,
evaluation should check:

\begin{enumerate}
\def\labelenumi{\arabic{enumi}.}
\tightlist
\item
  \textbf{symbolic correctness:} Are formulas syntactically and
  mathematically valid?
\item
  \textbf{logical correctness:} Does each conclusion follow from
  previous statements?
\item
  \textbf{numerical correctness:} Are computations correct?
\item
  \textbf{definition consistency:} Are terms used consistently?
\item
  \textbf{assumption tracking:} Are required assumptions stated?
\item
  \textbf{citation and source grounding:} For RAG, is every factual
  claim supported by evidence?
\end{enumerate}

A useful rubric for generated mathematical explanations may score each
answer using a vector

\[
r = (r_{\text{correct}}, r_{\text{complete}}, r_{\text{clear}}, r_{\text{rigorous}}, r_{\text{audience}}).
\]

The final evaluation may use a weighted score

\[
S = w^T r,
\]

where weights \(w\) should be chosen before evaluation.

\section{20.18 AI-assisted evaluation
workflows}\label{ai-assisted-evaluation-workflows}

Large language models can help with evaluation, but they should not be
treated as infallible judges. A good AI-assisted workflow is:

\begin{enumerate}
\def\labelenumi{\arabic{enumi}.}
\tightlist
\item
  ask the AI judge to identify specific claims;
\item
  ask it to compare each claim with the reference or retrieved evidence;
\item
  require a structured score with reasons;
\item
  audit a sample manually;
\item
  measure agreement between AI judgments and human judgments.
\end{enumerate}

For mathematical material, an AI judge should be asked to separate:

\begin{itemize}
\tightlist
\item
  algebraic correctness;
\item
  statistical interpretation;
\item
  notation consistency;
\item
  hidden assumptions;
\item
  relevance to the question.
\end{itemize}

\begin{tcolorbox}[enhanced jigsaw, arc=.35mm, bottomrule=.15mm, bottomtitle=1mm, breakable, colback=white, colbacktitle=quarto-callout-important-color!10!white, colframe=quarto-callout-important-color-frame, coltitle=black, left=2mm, leftrule=.75mm, opacityback=0, opacitybacktitle=0.6, rightrule=.15mm, title=\textcolor{quarto-callout-important-color}{\faExclamation}\hspace{0.5em}{AI judge caution}, titlerule=0mm, toprule=.15mm, toptitle=1mm]

Using a language model to evaluate another language model creates a
second modeling problem. The evaluator has its own biases, blind spots,
and uncertainty. Treat AI-based evaluation as an instrument that must
itself be validated.

\end{tcolorbox}

\subsection{Prompt pattern for AI-assisted
evaluation}\label{prompt-pattern-for-ai-assisted-evaluation}

\begin{Shaded}
\begin{Highlighting}[]
\NormalTok{You are evaluating an NLP model answer for a graduate applied mathematics audience.}
\NormalTok{Evaluate the answer using the following dimensions:}
\NormalTok{1. mathematical correctness,}
\NormalTok{2. statistical interpretation,}
\NormalTok{3. notation consistency,}
\NormalTok{4. completeness,}
\NormalTok{5. clarity.}
\NormalTok{For each dimension, give a score from 0 to 4 and cite the exact sentence or formula that supports your judgment.}
\NormalTok{Do not reward fluent but unsupported claims.}
\end{Highlighting}
\end{Shaded}

This prompt is not a substitute for human evaluation. It is a way to
make the evaluation process more structured.

\section{20.19 A compact evaluation
checklist}\label{a-compact-evaluation-checklist}

Before reporting an NLP evaluation score, answer the following
questions.

\begin{enumerate}
\def\labelenumi{\arabic{enumi}.}
\tightlist
\item
  What is the target population distribution?
\item
  What is the exact test set?
\item
  Was the test set used during development?
\item
  What is the loss or metric?
\item
  Does the metric match the desired behavior?
\item
  What is the uncertainty of the estimate?
\item
  Are subgroups or domains evaluated separately?
\item
  Are errors qualitatively inspected?
\item
  Are failure cases reported?
\item
  Could the model optimize the metric while failing the real task?
\end{enumerate}

\section{20.20 Worked example: evaluating a mini sentiment
classifier}\label{worked-example-evaluating-a-mini-sentiment-classifier}

Suppose a sentiment classifier outputs the following probabilities for
the positive class:

\[
(0.9,0.8,0.7,0.6,0.4,0.3,0.2,0.1)
\]

and the true labels are

\[
(1,1,0,1,0,0,1,0).
\]

At threshold \(\tau=0.5\), the predictions are

\[
(1,1,1,1,0,0,0,0).
\]

Thus

\[
TP=3, \quad FP=1, \quad FN=1, \quad TN=3.
\]

Accuracy is

\[
\frac{TP+TN}{8}=\frac{6}{8}=0.75.
\]

Precision is

\[
\frac{TP}{TP+FP}=\frac{3}{4}=0.75.
\]

Recall is

\[
\frac{TP}{TP+FN}=\frac{3}{4}=0.75.
\]

Therefore

\[
F_1 = 0.75.
\]

The example is small, but it illustrates the general idea: metrics are
functions of counts induced by a decision rule.

\section{20.21 Common mistakes}\label{common-mistakes-9}

\begin{enumerate}
\def\labelenumi{\arabic{enumi}.}
\tightlist
\item
  \textbf{Reporting one number without uncertainty.} A score on a small
  test set may not be stable.
\item
  \textbf{Comparing perplexity across different tokenizers.}
  Tokenization changes the unit of measurement.
\item
  \textbf{Using accuracy on imbalanced data.} High accuracy may hide
  poor minority-class performance.
\item
  \textbf{Treating overlap metrics as semantic correctness.} Word
  overlap does not guarantee meaning.
\item
  \textbf{Ignoring retrieval in RAG evaluation.} A good answer may be
  impossible if retrieval failed.
\item
  \textbf{Optimizing the test set.} Repeated test-set use biases
  reported performance.
\item
  \textbf{Not examining errors.} Aggregate metrics rarely reveal all
  failure modes.
\item
  \textbf{Using an AI judge without validation.} The evaluator must also
  be evaluated.
\end{enumerate}

\section{20.22 Python mini-project: evaluation report from
predictions}\label{python-mini-project-evaluation-report-from-predictions}

The following code produces a compact evaluation report for a binary
classifier. Students can extend it to multiclass classification,
retrieval, or generated-text evaluation.

\begin{Shaded}
\begin{Highlighting}[]
\KeywordTok{def}\NormalTok{ evaluation\_report(y\_true, p\_score, threshold}\OperatorTok{=}\FloatTok{0.5}\NormalTok{, bins}\OperatorTok{=}\DecValTok{5}\NormalTok{):}
\NormalTok{    base }\OperatorTok{=}\NormalTok{ binary\_metrics(y\_true, p\_score, threshold)}
\NormalTok{    rows, ece }\OperatorTok{=}\NormalTok{ calibration\_table(y\_true, p\_score, bins}\OperatorTok{=}\NormalTok{bins)}
\NormalTok{    brier }\OperatorTok{=} \BuiltInTok{float}\NormalTok{(np.mean((p\_score }\OperatorTok{{-}}\NormalTok{ y\_true) }\OperatorTok{**} \DecValTok{2}\NormalTok{))}
    \ControlFlowTok{return}\NormalTok{ \{}
        \OperatorTok{**}\NormalTok{base,}
        \StringTok{"ECE"}\NormalTok{: ece,}
        \StringTok{"Brier"}\NormalTok{: brier,}
        \StringTok{"calibration\_bins"}\NormalTok{: rows,}
\NormalTok{    \}}

\NormalTok{report }\OperatorTok{=}\NormalTok{ evaluation\_report(y\_true, p\_score, threshold}\OperatorTok{=}\FloatTok{0.5}\NormalTok{, bins}\OperatorTok{=}\DecValTok{5}\NormalTok{)}
\ControlFlowTok{for}\NormalTok{ key, value }\KeywordTok{in}\NormalTok{ report.items():}
    \BuiltInTok{print}\NormalTok{(key, }\StringTok{":"}\NormalTok{, value)}
\end{Highlighting}
\end{Shaded}

\begin{verbatim}
threshold : 0.5
TP : 3
FP : 1
FN : 2
TN : 4
accuracy : 0.7
precision : 0.75
recall : 0.6
f1 : 0.6666666666666665
ECE : 0.21000000000000005
Brier : 0.21642
calibration_bins : [(0.0, 0.2, 1, 0.1, 0.0), (0.2, 0.4, 3, 0.2933333333333333, 0.6666666666666666), (0.4, 0.6000000000000001, 3, 0.47666666666666674, 0.3333333333333333), (0.6000000000000001, 0.8, 2, 0.685, 0.5), (0.8, 1.0, 1, 0.92, 1.0)]
\end{verbatim}

\section{20.23 Exercises}\label{exercises-19}

\subsection{Conceptual exercises}\label{conceptual-exercises-8}

\begin{enumerate}
\def\labelenumi{\arabic{enumi}.}
\tightlist
\item
  Explain why evaluation is an estimation problem.
\item
  Give an example where accuracy is misleading but macro \(F_1\) is
  informative.
\item
  Explain why perplexity should not be compared directly across two
  models with different tokenizers.
\item
  Describe a situation where ROUGE-style recall rewards a poor summary.
\item
  Explain why a RAG system should be evaluated in separate retrieval and
  generation stages.
\item
  Describe one way an AI judge can be useful and one way it can be
  misleading.
\end{enumerate}

\subsection{Mathematical exercises}\label{mathematical-exercises-7}

\begin{enumerate}
\def\labelenumi{\arabic{enumi}.}
\tightlist
\item
  Prove that the empirical accuracy is an unbiased estimator of
  population accuracy when test examples are independent samples from
  the target distribution.
\item
  Derive the \(F_1\) formula from precision and recall as a harmonic
  mean.
\item
  Show that if \(\widehat H\) is average negative log-likelihood in
  nats, then perplexity is \(\exp(\widehat H)\).
\item
  Compute Precision@3, Recall@3, and MRR for the ranked list
  \((d_3,d_5,d_1,d_2)\) when the relevant set is \(\{d_1,d_2\}\).
\item
  For a binary classifier with scores \((0.9,0.6,0.4,0.2)\) and labels
  \((1,0,1,0)\), compute accuracy, precision, recall, and \(F_1\) at
  thresholds \(0.3\), \(0.5\), and \(0.7\).
\item
  Suppose item-level losses have sample mean \(0.42\), sample standard
  deviation \(0.16\), and \(n=100\). Compute an approximate \(95\%\)
  confidence interval for population loss.
\end{enumerate}

\subsection{Computational exercises}\label{computational-exercises-5}

\begin{enumerate}
\def\labelenumi{\arabic{enumi}.}
\tightlist
\item
  Implement macro, micro, and weighted \(F_1\) for a multiclass
  classifier from scratch.
\item
  Write a function that computes token-level cross-entropy and
  perplexity from logits and target token indices.
\item
  Implement a bootstrap confidence interval for \(F_1\).
\item
  Implement BLEU-like clipped precision for \(n=1,2,3,4\).
\item
  Implement Precision@\(k\), Recall@\(k\), MRR, and nDCG for a list of
  queries.
\item
  Build a small evaluation dashboard for a text classifier using Plotly.
\end{enumerate}

\section{20.24 Mini-project}\label{mini-project-4}

Choose one NLP system from the earlier chapters:

\begin{itemize}
\tightlist
\item
  Naive Bayes sentiment classifier;
\item
  TF-IDF retrieval system;
\item
  embedding-based retrieval system;
\item
  small language model;
\item
  RAG prototype.
\end{itemize}

Write a short evaluation report containing:

\begin{enumerate}
\def\labelenumi{\arabic{enumi}.}
\tightlist
\item
  task definition;
\item
  data split;
\item
  chosen metrics and why they are appropriate;
\item
  numerical results;
\item
  confidence intervals or bootstrap intervals;
\item
  at least five qualitative error examples;
\item
  one recommendation for improving the system.
\end{enumerate}

The report should distinguish between what the metric proves and what it
does not prove.

\section{20.25 Lab connection}\label{lab-connection-9}

The companion lab for this chapter should ask students to build an
evaluation toolkit from scratch. The lab should include:

\begin{enumerate}
\def\labelenumi{\arabic{enumi}.}
\tightlist
\item
  classification metrics;
\item
  calibration metrics;
\item
  perplexity calculation;
\item
  overlap metrics for generated text;
\item
  retrieval metrics;
\item
  bootstrap uncertainty;
\item
  a small Plotly dashboard.
\end{enumerate}

The main lesson is that evaluation is not a final afterthought. It is a
mathematical design choice that shapes the entire NLP workflow.

\section{20.26 Summary}\label{summary-4}

This chapter developed evaluation as a statistical and mathematical
problem. The central object is the population risk

\[
R(\theta)=E_P[L(f_\theta(X),Y)],
\]

and the practical object is an empirical estimate from a finite test
set. Classification metrics arise from confusion counts. Language-model
metrics arise from likelihood and cross-entropy. Calibration evaluates
probabilistic honesty. Overlap metrics approximate generated-text
similarity but do not fully measure meaning. Retrieval and RAG systems
require multi-stage evaluation. Human and AI-assisted evaluation can be
valuable, but both require clear rubrics and validation.

A mathematically mature NLP practitioner should always ask:

\begin{quote}
What quantity is this metric estimating, and is that quantity actually
what we care about?
\end{quote}

\part{Part VII. Advanced Mathematical Perspectives}

\chapter{Chapter 21. Information Theory and
Language}\label{chapter-21.-information-theory-and-language}

Entropy, Divergence, Mutual Information, and Compression

\hfill\break

\chapter{Chapter 21. Information Theory and
Language}\label{chapter-21.-information-theory-and-language-1}

Information theory gives a mathematical language for uncertainty,
prediction, compression, and dependence. For natural language
processing, these ideas are not optional background. They appear
directly in language-model training, perplexity, tokenization, decoding,
representation learning, retrieval, and alignment.

The central observation is simple:

\begin{quote}
A language model assigns probabilities to possible tokens and sequences.
Therefore, a language model can be studied using the mathematics of
probability distributions and information.
\end{quote}

Earlier chapters used cross-entropy as a loss function and perplexity as
an evaluation metric. This chapter explains why those quantities are
natural. The goal is to make explicit the connection between prediction
and compression:

\[
\text{good probabilistic prediction} \quad \Longleftrightarrow \quad \text{short code length on average}.
\]

For MA Applied Mathematics and MS Statistics students, information
theory is especially useful because it gives clean definitions that
apply to both classical statistical NLP and modern neural language
models.

\begin{tcolorbox}[enhanced jigsaw, arc=.35mm, bottomrule=.15mm, bottomtitle=1mm, breakable, colback=white, colbacktitle=quarto-callout-tip-color!10!white, colframe=quarto-callout-tip-color-frame, coltitle=black, left=2mm, leftrule=.75mm, opacityback=0, opacitybacktitle=0.6, rightrule=.15mm, title=\textcolor{quarto-callout-tip-color}{\faLightbulb}\hspace{0.5em}{Guiding principle}, titlerule=0mm, toprule=.15mm, toptitle=1mm]

Entropy measures uncertainty under a true distribution. Cross-entropy
measures the average loss when one distribution is used to predict
another. KL divergence measures the extra cost caused by using the wrong
distribution.

\end{tcolorbox}

\section{21.1 Learning goals}\label{learning-goals-20}

After studying this chapter, you should be able to:

\begin{enumerate}
\def\labelenumi{\arabic{enumi}.}
\tightlist
\item
  compute entropy, cross-entropy, KL divergence, and perplexity for
  discrete token distributions;
\item
  explain why cross-entropy is the standard training loss for language
  models;
\item
  interpret perplexity as an exponential average branching factor;
\item
  use the entropy chain rule to decompose uncertainty in a sequence;
\item
  define conditional entropy and mutual information;
\item
  connect pointwise mutual information with word co-occurrence and
  embeddings;
\item
  explain how tokenization changes entropy and code length;
\item
  interpret KL regularization in fine-tuning and alignment;
\item
  implement information-theoretic quantities from scratch in Python;
\item
  critique language-model evaluations using information-theoretic
  reasoning.
\end{enumerate}

\section{21.2 Discrete distributions over
tokens}\label{discrete-distributions-over-tokens}

Let \(V=\{1,2,\ldots,d\}\) be a finite vocabulary. A token distribution
is a vector

\[
p=(p_1,\ldots,p_d) \in \Delta^{d-1},
\]

where

\[
p_j \ge 0, \qquad \sum_{j=1}^d p_j=1.
\]

Here \(\Delta^{d-1}\) is the probability simplex. In NLP, \(p_j\) may
represent the empirical frequency of token \(j\) in a corpus, the
conditional probability assigned by a language model to token \(j\), or
the predicted class probability in a classifier.

Given counts \(c_1,\ldots,c_d\) from a corpus, the empirical
distribution is

\[
\widehat p_j = \frac{c_j}{\sum_{r=1}^d c_r}.
\]

This is the maximum likelihood estimator for a categorical distribution.
It is also the simplest bridge between text and information theory: a
corpus becomes a point in a probability simplex.

\begin{tcolorbox}[enhanced jigsaw, arc=.35mm, bottomrule=.15mm, bottomtitle=1mm, breakable, colback=white, colbacktitle=quarto-callout-note-color!10!white, colframe=quarto-callout-note-color-frame, coltitle=black, left=2mm, leftrule=.75mm, opacityback=0, opacitybacktitle=0.6, rightrule=.15mm, title=\textcolor{quarto-callout-note-color}{\faInfo}\hspace{0.5em}{Logarithm base}, titlerule=0mm, toprule=.15mm, toptitle=1mm]

If logs use base \(2\), entropy is measured in bits. If logs use base
\(e\), entropy is measured in nats. Machine learning usually uses
natural logarithms because gradients are simpler. This chapter mostly
uses natural logarithms.

\end{tcolorbox}

\section{21.3 Entropy}\label{entropy-1}

Let \(X\) be a discrete random variable with distribution \(p\) on
\(V\). The entropy of \(X\) is

\[
H(X)=H(p)=-\sum_{j=1}^d p_j \log p_j.
\]

By convention, \(0\log 0=0\). Entropy measures the average uncertainty
before observing \(X\). Equivalently, it is the average information
content of the outcome:

\[
H(X)=E[-\log p_X].
\]

The quantity

\[
I(x)=-\log p(x)
\]

is called the surprisal or self-information of the event \(X=x\). Rare
events have high surprisal; common events have low surprisal.

\subsection{Basic properties}\label{basic-properties}

For a distribution on \(d\) outcomes,

\[
0 \le H(p) \le \log d.
\]

The lower bound occurs when all probability mass is concentrated on one
outcome. The upper bound occurs at the uniform distribution \(p_j=1/d\).
Thus entropy is largest when all tokens are equally likely.

This upper bound is easy to prove using KL divergence. Let \(u_j=1/d\)
be the uniform distribution. Because \(D_{\mathrm{KL}}(p\|u)\ge 0\),

\[
0 \le \sum_{j=1}^d p_j\log\frac{p_j}{1/d}
= \sum_{j=1}^d p_j\log p_j + \log d.
\]

Therefore,

\[
H(p) = -\sum_{j=1}^d p_j\log p_j \le \log d.
\]

The inequality says that no distribution on \(d\) tokens has more
uncertainty than the uniform distribution.

\section{21.4 Cross-entropy and the log
score}\label{cross-entropy-and-the-log-score}

Suppose the true token distribution is \(p\), but a model predicts
distribution \(q\). The cross-entropy of \(p\) relative to \(q\) is

\[
H(p,q) = -\sum_{j=1}^d p_j \log q_j.
\]

It is the expected negative log-probability assigned by \(q\) when the
data are generated from \(p\):

\[
H(p,q)=E_{X\sim p}[-\log q_X].
\]

This is exactly the population version of the language-model loss. Given
tokens \(x_1,\ldots,x_n\), the empirical cross-entropy loss is

\[
\widehat H(p,q) = -\frac{1}{n}\sum_{i=1}^n \log q_{x_i}.
\]

If \(q=q_\theta(\cdot \mid \text{context})\) is a neural language model,
then minimizing cross-entropy is the same as maximizing likelihood.

\subsection{Why the negative log
score?}\label{why-the-negative-log-score}

The negative log score has several special properties. It rewards
assigning high probability to the observed token, but it also punishes
overconfident mistakes severely. If \(q_x\) is close to zero for the
observed token \(x\), then \(-\log q_x\) is very large.

For one-hot target vector \(e_y\) and predicted distribution \(q\), the
cross-entropy loss is

\[
\ell(q,y)=-\sum_{j=1}^d (e_y)_j\log q_j=-\log q_y.
\]

This is the loss used in next-token prediction.

\section{21.5 KL divergence}\label{kl-divergence-1}

The Kullback--Leibler divergence from \(q\) to \(p\) is

\[
D_{\mathrm{KL}}(p\|q)=\sum_{j=1}^d p_j\log\frac{p_j}{q_j}.
\]

It is not symmetric, so it is not a metric. However, it has a crucial
interpretation:

\[
D_{\mathrm{KL}}(p\|q) = H(p,q)-H(p).
\]

Thus KL divergence is the extra expected log loss incurred by predicting
with \(q\) when the true distribution is \(p\). Since \(H(p)\) does not
depend on \(q\), minimizing cross-entropy over \(q\) is equivalent to
minimizing KL divergence from the data distribution to the model
distribution.

\subsection{Gibbs inequality}\label{gibbs-inequality}

For distributions \(p\) and \(q\) with \(q_j>0\) whenever \(p_j>0\),

\[
D_{\mathrm{KL}}(p\|q)\ge 0,
\]

with equality if and only if \(p=q\). One proof uses the concavity of
\(\log\). Let \(Z=q_X/p_X\) when \(X\sim p\). Then

\[
-D_{\mathrm{KL}}(p\|q)
= \sum_j p_j\log\frac{q_j}{p_j}
\le \log\left(\sum_j p_j\frac{q_j}{p_j}\right)
=\log\left(\sum_j q_j\right)=0.
\]

Therefore \(D_{\mathrm{KL}}(p\|q)\ge 0\).

\begin{tcolorbox}[enhanced jigsaw, arc=.35mm, bottomrule=.15mm, bottomtitle=1mm, breakable, colback=white, colbacktitle=quarto-callout-warning-color!10!white, colframe=quarto-callout-warning-color-frame, coltitle=black, left=2mm, leftrule=.75mm, opacityback=0, opacitybacktitle=0.6, rightrule=.15mm, title=\textcolor{quarto-callout-warning-color}{\faExclamationTriangle}\hspace{0.5em}{Support mismatch}, titlerule=0mm, toprule=.15mm, toptitle=1mm]

If \(p_j>0\) but \(q_j=0\), then \(D_{\mathrm{KL}}(p\|q)=\infty\). This
is the information-theoretic version of the zero-probability problem in
n-gram language models. Smoothing is necessary because a model should
not assign zero probability to plausible future tokens.

\end{tcolorbox}

\section{21.6 Perplexity}\label{perplexity-1}

For a sequence of \(n\) observed tokens \(x_1,\ldots,x_n\), suppose a
language model assigns conditional probabilities

\[
q_i = q_\theta(x_i\mid x_{<i}).
\]

The average negative log-likelihood is

\[
\operatorname{NLL}_{\text{avg}} = -\frac{1}{n}\sum_{i=1}^n \log q_i.
\]

The perplexity is

\[
\operatorname{PPL}=\exp\left(\operatorname{NLL}_{\text{avg}}\right).
\]

When natural logarithms are used, perplexity is the exponential of
average cross-entropy in nats. When base-\(2\) logarithms are used,

\[
\operatorname{PPL}=2^{H_2},
\]

where \(H_2\) is measured in bits.

Perplexity can be interpreted as an average branching factor. A
perplexity of \(20\) means the model is, in an information-theoretic
sense, as uncertain as if it were choosing uniformly among about \(20\)
equally likely tokens at each step.

\begin{tcolorbox}[enhanced jigsaw, arc=.35mm, bottomrule=.15mm, bottomtitle=1mm, breakable, colback=white, colbacktitle=quarto-callout-note-color!10!white, colframe=quarto-callout-note-color-frame, coltitle=black, left=2mm, leftrule=.75mm, opacityback=0, opacitybacktitle=0.6, rightrule=.15mm, title=\textcolor{quarto-callout-note-color}{\faInfo}\hspace{0.5em}{Perplexity depends on tokenization}, titlerule=0mm, toprule=.15mm, toptitle=1mm]

Perplexity is not directly comparable across different tokenizers. A
character-level model, word-level model, and subword-level model define
different random variables. Changing the tokenization changes the unit
of prediction and therefore changes the entropy scale.

\end{tcolorbox}

\section{21.7 Sequence entropy and the chain
rule}\label{sequence-entropy-and-the-chain-rule}

Let \(W_{1:T}=(W_1,\ldots,W_T)\) be a sequence of tokens. The joint
entropy is

\[
H(W_{1:T})=-\sum_{w_{1:T}}p(w_{1:T})\log p(w_{1:T}).
\]

The entropy chain rule states that

\[
H(W_{1:T})=\sum_{t=1}^T H(W_t\mid W_{1:t-1}).
\]

This identity matches the autoregressive factorization

\[
p(w_{1:T})=\prod_{t=1}^T p(w_t\mid w_{1:t-1}).
\]

Taking negative logs and expectations gives the entropy chain rule. Thus
autoregressive language modeling is not just an engineering trick. It is
aligned with the fundamental decomposition of uncertainty in a sequence.

\subsection{Entropy rate}\label{entropy-rate}

For a stationary stochastic process \(W_1,W_2,\ldots\), the entropy rate
is

\[
\bar H = \lim_{T\to\infty}\frac{1}{T}H(W_{1:T}),
\]

when the limit exists. Equivalently, under suitable conditions,

\[
\bar H = \lim_{T\to\infty}H(W_T\mid W_{1:T-1}).
\]

For language, the entropy rate measures the long-run average uncertainty
per token after using all previous context. Better context modeling
should reduce conditional entropy.

\section{21.8 Conditional entropy and mutual
information}\label{conditional-entropy-and-mutual-information}

For random variables \(X\) and \(Y\), the conditional entropy is

\[
H(Y\mid X) = -\sum_x p(x)\sum_y p(y\mid x)\log p(y\mid x).
\]

It measures the remaining uncertainty in \(Y\) after observing \(X\).
The mutual information between \(X\) and \(Y\) is

\[
I(X;Y)=H(Y)-H(Y\mid X).
\]

It is also

\[
I(X;Y)=\sum_{x,y}p(x,y)\log\frac{p(x,y)}{p(x)p(y)}.
\]

Therefore,

\[
I(X;Y)=D_{\mathrm{KL}}(p(x,y)\|p(x)p(y)).
\]

Mutual information measures how far the joint distribution is from
independence. If \(X\) and \(Y\) are independent, then
\(p(x,y)=p(x)p(y)\) and \(I(X;Y)=0\).

\subsection{Mutual information in NLP}\label{mutual-information-in-nlp}

Mutual information appears naturally in several places:

\begin{itemize}
\tightlist
\item
  word co-occurrence: how much does seeing one word tell us about a
  nearby word?
\item
  classification: how informative is a token about a label?
\item
  representation learning: how much information does an embedding retain
  about input or target variables?
\item
  retrieval: how informative is a document about a query?
\item
  prompting: how much does an instruction reduce uncertainty about the
  desired answer?
\end{itemize}

\section{21.9 Pointwise mutual information and word
co-occurrence}\label{pointwise-mutual-information-and-word-co-occurrence}

For two events \(X=x\) and \(Y=y\), the pointwise mutual information is

\[
\operatorname{PMI}(x,y)=\log\frac{p(x,y)}{p(x)p(y)}.
\]

PMI is positive when \(x\) and \(y\) co-occur more often than expected
under independence. It is negative when they co-occur less often than
expected.

In word-context analysis, we may let \(X\) be a center word and \(Y\) be
a context word. Then PMI becomes

\[
\operatorname{PMI}(w,c)=\log\frac{p(w,c)}{p(w)p(c)}.
\]

The positive part is called positive PMI:

\[
\operatorname{PPMI}(w,c)=\max\{\operatorname{PMI}(w,c),0\}.
\]

Chapter 8 used PPMI matrices for word embeddings. This chapter explains
why the ratio is meaningful: it compares observed co-occurrence with
what independence would predict.

\section{21.10 Compression
interpretation}\label{compression-interpretation}

Information theory connects probability and coding. If an event \(x\)
has probability \(p(x)\), an ideal code assigns it length approximately

\[
\ell(x)=-\log_2 p(x)
\]

bits. Frequent events receive shorter codes; rare events receive longer
codes. The expected code length is

\[
E[\ell(X)] = H_2(X).
\]

If we encode data generated from \(p\) using a code designed for \(q\),
the expected code length becomes

\[
H_2(p,q)=H_2(p)+D_{\mathrm{KL},2}(p\|q).
\]

Thus KL divergence is the excess number of bits per token caused by
using the wrong probability model. This gives a concrete interpretation
of language-model loss:

\begin{quote}
Lower cross-entropy means the model gives a shorter average code for the
text.
\end{quote}

This viewpoint is useful for understanding why language modeling is such
a general pretraining task. To compress text well, a model must learn
spelling, syntax, semantics, style, factual regularities, and long-range
dependencies.

\section{21.11 Tokenization and
information}\label{tokenization-and-information}

Tokenization changes the random variable being modeled. Suppose the same
text can be represented as characters, words, or subwords. Each
representation gives a different sequence:

\[
C_{1:m}, \qquad W_{1:n}, \qquad S_{1:r}.
\]

The entropy per token is not directly comparable across these
representations because the token units are different. A word-level
model may have a large vocabulary and high entropy per token. A
character-level model may have a small vocabulary and lower entropy per
token, but many more tokens per sentence.

A more meaningful comparison is often code length per original character
or per original byte:

\[
\text{bits per character} = \frac{\text{total code length in bits}}{\text{number of characters}}.
\]

This is why compression-based comparisons must specify the coding unit.

\section{21.12 KL divergence in fine-tuning and
alignment}\label{kl-divergence-in-fine-tuning-and-alignment}

In Chapter 18, we saw objectives that include a KL penalty. A common
form is

\[
\max_\theta E_{y\sim \pi_\theta(\cdot\mid x)}[r(x,y)]
-\beta D_{\mathrm{KL}}(\pi_\theta(\cdot\mid x)\|\pi_0(\cdot\mid x)).
\]

Here \(\pi_0\) is a reference model and \(\pi_\theta\) is the fine-tuned
model. The reward term encourages helpful behavior. The KL term
discourages the new model from moving too far from the reference
distribution.

Mathematically, the KL penalty is a distributional trust region. It does
not simply penalize parameter distance. Two models may have very
different parameter vectors but similar output distributions, or similar
parameters but different output distributions. KL regularization
measures distance in probability behavior.

\section{21.13 Information bottleneck
viewpoint}\label{information-bottleneck-viewpoint}

Let \(X\) be an input sequence, \(Z\) be a learned representation, and
\(Y\) be a target. The information bottleneck principle asks for a
representation \(Z\) that is useful for predicting \(Y\) but compressed
relative to \(X\). One idealized objective is

\[
\min \ I(X;Z) - \beta I(Z;Y).
\]

The first term encourages compression: \(Z\) should not retain
irrelevant details of \(X\). The second term encourages prediction:
\(Z\) should retain information about \(Y\).

This viewpoint is not a complete theory of large language models, but it
provides useful language for asking mathematical questions:

\begin{itemize}
\tightlist
\item
  Which information about the input is preserved in an embedding?
\item
  Which information is discarded by pooling, truncation, or retrieval?
\item
  How does fine-tuning change the information available for a downstream
  task?
\item
  How can we measure whether a representation contains syntactic,
  semantic, or factual information?
\end{itemize}

\section{21.14 Interactive: entropy, KL divergence, and
perplexity}\label{interactive-entropy-kl-divergence-and-perplexity}

The following interactive changes a predicted distribution \(q\) over
four possible next tokens while keeping a fixed empirical distribution
\(p\). It shows how entropy, cross-entropy, KL divergence, and
perplexity respond.

\section{21.15 Python example: empirical
entropy}\label{python-example-empirical-entropy}

The first example computes token counts, an empirical distribution, and
entropy from scratch.

\begin{Shaded}
\begin{Highlighting}[]
\ImportTok{from}\NormalTok{ collections }\ImportTok{import}\NormalTok{ Counter}
\ImportTok{import}\NormalTok{ math}

\NormalTok{text }\OperatorTok{=} \StringTok{"machine learning models learn from data and language models learn from text data"}
\NormalTok{tokens }\OperatorTok{=}\NormalTok{ text.split()}
\NormalTok{counts }\OperatorTok{=}\NormalTok{ Counter(tokens)}
\NormalTok{N }\OperatorTok{=} \BuiltInTok{sum}\NormalTok{(counts.values())}
\NormalTok{probs }\OperatorTok{=}\NormalTok{ \{w: c }\OperatorTok{/}\NormalTok{ N }\ControlFlowTok{for}\NormalTok{ w, c }\KeywordTok{in}\NormalTok{ counts.items()\}}

\NormalTok{entropy }\OperatorTok{=} \OperatorTok{{-}}\BuiltInTok{sum}\NormalTok{(p }\OperatorTok{*}\NormalTok{ math.log(p) }\ControlFlowTok{for}\NormalTok{ p }\KeywordTok{in}\NormalTok{ probs.values())}
\NormalTok{perplexity }\OperatorTok{=}\NormalTok{ math.exp(entropy)}

\BuiltInTok{print}\NormalTok{(}\StringTok{"Counts:"}\NormalTok{, counts)}
\BuiltInTok{print}\NormalTok{(}\StringTok{"Empirical probabilities:"}\NormalTok{)}
\ControlFlowTok{for}\NormalTok{ w, p }\KeywordTok{in} \BuiltInTok{sorted}\NormalTok{(probs.items()):}
    \BuiltInTok{print}\NormalTok{(}\SpecialStringTok{f"  }\SpecialCharTok{\{}\NormalTok{w}\SpecialCharTok{:10s\}}\SpecialStringTok{ }\SpecialCharTok{\{}\NormalTok{p}\SpecialCharTok{:.3f\}}\SpecialStringTok{"}\NormalTok{)}
\BuiltInTok{print}\NormalTok{(}\StringTok{"Entropy in nats:"}\NormalTok{, }\BuiltInTok{round}\NormalTok{(entropy, }\DecValTok{4}\NormalTok{))}
\BuiltInTok{print}\NormalTok{(}\StringTok{"Entropy in bits:"}\NormalTok{, }\BuiltInTok{round}\NormalTok{(entropy }\OperatorTok{/}\NormalTok{ math.log(}\DecValTok{2}\NormalTok{), }\DecValTok{4}\NormalTok{))}
\BuiltInTok{print}\NormalTok{(}\StringTok{"Empirical perplexity:"}\NormalTok{, }\BuiltInTok{round}\NormalTok{(perplexity, }\DecValTok{4}\NormalTok{))}
\end{Highlighting}
\end{Shaded}

\begin{verbatim}
Counts: Counter({'models': 2, 'learn': 2, 'from': 2, 'data': 2, 'machine': 1, 'learning': 1, 'and': 1, 'language': 1, 'text': 1})
Empirical probabilities:
  and        0.077
  data       0.154
  from       0.154
  language   0.077
  learn      0.154
  learning   0.077
  machine    0.077
  models     0.154
  text       0.077
Entropy in nats: 2.1384
Entropy in bits: 3.0851
Empirical perplexity: 8.4858
\end{verbatim}

Entropy is not merely a vocabulary-size measure. It depends on the full
distribution. A corpus with many rare tokens can have high vocabulary
size but still lower entropy if a few tokens dominate.

\section{21.16 Python example: cross-entropy, KL divergence, and
smoothing}\label{python-example-cross-entropy-kl-divergence-and-smoothing}

Now compare an empirical distribution \(p\) with two model distributions
\(q\).

\begin{Shaded}
\begin{Highlighting}[]
\ImportTok{import}\NormalTok{ numpy }\ImportTok{as}\NormalTok{ np}

\NormalTok{vocab }\OperatorTok{=}\NormalTok{ [}\StringTok{"machine"}\NormalTok{, }\StringTok{"language"}\NormalTok{, }\StringTok{"data"}\NormalTok{, }\StringTok{"model"}\NormalTok{]}
\NormalTok{p }\OperatorTok{=}\NormalTok{ np.array([}\FloatTok{0.45}\NormalTok{, }\FloatTok{0.30}\NormalTok{, }\FloatTok{0.20}\NormalTok{, }\FloatTok{0.05}\NormalTok{])}
\NormalTok{q\_good }\OperatorTok{=}\NormalTok{ np.array([}\FloatTok{0.40}\NormalTok{, }\FloatTok{0.32}\NormalTok{, }\FloatTok{0.20}\NormalTok{, }\FloatTok{0.08}\NormalTok{])}
\NormalTok{q\_bad }\OperatorTok{=}\NormalTok{ np.array([}\FloatTok{0.10}\NormalTok{, }\FloatTok{0.20}\NormalTok{, }\FloatTok{0.30}\NormalTok{, }\FloatTok{0.40}\NormalTok{])}

\KeywordTok{def}\NormalTok{ entropy\_np(p):}
\NormalTok{    p }\OperatorTok{=}\NormalTok{ np.asarray(p, dtype}\OperatorTok{=}\BuiltInTok{float}\NormalTok{)}
    \ControlFlowTok{return} \OperatorTok{{-}}\NormalTok{np.}\BuiltInTok{sum}\NormalTok{(np.where(p }\OperatorTok{\textgreater{}} \DecValTok{0}\NormalTok{, p }\OperatorTok{*}\NormalTok{ np.log(p), }\FloatTok{0.0}\NormalTok{))}

\KeywordTok{def}\NormalTok{ cross\_entropy\_np(p, q):}
\NormalTok{    p }\OperatorTok{=}\NormalTok{ np.asarray(p, dtype}\OperatorTok{=}\BuiltInTok{float}\NormalTok{)}
\NormalTok{    q }\OperatorTok{=}\NormalTok{ np.asarray(q, dtype}\OperatorTok{=}\BuiltInTok{float}\NormalTok{)}
    \ControlFlowTok{return} \OperatorTok{{-}}\NormalTok{np.}\BuiltInTok{sum}\NormalTok{(p }\OperatorTok{*}\NormalTok{ np.log(q))}

\KeywordTok{def}\NormalTok{ kl\_np(p, q):}
    \ControlFlowTok{return}\NormalTok{ cross\_entropy\_np(p, q) }\OperatorTok{{-}}\NormalTok{ entropy\_np(p)}

\ControlFlowTok{for}\NormalTok{ name, q }\KeywordTok{in}\NormalTok{ [(}\StringTok{"good"}\NormalTok{, q\_good), (}\StringTok{"bad"}\NormalTok{, q\_bad)]:}
    \BuiltInTok{print}\NormalTok{(name)}
    \BuiltInTok{print}\NormalTok{(}\StringTok{"  cross entropy:"}\NormalTok{, }\BuiltInTok{round}\NormalTok{(cross\_entropy\_np(p, q), }\DecValTok{4}\NormalTok{))}
    \BuiltInTok{print}\NormalTok{(}\StringTok{"  KL divergence:"}\NormalTok{, }\BuiltInTok{round}\NormalTok{(kl\_np(p, q), }\DecValTok{4}\NormalTok{))}
    \BuiltInTok{print}\NormalTok{(}\StringTok{"  perplexity:"}\NormalTok{, }\BuiltInTok{round}\NormalTok{(np.exp(cross\_entropy\_np(p, q)), }\DecValTok{4}\NormalTok{))}
\end{Highlighting}
\end{Shaded}

\begin{verbatim}
good
  cross entropy: 1.2023
  KL divergence: 0.0101
  perplexity: 3.3279
bad
  cross entropy: 1.8056
  KL divergence: 0.6134
  perplexity: 6.0836
\end{verbatim}

If a model assigns zero probability to an event with positive empirical
probability, the cross-entropy becomes infinite. Additive smoothing
prevents this.

\begin{Shaded}
\begin{Highlighting}[]
\NormalTok{counts }\OperatorTok{=}\NormalTok{ np.array([}\DecValTok{9}\NormalTok{, }\DecValTok{6}\NormalTok{, }\DecValTok{4}\NormalTok{, }\DecValTok{1}\NormalTok{])}
\NormalTok{alpha }\OperatorTok{=} \FloatTok{0.5}
\NormalTok{p\_unsmoothed }\OperatorTok{=}\NormalTok{ counts }\OperatorTok{/}\NormalTok{ counts.}\BuiltInTok{sum}\NormalTok{()}
\NormalTok{p\_smoothed }\OperatorTok{=}\NormalTok{ (counts }\OperatorTok{+}\NormalTok{ alpha) }\OperatorTok{/}\NormalTok{ (counts.}\BuiltInTok{sum}\NormalTok{() }\OperatorTok{+}\NormalTok{ alpha }\OperatorTok{*} \BuiltInTok{len}\NormalTok{(counts))}

\BuiltInTok{print}\NormalTok{(}\StringTok{"Unsmoothed:"}\NormalTok{, np.}\BuiltInTok{round}\NormalTok{(p\_unsmoothed, }\DecValTok{3}\NormalTok{))}
\BuiltInTok{print}\NormalTok{(}\StringTok{"Smoothed:  "}\NormalTok{, np.}\BuiltInTok{round}\NormalTok{(p\_smoothed, }\DecValTok{3}\NormalTok{))}
\BuiltInTok{print}\NormalTok{(}\StringTok{"Entropy unsmoothed:"}\NormalTok{, }\BuiltInTok{round}\NormalTok{(entropy\_np(p\_unsmoothed), }\DecValTok{4}\NormalTok{))}
\BuiltInTok{print}\NormalTok{(}\StringTok{"Entropy smoothed:"}\NormalTok{, }\BuiltInTok{round}\NormalTok{(entropy\_np(p\_smoothed), }\DecValTok{4}\NormalTok{))}
\end{Highlighting}
\end{Shaded}

\begin{verbatim}
Unsmoothed: [0.45 0.3  0.2  0.05]
Smoothed:   [0.432 0.295 0.205 0.068]
Entropy unsmoothed: 1.1922
Entropy smoothed: 1.2306
\end{verbatim}

\section{21.17 Python example: sequence entropy for a bigram
model}\label{python-example-sequence-entropy-for-a-bigram-model}

For a bigram Markov chain, conditional entropy is computed row by row.
If \(P\) is a transition matrix and \(\pi\) is the stationary or
empirical distribution over previous tokens, then

\[
H(W_t\mid W_{t-1})=\sum_i \pi_i H(P_{i,\cdot}).
\]

\begin{Shaded}
\begin{Highlighting}[]
\ImportTok{import}\NormalTok{ numpy }\ImportTok{as}\NormalTok{ np}

\NormalTok{vocab }\OperatorTok{=}\NormalTok{ [}\StringTok{"\textless{}s\textgreater{}"}\NormalTok{, }\StringTok{"the"}\NormalTok{, }\StringTok{"model"}\NormalTok{, }\StringTok{"learns"}\NormalTok{]}
\NormalTok{P }\OperatorTok{=}\NormalTok{ np.array([}
\NormalTok{    [}\FloatTok{0.00}\NormalTok{, }\FloatTok{0.70}\NormalTok{, }\FloatTok{0.20}\NormalTok{, }\FloatTok{0.10}\NormalTok{],}
\NormalTok{    [}\FloatTok{0.00}\NormalTok{, }\FloatTok{0.05}\NormalTok{, }\FloatTok{0.80}\NormalTok{, }\FloatTok{0.15}\NormalTok{],}
\NormalTok{    [}\FloatTok{0.00}\NormalTok{, }\FloatTok{0.10}\NormalTok{, }\FloatTok{0.10}\NormalTok{, }\FloatTok{0.80}\NormalTok{],}
\NormalTok{    [}\FloatTok{0.00}\NormalTok{, }\FloatTok{0.60}\NormalTok{, }\FloatTok{0.20}\NormalTok{, }\FloatTok{0.20}\NormalTok{]}
\NormalTok{])}
\NormalTok{pi }\OperatorTok{=}\NormalTok{ np.array([}\FloatTok{0.10}\NormalTok{, }\FloatTok{0.35}\NormalTok{, }\FloatTok{0.30}\NormalTok{, }\FloatTok{0.25}\NormalTok{])}

\KeywordTok{def}\NormalTok{ row\_entropy(row):}
\NormalTok{    positive }\OperatorTok{=}\NormalTok{ row[row }\OperatorTok{\textgreater{}} \DecValTok{0}\NormalTok{]}
    \ControlFlowTok{return} \OperatorTok{{-}}\NormalTok{np.}\BuiltInTok{sum}\NormalTok{(positive }\OperatorTok{*}\NormalTok{ np.log(positive))}

\NormalTok{row\_entropies }\OperatorTok{=}\NormalTok{ np.array([row\_entropy(row) }\ControlFlowTok{for}\NormalTok{ row }\KeywordTok{in}\NormalTok{ P])}
\NormalTok{conditional\_entropy }\OperatorTok{=}\NormalTok{ np.}\BuiltInTok{sum}\NormalTok{(pi }\OperatorTok{*}\NormalTok{ row\_entropies)}

\ControlFlowTok{for}\NormalTok{ word, h }\KeywordTok{in} \BuiltInTok{zip}\NormalTok{(vocab, row\_entropies):}
    \BuiltInTok{print}\NormalTok{(}\SpecialStringTok{f"H(next | }\SpecialCharTok{\{}\NormalTok{word}\SpecialCharTok{:6s\}}\SpecialStringTok{) = }\SpecialCharTok{\{}\NormalTok{h}\SpecialCharTok{:.4f\}}\SpecialStringTok{ nats"}\NormalTok{)}
\BuiltInTok{print}\NormalTok{(}\StringTok{"Average conditional entropy:"}\NormalTok{, }\BuiltInTok{round}\NormalTok{(conditional\_entropy, }\DecValTok{4}\NormalTok{))}
\BuiltInTok{print}\NormalTok{(}\StringTok{"Conditional perplexity:"}\NormalTok{, }\BuiltInTok{round}\NormalTok{(np.exp(conditional\_entropy), }\DecValTok{4}\NormalTok{))}
\end{Highlighting}
\end{Shaded}

\begin{verbatim}
H(next | <s>   ) = 0.8018 nats
H(next | the   ) = 0.6129 nats
H(next | model ) = 0.6390 nats
H(next | learns) = 0.9503 nats
Average conditional entropy: 0.724
Conditional perplexity: 2.0626
\end{verbatim}

This calculation shows why context matters. If the previous token makes
the next-token distribution sharper, then conditional entropy decreases.

\section{21.18 Python example: PMI matrix from co-occurrence
counts}\label{python-example-pmi-matrix-from-co-occurrence-counts}

The following example computes a PMI matrix from a small word-context
count table.

\begin{Shaded}
\begin{Highlighting}[]
\ImportTok{import}\NormalTok{ numpy }\ImportTok{as}\NormalTok{ np}
\ImportTok{import}\NormalTok{ pandas }\ImportTok{as}\NormalTok{ pd}

\NormalTok{words }\OperatorTok{=}\NormalTok{ [}\StringTok{"king"}\NormalTok{, }\StringTok{"queen"}\NormalTok{, }\StringTok{"bank"}\NormalTok{, }\StringTok{"river"}\NormalTok{]}
\NormalTok{contexts }\OperatorTok{=}\NormalTok{ [}\StringTok{"royal"}\NormalTok{, }\StringTok{"money"}\NormalTok{, }\StringTok{"water"}\NormalTok{]}
\NormalTok{C }\OperatorTok{=}\NormalTok{ np.array([}
\NormalTok{    [}\DecValTok{9}\NormalTok{, }\DecValTok{1}\NormalTok{, }\DecValTok{0}\NormalTok{],}
\NormalTok{    [}\DecValTok{8}\NormalTok{, }\DecValTok{1}\NormalTok{, }\DecValTok{0}\NormalTok{],}
\NormalTok{    [}\DecValTok{0}\NormalTok{, }\DecValTok{7}\NormalTok{, }\DecValTok{4}\NormalTok{],}
\NormalTok{    [}\DecValTok{0}\NormalTok{, }\DecValTok{1}\NormalTok{, }\DecValTok{9}\NormalTok{]}
\NormalTok{], dtype}\OperatorTok{=}\BuiltInTok{float}\NormalTok{)}

\NormalTok{N }\OperatorTok{=}\NormalTok{ C.}\BuiltInTok{sum}\NormalTok{()}
\NormalTok{p\_wc }\OperatorTok{=}\NormalTok{ C }\OperatorTok{/}\NormalTok{ N}
\NormalTok{p\_w }\OperatorTok{=}\NormalTok{ p\_wc.}\BuiltInTok{sum}\NormalTok{(axis}\OperatorTok{=}\DecValTok{1}\NormalTok{, keepdims}\OperatorTok{=}\VariableTok{True}\NormalTok{)}
\NormalTok{p\_c }\OperatorTok{=}\NormalTok{ p\_wc.}\BuiltInTok{sum}\NormalTok{(axis}\OperatorTok{=}\DecValTok{0}\NormalTok{, keepdims}\OperatorTok{=}\VariableTok{True}\NormalTok{)}

\ControlFlowTok{with}\NormalTok{ np.errstate(divide}\OperatorTok{=}\StringTok{"ignore"}\NormalTok{, invalid}\OperatorTok{=}\StringTok{"ignore"}\NormalTok{):}
\NormalTok{    PMI }\OperatorTok{=}\NormalTok{ np.log(p\_wc }\OperatorTok{/}\NormalTok{ (p\_w }\OperatorTok{@}\NormalTok{ p\_c))}
\NormalTok{PMI }\OperatorTok{=}\NormalTok{ np.where(C }\OperatorTok{\textgreater{}} \DecValTok{0}\NormalTok{, PMI, np.nan)}
\NormalTok{PPMI }\OperatorTok{=}\NormalTok{ np.where(np.isnan(PMI), }\FloatTok{0.0}\NormalTok{, np.maximum(PMI, }\FloatTok{0.0}\NormalTok{))}

\BuiltInTok{print}\NormalTok{(}\StringTok{"PMI matrix"}\NormalTok{)}
\NormalTok{display(pd.DataFrame(PMI, index}\OperatorTok{=}\NormalTok{words, columns}\OperatorTok{=}\NormalTok{contexts).}\BuiltInTok{round}\NormalTok{(}\DecValTok{3}\NormalTok{))}
\BuiltInTok{print}\NormalTok{(}\StringTok{"PPMI matrix"}\NormalTok{)}
\NormalTok{display(pd.DataFrame(PPMI, index}\OperatorTok{=}\NormalTok{words, columns}\OperatorTok{=}\NormalTok{contexts).}\BuiltInTok{round}\NormalTok{(}\DecValTok{3}\NormalTok{))}
\end{Highlighting}
\end{Shaded}

\begin{verbatim}
PMI matrix
\end{verbatim}

{\def\LTcaptype{none} % do not increment counter
\begin{longtable}[]{@{}llll@{}}
\toprule\noalign{}
& royal & money & water \\
\midrule\noalign{}
\endhead
\bottomrule\noalign{}
\endlastfoot
king & 0.750 & -0.916 & NaN \\
queen & 0.738 & -0.811 & NaN \\
bank & NaN & 0.934 & 0.112 \\
river & NaN & -0.916 & 1.019 \\
\end{longtable}
}

\begin{verbatim}
PPMI matrix
\end{verbatim}

{\def\LTcaptype{none} % do not increment counter
\begin{longtable}[]{@{}llll@{}}
\toprule\noalign{}
& royal & money & water \\
\midrule\noalign{}
\endhead
\bottomrule\noalign{}
\endlastfoot
king & 0.750 & 0.000 & 0.000 \\
queen & 0.738 & 0.000 & 0.000 \\
bank & 0.000 & 0.934 & 0.112 \\
river & 0.000 & 0.000 & 1.019 \\
\end{longtable}
}

Positive PMI highlights associations beyond chance. This is one reason
PMI-like transformations are useful before low-rank factorization.

\section{21.19 Python example: mutual information between token and
label}\label{python-example-mutual-information-between-token-and-label}

Suppose \(X\) is the indicator of whether a token appears in a document
and \(Y\) is a binary sentiment label. We can compute \(I(X;Y)\) from a
\(2\times 2\) table.

\begin{Shaded}
\begin{Highlighting}[]
\ImportTok{import}\NormalTok{ numpy }\ImportTok{as}\NormalTok{ np}

\KeywordTok{def}\NormalTok{ mutual\_information\_from\_table(table):}
\NormalTok{    table }\OperatorTok{=}\NormalTok{ np.asarray(table, dtype}\OperatorTok{=}\BuiltInTok{float}\NormalTok{)}
\NormalTok{    pxy }\OperatorTok{=}\NormalTok{ table }\OperatorTok{/}\NormalTok{ table.}\BuiltInTok{sum}\NormalTok{()}
\NormalTok{    px }\OperatorTok{=}\NormalTok{ pxy.}\BuiltInTok{sum}\NormalTok{(axis}\OperatorTok{=}\DecValTok{1}\NormalTok{, keepdims}\OperatorTok{=}\VariableTok{True}\NormalTok{)}
\NormalTok{    py }\OperatorTok{=}\NormalTok{ pxy.}\BuiltInTok{sum}\NormalTok{(axis}\OperatorTok{=}\DecValTok{0}\NormalTok{, keepdims}\OperatorTok{=}\VariableTok{True}\NormalTok{)}
\NormalTok{    expected }\OperatorTok{=}\NormalTok{ px }\OperatorTok{@}\NormalTok{ py}
\NormalTok{    mask }\OperatorTok{=}\NormalTok{ pxy }\OperatorTok{\textgreater{}} \DecValTok{0}
    \ControlFlowTok{return}\NormalTok{ np.}\BuiltInTok{sum}\NormalTok{(pxy[mask] }\OperatorTok{*}\NormalTok{ np.log(pxy[mask] }\OperatorTok{/}\NormalTok{ expected[mask]))}

\CommentTok{\# rows: token absent/present, columns: negative/positive}
\NormalTok{excellent\_table }\OperatorTok{=}\NormalTok{ np.array([}
\NormalTok{    [}\DecValTok{40}\NormalTok{, }\DecValTok{20}\NormalTok{],}
\NormalTok{    [}\DecValTok{5}\NormalTok{, }\DecValTok{35}\NormalTok{]}
\NormalTok{])}
\NormalTok{the\_table }\OperatorTok{=}\NormalTok{ np.array([}
\NormalTok{    [}\DecValTok{10}\NormalTok{, }\DecValTok{12}\NormalTok{],}
\NormalTok{    [}\DecValTok{35}\NormalTok{, }\DecValTok{43}\NormalTok{]}
\NormalTok{])}

\BuiltInTok{print}\NormalTok{(}\StringTok{"I(excellent present; label):"}\NormalTok{, }\BuiltInTok{round}\NormalTok{(mutual\_information\_from\_table(excellent\_table), }\DecValTok{4}\NormalTok{))}
\BuiltInTok{print}\NormalTok{(}\StringTok{"I(the present; label):      "}\NormalTok{, }\BuiltInTok{round}\NormalTok{(mutual\_information\_from\_table(the\_table), }\DecValTok{4}\NormalTok{))}
\end{Highlighting}
\end{Shaded}

\begin{verbatim}
I(excellent present; label): 0.1555
I(the present; label):       0.0
\end{verbatim}

A sentiment-bearing word usually has higher mutual information with the
label than a function word. However, mutual information is still a
statistical estimate and can be biased upward for rare features.

\section{21.20 Python example: coding length and model
comparison}\label{python-example-coding-length-and-model-comparison}

For a held-out sequence, compare two language models by average code
length.

\begin{Shaded}
\begin{Highlighting}[]
\ImportTok{import}\NormalTok{ numpy }\ImportTok{as}\NormalTok{ np}

\NormalTok{observed }\OperatorTok{=}\NormalTok{ [}\StringTok{"data"}\NormalTok{, }\StringTok{"model"}\NormalTok{, }\StringTok{"data"}\NormalTok{, }\StringTok{"language"}\NormalTok{, }\StringTok{"model"}\NormalTok{, }\StringTok{"data"}\NormalTok{]}
\NormalTok{vocab }\OperatorTok{=}\NormalTok{ [}\StringTok{"data"}\NormalTok{, }\StringTok{"language"}\NormalTok{, }\StringTok{"model"}\NormalTok{]}

\NormalTok{q1 }\OperatorTok{=}\NormalTok{ \{}\StringTok{"data"}\NormalTok{: }\FloatTok{0.50}\NormalTok{, }\StringTok{"language"}\NormalTok{: }\FloatTok{0.20}\NormalTok{, }\StringTok{"model"}\NormalTok{: }\FloatTok{0.30}\NormalTok{\}}
\NormalTok{q2 }\OperatorTok{=}\NormalTok{ \{}\StringTok{"data"}\NormalTok{: }\FloatTok{0.25}\NormalTok{, }\StringTok{"language"}\NormalTok{: }\FloatTok{0.50}\NormalTok{, }\StringTok{"model"}\NormalTok{: }\FloatTok{0.25}\NormalTok{\}}

\KeywordTok{def}\NormalTok{ average\_code\_length\_bits(sequence, q):}
    \ControlFlowTok{return}\NormalTok{ np.mean([}\OperatorTok{{-}}\NormalTok{np.log2(q[token]) }\ControlFlowTok{for}\NormalTok{ token }\KeywordTok{in}\NormalTok{ sequence])}

\ControlFlowTok{for}\NormalTok{ name, q }\KeywordTok{in}\NormalTok{ [(}\StringTok{"model 1"}\NormalTok{, q1), (}\StringTok{"model 2"}\NormalTok{, q2)]:}
\NormalTok{    bits }\OperatorTok{=}\NormalTok{ average\_code\_length\_bits(observed, q)}
    \BuiltInTok{print}\NormalTok{(name)}
    \BuiltInTok{print}\NormalTok{(}\StringTok{"  average bits per token:"}\NormalTok{, }\BuiltInTok{round}\NormalTok{(bits, }\DecValTok{4}\NormalTok{))}
    \BuiltInTok{print}\NormalTok{(}\StringTok{"  perplexity:"}\NormalTok{, }\BuiltInTok{round}\NormalTok{(}\DecValTok{2} \OperatorTok{**}\NormalTok{ bits, }\DecValTok{4}\NormalTok{))}
\end{Highlighting}
\end{Shaded}

\begin{verbatim}
model 1
  average bits per token: 1.466
  perplexity: 2.7625
model 2
  average bits per token: 1.8333
  perplexity: 3.5636
\end{verbatim}

The model with shorter average code length is also the model with lower
held-out cross-entropy.

\section{21.21 AI-assisted learning
activities}\label{ai-assisted-learning-activities-4}

Use an AI assistant as a mathematical discussion partner, not as an
authority. Good prompts ask the model to expose assumptions and verify
calculations.

\subsection{Prompt 1: entropy audit}\label{prompt-1-entropy-audit}

\begin{Shaded}
\begin{Highlighting}[]
\NormalTok{I have a token distribution p = [0.5, 0.25, 0.15, 0.10].}
\NormalTok{Compute entropy in nats and bits.}
\NormalTok{Then explain which step uses the convention 0 log 0 = 0 and why that convention is mathematically reasonable.}
\end{Highlighting}
\end{Shaded}

\subsection{Prompt 2: cross-entropy versus KL
divergence}\label{prompt-2-cross-entropy-versus-kl-divergence}

\begin{Shaded}
\begin{Highlighting}[]
\NormalTok{Explain the identity H(p,q) = H(p) + KL(p || q) using a four{-}token example.}
\NormalTok{Do not skip the algebra. Then explain why minimizing cross{-}entropy over q is equivalent to minimizing KL divergence.}
\end{Highlighting}
\end{Shaded}

\subsection{Prompt 3: critique a perplexity
comparison}\label{prompt-3-critique-a-perplexity-comparison}

\begin{Shaded}
\begin{Highlighting}[]
\NormalTok{A paper reports that model A has perplexity 25 using a word tokenizer and model B has perplexity 12 using a subword tokenizer.}
\NormalTok{Can we conclude that model B is better? Give a careful mathematical answer.}
\end{Highlighting}
\end{Shaded}

\subsection{Prompt 4: PMI and
co-occurrence}\label{prompt-4-pmi-and-co-occurrence}

\begin{Shaded}
\begin{Highlighting}[]
\NormalTok{Construct a tiny word{-}context table where two words have high raw co{-}occurrence but low PMI because the context word is very frequent.}
\NormalTok{Explain the calculation.}
\end{Highlighting}
\end{Shaded}

\begin{tcolorbox}[enhanced jigsaw, arc=.35mm, bottomrule=.15mm, bottomtitle=1mm, breakable, colback=white, colbacktitle=quarto-callout-warning-color!10!white, colframe=quarto-callout-warning-color-frame, coltitle=black, left=2mm, leftrule=.75mm, opacityback=0, opacitybacktitle=0.6, rightrule=.15mm, title=\textcolor{quarto-callout-warning-color}{\faExclamationTriangle}\hspace{0.5em}{AI caution}, titlerule=0mm, toprule=.15mm, toptitle=1mm]

AI systems often confuse entropy, cross-entropy, and KL divergence
because the formulas look similar. Always check which distribution is
inside the expectation and which distribution appears inside the
logarithm.

\end{tcolorbox}

\section{21.22 Common mistakes}\label{common-mistakes-10}

\begin{enumerate}
\def\labelenumi{\arabic{enumi}.}
\tightlist
\item
  \textbf{Comparing perplexity across tokenizers.} Perplexity depends on
  the prediction unit.
\item
  \textbf{Treating KL divergence as a distance metric.} KL is asymmetric
  and does not satisfy the triangle inequality.
\item
  \textbf{Forgetting support conditions.} If \(q_j=0\) where \(p_j>0\),
  cross-entropy and KL divergence are infinite.
\item
  \textbf{Confusing entropy and cross-entropy.} Entropy uses the true
  distribution inside both the expectation and the logarithm.
  Cross-entropy uses the true distribution in the expectation and the
  model distribution in the logarithm.
\item
  \textbf{Overinterpreting PMI for rare events.} PMI can be large for
  very rare co-occurrences. Smoothing and minimum-count thresholds are
  often needed.
\item
  \textbf{Ignoring logarithm base.} Natural logs give nats. Base-\(2\)
  logs give bits.
\item
  \textbf{Assuming lower entropy always means better language.} Entropy
  measures uncertainty under a distribution, not semantic quality by
  itself.
\end{enumerate}

\section{21.23 Worked calculation}\label{worked-calculation-3}

Let

\[
p=(0.5,0.3,0.2), \qquad q=(0.4,0.4,0.2).
\]

Then

\[
H(p)=-0.5\log 0.5-0.3\log 0.3-0.2\log 0.2.
\]

Numerically,

\[
H(p) \approx 1.0297.
\]

The cross-entropy is

\[
H(p,q)=-0.5\log 0.4-0.3\log 0.4-0.2\log 0.2.
\]

Numerically,

\[
H(p,q)\approx 1.0559.
\]

Therefore,

\[
D_{\mathrm{KL}}(p\|q)=H(p,q)-H(p)\approx 0.0262.
\]

The perplexity associated with \(H(p,q)\) is

\[
\exp(1.0559)\approx 2.875.
\]

This means that under model \(q\), the average uncertainty is comparable
to choosing uniformly from about \(2.875\) options.

\section{21.24 Exercises}\label{exercises-20}

\subsection{Conceptual exercises}\label{conceptual-exercises-9}

\begin{enumerate}
\def\labelenumi{\arabic{enumi}.}
\tightlist
\item
  Explain why entropy is maximized by the uniform distribution on a
  finite vocabulary.
\item
  Explain why cross-entropy is always at least entropy.
\item
  Give an example showing that
  \(D_{\mathrm{KL}}(p\|q)\ne D_{\mathrm{KL}}(q\|p)\).
\item
  Explain why perplexity cannot be compared directly across different
  tokenizers.
\item
  Interpret \(I(X;Y)=0\) in terms of independence.
\item
  Explain why PMI can be unstable for rare events.
\end{enumerate}

\subsection{Computational exercises}\label{computational-exercises-6}

\begin{enumerate}
\def\labelenumi{\arabic{enumi}.}
\setcounter{enumi}{6}
\tightlist
\item
  Write a function that computes entropy, cross-entropy, KL divergence,
  and perplexity for two discrete distributions.
\item
  Given a small corpus, compute unigram entropy and bigram conditional
  entropy. Which is smaller? Why?
\item
  Build a word-context matrix for a small corpus and compute PMI and
  PPMI.
\item
  For a binary classification data set, compute mutual information
  between each token and the class label.
\item
  Compare two toy language models by average bits per token on a
  held-out sequence.
\item
  Simulate distributions with increasing vocabulary size and plot
  entropy as the distribution becomes more uniform.
\end{enumerate}

\subsection{Mathematical exercises}\label{mathematical-exercises-8}

\begin{enumerate}
\def\labelenumi{\arabic{enumi}.}
\setcounter{enumi}{12}
\tightlist
\item
  Prove Gibbs inequality using Jensen's inequality.
\item
  Prove the entropy chain rule for two variables:
  \(H(X,Y)=H(X)+H(Y\mid X)\).
\item
  Show that \(I(X;Y)=H(X)+H(Y)-H(X,Y)\).
\item
  Show that \(I(X;Y)=D_{\mathrm{KL}}(p(x,y)\|p(x)p(y))\).
\item
  Suppose \(X\to Z\to Y\) is a Markov chain. State and interpret the
  data processing inequality in the context of NLP representations.
\item
  For a Markov chain with transition matrix \(P\) and stationary
  distribution \(\pi\), derive the entropy rate formula
  \(\sum_i \pi_i H(P_{i,\cdot})\).
\end{enumerate}

\section{21.25 Mini-project: information profile of a
corpus}\label{mini-project-information-profile-of-a-corpus}

Choose a small corpus such as course descriptions, news headlines,
abstracts, or student advising questions. Build an information-theoretic
profile of the corpus.

Your report should include:

\begin{enumerate}
\def\labelenumi{\arabic{enumi}.}
\tightlist
\item
  unigram token distribution and entropy;
\item
  entropy as a function of vocabulary truncation;
\item
  bigram conditional entropy;
\item
  a PMI or PPMI word-context table;
\item
  a comparison of at least two simple language models by cross-entropy
  and perplexity;
\item
  a discussion of how tokenization changes the numerical values;
\item
  a short reflection on what information theory captures and what it
  misses.
\end{enumerate}

\section{21.26 Python lab connection}\label{python-lab-connection-10}

The associated lab for this chapter is planned as:

\begin{Shaded}
\begin{Highlighting}[]
\NormalTok{labs/chapter{-}21{-}information{-}theory.ipynb}
\end{Highlighting}
\end{Shaded}

The lab should be written for independent study and should include:

\begin{itemize}
\tightlist
\item
  a short review of discrete probability distributions;
\item
  entropy and cross-entropy from scratch;
\item
  KL divergence and support mismatch examples;
\item
  perplexity calculations for held-out text;
\item
  bigram conditional entropy;
\item
  PMI and PPMI matrices;
\item
  a small corpus-analysis mini-project.
\end{itemize}

\section{21.27 Summary}\label{summary-5}

Information theory gives a common mathematical vocabulary for many parts
of NLP. Entropy measures uncertainty, cross-entropy measures predictive
loss, KL divergence measures distributional mismatch, and mutual
information measures statistical dependence. Perplexity is the
exponential form of average cross-entropy and is central to
language-model evaluation. PMI connects co-occurrence statistics with
word embeddings. KL regularization connects information theory with
fine-tuning and alignment.

The most important lesson is that information-theoretic quantities are
always defined relative to random variables and probability
distributions. Before computing entropy, perplexity, or mutual
information, first identify the tokenization, the sample space, the data
distribution, and the model distribution.

\section{Notes and references}\label{notes-and-references-11}

Classical information theory begins with Shannon (1948) and is developed
systematically in Cover and Thomas (2006). For statistical NLP
connections, see Manning and Schuetze (1999) and Jurafsky and Martin
(2024). For neural language modeling and cross-entropy training, see
Bengio et al. (2003), Goodfellow et al. (2016b), and Vaswani et al.
(2017).

\chapter{Optimization for NLP}\label{optimization-for-nlp}

Gradients, stochasticity, and large-scale training

\hfill\break

Modern NLP models are trained by solving very large optimization
problems. The model may be a linear classifier, a recurrent network, a
Transformer, or a retrieval-augmented system with trainable components,
but the mathematical core is the same: choose parameters, define a loss,
estimate a gradient from data, and update the parameters. This chapter
develops optimization for NLP from a clean mathematical point of view.

\section{22.1 Learning goals}\label{learning-goals-21}

After reading this chapter, students should be able to:

\begin{itemize}
\tightlist
\item
  formulate NLP training as empirical risk minimization;
\item
  derive gradients for softmax cross-entropy losses;
\item
  distinguish full-batch gradients from stochastic mini-batch gradients;
\item
  explain how SGD, momentum, RMSProp, Adam, and AdamW update parameters;
\item
  interpret learning rate, batch size, gradient noise, clipping, and
  regularization mathematically;
\item
  understand why optimization in neural NLP is nonconvex but still
  useful in practice;
\item
  implement small optimization algorithms from scratch in Python;
\item
  use AI tools to audit an optimization setup without treating them as
  black boxes.
\end{itemize}

\begin{tcolorbox}[enhanced jigsaw, arc=.35mm, bottomrule=.15mm, bottomtitle=1mm, breakable, colback=white, colbacktitle=quarto-callout-note-color!10!white, colframe=quarto-callout-note-color-frame, coltitle=black, left=2mm, leftrule=.75mm, opacityback=0, opacitybacktitle=0.6, rightrule=.15mm, title=\textcolor{quarto-callout-note-color}{\faInfo}\hspace{0.5em}{Guiding question}, titlerule=0mm, toprule=.15mm, toptitle=1mm]

When an NLP model improves during training, what mathematical object is
changing?

The answer is not simply ``the model learns language.'' More precisely,
the parameter vector \(\theta\) changes so that a loss function, usually
an empirical estimate of an expected negative log-likelihood, becomes
smaller on the training distribution.

\end{tcolorbox}

\section{22.2 The optimization problem behind
NLP}\label{the-optimization-problem-behind-nlp}

Let \(z\) denote an NLP input. Depending on the task, \(z\) may be a
token, a document, a sentence pair, a prompt, or a sequence prefix. Let
\(y\) denote the target. A model with parameters \(\theta \in R^p\)
defines either a prediction function

\[
f_\theta(z)
\]

or a conditional probability distribution

\[
p_\theta(y \mid z).
\]

A loss function measures the cost of a prediction. In supervised text
classification, a common loss is

\[
\ell(\theta;z,y) = -\log p_\theta(y\mid z).
\]

Given training data

\[
D_n = \{(z_i,y_i)\}_{i=1}^n,
\]

the empirical risk is

\[
L_n(\theta)=\frac{1}{n}\sum_{i=1}^n \ell(\theta;z_i,y_i).
\]

Training means solving, approximately,

\[
\hat \theta \in \arg\min_{\theta \in R^p} L_n(\theta).
\]

The population risk is

\[
L(\theta)=E_{(Z,Y)\sim P}\left[\ell(\theta;Z,Y)\right],
\]

where \(P\) is the data-generating distribution. Since \(P\) is unknown,
the empirical risk \(L_n\) is used as an estimate.

\textbf{Optimization viewpoint.} NLP training is the numerical problem
of minimizing \(L_n(\theta)\), but statistical learning is about whether
small \(L_n(\theta)\) also implies small \(L(\theta)\).

\section{22.3 Token-level losses for language
models}\label{token-level-losses-for-language-models}

For autoregressive language modeling, a training sequence is

\[
w_{1:T}=(w_1,\ldots,w_T).
\]

The model defines next-token probabilities

\[
p_\theta(w_t\mid w_{<t}).
\]

The negative log-likelihood for one sequence is

\[
\ell(\theta;w_{1:T})
= -\sum_{t=1}^T \log p_\theta(w_t\mid w_{<t}).
\]

For a corpus with sequences \(w^{(1)},\ldots,w^{(n)}\), the average
token-level loss is

\[
L_n(\theta)
= -\frac{1}{N}\sum_{i=1}^n\sum_{t=1}^{T_i}
\log p_\theta\left(w_t^{(i)}\mid w_{<t}^{(i)}\right),
\]

where

\[
N = \sum_{i=1}^n T_i.
\]

The corresponding perplexity is

\[
\operatorname{PPL}(\theta)=\exp(L_n(\theta))
\]

when the loss uses natural logarithms.

This gives a direct connection between optimization and language
modeling:

\begin{quote}
Lower training loss means the model assigns higher probability to the
observed tokens under teacher forcing.
\end{quote}

\section{22.4 Softmax cross-entropy and its
gradient}\label{softmax-cross-entropy-and-its-gradient}

Many NLP objectives end with a softmax. Suppose a model produces logits

\[
z = (z_1,\ldots,z_K) \in R^K.
\]

The softmax probabilities are

\[
p_j = \frac{\exp(z_j)}{\sum_{k=1}^K \exp(z_k)}.
\]

For a one-hot target vector \(e_y\), the cross-entropy loss is

\[
\ell(z,y)=-\log p_y.
\]

Equivalently,

\[
\ell(z,y)= -z_y + \log\left(\sum_{k=1}^K \exp(z_k)\right).
\]

Differentiating gives the fundamental gradient identity

\[
\frac{\partial \ell}{\partial z_j}=p_j-1_{j=y}.
\]

In vector form,

\[
\nabla_z \ell = p-e_y.
\]

This identity appears throughout NLP:

\begin{itemize}
\tightlist
\item
  in text classifiers;
\item
  in masked language models;
\item
  in next-token prediction;
\item
  in instruction tuning;
\item
  in preference losses after reduction to logits.
\end{itemize}

\begin{tcolorbox}[enhanced jigsaw, arc=.35mm, bottomrule=.15mm, bottomtitle=1mm, breakable, colback=white, colbacktitle=quarto-callout-tip-color!10!white, colframe=quarto-callout-tip-color-frame, coltitle=black, left=2mm, leftrule=.75mm, opacityback=0, opacitybacktitle=0.6, rightrule=.15mm, title=\textcolor{quarto-callout-tip-color}{\faLightbulb}\hspace{0.5em}{Interpretation}, titlerule=0mm, toprule=.15mm, toptitle=1mm]

The gradient \(p-e_y\) says: decrease the logit of classes assigned too
much probability and increase the logit of the correct class. The size
of the correction depends on the model's current probability error.

\end{tcolorbox}

\section{22.5 From logits to parameter
gradients}\label{from-logits-to-parameter-gradients}

Let the logits depend on parameters \(\theta\) through a differentiable
map

\[
z=f_\theta(x).
\]

The chain rule gives

\[
\nabla_\theta \ell(\theta;x,y)
= J_\theta f_\theta(x)^T (p-e_y),
\]

where \(J_\theta f_\theta(x)\) is the Jacobian of the logits with
respect to the parameters.

For a linear softmax classifier,

\[
z = Wh+b,
\]

where \(h \in R^d\), \(W \in R^{K\times d}\), and \(b \in R^K\). Then

\[
\nabla_W \ell = (p-e_y)h^T,
\]

\[
\nabla_b \ell = p-e_y,
\]

and

\[
\nabla_h \ell = W^T(p-e_y).
\]

This small calculation is the local building block of backpropagation.

\section{22.6 Gradient descent}\label{gradient-descent}

Gradient descent updates parameters by moving in the direction of
steepest local decrease:

\[
\theta_{k+1}=\theta_k - \eta_k \nabla L_n(\theta_k),
\]

where \(\eta_k>0\) is the learning rate.

For a quadratic loss

\[
L(\theta)=\frac{1}{2}\theta^T A \theta - b^T \theta,
\]

with \(A\) symmetric positive definite, the gradient is

\[
\nabla L(\theta)=A \theta-b.
\]

Gradient descent converges when the learning rate is small enough. If
\(\lambda_{\max}(A)\) is the largest eigenvalue of \(A\), then a
constant learning rate satisfying

\[
0<\eta<\frac{2}{\lambda_{\max}(A)}
\]

ensures convergence for this quadratic problem.

This idealized result matters because neural-network losses are locally
approximated by quadratic functions through Taylor expansion:

\[
L(\theta+u)\approx L(\theta)+\nabla L(\theta)^T u + \frac{1}{2}u^T H(\theta)u.
\]

Large curvature directions require smaller steps.

\section{22.7 Stochastic gradients}\label{stochastic-gradients}

For large NLP corpora, the full gradient

\[
\nabla L_n(\theta)=\frac{1}{n}\sum_{i=1}^n \nabla \ell(\theta;z_i,y_i)
\]

is expensive. Mini-batch stochastic gradient descent uses a random
subset \(B_k\subset \{1,\ldots,n\}\):

\[
g_k=\frac{1}{|B_k|}\sum_{i\in B_k}\nabla \ell(\theta_k;z_i,y_i).
\]

If the mini-batch is sampled uniformly, then

\[
E[g_k\mid \theta_k]=\nabla L_n(\theta_k).
\]

Thus \(g_k\) is an unbiased estimator of the full gradient. Its variance
decreases as the batch size increases. Under idealized independent
sampling,

\[
\operatorname{Var}(g_k)\approx \frac{1}{B}\Sigma(\theta_k),
\]

where \(B=|B_k|\) and \(\Sigma(\theta_k)\) is the covariance matrix of
per-example gradients.

The SGD update is

\[
\theta_{k+1}=\theta_k-\eta_k g_k.
\]

\textbf{Statistical view of SGD.} SGD replaces an exact derivative by a
noisy but cheaper estimator. The learning rate controls how much of this
noise is converted into parameter motion.

\section{22.8 Learning rate as a scale
parameter}\label{learning-rate-as-a-scale-parameter}

The learning rate is not merely a software option. It determines the
scale of parameter movement:

\[
\|\theta_{k+1}-\theta_k\| = \eta_k \|g_k\|.
\]

If \(\eta_k\) is too small, learning is slow. If it is too large, the
iterates may oscillate or diverge. In neural NLP, useful learning-rate
schedules often include:

\begin{enumerate}
\def\labelenumi{\arabic{enumi}.}
\tightlist
\item
  \textbf{warmup}, where the learning rate increases gradually at the
  beginning;
\item
  \textbf{constant phase}, where the learning rate remains approximately
  stable;
\item
  \textbf{decay}, where the learning rate decreases to refine the
  solution.
\end{enumerate}

A simple inverse-square-root schedule is

\[
\eta_k = c\min\{k^{-1/2}, k w^{-3/2}\},
\]

where \(w\) is the warmup length. A cosine decay schedule is

\[
\eta_k = \eta_{\min}
+\frac{1}{2}(\eta_{\max}-\eta_{\min})
\left(1+\cos\left(\frac{\pi k}{K}\right)\right).
\]

The mathematical purpose of a schedule is to balance exploration and
stabilization.

\section{22.9 Momentum}\label{momentum}

Momentum keeps a running average of past gradients:

\[
v_{k+1}=\mu v_k + g_k,
\]

\[
\theta_{k+1}=\theta_k-\eta v_{k+1},
\]

where \(0\leq \mu <1\). When gradients point in a consistent direction,
momentum accelerates movement. When gradients oscillate, momentum can
average out some oscillation.

For a long sequence of gradients, the velocity is

\[
v_{k+1}=g_k+\mu g_{k-1}+\mu^2 g_{k-2}+\cdots.
\]

Thus momentum is an exponentially weighted history of gradients.

\section{22.10 Adaptive methods: AdaGrad, RMSProp, and
Adam}\label{adaptive-methods-adagrad-rmsprop-and-adam}

NLP models often have sparse or highly uneven gradients. Some parameters
update frequently; others update rarely. Adaptive methods rescale
coordinates based on historical gradient magnitudes.

\subsection{AdaGrad}\label{adagrad}

AdaGrad accumulates squared gradients coordinatewise:

\[
s_{k+1}=s_k+g_k\odot g_k,
\]

\[
\theta_{k+1}=\theta_k-\eta \frac{g_k}{\sqrt{s_{k+1}}+\epsilon}.
\]

Coordinates with large accumulated gradients receive smaller future
steps.

\subsection{RMSProp}\label{rmsprop}

RMSProp uses an exponential moving average instead of a cumulative sum:

\[
s_{k+1}=\rho s_k+(1-\rho)(g_k\odot g_k),
\]

\[
\theta_{k+1}=\theta_k-\eta \frac{g_k}{\sqrt{s_{k+1}}+\epsilon}.
\]

\subsection{Adam}\label{adam}

Adam combines momentum and coordinatewise scaling:

\[
m_{k+1}=\beta_1 m_k+(1-\beta_1)g_k,
\]

\[
v_{k+1}=\beta_2 v_k+(1-\beta_2)(g_k\odot g_k).
\]

Bias-corrected estimates are

\[
\hat m_{k+1}=\frac{m_{k+1}}{1-\beta_1^{k+1}},
\qquad
\hat v_{k+1}=\frac{v_{k+1}}{1-\beta_2^{k+1}}.
\]

The update is

\[
\theta_{k+1}=\theta_k-\eta\frac{\hat m_{k+1}}{\sqrt{\hat v_{k+1}}+\epsilon}.
\]

Adam is popular in neural NLP because it handles noisy,
high-dimensional, and unevenly scaled gradients well.

\section{22.11 Weight decay and AdamW}\label{weight-decay-and-adamw}

A common regularized objective is

\[
L_{\lambda}(\theta)=L_n(\theta)+\frac{\lambda}{2}\|\theta\|_2^2.
\]

For ordinary gradient descent, this gives

\[
\theta_{k+1}
=\theta_k-\eta \left(\nabla L_n(\theta_k)+\lambda \theta_k\right)
=(1-\eta\lambda)\theta_k-\eta\nabla L_n(\theta_k).
\]

The regularizer shrinks parameters toward zero.

For adaptive optimizers, directly adding \(\lambda \theta\) to the
gradient interacts with coordinatewise scaling. AdamW decouples weight
decay from the adaptive gradient step:

\[
\theta_{k+1}
=(1-\eta\lambda)\theta_k
-\eta\frac{\hat m_{k+1}}{\sqrt{\hat v_{k+1}}+\epsilon}.
\]

This is often a cleaner implementation of shrinkage for Transformer
models.

\section{22.12 Gradient clipping}\label{gradient-clipping-1}

RNNs and Transformers can produce very large gradients. Gradient
clipping limits the step size by replacing \(g\) with

\[
\tilde g = g\min\left\{1,\frac{c}{\|g\|}\right\}.
\]

If \(\|g\|\leq c\), then \(\tilde g=g\). If \(\|g\|>c\), then
\(\|\tilde g\|=c\).

Clipping is not the same as fixing the model. It is a numerical
stabilization method. It prevents one unusually large mini-batch
gradient from causing a destructive update.

\section{22.13 Optimization and
generalization}\label{optimization-and-generalization}

Optimization minimizes training loss, but the goal is usually good
performance on new data. A model with very low empirical risk can still
have high population risk if it overfits.

Important sources of implicit or explicit regularization include:

\begin{itemize}
\tightlist
\item
  weight decay;
\item
  dropout;
\item
  early stopping;
\item
  data augmentation;
\item
  label smoothing;
\item
  small-batch gradient noise;
\item
  learning-rate schedules;
\item
  architecture constraints;
\item
  parameter-efficient fine-tuning.
\end{itemize}

From a statistical point of view, the optimizer is part of the learning
algorithm. Two optimizers can reach different solutions with similar
training losses but different generalization behavior.

\section{22.14 Nonconvexity in neural
NLP}\label{nonconvexity-in-neural-nlp}

Linear regression with squared loss is convex. Logistic regression is
convex in its parameters. Neural NLP models are typically nonconvex
because they are compositions of many nonlinear maps:

\[
f_\theta = f_L\circ f_{L-1}\circ \cdots \circ f_1.
\]

The loss surface may contain many critical points where

\[
\nabla L_n(\theta)=0.
\]

These can be local minima, saddle points, or flat regions. In
high-dimensional neural networks, saddle points and nearly flat valleys
are especially important.

A second-order Taylor approximation near \(\theta\) is

\[
L_n(\theta+u)
\approx L_n(\theta)+\nabla L_n(\theta)^Tu+\frac{1}{2}u^T H(\theta)u.
\]

The Hessian eigenvalues describe local curvature. Large positive
eigenvalues indicate sharp directions. Small eigenvalues indicate flat
directions. Negative eigenvalues indicate directions of local decrease.

In modern NLP, the parameter dimension \(p\) is often much larger than
the number of human-interpretable modeling choices. Therefore
optimization analysis focuses less on finding the unique global
minimizer and more on stability, scaling, gradient quality, and
generalization.

\section{22.15 Numerical stability}\label{numerical-stability-1}

Optimization can fail because of numerical issues, not because the
mathematical model is conceptually wrong. Common stability tricks
include:

\subsection{Stable log-sum-exp}\label{stable-log-sum-exp}

Instead of computing

\[
\log\sum_j \exp(z_j),
\]

use

\[
m + \log\sum_j \exp(z_j-m),
\qquad m=\max_j z_j.
\]

This avoids overflow when logits are large.

\subsection{Loss masks}\label{loss-masks}

For padded sequences, the loss should ignore padding tokens. If
\(M_{i,t}\in\{0,1\}\) is a mask, then

\[
L(\theta)=
-\frac{1}{\sum_{i,t} M_{i,t}}
\sum_{i,t} M_{i,t}
\log p_\theta(w_{i,t}\mid w_{i,<t}).
\]

\subsection{Mixed precision}\label{mixed-precision}

Large models often use lower precision arithmetic to save memory and
increase speed. This introduces a tradeoff between numerical efficiency
and stability. Loss scaling and careful normalization help prevent
underflow and overflow.

\section{22.16 Interactive: optimization
dynamics}\label{interactive-optimization-dynamics}

The following interactive compares gradient descent, momentum, and Adam
on a two-dimensional quadratic loss

\[
L(x,y)=\frac{1}{2}(a x^2+b y^2).
\]

The parameter \(b/a\) controls the condition number. A large condition
number creates a narrow valley, where plain gradient descent may zigzag.

\section{22.17 Python example: softmax regression
gradient}\label{python-example-softmax-regression-gradient}

This example derives and checks the gradient for a small text
classifier. The model is

\[
z_i = W x_i + b,
\qquad
p_i=\operatorname{softmax}(z_i),
\]

with average cross-entropy loss.

\begin{Shaded}
\begin{Highlighting}[]
\ImportTok{import}\NormalTok{ numpy }\ImportTok{as}\NormalTok{ np}

\NormalTok{rng }\OperatorTok{=}\NormalTok{ np.random.default\_rng(}\DecValTok{7243}\NormalTok{)}

\NormalTok{X }\OperatorTok{=}\NormalTok{ np.array([}
\NormalTok{    [}\DecValTok{2}\NormalTok{, }\DecValTok{1}\NormalTok{, }\DecValTok{0}\NormalTok{, }\DecValTok{0}\NormalTok{],}
\NormalTok{    [}\DecValTok{1}\NormalTok{, }\DecValTok{2}\NormalTok{, }\DecValTok{0}\NormalTok{, }\DecValTok{0}\NormalTok{],}
\NormalTok{    [}\DecValTok{0}\NormalTok{, }\DecValTok{0}\NormalTok{, }\DecValTok{2}\NormalTok{, }\DecValTok{1}\NormalTok{],}
\NormalTok{    [}\DecValTok{0}\NormalTok{, }\DecValTok{0}\NormalTok{, }\DecValTok{1}\NormalTok{, }\DecValTok{2}\NormalTok{],}
\NormalTok{    [}\DecValTok{1}\NormalTok{, }\DecValTok{0}\NormalTok{, }\DecValTok{1}\NormalTok{, }\DecValTok{0}\NormalTok{],}
\NormalTok{    [}\DecValTok{0}\NormalTok{, }\DecValTok{1}\NormalTok{, }\DecValTok{0}\NormalTok{, }\DecValTok{1}\NormalTok{]}
\NormalTok{], dtype}\OperatorTok{=}\BuiltInTok{float}\NormalTok{)}
\NormalTok{y }\OperatorTok{=}\NormalTok{ np.array([}\DecValTok{0}\NormalTok{, }\DecValTok{0}\NormalTok{, }\DecValTok{1}\NormalTok{, }\DecValTok{1}\NormalTok{, }\DecValTok{0}\NormalTok{, }\DecValTok{1}\NormalTok{])}
\NormalTok{K }\OperatorTok{=} \DecValTok{2}
\NormalTok{n, d }\OperatorTok{=}\NormalTok{ X.shape}

\NormalTok{W }\OperatorTok{=}\NormalTok{ rng.normal(scale}\OperatorTok{=}\FloatTok{0.1}\NormalTok{, size}\OperatorTok{=}\NormalTok{(K, d))}
\NormalTok{b }\OperatorTok{=}\NormalTok{ np.zeros(K)}

\KeywordTok{def}\NormalTok{ softmax\_rows(Z):}
\NormalTok{    Z\_shift }\OperatorTok{=}\NormalTok{ Z }\OperatorTok{{-}}\NormalTok{ Z.}\BuiltInTok{max}\NormalTok{(axis}\OperatorTok{=}\DecValTok{1}\NormalTok{, keepdims}\OperatorTok{=}\VariableTok{True}\NormalTok{)}
\NormalTok{    expZ }\OperatorTok{=}\NormalTok{ np.exp(Z\_shift)}
    \ControlFlowTok{return}\NormalTok{ expZ }\OperatorTok{/}\NormalTok{ expZ.}\BuiltInTok{sum}\NormalTok{(axis}\OperatorTok{=}\DecValTok{1}\NormalTok{, keepdims}\OperatorTok{=}\VariableTok{True}\NormalTok{)}

\KeywordTok{def}\NormalTok{ loss\_and\_grad(W, b, X, y):}
\NormalTok{    n }\OperatorTok{=}\NormalTok{ X.shape[}\DecValTok{0}\NormalTok{]}
\NormalTok{    Z }\OperatorTok{=}\NormalTok{ X }\OperatorTok{@}\NormalTok{ W.T }\OperatorTok{+}\NormalTok{ b}
\NormalTok{    P }\OperatorTok{=}\NormalTok{ softmax\_rows(Z)}
\NormalTok{    loss }\OperatorTok{=} \OperatorTok{{-}}\NormalTok{np.mean(np.log(P[np.arange(n), y]))}
\NormalTok{    G }\OperatorTok{=}\NormalTok{ P.copy()}
\NormalTok{    G[np.arange(n), y] }\OperatorTok{{-}=} \FloatTok{1.0}
\NormalTok{    G }\OperatorTok{/=}\NormalTok{ n}
\NormalTok{    grad\_W }\OperatorTok{=}\NormalTok{ G.T }\OperatorTok{@}\NormalTok{ X}
\NormalTok{    grad\_b }\OperatorTok{=}\NormalTok{ G.}\BuiltInTok{sum}\NormalTok{(axis}\OperatorTok{=}\DecValTok{0}\NormalTok{)}
    \ControlFlowTok{return}\NormalTok{ loss, grad\_W, grad\_b}

\NormalTok{loss, grad\_W, grad\_b }\OperatorTok{=}\NormalTok{ loss\_and\_grad(W, b, X, y)}
\BuiltInTok{print}\NormalTok{(}\StringTok{"loss:"}\NormalTok{, }\BuiltInTok{round}\NormalTok{(loss, }\DecValTok{6}\NormalTok{))}
\BuiltInTok{print}\NormalTok{(}\StringTok{"grad\_W shape:"}\NormalTok{, grad\_W.shape)}
\BuiltInTok{print}\NormalTok{(}\StringTok{"grad\_b:"}\NormalTok{, np.}\BuiltInTok{round}\NormalTok{(grad\_b, }\DecValTok{6}\NormalTok{))}
\end{Highlighting}
\end{Shaded}

\begin{verbatim}
loss: 0.736831
grad_W shape: (2, 4)
grad_b: [-0.069829  0.069829]
\end{verbatim}

The shapes are important:

\[
X\in R^{n\times d},
\qquad
W\in R^{K\times d},
\qquad
\nabla_W L\in R^{K\times d}.
\]

\section{22.18 Python example: finite-difference gradient
check}\label{python-example-finite-difference-gradient-check}

A gradient check compares an analytic derivative with a numerical
approximation:

\[
\frac{\partial L}{\partial \theta_j}
\approx
\frac{L(\theta+\epsilon e_j)-L(\theta-\epsilon e_j)}{2\epsilon}.
\]

\begin{Shaded}
\begin{Highlighting}[]
\NormalTok{eps }\OperatorTok{=} \FloatTok{1e{-}5}
\NormalTok{j, k }\OperatorTok{=} \DecValTok{0}\NormalTok{, }\DecValTok{2}
\NormalTok{W\_plus }\OperatorTok{=}\NormalTok{ W.copy()}
\NormalTok{W\_minus }\OperatorTok{=}\NormalTok{ W.copy()}
\NormalTok{W\_plus[j, k] }\OperatorTok{+=}\NormalTok{ eps}
\NormalTok{W\_minus[j, k] }\OperatorTok{{-}=}\NormalTok{ eps}
\NormalTok{loss\_plus, \_, \_ }\OperatorTok{=}\NormalTok{ loss\_and\_grad(W\_plus, b, X, y)}
\NormalTok{loss\_minus, \_, \_ }\OperatorTok{=}\NormalTok{ loss\_and\_grad(W\_minus, b, X, y)}
\NormalTok{finite\_difference }\OperatorTok{=}\NormalTok{ (loss\_plus }\OperatorTok{{-}}\NormalTok{ loss\_minus) }\OperatorTok{/}\NormalTok{ (}\DecValTok{2} \OperatorTok{*}\NormalTok{ eps)}
\NormalTok{analytic }\OperatorTok{=}\NormalTok{ grad\_W[j, k]}

\BuiltInTok{print}\NormalTok{(}\StringTok{"finite difference:"}\NormalTok{, }\BuiltInTok{round}\NormalTok{(finite\_difference, }\DecValTok{8}\NormalTok{))}
\BuiltInTok{print}\NormalTok{(}\StringTok{"analytic gradient:"}\NormalTok{, }\BuiltInTok{round}\NormalTok{(analytic, }\DecValTok{8}\NormalTok{))}
\BuiltInTok{print}\NormalTok{(}\StringTok{"absolute error:"}\NormalTok{, }\BuiltInTok{abs}\NormalTok{(finite\_difference }\OperatorTok{{-}}\NormalTok{ analytic))}
\end{Highlighting}
\end{Shaded}

\begin{verbatim}
finite difference: 0.11785573
analytic gradient: 0.11785573
absolute error: 7.044365091246618e-13
\end{verbatim}

Gradient checking is useful when writing custom training code. It is too
expensive for normal training, but it is excellent for debugging.

\section{22.19 Python example: GD, momentum, and Adam on a
quadratic}\label{python-example-gd-momentum-and-adam-on-a-quadratic}

The next example mirrors the interactive figure.

\begin{Shaded}
\begin{Highlighting}[]
\ImportTok{import}\NormalTok{ numpy }\ImportTok{as}\NormalTok{ np}

\NormalTok{A }\OperatorTok{=}\NormalTok{ np.diag([}\FloatTok{1.0}\NormalTok{, }\FloatTok{20.0}\NormalTok{])}

\KeywordTok{def}\NormalTok{ quad\_loss(theta):}
    \ControlFlowTok{return} \FloatTok{0.5} \OperatorTok{*}\NormalTok{ theta }\OperatorTok{@}\NormalTok{ A }\OperatorTok{@}\NormalTok{ theta}

\KeywordTok{def}\NormalTok{ quad\_grad(theta):}
    \ControlFlowTok{return}\NormalTok{ A }\OperatorTok{@}\NormalTok{ theta}

\KeywordTok{def}\NormalTok{ optimize(method}\OperatorTok{=}\StringTok{"gd"}\NormalTok{, lr}\OperatorTok{=}\FloatTok{0.08}\NormalTok{, steps}\OperatorTok{=}\DecValTok{60}\NormalTok{, beta}\OperatorTok{=}\FloatTok{0.9}\NormalTok{):}
\NormalTok{    theta }\OperatorTok{=}\NormalTok{ np.array([}\FloatTok{2.8}\NormalTok{, }\FloatTok{2.2}\NormalTok{], dtype}\OperatorTok{=}\BuiltInTok{float}\NormalTok{)}
\NormalTok{    path }\OperatorTok{=}\NormalTok{ [theta.copy()]}
\NormalTok{    losses }\OperatorTok{=}\NormalTok{ [quad\_loss(theta)]}
\NormalTok{    m }\OperatorTok{=}\NormalTok{ np.zeros\_like(theta)}
\NormalTok{    v }\OperatorTok{=}\NormalTok{ np.zeros\_like(theta)}
\NormalTok{    beta1, beta2 }\OperatorTok{=} \FloatTok{0.9}\NormalTok{, }\FloatTok{0.99}
\NormalTok{    eps }\OperatorTok{=} \FloatTok{1e{-}8}
    \ControlFlowTok{for}\NormalTok{ t }\KeywordTok{in} \BuiltInTok{range}\NormalTok{(}\DecValTok{1}\NormalTok{, steps }\OperatorTok{+} \DecValTok{1}\NormalTok{):}
\NormalTok{        g }\OperatorTok{=}\NormalTok{ quad\_grad(theta)}
        \ControlFlowTok{if}\NormalTok{ method }\OperatorTok{==} \StringTok{"gd"}\NormalTok{:}
\NormalTok{            theta }\OperatorTok{=}\NormalTok{ theta }\OperatorTok{{-}}\NormalTok{ lr }\OperatorTok{*}\NormalTok{ g}
        \ControlFlowTok{elif}\NormalTok{ method }\OperatorTok{==} \StringTok{"momentum"}\NormalTok{:}
\NormalTok{            m }\OperatorTok{=}\NormalTok{ beta }\OperatorTok{*}\NormalTok{ m }\OperatorTok{+}\NormalTok{ g}
\NormalTok{            theta }\OperatorTok{=}\NormalTok{ theta }\OperatorTok{{-}}\NormalTok{ lr }\OperatorTok{*}\NormalTok{ m}
        \ControlFlowTok{elif}\NormalTok{ method }\OperatorTok{==} \StringTok{"adam"}\NormalTok{:}
\NormalTok{            m }\OperatorTok{=}\NormalTok{ beta1 }\OperatorTok{*}\NormalTok{ m }\OperatorTok{+}\NormalTok{ (}\DecValTok{1} \OperatorTok{{-}}\NormalTok{ beta1) }\OperatorTok{*}\NormalTok{ g}
\NormalTok{            v }\OperatorTok{=}\NormalTok{ beta2 }\OperatorTok{*}\NormalTok{ v }\OperatorTok{+}\NormalTok{ (}\DecValTok{1} \OperatorTok{{-}}\NormalTok{ beta2) }\OperatorTok{*}\NormalTok{ (g }\OperatorTok{*}\NormalTok{ g)}
\NormalTok{            m\_hat }\OperatorTok{=}\NormalTok{ m }\OperatorTok{/}\NormalTok{ (}\DecValTok{1} \OperatorTok{{-}}\NormalTok{ beta1 }\OperatorTok{**}\NormalTok{ t)}
\NormalTok{            v\_hat }\OperatorTok{=}\NormalTok{ v }\OperatorTok{/}\NormalTok{ (}\DecValTok{1} \OperatorTok{{-}}\NormalTok{ beta2 }\OperatorTok{**}\NormalTok{ t)}
\NormalTok{            theta }\OperatorTok{=}\NormalTok{ theta }\OperatorTok{{-}}\NormalTok{ lr }\OperatorTok{*}\NormalTok{ m\_hat }\OperatorTok{/}\NormalTok{ (np.sqrt(v\_hat) }\OperatorTok{+}\NormalTok{ eps)}
\NormalTok{        path.append(theta.copy())}
\NormalTok{        losses.append(quad\_loss(theta))}
    \ControlFlowTok{return}\NormalTok{ np.array(path), np.array(losses)}

\ControlFlowTok{for}\NormalTok{ method }\KeywordTok{in}\NormalTok{ [}\StringTok{"gd"}\NormalTok{, }\StringTok{"momentum"}\NormalTok{, }\StringTok{"adam"}\NormalTok{]:}
\NormalTok{    path, losses }\OperatorTok{=}\NormalTok{ optimize(method}\OperatorTok{=}\NormalTok{method)}
    \BuiltInTok{print}\NormalTok{(}\SpecialStringTok{f"}\SpecialCharTok{\{}\NormalTok{method}\SpecialCharTok{:8s\}}\SpecialStringTok{ final theta = }\SpecialCharTok{\{}\NormalTok{np}\SpecialCharTok{.}\BuiltInTok{round}\NormalTok{(path[}\OperatorTok{{-}}\DecValTok{1}\NormalTok{], }\DecValTok{4}\NormalTok{)}\SpecialCharTok{\}}\SpecialStringTok{, final loss = }\SpecialCharTok{\{}\NormalTok{losses[}\OperatorTok{{-}}\DecValTok{1}\NormalTok{]}\SpecialCharTok{:.6f\}}\SpecialStringTok{"}\NormalTok{)}
\end{Highlighting}
\end{Shaded}

\begin{verbatim}
gd       final theta = [0.0188 0.    ], final loss = 0.000177
momentum final theta = [-0.0146 -0.1004], final loss = 0.100870
adam     final theta = [-0.1226 -0.0936], final loss = 0.095036
\end{verbatim}

A quadratic example is simple, but it teaches three practical lessons:

\begin{enumerate}
\def\labelenumi{\arabic{enumi}.}
\tightlist
\item
  curvature affects stable step sizes;
\item
  momentum can reduce zigzagging;
\item
  adaptive scaling can help when coordinates have very different
  curvature.
\end{enumerate}

\section{22.20 Python example: mini-batch gradient
noise}\label{python-example-mini-batch-gradient-noise}

Mini-batch gradients are random estimates of the full gradient. The
variance decreases with batch size.

\begin{Shaded}
\begin{Highlighting}[]
\NormalTok{rng }\OperatorTok{=}\NormalTok{ np.random.default\_rng(}\DecValTok{2026}\NormalTok{)}
\NormalTok{N }\OperatorTok{=} \DecValTok{1000}
\NormalTok{true\_grad }\OperatorTok{=}\NormalTok{ np.array([}\FloatTok{1.0}\NormalTok{, }\OperatorTok{{-}}\FloatTok{0.5}\NormalTok{, }\FloatTok{0.25}\NormalTok{])}
\NormalTok{per\_example\_grads }\OperatorTok{=}\NormalTok{ true\_grad }\OperatorTok{+}\NormalTok{ rng.normal(scale}\OperatorTok{=}\NormalTok{[}\FloatTok{2.0}\NormalTok{, }\FloatTok{1.0}\NormalTok{, }\FloatTok{0.5}\NormalTok{], size}\OperatorTok{=}\NormalTok{(N, }\DecValTok{3}\NormalTok{))}
\NormalTok{full\_grad }\OperatorTok{=}\NormalTok{ per\_example\_grads.mean(axis}\OperatorTok{=}\DecValTok{0}\NormalTok{)}

\ControlFlowTok{for}\NormalTok{ B }\KeywordTok{in}\NormalTok{ [}\DecValTok{1}\NormalTok{, }\DecValTok{4}\NormalTok{, }\DecValTok{16}\NormalTok{, }\DecValTok{64}\NormalTok{, }\DecValTok{256}\NormalTok{]:}
\NormalTok{    estimates }\OperatorTok{=}\NormalTok{ []}
    \ControlFlowTok{for}\NormalTok{ \_ }\KeywordTok{in} \BuiltInTok{range}\NormalTok{(}\DecValTok{500}\NormalTok{):}
\NormalTok{        idx }\OperatorTok{=}\NormalTok{ rng.choice(N, size}\OperatorTok{=}\NormalTok{B, replace}\OperatorTok{=}\VariableTok{False}\NormalTok{)}
\NormalTok{        estimates.append(per\_example\_grads[idx].mean(axis}\OperatorTok{=}\DecValTok{0}\NormalTok{))}
\NormalTok{    estimates }\OperatorTok{=}\NormalTok{ np.array(estimates)}
\NormalTok{    avg\_error\_norm }\OperatorTok{=}\NormalTok{ np.linalg.norm(estimates }\OperatorTok{{-}}\NormalTok{ full\_grad, axis}\OperatorTok{=}\DecValTok{1}\NormalTok{).mean()}
    \BuiltInTok{print}\NormalTok{(}\SpecialStringTok{f"batch size }\SpecialCharTok{\{}\NormalTok{B}\SpecialCharTok{:3d\}}\SpecialStringTok{: mean gradient{-}estimation error norm = }\SpecialCharTok{\{}\NormalTok{avg\_error\_norm}\SpecialCharTok{:.4f\}}\SpecialStringTok{"}\NormalTok{)}
\end{Highlighting}
\end{Shaded}

\begin{verbatim}
batch size   1: mean gradient-estimation error norm = 2.0634
batch size   4: mean gradient-estimation error norm = 1.0484
batch size  16: mean gradient-estimation error norm = 0.5055
batch size  64: mean gradient-estimation error norm = 0.2485
batch size 256: mean gradient-estimation error norm = 0.1095
\end{verbatim}

Larger batches give more accurate gradient estimates, but they are more
expensive and may reduce the useful exploratory noise of SGD.

\section{22.21 Python example: gradient
clipping}\label{python-example-gradient-clipping}

\begin{Shaded}
\begin{Highlighting}[]
\KeywordTok{def}\NormalTok{ clip\_gradient(g, threshold):}
\NormalTok{    norm }\OperatorTok{=}\NormalTok{ np.linalg.norm(g)}
    \ControlFlowTok{if}\NormalTok{ norm }\OperatorTok{\textless{}=}\NormalTok{ threshold:}
        \ControlFlowTok{return}\NormalTok{ g.copy()}
    \ControlFlowTok{return}\NormalTok{ g }\OperatorTok{*}\NormalTok{ (threshold }\OperatorTok{/}\NormalTok{ norm)}

\NormalTok{gradients }\OperatorTok{=}\NormalTok{ [}
\NormalTok{    np.array([}\FloatTok{0.2}\NormalTok{, }\OperatorTok{{-}}\FloatTok{0.1}\NormalTok{, }\FloatTok{0.3}\NormalTok{]),}
\NormalTok{    np.array([}\FloatTok{5.0}\NormalTok{, }\OperatorTok{{-}}\FloatTok{4.0}\NormalTok{, }\FloatTok{3.0}\NormalTok{]),}
\NormalTok{    np.array([}\FloatTok{40.0}\NormalTok{, }\OperatorTok{{-}}\FloatTok{20.0}\NormalTok{, }\FloatTok{10.0}\NormalTok{])}
\NormalTok{]}
\NormalTok{threshold }\OperatorTok{=} \FloatTok{2.0}

\ControlFlowTok{for}\NormalTok{ g }\KeywordTok{in}\NormalTok{ gradients:}
\NormalTok{    clipped }\OperatorTok{=}\NormalTok{ clip\_gradient(g, threshold)}
    \BuiltInTok{print}\NormalTok{(}\StringTok{"original norm:"}\NormalTok{, }\BuiltInTok{round}\NormalTok{(np.linalg.norm(g), }\DecValTok{3}\NormalTok{),}
          \StringTok{" clipped norm:"}\NormalTok{, }\BuiltInTok{round}\NormalTok{(np.linalg.norm(clipped), }\DecValTok{3}\NormalTok{),}
          \StringTok{" clipped vector:"}\NormalTok{, np.}\BuiltInTok{round}\NormalTok{(clipped, }\DecValTok{3}\NormalTok{))}
\end{Highlighting}
\end{Shaded}

\begin{verbatim}
original norm: 0.374  clipped norm: 0.374  clipped vector: [ 0.2 -0.1  0.3]
original norm: 7.071  clipped norm: 2.0  clipped vector: [ 1.414 -1.131  0.849]
original norm: 45.826  clipped norm: 2.0  clipped vector: [ 1.746 -0.873  0.436]
\end{verbatim}

Clipping preserves direction but changes magnitude when the gradient is
too large.

\section{22.22 Python example: training a tiny text
classifier}\label{python-example-training-a-tiny-text-classifier}

The following code trains a softmax classifier using averaged
token-count features. This is not a state-of-the-art NLP model. Its
purpose is to show the complete optimization loop in a transparent way.

\begin{Shaded}
\begin{Highlighting}[]
\ImportTok{from}\NormalTok{ collections }\ImportTok{import}\NormalTok{ Counter}
\ImportTok{import}\NormalTok{ numpy }\ImportTok{as}\NormalTok{ np}

\NormalTok{train\_texts }\OperatorTok{=}\NormalTok{ [}
    \StringTok{"excellent clear helpful explanation"}\NormalTok{,}
    \StringTok{"useful clear lecture helpful"}\NormalTok{,}
    \StringTok{"confusing unclear poor explanation"}\NormalTok{,}
    \StringTok{"poor confusing lecture unclear"}\NormalTok{,}
    \StringTok{"clear useful examples"}\NormalTok{,}
    \StringTok{"unclear poor examples"}
\NormalTok{]}
\NormalTok{train\_y }\OperatorTok{=}\NormalTok{ np.array([}\DecValTok{1}\NormalTok{, }\DecValTok{1}\NormalTok{, }\DecValTok{0}\NormalTok{, }\DecValTok{0}\NormalTok{, }\DecValTok{1}\NormalTok{, }\DecValTok{0}\NormalTok{])}

\NormalTok{vocab }\OperatorTok{=} \BuiltInTok{sorted}\NormalTok{(}\BuiltInTok{set}\NormalTok{(word }\ControlFlowTok{for}\NormalTok{ text }\KeywordTok{in}\NormalTok{ train\_texts }\ControlFlowTok{for}\NormalTok{ word }\KeywordTok{in}\NormalTok{ text.split()))}
\NormalTok{word\_to\_id }\OperatorTok{=}\NormalTok{ \{w: i }\ControlFlowTok{for}\NormalTok{ i, w }\KeywordTok{in} \BuiltInTok{enumerate}\NormalTok{(vocab)\}}

\NormalTok{X }\OperatorTok{=}\NormalTok{ np.zeros((}\BuiltInTok{len}\NormalTok{(train\_texts), }\BuiltInTok{len}\NormalTok{(vocab)))}
\ControlFlowTok{for}\NormalTok{ i, text }\KeywordTok{in} \BuiltInTok{enumerate}\NormalTok{(train\_texts):}
\NormalTok{    counts }\OperatorTok{=}\NormalTok{ Counter(text.split())}
    \ControlFlowTok{for}\NormalTok{ word, count }\KeywordTok{in}\NormalTok{ counts.items():}
\NormalTok{        X[i, word\_to\_id[word]] }\OperatorTok{=}\NormalTok{ count}
\NormalTok{    X[i] }\OperatorTok{/=} \BuiltInTok{max}\NormalTok{(}\DecValTok{1}\NormalTok{, X[i].}\BuiltInTok{sum}\NormalTok{())}

\NormalTok{rng }\OperatorTok{=}\NormalTok{ np.random.default\_rng(}\DecValTok{11}\NormalTok{)}
\NormalTok{K, d }\OperatorTok{=} \DecValTok{2}\NormalTok{, X.shape[}\DecValTok{1}\NormalTok{]}
\NormalTok{W }\OperatorTok{=}\NormalTok{ rng.normal(scale}\OperatorTok{=}\FloatTok{0.01}\NormalTok{, size}\OperatorTok{=}\NormalTok{(K, d))}
\NormalTok{b }\OperatorTok{=}\NormalTok{ np.zeros(K)}

\NormalTok{lr }\OperatorTok{=} \FloatTok{0.8}
\NormalTok{loss\_history }\OperatorTok{=}\NormalTok{ []}
\ControlFlowTok{for}\NormalTok{ epoch }\KeywordTok{in} \BuiltInTok{range}\NormalTok{(}\DecValTok{100}\NormalTok{):}
\NormalTok{    loss, grad\_W, grad\_b }\OperatorTok{=}\NormalTok{ loss\_and\_grad(W, b, X, train\_y)}
\NormalTok{    W }\OperatorTok{{-}=}\NormalTok{ lr }\OperatorTok{*}\NormalTok{ grad\_W}
\NormalTok{    b }\OperatorTok{{-}=}\NormalTok{ lr }\OperatorTok{*}\NormalTok{ grad\_b}
\NormalTok{    loss\_history.append(loss)}

\NormalTok{P }\OperatorTok{=}\NormalTok{ softmax\_rows(X }\OperatorTok{@}\NormalTok{ W.T }\OperatorTok{+}\NormalTok{ b)}
\NormalTok{pred }\OperatorTok{=}\NormalTok{ P.argmax(axis}\OperatorTok{=}\DecValTok{1}\NormalTok{)}
\BuiltInTok{print}\NormalTok{(}\StringTok{"vocabulary:"}\NormalTok{, vocab)}
\BuiltInTok{print}\NormalTok{(}\StringTok{"final loss:"}\NormalTok{, }\BuiltInTok{round}\NormalTok{(loss\_history[}\OperatorTok{{-}}\DecValTok{1}\NormalTok{], }\DecValTok{4}\NormalTok{))}
\BuiltInTok{print}\NormalTok{(}\StringTok{"predictions:"}\NormalTok{, pred.tolist())}
\BuiltInTok{print}\NormalTok{(}\StringTok{"training accuracy:"}\NormalTok{, }\BuiltInTok{float}\NormalTok{(np.mean(pred }\OperatorTok{==}\NormalTok{ train\_y)))}

\NormalTok{weights\_for\_positive }\OperatorTok{=}\NormalTok{ W[}\DecValTok{1}\NormalTok{] }\OperatorTok{{-}}\NormalTok{ W[}\DecValTok{0}\NormalTok{]}
\ControlFlowTok{for}\NormalTok{ word, weight }\KeywordTok{in} \BuiltInTok{sorted}\NormalTok{(}\BuiltInTok{zip}\NormalTok{(vocab, weights\_for\_positive), key}\OperatorTok{=}\KeywordTok{lambda}\NormalTok{ x: }\OperatorTok{{-}}\NormalTok{x[}\DecValTok{1}\NormalTok{]):}
    \BuiltInTok{print}\NormalTok{(}\SpecialStringTok{f"}\SpecialCharTok{\{}\NormalTok{word}\SpecialCharTok{:12s\}}\SpecialStringTok{ }\SpecialCharTok{\{}\NormalTok{weight}\SpecialCharTok{: .3f\}}\SpecialStringTok{"}\NormalTok{)}
\end{Highlighting}
\end{Shaded}

\begin{verbatim}
vocabulary: ['clear', 'confusing', 'examples', 'excellent', 'explanation', 'helpful', 'lecture', 'poor', 'unclear', 'useful']
final loss: 0.0819
predictions: [1, 1, 0, 0, 1, 0]
training accuracy: 1.0
clear         4.177
useful        2.818
helpful       2.564
excellent     1.373
explanation   0.115
examples     -0.013
lecture      -0.053
confusing    -2.508
poor         -4.111
unclear      -4.118
\end{verbatim}

This example connects optimization with interpretability. The learned
weight difference \(W_1-W_0\) gives a linear score direction for the
positive class.

\section{22.23 AI-assisted learning
prompts}\label{ai-assisted-learning-prompts}

Use an AI assistant as a mathematical critic rather than a replacement
for derivation.

\begin{tcolorbox}[enhanced jigsaw, arc=.35mm, bottomrule=.15mm, bottomtitle=1mm, breakable, colback=white, colbacktitle=quarto-callout-tip-color!10!white, colframe=quarto-callout-tip-color-frame, coltitle=black, left=2mm, leftrule=.75mm, opacityback=0, opacitybacktitle=0.6, rightrule=.15mm, title=\textcolor{quarto-callout-tip-color}{\faLightbulb}\hspace{0.5em}{Prompt 1: audit an objective}, titlerule=0mm, toprule=.15mm, toptitle=1mm]

I am training a text classifier with softmax cross-entropy. Write the
empirical risk \(L_n(\theta)\) explicitly, identify the model
probabilities, and explain which terms depend on \(\theta\). Then check
whether my notation distinguishes logits, probabilities, and labels.

\end{tcolorbox}

\begin{tcolorbox}[enhanced jigsaw, arc=.35mm, bottomrule=.15mm, bottomtitle=1mm, breakable, colback=white, colbacktitle=quarto-callout-tip-color!10!white, colframe=quarto-callout-tip-color-frame, coltitle=black, left=2mm, leftrule=.75mm, opacityback=0, opacitybacktitle=0.6, rightrule=.15mm, title=\textcolor{quarto-callout-tip-color}{\faLightbulb}\hspace{0.5em}{Prompt 2: verify a gradient}, titlerule=0mm, toprule=.15mm, toptitle=1mm]

Here is my derivation of the softmax cross-entropy gradient. Check each
line and identify the first incorrect equality, if any. Do not skip
dimension checks.

\end{tcolorbox}

\begin{tcolorbox}[enhanced jigsaw, arc=.35mm, bottomrule=.15mm, bottomtitle=1mm, breakable, colback=white, colbacktitle=quarto-callout-tip-color!10!white, colframe=quarto-callout-tip-color-frame, coltitle=black, left=2mm, leftrule=.75mm, opacityback=0, opacitybacktitle=0.6, rightrule=.15mm, title=\textcolor{quarto-callout-tip-color}{\faLightbulb}\hspace{0.5em}{Prompt 3: diagnose optimization instability}, titlerule=0mm, toprule=.15mm, toptitle=1mm]

My training loss sometimes becomes \texttt{nan}. Give a mathematical
checklist involving log-sum-exp, learning rate, gradient norm, loss
masking, division by small numbers, and probability clipping. Separate
likely numerical causes from modeling causes.

\end{tcolorbox}

\begin{tcolorbox}[enhanced jigsaw, arc=.35mm, bottomrule=.15mm, bottomtitle=1mm, breakable, colback=white, colbacktitle=quarto-callout-warning-color!10!white, colframe=quarto-callout-warning-color-frame, coltitle=black, left=2mm, leftrule=.75mm, opacityback=0, opacitybacktitle=0.6, rightrule=.15mm, title=\textcolor{quarto-callout-warning-color}{\faExclamationTriangle}\hspace{0.5em}{AI caution}, titlerule=0mm, toprule=.15mm, toptitle=1mm]

An AI assistant may correctly name Adam or SGD while giving an
incomplete formula. Always ask for the update equations, the role of
each hyperparameter, and the dimensions of the quantities involved.

\end{tcolorbox}

\section{22.24 Common mistakes}\label{common-mistakes-11}

\begin{enumerate}
\def\labelenumi{\arabic{enumi}.}
\tightlist
\item
  \textbf{Confusing logits and probabilities.} Softmax converts logits
  into probabilities. Cross-entropy should usually be computed using a
  numerically stable log-softmax.
\item
  \textbf{Averaging over padding tokens.} Padding tokens should be
  masked out of token-level losses.
\item
  \textbf{Using a learning rate without checking gradient scale.} The
  product \(\eta\|g\|\) controls step size.
\item
  \textbf{Assuming smaller training loss means better model.} Training
  loss and test performance can move differently.
\item
  \textbf{Treating Adam as magic.} Adam rescales coordinates using
  gradient moments; it still requires learning-rate tuning and
  regularization.
\item
  \textbf{Ignoring random variation.} Mini-batch order, initialization,
  and hardware can change training trajectories.
\end{enumerate}

\section{22.25 Worked calculation}\label{worked-calculation-4}

Suppose a three-class model outputs probabilities

\[
p=(0.20,0.60,0.20)
\]

and the correct class is class \(2\). With one-hot target

\[
e_2=(0,1,0),
\]

the gradient with respect to logits is

\[
\nabla_z\ell=p-e_2=(0.20,-0.40,0.20).
\]

If the logit model is \(z=Wh+b\) with

\[
h=(1,2)^T,
\]

then

\[
\nabla_W\ell=(p-e_2)h^T
=
\begin{pmatrix}
0.20 \\
-0.40 \\
0.20
\end{pmatrix}
\begin{pmatrix}
1 & 2
\end{pmatrix}
=
\begin{pmatrix}
0.20 & 0.40 \\
-0.40 & -0.80 \\
0.20 & 0.40
\end{pmatrix}.
\]

Thus gradient descent subtracts a positive multiple of this matrix. The
row for the correct class increases in the direction of \(h\), while the
rows for incorrect classes decrease in that direction.

\section{22.26 Exercises}\label{exercises-21}

\begin{enumerate}
\def\labelenumi{\arabic{enumi}.}
\item
  \textbf{Softmax gradient.} Starting from \[
  \ell(z,y)=-z_y+\log\sum_k \exp(z_k),
  \] derive \(\partial \ell/\partial z_j=p_j-1_{j=y}\).
\item
  \textbf{Linear classifier gradient.} For \(z=Wh+b\), derive
  \(\nabla_W\ell\), \(\nabla_b\ell\), and \(\nabla_h\ell\).
\item
  \textbf{Learning-rate stability.} Consider
  \(L(\theta)=\frac{1}{2}\lambda \theta^2\). Show that gradient descent
  converges if \(0<\eta<2/\lambda\).
\item
  \textbf{Mini-batch variance.} Suppose per-example gradients have
  covariance \(\Sigma\). Under independent sampling, show that the
  mini-batch gradient covariance is approximately \(\Sigma/B\).
\item
  \textbf{Gradient clipping.} Prove that the clipped gradient \[
  \tilde g=g\min\{1,c/\|g\|\}
  \] has norm at most \(c\).
\item
  \textbf{Adam update.} Write one full Adam update for the scalar
  gradient sequence \(g_1=2\), \(g_2=1\), using \(\beta_1=0.9\),
  \(\beta_2=0.99\), and \(\epsilon=10^{-8}\).
\item
  \textbf{Loss masking.} Create a small padded batch of token labels and
  logits. Compute the masked average cross-entropy by hand.
\item
  \textbf{Optimization versus generalization.} Describe an experiment
  comparing two optimizers that achieve similar training loss but
  different validation loss.
\end{enumerate}

\section{22.27 Mini-project: optimizer comparison for text
classification}\label{mini-project-optimizer-comparison-for-text-classification}

Build a small text classifier using count vectors, TF-IDF vectors, or
learned embeddings. Compare at least three optimizers:

\begin{enumerate}
\def\labelenumi{\arabic{enumi}.}
\tightlist
\item
  SGD;
\item
  SGD with momentum;
\item
  Adam or AdamW.
\end{enumerate}

For each optimizer, report:

\begin{itemize}
\tightlist
\item
  training loss curve;
\item
  validation loss curve;
\item
  gradient norm curve;
\item
  final classification metrics;
\item
  sensitivity to learning rate;
\item
  an explanation of the observed behavior using the mathematics of this
  chapter.
\end{itemize}

The final report should separate three questions:

\begin{enumerate}
\def\labelenumi{\arabic{enumi}.}
\tightlist
\item
  Which optimizer minimized training loss fastest?
\item
  Which optimizer gave the best validation performance?
\item
  Which optimizer was most stable across random seeds?
\end{enumerate}

\section{22.28 Python lab connection}\label{python-lab-connection-11}

The associated lab for this chapter is planned as:

\begin{Shaded}
\begin{Highlighting}[]
\NormalTok{labs/chapter{-}22{-}optimization{-}for{-}nlp.ipynb}
\end{Highlighting}
\end{Shaded}

The lab should be designed for independent study and should include:

\begin{itemize}
\tightlist
\item
  a review of gradients and softmax cross-entropy;
\item
  implementation of SGD, momentum, RMSProp, Adam, and AdamW from
  scratch;
\item
  gradient checking;
\item
  gradient clipping;
\item
  training a small text classifier;
\item
  experiments with batch size and learning-rate schedules;
\item
  a short AI-assisted debugging activity.
\end{itemize}

\section{22.29 What to remember}\label{what-to-remember-7}

\begin{itemize}
\tightlist
\item
  Most NLP training objectives are empirical risks built from negative
  log-likelihood or cross-entropy.
\item
  The softmax cross-entropy gradient is \(p-e_y\).
\item
  Backpropagation is repeated use of the chain rule.
\item
  SGD uses noisy but inexpensive gradient estimates.
\item
  Momentum averages gradient directions over time.
\item
  Adam rescales coordinates using first and second gradient moments.
\item
  AdamW separates weight decay from adaptive gradient scaling.
\item
  Gradient clipping controls the norm of extreme updates.
\item
  Optimization performance and statistical generalization are related
  but not identical.
\item
  In mathematical NLP, every training run should specify the model, data
  distribution, loss, optimizer, schedule, and evaluation metric.
\end{itemize}

\section{22.30 Notes and references}\label{notes-and-references-12}

Optimization for NLP draws on numerical optimization, stochastic
approximation, statistical learning, and deep learning. For broader
background, see Boyd and Vandenberghe (2004) for convex optimization,
Nocedal and Wright (2006) for numerical optimization, Goodfellow et al.
(2016a) for deep learning, and Vaswani et al. (2017) for Transformer
training context. Adam was introduced in Kingma and Ba (2015a), and
decoupled weight decay in AdamW is associated with Loshchilov and Hutter
(2019a).

\chapter{Graphs and Language}\label{graphs-and-language-1}

Dependency, knowledge, and message passing

\hfill\break

Language is not only a sequence. A sentence has syntactic relations, a
document collection has citation and hyperlink structure, a knowledge
base has subject-relation-object triples, and a corpus has co-occurrence
patterns among words. Graphs provide a mathematical language for
representing these relational structures. This chapter develops
graph-based NLP from a clean mathematical point of view: nodes, edges,
matrices, Laplacians, walks, propagation, and message passing.

\section{23.1 Learning goals}\label{learning-goals-22}

After reading this chapter, students should be able to:

\begin{itemize}
\tightlist
\item
  represent linguistic objects as graphs;
\item
  convert graphs into adjacency, degree, incidence, and Laplacian
  matrices;
\item
  distinguish directed, undirected, weighted, bipartite, heterogeneous,
  and relational graphs;
\item
  construct word co-occurrence graphs, dependency graphs, document-term
  graphs, and knowledge graphs;
\item
  interpret graph walks, centrality, PageRank, and label propagation
  mathematically;
\item
  derive one layer of graph neural message passing;
\item
  explain graph convolution, graph attention, and relation-specific
  message passing;
\item
  implement small graph algorithms from scratch in Python;
\item
  use AI tools to audit graph construction choices rather than simply
  generate graph code.
\end{itemize}

\begin{tcolorbox}[enhanced jigsaw, arc=.35mm, bottomrule=.15mm, bottomtitle=1mm, breakable, colback=white, colbacktitle=quarto-callout-note-color!10!white, colframe=quarto-callout-note-color-frame, coltitle=black, left=2mm, leftrule=.75mm, opacityback=0, opacitybacktitle=0.6, rightrule=.15mm, title=\textcolor{quarto-callout-note-color}{\faInfo}\hspace{0.5em}{Guiding question}, titlerule=0mm, toprule=.15mm, toptitle=1mm]

When should an NLP problem be represented as a graph rather than as a
sequence or a bag of words?

A graph is useful when the relational structure is part of the signal.
The important mathematical decision is not ``use a graph neural
network'' but first: choose nodes, choose edges, choose weights, and
decide what invariances the representation should have.

\end{tcolorbox}

\section{23.2 From sequences to relational
objects}\label{from-sequences-to-relational-objects}

Earlier chapters represented a text as a sequence

\[
w_{1:T}=(w_1,\\ldots,w_T)
\]

over a vocabulary \(V\), or as a vector \(x\in R^d\). These
representations are powerful but not always natural for relational
information.

A graph representation starts with a set of nodes

\[
\mathcal V = \{1,2,\ldots,n\}
\]

and a set of edges

\[
\mathcal E \subseteq \mathcal V \times \mathcal V.
\]

In NLP, nodes may represent:

\begin{itemize}
\tightlist
\item
  tokens in a sentence;
\item
  word types in a vocabulary;
\item
  documents in a corpus;
\item
  entities in a knowledge base;
\item
  concepts in an ontology;
\item
  users, questions, courses, or resources in an advising system.
\end{itemize}

Edges may represent:

\begin{itemize}
\tightlist
\item
  adjacency in a sentence;
\item
  syntactic dependency;
\item
  word co-occurrence;
\item
  citation or hyperlink;
\item
  semantic relation;
\item
  membership of a word in a document;
\item
  retrieval similarity.
\end{itemize}

Thus a graph is not just a data structure. It is a modeling assumption
about which relations are relevant.

\textbf{Graph modeling decision.} A graph-based NLP model begins with a
map

\[
\text{text data} \longmapsto G=(\mathcal V,\mathcal E,W,X),
\]

where \(W\) stores edge weights and \(X\) stores node features. The
quality of the downstream model depends strongly on this construction.

\section{23.3 Basic graph definitions}\label{basic-graph-definitions}

An undirected weighted graph on \(n\) nodes is described by a symmetric
adjacency matrix

\[
A\in R^{n\times n},
\qquad
A_{ij}=A_{ji}\ge 0.
\]

Here \(A_{ij}>0\) means that nodes \(i\) and \(j\) are connected, and
\(A_{ij}\) gives the edge weight. For an unweighted graph,
\(A_{ij}\in \{0,1\}\).

A directed graph uses a possibly nonsymmetric adjacency matrix. Then
\(A_{ij}>0\) may mean an edge from node \(i\) to node \(j\). The
direction convention must be stated because different software packages
use different conventions.

The weighted degree of node \(i\) is

\[
d_i = \sum_{j=1}^n A_{ij}.
\]

The degree matrix is diagonal:

\[
D=\operatorname{diag}(d_1,\ldots,d_n).
\]

For undirected graphs, the combinatorial graph Laplacian is

\[
L=D-A.
\]

The normalized adjacency matrix is often written

\[
S=D^{-1/2} A D^{-1/2},
\]

and the normalized Laplacian is

\[
L_{\mathrm{norm}} = I-D^{-1/2}AD^{-1/2}=I-S.
\]

When isolated nodes occur, \(d_i=0\), and the inverse degree factors
must be handled carefully.

\section{23.4 Why the Laplacian
matters}\label{why-the-laplacian-matters}

For a scalar function \(f\in R^n\) on nodes, the quadratic form of the
Laplacian is

\[
f^T L f
= \frac{1}{2}\sum_{i,j=1}^n A_{ij}(f_i-f_j)^2.
\]

This identity is central. It says that \(f^T L f\) is small when
connected nodes have similar values. In NLP, \(f_i\) may be a label
score, a semantic score, or one coordinate of a node embedding.

\begin{tcolorbox}[enhanced jigsaw, arc=.35mm, bottomrule=.15mm, bottomtitle=1mm, breakable, colback=white, colbacktitle=quarto-callout-tip-color!10!white, colframe=quarto-callout-tip-color-frame, coltitle=black, left=2mm, leftrule=.75mm, opacityback=0, opacitybacktitle=0.6, rightrule=.15mm, title=\textcolor{quarto-callout-tip-color}{\faLightbulb}\hspace{0.5em}{Smoothness prior on graphs}, titlerule=0mm, toprule=.15mm, toptitle=1mm]

A graph Laplacian encodes the prior that connected nodes should have
similar predictions or representations. For example, if two documents
are very similar, a classifier may penalize assigning very different
label scores to them.

\end{tcolorbox}

The Laplacian also appears in spectral clustering. Eigenvectors
associated with small eigenvalues reveal low-frequency graph signals,
meaning functions that vary slowly across edges. This is the graph
analogue of using low-frequency Fourier modes to describe a smooth
signal.

\section{23.5 A sentence as a dependency
graph}\label{a-sentence-as-a-dependency-graph}

A dependency parse represents a sentence by directed grammatical
relations. For example, consider

\begin{Shaded}
\begin{Highlighting}[]
\NormalTok{language models predict tokens}
\end{Highlighting}
\end{Shaded}

A possible simplified dependency graph is:

\begin{Shaded}
\begin{Highlighting}[]
\NormalTok{predict {-}\textgreater{} models}
\NormalTok{models  {-}\textgreater{} language}
\NormalTok{predict {-}\textgreater{} tokens}
\end{Highlighting}
\end{Shaded}

The words are nodes. Directed edges represent syntactic relations. The
main verb \texttt{predict} is connected to subject-like and object-like
components.

Mathematically, for a sentence with \(T\) tokens, a dependency graph has
an adjacency matrix

\[
A\in \{0,1\}^{T\times T}.
\]

If \(A_{ij}=1\) means token \(i\) points to token \(j\), then row \(i\)
lists outgoing syntactic neighbors. A graph-based model may update a
token representation by aggregating information from its syntactic
neighbors rather than from all tokens.

Dependency graphs are useful because they remove some irrelevant
word-distance effects. In a long sentence, two syntactically related
words may be far apart in the sequence but close in the dependency
graph.

\section{23.6 Word co-occurrence
graphs}\label{word-co-occurrence-graphs}

A word co-occurrence graph has one node per word type. Two words are
connected if they occur near each other in a corpus. A simple
construction uses a window size \(r\). For every occurrence of word
\(w_t\), connect it to words \(w_s\) satisfying

\[
0 < |s-t| \le r.
\]

The edge weight may count co-occurrence frequency:

\[
A_{ij}=\#\{(s,t): w_s=i, w_t=j, 0<|s-t|\le r\}.
\]

After symmetrization, \(A\) becomes an undirected weighted graph. This
graph captures the distributional hypothesis: words appearing in similar
neighborhoods tend to have related meanings.

This connects graph methods to Chapter 8. Co-occurrence matrices, PMI
matrices, and word graphs are different views of the same empirical
object.

\section{23.7 Document-term graphs as bipartite
graphs}\label{document-term-graphs-as-bipartite-graphs}

A document-term matrix

\[
X\in R^{m\times d}
\]

can be interpreted as a bipartite graph. One part contains document
nodes

\[
\{D_1,\ldots,D_m\},
\]

and the other part contains term nodes

\[
\{t_1,\ldots,t_d\}.
\]

There is an edge between document \(D_i\) and term \(t_j\) with weight
\(X_{ij}\). The full adjacency matrix has block form

\[
A_{\\text{bip}}=
\begin{bmatrix}
0 & X \\
X^T & 0
\end{bmatrix}.
\]

This construction lets information propagate between documents through
shared terms and between terms through shared documents. For example,
two documents need not share a direct edge; they may become related
through paths of length two:

\[
D_i \to t_j \to D_k.
\]

The document-document similarity matrix is

\[
XX^T,
\]

and the term-term similarity matrix is

\[
X^T X.
\]

Thus the familiar Gram matrices from vector-space NLP are also two-step
walk matrices in a bipartite graph.

\section{23.8 Knowledge graphs}\label{knowledge-graphs}

A knowledge graph stores triples

\[
(h,r,t),
\]

where \(h\) is a head entity, \(r\) is a relation, and \(t\) is a tail
entity. For example:

\begin{Shaded}
\begin{Highlighting}[]
\NormalTok{(Transformer, is\_a, neural\_architecture)}
\NormalTok{(BERT, uses, self\_attention)}
\NormalTok{(NLP, uses, probability)}
\end{Highlighting}
\end{Shaded}

A knowledge graph is not merely a graph; it is usually a labeled
directed multigraph. Different relations should not always be treated as
the same type of edge.

For \(R\) relation types, one may define relation-specific adjacency
matrices

\[
A^{(1)},\ldots,A^{(R)}.
\]

The entry \(A^{(r)}_{ij}=1\) means that relation \(r\) connects entity
\(i\) to entity \(j\). This creates an adjacency tensor

\[
\mathcal A\in R^{R\times n\times n}.
\]

Relation-specific structure is important. The relation \texttt{is\_a} is
not the same as \texttt{part\_of}, and neither is the same as
\texttt{cites} or \texttt{teaches}.

\section{23.9 Random walks and
centrality}\label{random-walks-and-centrality}

Graphs support random walks. For an undirected weighted graph with
degree matrix \(D\), the row-stochastic transition matrix is

\[
P=D^{-1}A.
\]

If \(p_t\in R^n\) is a row vector distribution over nodes at time \(t\),
then

\[
p_{t+1}=p_t P.
\]

A stationary distribution satisfies

\[
\pi = \pi P.
\]

For connected non-bipartite undirected graphs, the stationary
distribution is proportional to degree:

\[
\pi_i = \frac{d_i}{\sum_j d_j}.
\]

Random walks appear in NLP through ranking, graph-based disambiguation,
label propagation, and representation learning. PageRank modifies the
walk by adding teleportation:

\[
p_{t+1}=\alpha p_t P + (1-\alpha)v,
\]

where \(v\) is a restart distribution and \(0<\alpha<1\).

\section{23.10 Label propagation}\label{label-propagation}

Suppose a graph has a small set of labeled nodes and many unlabeled
nodes. Let \(Y\in R^{n\times K}\) store label scores. A simple label
propagation iteration is

\[
F^{(k+1)} = \alpha S F^{(k)} + (1-\alpha)Y,
\]

where \(S\) is a normalized adjacency matrix. At convergence, if
\(0<\alpha<1\), the fixed point satisfies

\[
F^* = \alpha S F^* + (1-\alpha)Y,
\]

so

\[
F^* = (1-\alpha)(I-\alpha S)^{-1}Y,
\]

provided the inverse exists.

This formula shows that label propagation is not magic. It is repeated
averaging over a graph with a restart to the initial labels.

In NLP, this can be used to propagate sentiment labels through a word
graph, topic labels through a document graph, or resource categories
through a knowledge graph.

\section{23.11 Message passing as the common
abstraction}\label{message-passing-as-the-common-abstraction}

Many graph neural networks use message passing. Each node \(i\) has a
hidden state

\[
h_i^{(\ell)}\in R^{d_\ell}
\]

at layer \(\ell\). A generic message passing layer has the form

\[
m_i^{(\ell)} = \operatorname{AGG}^{(\ell)}\left(\{\psi^{(\ell)}(h_i^{(\ell)},h_j^{(\ell)},e_{ij}) : j\in N(i)\}\right),
\]

then

\[
h_i^{(\ell+1)} = \phi^{(\ell)}\left(h_i^{(\ell)},m_i^{(\ell)}\right).
\]

Here:

\begin{itemize}
\tightlist
\item
  \(N(i)\) is the neighborhood of node \(i\);
\item
  \(e_{ij}\) may store edge features;
\item
  \(\psi\) creates messages;
\item
  \(\operatorname{AGG}\) aggregates messages, often by sum, mean, or
  attention;
\item
  \(\phi\) updates the node state.
\end{itemize}

The aggregation should usually be permutation invariant with respect to
the ordering of neighbors. This is the graph analogue of a set function.

\section{23.12 A graph convolution
layer}\label{a-graph-convolution-layer}

A common simplified graph convolution layer is

\[
H^{(\ell+1)} = \sigma(\tilde S H^{(\ell)} W^{(\ell)}),
\]

where:

\begin{itemize}
\tightlist
\item
  \(H^{(\ell)}\in R^{n\times d_\ell}\) stores node features;
\item
  \(W^{(\ell)}\in R^{d_\ell\times d_{\ell+1}}\) is a trainable weight
  matrix;
\item
  \(\sigma\) is a nonlinear activation;
\item
  \(\tilde S\) is a normalized adjacency matrix, often built after
  adding self-loops.
\end{itemize}

With self-loops,

\[
\tilde A = A + I,
\qquad
\tilde D_{ii}=\sum_j \tilde A_{ij},
\]

and

\[
\tilde S = \tilde D^{-1/2}\tilde A\tilde D^{-1/2}.
\]

The layer has two operations:

\begin{enumerate}
\def\labelenumi{\arabic{enumi}.}
\tightlist
\item
  mix nearby node features by multiplying by \(\tilde S\);
\item
  transform feature coordinates by multiplying by \(W^{(\ell)}\).
\end{enumerate}

\textbf{Graph convolution as smoothing plus learning.} The matrix
\(\tilde S\) performs graph smoothing. The matrix \(W^{(\ell)}\) learns
feature combinations. Their product creates a trainable local averaging
operator on graph signals.

\section{23.13 Graph attention}\label{graph-attention}

Self-attention from Chapters 13 and 14 allows every token to attend to
every other token. Graph attention restricts attention to graph
neighborhoods. For node \(i\), compute a score for each neighbor
\(j\in N(i)\):

\[
e_{ij} = a(h_i,h_j).
\]

Then normalize over the neighborhood:

\[
\alpha_{ij} = \frac{\exp(e_{ij})}{\sum_{k\in N(i)}\exp(e_{ik})}.
\]

The update is

\[
h_i' = \sigma\left(\sum_{j\in N(i)} \alpha_{ij} W h_j\right).
\]

This is attention on a sparse graph. In NLP, it can be used on
dependency graphs, document graphs, entity graphs, and hybrid graphs
connecting documents, passages, entities, and concepts.

\section{23.14 Relational message
passing}\label{relational-message-passing}

Knowledge graphs require relation-specific updates. A simple relational
message passing layer is

\[
h_i' = \sigma\left(
W_0 h_i + \sum_{r=1}^R\sum_{j\in N_r(i)} \frac{1}{c_{i,r}} W_r h_j
\right),
\]

where \(N_r(i)\) is the set of neighbors connected to \(i\) by relation
\(r\), \(W_r\) is a relation-specific matrix, and \(c_{i,r}\) is a
normalization constant.

This is important because edge labels carry meaning. Treating all
relations as identical may lose the distinction between syntactic,
semantic, citation, and membership relations.

\section{23.15 Interactive: graph message
passing}\label{interactive-graph-message-passing}

The following interactive shows a small word graph. Node colors
represent a scalar graph signal. Each message-passing step replaces a
node's value by a convex combination of its current value and the
average value of its neighbors:

\[
f^{(k+1)} = (1-\alpha)f^{(0)} + \alpha P f^{(k)}.
\]

The parameter \(\alpha\) controls how strongly information propagates
through the graph.

\section{23.16 Python example: build a word co-occurrence
graph}\label{python-example-build-a-word-co-occurrence-graph}

The following example builds a weighted undirected word graph from a
tiny corpus. The code avoids external graph libraries so that the
matrices are visible.

\begin{Shaded}
\begin{Highlighting}[]
\ImportTok{import}\NormalTok{ numpy }\ImportTok{as}\NormalTok{ np}
\ImportTok{import}\NormalTok{ pandas }\ImportTok{as}\NormalTok{ pd}
\ImportTok{from}\NormalTok{ collections }\ImportTok{import}\NormalTok{ Counter}

\NormalTok{corpus }\OperatorTok{=}\NormalTok{ [}
    \StringTok{"graph models represent language structure"}\NormalTok{,}
    \StringTok{"language models use tokens and context"}\NormalTok{,}
    \StringTok{"probability models describe uncertainty"}\NormalTok{,}
    \StringTok{"graph neural networks pass messages"}\NormalTok{,}
    \StringTok{"tokens cooccur with context in language"}
\NormalTok{]}

\KeywordTok{def}\NormalTok{ tokenize(text):}
    \ControlFlowTok{return}\NormalTok{ text.lower().split()}

\NormalTok{tokenized }\OperatorTok{=}\NormalTok{ [tokenize(doc) }\ControlFlowTok{for}\NormalTok{ doc }\KeywordTok{in}\NormalTok{ corpus]}
\NormalTok{vocab }\OperatorTok{=} \BuiltInTok{sorted}\NormalTok{(}\BuiltInTok{set}\NormalTok{(tok }\ControlFlowTok{for}\NormalTok{ doc }\KeywordTok{in}\NormalTok{ tokenized }\ControlFlowTok{for}\NormalTok{ tok }\KeywordTok{in}\NormalTok{ doc))}
\NormalTok{index }\OperatorTok{=}\NormalTok{ \{w: i }\ControlFlowTok{for}\NormalTok{ i, w }\KeywordTok{in} \BuiltInTok{enumerate}\NormalTok{(vocab)\}}

\NormalTok{window }\OperatorTok{=} \DecValTok{2}
\NormalTok{A }\OperatorTok{=}\NormalTok{ np.zeros((}\BuiltInTok{len}\NormalTok{(vocab), }\BuiltInTok{len}\NormalTok{(vocab)), dtype}\OperatorTok{=}\BuiltInTok{float}\NormalTok{)}
\ControlFlowTok{for}\NormalTok{ doc }\KeywordTok{in}\NormalTok{ tokenized:}
    \ControlFlowTok{for}\NormalTok{ t, w }\KeywordTok{in} \BuiltInTok{enumerate}\NormalTok{(doc):}
\NormalTok{        i }\OperatorTok{=}\NormalTok{ index[w]}
        \ControlFlowTok{for}\NormalTok{ s }\KeywordTok{in} \BuiltInTok{range}\NormalTok{(}\BuiltInTok{max}\NormalTok{(}\DecValTok{0}\NormalTok{, t }\OperatorTok{{-}}\NormalTok{ window), }\BuiltInTok{min}\NormalTok{(}\BuiltInTok{len}\NormalTok{(doc), t }\OperatorTok{+}\NormalTok{ window }\OperatorTok{+} \DecValTok{1}\NormalTok{)):}
            \ControlFlowTok{if}\NormalTok{ s }\OperatorTok{==}\NormalTok{ t:}
                \ControlFlowTok{continue}
\NormalTok{            j }\OperatorTok{=}\NormalTok{ index[doc[s]]}
\NormalTok{            A[i, j] }\OperatorTok{+=} \DecValTok{1}

\CommentTok{\# Symmetrize the directed window counts.}
\NormalTok{A }\OperatorTok{=} \FloatTok{0.5} \OperatorTok{*}\NormalTok{ (A }\OperatorTok{+}\NormalTok{ A.T)}

\NormalTok{degrees }\OperatorTok{=}\NormalTok{ A.}\BuiltInTok{sum}\NormalTok{(axis}\OperatorTok{=}\DecValTok{1}\NormalTok{)}
\NormalTok{summary }\OperatorTok{=}\NormalTok{ pd.DataFrame(\{}\StringTok{"word"}\NormalTok{: vocab, }\StringTok{"degree"}\NormalTok{: degrees\}).sort\_values(}\StringTok{"degree"}\NormalTok{, ascending}\OperatorTok{=}\VariableTok{False}\NormalTok{)}
\BuiltInTok{print}\NormalTok{(summary.head(}\DecValTok{8}\NormalTok{).to\_string(index}\OperatorTok{=}\VariableTok{False}\NormalTok{))}
\end{Highlighting}
\end{Shaded}

\begin{verbatim}
     word  degree
   models     9.0
 language     7.0
   tokens     6.0
  context     6.0
 networks     4.0
      use     4.0
represent     4.0
     with     4.0
\end{verbatim}

The matrix \(A\) is a weighted adjacency matrix. The degree of a word is
the total co-occurrence weight attached to it.

\section{23.17 Python example: Laplacian and
smoothness}\label{python-example-laplacian-and-smoothness}

For an undirected graph, the Laplacian is \(L=D-A\). The quantity
\(f^TLf\) measures how much a signal varies across edges.

\begin{Shaded}
\begin{Highlighting}[]
\NormalTok{D }\OperatorTok{=}\NormalTok{ np.diag(A.}\BuiltInTok{sum}\NormalTok{(axis}\OperatorTok{=}\DecValTok{1}\NormalTok{))}
\NormalTok{L }\OperatorTok{=}\NormalTok{ D }\OperatorTok{{-}}\NormalTok{ A}

\NormalTok{f }\OperatorTok{=}\NormalTok{ np.zeros(}\BuiltInTok{len}\NormalTok{(vocab))}
\ControlFlowTok{for}\NormalTok{ word, value }\KeywordTok{in}\NormalTok{ \{}\StringTok{"graph"}\NormalTok{: }\FloatTok{1.0}\NormalTok{, }\StringTok{"models"}\NormalTok{: }\FloatTok{0.8}\NormalTok{, }\StringTok{"language"}\NormalTok{: }\FloatTok{0.5}\NormalTok{, }\StringTok{"uncertainty"}\NormalTok{: }\OperatorTok{{-}}\FloatTok{0.5}\NormalTok{\}.items():}
    \ControlFlowTok{if}\NormalTok{ word }\KeywordTok{in}\NormalTok{ index:}
\NormalTok{        f[index[word]] }\OperatorTok{=}\NormalTok{ value}

\NormalTok{energy\_matrix }\OperatorTok{=} \BuiltInTok{float}\NormalTok{(f.T }\OperatorTok{@}\NormalTok{ L }\OperatorTok{@}\NormalTok{ f)}
\NormalTok{energy\_sum }\OperatorTok{=} \FloatTok{0.0}
\ControlFlowTok{for}\NormalTok{ i }\KeywordTok{in} \BuiltInTok{range}\NormalTok{(}\BuiltInTok{len}\NormalTok{(vocab)):}
    \ControlFlowTok{for}\NormalTok{ j }\KeywordTok{in} \BuiltInTok{range}\NormalTok{(}\BuiltInTok{len}\NormalTok{(vocab)):}
\NormalTok{        energy\_sum }\OperatorTok{+=} \FloatTok{0.5} \OperatorTok{*}\NormalTok{ A[i, j] }\OperatorTok{*}\NormalTok{ (f[i] }\OperatorTok{{-}}\NormalTok{ f[j]) }\OperatorTok{**} \DecValTok{2}

\BuiltInTok{print}\NormalTok{(}\StringTok{"f\^{}T L f:"}\NormalTok{, }\BuiltInTok{round}\NormalTok{(energy\_matrix, }\DecValTok{6}\NormalTok{))}
\BuiltInTok{print}\NormalTok{(}\StringTok{"edge{-}sum formula:"}\NormalTok{, }\BuiltInTok{round}\NormalTok{(energy\_sum, }\DecValTok{6}\NormalTok{))}
\BuiltInTok{print}\NormalTok{(}\StringTok{"smallest Laplacian eigenvalues:"}\NormalTok{, np.}\BuiltInTok{round}\NormalTok{(np.linalg.eigvalsh(L)[:}\DecValTok{5}\NormalTok{], }\DecValTok{6}\NormalTok{))}
\end{Highlighting}
\end{Shaded}

\begin{verbatim}
f^T L f: 9.61
edge-sum formula: 9.61
smallest Laplacian eigenvalues: [-0.        0.245177  0.686777  1.208809  2.      ]
\end{verbatim}

The two values agree up to numerical precision. This verifies the
identity

\[
f^T L f = \frac{1}{2}\sum_{i,j} A_{ij}(f_i-f_j)^2.
\]

\section{23.18 Python example: random walk and
PageRank}\label{python-example-random-walk-and-pagerank}

We now convert the adjacency matrix to a row-stochastic transition
matrix.

\begin{Shaded}
\begin{Highlighting}[]
\KeywordTok{def}\NormalTok{ row\_stochastic(A):}
\NormalTok{    row\_sums }\OperatorTok{=}\NormalTok{ A.}\BuiltInTok{sum}\NormalTok{(axis}\OperatorTok{=}\DecValTok{1}\NormalTok{, keepdims}\OperatorTok{=}\VariableTok{True}\NormalTok{)}
\NormalTok{    P }\OperatorTok{=}\NormalTok{ np.zeros\_like(A, dtype}\OperatorTok{=}\BuiltInTok{float}\NormalTok{)}
\NormalTok{    mask }\OperatorTok{=}\NormalTok{ row\_sums[:, }\DecValTok{0}\NormalTok{] }\OperatorTok{\textgreater{}} \DecValTok{0}
\NormalTok{    P[mask] }\OperatorTok{=}\NormalTok{ A[mask] }\OperatorTok{/}\NormalTok{ row\_sums[mask]}
\NormalTok{    P[}\OperatorTok{\textasciitilde{}}\NormalTok{mask, }\OperatorTok{\textasciitilde{}}\NormalTok{mask] }\OperatorTok{=} \FloatTok{1.0}
    \ControlFlowTok{return}\NormalTok{ P}

\NormalTok{P }\OperatorTok{=}\NormalTok{ row\_stochastic(A)}

\NormalTok{alpha }\OperatorTok{=} \FloatTok{0.85}
\NormalTok{n }\OperatorTok{=}\NormalTok{ A.shape[}\DecValTok{0}\NormalTok{]}
\NormalTok{v }\OperatorTok{=}\NormalTok{ np.ones(n) }\OperatorTok{/}\NormalTok{ n}
\NormalTok{p }\OperatorTok{=}\NormalTok{ v.copy()}
\ControlFlowTok{for}\NormalTok{ \_ }\KeywordTok{in} \BuiltInTok{range}\NormalTok{(}\DecValTok{100}\NormalTok{):}
\NormalTok{    p }\OperatorTok{=}\NormalTok{ alpha }\OperatorTok{*}\NormalTok{ (p }\OperatorTok{@}\NormalTok{ P) }\OperatorTok{+}\NormalTok{ (}\DecValTok{1} \OperatorTok{{-}}\NormalTok{ alpha) }\OperatorTok{*}\NormalTok{ v}

\NormalTok{pagerank }\OperatorTok{=}\NormalTok{ pd.DataFrame(\{}\StringTok{"word"}\NormalTok{: vocab, }\StringTok{"pagerank"}\NormalTok{: p\}).sort\_values(}\StringTok{"pagerank"}\NormalTok{, ascending}\OperatorTok{=}\VariableTok{False}\NormalTok{)}
\BuiltInTok{print}\NormalTok{(pagerank.head(}\DecValTok{8}\NormalTok{).to\_string(index}\OperatorTok{=}\VariableTok{False}\NormalTok{))}
\BuiltInTok{print}\NormalTok{(}\StringTok{"sum of scores:"}\NormalTok{, }\BuiltInTok{round}\NormalTok{(p.}\BuiltInTok{sum}\NormalTok{(), }\DecValTok{6}\NormalTok{))}
\end{Highlighting}
\end{Shaded}

\begin{verbatim}
     word  pagerank
   models  0.112016
 language  0.084972
  context  0.072496
   tokens  0.072312
 networks  0.062037
    graph  0.056305
represent  0.053312
     with  0.050625
sum of scores: 1.0
\end{verbatim}

PageRank is a smoothed centrality measure. In a word graph, high
PageRank words are connected to many other important words, but this
should not be confused with semantic importance in every task.

\section{23.19 Python example: label
propagation}\label{python-example-label-propagation}

Suppose we label a few words as mathematical/technical and a few as
nontechnical. Label propagation spreads these scores across the graph.

\begin{Shaded}
\begin{Highlighting}[]
\NormalTok{Y }\OperatorTok{=}\NormalTok{ np.zeros(}\BuiltInTok{len}\NormalTok{(vocab))}
\ControlFlowTok{for}\NormalTok{ word, value }\KeywordTok{in}\NormalTok{ \{}
    \StringTok{"graph"}\NormalTok{: }\FloatTok{1.0}\NormalTok{,}
    \StringTok{"probability"}\NormalTok{: }\FloatTok{1.0}\NormalTok{,}
    \StringTok{"neural"}\NormalTok{: }\FloatTok{1.0}\NormalTok{,}
    \StringTok{"tokens"}\NormalTok{: }\FloatTok{0.4}\NormalTok{,}
    \StringTok{"uncertainty"}\NormalTok{: }\FloatTok{0.8}\NormalTok{,}
    \StringTok{"and"}\NormalTok{: }\OperatorTok{{-}}\FloatTok{0.5}\NormalTok{,}
    \StringTok{"with"}\NormalTok{: }\OperatorTok{{-}}\FloatTok{0.5}
\NormalTok{\}.items():}
    \ControlFlowTok{if}\NormalTok{ word }\KeywordTok{in}\NormalTok{ index:}
\NormalTok{        Y[index[word]] }\OperatorTok{=}\NormalTok{ value}

\NormalTok{alpha }\OperatorTok{=} \FloatTok{0.7}
\NormalTok{F }\OperatorTok{=}\NormalTok{ Y.copy()}
\ControlFlowTok{for}\NormalTok{ \_ }\KeywordTok{in} \BuiltInTok{range}\NormalTok{(}\DecValTok{20}\NormalTok{):}
\NormalTok{    F }\OperatorTok{=}\NormalTok{ alpha }\OperatorTok{*}\NormalTok{ (P }\OperatorTok{@}\NormalTok{ F) }\OperatorTok{+}\NormalTok{ (}\DecValTok{1} \OperatorTok{{-}}\NormalTok{ alpha) }\OperatorTok{*}\NormalTok{ Y}

\NormalTok{propagation }\OperatorTok{=}\NormalTok{ pd.DataFrame(\{}\StringTok{"word"}\NormalTok{: vocab, }\StringTok{"score"}\NormalTok{: F\}).sort\_values(}\StringTok{"score"}\NormalTok{, ascending}\OperatorTok{=}\VariableTok{False}\NormalTok{)}
\BuiltInTok{print}\NormalTok{(propagation.head(}\DecValTok{10}\NormalTok{).to\_string(index}\OperatorTok{=}\VariableTok{False}\NormalTok{))}
\end{Highlighting}
\end{Shaded}

\begin{verbatim}
       word    score
     neural 0.517939
      graph 0.480214
probability 0.426963
uncertainty 0.366963
   networks 0.239906
   describe 0.218831
       pass 0.213894
   messages 0.158832
     models 0.143921
  represent 0.128022
\end{verbatim}

The formula is simple, but the modeling choices matter. If the graph is
noisy, propagation spreads noise. If the graph is meaningful,
propagation spreads useful information.

\section{23.20 Python example: one graph convolution
layer}\label{python-example-one-graph-convolution-layer}

We implement one normalized graph convolution layer

\[
H' = \operatorname{ReLU}(\tilde S H W).
\]

\begin{Shaded}
\begin{Highlighting}[]
\NormalTok{rng }\OperatorTok{=}\NormalTok{ np.random.default\_rng(}\DecValTok{7243}\NormalTok{)}

\NormalTok{A\_tilde }\OperatorTok{=}\NormalTok{ A }\OperatorTok{+}\NormalTok{ np.eye(}\BuiltInTok{len}\NormalTok{(vocab))}
\NormalTok{d\_tilde }\OperatorTok{=}\NormalTok{ A\_tilde.}\BuiltInTok{sum}\NormalTok{(axis}\OperatorTok{=}\DecValTok{1}\NormalTok{)}
\NormalTok{D\_inv\_sqrt }\OperatorTok{=}\NormalTok{ np.diag(}\DecValTok{1} \OperatorTok{/}\NormalTok{ np.sqrt(np.maximum(d\_tilde, }\FloatTok{1e{-}12}\NormalTok{)))}
\NormalTok{S\_tilde }\OperatorTok{=}\NormalTok{ D\_inv\_sqrt }\OperatorTok{@}\NormalTok{ A\_tilde }\OperatorTok{@}\NormalTok{ D\_inv\_sqrt}

\CommentTok{\# Simple node features: degree, PageRank, and seeded label score.}
\NormalTok{H }\OperatorTok{=}\NormalTok{ np.column\_stack([}
\NormalTok{    degrees }\OperatorTok{/} \BuiltInTok{max}\NormalTok{(degrees.}\BuiltInTok{max}\NormalTok{(), }\DecValTok{1}\NormalTok{),}
\NormalTok{    p }\OperatorTok{/}\NormalTok{ p.}\BuiltInTok{max}\NormalTok{(),}
\NormalTok{    Y}
\NormalTok{])}

\NormalTok{W }\OperatorTok{=}\NormalTok{ rng.normal(scale}\OperatorTok{=}\FloatTok{0.5}\NormalTok{, size}\OperatorTok{=}\NormalTok{(H.shape[}\DecValTok{1}\NormalTok{], }\DecValTok{2}\NormalTok{))}
\NormalTok{Z }\OperatorTok{=}\NormalTok{ S\_tilde }\OperatorTok{@}\NormalTok{ H }\OperatorTok{@}\NormalTok{ W}
\NormalTok{H\_next }\OperatorTok{=}\NormalTok{ np.maximum(Z, }\FloatTok{0.0}\NormalTok{)}

\BuiltInTok{print}\NormalTok{(}\StringTok{"H shape:"}\NormalTok{, H.shape)}
\BuiltInTok{print}\NormalTok{(}\StringTok{"W shape:"}\NormalTok{, W.shape)}
\BuiltInTok{print}\NormalTok{(}\StringTok{"H\_next shape:"}\NormalTok{, H\_next.shape)}
\BuiltInTok{print}\NormalTok{(pd.DataFrame(H\_next, index}\OperatorTok{=}\NormalTok{vocab, columns}\OperatorTok{=}\NormalTok{[}\StringTok{"coord\_1"}\NormalTok{, }\StringTok{"coord\_2"}\NormalTok{]).head())}
\end{Highlighting}
\end{Shaded}

\begin{verbatim}
H shape: (19, 3)
W shape: (3, 2)
H_next shape: (19, 2)
          coord_1   coord_2
and           0.0  0.000000
context       0.0  0.000000
cooccur       0.0  0.000000
describe      0.0  0.429837
graph         0.0  0.341356
\end{verbatim}

This example is not meant to train a useful model. It is meant to make
the dimensions visible.

\section{23.21 Python example: relation-specific knowledge graph
matrices}\label{python-example-relation-specific-knowledge-graph-matrices}

For a small knowledge graph, we can build one adjacency matrix per
relation.

\begin{Shaded}
\begin{Highlighting}[]
\NormalTok{triples }\OperatorTok{=}\NormalTok{ [}
\NormalTok{    (}\StringTok{"BERT"}\NormalTok{, }\StringTok{"uses"}\NormalTok{, }\StringTok{"self\_attention"}\NormalTok{),}
\NormalTok{    (}\StringTok{"Transformer"}\NormalTok{, }\StringTok{"uses"}\NormalTok{, }\StringTok{"self\_attention"}\NormalTok{),}
\NormalTok{    (}\StringTok{"NLP"}\NormalTok{, }\StringTok{"uses"}\NormalTok{, }\StringTok{"probability"}\NormalTok{),}
\NormalTok{    (}\StringTok{"NLP"}\NormalTok{, }\StringTok{"uses"}\NormalTok{, }\StringTok{"linear\_algebra"}\NormalTok{),}
\NormalTok{    (}\StringTok{"BERT"}\NormalTok{, }\StringTok{"is\_a"}\NormalTok{, }\StringTok{"Transformer"}\NormalTok{),}
\NormalTok{    (}\StringTok{"Transformer"}\NormalTok{, }\StringTok{"is\_a"}\NormalTok{, }\StringTok{"neural\_model"}\NormalTok{),}
\NormalTok{    (}\StringTok{"GCN"}\NormalTok{, }\StringTok{"is\_a"}\NormalTok{, }\StringTok{"graph\_model"}\NormalTok{),}
\NormalTok{    (}\StringTok{"GAT"}\NormalTok{, }\StringTok{"is\_a"}\NormalTok{, }\StringTok{"graph\_model"}\NormalTok{)}
\NormalTok{]}

\NormalTok{entities }\OperatorTok{=} \BuiltInTok{sorted}\NormalTok{(}\BuiltInTok{set}\NormalTok{([h }\ControlFlowTok{for}\NormalTok{ h, r, t }\KeywordTok{in}\NormalTok{ triples] }\OperatorTok{+}\NormalTok{ [t }\ControlFlowTok{for}\NormalTok{ h, r, t }\KeywordTok{in}\NormalTok{ triples]))}
\NormalTok{relations }\OperatorTok{=} \BuiltInTok{sorted}\NormalTok{(}\BuiltInTok{set}\NormalTok{(r }\ControlFlowTok{for}\NormalTok{ h, r, t }\KeywordTok{in}\NormalTok{ triples))}
\NormalTok{ent\_index }\OperatorTok{=}\NormalTok{ \{e: i }\ControlFlowTok{for}\NormalTok{ i, e }\KeywordTok{in} \BuiltInTok{enumerate}\NormalTok{(entities)\}}
\NormalTok{rel\_index }\OperatorTok{=}\NormalTok{ \{r: i }\ControlFlowTok{for}\NormalTok{ i, r }\KeywordTok{in} \BuiltInTok{enumerate}\NormalTok{(relations)\}}

\NormalTok{A\_rel }\OperatorTok{=}\NormalTok{ np.zeros((}\BuiltInTok{len}\NormalTok{(relations), }\BuiltInTok{len}\NormalTok{(entities), }\BuiltInTok{len}\NormalTok{(entities)))}
\ControlFlowTok{for}\NormalTok{ h, r, t }\KeywordTok{in}\NormalTok{ triples:}
\NormalTok{    A\_rel[rel\_index[r], ent\_index[h], ent\_index[t]] }\OperatorTok{=} \DecValTok{1}

\BuiltInTok{print}\NormalTok{(}\StringTok{"relations:"}\NormalTok{, relations)}
\BuiltInTok{print}\NormalTok{(}\StringTok{"entities:"}\NormalTok{, entities)}
\BuiltInTok{print}\NormalTok{(}\StringTok{"adjacency tensor shape:"}\NormalTok{, A\_rel.shape)}
\ControlFlowTok{for}\NormalTok{ r }\KeywordTok{in}\NormalTok{ relations:}
    \BuiltInTok{print}\NormalTok{(}\StringTok{"}\CharTok{\textbackslash{}n}\StringTok{Relation:"}\NormalTok{, r)}
    \BuiltInTok{print}\NormalTok{(pd.DataFrame(A\_rel[rel\_index[r]], index}\OperatorTok{=}\NormalTok{entities, columns}\OperatorTok{=}\NormalTok{entities).astype(}\BuiltInTok{int}\NormalTok{))}
\end{Highlighting}
\end{Shaded}

\begin{verbatim}
relations: ['is_a', 'uses']
entities: ['BERT', 'GAT', 'GCN', 'NLP', 'Transformer', 'graph_model', 'linear_algebra', 'neural_model', 'probability', 'self_attention']
adjacency tensor shape: (2, 10, 10)

Relation: is_a
                BERT  GAT  GCN  NLP  Transformer  graph_model  linear_algebra   
BERT               0    0    0    0            1            0               0  \
GAT                0    0    0    0            0            1               0   
GCN                0    0    0    0            0            1               0   
NLP                0    0    0    0            0            0               0   
Transformer        0    0    0    0            0            0               0   
graph_model        0    0    0    0            0            0               0   
linear_algebra     0    0    0    0            0            0               0   
neural_model       0    0    0    0            0            0               0   
probability        0    0    0    0            0            0               0   
self_attention     0    0    0    0            0            0               0   

                neural_model  probability  self_attention  
BERT                       0            0               0  
GAT                        0            0               0  
GCN                        0            0               0  
NLP                        0            0               0  
Transformer                1            0               0  
graph_model                0            0               0  
linear_algebra             0            0               0  
neural_model               0            0               0  
probability                0            0               0  
self_attention             0            0               0  

Relation: uses
                BERT  GAT  GCN  NLP  Transformer  graph_model  linear_algebra   
BERT               0    0    0    0            0            0               0  \
GAT                0    0    0    0            0            0               0   
GCN                0    0    0    0            0            0               0   
NLP                0    0    0    0            0            0               1   
Transformer        0    0    0    0            0            0               0   
graph_model        0    0    0    0            0            0               0   
linear_algebra     0    0    0    0            0            0               0   
neural_model       0    0    0    0            0            0               0   
probability        0    0    0    0            0            0               0   
self_attention     0    0    0    0            0            0               0   

                neural_model  probability  self_attention  
BERT                       0            0               1  
GAT                        0            0               0  
GCN                        0            0               0  
NLP                        0            1               0  
Transformer                0            0               1  
graph_model                0            0               0  
linear_algebra             0            0               0  
neural_model               0            0               0  
probability                0            0               0  
self_attention             0            0               0  
\end{verbatim}

This is the matrix form behind relation-specific message passing.

\section{23.22 Python example: graph attention
weights}\label{python-example-graph-attention-weights}

The following code computes attention weights over graph neighbors for
one node. It mirrors the attention formulas from Chapters 13 and 14 but
restricts the normalization to neighbors.

\begin{Shaded}
\begin{Highlighting}[]
\CommentTok{\# Use H as node features from the previous example.}
\NormalTok{q }\OperatorTok{=}\NormalTok{ rng.normal(size}\OperatorTok{=}\NormalTok{H.shape[}\DecValTok{1}\NormalTok{])}
\NormalTok{k }\OperatorTok{=}\NormalTok{ rng.normal(size}\OperatorTok{=}\NormalTok{H.shape[}\DecValTok{1}\NormalTok{])}
\NormalTok{node }\OperatorTok{=}\NormalTok{ index[}\StringTok{"models"}\NormalTok{] }\ControlFlowTok{if} \StringTok{"models"} \KeywordTok{in}\NormalTok{ index }\ControlFlowTok{else} \DecValTok{0}
\NormalTok{neighbors }\OperatorTok{=}\NormalTok{ np.where(A[node] }\OperatorTok{\textgreater{}} \DecValTok{0}\NormalTok{)[}\DecValTok{0}\NormalTok{]}

\NormalTok{scores }\OperatorTok{=}\NormalTok{ []}
\ControlFlowTok{for}\NormalTok{ j }\KeywordTok{in}\NormalTok{ neighbors:}
\NormalTok{    score }\OperatorTok{=}\NormalTok{ (H[node] }\OperatorTok{@}\NormalTok{ q) }\OperatorTok{+}\NormalTok{ (H[j] }\OperatorTok{@}\NormalTok{ k)}
\NormalTok{    scores.append(score)}
\NormalTok{scores }\OperatorTok{=}\NormalTok{ np.array(scores)}
\NormalTok{weights }\OperatorTok{=}\NormalTok{ np.exp(scores }\OperatorTok{{-}}\NormalTok{ scores.}\BuiltInTok{max}\NormalTok{())}
\NormalTok{weights }\OperatorTok{=}\NormalTok{ weights }\OperatorTok{/}\NormalTok{ weights.}\BuiltInTok{sum}\NormalTok{()}

\NormalTok{attn\_table }\OperatorTok{=}\NormalTok{ pd.DataFrame(\{}
    \StringTok{"neighbor"}\NormalTok{: [vocab[j] }\ControlFlowTok{for}\NormalTok{ j }\KeywordTok{in}\NormalTok{ neighbors],}
    \StringTok{"score"}\NormalTok{: scores,}
    \StringTok{"attention\_weight"}\NormalTok{: weights}
\NormalTok{\}).sort\_values(}\StringTok{"attention\_weight"}\NormalTok{, ascending}\OperatorTok{=}\VariableTok{False}\NormalTok{)}
\BuiltInTok{print}\NormalTok{(}\StringTok{"central node:"}\NormalTok{, vocab[node])}
\BuiltInTok{print}\NormalTok{(attn\_table.to\_string(index}\OperatorTok{=}\VariableTok{False}\NormalTok{))}
\end{Highlighting}
\end{Shaded}

\begin{verbatim}
central node: models
   neighbor     score  attention_weight
   language -0.055853          0.187169
  represent -0.235786          0.156347
        use -0.254089          0.153511
   describe -0.282875          0.149156
     tokens -0.414035          0.130821
uncertainty -0.933685          0.077803
      graph -0.934197          0.077763
probability -1.076780          0.067430
\end{verbatim}

Graph attention is local self-attention on a specified graph. The graph
determines which attention scores are even allowed to exist.

\section{23.23 AI-assisted learning
activity}\label{ai-assisted-learning-activity}

Use an AI assistant as a graph-construction critic.

\begin{enumerate}
\def\labelenumi{\arabic{enumi}.}
\tightlist
\item
  Give the assistant a small NLP problem, such as course-resource
  recommendation, document search, sentiment lexicon expansion, or
  mathematical-concept mapping.
\item
  Ask for three possible graph constructions.
\item
  For each construction, require the assistant to state:

  \begin{itemize}
  \tightlist
  \item
    node types;
  \item
    edge types;
  \item
    edge weights;
  \item
    whether the graph is directed;
  \item
    what information is lost;
  \item
    what bias the graph introduces.
  \end{itemize}
\item
  Ask for the adjacency matrix form and the corresponding
  message-passing update.
\item
  Critique the answer manually.
\end{enumerate}

\begin{tcolorbox}[enhanced jigsaw, arc=.35mm, bottomrule=.15mm, bottomtitle=1mm, breakable, colback=white, colbacktitle=quarto-callout-warning-color!10!white, colframe=quarto-callout-warning-color-frame, coltitle=black, left=2mm, leftrule=.75mm, opacityback=0, opacitybacktitle=0.6, rightrule=.15mm, title=\textcolor{quarto-callout-warning-color}{\faExclamationTriangle}\hspace{0.5em}{AI caution}, titlerule=0mm, toprule=.15mm, toptitle=1mm]

An AI tool may suggest a graph because it sounds modern. Your job is to
ask whether the graph represents a real relation in the data and whether
the proposed edge weights estimate anything meaningful.

\end{tcolorbox}

\section{23.24 Common mistakes}\label{common-mistakes-12}

\begin{enumerate}
\def\labelenumi{\arabic{enumi}.}
\tightlist
\item
  \textbf{Confusing a graph with a drawing.} The graph is the pair
  \((\mathcal V,\mathcal E)\) and its weights, not the plotted layout.
\item
  \textbf{Forgetting direction.} A dependency relation and its reverse
  are not the same unless explicitly symmetrized.
\item
  \textbf{Using unnormalized adjacency blindly.} High-degree nodes can
  dominate message passing.
\item
  \textbf{Ignoring edge labels.} Knowledge graph relations often require
  relation-specific parameters.
\item
  \textbf{Over-smoothing.} Too many propagation layers can make node
  representations too similar.
\item
  \textbf{Assuming graph edges are true semantics.} Edges are
  constructed from data and assumptions; they may be noisy.
\end{enumerate}

\section{23.25 Worked calculation}\label{worked-calculation-5}

Consider the undirected path graph with three nodes:

\[
1-2-3.
\]

The adjacency matrix is

\[
A=\begin{bmatrix}
0 & 1 & 0\\
1 & 0 & 1\\
0 & 1 & 0
\end{bmatrix}.
\]

The degree matrix is

\[
D=\begin{bmatrix}
1 & 0 & 0\\
0 & 2 & 0\\
0 & 0 & 1
\end{bmatrix}.
\]

The Laplacian is

\[
L=D-A=
\begin{bmatrix}
1 & -1 & 0\\
-1 & 2 & -1\\
0 & -1 & 1
\end{bmatrix}.
\]

For a signal \(f=(1,0,-1)^T\),

\[
Lf = (1,0,-1)^T,
\]

so

\[
f^TLf = 2.
\]

Using the edge formula gives

\[
(1-0)^2 + (0-(-1))^2 = 2.
\]

Thus the Laplacian energy equals the sum of squared differences over
undirected edges.

\section{23.26 Exercises}\label{exercises-22}

\begin{enumerate}
\def\labelenumi{\arabic{enumi}.}
\tightlist
\item
  \textbf{Adjacency matrix.} Construct the adjacency matrix for a
  dependency graph with five tokens. State clearly whether \(A_{ij}=1\)
  means \(i\to j\) or \(j\to i\).
\item
  \textbf{Laplacian identity.} Prove that
  \(f^TLf = \frac{1}{2}\sum_{i,j}A_{ij}(f_i-f_j)^2\) for an undirected
  weighted graph.
\item
  \textbf{Bipartite text graph.} Given a \(3\times 4\) document-term
  matrix \(X\), write the block adjacency matrix of the document-term
  graph.
\item
  \textbf{Two-step walks.} Show that \(XX^T\) counts weighted two-step
  paths from documents to documents through terms.
\item
  \textbf{Random walk.} For a small graph, compute \(P=D^{-1}A\) and run
  two steps of a random walk by hand.
\item
  \textbf{Label propagation.} Derive the fixed-point formula for
  \(F^* = \alpha S F^* + (1-\alpha)Y\).
\item
  \textbf{Graph convolution.} Check the dimensions in
  \(H'=\sigma(\tilde S H W)\) when \(H\in R^{n\times d}\) and
  \(W\in R^{d\times m}\).
\item
  \textbf{Graph attention.} Explain how graph attention differs from
  Transformer self-attention.
\item
  \textbf{Knowledge graph.} Build relation-specific adjacency matrices
  for a set of five triples.
\item
  \textbf{Model critique.} Choose an NLP task and compare a sequence
  model, a bag-of-words model, and a graph model for that task.
\end{enumerate}

\section{23.27 Mini-project}\label{mini-project-5}

Build a small graph-based course-resource recommender.

Suggested node types:

\begin{itemize}
\tightlist
\item
  student question;
\item
  course;
\item
  program requirement;
\item
  webpage or policy document;
\item
  mathematical topic.
\end{itemize}

Suggested edge types:

\begin{itemize}
\tightlist
\item
  question mentions topic;
\item
  topic appears in course;
\item
  course satisfies requirement;
\item
  webpage explains policy;
\item
  course prerequisite relation.
\end{itemize}

Tasks:

\begin{enumerate}
\def\labelenumi{\arabic{enumi}.}
\tightlist
\item
  Define the graph formally.
\item
  Build a small example with at least 20 nodes.
\item
  Implement one random-walk or label-propagation method.
\item
  Compare it with TF-IDF retrieval.
\item
  Explain when graph structure helps and when it adds noise.
\end{enumerate}

\section{23.28 Python lab connection}\label{python-lab-connection-12}

The associated lab for this chapter is planned as:

\begin{Shaded}
\begin{Highlighting}[]
\NormalTok{labs/chapter{-}23{-}graphs{-}and{-}language.ipynb}
\end{Highlighting}
\end{Shaded}

The lab should be written for independent study and should include:

\begin{itemize}
\tightlist
\item
  graph construction from text;
\item
  adjacency, degree, and Laplacian matrices;
\item
  PageRank and label propagation;
\item
  graph convolution from scratch;
\item
  local graph attention;
\item
  a small knowledge graph exercise;
\item
  an AI-assisted graph-design audit.
\end{itemize}

\section{23.29 What to remember}\label{what-to-remember-8}

\begin{itemize}
\tightlist
\item
  A graph representation is a modeling decision about relations.
\item
  Matrices make graph structure computable: \(A\), \(D\), \(L\), \(P\),
  and \(S\).
\item
  The Laplacian measures variation of signals across edges.
\item
  Bipartite document-term graphs connect vector-space NLP with graph
  walks.
\item
  Message passing is local aggregation followed by transformation.
\item
  Graph attention is attention restricted by an edge set.
\item
  Knowledge graphs require relation-specific structure.
\item
  Graph methods are powerful when edges represent meaningful relations
  and risky when edges are noisy or arbitrary.
\end{itemize}

\section{23.30 Notes and references}\label{notes-and-references-13}

Classical graph-based thinking in NLP is connected to dependency parsing
and information retrieval in Manning and Schuetze (1999), Jurafsky and
Martin (2024), and Manning et al. (2008). Spectral graph ideas are
treated in Chung (1997). PageRank was introduced in Brin and Page
(1998). Graph convolution and graph neural message passing are
associated with Kipf and Welling (2017a) and Hamilton et al. (2017).
Graph attention is developed in Velickovic et al. (2018). Knowledge
graph embeddings are surveyed in Wang et al. (2017).

\chapter{Topology and Geometry of Language
Representations}\label{topology-and-geometry-of-language-representations}

Metric spaces, persistent homology, and representation diagnostics

\hfill\break

Modern NLP represents words, sentences, documents, and prompts as
high-dimensional vectors. Once language objects become vectors, a corpus
becomes a point cloud, a similarity graph, or a finite metric space.
Linear algebra explains directions, projections, and low-rank structure;
probability explains uncertainty and likelihood; topology asks a
complementary question: what qualitative shape is present in the
representation space? This chapter develops the topology and geometry of
language representations from a clean mathematical point of view, with
emphasis on finite metric spaces, neighborhood graphs, simplicial
complexes, persistent homology, and representation diagnostics.

\section{24.1 Learning goals}\label{learning-goals-23}

After reading this chapter, students should be able to:

\begin{itemize}
\tightlist
\item
  view a collection of language embeddings as a finite metric space;
\item
  construct neighborhood graphs and Vietoris--Rips complexes from
  pairwise distances;
\item
  compute and interpret Betti numbers in dimensions 0 and 1;
\item
  explain persistent homology as a multiscale summary of shape;
\item
  distinguish geometric summaries from topological summaries;
\item
  understand how clustering, loops, anisotropy, and intrinsic dimension
  appear in embedding spaces;
\item
  use persistent features carefully as diagnostics rather than as
  automatic semantic truths;
\item
  implement small topological summaries from scratch in Python;
\item
  use AI tools to critique representation geometry and topology
  responsibly.
\end{itemize}

\begin{tcolorbox}[enhanced jigsaw, arc=.35mm, bottomrule=.15mm, bottomtitle=1mm, breakable, colback=white, colbacktitle=quarto-callout-note-color!10!white, colframe=quarto-callout-note-color-frame, coltitle=black, left=2mm, leftrule=.75mm, opacityback=0, opacitybacktitle=0.6, rightrule=.15mm, title=\textcolor{quarto-callout-note-color}{\faInfo}\hspace{0.5em}{Guiding question}, titlerule=0mm, toprule=.15mm, toptitle=1mm]

What is the shape of a language representation space?

A linear-algebra answer might examine singular values, principal
components, and distances. A statistical answer might estimate class
separation or uncertainty. A topological answer asks whether the data
form connected components, branching structures, loops, voids, or other
multiscale patterns that persist across distance thresholds.

\end{tcolorbox}

\section{24.2 From embeddings to point
clouds}\label{from-embeddings-to-point-clouds}

Let \(T\) be a collection of language objects. Depending on the problem,
the objects may be:

\begin{itemize}
\tightlist
\item
  words in a vocabulary;
\item
  tokens in one sentence;
\item
  sentence embeddings from a corpus;
\item
  document embeddings;
\item
  contextual token embeddings from one layer of a Transformer;
\item
  prompts, questions, answers, or retrieved chunks.
\end{itemize}

An embedding model is a map

\[
f_\theta:T\to R^d.
\]

For objects \(t_1,\ldots,t_n\), define

\[
x_i=f_\theta(t_i)\in R^d,
\qquad
X=\{x_1,\ldots,x_n\}\subset R^d.
\]

Thus a finite collection of language objects becomes a finite point
cloud. The point cloud may be studied by Euclidean distances,

\[
d_E(x_i,x_j)=\|x_i-x_j\|_2,
\]

cosine distances,

\[
d_C(x_i,x_j)=1-\frac{x_i^T x_j}{\|x_i\|_2\|x_j\|_2},
\]

or other task-dependent dissimilarities.

\textbf{Finite metric-space viewpoint.} For topological data analysis,
the central object is often not the coordinate matrix
\(X\in R^{n\times d}\) itself, but the pairwise distance matrix

\[
D=(d_{ij}),
\qquad
d_{ij}=d(x_i,x_j).
\]

Many topological constructions use only \(D\), so they can be applied to
Euclidean embeddings, cosine embeddings, edit-distance spaces, or
learned similarity spaces.

\section{24.3 Geometry versus topology}\label{geometry-versus-topology}

Geometry studies quantities such as lengths, angles, curvature,
projections, and distances. Topology studies qualitative features that
are stable under continuous deformation, such as connected components
and holes.

For language embeddings, geometry asks questions such as:

\begin{itemize}
\tightlist
\item
  Which sentence is nearest to a query?
\item
  Which direction separates two topics?
\item
  How much variance is explained by the first few principal components?
\item
  Are embeddings anisotropic?
\end{itemize}

Topology asks questions such as:

\begin{itemize}
\tightlist
\item
  How many connected clusters appear at a given scale?
\item
  Do clusters merge gradually or suddenly?
\item
  Is there evidence of loop-like structure among semantic groups?
\item
  Does the same qualitative shape persist across several scales?
\end{itemize}

A simple example helps. Suppose sentence embeddings lie approximately on
a circle because the sentences describe a cyclic process such as
morning, afternoon, evening, night, and morning again. PCA may show a
two-dimensional plane. Clustering may find local groups. Persistent
homology can detect a one-dimensional loop.

\section{24.4 Neighborhood graphs}\label{neighborhood-graphs}

A first topological construction is the \(\varepsilon\)-neighborhood
graph. Given a finite metric space \((X,d)\) and a scale parameter
\(\varepsilon\ge 0\), define a graph

\[
G_\varepsilon=(V,E_\varepsilon),
\qquad
V=\{1,\ldots,n\},
\]

where

\[
(i,j)\in E_\varepsilon
\quad\Longleftrightarrow\quad
d(x_i,x_j)\le \varepsilon.
\]

As \(\varepsilon\) increases, the graph changes:

\begin{itemize}
\tightlist
\item
  for very small \(\varepsilon\), every point is isolated;
\item
  at moderate \(\varepsilon\), nearby points form local clusters;
\item
  at larger \(\varepsilon\), clusters merge;
\item
  for sufficiently large \(\varepsilon\), the graph becomes connected.
\end{itemize}

The number of connected components is the zeroth Betti number:

\[
\beta_0(G_\varepsilon)=\text{number of connected components of }G_\varepsilon.
\]

For a graph, the number of independent cycles is the graph cycle rank:

\[
\beta_1^{\mathrm{graph}}(G_\varepsilon)=|E_\varepsilon|-|V|+\beta_0(G_\varepsilon).
\]

This graph quantity counts cycles before considering whether triangular
cycles should be filled in. Persistent homology refines this idea by
using simplicial complexes.

\section{24.5 Simplicial complexes}\label{simplicial-complexes}

A simplicial complex is a collection of finite subsets closed under
taking subsets. Its elements are called simplices.

\begin{itemize}
\tightlist
\item
  A 0-simplex is a vertex.
\item
  A 1-simplex is an edge.
\item
  A 2-simplex is a filled triangle.
\item
  A 3-simplex is a filled tetrahedron.
\end{itemize}

If \(\sigma\) is a simplex and \(\tau\subseteq \sigma\), then \(\tau\)
must also be a simplex. This closure condition is important: if a filled
triangle is present, then its three edges and three vertices must be
present.

For a finite metric space, the most common construction is the
Vietoris--Rips complex.

\textbf{Vietoris--Rips complex.} Given \((X,d)\) and
\(\varepsilon\ge 0\), the Vietoris--Rips complex
\(\operatorname{VR}_\varepsilon(X)\) contains a simplex
\(\{i_0,\ldots,i_k\}\) whenever all pairwise distances are at most
\(\varepsilon\):

\[
d(x_{i_a},x_{i_b})\le \varepsilon
\quad
\text{for all }a,b.
\]

Equivalently, \(\operatorname{VR}_\varepsilon(X)\) is the clique complex
of the \(\varepsilon\)-neighborhood graph.

This construction is useful for language embeddings because it converts
similarity information into a multiscale combinatorial object.

\section{24.6 Betti numbers}\label{betti-numbers}

Betti numbers count independent topological features. Informally:

\[
\beta_0 = \text{number of connected components},
\]

\[
\beta_1 = \text{number of independent one-dimensional loops},
\]

\[
\beta_2 = \text{number of independent two-dimensional voids}.
\]

For a graph without filled triangles, the formula

\[
\beta_1=|E|-|V|+\beta_0
\]

counts independent graph cycles. For a simplicial complex, some graph
cycles become boundaries of filled triangles and therefore should not
count as holes.

For a 2-dimensional simplicial complex with vertices, edges, and
triangles, one can compute

\[
\beta_1 = \dim \ker \partial_1 - \operatorname{rank}\partial_2,
\]

where

\[
\partial_1:C_1\to C_0
\]

maps oriented edges to differences of endpoints, and

\[
\partial_2:C_2\to C_1
\]

maps oriented triangles to signed sums of boundary edges.

Since

\[
\dim\ker \partial_1 = |E| - \operatorname{rank}\partial_1,
\]

and for a graph

\[
\operatorname{rank}\partial_1 = |V|-\beta_0,
\]

we get

\[
\beta_1
=
|E|-|V|+\beta_0-\operatorname{rank}\partial_2.
\]

This formula shows why filled triangles matter: the rank of
\(\partial_2\) subtracts cycles that have become triangle boundaries.

\section{24.7 Filtrations and persistent
homology}\label{filtrations-and-persistent-homology}

A single value of \(\varepsilon\) is arbitrary. Persistent homology
studies all scales at once. Since

\[
\varepsilon_1\le \varepsilon_2
\quad\Longrightarrow\quad
\operatorname{VR}_{\varepsilon_1}(X)\subseteq \operatorname{VR}_{\varepsilon_2}(X),
\]

we obtain a nested sequence

\[
\operatorname{VR}_{\varepsilon_1}(X)
\subseteq
\operatorname{VR}_{\varepsilon_2}(X)
\subseteq
\cdots
\subseteq
\operatorname{VR}_{\varepsilon_m}(X).
\]

This is called a filtration.

As \(\varepsilon\) increases, features appear and disappear. A connected
component is born when a point appears and dies when it merges with an
older component. A loop is born when edges close a cycle and dies when
the loop becomes filled by triangles.

A persistent homology summary records intervals

\[
[b_i,d_i),
\]

where \(b_i\) is the birth scale and \(d_i\) is the death scale of a
feature. The persistence length is

\[
d_i-b_i.
\]

Longer intervals are usually interpreted as more stable multiscale
features. Short intervals are often interpreted as local noise, although
this interpretation depends on the problem.

\begin{tcolorbox}[enhanced jigsaw, arc=.35mm, bottomrule=.15mm, bottomtitle=1mm, breakable, colback=white, colbacktitle=quarto-callout-warning-color!10!white, colframe=quarto-callout-warning-color-frame, coltitle=black, left=2mm, leftrule=.75mm, opacityback=0, opacitybacktitle=0.6, rightrule=.15mm, title=\textcolor{quarto-callout-warning-color}{\faExclamationTriangle}\hspace{0.5em}{Interpretation warning}, titlerule=0mm, toprule=.15mm, toptitle=1mm]

A persistent loop in an embedding space is not automatically a
linguistic fact. It may come from semantic structure, sampling design,
normalization, model architecture, dimensionality reduction, or
numerical artifacts. Topology gives diagnostics; interpretation still
requires domain reasoning.

\end{tcolorbox}

\section{24.8 Why topology can help
NLP}\label{why-topology-can-help-nlp}

Topological summaries are useful because they are less tied to a
particular coordinate system than raw embedding coordinates. If two
embedding spaces are transformed by rotations or small metric
perturbations, their large-scale topological summaries may remain
similar.

In NLP, topology can help analyze:

\begin{itemize}
\tightlist
\item
  whether sentence embeddings form topic clusters;
\item
  whether generated texts cover the same representation shape as human
  texts;
\item
  whether contextual embeddings collapse into a narrow cone;
\item
  whether different layers of a Transformer gradually change
  representation geometry;
\item
  whether retrieval chunks form disconnected semantic islands;
\item
  whether class labels align with connected components or cut across
  them.
\end{itemize}

This is especially useful for mathematical NLP because it gives students
a way to ask structural questions beyond accuracy.

\section{24.9 Local dimension and
anisotropy}\label{local-dimension-and-anisotropy}

Topology should be combined with geometry. Two important geometric
diagnostics are local intrinsic dimension and anisotropy.

Let \(x_i\) be a point and let \(N_k(i)\) be its \(k\) nearest
neighbors. Let

\[
Y_i = \{x_j-x_i:j\in N_k(i)\}.
\]

The local covariance matrix is

\[
C_i = \frac{1}{k}\sum_{j\in N_k(i)} (x_j-x_i)(x_j-x_i)^T.
\]

If the eigenvalues of \(C_i\) are

\[
\lambda_1\ge \lambda_2\ge\cdots\ge\lambda_d\ge 0,
\]

then the number of large eigenvalues gives a local dimension heuristic.
For example, a curve has approximately one large local eigenvalue, a
surface has approximately two, and a full-dimensional cloud has many.

Anisotropy means that vectors concentrate in a narrow cone or a few
dominant directions. One diagnostic is the spectrum of the centered
covariance matrix

\[
C=\frac{1}{n}\sum_{i=1}^n (x_i-\bar x)(x_i-\bar x)^T.
\]

A very steep spectrum indicates that the cloud is dominated by a few
directions. This can affect nearest-neighbor retrieval and apparent
topological structure.

\section{24.10 Persistent homology and representation
layers}\label{persistent-homology-and-representation-layers}

For a Transformer, each layer produces a representation matrix

\[
H^{(\ell)}\in R^{T\times d},
\]

for tokens in a sentence, or a collection of sentence embeddings across
a dataset. We can treat the rows of \(H^{(\ell)}\) as a point cloud and
compare topological summaries across layers.

For each layer \(\ell\), one may compute a persistence summary

\[
\operatorname{PH}_k(H^{(\ell)}),
\]

where \(k=0\) studies components and \(k=1\) studies loops. This gives a
layerwise diagnostic:

\[
\ell \mapsto \operatorname{PH}_k(H^{(\ell)}).
\]

Possible observations include:

\begin{itemize}
\tightlist
\item
  early layers preserve local lexical neighborhoods;
\item
  middle layers separate syntactic or semantic roles;
\item
  later layers may collapse task-relevant examples into decision
  regions;
\item
  instruction-tuned models may reshape prompt and answer
  representations.
\end{itemize}

These are hypotheses to test, not universal laws.

\section{24.11 Stability viewpoint}\label{stability-viewpoint}

A major reason persistent homology is attractive is stability.
Informally, if two point clouds are close as metric spaces, then their
persistence diagrams are close. The precise form of stability depends on
the construction and the metric used on diagrams.

Let \(D_X\) and \(D_Y\) be persistence diagrams for two point clouds. A
common diagram distance is the bottleneck distance

\[
d_B(D_X,D_Y),
\]

which matches points in one diagram to points in the other, allowing
points to be matched to the diagonal. The diagonal represents
zero-persistence features.

The stability principle says, informally,

\[
\text{small perturbation of distances}
\quad\Longrightarrow\quad
\text{small perturbation of persistence diagrams}.
\]

For NLP, this matters because embeddings are noisy and model-dependent.
A diagnostic that changes wildly under tiny perturbations is not
reliable.

\section{24.12 Mapper and semantic
organization}\label{mapper-and-semantic-organization}

Another topological method is Mapper. Mapper starts with:

\begin{enumerate}
\def\labelenumi{\arabic{enumi}.}
\tightlist
\item
  a point cloud \(X\);
\item
  a filter function \(g:X\to R^q\);
\item
  a cover of the filter range;
\item
  a clustering method inside each cover element.
\end{enumerate}

The output is a graph whose nodes are local clusters and whose edges
indicate overlap. In NLP, possible filter functions include:

\begin{itemize}
\tightlist
\item
  first principal component score;
\item
  sentiment score;
\item
  document length;
\item
  perplexity under a language model;
\item
  retrieval score;
\item
  topic probability;
\item
  projection onto a semantic direction.
\end{itemize}

Mapper can reveal branching or circular organization, but it has tuning
choices: cover size, overlap, clustering method, and distance metric.

\section{24.13 A worked example: six sentence
embeddings}\label{a-worked-example-six-sentence-embeddings}

Suppose six short sentence embeddings have pairwise distances arranged
so that points \(1,2,3\) are close to each other and points \(4,5,6\)
are close to each other, with a larger distance between the two groups.

At a small scale \(\varepsilon\), all six points are disconnected:

\[
\beta_0=6,
\qquad
\beta_1=0.
\]

At a moderate scale, each group forms a triangle. If the triangles are
filled in the Rips complex, each group is contractible:

\[
\beta_0=2,
\qquad
\beta_1=0.
\]

At a larger scale, the two groups connect by at least one edge:

\[
\beta_0=1.
\]

This is a topological version of clustering. The death times of
connected components encode inter-cluster distances. In dimension 0,
persistent homology is closely related to single-linkage hierarchical
clustering.

\section{24.14 Interactive: Vietoris--Rips complex for a toy embedding
cloud}\label{interactive-vietorisrips-complex-for-a-toy-embedding-cloud}

The interactive below uses a two-dimensional toy embedding cloud with
two semantic arcs and a bridge point. Move the distance threshold
\(\varepsilon\). The plot draws the neighborhood graph. The summary
reports connected components, edges, filled triangles, and an exact
\(\beta_1\) value for the clique complex up to dimension 2.

\section{24.15 Python example: pairwise
distances}\label{python-example-pairwise-distances}

We begin with a tiny collection of two-dimensional embeddings. In a real
NLP application these would come from TF-IDF, SVD, Word2Vec, a
sentence-transformer model, or a Transformer layer.

\begin{Shaded}
\begin{Highlighting}[]
\ImportTok{import}\NormalTok{ numpy }\ImportTok{as}\NormalTok{ np}
\ImportTok{import}\NormalTok{ pandas }\ImportTok{as}\NormalTok{ pd}

\NormalTok{labels }\OperatorTok{=}\NormalTok{ [}\StringTok{"morning"}\NormalTok{, }\StringTok{"class"}\NormalTok{, }\StringTok{"lunch"}\NormalTok{, }\StringTok{"office"}\NormalTok{, }\StringTok{"evening"}\NormalTok{, }\StringTok{"reading"}\NormalTok{, }\StringTok{"sleep"}\NormalTok{, }\StringTok{"dream"}\NormalTok{]}
\NormalTok{X }\OperatorTok{=}\NormalTok{ np.array([}
\NormalTok{    [}\FloatTok{1.00}\NormalTok{, }\FloatTok{0.00}\NormalTok{],}
\NormalTok{    [}\FloatTok{0.70}\NormalTok{, }\FloatTok{0.72}\NormalTok{],}
\NormalTok{    [}\OperatorTok{{-}}\FloatTok{0.05}\NormalTok{, }\FloatTok{1.02}\NormalTok{],}
\NormalTok{    [}\OperatorTok{{-}}\FloatTok{0.78}\NormalTok{, }\FloatTok{0.65}\NormalTok{],}
\NormalTok{    [}\OperatorTok{{-}}\FloatTok{1.02}\NormalTok{, }\OperatorTok{{-}}\FloatTok{0.05}\NormalTok{],}
\NormalTok{    [}\OperatorTok{{-}}\FloatTok{0.62}\NormalTok{, }\OperatorTok{{-}}\FloatTok{0.82}\NormalTok{],}
\NormalTok{    [}\FloatTok{0.15}\NormalTok{, }\OperatorTok{{-}}\FloatTok{0.98}\NormalTok{],}
\NormalTok{    [}\FloatTok{0.86}\NormalTok{, }\OperatorTok{{-}}\FloatTok{0.48}\NormalTok{],}
\NormalTok{])}

\NormalTok{D }\OperatorTok{=}\NormalTok{ np.sqrt(((X[:, }\VariableTok{None}\NormalTok{, :] }\OperatorTok{{-}}\NormalTok{ X[}\VariableTok{None}\NormalTok{, :, :]) }\OperatorTok{**} \DecValTok{2}\NormalTok{).}\BuiltInTok{sum}\NormalTok{(axis}\OperatorTok{=}\DecValTok{2}\NormalTok{))}
\NormalTok{pd.DataFrame(np.}\BuiltInTok{round}\NormalTok{(D, }\DecValTok{2}\NormalTok{), index}\OperatorTok{=}\NormalTok{labels, columns}\OperatorTok{=}\NormalTok{labels)}
\end{Highlighting}
\end{Shaded}

{\def\LTcaptype{none} % do not increment counter
\begin{longtable}[]{@{}lllllllll@{}}
\toprule\noalign{}
& morning & class & lunch & office & evening & reading & sleep &
dream \\
\midrule\noalign{}
\endhead
\bottomrule\noalign{}
\endlastfoot
morning & 0.00 & 0.78 & 1.46 & 1.89 & 2.02 & 1.82 & 1.30 & 0.50 \\
class & 0.78 & 0.00 & 0.81 & 1.48 & 1.88 & 2.03 & 1.79 & 1.21 \\
lunch & 1.46 & 0.81 & 0.00 & 0.82 & 1.44 & 1.93 & 2.01 & 1.75 \\
office & 1.89 & 1.48 & 0.82 & 0.00 & 0.74 & 1.48 & 1.88 & 1.99 \\
evening & 2.02 & 1.88 & 1.44 & 0.74 & 0.00 & 0.87 & 1.49 & 1.93 \\
reading & 1.82 & 2.03 & 1.93 & 1.48 & 0.87 & 0.00 & 0.79 & 1.52 \\
sleep & 1.30 & 1.79 & 2.01 & 1.88 & 1.49 & 0.79 & 0.00 & 0.87 \\
dream & 0.50 & 1.21 & 1.75 & 1.99 & 1.93 & 1.52 & 0.87 & 0.00 \\
\end{longtable}
}

The distance matrix is the input to the Vietoris--Rips construction.

\section{24.16 Python example: neighborhood graph and
components}\label{python-example-neighborhood-graph-and-components}

The following code builds the \(\varepsilon\)-neighborhood graph and
counts connected components.

\begin{Shaded}
\begin{Highlighting}[]
\KeywordTok{def}\NormalTok{ epsilon\_edges(D, eps):}
\NormalTok{    n }\OperatorTok{=}\NormalTok{ D.shape[}\DecValTok{0}\NormalTok{]}
\NormalTok{    edges }\OperatorTok{=}\NormalTok{ []}
    \ControlFlowTok{for}\NormalTok{ i }\KeywordTok{in} \BuiltInTok{range}\NormalTok{(n):}
        \ControlFlowTok{for}\NormalTok{ j }\KeywordTok{in} \BuiltInTok{range}\NormalTok{(i }\OperatorTok{+} \DecValTok{1}\NormalTok{, n):}
            \ControlFlowTok{if}\NormalTok{ D[i, j] }\OperatorTok{\textless{}=}\NormalTok{ eps:}
\NormalTok{                edges.append((i, j))}
    \ControlFlowTok{return}\NormalTok{ edges}

\KeywordTok{class}\NormalTok{ UnionFind:}
    \KeywordTok{def} \FunctionTok{\_\_init\_\_}\NormalTok{(}\VariableTok{self}\NormalTok{, n):}
        \VariableTok{self}\NormalTok{.parent }\OperatorTok{=} \BuiltInTok{list}\NormalTok{(}\BuiltInTok{range}\NormalTok{(n))}

    \KeywordTok{def}\NormalTok{ find(}\VariableTok{self}\NormalTok{, a):}
        \ControlFlowTok{while} \VariableTok{self}\NormalTok{.parent[a] }\OperatorTok{!=}\NormalTok{ a:}
            \VariableTok{self}\NormalTok{.parent[a] }\OperatorTok{=} \VariableTok{self}\NormalTok{.parent[}\VariableTok{self}\NormalTok{.parent[a]]}
\NormalTok{            a }\OperatorTok{=} \VariableTok{self}\NormalTok{.parent[a]}
        \ControlFlowTok{return}\NormalTok{ a}

    \KeywordTok{def}\NormalTok{ union(}\VariableTok{self}\NormalTok{, a, b):}
\NormalTok{        ra, rb }\OperatorTok{=} \VariableTok{self}\NormalTok{.find(a), }\VariableTok{self}\NormalTok{.find(b)}
        \ControlFlowTok{if}\NormalTok{ ra }\OperatorTok{!=}\NormalTok{ rb:}
            \VariableTok{self}\NormalTok{.parent[rb] }\OperatorTok{=}\NormalTok{ ra}

    \KeywordTok{def}\NormalTok{ count\_components(}\VariableTok{self}\NormalTok{):}
        \ControlFlowTok{return} \BuiltInTok{len}\NormalTok{(\{}\VariableTok{self}\NormalTok{.find(i) }\ControlFlowTok{for}\NormalTok{ i }\KeywordTok{in} \BuiltInTok{range}\NormalTok{(}\BuiltInTok{len}\NormalTok{(}\VariableTok{self}\NormalTok{.parent))\})}

\KeywordTok{def}\NormalTok{ beta0\_graph(D, eps):}
\NormalTok{    n }\OperatorTok{=}\NormalTok{ D.shape[}\DecValTok{0}\NormalTok{]}
\NormalTok{    uf }\OperatorTok{=}\NormalTok{ UnionFind(n)}
\NormalTok{    edges }\OperatorTok{=}\NormalTok{ epsilon\_edges(D, eps)}
    \ControlFlowTok{for}\NormalTok{ i, j }\KeywordTok{in}\NormalTok{ edges:}
\NormalTok{        uf.union(i, j)}
    \ControlFlowTok{return}\NormalTok{ uf.count\_components(), edges}

\ControlFlowTok{for}\NormalTok{ eps }\KeywordTok{in}\NormalTok{ [}\FloatTok{0.4}\NormalTok{, }\FloatTok{0.8}\NormalTok{, }\FloatTok{1.1}\NormalTok{, }\FloatTok{1.5}\NormalTok{]:}
\NormalTok{    b0, edges }\OperatorTok{=}\NormalTok{ beta0\_graph(D, eps)}
    \BuiltInTok{print}\NormalTok{(}\SpecialStringTok{f"epsilon=}\SpecialCharTok{\{}\NormalTok{eps}\SpecialCharTok{:.1f\}}\SpecialStringTok{, edges=}\SpecialCharTok{\{}\BuiltInTok{len}\NormalTok{(edges)}\SpecialCharTok{\}}\SpecialStringTok{, beta0=}\SpecialCharTok{\{}\NormalTok{b0}\SpecialCharTok{\}}\SpecialStringTok{"}\NormalTok{)}
\end{Highlighting}
\end{Shaded}

\begin{verbatim}
epsilon=0.4, edges=0, beta0=8
epsilon=0.8, edges=4, beta0=4
epsilon=1.1, edges=8, beta0=1
epsilon=1.5, edges=15, beta0=1
\end{verbatim}

\section{24.17 Python example: Betti numbers for a Rips complex up to
dimension
2}\label{python-example-betti-numbers-for-a-rips-complex-up-to-dimension-2}

Now we include filled triangles. The function below computes \(\beta_0\)
and \(\beta_1\) for the clique complex up to dimension 2.

\begin{Shaded}
\begin{Highlighting}[]
\KeywordTok{def}\NormalTok{ rips\_edges\_triangles(D, eps):}
\NormalTok{    n }\OperatorTok{=}\NormalTok{ D.shape[}\DecValTok{0}\NormalTok{]}
\NormalTok{    edges }\OperatorTok{=}\NormalTok{ epsilon\_edges(D, eps)}
\NormalTok{    edge\_index }\OperatorTok{=}\NormalTok{ \{edge: k }\ControlFlowTok{for}\NormalTok{ k, edge }\KeywordTok{in} \BuiltInTok{enumerate}\NormalTok{(edges)\}}
\NormalTok{    triangles }\OperatorTok{=}\NormalTok{ []}
    \ControlFlowTok{for}\NormalTok{ i }\KeywordTok{in} \BuiltInTok{range}\NormalTok{(n):}
        \ControlFlowTok{for}\NormalTok{ j }\KeywordTok{in} \BuiltInTok{range}\NormalTok{(i }\OperatorTok{+} \DecValTok{1}\NormalTok{, n):}
            \ControlFlowTok{for}\NormalTok{ k }\KeywordTok{in} \BuiltInTok{range}\NormalTok{(j }\OperatorTok{+} \DecValTok{1}\NormalTok{, n):}
                \ControlFlowTok{if}\NormalTok{ ((i, j) }\KeywordTok{in}\NormalTok{ edge\_index }\KeywordTok{and}
\NormalTok{                    (i, k) }\KeywordTok{in}\NormalTok{ edge\_index }\KeywordTok{and}
\NormalTok{                    (j, k) }\KeywordTok{in}\NormalTok{ edge\_index):}
\NormalTok{                    triangles.append((i, j, k))}
    \ControlFlowTok{return}\NormalTok{ edges, edge\_index, triangles}

\KeywordTok{def}\NormalTok{ rank\_matrix(A, tol}\OperatorTok{=}\FloatTok{1e{-}10}\NormalTok{):}
    \ControlFlowTok{if}\NormalTok{ A.size }\OperatorTok{==} \DecValTok{0}\NormalTok{:}
        \ControlFlowTok{return} \DecValTok{0}
    \ControlFlowTok{return}\NormalTok{ np.linalg.matrix\_rank(A, tol}\OperatorTok{=}\NormalTok{tol)}

\KeywordTok{def}\NormalTok{ betti\_0\_1\_rips(D, eps):}
\NormalTok{    n }\OperatorTok{=}\NormalTok{ D.shape[}\DecValTok{0}\NormalTok{]}
\NormalTok{    beta0, edges }\OperatorTok{=}\NormalTok{ beta0\_graph(D, eps)}
\NormalTok{    edges, edge\_index, triangles }\OperatorTok{=}\NormalTok{ rips\_edges\_triangles(D, eps)}

\NormalTok{    graph\_cycle\_rank }\OperatorTok{=} \BuiltInTok{len}\NormalTok{(edges) }\OperatorTok{{-}}\NormalTok{ n }\OperatorTok{+}\NormalTok{ beta0}

    \CommentTok{\# Boundary matrix partial\_2: C\_2 {-}\textgreater{} C\_1.}
\NormalTok{    B2 }\OperatorTok{=}\NormalTok{ np.zeros((}\BuiltInTok{len}\NormalTok{(edges), }\BuiltInTok{len}\NormalTok{(triangles)))}
    \ControlFlowTok{for}\NormalTok{ col, (i, j, k) }\KeywordTok{in} \BuiltInTok{enumerate}\NormalTok{(triangles):}
        \CommentTok{\# oriented boundary: [j,k] {-} [i,k] + [i,j]}
\NormalTok{        B2[edge\_index[(j, k)], col] }\OperatorTok{=} \DecValTok{1}
\NormalTok{        B2[edge\_index[(i, k)], col] }\OperatorTok{=} \OperatorTok{{-}}\DecValTok{1}
\NormalTok{        B2[edge\_index[(i, j)], col] }\OperatorTok{=} \DecValTok{1}

\NormalTok{    rank\_b2 }\OperatorTok{=}\NormalTok{ rank\_matrix(B2)}
\NormalTok{    beta1 }\OperatorTok{=}\NormalTok{ graph\_cycle\_rank }\OperatorTok{{-}}\NormalTok{ rank\_b2}
    \ControlFlowTok{return}\NormalTok{ \{}
        \StringTok{"epsilon"}\NormalTok{: eps,}
        \StringTok{"vertices"}\NormalTok{: n,}
        \StringTok{"edges"}\NormalTok{: }\BuiltInTok{len}\NormalTok{(edges),}
        \StringTok{"triangles"}\NormalTok{: }\BuiltInTok{len}\NormalTok{(triangles),}
        \StringTok{"beta0"}\NormalTok{: beta0,}
        \StringTok{"graph\_cycle\_rank"}\NormalTok{: graph\_cycle\_rank,}
        \StringTok{"rank\_boundary\_2"}\NormalTok{: rank\_b2,}
        \StringTok{"beta1"}\NormalTok{: }\BuiltInTok{int}\NormalTok{(}\BuiltInTok{round}\NormalTok{(beta1)),}
\NormalTok{    \}}

\NormalTok{summary }\OperatorTok{=}\NormalTok{ [betti\_0\_1\_rips(D, eps) }\ControlFlowTok{for}\NormalTok{ eps }\KeywordTok{in}\NormalTok{ [}\FloatTok{0.4}\NormalTok{, }\FloatTok{0.8}\NormalTok{, }\FloatTok{1.1}\NormalTok{, }\FloatTok{1.5}\NormalTok{]]}
\NormalTok{pd.DataFrame(summary)}
\end{Highlighting}
\end{Shaded}

{\def\LTcaptype{none} % do not increment counter
\begin{longtable}[]{@{}lllllllll@{}}
\toprule\noalign{}
& epsilon & vertices & edges & triangles & beta0 & graph\_cycle\_rank &
rank\_boundary\_2 & beta1 \\
\midrule\noalign{}
\endhead
\bottomrule\noalign{}
\endlastfoot
0 & 0.4 & 8 & 0 & 0 & 8 & 0 & 0 & 0 \\
1 & 0.8 & 8 & 4 & 0 & 4 & 0 & 0 & 0 \\
2 & 1.1 & 8 & 8 & 0 & 1 & 1 & 0 & 1 \\
3 & 1.5 & 8 & 15 & 7 & 1 & 8 & 7 & 1 \\
\end{longtable}
}

The graph cycle rank can be positive even when \(\beta_1=0\), because
filled triangles may kill cycles.

\section{24.18 Python example: a small persistence
table}\label{python-example-a-small-persistence-table}

A persistence algorithm tracks births and deaths exactly. Here we build
a simple scale table for \(\beta_0\) and \(\beta_1\) by evaluating
several thresholds.

\begin{Shaded}
\begin{Highlighting}[]
\NormalTok{eps\_grid }\OperatorTok{=}\NormalTok{ np.linspace(}\FloatTok{0.25}\NormalTok{, }\FloatTok{2.0}\NormalTok{, }\DecValTok{15}\NormalTok{)}
\NormalTok{rows }\OperatorTok{=}\NormalTok{ [betti\_0\_1\_rips(D, }\BuiltInTok{float}\NormalTok{(eps)) }\ControlFlowTok{for}\NormalTok{ eps }\KeywordTok{in}\NormalTok{ eps\_grid]}
\NormalTok{betti\_table }\OperatorTok{=}\NormalTok{ pd.DataFrame(rows)}
\NormalTok{betti\_table[[}\StringTok{"epsilon"}\NormalTok{, }\StringTok{"edges"}\NormalTok{, }\StringTok{"triangles"}\NormalTok{, }\StringTok{"beta0"}\NormalTok{, }\StringTok{"beta1"}\NormalTok{]]}
\end{Highlighting}
\end{Shaded}

{\def\LTcaptype{none} % do not increment counter
\begin{longtable}[]{@{}llllll@{}}
\toprule\noalign{}
& epsilon & edges & triangles & beta0 & beta1 \\
\midrule\noalign{}
\endhead
\bottomrule\noalign{}
\endlastfoot
0 & 0.250 & 0 & 0 & 8 & 0 \\
1 & 0.375 & 0 & 0 & 8 & 0 \\
2 & 0.500 & 1 & 0 & 7 & 0 \\
3 & 0.625 & 1 & 0 & 7 & 0 \\
4 & 0.750 & 2 & 0 & 6 & 0 \\
5 & 0.875 & 8 & 0 & 1 & 1 \\
6 & 1.000 & 8 & 0 & 1 & 1 \\
7 & 1.125 & 8 & 0 & 1 & 1 \\
8 & 1.250 & 9 & 1 & 1 & 1 \\
9 & 1.375 & 10 & 2 & 1 & 1 \\
10 & 1.500 & 15 & 7 & 1 & 1 \\
11 & 1.625 & 16 & 8 & 1 & 1 \\
12 & 1.750 & 16 & 8 & 1 & 1 \\
13 & 1.875 & 19 & 14 & 1 & 1 \\
14 & 2.000 & 25 & 38 & 1 & 0 \\
\end{longtable}
}

This table is not a replacement for a full persistent homology library,
but it shows the idea: topology changes as the scale changes.

\section{24.19 Python example: local intrinsic
dimension}\label{python-example-local-intrinsic-dimension}

We estimate local dimension by looking at the spectrum of local
covariance matrices.

\begin{Shaded}
\begin{Highlighting}[]
\KeywordTok{def}\NormalTok{ local\_dimension\_scores(X, k}\OperatorTok{=}\DecValTok{3}\NormalTok{):}
\NormalTok{    D }\OperatorTok{=}\NormalTok{ np.sqrt(((X[:, }\VariableTok{None}\NormalTok{, :] }\OperatorTok{{-}}\NormalTok{ X[}\VariableTok{None}\NormalTok{, :, :]) }\OperatorTok{**} \DecValTok{2}\NormalTok{).}\BuiltInTok{sum}\NormalTok{(axis}\OperatorTok{=}\DecValTok{2}\NormalTok{))}
\NormalTok{    scores }\OperatorTok{=}\NormalTok{ []}
    \ControlFlowTok{for}\NormalTok{ i }\KeywordTok{in} \BuiltInTok{range}\NormalTok{(X.shape[}\DecValTok{0}\NormalTok{]):}
\NormalTok{        nn }\OperatorTok{=}\NormalTok{ np.argsort(D[i])[}\DecValTok{1}\NormalTok{:k}\OperatorTok{+}\DecValTok{1}\NormalTok{]}
\NormalTok{        Y }\OperatorTok{=}\NormalTok{ X[nn] }\OperatorTok{{-}}\NormalTok{ X[i]}
\NormalTok{        C }\OperatorTok{=}\NormalTok{ (Y.T }\OperatorTok{@}\NormalTok{ Y) }\OperatorTok{/}\NormalTok{ k}
\NormalTok{        eigvals }\OperatorTok{=}\NormalTok{ np.linalg.eigvalsh(C)[::}\OperatorTok{{-}}\DecValTok{1}\NormalTok{]}
\NormalTok{        ratio }\OperatorTok{=}\NormalTok{ eigvals[}\DecValTok{0}\NormalTok{] }\OperatorTok{/}\NormalTok{ eigvals.}\BuiltInTok{sum}\NormalTok{() }\ControlFlowTok{if}\NormalTok{ eigvals.}\BuiltInTok{sum}\NormalTok{() }\OperatorTok{\textgreater{}} \DecValTok{0} \ControlFlowTok{else}\NormalTok{ np.nan}
\NormalTok{        scores.append(\{}
            \StringTok{"point"}\NormalTok{: labels[i],}
            \StringTok{"lambda1"}\NormalTok{: eigvals[}\DecValTok{0}\NormalTok{],}
            \StringTok{"lambda2"}\NormalTok{: eigvals[}\DecValTok{1}\NormalTok{],}
            \StringTok{"lambda1\_fraction"}\NormalTok{: ratio,}
\NormalTok{        \})}
    \ControlFlowTok{return}\NormalTok{ pd.DataFrame(scores)}

\NormalTok{local\_dimension\_scores(X, k}\OperatorTok{=}\DecValTok{3}\NormalTok{)}
\end{Highlighting}
\end{Shaded}

{\def\LTcaptype{none} % do not increment counter
\begin{longtable}[]{@{}lllll@{}}
\toprule\noalign{}
& point & lambda1 & lambda2 & lambda1\_fraction \\
\midrule\noalign{}
\endhead
\bottomrule\noalign{}
\endlastfoot
0 & morning & 0.694445 & 0.152655 & 0.819791 \\
1 & class & 0.765351 & 0.143483 & 0.842124 \\
2 & lunch & 0.945623 & 0.190410 & 0.832390 \\
3 & office & 0.935531 & 0.199102 & 0.824523 \\
4 & evening & 0.912725 & 0.216042 & 0.808603 \\
5 & reading & 0.993676 & 0.192291 & 0.837861 \\
6 & sleep & 0.877265 & 0.141235 & 0.861331 \\
7 & dream & 0.652672 & 0.170561 & 0.792816 \\
\end{longtable}
}

A point with most local variance in one direction has a curve-like
neighborhood. A point with substantial variance in two directions has a
more surface-like neighborhood.

\section{24.20 Python example: topology as a representation
diagnostic}\label{python-example-topology-as-a-representation-diagnostic}

The next code compares a circular cloud and a two-cluster cloud. Both
may have similar variance, but their topological summaries behave
differently.

\begin{Shaded}
\begin{Highlighting}[]
\NormalTok{rng }\OperatorTok{=}\NormalTok{ np.random.default\_rng(}\DecValTok{24}\NormalTok{)}

\NormalTok{angles }\OperatorTok{=}\NormalTok{ np.linspace(}\DecValTok{0}\NormalTok{, }\DecValTok{2}\OperatorTok{*}\NormalTok{np.pi, }\DecValTok{18}\NormalTok{, endpoint}\OperatorTok{=}\VariableTok{False}\NormalTok{)}
\NormalTok{circle }\OperatorTok{=}\NormalTok{ np.c\_[np.cos(angles), np.sin(angles)] }\OperatorTok{+} \FloatTok{0.05} \OperatorTok{*}\NormalTok{ rng.normal(size}\OperatorTok{=}\NormalTok{(}\DecValTok{18}\NormalTok{, }\DecValTok{2}\NormalTok{))}

\NormalTok{cluster\_a }\OperatorTok{=}\NormalTok{ rng.normal(loc}\OperatorTok{=}\NormalTok{[}\OperatorTok{{-}}\FloatTok{1.0}\NormalTok{, }\FloatTok{0.0}\NormalTok{], scale}\OperatorTok{=}\FloatTok{0.18}\NormalTok{, size}\OperatorTok{=}\NormalTok{(}\DecValTok{9}\NormalTok{, }\DecValTok{2}\NormalTok{))}
\NormalTok{cluster\_b }\OperatorTok{=}\NormalTok{ rng.normal(loc}\OperatorTok{=}\NormalTok{[}\FloatTok{1.0}\NormalTok{, }\FloatTok{0.0}\NormalTok{], scale}\OperatorTok{=}\FloatTok{0.18}\NormalTok{, size}\OperatorTok{=}\NormalTok{(}\DecValTok{9}\NormalTok{, }\DecValTok{2}\NormalTok{))}
\NormalTok{clusters }\OperatorTok{=}\NormalTok{ np.vstack([cluster\_a, cluster\_b])}

\KeywordTok{def}\NormalTok{ distance\_matrix(X):}
    \ControlFlowTok{return}\NormalTok{ np.sqrt(((X[:, }\VariableTok{None}\NormalTok{, :] }\OperatorTok{{-}}\NormalTok{ X[}\VariableTok{None}\NormalTok{, :, :]) }\OperatorTok{**} \DecValTok{2}\NormalTok{).}\BuiltInTok{sum}\NormalTok{(axis}\OperatorTok{=}\DecValTok{2}\NormalTok{))}

\KeywordTok{def}\NormalTok{ betti\_curve(X, eps\_grid):}
\NormalTok{    D }\OperatorTok{=}\NormalTok{ distance\_matrix(X)}
\NormalTok{    rows }\OperatorTok{=}\NormalTok{ []}
    \ControlFlowTok{for}\NormalTok{ eps }\KeywordTok{in}\NormalTok{ eps\_grid:}
\NormalTok{        rows.append(betti\_0\_1\_rips(D, }\BuiltInTok{float}\NormalTok{(eps)))}
    \ControlFlowTok{return}\NormalTok{ pd.DataFrame(rows)}

\NormalTok{eps\_grid }\OperatorTok{=}\NormalTok{ np.linspace(}\FloatTok{0.15}\NormalTok{, }\FloatTok{0.75}\NormalTok{, }\DecValTok{7}\NormalTok{)}
\NormalTok{circle\_curve }\OperatorTok{=}\NormalTok{ betti\_curve(circle, eps\_grid)}
\NormalTok{cluster\_curve }\OperatorTok{=}\NormalTok{ betti\_curve(clusters, eps\_grid)}

\NormalTok{comparison }\OperatorTok{=}\NormalTok{ pd.DataFrame(\{}
    \StringTok{"epsilon"}\NormalTok{: eps\_grid,}
    \StringTok{"circle\_beta0"}\NormalTok{: circle\_curve[}\StringTok{"beta0"}\NormalTok{],}
    \StringTok{"circle\_beta1"}\NormalTok{: circle\_curve[}\StringTok{"beta1"}\NormalTok{],}
    \StringTok{"clusters\_beta0"}\NormalTok{: cluster\_curve[}\StringTok{"beta0"}\NormalTok{],}
    \StringTok{"clusters\_beta1"}\NormalTok{: cluster\_curve[}\StringTok{"beta1"}\NormalTok{],}
\NormalTok{\})}
\NormalTok{comparison}
\end{Highlighting}
\end{Shaded}

{\def\LTcaptype{none} % do not increment counter
\begin{longtable}[]{@{}llllll@{}}
\toprule\noalign{}
& epsilon & circle\_beta0 & circle\_beta1 & clusters\_beta0 &
clusters\_beta1 \\
\midrule\noalign{}
\endhead
\bottomrule\noalign{}
\endlastfoot
0 & 0.15 & 18 & 0 & 8 & 0 \\
1 & 0.25 & 18 & 0 & 7 & 0 \\
2 & 0.35 & 7 & 0 & 4 & 0 \\
3 & 0.45 & 1 & 0 & 2 & 0 \\
4 & 0.55 & 1 & 1 & 2 & 0 \\
5 & 0.65 & 1 & 1 & 2 & 0 \\
6 & 0.75 & 1 & 1 & 2 & 0 \\
\end{longtable}
}

The circular cloud tends to develop a one-dimensional loop over an
intermediate range of scales. The two-cluster cloud emphasizes component
merging instead.

\section{24.21 Python Plotly
visualization}\label{python-plotly-visualization}

The following code produces a Plotly visualization in Python. It is
useful when students work in Jupyter or Google Colab.

\begin{Shaded}
\begin{Highlighting}[]
\ImportTok{import}\NormalTok{ plotly.graph\_objects }\ImportTok{as}\NormalTok{ go}

\NormalTok{eps }\OperatorTok{=} \FloatTok{0.8}
\NormalTok{b0, edges }\OperatorTok{=}\NormalTok{ beta0\_graph(D, eps)}
\NormalTok{edge\_x, edge\_y }\OperatorTok{=}\NormalTok{ [], []}
\ControlFlowTok{for}\NormalTok{ i, j }\KeywordTok{in}\NormalTok{ edges:}
\NormalTok{    edge\_x }\OperatorTok{+=}\NormalTok{ [X[i, }\DecValTok{0}\NormalTok{], X[j, }\DecValTok{0}\NormalTok{], }\VariableTok{None}\NormalTok{]}
\NormalTok{    edge\_y }\OperatorTok{+=}\NormalTok{ [X[i, }\DecValTok{1}\NormalTok{], X[j, }\DecValTok{1}\NormalTok{], }\VariableTok{None}\NormalTok{]}

\NormalTok{fig }\OperatorTok{=}\NormalTok{ go.Figure()}
\NormalTok{fig.add\_trace(go.Scatter(x}\OperatorTok{=}\NormalTok{edge\_x, y}\OperatorTok{=}\NormalTok{edge\_y, mode}\OperatorTok{=}\StringTok{"lines"}\NormalTok{, name}\OperatorTok{=}\StringTok{"edges"}\NormalTok{))}
\NormalTok{fig.add\_trace(go.Scatter(}
\NormalTok{    x}\OperatorTok{=}\NormalTok{X[:, }\DecValTok{0}\NormalTok{], y}\OperatorTok{=}\NormalTok{X[:, }\DecValTok{1}\NormalTok{], mode}\OperatorTok{=}\StringTok{"markers+text"}\NormalTok{,}
\NormalTok{    text}\OperatorTok{=}\NormalTok{labels, textposition}\OperatorTok{=}\StringTok{"top center"}\NormalTok{, name}\OperatorTok{=}\StringTok{"points"}
\NormalTok{))}
\NormalTok{fig.update\_layout(}
\NormalTok{    title}\OperatorTok{=}\SpecialStringTok{f"Epsilon graph with beta0=}\SpecialCharTok{\{}\NormalTok{b0}\SpecialCharTok{\}}\SpecialStringTok{"}\NormalTok{,}
\NormalTok{    xaxis}\OperatorTok{=}\BuiltInTok{dict}\NormalTok{(scaleanchor}\OperatorTok{=}\StringTok{"y"}\NormalTok{),}
\NormalTok{    width}\OperatorTok{=}\DecValTok{700}\NormalTok{,}
\NormalTok{    height}\OperatorTok{=}\DecValTok{500}
\NormalTok{)}
\NormalTok{fig.show()}
\end{Highlighting}
\end{Shaded}

\section{24.22 AI-assisted learning: topology audit
prompts}\label{ai-assisted-learning-topology-audit-prompts}

AI tools can be helpful for brainstorming diagnostics, but they should
not replace mathematical checks. Students can use the following prompts:

\textbf{Prompt 1: representation-space audit}

I have sentence embeddings from two models. Suggest a mathematical plan
for comparing their geometry and topology. Include distance choices,
normalization choices, PCA/SVD diagnostics, connected-component
persistence, and possible failure modes.

\textbf{Prompt 2: interpretation critique}

A persistent loop appears in a sentence-embedding point cloud. Give at
least five possible explanations, including semantic explanations and
artifact explanations. For each explanation, propose a test.

\textbf{Prompt 3: code review}

Review my Python code for a Vietoris--Rips complex. Check whether I
incorrectly count graph cycles as topological loops after triangles are
filled. Explain the role of the boundary matrix \(\partial_2\).

\section{24.23 Common mistakes}\label{common-mistakes-13}

\begin{enumerate}
\def\labelenumi{\arabic{enumi}.}
\item
  \textbf{Confusing a graph cycle with a topological loop.} A triangle
  in a graph is a cycle, but in a Rips complex it is usually filled, so
  it does not contribute to \(\beta_1\).
\item
  \textbf{Using dimensionality reduction as evidence of topology.} A
  loop in a two-dimensional visualization may be introduced or destroyed
  by projection. Compute topology from the original metric when
  possible.
\item
  \textbf{Ignoring normalization.} Euclidean distance and cosine
  distance can give very different filtrations.
\item
  \textbf{Overinterpreting short bars.} Low-persistence features may be
  noise or sampling artifacts.
\item
  \textbf{Forgetting sample-size effects.} Persistent homology depends
  on how densely the underlying shape is sampled.
\item
  \textbf{Treating persistent homology as a black box.} Students should
  understand vertices, edges, triangles, boundary maps, and Betti
  numbers before using a library.
\end{enumerate}

\section{24.24 Exercises}\label{exercises-23}

\begin{enumerate}
\def\labelenumi{\arabic{enumi}.}
\item
  Let \(X\) be four points forming the corners of a square. Compute
  \(\beta_0\) and graph cycle rank for the neighborhood graph at a scale
  where adjacent corners are connected but diagonals are not.
\item
  For the same square, decide whether the Vietoris--Rips complex has a
  one-dimensional loop at that scale.
\item
  Increase the scale until the diagonals appear. What happens to the
  Rips complex and \(\beta_1\)?
\item
  Let \(X\) be three points forming a triangle. Explain why the graph
  cycle rank is 1 but the Rips complex has \(\beta_1=0\) once the filled
  triangle appears.
\item
  Suppose two embedding models produce the same nearest-neighbor
  classification accuracy. Design a topological diagnostic that could
  still distinguish them.
\item
  Explain why cosine distance may be preferable to Euclidean distance
  for sentence embeddings that vary strongly in norm.
\item
  Derive the formula
  \(\beta_1=|E|-|V|+\beta_0-\operatorname{rank}\partial_2\) for a
  2-dimensional simplicial complex.
\item
  Construct a toy dataset of 12 points sampled from two circles. Predict
  how \(\beta_0\) and \(\beta_1\) should behave as \(\varepsilon\)
  increases.
\item
  In a retrieval system, suppose retrieved document chunks form three
  disconnected components in embedding space. Give two possible good
  interpretations and two possible bad interpretations.
\item
  Use an AI tool to generate a topological interpretation of a point
  cloud, then write a critique explaining which parts are supported by
  mathematics and which parts are speculative.
\end{enumerate}

\section{24.25 Mini-project: topology of course-advising
embeddings}\label{mini-project-topology-of-course-advising-embeddings}

Build a small corpus of 30--60 short documents about university advising
topics such as registration, transfer credit, degree requirements,
co-op, graduate electives, and academic standing. Compute embeddings
using any available method, for example TF-IDF/SVD or a sentence
embedding model. Then:

\begin{enumerate}
\def\labelenumi{\arabic{enumi}.}
\tightlist
\item
  choose a distance metric;
\item
  compute a distance matrix;
\item
  build \(\varepsilon\)-neighborhood graphs over a range of thresholds;
\item
  plot \(\beta_0\) as a function of \(\varepsilon\);
\item
  compute simple graph cycle ranks and, if possible, Rips \(\beta_1\)
  values;
\item
  compare the topology with known document categories;
\item
  write a short interpretation and a cautionary note.
\end{enumerate}

This project connects topology, retrieval, and responsible
interpretation.

\section{24.26 Lab connection}\label{lab-connection-10}

The accompanying lab should guide students through:

\begin{itemize}
\tightlist
\item
  constructing a point cloud from small text embeddings;
\item
  computing Euclidean and cosine distance matrices;
\item
  building neighborhood graphs;
\item
  computing connected components;
\item
  computing \(\beta_0\) and \(\beta_1\) for small Rips complexes;
\item
  plotting Betti curves;
\item
  comparing circular, clustered, and noisy point clouds;
\item
  applying the same tools to a small NLP corpus.
\end{itemize}

\section{24.27 Summary}\label{summary-6}

This chapter introduced a topological viewpoint on NLP representation
spaces. The central idea is simple: embeddings turn language objects
into points, distances turn point clouds into finite metric spaces, and
filtrations turn metric spaces into multiscale topological summaries.
Betti numbers and persistent homology provide tools for studying
connected components, loops, and other qualitative structures. These
tools are most valuable when used together with linear algebra,
probability, and domain knowledge. In mathematical NLP, topology is not
a replacement for modeling or evaluation; it is a structural lens for
asking better questions about representation spaces.

\chapter{Ethics, Bias, and Mathematical
Responsibility}\label{ethics-bias-and-mathematical-responsibility}

Fairness, uncertainty, privacy, and responsible NLP deployment

\hfill\break

Modern NLP systems are mathematical systems deployed in social
environments. A language model is trained on data, optimized under a
loss function, evaluated by metrics, and then used by people. Ethical
analysis therefore cannot be separated from mathematical analysis.
Before we can ask whether a system is fair, safe, or trustworthy, we
must specify the random variables, the population, the representation,
the target, the loss, the metric, and the deployment setting.

This chapter develops a mathematical language for responsibility in NLP.
The goal is not to reduce ethics to formulas. Rather, the goal is to
make assumptions visible. A formula can clarify what is being measured,
what is being ignored, and where statistical uncertainty enters.

\section{25.1 Learning goals}\label{learning-goals-24}

After reading this chapter, students should be able to:

\begin{itemize}
\tightlist
\item
  formulate bias, fairness, uncertainty, and privacy questions using
  random variables and distributions;
\item
  distinguish sampling bias, measurement bias, label bias,
  representation bias, optimization bias, and deployment bias;
\item
  compute common group fairness metrics from a confusion matrix;
\item
  explain why different fairness metrics can conflict;
\item
  analyze distribution shift using training and deployment
  distributions;
\item
  compute calibration error and understand why accuracy alone is
  insufficient;
\item
  use embedding association scores to audit representation bias;
\item
  describe selective prediction, abstention, and uncertainty-aware
  deployment;
\item
  connect mathematical clarity with responsible AI-assisted NLP
  practice.
\end{itemize}

\section{25.2 The statistical object behind a responsible NLP
system}\label{the-statistical-object-behind-a-responsible-nlp-system}

Let \(X\) denote an input text, such as a document, prompt, search
query, or conversation. Let \(Y\) denote the target variable, such as a
class label, relevance judgment, next token, toxicity rating, or
answer-quality score. Let \(A\) denote a group attribute or context
variable relevant for auditing. Depending on the application, \(A\) may
represent language variety, geography, institution type, dialect,
writing style, document source, demographic group, or another variable
that may affect model behavior.

A responsible statistical analysis begins with a joint distribution

\[
P(X,A,Y).
\]

A trained NLP model is a function, or a conditional distribution, built
from data:

\[
f_\theta:X \mapsto \widehat{Y},
\qquad
p_\theta(y\mid x).
\]

A deployment decision is usually not just the model output. It also
includes a threshold, abstention rule, retrieval rule, or human-review
rule. We therefore write the full deployed rule as

\[
d(x)=\text{decision made after observing }x.
\]

A performance claim has mathematical meaning only relative to a
distribution and a loss. If \(L(d(X),Y)\) is the loss, then the
population risk is

\[
R(d)=E_P[L(d(X),Y)].
\]

The empirical risk computed on a test set is an estimator:

\[
\widehat{R}(d)=\frac{1}{n}\sum_{i=1}^n L(d(x_i),y_i).
\]

The key point is simple but important: \textbf{a metric is not an
ethical conclusion}. A metric is a statistic computed under assumptions.

\textbf{Responsible NLP as a mathematical workflow.}

\begin{enumerate}
\def\labelenumi{\arabic{enumi}.}
\tightlist
\item
  Specify the population \(P(X,A,Y)\).
\item
  Specify the model or decision rule \(d\).
\item
  Specify the loss \(L\).
\item
  Estimate performance and uncertainty.
\item
  Audit performance conditionally on relevant subpopulations.
\item
  Study behavior under distribution shift.
\item
  Decide whether the system should be deployed, constrained, redesigned,
  or rejected.
\end{enumerate}

\section{25.3 Sources of bias in NLP}\label{sources-of-bias-in-nlp}

The word \textbf{bias} has several meanings. In statistics, bias often
refers to

\[
\operatorname{Bias}(\widehat{\theta})=E[\widehat{\theta}]-\theta.
\]

In responsible NLP, bias is broader. It often refers to systematic
error, unequal performance, stereotyped representation, or harmful
behavior. These notions are related, but not identical.

A useful mathematical decomposition is the following.

{\def\LTcaptype{none} % do not increment counter
\begin{longtable}[]{@{}
  >{\raggedright\arraybackslash}p{(\linewidth - 4\tabcolsep) * \real{0.3333}}
  >{\raggedright\arraybackslash}p{(\linewidth - 4\tabcolsep) * \real{0.3333}}
  >{\raggedright\arraybackslash}p{(\linewidth - 4\tabcolsep) * \real{0.3333}}@{}}
\toprule\noalign{}
\begin{minipage}[b]{\linewidth}\raggedright
Source
\end{minipage} & \begin{minipage}[b]{\linewidth}\raggedright
Mathematical form
\end{minipage} & \begin{minipage}[b]{\linewidth}\raggedright
NLP example
\end{minipage} \\
\midrule\noalign{}
\endhead
\bottomrule\noalign{}
\endlastfoot
Sampling bias & \(P_{\text{train}}(X,A,Y) \ne P_{\text{target}}(X,A,Y)\)
& training corpus overrepresents certain genres \\
Measurement bias & observed \(\widetilde{Y}\) differs from intended
\(Y\) & noisy toxicity labels or vague annotation instructions \\
Representation bias & \(\phi(X)\) loses relevant information &
tokenization fails for names, dialects, or multilingual text \\
Label bias & \(P(\widetilde{Y}\mid X,A)\) reflects annotator/systematic
bias & same text rated differently depending on dialect \\
Optimization bias & objective weights some errors more than others &
minimizing average loss hides subgroup error \\
Deployment bias & use differs from training/evaluation assumptions &
classifier used as a gatekeeping tool outside original context \\
Feedback bias & future data depends on earlier model decisions & search
ranking affects what users click next \\
\end{longtable}
}

For MA Applied Math and MS Statistics students, the important habit is
to ask:

\[
\text{Which distribution generated the data, and which distribution will use the model?}
\]

If those distributions differ, a test-set score may not estimate
deployment behavior.

\section{25.4 Group fairness metrics}\label{group-fairness-metrics}

Consider a binary decision problem with label \(Y\in\{0,1\}\) and
prediction \(\widehat{Y}\in\{0,1\}\). Suppose \(A\) is a group variable.
For a fixed group \(a\), define:

\[
\operatorname{TPR}_a=P(\widehat{Y}=1\mid Y=1,A=a),
\]

\[
\operatorname{FPR}_a=P(\widehat{Y}=1\mid Y=0,A=a),
\]

\[
\operatorname{PPV}_a=P(Y=1\mid \widehat{Y}=1,A=a),
\]

\[
\operatorname{SR}_a=P(\widehat{Y}=1\mid A=a),
\]

where TPR is true positive rate, FPR is false positive rate, PPV is
positive predictive value, and SR is selection rate.

Common fairness criteria include:

{\def\LTcaptype{none} % do not increment counter
\begin{longtable}[]{@{}
  >{\raggedright\arraybackslash}p{(\linewidth - 4\tabcolsep) * \real{0.3333}}
  >{\raggedright\arraybackslash}p{(\linewidth - 4\tabcolsep) * \real{0.3333}}
  >{\raggedright\arraybackslash}p{(\linewidth - 4\tabcolsep) * \real{0.3333}}@{}}
\toprule\noalign{}
\begin{minipage}[b]{\linewidth}\raggedright
Criterion
\end{minipage} & \begin{minipage}[b]{\linewidth}\raggedright
Mathematical condition
\end{minipage} & \begin{minipage}[b]{\linewidth}\raggedright
Interpretation
\end{minipage} \\
\midrule\noalign{}
\endhead
\bottomrule\noalign{}
\endlastfoot
Demographic parity &
\(P(\widehat{Y}=1\mid A=a)=P(\widehat{Y}=1\mid A=b)\) & equal selection
rates \\
Equal opportunity &
\(P(\widehat{Y}=1\mid Y=1,A=a)=P(\widehat{Y}=1\mid Y=1,A=b)\) & equal
true positive rates \\
Equalized odds & \(\widehat{Y}\perp A\mid Y\) & equal TPR and FPR \\
Predictive parity &
\(P(Y=1\mid \widehat{Y}=1,A=a)=P(Y=1\mid \widehat{Y}=1,A=b)\) & equal
precision among selected cases \\
Calibration by group & \(P(Y=1\mid S=s,A=a)=s\) & scores mean the same
probability in each group \\
\end{longtable}
}

Here \(S\in[0,1]\) is a model score, such as a predicted probability.

These criteria are not merely technical. They encode different normative
choices. Equalizing selection rates, equalizing error rates, and
equalizing predictive values may lead to different thresholds and
different social consequences.

\section{25.5 Why fairness metrics can
conflict}\label{why-fairness-metrics-can-conflict}

Suppose the base rates differ:

\[
P(Y=1\mid A=a)\ne P(Y=1\mid A=b).
\]

If a classifier is calibrated within each group and is not perfect, then
it is generally impossible to also equalize all error rates across
groups. This is not a paradox. It is a constraint imposed by conditional
probability.

To see the tension, write the selection rate for group \(a\) as

\[
P(\widehat{Y}=1\mid A=a)
=
\operatorname{TPR}_a P(Y=1\mid A=a)
+
\operatorname{FPR}_a P(Y=0\mid A=a).
\]

If the base rate \(P(Y=1\mid A=a)\) changes across groups, then equal
TPR and equal FPR generally imply different selection rates. Conversely,
forcing equal selection rates may require unequal error rates.

\begin{tcolorbox}[enhanced jigsaw, arc=.35mm, bottomrule=.15mm, bottomtitle=1mm, breakable, colback=white, colbacktitle=quarto-callout-important-color!10!white, colframe=quarto-callout-important-color-frame, coltitle=black, left=2mm, leftrule=.75mm, opacityback=0, opacitybacktitle=0.6, rightrule=.15mm, title=\textcolor{quarto-callout-important-color}{\faExclamation}\hspace{0.5em}{Important}, titlerule=0mm, toprule=.15mm, toptitle=1mm]

Fairness metrics should not be selected mechanically. In a real
application, one must ask what the decision means, what errors cost, who
is affected, what appeal process exists, and whether the model should be
used at all.

\end{tcolorbox}

\section{25.6 Confusion matrices by
group}\label{confusion-matrices-by-group}

For a fixed group \(a\), the confusion matrix counts are

\[
\begin{array}{c|cc}
& Y=1 & Y=0 \\
\hline
\widehat{Y}=1 & TP_a & FP_a \\
\widehat{Y}=0 & FN_a & TN_a
\end{array}
\]

From these counts,

\[
\operatorname{TPR}_a=\frac{TP_a}{TP_a+FN_a},
\qquad
\operatorname{FPR}_a=\frac{FP_a}{FP_a+TN_a},
\]

\[
\operatorname{PPV}_a=\frac{TP_a}{TP_a+FP_a},
\qquad
\operatorname{SR}_a=\frac{TP_a+FP_a}{TP_a+FP_a+FN_a+TN_a}.
\]

The following Python example computes these quantities.

\begin{Shaded}
\begin{Highlighting}[]
\ImportTok{import}\NormalTok{ numpy }\ImportTok{as}\NormalTok{ np}
\ImportTok{import}\NormalTok{ pandas }\ImportTok{as}\NormalTok{ pd}

\NormalTok{data }\OperatorTok{=}\NormalTok{ pd.DataFrame(\{}
    \StringTok{"group"}\NormalTok{: [}\StringTok{"A"}\NormalTok{] }\OperatorTok{*} \DecValTok{8} \OperatorTok{+}\NormalTok{ [}\StringTok{"B"}\NormalTok{] }\OperatorTok{*} \DecValTok{8}\NormalTok{,}
    \StringTok{"y"}\NormalTok{:     [}\DecValTok{1}\NormalTok{, }\DecValTok{1}\NormalTok{, }\DecValTok{1}\NormalTok{, }\DecValTok{1}\NormalTok{, }\DecValTok{0}\NormalTok{, }\DecValTok{0}\NormalTok{, }\DecValTok{0}\NormalTok{, }\DecValTok{0}\NormalTok{,  }\DecValTok{1}\NormalTok{, }\DecValTok{1}\NormalTok{, }\DecValTok{0}\NormalTok{, }\DecValTok{0}\NormalTok{, }\DecValTok{0}\NormalTok{, }\DecValTok{0}\NormalTok{, }\DecValTok{0}\NormalTok{, }\DecValTok{0}\NormalTok{],}
    \StringTok{"yhat"}\NormalTok{:  [}\DecValTok{1}\NormalTok{, }\DecValTok{1}\NormalTok{, }\DecValTok{0}\NormalTok{, }\DecValTok{0}\NormalTok{, }\DecValTok{1}\NormalTok{, }\DecValTok{0}\NormalTok{, }\DecValTok{0}\NormalTok{, }\DecValTok{0}\NormalTok{,  }\DecValTok{1}\NormalTok{, }\DecValTok{0}\NormalTok{, }\DecValTok{1}\NormalTok{, }\DecValTok{0}\NormalTok{, }\DecValTok{0}\NormalTok{, }\DecValTok{0}\NormalTok{, }\DecValTok{0}\NormalTok{, }\DecValTok{0}\NormalTok{]}
\NormalTok{\})}

\KeywordTok{def}\NormalTok{ group\_metrics(frame):}
\NormalTok{    rows }\OperatorTok{=}\NormalTok{ []}
    \ControlFlowTok{for}\NormalTok{ g, df }\KeywordTok{in}\NormalTok{ frame.groupby(}\StringTok{"group"}\NormalTok{):}
\NormalTok{        y }\OperatorTok{=}\NormalTok{ df[}\StringTok{"y"}\NormalTok{].to\_numpy()}
\NormalTok{        yhat }\OperatorTok{=}\NormalTok{ df[}\StringTok{"yhat"}\NormalTok{].to\_numpy()}
\NormalTok{        tp }\OperatorTok{=} \BuiltInTok{int}\NormalTok{(((yhat }\OperatorTok{==} \DecValTok{1}\NormalTok{) }\OperatorTok{\&}\NormalTok{ (y }\OperatorTok{==} \DecValTok{1}\NormalTok{)).}\BuiltInTok{sum}\NormalTok{())}
\NormalTok{        fp }\OperatorTok{=} \BuiltInTok{int}\NormalTok{(((yhat }\OperatorTok{==} \DecValTok{1}\NormalTok{) }\OperatorTok{\&}\NormalTok{ (y }\OperatorTok{==} \DecValTok{0}\NormalTok{)).}\BuiltInTok{sum}\NormalTok{())}
\NormalTok{        fn }\OperatorTok{=} \BuiltInTok{int}\NormalTok{(((yhat }\OperatorTok{==} \DecValTok{0}\NormalTok{) }\OperatorTok{\&}\NormalTok{ (y }\OperatorTok{==} \DecValTok{1}\NormalTok{)).}\BuiltInTok{sum}\NormalTok{())}
\NormalTok{        tn }\OperatorTok{=} \BuiltInTok{int}\NormalTok{(((yhat }\OperatorTok{==} \DecValTok{0}\NormalTok{) }\OperatorTok{\&}\NormalTok{ (y }\OperatorTok{==} \DecValTok{0}\NormalTok{)).}\BuiltInTok{sum}\NormalTok{())}
\NormalTok{        rows.append(\{}
            \StringTok{"group"}\NormalTok{: g,}
            \StringTok{"TP"}\NormalTok{: tp,}
            \StringTok{"FP"}\NormalTok{: fp,}
            \StringTok{"FN"}\NormalTok{: fn,}
            \StringTok{"TN"}\NormalTok{: tn,}
            \StringTok{"TPR"}\NormalTok{: tp }\OperatorTok{/}\NormalTok{ (tp }\OperatorTok{+}\NormalTok{ fn) }\ControlFlowTok{if}\NormalTok{ tp }\OperatorTok{+}\NormalTok{ fn }\OperatorTok{\textgreater{}} \DecValTok{0} \ControlFlowTok{else}\NormalTok{ np.nan,}
            \StringTok{"FPR"}\NormalTok{: fp }\OperatorTok{/}\NormalTok{ (fp }\OperatorTok{+}\NormalTok{ tn) }\ControlFlowTok{if}\NormalTok{ fp }\OperatorTok{+}\NormalTok{ tn }\OperatorTok{\textgreater{}} \DecValTok{0} \ControlFlowTok{else}\NormalTok{ np.nan,}
            \StringTok{"PPV"}\NormalTok{: tp }\OperatorTok{/}\NormalTok{ (tp }\OperatorTok{+}\NormalTok{ fp) }\ControlFlowTok{if}\NormalTok{ tp }\OperatorTok{+}\NormalTok{ fp }\OperatorTok{\textgreater{}} \DecValTok{0} \ControlFlowTok{else}\NormalTok{ np.nan,}
            \StringTok{"selection\_rate"}\NormalTok{: (tp }\OperatorTok{+}\NormalTok{ fp) }\OperatorTok{/} \BuiltInTok{len}\NormalTok{(df),}
            \StringTok{"base\_rate"}\NormalTok{: y.mean()}
\NormalTok{        \})}
    \ControlFlowTok{return}\NormalTok{ pd.DataFrame(rows)}

\NormalTok{group\_metrics(data)}
\end{Highlighting}
\end{Shaded}

{\def\LTcaptype{none} % do not increment counter
\begin{longtable}[]{@{}lllllllllll@{}}
\toprule\noalign{}
& group & TP & FP & FN & TN & TPR & FPR & PPV & selection\_rate &
base\_rate \\
\midrule\noalign{}
\endhead
\bottomrule\noalign{}
\endlastfoot
0 & A & 2 & 1 & 2 & 3 & 0.5 & 0.250000 & 0.666667 & 0.375 & 0.50 \\
1 & B & 1 & 1 & 1 & 5 & 0.5 & 0.166667 & 0.500000 & 0.250 & 0.25 \\
\end{longtable}
}

The output is not the end of the analysis. It is the beginning. A
subgroup with a small sample size may have high uncertainty, so
confidence intervals or bootstrap intervals are often needed.

\section{25.7 Interactive: threshold fairness
explorer}\label{interactive-threshold-fairness-explorer}

The following interactive illustrates how a single threshold changes
accuracy, selection rate, TPR, FPR, PPV, and fairness gaps across two
groups. The scores are synthetic. The point is to understand the
geometry of thresholding, not to claim a universal fairness rule.

Thresholding, error rates, and group gaps

Move the threshold and compare group-wise rates. The model score is a
one-dimensional summary \(S\in[0,1]\); the deployed decision is
\(\widehat{Y}=1\{S\ge t\}\).

Decision threshold t:
\protect\phantomsection\label{ch25-threshold-label}{0.50} Group B score
shift: \protect\phantomsection\label{ch25-shift-label}{-0.05} Show
metric: TPR gap FPR gap selection-rate gap PPV gap

\protect\phantomsection\label{ch25-status}
Loading local Plotly\ldots{}

\protect\phantomsection\label{ch25-score-plot}

\section{25.8 Python example: threshold metrics by
group}\label{python-example-threshold-metrics-by-group}

The interactive above uses the same mathematical idea as the following
Python simulation. We generate scores \(S\), choose a threshold \(t\),
and compute group-wise rates.

\begin{Shaded}
\begin{Highlighting}[]
\ImportTok{import}\NormalTok{ numpy }\ImportTok{as}\NormalTok{ np}
\ImportTok{import}\NormalTok{ pandas }\ImportTok{as}\NormalTok{ pd}

\NormalTok{rng }\OperatorTok{=}\NormalTok{ np.random.default\_rng(}\DecValTok{2025}\NormalTok{)}

\NormalTok{nA }\OperatorTok{=} \DecValTok{400}
\NormalTok{nB }\OperatorTok{=} \DecValTok{400}
\NormalTok{YA }\OperatorTok{=}\NormalTok{ rng.binomial(}\DecValTok{1}\NormalTok{, }\FloatTok{0.45}\NormalTok{, size}\OperatorTok{=}\NormalTok{nA)}
\NormalTok{YB }\OperatorTok{=}\NormalTok{ rng.binomial(}\DecValTok{1}\NormalTok{, }\FloatTok{0.30}\NormalTok{, size}\OperatorTok{=}\NormalTok{nB)}

\KeywordTok{def}\NormalTok{ clipped\_normal(mean, sd, size):}
    \ControlFlowTok{return}\NormalTok{ np.clip(rng.normal(mean, sd, size}\OperatorTok{=}\NormalTok{size), }\FloatTok{0.01}\NormalTok{, }\FloatTok{0.99}\NormalTok{)}

\NormalTok{SA }\OperatorTok{=}\NormalTok{ np.where(YA }\OperatorTok{==} \DecValTok{1}\NormalTok{, clipped\_normal(}\FloatTok{0.70}\NormalTok{, }\FloatTok{0.14}\NormalTok{, nA), clipped\_normal(}\FloatTok{0.32}\NormalTok{, }\FloatTok{0.15}\NormalTok{, nA))}
\NormalTok{SB }\OperatorTok{=}\NormalTok{ np.where(YB }\OperatorTok{==} \DecValTok{1}\NormalTok{, clipped\_normal(}\FloatTok{0.64}\NormalTok{, }\FloatTok{0.15}\NormalTok{, nB), clipped\_normal(}\FloatTok{0.30}\NormalTok{, }\FloatTok{0.15}\NormalTok{, nB))}

\NormalTok{scores }\OperatorTok{=}\NormalTok{ pd.DataFrame(\{}
    \StringTok{"group"}\NormalTok{: np.r\_[np.repeat(}\StringTok{"A"}\NormalTok{, nA), np.repeat(}\StringTok{"B"}\NormalTok{, nB)],}
    \StringTok{"y"}\NormalTok{: np.r\_[YA, YB],}
    \StringTok{"score"}\NormalTok{: np.r\_[SA, SB]}
\NormalTok{\})}

\KeywordTok{def}\NormalTok{ threshold\_report(df, threshold):}
\NormalTok{    out }\OperatorTok{=}\NormalTok{ []}
    \ControlFlowTok{for}\NormalTok{ g, sub }\KeywordTok{in}\NormalTok{ df.groupby(}\StringTok{"group"}\NormalTok{):}
\NormalTok{        y }\OperatorTok{=}\NormalTok{ sub[}\StringTok{"y"}\NormalTok{].to\_numpy()}
\NormalTok{        yhat }\OperatorTok{=}\NormalTok{ (sub[}\StringTok{"score"}\NormalTok{].to\_numpy() }\OperatorTok{\textgreater{}=}\NormalTok{ threshold).astype(}\BuiltInTok{int}\NormalTok{)}
\NormalTok{        tp }\OperatorTok{=}\NormalTok{ ((yhat }\OperatorTok{==} \DecValTok{1}\NormalTok{) }\OperatorTok{\&}\NormalTok{ (y }\OperatorTok{==} \DecValTok{1}\NormalTok{)).}\BuiltInTok{sum}\NormalTok{()}
\NormalTok{        fp }\OperatorTok{=}\NormalTok{ ((yhat }\OperatorTok{==} \DecValTok{1}\NormalTok{) }\OperatorTok{\&}\NormalTok{ (y }\OperatorTok{==} \DecValTok{0}\NormalTok{)).}\BuiltInTok{sum}\NormalTok{()}
\NormalTok{        fn }\OperatorTok{=}\NormalTok{ ((yhat }\OperatorTok{==} \DecValTok{0}\NormalTok{) }\OperatorTok{\&}\NormalTok{ (y }\OperatorTok{==} \DecValTok{1}\NormalTok{)).}\BuiltInTok{sum}\NormalTok{()}
\NormalTok{        tn }\OperatorTok{=}\NormalTok{ ((yhat }\OperatorTok{==} \DecValTok{0}\NormalTok{) }\OperatorTok{\&}\NormalTok{ (y }\OperatorTok{==} \DecValTok{0}\NormalTok{)).}\BuiltInTok{sum}\NormalTok{()}
\NormalTok{        out.append(\{}
            \StringTok{"group"}\NormalTok{: g,}
            \StringTok{"base\_rate"}\NormalTok{: y.mean(),}
            \StringTok{"selection\_rate"}\NormalTok{: yhat.mean(),}
            \StringTok{"TPR"}\NormalTok{: tp }\OperatorTok{/}\NormalTok{ (tp }\OperatorTok{+}\NormalTok{ fn),}
            \StringTok{"FPR"}\NormalTok{: fp }\OperatorTok{/}\NormalTok{ (fp }\OperatorTok{+}\NormalTok{ tn),}
            \StringTok{"PPV"}\NormalTok{: tp }\OperatorTok{/}\NormalTok{ (tp }\OperatorTok{+}\NormalTok{ fp) }\ControlFlowTok{if}\NormalTok{ tp }\OperatorTok{+}\NormalTok{ fp }\OperatorTok{\textgreater{}} \DecValTok{0} \ControlFlowTok{else}\NormalTok{ np.nan,}
            \StringTok{"accuracy"}\NormalTok{: (tp }\OperatorTok{+}\NormalTok{ tn) }\OperatorTok{/} \BuiltInTok{len}\NormalTok{(sub)}
\NormalTok{        \})}
    \ControlFlowTok{return}\NormalTok{ pd.DataFrame(out)}

\NormalTok{threshold\_report(scores, threshold}\OperatorTok{=}\FloatTok{0.50}\NormalTok{)}
\end{Highlighting}
\end{Shaded}

{\def\LTcaptype{none} % do not increment counter
\begin{longtable}[]{@{}llllllll@{}}
\toprule\noalign{}
& group & base\_rate & selection\_rate & TPR & FPR & PPV & accuracy \\
\midrule\noalign{}
\endhead
\bottomrule\noalign{}
\endlastfoot
0 & A & 0.455 & 0.4725 & 0.934066 & 0.087156 & 0.899471 & 0.9225 \\
1 & B & 0.305 & 0.2850 & 0.803279 & 0.057554 & 0.859649 & 0.9000 \\
\end{longtable}
}

A useful classroom activity is to search over thresholds and compare the
threshold that maximizes accuracy with the threshold that reduces a
chosen fairness gap.

\begin{Shaded}
\begin{Highlighting}[]
\NormalTok{thresholds }\OperatorTok{=}\NormalTok{ np.linspace(}\FloatTok{0.05}\NormalTok{, }\FloatTok{0.95}\NormalTok{, }\DecValTok{91}\NormalTok{)}
\NormalTok{rows }\OperatorTok{=}\NormalTok{ []}

\ControlFlowTok{for}\NormalTok{ t }\KeywordTok{in}\NormalTok{ thresholds:}
\NormalTok{    rep }\OperatorTok{=}\NormalTok{ threshold\_report(scores, t).set\_index(}\StringTok{"group"}\NormalTok{)}
\NormalTok{    rows.append(\{}
        \StringTok{"threshold"}\NormalTok{: t,}
        \StringTok{"mean\_accuracy"}\NormalTok{: rep[}\StringTok{"accuracy"}\NormalTok{].mean(),}
        \StringTok{"TPR\_gap"}\NormalTok{: }\BuiltInTok{abs}\NormalTok{(rep.loc[}\StringTok{"A"}\NormalTok{, }\StringTok{"TPR"}\NormalTok{] }\OperatorTok{{-}}\NormalTok{ rep.loc[}\StringTok{"B"}\NormalTok{, }\StringTok{"TPR"}\NormalTok{]),}
        \StringTok{"FPR\_gap"}\NormalTok{: }\BuiltInTok{abs}\NormalTok{(rep.loc[}\StringTok{"A"}\NormalTok{, }\StringTok{"FPR"}\NormalTok{] }\OperatorTok{{-}}\NormalTok{ rep.loc[}\StringTok{"B"}\NormalTok{, }\StringTok{"FPR"}\NormalTok{]),}
        \StringTok{"selection\_gap"}\NormalTok{: }\BuiltInTok{abs}\NormalTok{(rep.loc[}\StringTok{"A"}\NormalTok{, }\StringTok{"selection\_rate"}\NormalTok{] }\OperatorTok{{-}}\NormalTok{ rep.loc[}\StringTok{"B"}\NormalTok{, }\StringTok{"selection\_rate"}\NormalTok{]),}
        \StringTok{"PPV\_gap"}\NormalTok{: }\BuiltInTok{abs}\NormalTok{(rep.loc[}\StringTok{"A"}\NormalTok{, }\StringTok{"PPV"}\NormalTok{] }\OperatorTok{{-}}\NormalTok{ rep.loc[}\StringTok{"B"}\NormalTok{, }\StringTok{"PPV"}\NormalTok{])}
\NormalTok{    \})}

\NormalTok{scan }\OperatorTok{=}\NormalTok{ pd.DataFrame(rows)}
\NormalTok{scan.sort\_values(}\StringTok{"mean\_accuracy"}\NormalTok{, ascending}\OperatorTok{=}\VariableTok{False}\NormalTok{).head()}
\end{Highlighting}
\end{Shaded}

{\def\LTcaptype{none} % do not increment counter
\begin{longtable}[]{@{}lllllll@{}}
\toprule\noalign{}
& threshold & mean\_accuracy & TPR\_gap & FPR\_gap & selection\_gap &
PPV\_gap \\
\midrule\noalign{}
\endhead
\bottomrule\noalign{}
\endlastfoot
44 & 0.49 & 0.91625 & 0.122681 & 0.022408 & 0.1825 & 0.052302 \\
45 & 0.50 & 0.91125 & 0.130787 & 0.029602 & 0.1875 & 0.039822 \\
42 & 0.47 & 0.90875 & 0.081787 & 0.022144 & 0.1675 & 0.071287 \\
46 & 0.51 & 0.90875 & 0.127995 & 0.033199 & 0.1875 & 0.032324 \\
43 & 0.48 & 0.90750 & 0.119978 & 0.022771 & 0.1800 & 0.067249 \\
\end{longtable}
}

\begin{Shaded}
\begin{Highlighting}[]
\NormalTok{scan.sort\_values(}\StringTok{"TPR\_gap"}\NormalTok{).head()}
\end{Highlighting}
\end{Shaded}

{\def\LTcaptype{none} % do not increment counter
\begin{longtable}[]{@{}lllllll@{}}
\toprule\noalign{}
& threshold & mean\_accuracy & TPR\_gap & FPR\_gap & selection\_gap &
PPV\_gap \\
\midrule\noalign{}
\endhead
\bottomrule\noalign{}
\endlastfoot
0 & 0.05 & 0.41375 & 0.0 & 0.056201 & 0.0425 & 0.138008 \\
12 & 0.17 & 0.48250 & 0.0 & 0.065804 & 0.0650 & 0.136646 \\
10 & 0.15 & 0.46625 & 0.0 & 0.051779 & 0.0525 & 0.140381 \\
9 & 0.14 & 0.46125 & 0.0 & 0.053759 & 0.0525 & 0.139745 \\
8 & 0.13 & 0.45000 & 0.0 & 0.037753 & 0.0400 & 0.143783 \\
\end{longtable}
}

These two tables usually point to different thresholds. This is a
mathematical signal that the design decision is not purely technical.

\section{25.9 Calibration and
uncertainty}\label{calibration-and-uncertainty}

A model score \(S\in[0,1]\) is calibrated if

\[
P(Y=1\mid S=s)=s.
\]

In practice we approximate this condition using bins. If bin \(B_m\)
contains examples whose scores lie in an interval, define

\[
\operatorname{acc}(B_m)=\frac{1}{|B_m|}\sum_{i\in B_m} y_i,
\qquad
\operatorname{conf}(B_m)=\frac{1}{|B_m|}\sum_{i\in B_m} s_i.
\]

The expected calibration error is commonly estimated by

\[
\operatorname{ECE}=\sum_{m=1}^M \frac{|B_m|}{n}\left|\operatorname{acc}(B_m)-\operatorname{conf}(B_m)\right|.
\]

Calibration matters because a score of \(0.9\) should not merely rank
higher than a score of \(0.6\); it should mean something probabilistic.
In NLP, calibration is important for moderation, retrieval confidence,
question-answering uncertainty, and decision support.

\begin{Shaded}
\begin{Highlighting}[]
\KeywordTok{def}\NormalTok{ calibration\_table(y, s, bins}\OperatorTok{=}\DecValTok{8}\NormalTok{):}
\NormalTok{    y }\OperatorTok{=}\NormalTok{ np.asarray(y)}
\NormalTok{    s }\OperatorTok{=}\NormalTok{ np.asarray(s)}
\NormalTok{    edges }\OperatorTok{=}\NormalTok{ np.linspace(}\DecValTok{0}\NormalTok{, }\DecValTok{1}\NormalTok{, bins }\OperatorTok{+} \DecValTok{1}\NormalTok{)}
\NormalTok{    rows }\OperatorTok{=}\NormalTok{ []}
\NormalTok{    ece }\OperatorTok{=} \FloatTok{0.0}
    \ControlFlowTok{for}\NormalTok{ j }\KeywordTok{in} \BuiltInTok{range}\NormalTok{(bins):}
\NormalTok{        left, right }\OperatorTok{=}\NormalTok{ edges[j], edges[j }\OperatorTok{+} \DecValTok{1}\NormalTok{]}
        \ControlFlowTok{if}\NormalTok{ j }\OperatorTok{==}\NormalTok{ bins }\OperatorTok{{-}} \DecValTok{1}\NormalTok{:}
\NormalTok{            mask }\OperatorTok{=}\NormalTok{ (s }\OperatorTok{\textgreater{}=}\NormalTok{ left) }\OperatorTok{\&}\NormalTok{ (s }\OperatorTok{\textless{}=}\NormalTok{ right)}
        \ControlFlowTok{else}\NormalTok{:}
\NormalTok{            mask }\OperatorTok{=}\NormalTok{ (s }\OperatorTok{\textgreater{}=}\NormalTok{ left) }\OperatorTok{\&}\NormalTok{ (s }\OperatorTok{\textless{}}\NormalTok{ right)}
        \ControlFlowTok{if}\NormalTok{ mask.}\BuiltInTok{sum}\NormalTok{() }\OperatorTok{==} \DecValTok{0}\NormalTok{:}
            \ControlFlowTok{continue}
\NormalTok{        acc }\OperatorTok{=}\NormalTok{ y[mask].mean()}
\NormalTok{        conf }\OperatorTok{=}\NormalTok{ s[mask].mean()}
\NormalTok{        weight }\OperatorTok{=}\NormalTok{ mask.mean()}
\NormalTok{        ece }\OperatorTok{+=}\NormalTok{ weight }\OperatorTok{*} \BuiltInTok{abs}\NormalTok{(acc }\OperatorTok{{-}}\NormalTok{ conf)}
\NormalTok{        rows.append(\{}
            \StringTok{"bin"}\NormalTok{: }\SpecialStringTok{f"[}\SpecialCharTok{\{}\NormalTok{left}\SpecialCharTok{:.2f\}}\SpecialStringTok{, }\SpecialCharTok{\{}\NormalTok{right}\SpecialCharTok{:.2f\}}\SpecialStringTok{]"}\NormalTok{,}
            \StringTok{"count"}\NormalTok{: }\BuiltInTok{int}\NormalTok{(mask.}\BuiltInTok{sum}\NormalTok{()),}
            \StringTok{"accuracy"}\NormalTok{: acc,}
            \StringTok{"mean\_score"}\NormalTok{: conf,}
            \StringTok{"abs\_gap"}\NormalTok{: }\BuiltInTok{abs}\NormalTok{(acc }\OperatorTok{{-}}\NormalTok{ conf)}
\NormalTok{        \})}
    \ControlFlowTok{return}\NormalTok{ pd.DataFrame(rows), ece}

\NormalTok{cal\_table, ece }\OperatorTok{=}\NormalTok{ calibration\_table(scores[}\StringTok{"y"}\NormalTok{], scores[}\StringTok{"score"}\NormalTok{], bins}\OperatorTok{=}\DecValTok{8}\NormalTok{)}
\NormalTok{cal\_table, ece}
\end{Highlighting}
\end{Shaded}

\begin{verbatim}
(            bin  count  accuracy  mean_score   abs_gap
 0  [0.00, 0.12]     54  0.000000    0.057371  0.057371
 1  [0.12, 0.25]    133  0.015038    0.199292  0.184255
 2  [0.25, 0.38]    173  0.023121    0.316975  0.293854
 3  [0.38, 0.50]    137  0.218978    0.433008  0.214030
 4  [0.50, 0.62]    116  0.750000    0.567332  0.182668
 5  [0.62, 0.75]    108  0.953704    0.682537  0.271166
 6  [0.75, 0.88]     58  0.982759    0.796844  0.185914
 7  [0.88, 1.00]     21  1.000000    0.939546  0.060454,
 0.21286359605060212)
\end{verbatim}

A classifier can have high accuracy and poor calibration. Conversely, a
calibrated classifier may not have the best thresholded accuracy. These
are different mathematical properties.

\section{25.10 Distribution shift}\label{distribution-shift}

Let \(P\) be the training distribution and \(Q\) be the deployment
distribution. Distribution shift means

\[
P(X,A,Y)\ne Q(X,A,Y).
\]

Important special cases include:

\begin{itemize}
\tightlist
\item
  \textbf{Covariate shift:} \(P(X)\ne Q(X)\) but
  \(P(Y\mid X)=Q(Y\mid X)\).
\item
  \textbf{Label shift:} \(P(Y)\ne Q(Y)\) but \(P(X\mid Y)=Q(X\mid Y)\).
\item
  \textbf{Concept shift:} \(P(Y\mid X)\ne Q(Y\mid X)\).
\end{itemize}

In text applications, distribution shift appears when a model trained on
one source is used on another source: social media versus formal essays,
one university's policy pages versus another's, one time period versus
another, or one language variety versus another.

Under covariate shift, if the density ratio \(w(x)=q(x)/p(x)\) is known
or estimated, then

\[
R_Q(d)=E_Q[L(d(X),Y)]
=E_P\left[\frac{q(X)}{p(X)}L(d(X),Y)\right].
\]

Thus weighted empirical risk can estimate deployment risk:

\[
\widehat{R}_Q(d)=\frac{1}{n}\sum_{i=1}^n w(x_i)L(d(x_i),y_i).
\]

\begin{Shaded}
\begin{Highlighting}[]
\CommentTok{\# A toy covariate{-}shift example with two text sources.}
\NormalTok{shift\_df }\OperatorTok{=}\NormalTok{ pd.DataFrame(\{}
    \StringTok{"source"}\NormalTok{: [}\StringTok{"train"}\NormalTok{] }\OperatorTok{*} \DecValTok{6} \OperatorTok{+}\NormalTok{ [}\StringTok{"deploy"}\NormalTok{] }\OperatorTok{*} \DecValTok{6}\NormalTok{,}
    \StringTok{"text\_type"}\NormalTok{: [}\StringTok{"short"}\NormalTok{, }\StringTok{"short"}\NormalTok{, }\StringTok{"short"}\NormalTok{, }\StringTok{"long"}\NormalTok{, }\StringTok{"long"}\NormalTok{, }\StringTok{"long"}\NormalTok{, }\StringTok{"short"}\NormalTok{, }\StringTok{"short"}\NormalTok{, }\StringTok{"long"}\NormalTok{, }\StringTok{"long"}\NormalTok{, }\StringTok{"long"}\NormalTok{, }\StringTok{"long"}\NormalTok{],}
    \StringTok{"loss"}\NormalTok{: [}\DecValTok{0}\NormalTok{, }\DecValTok{1}\NormalTok{, }\DecValTok{0}\NormalTok{, }\DecValTok{0}\NormalTok{, }\DecValTok{0}\NormalTok{, }\DecValTok{1}\NormalTok{, }\DecValTok{0}\NormalTok{, }\DecValTok{1}\NormalTok{, }\DecValTok{1}\NormalTok{, }\DecValTok{0}\NormalTok{, }\DecValTok{1}\NormalTok{, }\DecValTok{1}\NormalTok{]}
\NormalTok{\})}

\NormalTok{train }\OperatorTok{=}\NormalTok{ shift\_df[shift\_df[}\StringTok{"source"}\NormalTok{] }\OperatorTok{==} \StringTok{"train"}\NormalTok{]}
\NormalTok{deploy }\OperatorTok{=}\NormalTok{ shift\_df[shift\_df[}\StringTok{"source"}\NormalTok{] }\OperatorTok{==} \StringTok{"deploy"}\NormalTok{]}

\NormalTok{p\_train }\OperatorTok{=}\NormalTok{ train[}\StringTok{"text\_type"}\NormalTok{].value\_counts(normalize}\OperatorTok{=}\VariableTok{True}\NormalTok{)}
\NormalTok{p\_deploy }\OperatorTok{=}\NormalTok{ deploy[}\StringTok{"text\_type"}\NormalTok{].value\_counts(normalize}\OperatorTok{=}\VariableTok{True}\NormalTok{)}

\NormalTok{weights }\OperatorTok{=}\NormalTok{ train[}\StringTok{"text\_type"}\NormalTok{].}\BuiltInTok{map}\NormalTok{(}\KeywordTok{lambda}\NormalTok{ z: p\_deploy[z] }\OperatorTok{/}\NormalTok{ p\_train[z])}
\NormalTok{unweighted\_risk }\OperatorTok{=}\NormalTok{ train[}\StringTok{"loss"}\NormalTok{].mean()}
\NormalTok{weighted\_risk }\OperatorTok{=}\NormalTok{ np.average(train[}\StringTok{"loss"}\NormalTok{], weights}\OperatorTok{=}\NormalTok{weights)}
\NormalTok{deploy\_risk }\OperatorTok{=}\NormalTok{ deploy[}\StringTok{"loss"}\NormalTok{].mean()}

\NormalTok{pd.DataFrame(\{}
    \StringTok{"estimate"}\NormalTok{: [}\StringTok{"unweighted train risk"}\NormalTok{, }\StringTok{"weighted train risk"}\NormalTok{, }\StringTok{"observed deploy risk"}\NormalTok{],}
    \StringTok{"risk"}\NormalTok{: [unweighted\_risk, weighted\_risk, deploy\_risk]}
\NormalTok{\})}
\end{Highlighting}
\end{Shaded}

{\def\LTcaptype{none} % do not increment counter
\begin{longtable}[]{@{}lll@{}}
\toprule\noalign{}
& estimate & risk \\
\midrule\noalign{}
\endhead
\bottomrule\noalign{}
\endlastfoot
0 & unweighted train risk & 0.333333 \\
1 & weighted train risk & 0.333333 \\
2 & observed deploy risk & 0.666667 \\
\end{longtable}
}

The weighted estimate moves toward the deployment risk because it
accounts for the changed mixture of text types. In real NLP, estimating
\(q(x)/p(x)\) is itself a hard statistical problem.

\section{25.11 Embedding association and representation
bias}\label{embedding-association-and-representation-bias}

Embedding models represent words, sentences, or documents as vectors.
Bias can appear as geometric association. Let \(X\) and \(Y\) be two
sets of target words, and let \(A\) and \(B\) be two sets of attribute
words. A simple association score is

\[
s(w,A,B)=\frac{1}{|A|}\sum_{a\in A}\cos(e_w,e_a)
-
\frac{1}{|B|}\sum_{b\in B}\cos(e_w,e_b).
\]

A set-level effect score is

\[
S(X,Y,A,B)=\sum_{x\in X}s(x,A,B)-\sum_{y\in Y}s(y,A,B).
\]

This type of score is inspired by word-embedding association tests and
related representation-auditing methods Caliskan et al. (2017),
Bolukbasi, Chang, Zou, Saligrama, and Adam Tauman Kalai (2016).

\begin{Shaded}
\begin{Highlighting}[]
\CommentTok{\# A small toy embedding audit.  These numbers are artificial and for pedagogy only.}
\NormalTok{emb }\OperatorTok{=}\NormalTok{ \{}
    \StringTok{"nurse"}\NormalTok{: np.array([}\FloatTok{0.10}\NormalTok{, }\FloatTok{0.80}\NormalTok{]),}
    \StringTok{"engineer"}\NormalTok{: np.array([}\FloatTok{0.80}\NormalTok{, }\FloatTok{0.10}\NormalTok{]),}
    \StringTok{"teacher"}\NormalTok{: np.array([}\FloatTok{0.20}\NormalTok{, }\FloatTok{0.70}\NormalTok{]),}
    \StringTok{"programmer"}\NormalTok{: np.array([}\FloatTok{0.75}\NormalTok{, }\FloatTok{0.20}\NormalTok{]),}
    \StringTok{"care"}\NormalTok{: np.array([}\FloatTok{0.00}\NormalTok{, }\FloatTok{1.00}\NormalTok{]),}
    \StringTok{"family"}\NormalTok{: np.array([}\FloatTok{0.10}\NormalTok{, }\FloatTok{0.95}\NormalTok{]),}
    \StringTok{"logic"}\NormalTok{: np.array([}\FloatTok{1.00}\NormalTok{, }\FloatTok{0.00}\NormalTok{]),}
    \StringTok{"machine"}\NormalTok{: np.array([}\FloatTok{0.95}\NormalTok{, }\FloatTok{0.10}\NormalTok{]),}
\NormalTok{\}}

\KeywordTok{def}\NormalTok{ cosine(u, v):}
\NormalTok{    u }\OperatorTok{=}\NormalTok{ np.asarray(u)}
\NormalTok{    v }\OperatorTok{=}\NormalTok{ np.asarray(v)}
    \ControlFlowTok{return} \BuiltInTok{float}\NormalTok{(u }\OperatorTok{@}\NormalTok{ v }\OperatorTok{/}\NormalTok{ (np.linalg.norm(u) }\OperatorTok{*}\NormalTok{ np.linalg.norm(v)))}

\KeywordTok{def}\NormalTok{ assoc(word, A, B):}
    \ControlFlowTok{return}\NormalTok{ np.mean([cosine(emb[word], emb[a]) }\ControlFlowTok{for}\NormalTok{ a }\KeywordTok{in}\NormalTok{ A]) }\OperatorTok{{-}}\NormalTok{ np.mean([cosine(emb[word], emb[b]) }\ControlFlowTok{for}\NormalTok{ b }\KeywordTok{in}\NormalTok{ B])}

\NormalTok{X }\OperatorTok{=}\NormalTok{ [}\StringTok{"nurse"}\NormalTok{, }\StringTok{"teacher"}\NormalTok{]}
\NormalTok{Y }\OperatorTok{=}\NormalTok{ [}\StringTok{"engineer"}\NormalTok{, }\StringTok{"programmer"}\NormalTok{]}
\NormalTok{A }\OperatorTok{=}\NormalTok{ [}\StringTok{"care"}\NormalTok{, }\StringTok{"family"}\NormalTok{]}
\NormalTok{B }\OperatorTok{=}\NormalTok{ [}\StringTok{"logic"}\NormalTok{, }\StringTok{"machine"}\NormalTok{]}

\NormalTok{pd.DataFrame(\{}
    \StringTok{"word"}\NormalTok{: X }\OperatorTok{+}\NormalTok{ Y,}
    \StringTok{"association\_to\_A\_minus\_B"}\NormalTok{: [assoc(w, A, B) }\ControlFlowTok{for}\NormalTok{ w }\KeywordTok{in}\NormalTok{ X }\OperatorTok{+}\NormalTok{ Y]}
\NormalTok{\})}
\end{Highlighting}
\end{Shaded}

{\def\LTcaptype{none} % do not increment counter
\begin{longtable}[]{@{}lll@{}}
\toprule\noalign{}
& word & association\_to\_A\_minus\_B \\
\midrule\noalign{}
\endhead
\bottomrule\noalign{}
\endlastfoot
0 & nurse & 0.820412 \\
1 & teacher & 0.648967 \\
2 & engineer & -0.820412 \\
3 & programmer & -0.669537 \\
\end{longtable}
}

A linear debiasing operation can be written as orthogonal projection. If
\(b\) is a unit direction representing an unwanted association, then the
projection removing that direction is

\[
\Pi_{b^\perp}(v)=v-(v^Tb)b.
\]

This operation is mathematically clean, but it is not a complete
solution. Bias may be nonlinear, contextual, or encoded in many
directions. Projection can also remove useful information or create a
false sense of safety.

\begin{Shaded}
\begin{Highlighting}[]
\CommentTok{\# Remove the direction from "logic/machine" to "care/family" in the toy example.}
\NormalTok{mean\_A }\OperatorTok{=}\NormalTok{ np.mean([emb[a] }\ControlFlowTok{for}\NormalTok{ a }\KeywordTok{in}\NormalTok{ A], axis}\OperatorTok{=}\DecValTok{0}\NormalTok{)}
\NormalTok{mean\_B }\OperatorTok{=}\NormalTok{ np.mean([emb[b] }\ControlFlowTok{for}\NormalTok{ b }\KeywordTok{in}\NormalTok{ B], axis}\OperatorTok{=}\DecValTok{0}\NormalTok{)}
\NormalTok{bias\_direction }\OperatorTok{=}\NormalTok{ mean\_A }\OperatorTok{{-}}\NormalTok{ mean\_B}
\NormalTok{bias\_direction }\OperatorTok{=}\NormalTok{ bias\_direction }\OperatorTok{/}\NormalTok{ np.linalg.norm(bias\_direction)}

\KeywordTok{def}\NormalTok{ remove\_direction(v, b):}
    \ControlFlowTok{return}\NormalTok{ v }\OperatorTok{{-}}\NormalTok{ (v }\OperatorTok{@}\NormalTok{ b) }\OperatorTok{*}\NormalTok{ b}

\NormalTok{emb\_projected }\OperatorTok{=}\NormalTok{ \{w: remove\_direction(v, bias\_direction) }\ControlFlowTok{for}\NormalTok{ w, v }\KeywordTok{in}\NormalTok{ emb.items()\}}

\KeywordTok{def}\NormalTok{ assoc\_projected(word, A, B):}
    \ControlFlowTok{return}\NormalTok{ np.mean([cosine(emb\_projected[word], emb\_projected[a]) }\ControlFlowTok{for}\NormalTok{ a }\KeywordTok{in}\NormalTok{ A]) }\OperatorTok{{-}}\NormalTok{ np.mean([cosine(emb\_projected[word], emb\_projected[b]) }\ControlFlowTok{for}\NormalTok{ b }\KeywordTok{in}\NormalTok{ B])}

\NormalTok{pd.DataFrame(\{}
    \StringTok{"word"}\NormalTok{: X }\OperatorTok{+}\NormalTok{ Y,}
    \StringTok{"before"}\NormalTok{: [assoc(w, A, B) }\ControlFlowTok{for}\NormalTok{ w }\KeywordTok{in}\NormalTok{ X }\OperatorTok{+}\NormalTok{ Y],}
    \StringTok{"after\_projection"}\NormalTok{: [assoc\_projected(w, A, B) }\ControlFlowTok{for}\NormalTok{ w }\KeywordTok{in}\NormalTok{ X }\OperatorTok{+}\NormalTok{ Y]}
\NormalTok{\})}
\end{Highlighting}
\end{Shaded}

{\def\LTcaptype{none} % do not increment counter
\begin{longtable}[]{@{}llll@{}}
\toprule\noalign{}
& word & before & after\_projection \\
\midrule\noalign{}
\endhead
\bottomrule\noalign{}
\endlastfoot
0 & nurse & 0.820412 & 0.000000e+00 \\
1 & teacher & 0.648967 & 1.110223e-16 \\
2 & engineer & -0.820412 & 0.000000e+00 \\
3 & programmer & -0.669537 & 0.000000e+00 \\
\end{longtable}
}

\section{25.12 Selective prediction and
abstention}\label{selective-prediction-and-abstention}

Responsible deployment often requires a model not to answer every query.
A selective classifier returns either a prediction or an abstention
symbol \(\bot\):

\[
g(x)\in\{0,1,\bot\}.
\]

For a confidence score \(c(x)\) and threshold \(\tau\), a simple rule is

\[
g_\tau(x)=
\begin{cases}
\widehat{y}(x), & c(x)\ge \tau,\\
\bot, & c(x)<\tau.
\end{cases}
\]

Two quantities are central:

\[
\operatorname{coverage}(\tau)=P(c(X)\ge \tau),
\]

\[
\operatorname{selective\ risk}(\tau)
=E[L(\widehat{Y},Y)\mid c(X)\ge \tau].
\]

Increasing \(\tau\) often decreases risk among answered cases but also
decreases coverage.

\begin{Shaded}
\begin{Highlighting}[]
\NormalTok{conf }\OperatorTok{=}\NormalTok{ np.maximum(scores[}\StringTok{"score"}\NormalTok{].to\_numpy(), }\DecValTok{1} \OperatorTok{{-}}\NormalTok{ scores[}\StringTok{"score"}\NormalTok{].to\_numpy())}
\NormalTok{y }\OperatorTok{=}\NormalTok{ scores[}\StringTok{"y"}\NormalTok{].to\_numpy()}
\NormalTok{yhat }\OperatorTok{=}\NormalTok{ (scores[}\StringTok{"score"}\NormalTok{].to\_numpy() }\OperatorTok{\textgreater{}=} \FloatTok{0.50}\NormalTok{).astype(}\BuiltInTok{int}\NormalTok{)}
\NormalTok{loss }\OperatorTok{=}\NormalTok{ (yhat }\OperatorTok{!=}\NormalTok{ y).astype(}\BuiltInTok{float}\NormalTok{)}

\NormalTok{selective\_rows }\OperatorTok{=}\NormalTok{ []}
\ControlFlowTok{for}\NormalTok{ tau }\KeywordTok{in}\NormalTok{ np.linspace(}\FloatTok{0.50}\NormalTok{, }\FloatTok{0.95}\NormalTok{, }\DecValTok{10}\NormalTok{):}
\NormalTok{    keep }\OperatorTok{=}\NormalTok{ conf }\OperatorTok{\textgreater{}=}\NormalTok{ tau}
\NormalTok{    selective\_rows.append(\{}
        \StringTok{"confidence\_threshold"}\NormalTok{: tau,}
        \StringTok{"coverage"}\NormalTok{: keep.mean(),}
        \StringTok{"selective\_error"}\NormalTok{: loss[keep].mean() }\ControlFlowTok{if}\NormalTok{ keep.}\BuiltInTok{sum}\NormalTok{() }\OperatorTok{\textgreater{}} \DecValTok{0} \ControlFlowTok{else}\NormalTok{ np.nan}
\NormalTok{    \})}

\NormalTok{pd.DataFrame(selective\_rows)}
\end{Highlighting}
\end{Shaded}

{\def\LTcaptype{none} % do not increment counter
\begin{longtable}[]{@{}llll@{}}
\toprule\noalign{}
& confidence\_threshold & coverage & selective\_error \\
\midrule\noalign{}
\endhead
\bottomrule\noalign{}
\endlastfoot
0 & 0.50 & 1.00000 & 0.088750 \\
1 & 0.55 & 0.88125 & 0.052482 \\
2 & 0.60 & 0.76000 & 0.026316 \\
3 & 0.65 & 0.60250 & 0.014523 \\
4 & 0.70 & 0.45000 & 0.011111 \\
5 & 0.75 & 0.33250 & 0.011278 \\
6 & 0.80 & 0.20625 & 0.012121 \\
7 & 0.85 & 0.11750 & 0.010638 \\
8 & 0.90 & 0.07000 & 0.000000 \\
9 & 0.95 & 0.04375 & 0.000000 \\
\end{longtable}
}

In high-stakes applications, abstention should usually trigger a defined
alternative process, such as human review, retrieval of source
documents, or refusal to answer.

\section{25.13 Privacy and memorization}\label{privacy-and-memorization}

NLP systems can leak information from training data, prompts, logs, or
retrieved documents. A mathematical privacy guarantee should describe
how the output distribution changes when one individual's data is
changed.

A randomized algorithm \(M\) is \(\varepsilon\)-differentially private
if for all neighboring datasets \(D\) and \(D'\) differing in one
individual, and for all measurable output sets \(S\),

\[
P(M(D)\in S)\le e^\varepsilon P(M(D')\in S).
\]

Smaller \(\varepsilon\) means stronger privacy. Differential privacy is
not the only privacy concept, but it is mathematically precise Dwork and
Roth (2014).

A simple example is randomized response. Suppose a private binary answer
is \(Y\in\{0,1\}\). The released answer \(Z\) equals \(Y\) with
probability \(p\) and is flipped with probability \(1-p\). The privacy
parameter is

\[
\varepsilon=\log\frac{p}{1-p}.
\]

\begin{Shaded}
\begin{Highlighting}[]
\NormalTok{ps }\OperatorTok{=}\NormalTok{ np.array([}\FloatTok{0.55}\NormalTok{, }\FloatTok{0.60}\NormalTok{, }\FloatTok{0.70}\NormalTok{, }\FloatTok{0.80}\NormalTok{, }\FloatTok{0.90}\NormalTok{])}
\NormalTok{eps }\OperatorTok{=}\NormalTok{ np.log(ps }\OperatorTok{/}\NormalTok{ (}\DecValTok{1} \OperatorTok{{-}}\NormalTok{ ps))}
\NormalTok{pd.DataFrame(\{}\StringTok{"truth\_probability\_p"}\NormalTok{: ps, }\StringTok{"epsilon"}\NormalTok{: eps\})}
\end{Highlighting}
\end{Shaded}

{\def\LTcaptype{none} % do not increment counter
\begin{longtable}[]{@{}lll@{}}
\toprule\noalign{}
& truth\_probability\_p & epsilon \\
\midrule\noalign{}
\endhead
\bottomrule\noalign{}
\endlastfoot
0 & 0.55 & 0.200671 \\
1 & 0.60 & 0.405465 \\
2 & 0.70 & 0.847298 \\
3 & 0.80 & 1.386294 \\
4 & 0.90 & 2.197225 \\
\end{longtable}
}

Privacy has a utility tradeoff. As \(p\) approaches \(1/2\), the
released answer becomes more private but less informative. As \(p\)
approaches \(1\), the answer is accurate but less private.

\section{25.14 Hallucination, factuality, and RAG
responsibility}\label{hallucination-factuality-and-rag-responsibility}

For a generated answer \(Z\) to a prompt \(X\), factuality cannot be
determined from fluency alone. A response may be syntactically excellent
and factually false. A clean way to describe a retrieval-augmented
system is

\[
D_1,\ldots,D_k \sim r_\eta(d\mid x),
\qquad
Z\sim p_\theta(z\mid x,D_1,\ldots,D_k).
\]

Errors can enter through retrieval or generation:

\begin{itemize}
\tightlist
\item
  relevant documents may not be retrieved;
\item
  retrieved documents may be outdated, ambiguous, or contradictory;
\item
  the generator may ignore retrieved evidence;
\item
  the generator may blend evidence with unsupported prior knowledge;
\item
  the final answer may omit uncertainty or citations.
\end{itemize}

For mathematical responsibility, it is useful to separate:

\[
P(\text{correct answer})
=P(\text{retrieval sufficient})P(\text{generation faithful}\mid \text{retrieval sufficient})+
\cdots.
\]

This is not a complete causal model, but it forces students to separate
retrieval failure from generation failure.

\section{25.15 AI-assisted auditing
workflow}\label{ai-assisted-auditing-workflow}

AI tools can help audit NLP systems, but they should not replace
statistical reasoning. A useful AI-assisted workflow is:

\begin{enumerate}
\def\labelenumi{\arabic{enumi}.}
\tightlist
\item
  Ask the AI system to state the task, target variable, decision rule,
  and possible harms.
\item
  Ask it to propose subgroup variables that might matter.
\item
  Ask it to list plausible distribution shifts.
\item
  Convert the suggestions into mathematical quantities.
\item
  Compute metrics and uncertainty using data.
\item
  Ask the AI system to critique the audit for missing assumptions.
\item
  Document the final human judgment.
\end{enumerate}

\begin{tcolorbox}[enhanced jigsaw, arc=.35mm, bottomrule=.15mm, bottomtitle=1mm, breakable, colback=white, colbacktitle=quarto-callout-note-color!10!white, colframe=quarto-callout-note-color-frame, coltitle=black, left=2mm, leftrule=.75mm, opacityback=0, opacitybacktitle=0.6, rightrule=.15mm, title=\textcolor{quarto-callout-note-color}{\faInfo}\hspace{0.5em}{Note}, titlerule=0mm, toprule=.15mm, toptitle=1mm]

AI assistance is useful for brainstorming failure modes and drafting
audit checklists. It is not evidence. Evidence comes from data, source
documents, mathematical definitions, and careful evaluation.

\end{tcolorbox}

\section{25.16 Worked example: two classifiers with the same
accuracy}\label{worked-example-two-classifiers-with-the-same-accuracy}

Suppose two classifiers have the same overall accuracy, but their
group-specific errors differ.

\begin{Shaded}
\begin{Highlighting}[]
\NormalTok{comparison }\OperatorTok{=}\NormalTok{ pd.DataFrame(\{}
    \StringTok{"model"}\NormalTok{: [}\StringTok{"M1"}\NormalTok{, }\StringTok{"M1"}\NormalTok{, }\StringTok{"M2"}\NormalTok{, }\StringTok{"M2"}\NormalTok{],}
    \StringTok{"group"}\NormalTok{: [}\StringTok{"A"}\NormalTok{, }\StringTok{"B"}\NormalTok{, }\StringTok{"A"}\NormalTok{, }\StringTok{"B"}\NormalTok{],}
    \StringTok{"n"}\NormalTok{: [}\DecValTok{100}\NormalTok{, }\DecValTok{100}\NormalTok{, }\DecValTok{100}\NormalTok{, }\DecValTok{100}\NormalTok{],}
    \StringTok{"accuracy"}\NormalTok{: [}\FloatTok{0.90}\NormalTok{, }\FloatTok{0.70}\NormalTok{, }\FloatTok{0.80}\NormalTok{, }\FloatTok{0.80}\NormalTok{],}
    \StringTok{"TPR"}\NormalTok{: [}\FloatTok{0.92}\NormalTok{, }\FloatTok{0.58}\NormalTok{, }\FloatTok{0.80}\NormalTok{, }\FloatTok{0.79}\NormalTok{],}
    \StringTok{"FPR"}\NormalTok{: [}\FloatTok{0.10}\NormalTok{, }\FloatTok{0.22}\NormalTok{, }\FloatTok{0.16}\NormalTok{, }\FloatTok{0.15}\NormalTok{]}
\NormalTok{\})}
\NormalTok{comparison}
\end{Highlighting}
\end{Shaded}

{\def\LTcaptype{none} % do not increment counter
\begin{longtable}[]{@{}lllllll@{}}
\toprule\noalign{}
& model & group & n & accuracy & TPR & FPR \\
\midrule\noalign{}
\endhead
\bottomrule\noalign{}
\endlastfoot
0 & M1 & A & 100 & 0.9 & 0.92 & 0.10 \\
1 & M1 & B & 100 & 0.7 & 0.58 & 0.22 \\
2 & M2 & A & 100 & 0.8 & 0.80 & 0.16 \\
3 & M2 & B & 100 & 0.8 & 0.79 & 0.15 \\
\end{longtable}
}

Both models have average accuracy \(0.80\), but Model M1 has a much
larger group performance disparity. Whether M2 is preferable depends on
the application, but the overall metric alone hides the difference.

\begin{Shaded}
\begin{Highlighting}[]
\NormalTok{summary }\OperatorTok{=}\NormalTok{ comparison.groupby(}\StringTok{"model"}\NormalTok{).agg(}
\NormalTok{    mean\_accuracy}\OperatorTok{=}\NormalTok{(}\StringTok{"accuracy"}\NormalTok{, }\StringTok{"mean"}\NormalTok{),}
\NormalTok{    TPR\_gap}\OperatorTok{=}\NormalTok{(}\StringTok{"TPR"}\NormalTok{, }\KeywordTok{lambda}\NormalTok{ z: z.}\BuiltInTok{max}\NormalTok{() }\OperatorTok{{-}}\NormalTok{ z.}\BuiltInTok{min}\NormalTok{()),}
\NormalTok{    FPR\_gap}\OperatorTok{=}\NormalTok{(}\StringTok{"FPR"}\NormalTok{, }\KeywordTok{lambda}\NormalTok{ z: z.}\BuiltInTok{max}\NormalTok{() }\OperatorTok{{-}}\NormalTok{ z.}\BuiltInTok{min}\NormalTok{())}
\NormalTok{)}
\NormalTok{summary}
\end{Highlighting}
\end{Shaded}

{\def\LTcaptype{none} % do not increment counter
\begin{longtable}[]{@{}llll@{}}
\toprule\noalign{}
& mean\_accuracy & TPR\_gap & FPR\_gap \\
model & & & \\
\midrule\noalign{}
\endhead
\bottomrule\noalign{}
\endlastfoot
M1 & 0.8 & 0.34 & 0.12 \\
M2 & 0.8 & 0.01 & 0.01 \\
\end{longtable}
}

The mathematical lesson is that aggregation is a projection. It can
remove information. Responsible evaluation often requires conditional
metrics.

\section{25.17 Common mistakes}\label{common-mistakes-14}

\begin{enumerate}
\def\labelenumi{\arabic{enumi}.}
\tightlist
\item
  \textbf{Treating accuracy as responsibility.} Accuracy is one
  statistic, not an ethical guarantee.
\item
  \textbf{Ignoring the deployment distribution.} A test set estimates
  performance only for the distribution it represents.
\item
  \textbf{Using group metrics without uncertainty.} Small subgroups can
  produce unstable estimates.
\item
  \textbf{Assuming one fairness metric is always correct.} Different
  metrics encode different choices.
\item
  \textbf{Assuming debiasing projection solves bias.} Geometric
  correction can help but is incomplete.
\item
  \textbf{Confusing fluency with factuality.} Language quality does not
  imply truth.
\item
  \textbf{Using AI critiques as evidence.} AI-generated audits must be
  checked against data and definitions.
\end{enumerate}

\section{25.18 Exercises}\label{exercises-24}

\begin{enumerate}
\def\labelenumi{\arabic{enumi}.}
\item
  \textbf{Group metrics.} Given the following counts for groups A and B,
  compute TPR, FPR, PPV, selection rate, and accuracy. Which fairness
  criteria are approximately satisfied?
\item
  \textbf{Threshold tradeoff.} For a score-based classifier, write code
  that searches over thresholds and plots TPR gap, FPR gap,
  selection-rate gap, and accuracy.
\item
  \textbf{Calibration.} Construct a model that has high accuracy but
  poor calibration. Compute ECE.
\item
  \textbf{Distribution shift.} Suppose the training data contains
  \(70\%\) formal text and \(30\%\) informal text, but deployment
  contains \(30\%\) formal text and \(70\%\) informal text. If error
  rates differ by text type, compute the shifted risk.
\item
  \textbf{Embedding association.} Build a small toy embedding example
  and compute association scores before and after projection.
\item
  \textbf{Selective prediction.} For a classifier with confidence
  scores, compute coverage and selective risk as functions of the
  abstention threshold.
\item
  \textbf{Privacy.} For randomized response with \(p=0.75\), compute
  \(\varepsilon\). Interpret the result.
\item
  \textbf{RAG audit.} Design separate metrics for retrieval sufficiency
  and generation faithfulness.
\end{enumerate}

\section{25.19 Mini-project: mathematical audit of an NLP
classifier}\label{mini-project-mathematical-audit-of-an-nlp-classifier}

Choose a small text classification task, such as sentiment
classification, topic classification, helpfulness prediction, or
moderation. Your project should include:

\begin{enumerate}
\def\labelenumi{\arabic{enumi}.}
\tightlist
\item
  a precise definition of \(X\), \(Y\), and any auditing variable \(A\);
\item
  a baseline model and decision rule;
\item
  group-wise confusion matrices;
\item
  at least two fairness metrics;
\item
  calibration analysis;
\item
  a distribution-shift experiment;
\item
  an uncertainty discussion;
\item
  a written statement about whether the model should be used and under
  what constraints.
\end{enumerate}

The goal is not to produce the highest-performing model. The goal is to
practice mathematically responsible evaluation.

\section{25.20 Python lab connection}\label{python-lab-connection-13}

The associated independent-study lab is planned as:

\begin{Shaded}
\begin{Highlighting}[]
\NormalTok{labs/chapter{-}25{-}ethics{-}bias{-}responsibility.ipynb}
\end{Highlighting}
\end{Shaded}

The lab should contain:

\begin{itemize}
\tightlist
\item
  a self-contained explanation of group fairness metrics;
\item
  Python implementations from scratch;
\item
  threshold scanning and visualization;
\item
  calibration and ECE;
\item
  distribution-shift weighting;
\item
  embedding association and projection;
\item
  selective prediction and abstention;
\item
  a final audit report template.
\end{itemize}

\section{25.21 Chapter summary}\label{chapter-summary-4}

Ethics in NLP is not only a philosophical add-on after model training.
It begins when we choose the data distribution, target variable,
representation, model class, objective function, metric, and deployment
rule. Mathematics helps because it forces these choices to become
explicit.

The core lessons are:

\begin{itemize}
\tightlist
\item
  responsible NLP starts with random variables and distributions;
\item
  bias can enter through data, measurement, representation,
  optimization, and deployment;
\item
  fairness metrics are conditional probabilities, not universal moral
  truths;
\item
  different fairness metrics can conflict when base rates differ;
\item
  calibration, uncertainty, and distribution shift are central to
  deployment;
\item
  privacy and factuality require separate mathematical treatment;
\item
  AI assistance can support auditing, but evidence must come from
  definitions, data, and transparent evaluation.
\end{itemize}

\section{25.22 Notes and references}\label{notes-and-references-14}

This chapter connects statistical evaluation, fairness in machine
learning, privacy, and representation auditing. Useful background
includes the fairness text of Barocas, Hardt, and Narayanan Barocas et
al. (2019a), the equalized-odds framework of Hardt, Price, and Srebro
Hardt et al. (2016a), differential privacy foundations Dwork and Roth
(2014), embedding bias studies Bolukbasi, Chang, Zou, Saligrama, and
Adam Tauman Kalai (2016), Caliskan et al. (2017), and broader
discussions of dataset and language-model risk Buolamwini and Gebru
(2018), Bender et al. (2021a).

\bookmarksetup{startatroot}

\chapter{Computer Labs}\label{computer-labs-1}

Google Colab Notebooks for Independent Study

\hfill\break

The computer labs for \emph{Mathematical Foundations of Natural Language
Processing} are designed for independent study. Each lab begins with the
necessary mathematical and conceptual background before programming,
then develops the ideas through small transparent examples in Python.

The labs are stored in the \texttt{labs/} folder of the GitHub
repository:

\url{https://github.com/wanghemath/Book-MathNLP}

Each notebook can be opened directly in Google Colab. Students do not
need to install Python locally, although the notebooks can also be
downloaded and run in JupyterLab.

\section{How to use the labs}\label{how-to-use-the-labs}

A good workflow for each chapter is:

\begin{enumerate}
\def\labelenumi{\arabic{enumi}.}
\tightlist
\item
  read the corresponding textbook chapter;
\item
  open the lab notebook in Google Colab;
\item
  run the notebook from top to bottom once without changing code;
\item
  return to each section and explain the mathematical idea in your own
  words;
\item
  complete the exercises;
\item
  complete the mini-project or extension task;
\item
  write a short reflection on assumptions, limitations, and possible
  improvements.
\end{enumerate}

The labs emphasize understanding before automation. Many examples
intentionally use small toy corpora so that students can inspect the
data, compute intermediate quantities by hand, and see how the
mathematics becomes code.

\section{Repository and Colab
pattern}\label{repository-and-colab-pattern}

The general Colab link format is:

\begin{Shaded}
\begin{Highlighting}[]
\NormalTok{https://colab.research.google.com/github/wanghemath/Book{-}MathNLP/blob/main/labs/\textless{}notebook{-}name\textgreater{}.ipynb}
\end{Highlighting}
\end{Shaded}

For example, Lab 1 opens with:

\href{https://colab.research.google.com/github/wanghemath/Book-MathNLP/blob/main/labs/chapter-01-what-is-nlp.ipynb}{Open
Lab 1 in Colab}

\section{Lab overview by part}\label{lab-overview-by-part}

\subsection{Part I. Language as Data}\label{part-i.-language-as-data-4}

These labs introduce texts, tokens, vocabularies, and document vectors.
The central question is how raw language becomes a mathematical object.

{\def\LTcaptype{none} % do not increment counter
\begin{longtable}[]{@{}
  >{\raggedleft\arraybackslash}p{(\linewidth - 6\tabcolsep) * \real{0.3077}}
  >{\raggedright\arraybackslash}p{(\linewidth - 6\tabcolsep) * \real{0.2308}}
  >{\raggedright\arraybackslash}p{(\linewidth - 6\tabcolsep) * \real{0.2308}}
  >{\raggedright\arraybackslash}p{(\linewidth - 6\tabcolsep) * \real{0.2308}}@{}}
\toprule\noalign{}
\begin{minipage}[b]{\linewidth}\raggedleft
Lab
\end{minipage} & \begin{minipage}[b]{\linewidth}\raggedright
Chapter
\end{minipage} & \begin{minipage}[b]{\linewidth}\raggedright
Notebook
\end{minipage} & \begin{minipage}[b]{\linewidth}\raggedright
Main ideas
\end{minipage} \\
\midrule\noalign{}
\endhead
\bottomrule\noalign{}
\endlastfoot
1 & What Is Natural Language Processing? &
\href{https://colab.research.google.com/github/wanghemath/Book-MathNLP/blob/main/labs/chapter-01-what-is-nlp.ipynb}{Open
in Colab} & NLP pipeline, tokens, vocabulary, vectors, document-term
matrices, first classifier intuition \\
2 & Text, Tokens, and Vocabularies &
\href{https://colab.research.google.com/github/wanghemath/Book-MathNLP/blob/main/labs/chapter-02-text-tokens-vocabularies.ipynb}{Open
in Colab} & tokenization, normalization, unknown tokens, vocabulary
size, entropy, simple BPE \\
3 & Documents as Vectors &
\href{https://colab.research.google.com/github/wanghemath/Book-MathNLP/blob/main/labs/chapter-03-documents-as-vectors.ipynb}{Open
in Colab} & count vectors, binary vectors, relative frequency, TF-IDF,
cosine similarity, SVD retrieval \\
\end{longtable}
}

\subsection{Part II. Probability Models for
Language}\label{part-ii.-probability-models-for-language-4}

These labs develop language modeling and classification from finite
probability distributions.

{\def\LTcaptype{none} % do not increment counter
\begin{longtable}[]{@{}
  >{\raggedleft\arraybackslash}p{(\linewidth - 6\tabcolsep) * \real{0.3077}}
  >{\raggedright\arraybackslash}p{(\linewidth - 6\tabcolsep) * \real{0.2308}}
  >{\raggedright\arraybackslash}p{(\linewidth - 6\tabcolsep) * \real{0.2308}}
  >{\raggedright\arraybackslash}p{(\linewidth - 6\tabcolsep) * \real{0.2308}}@{}}
\toprule\noalign{}
\begin{minipage}[b]{\linewidth}\raggedleft
Lab
\end{minipage} & \begin{minipage}[b]{\linewidth}\raggedright
Chapter
\end{minipage} & \begin{minipage}[b]{\linewidth}\raggedright
Notebook
\end{minipage} & \begin{minipage}[b]{\linewidth}\raggedright
Main ideas
\end{minipage} \\
\midrule\noalign{}
\endhead
\bottomrule\noalign{}
\endlastfoot
4 & Probability and Language &
\href{https://colab.research.google.com/github/wanghemath/Book-MathNLP/blob/main/labs/chapter-04-probability-and-language.ipynb}{Open
in Colab} & empirical distributions, likelihood, MLE, entropy, KL
divergence, perplexity, bigram preview \\
5 & N-gram Language Models &
\href{https://colab.research.google.com/github/wanghemath/Book-MathNLP/blob/main/labs/chapter-05-ngram-language-models.ipynb}{Open
in Colab} & Markov assumption, boundary tokens, smoothing,
log-probability, generation, interpolation \\
6 & Naive Bayes Text Classification &
\href{https://colab.research.google.com/github/wanghemath/Book-MathNLP/blob/main/labs/chapter-06-naive-bayes-text-classification.ipynb}{Open
in Colab} & Bayes rule, multinomial Naive Bayes, log-scores, smoothing,
posterior probabilities, evaluation \\
\end{longtable}
}

\subsection{Part III. Linear Algebra of
Language}\label{part-iii.-linear-algebra-of-language-4}

These labs study language as geometry in high-dimensional vector spaces.

{\def\LTcaptype{none} % do not increment counter
\begin{longtable}[]{@{}
  >{\raggedleft\arraybackslash}p{(\linewidth - 6\tabcolsep) * \real{0.3077}}
  >{\raggedright\arraybackslash}p{(\linewidth - 6\tabcolsep) * \real{0.2308}}
  >{\raggedright\arraybackslash}p{(\linewidth - 6\tabcolsep) * \real{0.2308}}
  >{\raggedright\arraybackslash}p{(\linewidth - 6\tabcolsep) * \real{0.2308}}@{}}
\toprule\noalign{}
\begin{minipage}[b]{\linewidth}\raggedleft
Lab
\end{minipage} & \begin{minipage}[b]{\linewidth}\raggedright
Chapter
\end{minipage} & \begin{minipage}[b]{\linewidth}\raggedright
Notebook
\end{minipage} & \begin{minipage}[b]{\linewidth}\raggedright
Main ideas
\end{minipage} \\
\midrule\noalign{}
\endhead
\bottomrule\noalign{}
\endlastfoot
7 & Vector Space Models &
\href{https://colab.research.google.com/github/wanghemath/Book-MathNLP/blob/main/labs/chapter-07-vector-space-models.ipynb}{Open
in Colab} & document-term matrices, Gram matrices, cosine retrieval,
TF-IDF scaling, SVD, LSA \\
8 & Word Embeddings &
\href{https://colab.research.google.com/github/wanghemath/Book-MathNLP/blob/main/labs/chapter-08-word-embeddings.ipynb}{Open
in Colab} & one-hot vectors, embedding lookup, co-occurrence, PMI, PPMI,
SVD embeddings, negative sampling \\
9 & Geometry of Embeddings &
\href{https://colab.research.google.com/github/wanghemath/Book-MathNLP/blob/main/labs/chapter-09-geometry-of-embeddings.ipynb}{Open
in Colab} & metrics, cosine similarity, PCA, anisotropy, hubness,
clustering, semantic directions \\
\end{longtable}
}

\subsection{Part IV. Neural Networks for
Sequences}\label{part-iv.-neural-networks-for-sequences-4}

These labs introduce neural models for text and show why order matters.

{\def\LTcaptype{none} % do not increment counter
\begin{longtable}[]{@{}
  >{\raggedleft\arraybackslash}p{(\linewidth - 6\tabcolsep) * \real{0.3077}}
  >{\raggedright\arraybackslash}p{(\linewidth - 6\tabcolsep) * \real{0.2308}}
  >{\raggedright\arraybackslash}p{(\linewidth - 6\tabcolsep) * \real{0.2308}}
  >{\raggedright\arraybackslash}p{(\linewidth - 6\tabcolsep) * \real{0.2308}}@{}}
\toprule\noalign{}
\begin{minipage}[b]{\linewidth}\raggedleft
Lab
\end{minipage} & \begin{minipage}[b]{\linewidth}\raggedright
Chapter
\end{minipage} & \begin{minipage}[b]{\linewidth}\raggedright
Notebook
\end{minipage} & \begin{minipage}[b]{\linewidth}\raggedright
Main ideas
\end{minipage} \\
\midrule\noalign{}
\endhead
\bottomrule\noalign{}
\endlastfoot
10 & Neural Networks for Text &
\href{https://colab.research.google.com/github/wanghemath/Book-MathNLP/blob/main/labs/chapter-10-neural-networks-for-text.ipynb}{Open
in Colab} & embeddings, pooling, softmax, cross-entropy, gradients,
backpropagation through embeddings \\
11 & Recurrent Neural Networks &
\href{https://colab.research.google.com/github/wanghemath/Book-MathNLP/blob/main/labs/chapter-11-recurrent-neural-networks.ipynb}{Open
in Colab} & recurrence, hidden states, many-to-one classification, BPTT,
gradient clipping, vanishing gradients \\
12 & LSTM and GRU Models &
\href{https://colab.research.google.com/github/wanghemath/Book-MathNLP/blob/main/labs/chapter-12-lstm-gru-models.ipynb}{Open
in Colab} & gated memory, LSTM equations, GRU equations, forget/update
gates, long-range dependencies \\
\end{longtable}
}

\subsection{Part V. Attention and
Transformers}\label{part-v.-attention-and-transformers-4}

These labs build the core computations behind modern transformer-based
NLP.

{\def\LTcaptype{none} % do not increment counter
\begin{longtable}[]{@{}
  >{\raggedleft\arraybackslash}p{(\linewidth - 6\tabcolsep) * \real{0.3077}}
  >{\raggedright\arraybackslash}p{(\linewidth - 6\tabcolsep) * \real{0.2308}}
  >{\raggedright\arraybackslash}p{(\linewidth - 6\tabcolsep) * \real{0.2308}}
  >{\raggedright\arraybackslash}p{(\linewidth - 6\tabcolsep) * \real{0.2308}}@{}}
\toprule\noalign{}
\begin{minipage}[b]{\linewidth}\raggedleft
Lab
\end{minipage} & \begin{minipage}[b]{\linewidth}\raggedright
Chapter
\end{minipage} & \begin{minipage}[b]{\linewidth}\raggedright
Notebook
\end{minipage} & \begin{minipage}[b]{\linewidth}\raggedright
Main ideas
\end{minipage} \\
\midrule\noalign{}
\endhead
\bottomrule\noalign{}
\endlastfoot
13 & Attention Mechanisms &
\href{https://colab.research.google.com/github/wanghemath/Book-MathNLP/blob/main/labs/chapter-13-attention-mechanisms.ipynb}{Open
in Colab} & queries, keys, values, scaled dot-product attention,
masking, attention heatmaps, entropy \\
14 & Self-Attention &
\href{https://colab.research.google.com/github/wanghemath/Book-MathNLP/blob/main/labs/chapter-14-self-attention.ipynb}{Open
in Colab} & contextual embeddings, permutation equivariance, positional
encodings, padding masks, causal masks \\
15 & Transformers &
\href{https://colab.research.google.com/github/wanghemath/Book-MathNLP/blob/main/labs/chapter-15-transformers.ipynb}{Open
in Colab} & multi-head attention, residual connections, layer
normalization, feedforward networks, encoder block \\
16 & Positional Encoding and Sequence Geometry &
\href{https://colab.research.google.com/github/wanghemath/Book-MathNLP/blob/main/labs/chapter-16-positional-encoding.ipynb}{Open
in Colab} & learned position embeddings, sinusoidal encodings, relative
bias, RoPE, causal position geometry \\
\end{longtable}
}

\subsection{Part VI. Large Language Models and
Retrieval}\label{part-vi.-large-language-models-and-retrieval-4}

These labs connect autoregressive modeling, LLM training objectives,
retrieval systems, RAG, and evaluation.

{\def\LTcaptype{none} % do not increment counter
\begin{longtable}[]{@{}
  >{\raggedleft\arraybackslash}p{(\linewidth - 6\tabcolsep) * \real{0.3077}}
  >{\raggedright\arraybackslash}p{(\linewidth - 6\tabcolsep) * \real{0.2308}}
  >{\raggedright\arraybackslash}p{(\linewidth - 6\tabcolsep) * \real{0.2308}}
  >{\raggedright\arraybackslash}p{(\linewidth - 6\tabcolsep) * \real{0.2308}}@{}}
\toprule\noalign{}
\begin{minipage}[b]{\linewidth}\raggedleft
Lab
\end{minipage} & \begin{minipage}[b]{\linewidth}\raggedright
Chapter
\end{minipage} & \begin{minipage}[b]{\linewidth}\raggedright
Notebook
\end{minipage} & \begin{minipage}[b]{\linewidth}\raggedright
Main ideas
\end{minipage} \\
\midrule\noalign{}
\endhead
\bottomrule\noalign{}
\endlastfoot
17 & Autoregressive Language Models &
\href{https://colab.research.google.com/github/wanghemath/Book-MathNLP/blob/main/labs/chapter-17-autoregressive-language-models.ipynb}{Open
in Colab} & chain rule, teacher forcing, next-token probability,
perplexity, top-\(k\) sampling, nucleus sampling \\
18 & Training Large Language Models &
\href{https://colab.research.google.com/github/wanghemath/Book-MathNLP/blob/main/labs/chapter-18-training-large-language-models.ipynb}{Open
in Colab} & pretraining, SFT masking, AdamW, validation loss, KL
regularization, preference loss, low-rank adaptation \\
19 & Embeddings, Retrieval, and RAG &
\href{https://colab.research.google.com/github/wanghemath/Book-MathNLP/blob/main/labs/chapter-19-embeddings-retrieval-rag.ipynb}{Open
in Colab} & sparse retrieval, dense retrieval, hybrid search, reranking,
chunking, retrieval metrics, toy RAG \\
20 & Evaluation &
\href{https://colab.research.google.com/github/wanghemath/Book-MathNLP/blob/main/labs/chapter-20-evaluation-of-language-models.ipynb}{Open
in Colab} & classification metrics, calibration, ECE, retrieval metrics,
RAG evaluation, bootstrap intervals, subgroup analysis \\
\end{longtable}
}

\subsection{Part VII. Advanced Mathematical
Perspectives}\label{part-vii.-advanced-mathematical-perspectives-4}

These labs connect NLP to information theory, optimization, graphs,
topology, geometry, and responsible AI.

{\def\LTcaptype{none} % do not increment counter
\begin{longtable}[]{@{}
  >{\raggedleft\arraybackslash}p{(\linewidth - 6\tabcolsep) * \real{0.3077}}
  >{\raggedright\arraybackslash}p{(\linewidth - 6\tabcolsep) * \real{0.2308}}
  >{\raggedright\arraybackslash}p{(\linewidth - 6\tabcolsep) * \real{0.2308}}
  >{\raggedright\arraybackslash}p{(\linewidth - 6\tabcolsep) * \real{0.2308}}@{}}
\toprule\noalign{}
\begin{minipage}[b]{\linewidth}\raggedleft
Lab
\end{minipage} & \begin{minipage}[b]{\linewidth}\raggedright
Chapter
\end{minipage} & \begin{minipage}[b]{\linewidth}\raggedright
Notebook
\end{minipage} & \begin{minipage}[b]{\linewidth}\raggedright
Main ideas
\end{minipage} \\
\midrule\noalign{}
\endhead
\bottomrule\noalign{}
\endlastfoot
21 & Information Theory &
\href{https://colab.research.google.com/github/wanghemath/Book-MathNLP/blob/main/labs/chapter-21-information-theory.ipynb}{Open
in Colab} & surprisal, entropy, cross-entropy, KL divergence,
perplexity, mutual information, PMI/PPMI \\
22 & Optimization for NLP &
\href{https://colab.research.google.com/github/wanghemath/Book-MathNLP/blob/main/labs/chapter-22-optimization-for-nlp.ipynb}{Open
in Colab} & stable cross-entropy, gradient checking, SGD, momentum,
Adam, AdamW, clipping, masked losses \\
23 & Graphs and Language &
\href{https://colab.research.google.com/github/wanghemath/Book-MathNLP/blob/main/labs/chapter-23-graphs-and-language.ipynb}{Open
in Colab} & adjacency matrices, Laplacians, dependency graphs,
co-occurrence graphs, PageRank, GCN, graph attention \\
24 & Topology, Geometry, and Language &
\href{https://colab.research.google.com/github/wanghemath/Book-MathNLP/blob/main/labs/chapter-24-topology-geometry-language.ipynb}{Open
in Colab} & metric spaces, neighborhood graphs, Vietoris--Rips
complexes, Betti numbers, persistence, Mapper summaries \\
25 & Ethics, Bias, and Responsibility &
\href{https://colab.research.google.com/github/wanghemath/Book-MathNLP/blob/main/labs/chapter-25-ethics-bias-responsibility.ipynb}{Open
in Colab} & fairness metrics, calibration, distribution shift, embedding
audits, abstention, privacy, RAG factuality auditing \\
\end{longtable}
}

\section{Skills developed across the
labs}\label{skills-developed-across-the-labs}

The labs are cumulative. By completing all 25 labs, students practice
the following skills.

{\def\LTcaptype{none} % do not increment counter
\begin{longtable}[]{@{}
  >{\raggedright\arraybackslash}p{(\linewidth - 2\tabcolsep) * \real{0.5000}}
  >{\raggedright\arraybackslash}p{(\linewidth - 2\tabcolsep) * \real{0.5000}}@{}}
\toprule\noalign{}
\begin{minipage}[b]{\linewidth}\raggedright
Skill area
\end{minipage} & \begin{minipage}[b]{\linewidth}\raggedright
Examples from the labs
\end{minipage} \\
\midrule\noalign{}
\endhead
\bottomrule\noalign{}
\endlastfoot
Mathematical representation & finite vocabularies, probability
distributions, vectors, matrices, graphs, metric spaces \\
Numerical computation & NumPy arrays, matrix operations, stable softmax,
SVD, gradient checking \\
Modeling & n-gram models, Naive Bayes, neural classifiers, RNNs, gated
networks, transformers \\
Retrieval & TF-IDF search, dense similarity search, hybrid retrieval,
reranking, chunking, RAG diagnostics \\
Evaluation & accuracy, precision, recall, \(F_1\), perplexity,
calibration, bootstrap uncertainty, subgroup analysis \\
Responsible AI & bias audits, threshold tradeoffs, factuality checks,
privacy, abstention, documentation \\
\end{longtable}
}

\section{Suggested grading or self-study
structure}\label{suggested-grading-or-self-study-structure}

For a course, each lab can be used in one of three ways.

\subsection{Option A. Guided practice}\label{option-a.-guided-practice}

Students run the notebook, answer the short exercises, and submit
selected outputs. This is appropriate for weekly low-stakes assignments.

\subsection{Option B. Independent lab
report}\label{option-b.-independent-lab-report}

Students submit a short report containing:

\begin{itemize}
\tightlist
\item
  a summary of the main mathematical idea;
\item
  one key equation and its interpretation;
\item
  selected code output;
\item
  answers to exercises;
\item
  one limitation or failure mode;
\item
  one proposed extension.
\end{itemize}

\subsection{Option C. Mini-project
pathway}\label{option-c.-mini-project-pathway}

Students use the mini-project section as a starting point for a longer
project. This is appropriate for final projects, research preparation,
or independent study.

\section{Recommended pacing}\label{recommended-pacing}

A semester course can use the labs at different depths.

{\def\LTcaptype{none} % do not increment counter
\begin{longtable}[]{@{}
  >{\raggedright\arraybackslash}p{(\linewidth - 2\tabcolsep) * \real{0.5000}}
  >{\raggedright\arraybackslash}p{(\linewidth - 2\tabcolsep) * \real{0.5000}}@{}}
\toprule\noalign{}
\begin{minipage}[b]{\linewidth}\raggedright
Course length
\end{minipage} & \begin{minipage}[b]{\linewidth}\raggedright
Suggested use
\end{minipage} \\
\midrule\noalign{}
\endhead
\bottomrule\noalign{}
\endlastfoot
7 weeks & Labs 1--3, 5--6, 8, 13--15, 17, 19--20, 25 \\
10 weeks & Labs 1--10, 13--15, 17--20, 21--22, 25 \\
14 weeks & Labs 1--25, with selected mini-projects replacing some
exercises \\
\end{longtable}
}

A slower mathematical course may spend more time on Labs 4, 7, 8, 13,
15, 21, and 22. A more applied NLP course may spend more time on Labs 6,
10, 17, 19, 20, and 25.

\section{Technical notes}\label{technical-notes}

The labs are intended to run in Google Colab with common scientific
Python libraries:

\begin{itemize}
\tightlist
\item
  \texttt{numpy}
\item
  \texttt{pandas}
\item
  \texttt{matplotlib}
\item
  \texttt{scikit-learn}
\end{itemize}

Some notebooks use optional libraries when available, but the core
examples are designed to remain small and transparent. The emphasis is
on understanding the mathematical structure rather than training large
models.

\section{Student expectations}\label{student-expectations}

Students should keep the following principles in mind.

\begin{tcolorbox}[enhanced jigsaw, arc=.35mm, bottomrule=.15mm, bottomtitle=1mm, breakable, colback=white, colbacktitle=quarto-callout-tip-color!10!white, colframe=quarto-callout-tip-color-frame, coltitle=black, left=2mm, leftrule=.75mm, opacityback=0, opacitybacktitle=0.6, rightrule=.15mm, title=\textcolor{quarto-callout-tip-color}{\faLightbulb}\hspace{0.5em}{Tip}, titlerule=0mm, toprule=.15mm, toptitle=1mm]

A lab is not complete when the code runs. A lab is complete when you can
explain what mathematical object each variable represents and why the
output makes sense.

\end{tcolorbox}

For every model, ask:

\begin{itemize}
\tightlist
\item
  What is the input space?
\item
  What is the output space?
\item
  What parameters are learned or chosen?
\item
  What loss, score, or criterion is being optimized?
\item
  What assumptions does the method make about language?
\item
  What failure modes should be checked?
\end{itemize}

\section{Instructor notes}\label{instructor-notes}

The labs can be used as weekly assignments, recitation materials,
independent-study notebooks, or project preparation. Because each lab
includes background before programming, students can also use them as
review materials after completing the corresponding chapter.

When assigning labs, it is useful to ask students to submit both code
output and written interpretation. The written interpretation helps
students connect Python arrays back to the mathematical ideas in the
textbook.

\bookmarksetup{startatroot}

\chapter{Interactive HTML Modules}\label{interactive-html-modules}

Interactive HTML modules can be used for visual explanations of NLP
mathematics. Good candidates include tokenization, cosine similarity,
smoothing, SVD, embedding geometry, attention matrices, positional
encoding, sampling temperature, perplexity, and persistent homology.

\section{Design principles}\label{design-principles}

\begin{itemize}
\tightlist
\item
  Each module should teach one mathematical concept clearly.
\item
  Use Plotly.js or lightweight JavaScript visualizations when
  interactivity adds value.
\item
  Keep formulas in standard LaTeX format using MathJax.
\item
  Do not make the interactive component replace the explanation; it
  should support the explanation.
\end{itemize}

\section{Planned modules}\label{planned-modules}

{\def\LTcaptype{none} % do not increment counter
\begin{longtable}[]{@{}
  >{\raggedleft\arraybackslash}p{(\linewidth - 4\tabcolsep) * \real{0.4000}}
  >{\raggedright\arraybackslash}p{(\linewidth - 4\tabcolsep) * \real{0.3000}}
  >{\raggedright\arraybackslash}p{(\linewidth - 4\tabcolsep) * \real{0.3000}}@{}}
\toprule\noalign{}
\begin{minipage}[b]{\linewidth}\raggedleft
Chapter
\end{minipage} & \begin{minipage}[b]{\linewidth}\raggedright
Topic
\end{minipage} & \begin{minipage}[b]{\linewidth}\raggedright
HTML file
\end{minipage} \\
\midrule\noalign{}
\endhead
\bottomrule\noalign{}
\endlastfoot
1 & What Is Natural Language Processing? &
\texttt{interactive-htmls/chapter-01-what-is-nlp-interactive.html} \\
2 & Text, Tokens, and Vocabularies &
\texttt{interactive-htmls/chapter-02-text-tokens-vocabularies-interactive.html} \\
3 & Documents as Vectors &
\texttt{interactive-htmls/chapter-03-documents-as-vectors-interactive.html} \\
4 & Probability and Language &
\texttt{interactive-htmls/chapter-04-probability-and-language-interactive.html} \\
5 & N-Gram Language Models &
\texttt{interactive-htmls/chapter-05-ngram-language-models-interactive.html} \\
6 & Naive Bayes and Text Classification &
\texttt{interactive-htmls/chapter-06-naive-bayes-text-classification-interactive.html} \\
7 & Vector Space Models &
\texttt{interactive-htmls/chapter-07-vector-space-models-interactive.html} \\
8 & Word Embeddings &
\texttt{interactive-htmls/chapter-08-word-embeddings-interactive.html} \\
9 & Geometry of Embedding Spaces &
\texttt{interactive-htmls/chapter-09-geometry-of-embeddings-interactive.html} \\
10 & Neural Networks for Text &
\texttt{interactive-htmls/chapter-10-neural-networks-for-text-interactive.html} \\
11 & Recurrent Neural Networks &
\texttt{interactive-htmls/chapter-11-recurrent-neural-networks-interactive.html} \\
12 & LSTM and GRU Models &
\texttt{interactive-htmls/chapter-12-lstm-gru-models-interactive.html} \\
13 & Attention Mechanisms &
\texttt{interactive-htmls/chapter-13-attention-mechanisms-interactive.html} \\
14 & Self-Attention &
\texttt{interactive-htmls/chapter-14-self-attention-interactive.html} \\
15 & The Transformer Architecture &
\texttt{interactive-htmls/chapter-15-transformers-interactive.html} \\
16 & Positional Encoding and Sequence Geometry &
\texttt{interactive-htmls/chapter-16-positional-encoding-interactive.html} \\
17 & Autoregressive Language Models &
\texttt{interactive-htmls/chapter-17-autoregressive-language-models-interactive.html} \\
18 & Training Large Language Models &
\texttt{interactive-htmls/chapter-18-training-large-language-models-interactive.html} \\
19 & Embeddings, Retrieval, and RAG &
\texttt{interactive-htmls/chapter-19-embeddings-retrieval-rag-interactive.html} \\
20 & Evaluation of Language Models &
\texttt{interactive-htmls/chapter-20-evaluation-interactive.html} \\
21 & Information Theory and Language &
\texttt{interactive-htmls/chapter-21-information-theory-interactive.html} \\
22 & Optimization for NLP &
\texttt{interactive-htmls/chapter-22-optimization-for-nlp-interactive.html} \\
23 & Graphs and Language &
\texttt{interactive-htmls/chapter-23-graphs-and-language-interactive.html} \\
24 & Topology and Geometry of Language Representations &
\texttt{interactive-htmls/chapter-24-topology-geometry-language-interactive.html} \\
25 & Ethics, Bias, and Mathematical Responsibility &
\texttt{interactive-htmls/chapter-25-ethics-bias-responsibility-interactive.html} \\
\end{longtable}
}

\bookmarksetup{startatroot}

\chapter{Interactive Diagnostics}\label{interactive-diagnostics}

\bookmarksetup{startatroot}

\chapter{Interactive diagnostics}\label{interactive-diagnostics-1}

This page helps diagnose whether the HTML interactives in the NLP book
can load Plotly.js.

\section{What should work}\label{what-should-work}

The following chapters contain embedded Plotly.js interactives:

\begin{itemize}
\tightlist
\item
  Chapter 1: classifier boundary explorer
\item
  Chapter 2: vocabulary threshold explorer
\item
  Chapter 3: TF-IDF retrieval explorer
\item
  Chapter 4: temperature/entropy explorer
\item
  Chapter 5: smoothing explorer
\item
  Chapter 6: Naive Bayes log-odds explorer
\item
  Chapter 7: rank-k latent semantic map
\item
  Chapter 8: embedding analogy explorer
\item
  Chapter 9: embedding geometry explorer
\item
  Chapter 10: softmax classifier explorer
\item
  Chapter 11: recurrent memory explorer
\end{itemize}

\section{Most common causes of a blank
interactive}\label{most-common-causes-of-a-blank-interactive}

\begin{enumerate}
\def\labelenumi{\arabic{enumi}.}
\tightlist
\item
  The local file \texttt{assets/plotly.min.js} was not copied into
  \texttt{\_book/assets/}.
\item
  A browser extension blocked JavaScript.
\item
  The page was rendered to PDF instead of HTML.
\item
  The output site is opened from a path where local scripts are blocked.
\item
  An older cached version of the page is being displayed.
\end{enumerate}

\section{Quick browser check}\label{quick-browser-check}

Open the HTML page and press \textbf{Option-Command-I} on macOS or
\textbf{Ctrl-Shift-I} on Windows/Linux. In the Console tab, run:

\begin{Shaded}
\begin{Highlighting}[]
\BuiltInTok{window}\OperatorTok{.}\AttributeTok{Plotly} \OperatorTok{\&\&}\NormalTok{ Plotly}\OperatorTok{.}\AttributeTok{version}
\end{Highlighting}
\end{Shaded}

A working page should print a Plotly version, such as \texttt{2.35.2}.

\section{Minimal Plotly test}\label{minimal-plotly-test}

\bookmarksetup{startatroot}

\chapter{Final Projects}\label{final-projects}

Final projects should combine mathematical explanation, computation, and
statistical evaluation. A strong project should state the
representation, model, loss function, dataset, metric, and limitations.

\section{Suggested projects}\label{suggested-projects-1}

\begin{enumerate}
\def\labelenumi{\arabic{enumi}.}
\tightlist
\item
  \textbf{A small mathematical search engine.} Build TF-IDF and
  embedding-based search and compare the geometry of retrieval.
\item
  \textbf{Naive Bayes versus logistic regression versus transformers.}
  Compare assumptions, losses, and empirical behavior on a text
  classification task.
\item
  \textbf{Attention from scratch.} Implement scaled dot-product
  attention and analyze how attention weights change under
  perturbations.
\item
  \textbf{Geometry of sentence embeddings.} Study clustering, PCA,
  nearest neighbors, and bias directions.
\item
  \textbf{A small retrieval-augmented generation system.} Build a RAG
  prototype over official course or program documents and evaluate
  retrieval accuracy.
\item
  \textbf{Persistent homology of embedding spaces.} Compare topological
  summaries of sentence embeddings across genres, topics, or models.
\item
  \textbf{Information theory of language models.} Compare entropy,
  cross-entropy, KL divergence, and perplexity for small language
  models.
\item
  \textbf{Bias and distribution shift in text classification.} Study how
  model performance changes across groups, topics, or time periods.
\end{enumerate}

\section{Project report structure}\label{project-report-structure}

A recommended report has the following sections:

\begin{enumerate}
\def\labelenumi{\arabic{enumi}.}
\tightlist
\item
  Research question
\item
  Mathematical formulation
\item
  Data and preprocessing
\item
  Model or algorithm
\item
  Experiments
\item
  Evaluation and uncertainty
\item
  Limitations
\item
  Reproducible code appendix
\end{enumerate}

\bookmarksetup{startatroot}

\chapter{Python Reference}\label{python-reference}

This reference collects Python patterns used throughout the labs. The
goal is to keep code readable and mathematically connected.

\section{Core libraries}\label{core-libraries}

\begin{Shaded}
\begin{Highlighting}[]
\ImportTok{import}\NormalTok{ numpy }\ImportTok{as}\NormalTok{ np}
\ImportTok{import}\NormalTok{ pandas }\ImportTok{as}\NormalTok{ pd}
\ImportTok{import}\NormalTok{ matplotlib.pyplot }\ImportTok{as}\NormalTok{ plt}
\end{Highlighting}
\end{Shaded}

\section{Count vectors}\label{count-vectors-1}

\begin{Shaded}
\begin{Highlighting}[]
\NormalTok{vocab }\OperatorTok{=}\NormalTok{ [}\StringTok{"data"}\NormalTok{, }\StringTok{"language"}\NormalTok{, }\StringTok{"model"}\NormalTok{, }\StringTok{"statistics"}\NormalTok{]}
\NormalTok{doc }\OperatorTok{=}\NormalTok{ [}\StringTok{"language"}\NormalTok{, }\StringTok{"model"}\NormalTok{, }\StringTok{"model"}\NormalTok{, }\StringTok{"data"}\NormalTok{]}
\NormalTok{counts }\OperatorTok{=}\NormalTok{ np.array([doc.count(token) }\ControlFlowTok{for}\NormalTok{ token }\KeywordTok{in}\NormalTok{ vocab])}
\NormalTok{counts}
\end{Highlighting}
\end{Shaded}

\begin{verbatim}
array([1, 1, 2, 0])
\end{verbatim}

\section{Cosine similarity}\label{cosine-similarity-2}

\begin{Shaded}
\begin{Highlighting}[]
\KeywordTok{def}\NormalTok{ cosine\_similarity(x, y):}
\NormalTok{    x }\OperatorTok{=}\NormalTok{ np.asarray(x, dtype}\OperatorTok{=}\BuiltInTok{float}\NormalTok{)}
\NormalTok{    y }\OperatorTok{=}\NormalTok{ np.asarray(y, dtype}\OperatorTok{=}\BuiltInTok{float}\NormalTok{)}
\NormalTok{    denom }\OperatorTok{=}\NormalTok{ np.linalg.norm(x) }\OperatorTok{*}\NormalTok{ np.linalg.norm(y)}
    \ControlFlowTok{if}\NormalTok{ denom }\OperatorTok{==} \DecValTok{0}\NormalTok{:}
        \ControlFlowTok{return} \FloatTok{0.0}
    \ControlFlowTok{return} \BuiltInTok{float}\NormalTok{(np.dot(x, y) }\OperatorTok{/}\NormalTok{ denom)}
\end{Highlighting}
\end{Shaded}

\section{Softmax}\label{softmax}

\begin{Shaded}
\begin{Highlighting}[]
\KeywordTok{def}\NormalTok{ softmax(z):}
\NormalTok{    z }\OperatorTok{=}\NormalTok{ np.asarray(z, dtype}\OperatorTok{=}\BuiltInTok{float}\NormalTok{)}
\NormalTok{    z }\OperatorTok{=}\NormalTok{ z }\OperatorTok{{-}}\NormalTok{ np.}\BuiltInTok{max}\NormalTok{(z)}
\NormalTok{    exp\_z }\OperatorTok{=}\NormalTok{ np.exp(z)}
    \ControlFlowTok{return}\NormalTok{ exp\_z }\OperatorTok{/}\NormalTok{ exp\_z.}\BuiltInTok{sum}\NormalTok{()}
\end{Highlighting}
\end{Shaded}

\section{Scaled dot-product
attention}\label{scaled-dot-product-attention}

\begin{Shaded}
\begin{Highlighting}[]
\KeywordTok{def}\NormalTok{ attention(Q, K, V, mask}\OperatorTok{=}\VariableTok{None}\NormalTok{):}
\NormalTok{    Q }\OperatorTok{=}\NormalTok{ np.asarray(Q, dtype}\OperatorTok{=}\BuiltInTok{float}\NormalTok{)}
\NormalTok{    K }\OperatorTok{=}\NormalTok{ np.asarray(K, dtype}\OperatorTok{=}\BuiltInTok{float}\NormalTok{)}
\NormalTok{    V }\OperatorTok{=}\NormalTok{ np.asarray(V, dtype}\OperatorTok{=}\BuiltInTok{float}\NormalTok{)}
\NormalTok{    d\_k }\OperatorTok{=}\NormalTok{ Q.shape[}\OperatorTok{{-}}\DecValTok{1}\NormalTok{]}
\NormalTok{    scores }\OperatorTok{=}\NormalTok{ Q }\OperatorTok{@}\NormalTok{ K.T }\OperatorTok{/}\NormalTok{ np.sqrt(d\_k)}
    \ControlFlowTok{if}\NormalTok{ mask }\KeywordTok{is} \KeywordTok{not} \VariableTok{None}\NormalTok{:}
\NormalTok{        scores }\OperatorTok{=}\NormalTok{ scores }\OperatorTok{+}\NormalTok{ mask}
\NormalTok{    weights }\OperatorTok{=}\NormalTok{ np.apply\_along\_axis(softmax, }\DecValTok{1}\NormalTok{, scores)}
    \ControlFlowTok{return}\NormalTok{ weights }\OperatorTok{@}\NormalTok{ V, weights}
\end{Highlighting}
\end{Shaded}

\bookmarksetup{startatroot}

\chapter{References and Further
Reading}\label{references-and-further-reading}

This page collects the main references used throughout
\emph{Mathematical Foundations of Natural Language Processing}. The list
is intentionally broad: it includes classical NLP, probability and
information theory, linear algebra and retrieval, neural language
models, attention and transformers, retrieval-augmented generation,
evaluation, optimization, graph methods, topology, and responsible NLP.

\section{Core NLP references}\label{core-nlp-references}

For a broad introduction to natural language processing, the most useful
starting points are the classical statistical treatment of Manning and
Schuetze, the modern textbook by Jurafsky and Martin, the mathematically
clear introduction by Eisenstein, and the neural NLP text by Goldberg
(Manning and Schuetze 1999; Jurafsky and Martin 2024; Eisenstein 2019;
Goldberg 2017).

\section{Probability and information
theory}\label{probability-and-information-theory}

The probability, likelihood, entropy, cross-entropy, KL divergence, and
perplexity ideas used throughout the book are rooted in Shannon's
communication theory and standard information-theory treatments (Shannon
1948; Cover and Thomas 2006; MacKay 2003). For probabilistic modeling
and statistical learning background, see Bishop, Murphy, and
Hastie--Tibshirani--Friedman (Bishop 2006b; Murphy 2022; Hastie et al.
2009).

\section{Vector spaces, retrieval, and document
geometry}\label{vector-spaces-retrieval-and-document-geometry}

The document-vector view of language connects vector-space indexing,
latent semantic analysis, and probabilistic retrieval (Salton et al.
1975; Deerwester et al. 1990; Landauer and Dumais 1997; Robertson and
Zaragoza 2009; Manning et al. 2008). These sources support the chapters
on document-term matrices, TF-IDF, cosine similarity, low-rank
approximation, and retrieval.

\section{Tokenization and word
embeddings}\label{tokenization-and-word-embeddings}

Subword tokenization is central to modern NLP systems (Sennrich et al.
2016; Kudo and Richardson 2018). Word representations connect neural
prediction, matrix factorization, and distributional semantics (Bengio
et al. 2003; Mikolov, Chen, et al. 2013; Mikolov, Sutskever, et al.
2013; Pennington et al. 2014; Levy and Goldberg 2014).

\section{Sequence models, attention, and
transformers}\label{sequence-models-attention-and-transformers}

The sequence-modeling arc begins with recurrent networks, LSTM/GRU
memory mechanisms, and attention-based encoder-decoder models (Elman
1990; Hochreiter and Schmidhuber 1997; Cho et al. 2014; Bahdanau et al.
2015; Luong et al. 2015). Transformers and large language models build
on self-attention, positional structure, pretraining, scaling, and
instruction tuning (Vaswani et al. 2017; Devlin et al. 2019; Radford et
al. 2019; {Brown et al.} 2020; Kaplan et al. 2020; {Ouyang et al.}
2022b).

\section{Retrieval-augmented generation and semantic
search}\label{retrieval-augmented-generation-and-semantic-search}

The retrieval and RAG chapters use ideas from sentence embeddings, dense
retrieval, late interaction, large-scale similarity search, and
retrieval-augmented generation (Reimers and Gurevych 2019; Karpukhin et
al. 2020; Khattab and Zaharia 2020; Johnson et al. 2019; {Lewis et al.}
2020).

\section{Optimization and training}\label{optimization-and-training}

Gradient-based training, stochastic optimization, recurrent-gradient
pathologies, Adam, AdamW, and preference-based fine-tuning are supported
by the following sources (Goodfellow et al. 2016b; Boyd and Vandenberghe
2004; Pascanu et al. 2013; Kingma and Ba 2015b; Loshchilov and Hutter
2019b; Christiano et al. 2017; Rafailov et al. 2023b).

\section{Evaluation and calibration}\label{evaluation-and-calibration}

The evaluation chapter uses classical overlap metrics, probabilistic
scoring, calibration diagnostics, and modern model comparison ideas
(Papineni et al. 2002; Lin 2004; Brier 1950; Guo et al. 2017). These
references are not a substitute for task-specific human evaluation, but
they provide a mathematical starting point.

\section{Graphs and topology}\label{graphs-and-topology}

Graph-based language models, PageRank, graph neural networks, and graph
attention are connected to matrix operators on graphs (Page et al. 1999;
Kipf and Welling 2017b; Hamilton et al. 2017; Velickovic et al. 2018).
The topological viewpoint uses computational topology, persistent
homology, and Mapper-style summaries (Edelsbrunner and Harer 2010;
Carlsson 2009; Ghrist 2008; Singh et al. 2007).

\section{Ethics, bias, and
responsibility}\label{ethics-bias-and-responsibility}

The final chapter connects technical NLP modeling with fairness, bias
measurement, documentation, and responsible deployment (Barocas et al.
2019b; Dwork et al. 2012; Hardt et al. 2016b; Bolukbasi, Chang, Zou,
Saligrama, and Adam T. Kalai 2016; Caliskan et al. 2017; Blodgett et al.
2020; Bender et al. 2021b; Gebru et al. 2021; Mitchell et al. 2019).

\section{Bibliography}\label{bibliography}

\protect\phantomsection\label{refs}
\begin{CSLReferences}{1}{1}
\bibitem[\citeproctext]{ref-bahdanau_attention_2015}
Bahdanau, Dzmitry, Kyunghyun Cho, and Yoshua Bengio. 2015. {``Neural
Machine Translation by Jointly Learning to Align and Translate.''}
\emph{International Conference on Learning Representations}.

\bibitem[\citeproctext]{ref-barocas_fairness_2019}
Barocas, Solon, Moritz Hardt, and Arvind Narayanan. 2019b.
\emph{Fairness and Machine Learning: Limitations and Opportunities}.
Fairmlbook.org.

\bibitem[\citeproctext]{ref-barocas_hardt_narayanan_2019}
Barocas, Solon, Moritz Hardt, and Arvind Narayanan. 2019a.
\emph{Fairness and Machine Learning: Limitations and Opportunities}.
Fairmlbook.org.

\bibitem[\citeproctext]{ref-bender_stochastic_2021}
Bender, Emily M., Timnit Gebru, Angelina McMillan-Major, and Shmargaret
Shmitchell. 2021b. {``On the Dangers of Stochastic Parrots: Can Language
Models Be Too Big?''} \emph{Proceedings of FAccT}, 610--23.

\bibitem[\citeproctext]{ref-bender_stochastic_parrots_2021}
Bender, Emily M., Timnit Gebru, Angelina McMillan-Major, and Shmargaret
Shmitchell. 2021a. {``On the Dangers of Stochastic Parrots: Can Language
Models Be Too Big?''} \emph{Proceedings of FAccT}, 610--23.

\bibitem[\citeproctext]{ref-bengio_neural_2003}
Bengio, Yoshua, Rejean Ducharme, Pascal Vincent, and Christian Janvin.
2003. {``A Neural Probabilistic Language Model.''} \emph{Journal of
Machine Learning Research} 3: 1137--55.

\bibitem[\citeproctext]{ref-bishop_prml_2006}
Bishop, Christopher M. 2006b. \emph{Pattern Recognition and Machine
Learning}. Springer.

\bibitem[\citeproctext]{ref-bishop_2006}
Bishop, Christopher M. 2006a. \emph{Pattern Recognition and Machine
Learning}. Springer.

\bibitem[\citeproctext]{ref-blodgett_bias_2020}
Blodgett, Su Lin, Solon Barocas, Hal Daume III, and Hanna Wallach. 2020.
{``Language (Technology) Is Power: A Critical Survey of Bias in NLP.''}
\emph{Proceedings of ACL}, 5454--76.

\bibitem[\citeproctext]{ref-bolukbasi_debiasing_2016}
Bolukbasi, Tolga, Kai-Wei Chang, James Y. Zou, Venkatesh Saligrama, and
Adam T. Kalai. 2016. {``Man Is to Computer Programmer as Woman Is to
Homemaker? Debiasing Word Embeddings.''} \emph{Advances in Neural
Information Processing Systems} 29.

\bibitem[\citeproctext]{ref-bolukbasi_man_is_2016}
Bolukbasi, Tolga, Kai-Wei Chang, James Y. Zou, Venkatesh Saligrama, and
Adam Tauman Kalai. 2016. {``Man Is to Computer Programmer as Woman Is to
Homemaker? Debiasing Word Embeddings.''} \emph{Advances in Neural
Information Processing Systems} 29.

\bibitem[\citeproctext]{ref-boyd_vandenberghe_2004}
Boyd, Stephen, and Lieven Vandenberghe. 2004. \emph{Convex
Optimization}. Cambridge University Press.

\bibitem[\citeproctext]{ref-brier_1950}
Brier, Glenn W. 1950. {``Verification of Forecasts Expressed in Terms of
Probability.''} \emph{Monthly Weather Review} 78 (1): 1--3.

\bibitem[\citeproctext]{ref-brin_page_1998}
Brin, Sergey, and Lawrence Page. 1998. {``The Anatomy of a Large-Scale
Hypertextual Web Search Engine.''} \emph{Computer Networks and ISDN
Systems} 30 (1--7): 107--17.

\bibitem[\citeproctext]{ref-brown_gpt3_2020}
{Brown, Tom B., Benjamin Mann, Nick Ryder, et al.} 2020. {``Language
Models Are Few-Shot Learners.''} \emph{Advances in Neural Information
Processing Systems} 33: 1877--901.

\bibitem[\citeproctext]{ref-buolamwini_gebru_2018}
Buolamwini, Joy, and Timnit Gebru. 2018. {``Gender Shades:
Intersectional Accuracy Disparities in Commercial Gender
Classification.''} \emph{Proceedings of the Conference on Fairness,
Accountability and Transparency}, Proceedings of machine learning
research, vol. 81: 77--91.

\bibitem[\citeproctext]{ref-caliskan_semantics_2017}
Caliskan, Aylin, Joanna J. Bryson, and Arvind Narayanan. 2017.
{``Semantics Derived Automatically from Language Corpora Contain
Human-Like Biases.''} \emph{Science} 356 (6334): 183--86.

\bibitem[\citeproctext]{ref-carlsson_topology_2009}
Carlsson, Gunnar. 2009. {``Topology and Data.''} \emph{Bulletin of the
American Mathematical Society} 46 (2): 255--308.

\bibitem[\citeproctext]{ref-cho_gru_2014}
Cho, Kyunghyun, Bart van Merrienboer, Caglar Gulcehre, et al. 2014.
{``Learning Phrase Representations Using RNN Encoder-Decoder for
Statistical Machine Translation.''} \emph{Proceedings of EMNLP},
1724--34.

\bibitem[\citeproctext]{ref-christiano_preferences_2017}
Christiano, Paul F., Jan Leike, Tom B. Brown, Miljan Martic, Shane Legg,
and Dario Amodei. 2017. {``Deep Reinforcement Learning from Human
Preferences.''} \emph{Advances in Neural Information Processing Systems}
30.

\bibitem[\citeproctext]{ref-chung_spectral_graph_1997}
Chung, Fan R. K. 1997. \emph{Spectral Graph Theory}. Vol. 92. CBMS
Regional Conference Series in Mathematics. American Mathematical
Society.

\bibitem[\citeproctext]{ref-cover_thomas_2006}
Cover, Thomas M., and Joy A. Thomas. 2006. \emph{Elements of Information
Theory}. 2nd ed. Wiley.

\bibitem[\citeproctext]{ref-deerwester_lsa_1990}
Deerwester, Scott, Susan T. Dumais, George W. Furnas, Thomas K.
Landauer, and Richard Harshman. 1990. {``Indexing by Latent Semantic
Analysis.''} \emph{Journal of the American Society for Information
Science} 41 (6): 391--407.

\bibitem[\citeproctext]{ref-devlin_bert_2019}
Devlin, Jacob, Ming-Wei Chang, Kenton Lee, and Kristina Toutanova. 2019.
{``BERT: Pre-Training of Deep Bidirectional Transformers for Language
Understanding.''} \emph{Proceedings of NAACL-HLT}, 4171--86.

\bibitem[\citeproctext]{ref-dwork_fairness_2012}
Dwork, Cynthia, Moritz Hardt, Toniann Pitassi, Omer Reingold, and
Richard Zemel. 2012. {``Fairness Through Awareness.''} \emph{Proceedings
of ITCS}, 214--26.

\bibitem[\citeproctext]{ref-dwork_roth_2014}
Dwork, Cynthia, and Aaron Roth. 2014. \emph{The Algorithmic Foundations
of Differential Privacy}. Foundations and Trends in Theoretical Computer
Science. Now Publishers.

\bibitem[\citeproctext]{ref-edelsbrunner_harer_2010}
Edelsbrunner, Herbert, and John Harer. 2010. \emph{Computational
Topology: An Introduction}. American Mathematical Society.

\bibitem[\citeproctext]{ref-eisenstein_2019}
Eisenstein, Jacob. 2019. \emph{Introduction to Natural Language
Processing}. MIT Press.

\bibitem[\citeproctext]{ref-elman_1990}
Elman, Jeffrey L. 1990. {``Finding Structure in Time.''} \emph{Cognitive
Science} 14 (2): 179--211.

\bibitem[\citeproctext]{ref-gebru_datasheets_2021}
Gebru, Timnit, Jamie Morgenstern, Briana Vecchione, et al. 2021.
{``Datasheets for Datasets.''} \emph{Communications of the ACM} 64 (12):
86--92.

\bibitem[\citeproctext]{ref-ghrist_barcodes_2008}
Ghrist, Robert. 2008. {``Barcodes: The Persistent Topology of Data.''}
\emph{Bulletin of the American Mathematical Society} 45 (1): 61--75.

\bibitem[\citeproctext]{ref-goldberg_2017}
Goldberg, Yoav. 2017. \emph{Neural Network Methods for Natural Language
Processing}. Morgan; Claypool.

\bibitem[\citeproctext]{ref-goodfellow_deep_learning_2016}
Goodfellow, Ian, Yoshua Bengio, and Aaron Courville. 2016b. \emph{Deep
Learning}. MIT Press.

\bibitem[\citeproctext]{ref-goodfellow_bengio_courville_2016}
Goodfellow, Ian, Yoshua Bengio, and Aaron Courville. 2016a. \emph{Deep
Learning}. MIT Press.

\bibitem[\citeproctext]{ref-guo_calibration_2017}
Guo, Chuan, Geoff Pleiss, Yu Sun, and Kilian Q. Weinberger. 2017. {``On
Calibration of Modern Neural Networks.''} \emph{Proceedings of ICML},
1321--30.

\bibitem[\citeproctext]{ref-hamilton_graphsage_2017}
Hamilton, William L., Rex Ying, and Jure Leskovec. 2017. {``Inductive
Representation Learning on Large Graphs.''} \emph{Advances in Neural
Information Processing Systems} 30.

\bibitem[\citeproctext]{ref-hardt_equality_2016}
Hardt, Moritz, Eric Price, and Nathan Srebro. 2016b. {``Equality of
Opportunity in Supervised Learning.''} \emph{Advances in Neural
Information Processing Systems} 29.

\bibitem[\citeproctext]{ref-hardt_price_srebro_2016}
Hardt, Moritz, Eric Price, and Nathan Srebro. 2016a. {``Equality of
Opportunity in Supervised Learning.''} \emph{Advances in Neural
Information Processing Systems} 29.

\bibitem[\citeproctext]{ref-hastie_tibshirani_friedman_2009}
Hastie, Trevor, Robert Tibshirani, and Jerome Friedman. 2009. \emph{The
Elements of Statistical Learning}. 2nd ed. Springer.

\bibitem[\citeproctext]{ref-hochreiter_lstm_1997}
Hochreiter, Sepp, and Juergen Schmidhuber. 1997. {``Long Short-Term
Memory.''} \emph{Neural Computation} 9 (8): 1735--80.

\bibitem[\citeproctext]{ref-johnson_faiss_2019}
Johnson, Jeff, Matthijs Douze, and Herve Jegou. 2019. {``Billion-Scale
Similarity Search with GPUs.''} \emph{IEEE Transactions on Big Data} 7
(3): 535--47.

\bibitem[\citeproctext]{ref-jurafsky_martin_slp3}
Jurafsky, Dan, and James H. Martin. 2024. \emph{Speech and Language
Processing}. \url{https://web.stanford.edu/~jurafsky/slp3/}.

\bibitem[\citeproctext]{ref-kaplan_scaling_2020}
Kaplan, Jared, Sam McCandlish, Tom Henighan, et al. 2020. {``Scaling
Laws for Neural Language Models.''} \emph{arXiv Preprint
arXiv:2001.08361}.

\bibitem[\citeproctext]{ref-karpukhin_dpr_2020}
Karpukhin, Vladimir, Barlas Oguz, Sewon Min, et al. 2020. {``Dense
Passage Retrieval for Open-Domain Question Answering.''}
\emph{Proceedings of EMNLP}, 6769--81.

\bibitem[\citeproctext]{ref-khattab_colbert_2020}
Khattab, Omar, and Matei Zaharia. 2020. {``ColBERT: Efficient and
Effective Passage Search via Contextualized Late Interaction over
BERT.''} \emph{Proceedings of SIGIR}, 39--48.

\bibitem[\citeproctext]{ref-kingma_adam_2015}
Kingma, Diederik P., and Jimmy Ba. 2015b. {``Adam: A Method for
Stochastic Optimization.''} \emph{International Conference on Learning
Representations}.

\bibitem[\citeproctext]{ref-kingma_ba_2015}
Kingma, Diederik P., and Jimmy Ba. 2015a. {``Adam: A Method for
Stochastic Optimization.''} \emph{International Conference on Learning
Representations}.

\bibitem[\citeproctext]{ref-kipf_gcn_2017}
Kipf, Thomas N., and Max Welling. 2017b. {``Semi-Supervised
Classification with Graph Convolutional Networks.''} \emph{International
Conference on Learning Representations}.

\bibitem[\citeproctext]{ref-kipf_welling_2017}
Kipf, Thomas N., and Max Welling. 2017a. {``Semi-Supervised
Classification with Graph Convolutional Networks.''} \emph{International
Conference on Learning Representations}.

\bibitem[\citeproctext]{ref-kudo_sentencepiece_2018}
Kudo, Taku, and John Richardson. 2018. {``SentencePiece: A Simple and
Language Independent Subword Tokenizer and Detokenizer for Neural Text
Processing.''} \emph{Proceedings of EMNLP: System Demonstrations},
66--71.

\bibitem[\citeproctext]{ref-landauer_dumais_1997}
Landauer, Thomas K., and Susan T. Dumais. 1997. {``A Solution to Plato's
Problem: The Latent Semantic Analysis Theory of Acquisition, Induction,
and Representation of Knowledge.''} \emph{Psychological Review} 104 (2):
211--40.

\bibitem[\citeproctext]{ref-levy_neural_2014}
Levy, Omer, and Yoav Goldberg. 2014. {``Neural Word Embedding as
Implicit Matrix Factorization.''} \emph{Advances in Neural Information
Processing Systems} 27.

\bibitem[\citeproctext]{ref-lewis_rag_2020}
{Lewis, Patrick, Ethan Perez, Aleksandra Piktus, et al.} 2020.
{``Retrieval-Augmented Generation for Knowledge-Intensive NLP Tasks.''}
\emph{Advances in Neural Information Processing Systems} 33: 9459--74.

\bibitem[\citeproctext]{ref-lin_rouge_2004}
Lin, Chin-Yew. 2004. {``ROUGE: A Package for Automatic Evaluation of
Summaries.''} \emph{Text Summarization Branches Out}, 74--81.

\bibitem[\citeproctext]{ref-loshchilov_adamw_2019}
Loshchilov, Ilya, and Frank Hutter. 2019b. {``Decoupled Weight Decay
Regularization.''} \emph{International Conference on Learning
Representations}.

\bibitem[\citeproctext]{ref-loshchilov_hutter_2019}
Loshchilov, Ilya, and Frank Hutter. 2019a. {``Decoupled Weight Decay
Regularization.''} \emph{International Conference on Learning
Representations}.

\bibitem[\citeproctext]{ref-luong_attention_2015}
Luong, Minh-Thang, Hieu Pham, and Christopher D. Manning. 2015.
{``Effective Approaches to Attention-Based Neural Machine
Translation.''} \emph{Proceedings of EMNLP}, 1412--21.

\bibitem[\citeproctext]{ref-mackay_2003}
MacKay, David J. C. 2003. \emph{Information Theory, Inference, and
Learning Algorithms}. Cambridge University Press.

\bibitem[\citeproctext]{ref-manning_ir_2008}
Manning, Christopher D., Prabhakar Raghavan, and Hinrich Schuetze. 2008.
\emph{Introduction to Information Retrieval}. Cambridge University
Press.

\bibitem[\citeproctext]{ref-manning_schuetze_1999}
Manning, Christopher D., and Hinrich Schuetze. 1999. \emph{Foundations
of Statistical Natural Language Processing}. MIT Press.

\bibitem[\citeproctext]{ref-mikolov_word2vec_2013}
Mikolov, Tomas, Kai Chen, Greg Corrado, and Jeffrey Dean. 2013.
{``Efficient Estimation of Word Representations in Vector Space.''}
\emph{International Conference on Learning Representations Workshop}.

\bibitem[\citeproctext]{ref-mikolov_distributed_2013}
Mikolov, Tomas, Ilya Sutskever, Kai Chen, Greg Corrado, and Jeffrey
Dean. 2013. {``Distributed Representations of Words and Phrases and
Their Compositionality.''} \emph{Advances in Neural Information
Processing Systems} 26.

\bibitem[\citeproctext]{ref-mitchell_modelcards_2019}
Mitchell, Margaret, Simone Wu, Andrew Zaldivar, et al. 2019. {``Model
Cards for Model Reporting.''} \emph{Proceedings of FAccT}, 220--29.

\bibitem[\citeproctext]{ref-murphy_probabilistic_2022}
Murphy, Kevin P. 2022. \emph{Probabilistic Machine Learning: An
Introduction}. MIT Press.

\bibitem[\citeproctext]{ref-nocedal_wright_2006}
Nocedal, Jorge, and Stephen J. Wright. 2006. \emph{Numerical
Optimization}. 2nd ed. Springer.

\bibitem[\citeproctext]{ref-ouyang_instructgpt_2022}
{Ouyang, Long, Jeffrey Wu, Xu Jiang, et al.} 2022b. {``Training Language
Models to Follow Instructions with Human Feedback.''} \emph{Advances in
Neural Information Processing Systems} 35: 27730--44.

\bibitem[\citeproctext]{ref-ouyang_2022_training}
{Ouyang, Long, Jeffrey Wu, Xu Jiang, et al.} 2022a. {``Training Language
Models to Follow Instructions with Human Feedback.''} \emph{Advances in
Neural Information Processing Systems} 35: 27730--44.

\bibitem[\citeproctext]{ref-page_pagerank_1999}
Page, Lawrence, Sergey Brin, Rajeev Motwani, and Terry Winograd. 1999.
{``The PageRank Citation Ranking: Bringing Order to the Web.''}
\emph{Technical Report, Stanford InfoLab}.

\bibitem[\citeproctext]{ref-papineni_bleu_2002}
Papineni, Kishore, Salim Roukos, Todd Ward, and Wei-Jing Zhu. 2002.
{``BLEU: A Method for Automatic Evaluation of Machine Translation.''}
\emph{Proceedings of ACL}, 311--18.

\bibitem[\citeproctext]{ref-pascanu_gradient_2013}
Pascanu, Razvan, Tomas Mikolov, and Yoshua Bengio. 2013. {``On the
Difficulty of Training Recurrent Neural Networks.''} \emph{Proceedings
of ICML}, 1310--18.

\bibitem[\citeproctext]{ref-pennington_glove_2014}
Pennington, Jeffrey, Richard Socher, and Christopher D. Manning. 2014.
{``GloVe: Global Vectors for Word Representation.''} \emph{Proceedings
of EMNLP}, 1532--43.

\bibitem[\citeproctext]{ref-radford_gpt2_2019}
Radford, Alec, Jeffrey Wu, Rewon Child, David Luan, Dario Amodei, and
Ilya Sutskever. 2019. {``Language Models Are Unsupervised Multitask
Learners.''} \emph{OpenAI Technical Report}.

\bibitem[\citeproctext]{ref-rafailov_dpo_2023}
Rafailov, Rafael, Archit Sharma, Eric Mitchell, Christopher D. Manning,
Stefano Ermon, and Chelsea Finn. 2023b. {``Direct Preference
Optimization: Your Language Model Is Secretly a Reward Model.''}
\emph{arXiv Preprint arXiv:2305.18290}.

\bibitem[\citeproctext]{ref-rafailov_2023_dpo}
Rafailov, Rafael, Archit Sharma, Eric Mitchell, Christopher D. Manning,
Stefano Ermon, and Chelsea Finn. 2023a. {``Direct Preference
Optimization: Your Language Model Is Secretly a Reward Model.''}
\emph{arXiv Preprint arXiv:2305.18290}.

\bibitem[\citeproctext]{ref-reimers_sentencebert_2019}
Reimers, Nils, and Iryna Gurevych. 2019. {``Sentence-BERT: Sentence
Embeddings Using Siamese BERT-Networks.''} \emph{Proceedings of
EMNLP-IJCNLP}, 3982--92.

\bibitem[\citeproctext]{ref-robertson_zaragoza_2009}
Robertson, Stephen, and Hugo Zaragoza. 2009. {``The Probabilistic
Relevance Framework: BM25 and Beyond.''} \emph{Foundations and Trends in
Information Retrieval} 3 (4): 333--89.

\bibitem[\citeproctext]{ref-salton_wong_yang_1975}
Salton, Gerard, Anita Wong, and Chung-Shu Yang. 1975. {``A Vector Space
Model for Automatic Indexing.''} \emph{Communications of the ACM} 18
(11): 613--20.

\bibitem[\citeproctext]{ref-sennrich_bpe_2016}
Sennrich, Rico, Barry Haddow, and Alexandra Birch. 2016. {``Neural
Machine Translation of Rare Words with Subword Units.''}
\emph{Proceedings of ACL}, 1715--25.

\bibitem[\citeproctext]{ref-shannon_1948}
Shannon, Claude E. 1948. {``A Mathematical Theory of Communication.''}
\emph{Bell System Technical Journal} 27 (3): 379--423.

\bibitem[\citeproctext]{ref-singh_mapper_2007}
Singh, Gurjeet, Facundo Memoli, and Gunnar Carlsson. 2007.
{``Topological Methods for the Analysis of High Dimensional Data Sets
and 3D Object Recognition.''} \emph{Eurographics Symposium on
Point-Based Graphics}, 91--100.

\bibitem[\citeproctext]{ref-vaswani_attention_2017}
Vaswani, Ashish, Noam Shazeer, Niki Parmar, et al. 2017. {``Attention Is
All You Need.''} \emph{Advances in Neural Information Processing
Systems} 30.

\bibitem[\citeproctext]{ref-velickovic_gat_2018}
Velickovic, Petar, Guillem Cucurull, Arantxa Casanova, Adriana Romero,
Pietro Lio, and Yoshua Bengio. 2018. {``Graph Attention Networks.''}
\emph{International Conference on Learning Representations}.

\bibitem[\citeproctext]{ref-wang_kg_embedding_2017}
Wang, Quan, Zhendong Mao, Bin Wang, and Li Guo. 2017. {``Knowledge Graph
Embedding: A Survey of Approaches and Applications.''} \emph{IEEE
Transactions on Knowledge and Data Engineering} 29 (12): 2724--43.

\end{CSLReferences}




\end{document}
